%% The manual of GAOL v5: NOT Just Another Interval Arithmetic Library.
%%
%% It describes GAOL v5, by Jordan Ninin, which continues the GAOL of Frederic
%% Goualard: its installation, its use, each of its operations with examples,
%% the names of IEEE 1788-2015, the accuracy of the operations, how the
%% intervals are computed, the tests, and what a program written for GAOL 4
%% has to change. The pages of doc/ of the sources give each change with its
%% measures, and the manual takes their content up. The text of the manual of
%% GAOL 4 (manual/v4/gaol.tex) is kept where GAOL v5 behaves as GAOL 4.
%%
%% The layout and the writing follow the PhD thesis of Jordan Ninin (2010).
%% This file holds the layout and nothing else: the packages, the page layout
%% and the macros are in gaol_def.tex, the page styles, the front matter, the
%% bibliography, the appendices and the index are set here, and the text of
%% each part is in its own file of chapters/, which this file \input's:
%% one file per chapter. Every output shown in an example was obtained from
%% GAOL v5, and manual/check_examples.py checks them against the library: the
%% examples are kept as `example' environments with their lines of outputs.
%%
%% Copyright (c) 2026 ENSTA, France
%%
%% Created 2026-09-22 by Jordan NININ
%% Rewritten 2026-11-05 in the layout of the thesis
\documentclass[12pt,a4paper,twoside,openright]{book}

%% gaol_def.tex: the packages, the page layout and the macros of the manual
%% of GAOL v5 (manual/v5/gaol.tex), which %% gaol_def.tex: the packages, the page layout and the macros of the manual
%% of GAOL v5 (manual/v5/gaol.tex), which %% gaol_def.tex: the packages, the page layout and the macros of the manual
%% of GAOL v5 (manual/v5/gaol.tex), which \input{gaol_def} loads.
%%
%% The layout is the one of the PhD thesis of Jordan Ninin (Universite de
%% Toulouse, 2010): the class book in 12 points on A4 paper, printed on both
%% sides, the margins of 2.5 cm (top, bottom, outer) and 3.5 cm (inner), the
%% running heads of fancyhdr (the chapter on the even pages, the section on the
%% odd ones, a rule of 1 point under them, the page number between dashes at
%% the foot), the links of hyperref in blue, and the numbered environments
%% Definition, Example, Remark... of the thesis, numbered within the sections.
%% The macros of the reference (\code, \file, \pxref, the environments deffun,
%% defmethod, example, onscreen...) are those of the manual of GAOL 4 by
%% Frederic Goualard (manual/v4/manual.cls), kept under their names so that the
%% text of GAOL 4 the manual takes up, and manual/check_examples.py, which reads
%% the examples, go on working.
%%
%% Copyright (c) 2026 ENSTA, France
%%
%% Created 2026-11-05 by Jordan NININ

%%--------------------------------------------------------------------------
%% Packages
%%--------------------------------------------------------------------------
\usepackage[T1]{fontenc}
\usepackage[utf8]{inputenc}
%% Latin Modern, the outlines of Computer Modern in the T1 encoding: without
%% them pdfTeX embeds bitmaps
\usepackage{lmodern}
\usepackage{textcomp}
\usepackage[english]{babel}
\usepackage{amssymb,amsmath}
\usepackage{graphicx}
\usepackage{xcolor}
\usepackage{colortbl}
\usepackage{multirow}
\usepackage{array}
\usepackage{longtable}
\usepackage{multicol}
\usepackage{fancyvrb}
\usepackage{pifont}
\usepackage[labelsep=endash,labelfont=sc]{caption}
%% The index: \index[idx]{...}, printed by \printindex[idx]
\usepackage{index}
\newindex{idx}{idx}{ind}{Index}

%%--------------------------------------------------------------------------
%% The page
%%--------------------------------------------------------------------------
\raggedbottom
\usepackage[a4paper,top=2.5cm,bottom=2.5cm,left=3.5cm,right=2.5cm,includefoot,includehead,headheight=15pt]{geometry}

\usepackage{fancyhdr}
\pagestyle{fancy}
\fancyhead{}
\fancyfoot{}
\renewcommand{\chaptermark}[1]{\markboth{\chaptername\ \thechapter\ - \ #1}{}}
\renewcommand{\sectionmark}[1]{\markright{\thesection\ - \ #1}}
\renewcommand{\headrulewidth}{1pt}
\fancyhead[LE]{\bfseries\nouppercase{\leftmark}}
\fancyhead[RO]{\bfseries\nouppercase{\rightmark}}
\fancyfoot[CE,CO]{\bfseries--\ \thepage\ --}

%% The first page of a chapter: no running head, the page number at the foot
\fancypagestyle{plain}{%
  \fancyhead{}%
  \fancyfoot[CE,CO]{\bfseries--\ \thepage\ --}%
}
%% A page with nothing
\fancypagestyle{vide}{%
  \fancyhead{}%
  \fancyfoot{}%
  \renewcommand{\headrulewidth}{0pt}%
}
%% The page number alone, under the rule
\fancypagestyle{rien}{%
  \fancyhead{}%
  \fancyfoot[CE,CO]{\bfseries--\ \thepage\ --}%
  \renewcommand{\headrulewidth}{1pt}%
}
%% The running heads of the chapters
\fancypagestyle{YO}{%
  \fancyhead[LE]{\bfseries\nouppercase{\leftmark}}%
  \fancyhead[RO]{\bfseries\nouppercase{\rightmark}}%
  \fancyfoot[CE,CO]{\bfseries--\ \thepage\ --}%
  \renewcommand{\headrulewidth}{1pt}%
}
%% The unnumbered chapters: the same text in both running heads
\newcommand*{\fixedheads}[1]{%
  \fancyhead[LE]{\bfseries\nouppercase{#1}}%
  \fancyhead[RO]{\bfseries\nouppercase{#1}}%
  \fancyfoot[CE,CO]{\bfseries--\ \thepage\ --}%
  \renewcommand{\headrulewidth}{1pt}%
}
\fancypagestyle{INTRO}{\fixedheads{Introduction}}
\fancypagestyle{CONCLU}{\fixedheads{Conclusion}}
\fancypagestyle{BIBLIO}{\fixedheads{Bibliography}}
\fancypagestyle{INDEX}{\fixedheads{Index}}

%% The blank page before a chapter carries nothing
\makeatletter
\def\cleardoublepage{\clearpage
  \if@twoside\ifodd\c@page\else
    \hbox{}\thispagestyle{vide}\newpage\if@twocolumn\hbox{}\newpage\fi\fi\fi}
\makeatother

%% An unnumbered chapter listed in the table of contents
\newcommand{\schapter}[1]{%
  \chapter*{#1}%
  \addcontentsline{toc}{chapter}{#1}%
  \markboth{#1}{#1}%
}

%%--------------------------------------------------------------------------
%% The numbered environments of the thesis: Definition 1.2.1, Example 1.2.2...
%% numbered within the sections, with one counter for all of them
%%--------------------------------------------------------------------------
\newtheorem{Def}{Definition}[section]
\newtheorem{Lem}[Def]{Lemma}
\newtheorem{Prop}[Def]{Proposition}
\newtheorem{Th}[Def]{Theorem}
\newtheorem{Cor}[Def]{Corollary}
\newtheorem{Ex}[Def]{Example}
\newtheorem{Rem}[Def]{Remark}
\def\demo{{\em Proof. }}
\def\boite{\hfill{$\hbox{\vrule height4pt width4pt}$}}

%%--------------------------------------------------------------------------
%% Mathematics
%%--------------------------------------------------------------------------
\newcommand{\RSet}{\ensuremath{\mathbb{R}}}
\newcommand{\FSet}{\ensuremath{\mathbb{F}}}
\newcommand{\ISet}{\ensuremath{\mathbb{I}}}
\newcommand{\itv}[2]{\ensuremath{[#1,\,#2]}}
\newcommand{\roundDn}[1]{\ensuremath{\downarrow#1\downarrow}}
\newcommand{\roundUp}[1]{\ensuremath{\uparrow#1\uparrow}}
\newcommand{\roundNearest}[1]{\ensuremath{\updownarrow#1\updownarrow}}
\newcommand{\hull}[1]{\ensuremath{\mathop{\Box} #1}}
\newcommand{\leftBound}[1]{\ensuremath{\underline{#1}}}
\newcommand{\rightBound}[1]{\ensuremath{\overline{#1}}}
\newcommand*{\prevdouble}[1]{\ensuremath{{#1}^-}}
\newcommand*{\nextdouble}[1]{\ensuremath{{#1}^+}}
\newcommand{\returns}{\hbox{\ding{229}}}
\let\leq=\leqslant
\let\geq=\geqslant
\let\emptyset=\varnothing

%%--------------------------------------------------------------------------
%% Names in the text
%%--------------------------------------------------------------------------
%% Code, a name of the program
\newcommand*{\code}[1]{\textup{\texttt{#1}}}
%% A name to stand out: a value of an enumeration, a format
\newcommand*{\element}[1]{\textup{\textbf{\texttt{#1}}}}
%% A definition, in the text
\newcommand{\dfn}[1]{\textsl{#1}}
%% The name of a command or a program
\newcommand{\cmd}[1]{\textup{\texttt{#1}}}
%% The name of an option, and of a long option (-{}-name)
\newcommand{\option}[1]{\textup{\texttt{#1}}}
\newcommand{\loption}[1]{\textup{\texttt{-{}-#1}}}
%% The name of a variable or a parameter
\newcommand{\var}[1]{\textsl{#1}}
%% The name of a file
\newcommand{\file}[1]{\textup{\texttt{#1}}}
%% The name of an environment variable
\newcommand{\env}[1]{\textup{\texttt{#1}}\index[idx]{#1@\texttt{#1}}}
%% A name of IEEE 1788-2015
\newcommand*{\ieee}[1]{\textsf{#1}}
%% A non terminal and a terminal of a grammar
\newcommand*{\NT}[1]{\uppercase{\textit{#1}}}
\newcommand*{\TM}[1]{\textsf{#1}}
%% What GAOL v5 changes: a mark in the text (process.md, section 4)
\newcommand*{\newinvfive}{\unskip\nobreak\ \textcolor{blue}{\textsf{\scriptsize[GAOL~v5]}}}
%% A comment in a program, and what a program prints
\definecolor{remark}{gray}{0.40}
\newcommand{\rem}[1]{\textcolor{remark}{\textsl{#1}}}
\newcommand{\outputs}[1]{\textcolor{remark}{\textsl{\returns{} #1}}}
%% A warning, in the text: the manual writes its warnings as Remarks, this
%% command stays for the sentences that need a word
\newcommand*{\important}{\textbf{Caution.}~}

%%--------------------------------------------------------------------------
%% Cross references and citations
%%--------------------------------------------------------------------------
\newcommand*{\pxref}[1]{see Section~\hyperref[#1]{\ref{#1}}, page~\pageref{#1}}
\newcommand*{\spxref}[1]{see Section~\hyperref[#1]{[\nameref{#1}]}, page~\pageref{#1}}
\newcommand*{\xref}[1]{See Section~\hyperref[#1]{\ref{#1}}, page~\pageref{#1}}
\newcommand*{\sxref}[1]{See Section~\hyperref[#1]{[\nameref{#1}]}, page~\pageref{#1}}
\newcommand*{\exref}[1]{Section~\hyperref[#1]{\ref{#1}}, page~\pageref{#1}}
\newcommand*{\sexref}[1]{Section~\hyperref[#1]{[\nameref{#1}]}, page~\pageref{#1}}
%% The citations of the manual of GAOL 4 were printed in the margin; the
%% manual of GAOL v5 cites in the text, as the thesis does
\makeatletter
\newcommand*{\margincite}{\@ifstar{\margincite@}{\margincite@}}
\newcommand*{\margincite@}[2][]{%
  \def\@tempa{#1}%
  \ifx\@tempa\@empty\cite{#2}\else\cite[#1]{#2}\fi}
\newcommand*{\marginnocite}[1]{\nocite{#1}}
\makeatother

%%--------------------------------------------------------------------------
%% The index
%%--------------------------------------------------------------------------
\newcommand*{\findex}[1]{\index[idx]{#1@\texttt{#1}}}
\newcommand*{\cindex}[1]{\index[idx]{#1}}
\newcommand*{\idxself}{\~{}}
\newcommand*{\hyperpagebf}[1]{\textbf{\hyperpage{#1}}}
%% The main entry of a name, whose page number is in bold
\newcommand*{\defindex}[1]{\index[idx]{#1@\texttt{#1}|textbf}}

%%--------------------------------------------------------------------------
%% Programs, commands and outputs (fancyvrb): @name^argument~ is the command
%% \name{argument}, and manual/check_examples.py reads the lines @outputs^...~
%% of the examples
%%--------------------------------------------------------------------------
%% What a program prints
\DefineVerbatimEnvironment{onscreen}{Verbatim}{commandchars=\@\^\~,frame=single,label={\textsf{Output}},fontsize=\footnotesize,formatcom=\upshape}
%% An example program
\def\example{\medskip\examplevb}
\def\endexample{\endexamplevb}
\DefineVerbatimEnvironment{examplevb}{Verbatim}{fontsize=\footnotesize,commandchars=\@\^\~,frame=single,label={\textsf{Example}},numbers=left,numbersep=2pt,formatcom=\upshape}
%% A command of the shell
\DefineVerbatimEnvironment{cmdshell}{Verbatim}{commandchars=\@\^\~,fontsize=\footnotesize,formatcom=\upshape}

%%--------------------------------------------------------------------------
%% The reference: the definition of a function, a method, a type, a constant
%% or a macro, its signature on a line and its description indented under it
%%--------------------------------------------------------------------------
\makeatletter
\let\@firsttime=\@undefined
\newlength{\defwidth}
\setlength{\defwidth}{.94\textwidth}
%% A signature: the type, the name in bold, the parameters
\newcommand*{\signature}[3]{%
  \noindent{\ttfamily\raggedright\frenchspacing\hangindent=2em\hangafter=1\relax #1 \textbf{#2}#3\par}}

%% A function: #1 = return type, #2 = name, #3 = parameters
\newcommand{\deffunx}[3]{%
  \defindex{#2}%
  \ifx\@firsttime\@undefined\medskip\fi
  \gdef\@firsttime{}%
  \signature{#1}{#2}{#3}}
\newenvironment{deffun}[3]{%
  \deffunx{#1}{#2}{#3}%
  \begin{flushright}%
  \begin{minipage}{\defwidth}%
}{%
  \end{minipage}%
  \end{flushright}%
  \medskip
  \global\let\@firsttime=\@undefined
}

%% A function template: #1 = template, #2 = return type, #3 = name,
%% #4 = parameters
\newcommand{\defTfunx}[4]{%
  \defindex{#3}%
  \ifx\@firsttime\@undefined\medskip\fi
  \gdef\@firsttime{}%
  \noindent{\ttfamily\raggedright\frenchspacing #1\par}
  \signature{#2}{#3}{#4}}
\newenvironment{defTfun}[4]{%
  \defTfunx{#1}{#2}{#3}{#4}%
  \begin{flushright}%
  \begin{minipage}{\defwidth}%
}{%
  \end{minipage}%
  \end{flushright}%
  \medskip
  \global\let\@firsttime=\@undefined
}

%% A method: #1 = class, #2 = return type, #3 = name, #4 = parameters
\newcommand{\defmethodx}[4]{%
  \defindex{#3}%
  \ifx\@firsttime\@undefined\medskip\fi
  \gdef\@firsttime{}%
  \noindent{\ttfamily\raggedright\frenchspacing\hangindent=2em\hangafter=1\relax\if#2\@empty\else #2 \fi#1::\textbf{#3}#4\par}}
\newenvironment{defmethod}[4]{%
  \defmethodx{#1}{#2}{#3}{#4}%
  \begin{flushright}%
  \begin{minipage}{\defwidth}%
}{%
  \end{minipage}%
  \end{flushright}%
  \medskip
  \global\let\@firsttime=\@undefined
}

%% A static method: #1 = class, #2 = return type, #3 = name, #4 = parameters
\newcommand{\defstaticmethodx}[4]{%
  \defindex{#3}%
  \ifx\@firsttime\@undefined\medskip\fi
  \gdef\@firsttime{}%
  \noindent{\ttfamily\frenchspacing #2 #1::\textbf{#3}#4\hfill\textrm{\textit{static}}}\par}
\newenvironment{defstaticmethod}[4]{%
  \defstaticmethodx{#1}{#2}{#3}{#4}%
  \begin{flushright}%
  \begin{minipage}{\defwidth}%
}{%
  \end{minipage}%
  \end{flushright}%
  \medskip
  \global\let\@firsttime=\@undefined
}

%% A type: #1 = kind (struct, class, enum), #2 = name
\newenvironment{deftype}[2]{%
  \defindex{#2}%
  \medskip
  \noindent{\ttfamily\textbf{#2}\hfill\textrm{\textit{#1}}}\par
  \begin{flushright}%
  \begin{minipage}{\defwidth}%
}{%
  \end{minipage}%
  \end{flushright}%
  \medskip
}

%% A constant: #1 = type, #2 = name
\newenvironment{defconst}[2]{%
  \defindex{#2}%
  \medskip
  \noindent{\ttfamily\frenchspacing #1 \textbf{#2}\hfill\textrm{\textit{const}}}\par
  \begin{flushright}%
  \begin{minipage}{\defwidth}%
}{%
  \end{minipage}%
  \end{flushright}%
  \medskip
}

%% A macro without parameter: #1 = name
\newenvironment{defmacro}[1]{%
  \defindex{#1}%
  \medskip
  \noindent{\ttfamily\textbf{#1}\hfill\textrm{\textit{macro}}}\par
  \begin{flushright}%
  \begin{minipage}{\defwidth}%
}{%
  \end{minipage}%
  \end{flushright}%
  \medskip
}

%% A macro with parameters: #1 = name, #2 = parameters
\newenvironment{defmacroA}[2]{%
  \defindex{#1}%
  \medskip
  \noindent{\ttfamily\textbf{#1}#2\hfill\textrm{\textit{macro}}}\par
  \begin{flushright}%
  \begin{minipage}{\defwidth}%
}{%
  \end{minipage}%
  \end{flushright}%
  \medskip
}

%% The mathematical definition of an operation, one \op per line
\newenvironment{operation}{%
  \vspace*{-5pt}%
  \def\op##1{\noindent{\small\ensuremath{\blacksquare}}\ ##1\par}%
}{%
  \par\smallskip
}
\makeatother

%% The page of the copyright notice
\newenvironment{copyrighttxt}{%
  \parindent0pt%
  \newpage
  \thispagestyle{empty}%
  \null\vfill
}{%
  \newpage
}

%% The numbers of the table of contents and of the lists are wider than in the
%% thesis (3.10.10, 12.13)
\makeatletter
\renewcommand*{\l@section}{\@dottedtocline{1}{1.5em}{2.8em}}
\renewcommand*{\l@subsection}{\@dottedtocline{2}{4.3em}{3.6em}}
\renewcommand*{\l@table}{\@dottedtocline{1}{1.5em}{3.0em}}
\renewcommand*{\l@figure}{\@dottedtocline{1}{1.5em}{3.0em}}
\makeatother

%% A column of a table whose text is set ragged right: L{width}
\newcolumntype{L}[1]{>{\raggedright\arraybackslash}p{#1}}
%% The names of functions and options, in typewriter, are not hyphenated: let
%% the lines holding them stretch rather than overflow into the margin
\setlength{\emergencystretch}{3em}
%% A box that overflows the margin by less than 5 pt, which is 2.5 cm wide, is
%% not worth a warning
\hfuzz=5pt

%%--------------------------------------------------------------------------
%% hyperref comes last
%%--------------------------------------------------------------------------
\usepackage{hyperref}
%% The names of macros and functions, in typewriter, are long and are not
%% hyphenated: a line may break after an underscore (GAOL_PRESERVE_ROUNDING),
%% without a hyphen. In the strings of the PDF the underscore stays itself.
\makeatletter
\let\gaol@underscore\_
\renewcommand*{\_}{\texorpdfstring{\gaol@underscore\hspace{0pt}}{\gaol@underscore}}
\makeatother
\hypersetup{
    unicode      = false,
    pdftoolbar   = true,
    pdfmenubar   = true,
    pdffitwindow = true,
    pdftitle     = {GAOL v5: NOT Just Another Interval Arithmetic Library},
    pdfauthor    = {Jordan Ninin, Frederic Goualard},
    pdfsubject   = {The manual of GAOL v5},
    pdfnewwindow = true,
    pdfkeywords  = {interval arithmetic, IEEE 1788-2015, CORE-MATH,
                    constraint programming, C++},
    colorlinks   = true,
    linkcolor    = blue,
    citecolor    = blue,
    filecolor    = blue,
    urlcolor     = blue,
    breaklinks   = true
}
 loads.
%%
%% The layout is the one of the PhD thesis of Jordan Ninin (Universite de
%% Toulouse, 2010): the class book in 12 points on A4 paper, printed on both
%% sides, the margins of 2.5 cm (top, bottom, outer) and 3.5 cm (inner), the
%% running heads of fancyhdr (the chapter on the even pages, the section on the
%% odd ones, a rule of 1 point under them, the page number between dashes at
%% the foot), the links of hyperref in blue, and the numbered environments
%% Definition, Example, Remark... of the thesis, numbered within the sections.
%% The macros of the reference (\code, \file, \pxref, the environments deffun,
%% defmethod, example, onscreen...) are those of the manual of GAOL 4 by
%% Frederic Goualard (manual/v4/manual.cls), kept under their names so that the
%% text of GAOL 4 the manual takes up, and manual/check_examples.py, which reads
%% the examples, go on working.
%%
%% Copyright (c) 2026 ENSTA, France
%%
%% Created 2026-11-05 by Jordan NININ

%%--------------------------------------------------------------------------
%% Packages
%%--------------------------------------------------------------------------
\usepackage[T1]{fontenc}
\usepackage[utf8]{inputenc}
%% Latin Modern, the outlines of Computer Modern in the T1 encoding: without
%% them pdfTeX embeds bitmaps
\usepackage{lmodern}
\usepackage{textcomp}
\usepackage[english]{babel}
\usepackage{amssymb,amsmath}
\usepackage{graphicx}
\usepackage{xcolor}
\usepackage{colortbl}
\usepackage{multirow}
\usepackage{array}
\usepackage{longtable}
\usepackage{multicol}
\usepackage{fancyvrb}
\usepackage{pifont}
\usepackage[labelsep=endash,labelfont=sc]{caption}
%% The index: \index[idx]{...}, printed by \printindex[idx]
\usepackage{index}
\newindex{idx}{idx}{ind}{Index}

%%--------------------------------------------------------------------------
%% The page
%%--------------------------------------------------------------------------
\raggedbottom
\usepackage[a4paper,top=2.5cm,bottom=2.5cm,left=3.5cm,right=2.5cm,includefoot,includehead,headheight=15pt]{geometry}

\usepackage{fancyhdr}
\pagestyle{fancy}
\fancyhead{}
\fancyfoot{}
\renewcommand{\chaptermark}[1]{\markboth{\chaptername\ \thechapter\ - \ #1}{}}
\renewcommand{\sectionmark}[1]{\markright{\thesection\ - \ #1}}
\renewcommand{\headrulewidth}{1pt}
\fancyhead[LE]{\bfseries\nouppercase{\leftmark}}
\fancyhead[RO]{\bfseries\nouppercase{\rightmark}}
\fancyfoot[CE,CO]{\bfseries--\ \thepage\ --}

%% The first page of a chapter: no running head, the page number at the foot
\fancypagestyle{plain}{%
  \fancyhead{}%
  \fancyfoot[CE,CO]{\bfseries--\ \thepage\ --}%
}
%% A page with nothing
\fancypagestyle{vide}{%
  \fancyhead{}%
  \fancyfoot{}%
  \renewcommand{\headrulewidth}{0pt}%
}
%% The page number alone, under the rule
\fancypagestyle{rien}{%
  \fancyhead{}%
  \fancyfoot[CE,CO]{\bfseries--\ \thepage\ --}%
  \renewcommand{\headrulewidth}{1pt}%
}
%% The running heads of the chapters
\fancypagestyle{YO}{%
  \fancyhead[LE]{\bfseries\nouppercase{\leftmark}}%
  \fancyhead[RO]{\bfseries\nouppercase{\rightmark}}%
  \fancyfoot[CE,CO]{\bfseries--\ \thepage\ --}%
  \renewcommand{\headrulewidth}{1pt}%
}
%% The unnumbered chapters: the same text in both running heads
\newcommand*{\fixedheads}[1]{%
  \fancyhead[LE]{\bfseries\nouppercase{#1}}%
  \fancyhead[RO]{\bfseries\nouppercase{#1}}%
  \fancyfoot[CE,CO]{\bfseries--\ \thepage\ --}%
  \renewcommand{\headrulewidth}{1pt}%
}
\fancypagestyle{INTRO}{\fixedheads{Introduction}}
\fancypagestyle{CONCLU}{\fixedheads{Conclusion}}
\fancypagestyle{BIBLIO}{\fixedheads{Bibliography}}
\fancypagestyle{INDEX}{\fixedheads{Index}}

%% The blank page before a chapter carries nothing
\makeatletter
\def\cleardoublepage{\clearpage
  \if@twoside\ifodd\c@page\else
    \hbox{}\thispagestyle{vide}\newpage\if@twocolumn\hbox{}\newpage\fi\fi\fi}
\makeatother

%% An unnumbered chapter listed in the table of contents
\newcommand{\schapter}[1]{%
  \chapter*{#1}%
  \addcontentsline{toc}{chapter}{#1}%
  \markboth{#1}{#1}%
}

%%--------------------------------------------------------------------------
%% The numbered environments of the thesis: Definition 1.2.1, Example 1.2.2...
%% numbered within the sections, with one counter for all of them
%%--------------------------------------------------------------------------
\newtheorem{Def}{Definition}[section]
\newtheorem{Lem}[Def]{Lemma}
\newtheorem{Prop}[Def]{Proposition}
\newtheorem{Th}[Def]{Theorem}
\newtheorem{Cor}[Def]{Corollary}
\newtheorem{Ex}[Def]{Example}
\newtheorem{Rem}[Def]{Remark}
\def\demo{{\em Proof. }}
\def\boite{\hfill{$\hbox{\vrule height4pt width4pt}$}}

%%--------------------------------------------------------------------------
%% Mathematics
%%--------------------------------------------------------------------------
\newcommand{\RSet}{\ensuremath{\mathbb{R}}}
\newcommand{\FSet}{\ensuremath{\mathbb{F}}}
\newcommand{\ISet}{\ensuremath{\mathbb{I}}}
\newcommand{\itv}[2]{\ensuremath{[#1,\,#2]}}
\newcommand{\roundDn}[1]{\ensuremath{\downarrow#1\downarrow}}
\newcommand{\roundUp}[1]{\ensuremath{\uparrow#1\uparrow}}
\newcommand{\roundNearest}[1]{\ensuremath{\updownarrow#1\updownarrow}}
\newcommand{\hull}[1]{\ensuremath{\mathop{\Box} #1}}
\newcommand{\leftBound}[1]{\ensuremath{\underline{#1}}}
\newcommand{\rightBound}[1]{\ensuremath{\overline{#1}}}
\newcommand*{\prevdouble}[1]{\ensuremath{{#1}^-}}
\newcommand*{\nextdouble}[1]{\ensuremath{{#1}^+}}
\newcommand{\returns}{\hbox{\ding{229}}}
\let\leq=\leqslant
\let\geq=\geqslant
\let\emptyset=\varnothing

%%--------------------------------------------------------------------------
%% Names in the text
%%--------------------------------------------------------------------------
%% Code, a name of the program
\newcommand*{\code}[1]{\textup{\texttt{#1}}}
%% A name to stand out: a value of an enumeration, a format
\newcommand*{\element}[1]{\textup{\textbf{\texttt{#1}}}}
%% A definition, in the text
\newcommand{\dfn}[1]{\textsl{#1}}
%% The name of a command or a program
\newcommand{\cmd}[1]{\textup{\texttt{#1}}}
%% The name of an option, and of a long option (-{}-name)
\newcommand{\option}[1]{\textup{\texttt{#1}}}
\newcommand{\loption}[1]{\textup{\texttt{-{}-#1}}}
%% The name of a variable or a parameter
\newcommand{\var}[1]{\textsl{#1}}
%% The name of a file
\newcommand{\file}[1]{\textup{\texttt{#1}}}
%% The name of an environment variable
\newcommand{\env}[1]{\textup{\texttt{#1}}\index[idx]{#1@\texttt{#1}}}
%% A name of IEEE 1788-2015
\newcommand*{\ieee}[1]{\textsf{#1}}
%% A non terminal and a terminal of a grammar
\newcommand*{\NT}[1]{\uppercase{\textit{#1}}}
\newcommand*{\TM}[1]{\textsf{#1}}
%% What GAOL v5 changes: a mark in the text (process.md, section 4)
\newcommand*{\newinvfive}{\unskip\nobreak\ \textcolor{blue}{\textsf{\scriptsize[GAOL~v5]}}}
%% A comment in a program, and what a program prints
\definecolor{remark}{gray}{0.40}
\newcommand{\rem}[1]{\textcolor{remark}{\textsl{#1}}}
\newcommand{\outputs}[1]{\textcolor{remark}{\textsl{\returns{} #1}}}
%% A warning, in the text: the manual writes its warnings as Remarks, this
%% command stays for the sentences that need a word
\newcommand*{\important}{\textbf{Caution.}~}

%%--------------------------------------------------------------------------
%% Cross references and citations
%%--------------------------------------------------------------------------
\newcommand*{\pxref}[1]{see Section~\hyperref[#1]{\ref{#1}}, page~\pageref{#1}}
\newcommand*{\spxref}[1]{see Section~\hyperref[#1]{[\nameref{#1}]}, page~\pageref{#1}}
\newcommand*{\xref}[1]{See Section~\hyperref[#1]{\ref{#1}}, page~\pageref{#1}}
\newcommand*{\sxref}[1]{See Section~\hyperref[#1]{[\nameref{#1}]}, page~\pageref{#1}}
\newcommand*{\exref}[1]{Section~\hyperref[#1]{\ref{#1}}, page~\pageref{#1}}
\newcommand*{\sexref}[1]{Section~\hyperref[#1]{[\nameref{#1}]}, page~\pageref{#1}}
%% The citations of the manual of GAOL 4 were printed in the margin; the
%% manual of GAOL v5 cites in the text, as the thesis does
\makeatletter
\newcommand*{\margincite}{\@ifstar{\margincite@}{\margincite@}}
\newcommand*{\margincite@}[2][]{%
  \def\@tempa{#1}%
  \ifx\@tempa\@empty\cite{#2}\else\cite[#1]{#2}\fi}
\newcommand*{\marginnocite}[1]{\nocite{#1}}
\makeatother

%%--------------------------------------------------------------------------
%% The index
%%--------------------------------------------------------------------------
\newcommand*{\findex}[1]{\index[idx]{#1@\texttt{#1}}}
\newcommand*{\cindex}[1]{\index[idx]{#1}}
\newcommand*{\idxself}{\~{}}
\newcommand*{\hyperpagebf}[1]{\textbf{\hyperpage{#1}}}
%% The main entry of a name, whose page number is in bold
\newcommand*{\defindex}[1]{\index[idx]{#1@\texttt{#1}|textbf}}

%%--------------------------------------------------------------------------
%% Programs, commands and outputs (fancyvrb): @name^argument~ is the command
%% \name{argument}, and manual/check_examples.py reads the lines @outputs^...~
%% of the examples
%%--------------------------------------------------------------------------
%% What a program prints
\DefineVerbatimEnvironment{onscreen}{Verbatim}{commandchars=\@\^\~,frame=single,label={\textsf{Output}},fontsize=\footnotesize,formatcom=\upshape}
%% An example program
\def\example{\medskip\examplevb}
\def\endexample{\endexamplevb}
\DefineVerbatimEnvironment{examplevb}{Verbatim}{fontsize=\footnotesize,commandchars=\@\^\~,frame=single,label={\textsf{Example}},numbers=left,numbersep=2pt,formatcom=\upshape}
%% A command of the shell
\DefineVerbatimEnvironment{cmdshell}{Verbatim}{commandchars=\@\^\~,fontsize=\footnotesize,formatcom=\upshape}

%%--------------------------------------------------------------------------
%% The reference: the definition of a function, a method, a type, a constant
%% or a macro, its signature on a line and its description indented under it
%%--------------------------------------------------------------------------
\makeatletter
\let\@firsttime=\@undefined
\newlength{\defwidth}
\setlength{\defwidth}{.94\textwidth}
%% A signature: the type, the name in bold, the parameters
\newcommand*{\signature}[3]{%
  \noindent{\ttfamily\raggedright\frenchspacing\hangindent=2em\hangafter=1\relax #1 \textbf{#2}#3\par}}

%% A function: #1 = return type, #2 = name, #3 = parameters
\newcommand{\deffunx}[3]{%
  \defindex{#2}%
  \ifx\@firsttime\@undefined\medskip\fi
  \gdef\@firsttime{}%
  \signature{#1}{#2}{#3}}
\newenvironment{deffun}[3]{%
  \deffunx{#1}{#2}{#3}%
  \begin{flushright}%
  \begin{minipage}{\defwidth}%
}{%
  \end{minipage}%
  \end{flushright}%
  \medskip
  \global\let\@firsttime=\@undefined
}

%% A function template: #1 = template, #2 = return type, #3 = name,
%% #4 = parameters
\newcommand{\defTfunx}[4]{%
  \defindex{#3}%
  \ifx\@firsttime\@undefined\medskip\fi
  \gdef\@firsttime{}%
  \noindent{\ttfamily\raggedright\frenchspacing #1\par}
  \signature{#2}{#3}{#4}}
\newenvironment{defTfun}[4]{%
  \defTfunx{#1}{#2}{#3}{#4}%
  \begin{flushright}%
  \begin{minipage}{\defwidth}%
}{%
  \end{minipage}%
  \end{flushright}%
  \medskip
  \global\let\@firsttime=\@undefined
}

%% A method: #1 = class, #2 = return type, #3 = name, #4 = parameters
\newcommand{\defmethodx}[4]{%
  \defindex{#3}%
  \ifx\@firsttime\@undefined\medskip\fi
  \gdef\@firsttime{}%
  \noindent{\ttfamily\raggedright\frenchspacing\hangindent=2em\hangafter=1\relax\if#2\@empty\else #2 \fi#1::\textbf{#3}#4\par}}
\newenvironment{defmethod}[4]{%
  \defmethodx{#1}{#2}{#3}{#4}%
  \begin{flushright}%
  \begin{minipage}{\defwidth}%
}{%
  \end{minipage}%
  \end{flushright}%
  \medskip
  \global\let\@firsttime=\@undefined
}

%% A static method: #1 = class, #2 = return type, #3 = name, #4 = parameters
\newcommand{\defstaticmethodx}[4]{%
  \defindex{#3}%
  \ifx\@firsttime\@undefined\medskip\fi
  \gdef\@firsttime{}%
  \noindent{\ttfamily\frenchspacing #2 #1::\textbf{#3}#4\hfill\textrm{\textit{static}}}\par}
\newenvironment{defstaticmethod}[4]{%
  \defstaticmethodx{#1}{#2}{#3}{#4}%
  \begin{flushright}%
  \begin{minipage}{\defwidth}%
}{%
  \end{minipage}%
  \end{flushright}%
  \medskip
  \global\let\@firsttime=\@undefined
}

%% A type: #1 = kind (struct, class, enum), #2 = name
\newenvironment{deftype}[2]{%
  \defindex{#2}%
  \medskip
  \noindent{\ttfamily\textbf{#2}\hfill\textrm{\textit{#1}}}\par
  \begin{flushright}%
  \begin{minipage}{\defwidth}%
}{%
  \end{minipage}%
  \end{flushright}%
  \medskip
}

%% A constant: #1 = type, #2 = name
\newenvironment{defconst}[2]{%
  \defindex{#2}%
  \medskip
  \noindent{\ttfamily\frenchspacing #1 \textbf{#2}\hfill\textrm{\textit{const}}}\par
  \begin{flushright}%
  \begin{minipage}{\defwidth}%
}{%
  \end{minipage}%
  \end{flushright}%
  \medskip
}

%% A macro without parameter: #1 = name
\newenvironment{defmacro}[1]{%
  \defindex{#1}%
  \medskip
  \noindent{\ttfamily\textbf{#1}\hfill\textrm{\textit{macro}}}\par
  \begin{flushright}%
  \begin{minipage}{\defwidth}%
}{%
  \end{minipage}%
  \end{flushright}%
  \medskip
}

%% A macro with parameters: #1 = name, #2 = parameters
\newenvironment{defmacroA}[2]{%
  \defindex{#1}%
  \medskip
  \noindent{\ttfamily\textbf{#1}#2\hfill\textrm{\textit{macro}}}\par
  \begin{flushright}%
  \begin{minipage}{\defwidth}%
}{%
  \end{minipage}%
  \end{flushright}%
  \medskip
}

%% The mathematical definition of an operation, one \op per line
\newenvironment{operation}{%
  \vspace*{-5pt}%
  \def\op##1{\noindent{\small\ensuremath{\blacksquare}}\ ##1\par}%
}{%
  \par\smallskip
}
\makeatother

%% The page of the copyright notice
\newenvironment{copyrighttxt}{%
  \parindent0pt%
  \newpage
  \thispagestyle{empty}%
  \null\vfill
}{%
  \newpage
}

%% The numbers of the table of contents and of the lists are wider than in the
%% thesis (3.10.10, 12.13)
\makeatletter
\renewcommand*{\l@section}{\@dottedtocline{1}{1.5em}{2.8em}}
\renewcommand*{\l@subsection}{\@dottedtocline{2}{4.3em}{3.6em}}
\renewcommand*{\l@table}{\@dottedtocline{1}{1.5em}{3.0em}}
\renewcommand*{\l@figure}{\@dottedtocline{1}{1.5em}{3.0em}}
\makeatother

%% A column of a table whose text is set ragged right: L{width}
\newcolumntype{L}[1]{>{\raggedright\arraybackslash}p{#1}}
%% The names of functions and options, in typewriter, are not hyphenated: let
%% the lines holding them stretch rather than overflow into the margin
\setlength{\emergencystretch}{3em}
%% A box that overflows the margin by less than 5 pt, which is 2.5 cm wide, is
%% not worth a warning
\hfuzz=5pt

%%--------------------------------------------------------------------------
%% hyperref comes last
%%--------------------------------------------------------------------------
\usepackage{hyperref}
%% The names of macros and functions, in typewriter, are long and are not
%% hyphenated: a line may break after an underscore (GAOL_PRESERVE_ROUNDING),
%% without a hyphen. In the strings of the PDF the underscore stays itself.
\makeatletter
\let\gaol@underscore\_
\renewcommand*{\_}{\texorpdfstring{\gaol@underscore\hspace{0pt}}{\gaol@underscore}}
\makeatother
\hypersetup{
    unicode      = false,
    pdftoolbar   = true,
    pdfmenubar   = true,
    pdffitwindow = true,
    pdftitle     = {GAOL v5: NOT Just Another Interval Arithmetic Library},
    pdfauthor    = {Jordan Ninin, Frederic Goualard},
    pdfsubject   = {The manual of GAOL v5},
    pdfnewwindow = true,
    pdfkeywords  = {interval arithmetic, IEEE 1788-2015, CORE-MATH,
                    constraint programming, C++},
    colorlinks   = true,
    linkcolor    = blue,
    citecolor    = blue,
    filecolor    = blue,
    urlcolor     = blue,
    breaklinks   = true
}
 loads.
%%
%% The layout is the one of the PhD thesis of Jordan Ninin (Universite de
%% Toulouse, 2010): the class book in 12 points on A4 paper, printed on both
%% sides, the margins of 2.5 cm (top, bottom, outer) and 3.5 cm (inner), the
%% running heads of fancyhdr (the chapter on the even pages, the section on the
%% odd ones, a rule of 1 point under them, the page number between dashes at
%% the foot), the links of hyperref in blue, and the numbered environments
%% Definition, Example, Remark... of the thesis, numbered within the sections.
%% The macros of the reference (\code, \file, \pxref, the environments deffun,
%% defmethod, example, onscreen...) are those of the manual of GAOL 4 by
%% Frederic Goualard (manual/v4/manual.cls), kept under their names so that the
%% text of GAOL 4 the manual takes up, and manual/check_examples.py, which reads
%% the examples, go on working.
%%
%% Copyright (c) 2026 ENSTA, France
%%
%% Created 2026-11-05 by Jordan NININ

%%--------------------------------------------------------------------------
%% Packages
%%--------------------------------------------------------------------------
\usepackage[T1]{fontenc}
\usepackage[utf8]{inputenc}
%% Latin Modern, the outlines of Computer Modern in the T1 encoding: without
%% them pdfTeX embeds bitmaps
\usepackage{lmodern}
\usepackage{textcomp}
\usepackage[english]{babel}
\usepackage{amssymb,amsmath}
\usepackage{graphicx}
\usepackage{xcolor}
\usepackage{colortbl}
\usepackage{multirow}
\usepackage{array}
\usepackage{longtable}
\usepackage{multicol}
\usepackage{fancyvrb}
\usepackage{pifont}
\usepackage[labelsep=endash,labelfont=sc]{caption}
%% The index: \index[idx]{...}, printed by \printindex[idx]
\usepackage{index}
\newindex{idx}{idx}{ind}{Index}

%%--------------------------------------------------------------------------
%% The page
%%--------------------------------------------------------------------------
\raggedbottom
\usepackage[a4paper,top=2.5cm,bottom=2.5cm,left=3.5cm,right=2.5cm,includefoot,includehead,headheight=15pt]{geometry}

\usepackage{fancyhdr}
\pagestyle{fancy}
\fancyhead{}
\fancyfoot{}
\renewcommand{\chaptermark}[1]{\markboth{\chaptername\ \thechapter\ - \ #1}{}}
\renewcommand{\sectionmark}[1]{\markright{\thesection\ - \ #1}}
\renewcommand{\headrulewidth}{1pt}
\fancyhead[LE]{\bfseries\nouppercase{\leftmark}}
\fancyhead[RO]{\bfseries\nouppercase{\rightmark}}
\fancyfoot[CE,CO]{\bfseries--\ \thepage\ --}

%% The first page of a chapter: no running head, the page number at the foot
\fancypagestyle{plain}{%
  \fancyhead{}%
  \fancyfoot[CE,CO]{\bfseries--\ \thepage\ --}%
}
%% A page with nothing
\fancypagestyle{vide}{%
  \fancyhead{}%
  \fancyfoot{}%
  \renewcommand{\headrulewidth}{0pt}%
}
%% The page number alone, under the rule
\fancypagestyle{rien}{%
  \fancyhead{}%
  \fancyfoot[CE,CO]{\bfseries--\ \thepage\ --}%
  \renewcommand{\headrulewidth}{1pt}%
}
%% The running heads of the chapters
\fancypagestyle{YO}{%
  \fancyhead[LE]{\bfseries\nouppercase{\leftmark}}%
  \fancyhead[RO]{\bfseries\nouppercase{\rightmark}}%
  \fancyfoot[CE,CO]{\bfseries--\ \thepage\ --}%
  \renewcommand{\headrulewidth}{1pt}%
}
%% The unnumbered chapters: the same text in both running heads
\newcommand*{\fixedheads}[1]{%
  \fancyhead[LE]{\bfseries\nouppercase{#1}}%
  \fancyhead[RO]{\bfseries\nouppercase{#1}}%
  \fancyfoot[CE,CO]{\bfseries--\ \thepage\ --}%
  \renewcommand{\headrulewidth}{1pt}%
}
\fancypagestyle{INTRO}{\fixedheads{Introduction}}
\fancypagestyle{CONCLU}{\fixedheads{Conclusion}}
\fancypagestyle{BIBLIO}{\fixedheads{Bibliography}}
\fancypagestyle{INDEX}{\fixedheads{Index}}

%% The blank page before a chapter carries nothing
\makeatletter
\def\cleardoublepage{\clearpage
  \if@twoside\ifodd\c@page\else
    \hbox{}\thispagestyle{vide}\newpage\if@twocolumn\hbox{}\newpage\fi\fi\fi}
\makeatother

%% An unnumbered chapter listed in the table of contents
\newcommand{\schapter}[1]{%
  \chapter*{#1}%
  \addcontentsline{toc}{chapter}{#1}%
  \markboth{#1}{#1}%
}

%%--------------------------------------------------------------------------
%% The numbered environments of the thesis: Definition 1.2.1, Example 1.2.2...
%% numbered within the sections, with one counter for all of them
%%--------------------------------------------------------------------------
\newtheorem{Def}{Definition}[section]
\newtheorem{Lem}[Def]{Lemma}
\newtheorem{Prop}[Def]{Proposition}
\newtheorem{Th}[Def]{Theorem}
\newtheorem{Cor}[Def]{Corollary}
\newtheorem{Ex}[Def]{Example}
\newtheorem{Rem}[Def]{Remark}
\def\demo{{\em Proof. }}
\def\boite{\hfill{$\hbox{\vrule height4pt width4pt}$}}

%%--------------------------------------------------------------------------
%% Mathematics
%%--------------------------------------------------------------------------
\newcommand{\RSet}{\ensuremath{\mathbb{R}}}
\newcommand{\FSet}{\ensuremath{\mathbb{F}}}
\newcommand{\ISet}{\ensuremath{\mathbb{I}}}
\newcommand{\itv}[2]{\ensuremath{[#1,\,#2]}}
\newcommand{\roundDn}[1]{\ensuremath{\downarrow#1\downarrow}}
\newcommand{\roundUp}[1]{\ensuremath{\uparrow#1\uparrow}}
\newcommand{\roundNearest}[1]{\ensuremath{\updownarrow#1\updownarrow}}
\newcommand{\hull}[1]{\ensuremath{\mathop{\Box} #1}}
\newcommand{\leftBound}[1]{\ensuremath{\underline{#1}}}
\newcommand{\rightBound}[1]{\ensuremath{\overline{#1}}}
\newcommand*{\prevdouble}[1]{\ensuremath{{#1}^-}}
\newcommand*{\nextdouble}[1]{\ensuremath{{#1}^+}}
\newcommand{\returns}{\hbox{\ding{229}}}
\let\leq=\leqslant
\let\geq=\geqslant
\let\emptyset=\varnothing

%%--------------------------------------------------------------------------
%% Names in the text
%%--------------------------------------------------------------------------
%% Code, a name of the program
\newcommand*{\code}[1]{\textup{\texttt{#1}}}
%% A name to stand out: a value of an enumeration, a format
\newcommand*{\element}[1]{\textup{\textbf{\texttt{#1}}}}
%% A definition, in the text
\newcommand{\dfn}[1]{\textsl{#1}}
%% The name of a command or a program
\newcommand{\cmd}[1]{\textup{\texttt{#1}}}
%% The name of an option, and of a long option (-{}-name)
\newcommand{\option}[1]{\textup{\texttt{#1}}}
\newcommand{\loption}[1]{\textup{\texttt{-{}-#1}}}
%% The name of a variable or a parameter
\newcommand{\var}[1]{\textsl{#1}}
%% The name of a file
\newcommand{\file}[1]{\textup{\texttt{#1}}}
%% The name of an environment variable
\newcommand{\env}[1]{\textup{\texttt{#1}}\index[idx]{#1@\texttt{#1}}}
%% A name of IEEE 1788-2015
\newcommand*{\ieee}[1]{\textsf{#1}}
%% A non terminal and a terminal of a grammar
\newcommand*{\NT}[1]{\uppercase{\textit{#1}}}
\newcommand*{\TM}[1]{\textsf{#1}}
%% What GAOL v5 changes: a mark in the text (process.md, section 4)
\newcommand*{\newinvfive}{\unskip\nobreak\ \textcolor{blue}{\textsf{\scriptsize[GAOL~v5]}}}
%% A comment in a program, and what a program prints
\definecolor{remark}{gray}{0.40}
\newcommand{\rem}[1]{\textcolor{remark}{\textsl{#1}}}
\newcommand{\outputs}[1]{\textcolor{remark}{\textsl{\returns{} #1}}}
%% A warning, in the text: the manual writes its warnings as Remarks, this
%% command stays for the sentences that need a word
\newcommand*{\important}{\textbf{Caution.}~}

%%--------------------------------------------------------------------------
%% Cross references and citations
%%--------------------------------------------------------------------------
\newcommand*{\pxref}[1]{see Section~\hyperref[#1]{\ref{#1}}, page~\pageref{#1}}
\newcommand*{\spxref}[1]{see Section~\hyperref[#1]{[\nameref{#1}]}, page~\pageref{#1}}
\newcommand*{\xref}[1]{See Section~\hyperref[#1]{\ref{#1}}, page~\pageref{#1}}
\newcommand*{\sxref}[1]{See Section~\hyperref[#1]{[\nameref{#1}]}, page~\pageref{#1}}
\newcommand*{\exref}[1]{Section~\hyperref[#1]{\ref{#1}}, page~\pageref{#1}}
\newcommand*{\sexref}[1]{Section~\hyperref[#1]{[\nameref{#1}]}, page~\pageref{#1}}
%% The citations of the manual of GAOL 4 were printed in the margin; the
%% manual of GAOL v5 cites in the text, as the thesis does
\makeatletter
\newcommand*{\margincite}{\@ifstar{\margincite@}{\margincite@}}
\newcommand*{\margincite@}[2][]{%
  \def\@tempa{#1}%
  \ifx\@tempa\@empty\cite{#2}\else\cite[#1]{#2}\fi}
\newcommand*{\marginnocite}[1]{\nocite{#1}}
\makeatother

%%--------------------------------------------------------------------------
%% The index
%%--------------------------------------------------------------------------
\newcommand*{\findex}[1]{\index[idx]{#1@\texttt{#1}}}
\newcommand*{\cindex}[1]{\index[idx]{#1}}
\newcommand*{\idxself}{\~{}}
\newcommand*{\hyperpagebf}[1]{\textbf{\hyperpage{#1}}}
%% The main entry of a name, whose page number is in bold
\newcommand*{\defindex}[1]{\index[idx]{#1@\texttt{#1}|textbf}}

%%--------------------------------------------------------------------------
%% Programs, commands and outputs (fancyvrb): @name^argument~ is the command
%% \name{argument}, and manual/check_examples.py reads the lines @outputs^...~
%% of the examples
%%--------------------------------------------------------------------------
%% What a program prints
\DefineVerbatimEnvironment{onscreen}{Verbatim}{commandchars=\@\^\~,frame=single,label={\textsf{Output}},fontsize=\footnotesize,formatcom=\upshape}
%% An example program
\def\example{\medskip\examplevb}
\def\endexample{\endexamplevb}
\DefineVerbatimEnvironment{examplevb}{Verbatim}{fontsize=\footnotesize,commandchars=\@\^\~,frame=single,label={\textsf{Example}},numbers=left,numbersep=2pt,formatcom=\upshape}
%% A command of the shell
\DefineVerbatimEnvironment{cmdshell}{Verbatim}{commandchars=\@\^\~,fontsize=\footnotesize,formatcom=\upshape}

%%--------------------------------------------------------------------------
%% The reference: the definition of a function, a method, a type, a constant
%% or a macro, its signature on a line and its description indented under it
%%--------------------------------------------------------------------------
\makeatletter
\let\@firsttime=\@undefined
\newlength{\defwidth}
\setlength{\defwidth}{.94\textwidth}
%% A signature: the type, the name in bold, the parameters
\newcommand*{\signature}[3]{%
  \noindent{\ttfamily\raggedright\frenchspacing\hangindent=2em\hangafter=1\relax #1 \textbf{#2}#3\par}}

%% A function: #1 = return type, #2 = name, #3 = parameters
\newcommand{\deffunx}[3]{%
  \defindex{#2}%
  \ifx\@firsttime\@undefined\medskip\fi
  \gdef\@firsttime{}%
  \signature{#1}{#2}{#3}}
\newenvironment{deffun}[3]{%
  \deffunx{#1}{#2}{#3}%
  \begin{flushright}%
  \begin{minipage}{\defwidth}%
}{%
  \end{minipage}%
  \end{flushright}%
  \medskip
  \global\let\@firsttime=\@undefined
}

%% A function template: #1 = template, #2 = return type, #3 = name,
%% #4 = parameters
\newcommand{\defTfunx}[4]{%
  \defindex{#3}%
  \ifx\@firsttime\@undefined\medskip\fi
  \gdef\@firsttime{}%
  \noindent{\ttfamily\raggedright\frenchspacing #1\par}
  \signature{#2}{#3}{#4}}
\newenvironment{defTfun}[4]{%
  \defTfunx{#1}{#2}{#3}{#4}%
  \begin{flushright}%
  \begin{minipage}{\defwidth}%
}{%
  \end{minipage}%
  \end{flushright}%
  \medskip
  \global\let\@firsttime=\@undefined
}

%% A method: #1 = class, #2 = return type, #3 = name, #4 = parameters
\newcommand{\defmethodx}[4]{%
  \defindex{#3}%
  \ifx\@firsttime\@undefined\medskip\fi
  \gdef\@firsttime{}%
  \noindent{\ttfamily\raggedright\frenchspacing\hangindent=2em\hangafter=1\relax\if#2\@empty\else #2 \fi#1::\textbf{#3}#4\par}}
\newenvironment{defmethod}[4]{%
  \defmethodx{#1}{#2}{#3}{#4}%
  \begin{flushright}%
  \begin{minipage}{\defwidth}%
}{%
  \end{minipage}%
  \end{flushright}%
  \medskip
  \global\let\@firsttime=\@undefined
}

%% A static method: #1 = class, #2 = return type, #3 = name, #4 = parameters
\newcommand{\defstaticmethodx}[4]{%
  \defindex{#3}%
  \ifx\@firsttime\@undefined\medskip\fi
  \gdef\@firsttime{}%
  \noindent{\ttfamily\frenchspacing #2 #1::\textbf{#3}#4\hfill\textrm{\textit{static}}}\par}
\newenvironment{defstaticmethod}[4]{%
  \defstaticmethodx{#1}{#2}{#3}{#4}%
  \begin{flushright}%
  \begin{minipage}{\defwidth}%
}{%
  \end{minipage}%
  \end{flushright}%
  \medskip
  \global\let\@firsttime=\@undefined
}

%% A type: #1 = kind (struct, class, enum), #2 = name
\newenvironment{deftype}[2]{%
  \defindex{#2}%
  \medskip
  \noindent{\ttfamily\textbf{#2}\hfill\textrm{\textit{#1}}}\par
  \begin{flushright}%
  \begin{minipage}{\defwidth}%
}{%
  \end{minipage}%
  \end{flushright}%
  \medskip
}

%% A constant: #1 = type, #2 = name
\newenvironment{defconst}[2]{%
  \defindex{#2}%
  \medskip
  \noindent{\ttfamily\frenchspacing #1 \textbf{#2}\hfill\textrm{\textit{const}}}\par
  \begin{flushright}%
  \begin{minipage}{\defwidth}%
}{%
  \end{minipage}%
  \end{flushright}%
  \medskip
}

%% A macro without parameter: #1 = name
\newenvironment{defmacro}[1]{%
  \defindex{#1}%
  \medskip
  \noindent{\ttfamily\textbf{#1}\hfill\textrm{\textit{macro}}}\par
  \begin{flushright}%
  \begin{minipage}{\defwidth}%
}{%
  \end{minipage}%
  \end{flushright}%
  \medskip
}

%% A macro with parameters: #1 = name, #2 = parameters
\newenvironment{defmacroA}[2]{%
  \defindex{#1}%
  \medskip
  \noindent{\ttfamily\textbf{#1}#2\hfill\textrm{\textit{macro}}}\par
  \begin{flushright}%
  \begin{minipage}{\defwidth}%
}{%
  \end{minipage}%
  \end{flushright}%
  \medskip
}

%% The mathematical definition of an operation, one \op per line
\newenvironment{operation}{%
  \vspace*{-5pt}%
  \def\op##1{\noindent{\small\ensuremath{\blacksquare}}\ ##1\par}%
}{%
  \par\smallskip
}
\makeatother

%% The page of the copyright notice
\newenvironment{copyrighttxt}{%
  \parindent0pt%
  \newpage
  \thispagestyle{empty}%
  \null\vfill
}{%
  \newpage
}

%% The numbers of the table of contents and of the lists are wider than in the
%% thesis (3.10.10, 12.13)
\makeatletter
\renewcommand*{\l@section}{\@dottedtocline{1}{1.5em}{2.8em}}
\renewcommand*{\l@subsection}{\@dottedtocline{2}{4.3em}{3.6em}}
\renewcommand*{\l@table}{\@dottedtocline{1}{1.5em}{3.0em}}
\renewcommand*{\l@figure}{\@dottedtocline{1}{1.5em}{3.0em}}
\makeatother

%% A column of a table whose text is set ragged right: L{width}
\newcolumntype{L}[1]{>{\raggedright\arraybackslash}p{#1}}
%% The names of functions and options, in typewriter, are not hyphenated: let
%% the lines holding them stretch rather than overflow into the margin
\setlength{\emergencystretch}{3em}
%% A box that overflows the margin by less than 5 pt, which is 2.5 cm wide, is
%% not worth a warning
\hfuzz=5pt

%%--------------------------------------------------------------------------
%% hyperref comes last
%%--------------------------------------------------------------------------
\usepackage{hyperref}
%% The names of macros and functions, in typewriter, are long and are not
%% hyphenated: a line may break after an underscore (GAOL_PRESERVE_ROUNDING),
%% without a hyphen. In the strings of the PDF the underscore stays itself.
\makeatletter
\let\gaol@underscore\_
\renewcommand*{\_}{\texorpdfstring{\gaol@underscore\hspace{0pt}}{\gaol@underscore}}
\makeatother
\hypersetup{
    unicode      = false,
    pdftoolbar   = true,
    pdfmenubar   = true,
    pdffitwindow = true,
    pdftitle     = {GAOL v5: NOT Just Another Interval Arithmetic Library},
    pdfauthor    = {Jordan Ninin, Frederic Goualard},
    pdfsubject   = {The manual of GAOL v5},
    pdfnewwindow = true,
    pdfkeywords  = {interval arithmetic, IEEE 1788-2015, CORE-MATH,
                    constraint programming, C++},
    colorlinks   = true,
    linkcolor    = blue,
    citecolor    = blue,
    filecolor    = blue,
    urlcolor     = blue,
    breaklinks   = true
}

\input{gaol_version}

\begin{document}

%%==========================================================================
%% Front matter
%%==========================================================================
\frontmatter
\pagestyle{plain}

%% !TEX root = ../manual_gaolv5.tex
% The title page.
\begin{titlepage}
\thispagestyle{vide}
\null\vspace{1cm}
\begin{center}
\rule{\textwidth}{1pt}\\[1.4cm]
{\fontsize{40pt}{44pt}\selectfont\bfseries GAOL v5\par}
\vspace{0.8cm}
{\LARGE\bfseries NOT Just Another\\[0.2cm] Interval Arithmetic Library\par}
\vspace{1.4cm}
\rule{\textwidth}{1pt}\\[1.6cm]
{\Large User's manual and reference\par}
\vspace{0.6cm}
{\large Edition \edition, \updated\\[0.2cm]
for version \version{} of the library\par}
\vfill
{\Large Jordan \textsc{Ninin}\par}
\vspace{0.2cm}
{\large ENSTA, France\par}
\vspace{0.9cm}
{\Large Fr\'ed\'eric \textsc{Goualard}\par}
\vspace{0.2cm}
{\large LS2N, Nantes Universit\'e, France\par}
\vspace{1.4cm}
\rule{\textwidth}{1pt}
\end{center}
\end{titlepage}

%% !TEX root = ../manual_gaolv5.tex
% The copyright notice.
\begin{copyrighttxt}
Copyright \copyright{} 2001 Swiss Federal Institute of Technology, Switzerland

Copyright \copyright{} 2002 Institut de Recherche en Informatique de Nantes (IRIN), France

Copyright \copyright{} 2002--2016 LINA UMR CNRS 6241, France

Copyright \copyright{} 2017--2020 LS2N UMR CNRS 6004, France

Copyright \copyright{} 2026 ENSTA, France

\medskip
All trademarks, copyrights and trade names are the property of their
respective owners, even if they are not specified below.

Part of the work was done while Fr\'ed\'eric Goualard was a postdoctoral fellow at
the \emph{Swiss Federal Institute of Technology, Lausanne, Switzerland},
supported by the \emph{European Research Consortium for Informatics and
Mathematics} fellowship programme.

This is edition \edition{} of the documentation of GAOL v5, which follows the
manual of GAOL~4 by Fr\'ed\'eric Goualard. It is consistent with version
\version{} of the \texttt{gaol} library.

GAOL is distributed under the GNU Library General Public License, version~2,
and CORE-MATH, whose sources GAOL compiles into its library, under the MIT
licence; both are reproduced in Appendix~\ref{chap:copying}. The files GAOL
v5 adds carry the copyright of ENSTA, France, and are distributed under the
same licence as GAOL; the files of GAOL that GAOL v5 changes keep the
copyright of their authors.

Permission is granted to make and distribute verbatim copies of this manual
provided the copyright notice and this permission notice are preserved on all
copies.

\smallskip
GAOL IS PROVIDED ``AS IS'' WITHOUT WARRANTY OF ANY KIND, EITHER EXPRESSED OR
IMPLIED, INCLUDING, BUT NOT LIMITED TO, THE IMPLIED WARRANTIES OF
MERCHAN\-TA\-BILITY AND FITNESS FOR A PARTICULAR PURPOSE. THE ENTIRE RISK AS
TO THE QUALITY AND PERFORMANCE OF THIS SOFTWARE IS WITH YOU. SHOULD THIS
SOFTWARE PROVE DEFECTIVE, YOU ASSUME THE COST OF ALL NECESSARY SERVICING,
REPAIR OR CORRECTION.
\end{copyrighttxt}

%% !TEX root = ../manual_gaolv5.tex
% The contributors and the acknowledgements (unnumbered).
\chapter*{Contributors and acknowledgements}
\thispagestyle{plain}

GAOL was designed and written by Fr\'ed\'eric Goualard
(\url{https://frederic.goualard.net}), who is its main implementor from its
first version, named \code{cell} (2001), to GAOL~4.

GAOL v5 is written by Jordan Ninin, at ENSTA, for Codac, whose intervals are
built upon GAOL. It continues GAOL from its version 4.2.2, and brings the
changes GAOL needs to compile and compute right on every system it can.

\medskip
The function \code{interval pow(const interval\&, const interval\&)} was
designed by Marc Christie. The code for the multiplication and the division
is largely inspired from the one presented by Tim Hickey, Qun Ju and Maarten
Van Emden in \emph{Interval Arithmetic: from Principles to
Implementation}~\cite{Hickey-Ju-vanEmden:2001}.

\medskip
GAOL v5 takes up:
\begin{itemize}
\item[$\bullet$] the patch IBEX applies to GAOL, by Gilles Chabert
  (initialization of the floating-point unit, rounding direction of
  \code{midpoint()} and \code{operator<{}<});
\item[$\bullet$] the fork of GAOL by Fabrice Le Bars, for Visual C++, MinGW
  and the systems the \cmd{configure} of GAOL does not know;
\item[$\bullet$] the CMake build of GAOL in IBEX, by Cyril Bouvier and Gilles
  Chabert, with the compilation flags of the IBEX fork of Fabrice Le Bars;
\item[$\bullet$] the rounding tests of Codac, which its tests follow.
\end{itemize}

\medskip
\noindent
Its elementary functions are those of CORE-MATH~\cite{CoreMath:2022}, by
Alexei Sibidanov, Paul Zimmermann and the other authors of the project
(Inria), whose sources are distributed with GAOL and compiled into its
library.

\medskip
\noindent
The layout and the writing of this manual follow the PhD thesis of Jordan
Ninin (Universit\'e de Toulouse, 2010). The text of the manual of GAOL~4 by
Fr\'ed\'eric Goualard is kept wherever GAOL v5 behaves as GAOL~4.

%% !TEX root = ../manual_gaolv5.tex
% The abstract and the keywords (unnumbered).
%%
%% Copyright (c) 2026 ENSTA, France
%%
%% By: Jordan NININ (created 2026-11-06)
%%
%% GAOL is a software distributed WITHOUT ANY WARRANTY. Read the associated
%% COPYING.LIB file for information.
\chapter*{Abstract}
\thispagestyle{plain}

GAOL is a C++ library for arithmetic on floating-point intervals. Its
intervals enclose real numbers between two doubles, and every operation
returns an interval that contains the exact result, whatever the rounding
errors of the computer: this is the containment property, on which interval
algorithms rest to prove what they compute. GAOL also provides the
\emph{relational} versions of its operations, the reverse functions of IEEE
1788-2015, which interval constraint programming needs. GAOL v5, which this
manual describes, continues GAOL as Fr\'ed\'eric Goualard developed it, and
is written so that it compiles and computes right on every system it can:
Linux, macOS and Windows, on x86, x86-64, 32-bit and 64-bit ARM, POWER
(ppc64le), s390x and RISC-V (riscv64).

\begin{itemize}
\item[$\bullet$] Every elementary function that CORE-MATH provides is bounded
  with it, which is correctly rounded in the rounding direction in effect: the
  value computed at a bound, rounding upward, is the upper bound, and the
  double below it the lower one, the tightest bounds. The integer powers, the
  roots of an order other than 2 and 3 and \code{compoundm1}, which GAOL computes itself, are the tightest
  bounds or within a few doubles of them (Chapter~\ref{chap:accuracy}). The
  bounds are the same on every architecture and with every compiler, and no
  mathematical library is to be installed or linked: CORE-MATH is compiled
  into GAOL.
\item[$\bullet$] GAOL v5 follows IEEE 1788-2015 in the operations GAOL
  provides: the empty set, the infinite bounds, the domains of the functions,
  the comparisons and the interval literals. The namespace
  \code{gaol\_ieee1788} gives each operation of the standard the name and the
  argument order the standard gives it, and the tightness of every operation
  is documented, as the standard requires.
\item[$\bullet$] The rounding direction is the business of GAOL: each
  operation that computes a bound sets the direction upward when it is not,
  whichever direction the calling code left, and GAOL can be built to restore
  the direction of the program after each operation; the operations of the
  AVX-512 path, which that build takes where the processor has the
  instructions, set no direction, their instructions carrying it.
\item[$\bullet$] Numbers are read exactly, and intervals written in decimal
  are rounded outward, whatever the C library; written in hexadecimal, an
  interval is read back as the same doubles.
\item[$\bullet$] Three builds, CMake, autotools and meson, configure GAOL the
  same way; the tests compare the bounds GAOL computes with their exact
  values, and a continuous integration builds and tests GAOL on Linux, macOS
  and Windows, with the compilers and on the processors it supports.
\end{itemize}

\medskip
\noindent
This manual is both a user's guide and a reference. It presents the notions
of interval arithmetic GAOL relies on, the installation and the use of the
library, each of its operations with examples whose outputs were produced by
GAOL v5, the names of the standard, the accuracy of the operations, the way
the intervals are computed and measured, the tests and the continuous
integration, and what a program written for GAOL~4 has to change.

\bigskip
\noindent\textbf{Keywords:} interval arithmetic, floating-point arithmetic,
IEEE 1788-2015, correctly rounded elementary functions, CORE-MATH, reverse
functions, constraint programming, C++.


\tableofcontents
\listoftables

%%==========================================================================
%% Introduction
%%==========================================================================
\mainmatter
\pagestyle{INTRO}
%% !TEX root = ../manual_gaolv5.tex
% The introduction (unnumbered).
%%
%% Copyright (c) 2001 Swiss Federal Institute of Technology, Switzerland
%% Copyright (c) 2002-2016 Laboratoire d'Informatique de
%%                         Nantes-Atlantique, France
%% Copyright (c) 2017-2020 Laboratoire des Sciences du Numérique de Nantes,
%%                         France
%% Copyright (c) 2026 ENSTA, France
%%
%% By: Frederic Goualard, Jordan NININ
%%
%% GAOL is a software distributed WITHOUT ANY WARRANTY. Read the associated
%% COPYING.LIB file for information.
\chapter*{Introduction}\label{chap:introduction}
\addcontentsline{toc}{chapter}{Introduction}
\thispagestyle{plain}

Computers compute with floating-point numbers, and a floating-point number
is rarely the real number a program means: $0.1$ has no finite binary
expansion, the quotient $1/3$ is rounded, and the rounding errors of a long
computation add up in ways that are hard to follow. Interval arithmetic
answers this difficulty by computing with pairs of floating-point numbers:
an \emph{interval} $[a, b]$ stands for every real number between $a$ and
$b$, and each operation on intervals returns an interval that contains all
the results the operation can take on the reals it encloses. A program that
computes with intervals proves, at the end of its computation, that the
exact result lies within the interval it obtained. This is the containment
property, which the first chapter defines, and which every operation of GAOL
honours.

\medskip
GAOL\footnote{For those readers who are not native English speakers,
``gaol'' should be pronounced \textsf{j\={a}l}, like the word ``jail'', of
which it is a, chiefly British, variant.}
\cindex{gaol!pronunciation}%
is a C++ library to perform arithmetic with floating-point intervals. Its
development was initiated at the \emph{Swiss Federal Institute of
Technology}, Lausanne, Switzerland, while Fr\'ed\'eric Goualard was a
postdoctoral fellow supported by the \emph{Swiss National Science
Foundation}. It started as a limited version of \emph{Jail} (now
\emph{halloween}), a templated C++ interval library developed during
Goualard's PhD. To our knowledge, a unique feature of GAOL among the C++
interval arithmetic libraries available when it was written is the
implementation of \emph{relational arithmetic operators}, which interval
constraint programming requires (Chapter~\ref{chap:arithmetic}).
\cindex{operator!relational}\cindex{relational operator}%
Hence the pun of the name: GAOL is \emph{not} JAIL (\emph{Just Another
Interval Library}). The standard for interval arithmetic, IEEE
1788-2015~\cite{IEEE1788:2015}, has since required such operators, as its
\emph{reverse functions}.
\cindex{IEEE 1788-2015}%

\medskip
GAOL v5, written by Jordan Ninin at ENSTA, continues GAOL as Fr\'ed\'eric
Goualard developed it, from its version 4.2.2. It was written for Codac and
IBEX, whose intervals are built upon GAOL, and brings the changes GAOL needs to
compile and compute right on every system it can: Linux, macOS and Windows,
on x86, x86-64, 32-bit and 64-bit ARM, POWER (ppc64le), s390x and RISC-V
(riscv64).
\cindex{GAOL v5}%
What it brings, and this manual describes, is the following.

\begin{itemize}
\item[$\bullet$] \textbf{Correctly rounded elementary functions.} Every
  elementary function that CORE-MATH~\cite{CoreMath:2022} provides is bounded
  with it, which is correctly rounded in the rounding direction in effect: the
  value GAOL computes at a bound, rounding upward, \emph{is} the upper bound,
  and the double below it the lower one, the tightest bounds. The integer
  powers, the roots of an order other than 2 and 3 and \code{compoundm1}, which GAOL computes itself, are the
  tightest bounds or within a few doubles of them
  (Chapter~\ref{chap:accuracy}). The bounds are the same on every architecture
  and with every compiler, and there is no mathematical library to install, to
  find or to link along with GAOL: CORE-MATH is compiled into it.
\item[$\bullet$] \textbf{IEEE 1788-2015.} GAOL v5 follows the standard for
  interval arithmetic in the operations GAOL provides: the empty set, the
  infinite bounds, the domains of the functions, the comparisons and the
  interval literals. The few cases where it differs, mostly the \code{pow}
  of the namespace \code{gaol}, are stated where the operations are
  described, and the tightness of each operation is documented, as the
  standard requires (Chapter~\ref{chap:accuracy}).
\item[$\bullet$] \textbf{New operations}, which the standard requires or
  recommends: \code{exp2}, \code{exp10}, \code{log2}, \code{log10},
  \code{fma}, the integer functions \code{sign}, \code{trunc},
  \code{round\_ties\_to\_even} and \code{round\_ties\_to\_away}, sixteen
  recommended functions (\code{expm1}, \code{log1p}, \code{hypot},
  \code{compoundm1}, \code{sinpi}\ldots), the cancellative subtraction and
  addition, the radius, the comparisons \code{less} and
  \code{strictly\_less}, the boolean functions \code{is\_entire} and
  \code{is\_common\_interval}, \code{overlap()}, the reverse functions of the
  power and of \code{atan2} (\ieee{powRev1}, \ieee{powRev2}, \ieee{atan2Rev1},
  \ieee{atan2Rev2}), \ieee{pownRev} for a non-positive exponent and
  \ieee{mulRevToPair}; \code{atan2}, which GAOL declared, is implemented.
\item[$\bullet$] \textbf{The names of the standard.} The namespace
  \code{gaol\_ieee1788} gives each operation of the standard that GAOL
  provides the name and the argument order the standard gives it
  (Chapter~\ref{chap:ieee1788}).
\item[$\bullet$] \textbf{The rounding direction is the business of GAOL.}
  Each operation that computes a bound sets the rounding direction upward
  when it is not, whichever direction the calling code left; the
  operations of the AVX-512 path set none, their instructions carrying the
  direction. The user no longer has to keep it upward
  (Chapter~\ref{chap:rounding}).
\item[$\bullet$] \textbf{Exact input and output.} Numbers are read exactly,
  and intervals written in decimal are rounded outward, whatever the C
  library; written in hexadecimal, an interval is read back as the same
  doubles (Chapter~\ref{chap:io}).
\item[$\bullet$] \textbf{Three builds that agree.} GAOL has a CMake build,
  which is the one to use, and the autotools and meson builds, and the three
  configure GAOL the same way. Tests compare the bounds GAOL computes with
  their exact values, and a continuous integration builds and tests GAOL on
  Linux, macOS and Windows, with the compilers and on the processors listed in
  Chapter~\ref{chap:installation}.
\end{itemize}

\noindent
A program written for GAOL~4 needs a few changes, which
Chapter~\ref{chap:changes} lists: first of all, \file{gaol/gaol.h} no longer
opens the namespace \code{gaol}.

\section*{About this manual}

This document is both a manual and a reference to use GAOL. It assumes a
prior knowledge of interval arithmetic, whose notions the first chapter
recalls in the form GAOL uses them; the books and papers by Goldberg,
Neumaier, Moore, Alefeld and Herzberger, and
others~\cite{Goldberg:01,IEEE754:2008,Moore:Book1966,Alefeld-Herzberger:83,Neumaier:Book1990}
give a full presentation of floating-point arithmetic, interval arithmetic
and the use thereof.
\cindex{references!on interval arithmetic}%
The main entry point for interval arithmetic on the Web is the \emph{Interval
Computation site} of Vladik Kreino\-vich,
\url{http://www.cs.utep.edu/interval-comp/}.
\cindex{interval!computation site (web)}

\medskip
The pages of the directory \file{doc} of the sources give the details of each
change GAOL v5 brings, with its measures, and this manual takes their content
up; Appendix~\ref{chap:documentation} lists them. The mark
\textcolor{blue}{\textsf{\scriptsize[GAOL~v5]}} follows, throughout the
text, what behaves otherwise than in GAOL~4. A clause number in parentheses,
as (12.10.3), and a table cited as Table~9.1 or Table~10.7 are those of
IEEE 1788-2015 (the tables of this manual are numbered by chapter and linked).

\medskip
The examples of the reference show what they print, and every output shown
was produced by GAOL v5. A script of the sources compiles each example that
shows its output under its code, runs it and compares what it prints with the
manual (Chapter~\ref{chap:tests}); the outputs shown apart, in a frame
Output, are not compared by it. An
example that writes a \code{bool} sets \code{boolalpha} first: without it,
\code{cout} writes 1 and 0. The examples assume \code{using namespace gaol;}
and \code{using namespace std;}, except those of
Chapter~\ref{chap:ieee1788}, which open \code{gaol\_ieee1788} instead, and
each example is a program of its own.

\medskip
Classes, methods, functions, macros, constants and variables available in
the library but not described in this document are likely to change or to be
removed. Consequently, they should be used with caution, if at all. Among
them are the intervals of floats, \code{gaol::intervalf} and
\code{gaol::interval2f}, which are unfinished and which none of the builds
compiles: only a developer of GAOL does, defining
\code{GAOL\_FLOAT\_INTERVALS} in \file{gaol/gaol\_config.h}.
\cindex{undocumented functions}%
\cindex{function!undocumented}%

\section*{Outline}

The first chapter recalls the notions this manual relies on: floating-point
numbers and their roundings, intervals and the containment property, the
dependency problem, what IEEE 1788-2015 fixes, the idea of the relational
functions, the mathematical libraries from which correctly rounded
elementary functions came, and the accuracy modes of the standard.

\medskip
Chapters~\ref{chap:installation} and~\ref{chap:overview} are the user's
guide. The second chapter describes the installation of GAOL: the platforms
it supports, the compilers and options it refuses, the three builds and their
options, and what the configuration decides by itself. The third chapter
shows how a program uses GAOL: a first program, the flags of interval
arithmetic the code including the headers of GAOL is compiled with, the two
namespaces and their powers, the tools of interval algorithms, the common
errors, and the correspondence with the names of other interval libraries.

\medskip
Chapter~\ref{chap:rounding} is devoted to the rounding direction and to the
floating-point environment, which GAOL v5 takes charge of: what the upward
rounding does to the rest of the program, how the direction can be
preserved, the AVX-512 path, the NEON intervals of 64-bit ARM processors, the
FPU intervals, the modes that flush the subnormal numbers to zero, the
floating-point exceptions, and the initialization and the cleanup of GAOL.

\medskip
Chapters~\ref{chap:creation} to~\ref{chap:io} are the reference of the class
\code{interval}: the constructors and the assignments, the relations, the
arithmetic operations, functional and relational, the numeric and integer
functions, the splitting methods and the set operations, and the reading and
writing of intervals. Chapter~\ref{chap:ieee1788} gives the names of IEEE
1788-2015, and where the standard and GAOL differ. Chapter~\ref{chap:accuracy}
states the tightness of each operation, as the standard requires.

\medskip
Chapter~\ref{chap:computing} explains how GAOL computes its intervals: the
representation of an interval in a register, the SSE2, AVX-512, NEON and FPU
kernels, where the time of a program goes, the cost of the rounding direction
on each processor, what the compilers give, how GAOL compares with the other
interval libraries, with ITF1788, the test suite of the standard, and with the
benchmark of Tang et al., and the verdicts
on new paths. Chapter~\ref{chap:errors} describes the floating-point constants
and functions, the version information, the errors and the debugging
facilities. Chapter~\ref{chap:tests} presents the tests, the examples, the
continuous integration and the scripts of the repository.

\medskip
The last chapter, Chapter~\ref{chap:changes}, is written for the programs
written for GAOL~4: what they have to change to compile with GAOL v5, and what they will
see differently. The conclusion sums up what GAOL v5 established and what is
left open. The bibliography comes before two appendices, which list the
documentation and the references of GAOL, with the way to report a bug, and
reproduce the licences of GAOL and of CORE-MATH; an index ends the manual.

\clearpage
\pagestyle{YO}

%%==========================================================================
%% The chapters
%%==========================================================================
%% !TEX root = ../manual_gaolv5.tex
% Chapter 1.
%% Chapter 1: Interval arithmetic and floating-point numbers. Sources:
%% gaol.tex of f5bf180 141-257, 840-938, 1109-1150, 1913-1996, 2840-2846, 3493-3546,
%% 4070-4099, 5532-5546, 5574-5620, 5781-5788, 5881-5931;
%% doc/math-libraries.md (whole); doc/accuracy.md 1-47.
\chapter{Interval arithmetic and floating-point numbers}\label{chap:concepts}
\cindex{interval arithmetic}\cindex{floating-point numbers}%

A computer holds a real number as a floating-point number, a \code{double},
and a computation on doubles rounds each of its results: the quotient $1/3$,
the constant $0.1$ and most values of the elementary functions have no exact
representation, and a program that computes with them accumulates errors it
does not see. Interval arithmetic replaces each real number by two doubles
that enclose it, and each operation by an operation on intervals whose
result encloses every value the exact operation can take. The errors are
then not avoided but bounded: the result of an interval computation is a set
that contains the exact result, whatever the roundings along the way. GAOL
is a C++ library to perform arithmetic with floating-point intervals, and
this chapter gathers the notions on which the rest of the manual relies.

\medskip
In this chapter, we define the floating-point numbers and their roundings,
the floating-point intervals and the operations on them, and the vocabulary
of IEEE 1788-2015, the standard for interval arithmetic that GAOL v5
follows. The first section describes the doubles and the roundings
downward, upward and to nearest. The second defines the intervals, the
empty set, the unbounded intervals and the canonical intervals. The third
states the containment property and the outward rounding that enforces it,
with the property on which GAOL computes both bounds in a single rounding
direction. The fourth presents the dependency problem, which makes an
enclosure wider than the range. The fifth gives an overview of IEEE
1788-2015 as GAOL v5 follows it, and the sixth the idea of the relational
functions, which the standard calls reverse functions. The seventh explains
why the bounds of an elementary function depend on the mathematical library
that computes it, from the libm of the systems to CORE-MATH; the eighth
defines the three accuracy modes of the standard, tightest, accurate and
valid. The chapter ends with a discussion.

\section{Floating-point numbers and rounding}\label{sec:rounding-definitions}
%%------------------------------------------
\cindex{F@\FSet}\cindex{R@\RSet}\cindex{double}%

The numbers GAOL computes with are the \code{double}s of C++: the
floating-point numbers of the double format of IEEE
754~\cite{IEEE754:2008,Goldberg:01}, called binary64 since the 2008 revision
of the standard. A double is written on 64 bits: a sign, a significand of
53 bits, 52 of them stored, and an exponent of 11 bits; its value is
$\pm m\cdot2^{e}$ with $m$ an integer below $2^{53}$. We write \FSet{} for
the set of the doubles, the two infinities included, and \RSet{} for the set
of the real numbers. \FSet{} is finite and symmetric: $y\in\FSet$ if and
only if $-y\in\FSet$.

\medskip
The doubles are not spread evenly on the real line. Between two consecutive
powers of two, $2^{e}$ and $2^{e+1}$ with $-1022\leq e\leq1023$, they are $2^{52}$ in number, spaced by
$2^{e-52}$: this spacing is the \dfn{unit in the last place} (ulp) of the
doubles between these powers. The \dfn{normal numbers} are those of
magnitude between $2^{-1022}$, the smallest normal double, and
$\mathrm{MAX}=(2-2^{-52})\cdot2^{1023}\approx1.797693\cdot10^{308}$, the
largest double, which C names \code{DBL\_MAX}. Below $2^{-1022}$ lie the
\dfn{subnormal numbers}, the multiples of $2^{-1074}$, which fill the gap
between the smallest normal double and zero with the same spacing. The
bounds GAOL computes may be subnormal, and the modes of the processors that
flush the subnormal numbers to zero break them (\pxref{sec:ftz}).
\cindex{subnormal numbers}\cindex{unit in the last place}%

\medskip
Three kinds of values are special. The two zeros, $-0$ and $+0$, are equal in the
comparisons and stand for the same number; IEEE 1788-2015 and GAOL take them
for the same bound, and the sign a zero bound gets differs between the
builds of GAOL (\pxref{sec:zero-sign}). The two infinities, $-\infty$ and
$+\infty$, which an overflow or a division by zero produces, are the bounds
of the unbounded intervals. NaN (\emph{not a number}) is the result of an
invalid operation, as $0/0$ or $\sqrt{-1}$; it compares unequal to
everything, itself included, and GAOL holds the empty set as an interval
whose two bounds are NaN (\exref{sec:intervals}).
\cindex{NaN}\cindex{zero!sign of a bound}\cindex{infinity}%

\medskip
A real number that is no double has to be rounded. IEEE 754 defines four
rounding directions for the binary formats of the processors: to nearest, toward $-\infty$, toward $+\infty$ and
toward zero. The processor holds the direction in effect in a control
register, and every operation of its floating-point unit rounds its result
in that direction; a program starts in the rounding to nearest. Interval
arithmetic uses the first three, which we define now.
\cindex{rounding!direction}%

\begin{Def}[Round down, round up]\label{def:round-down-up}
\cindex{rounding!downward|hyperpagebf}\cindex{rounding!upward|hyperpagebf}%
\cindex{downward rounding ($\roundDn{\:}$)}%
\cindex{upward rounding ($\roundUp{\:}$)}%
\index[idx]{aaa@$\roundDn{\:}$|see{downward rounding}}%
\index[idx]{aab@$\roundUp{\:}$|see{upward rounding}}%
Given \RSet{} the set of real numbers and \FSet{} the set of floating-point
numbers (\code{double}), we have, for all $x\in\RSet$:
\begin{align*}
  \roundDn{x} &= \max\{y\in\FSet\mid y\leq x\},\\
  \roundUp{x} &= \min\{y\in\FSet\mid y\geq x\}.
\end{align*}
$\roundDn{x}$ is the rounding of $x$ \dfn{downward}, toward $-\infty$, and
$\roundUp{x}$ its rounding \dfn{upward}, toward $+\infty$.
\end{Def}

\begin{Def}[Round to nearest]\label{def:round-nearest}
\cindex{rounding!to nearest|hyperpagebf}%
Given \RSet{} the set of real numbers and \FSet{} the set of floating-point
numbers (\code{double}), $\roundNearest{x}$ is, for $x\in\RSet$, the
$y\in\FSet$ that minimizes $|x-y|$:
\begin{equation*}
  |x-\roundNearest{x}| = \min\{|x-y|\mid y\in\FSet\},
\end{equation*}
the one whose significand is even when two are equally near; a real of magnitude
at least $2^{1024}-2^{970}$ is rounded to the infinity of its sign, as IEEE 754
does.
\end{Def}

For a double $d$, the three roundings return $d$ itself: a double is never
moved. For a real $x$ that is no double, $\roundDn{x}$ and $\roundUp{x}$
are two consecutive doubles, and $\roundNearest{x}$ is one of them. Beyond
the largest double, the definitions give $\roundUp{x}=+\infty$ and
$\roundDn{x}=\mathrm{MAX}$ for $x>\mathrm{MAX}$.

\begin{Def}[Neighbours of a double]\label{def:neighbours}
\cindex{double!neighbours}%
For a double $d$, $\nextdouble{d}=\min\{y\in\FSet\mid y>d\}$ is the double
that follows $d$, and $\prevdouble{d}=\max\{y\in\FSet\mid y<d\}$ the double
that precedes it. For a real $x$ that is no double,
$\roundUp{x}=\nextdouble{(\roundDn{x})}$. The function \code{next\_float()}
of GAOL computes the successor of a double (Chapter~\ref{chap:errors}).
\end{Def}

\begin{Ex}\label{ex:one-tenth}
Programming with floating-point numbers is one of the few activities where
the compiler has to be considered with distrust. Let us consider the
following piece of code:
\begin{example}
#include <gaol/gaol>  @rem^// We do not import the gaol namespace~

using gaol::interval;

int main(void) {
  interval one_tenth(0.1); @rem^// BEWARE: wrong!~

  @emph^[some code using one_tenth]~
  gaol::cleanup();
}
\end{example}
Though a rational perfectly representable in decimal, \code{0.1} is not
representable in binary, at least not with a finite number of bits, and
thus requires rounding. The purpose of the programmer was to define an
interval containing this value. However, \code{0.1} is rounded \emph{at
compile time}, most certainly to the nearest representable floating-point
number $\roundNearest{0.1}$. As a consequence, \code{one\_tenth} is a
degenerate interval containing only $\roundNearest{0.1}$, and the
containment property (\exref{sec:containment}) is violated. The right way
to write such a constant is given in Example~\ref{ex:one-tenth-right}.
\end{Ex}

\section{Intervals}\label{sec:intervals}
%%-------------------
\cindex{interval!floating-point}\cindex{I@\ISet}%

\begin{Def}[Floating-point interval]\label{def:interval}
A \dfn{floating-point interval} $I=\itv{a}{b}$, with $a$ and $b$ in \FSet,
$a\leq b$, $a<+\infty$ and $b>-\infty$, is the set of the real numbers
between its bounds:
\begin{equation*}
  \itv{a}{b}=\{x\in\RSet\mid a\leq x\leq b\}.
\end{equation*}
$a=\leftBound{I}$ is its \dfn{lower bound}, or left bound, and
$b=\rightBound{I}$ its \dfn{upper bound}, or right bound. The empty set
$\emptyset$ is an interval as well. We write \ISet{} for the set of the
floating-point intervals.
\end{Def}

An interval is a set of real numbers: an infinite bound marks that the
interval is unbounded on that side, and is no member of it. The intervals
with an infinite bound, as \itv{1}{+\infty}, are the \dfn{unbounded}
intervals, intervals as any other; \itv{-\infty}{+\infty} is the whole real
line, which IEEE 1788-2015 calls \ieee{entire} and GAOL
\code{interval::universe()}. The point intervals $\itv{+\infty}{+\infty}$
and $\itv{-\infty}{-\infty}$ do not exist: the standard has no such
interval, and the constructors of GAOL v5 give the empty set for them
(\pxref{par:inf_inf_note}).\newinvfive{} An interval $\itv{a}{a}$ reduced to
one double is a \dfn{degenerate} interval, or point interval, which the
literals of the standard write $[a]$. The empty set is the result of an
operation whose argument has no point in the domain of the function
(\exref{sec:ieee1788-overview}); GAOL holds it as an interval whose two
bounds are NaN, and the method \code{is\_empty()} tells it.
\cindex{interval!unbounded}\cindex{interval!degenerate}\cindex{empty set}%

\begin{Def}[Hull]\label{def:hull}
\cindex{hull}%
The \dfn{hull} of a set $S$ of real numbers, written $\hull{S}$, is the
smallest floating-point interval containing $S$: $\hull{\emptyset}=
\emptyset$, and for a non-empty $S$
\begin{equation*}
  \hull{S}=\itv{\roundDn{\inf S}}{\roundUp{\sup S}},
\end{equation*}
where $\inf S=-\infty$ when $S$ has no lower bound and $\sup S=+\infty$ when
it has no upper bound.
\end{Def}

\begin{Def}[Canonical interval]\label{def:canonical}
\cindex{interval!canonical|see{canonical interval}}%
\cindex{canonical interval|hyperpagebf}%
A non-empty interval $I=\itv{a}{b}$ is \dfn{canonical} if and only if
$a\geq\prevdouble{b}$: it contains at most two consecutive floating-point
numbers.
\end{Def}

The hull of a real number $x$ that is no double,
$\hull{\{x\}}=\itv{\roundDn{x}}{\roundUp{x}}$, is canonical: it is the
tightest interval of doubles that encloses $x$, and no element of \ISet{}
does better.

\begin{Ex}\label{ex:one-tenth-right}
The right way to deal with rational constants that might not be perfectly
representable as floating-point numbers is to write them as strings, and to
read them with \code{gaol::textToInterval()} (\pxref{sec:constructors}), so
that they are correctly rounded downward and upward at \emph{run time}:
\begin{example}
#include <gaol/gaol>

using gaol::interval;

int main(void)
{
  @rem^// OK: this is the right way~
  interval one_tenth = gaol::textToInterval("0.1");
  @rem^// Another possible way~
  interval one_tenth2 = gaol::textToInterval("1/10");

  @emph^[some code using one_tenth]~
  gaol::cleanup();
  return 0;
}
\end{example}
Now, \code{one\_tenth} is the smallest floating-point interval enclosing
$0.1$: a canonical interval. Written in hexadecimal
(\pxref{sec:output-format}), which shows the bounds exactly, the interval of
Example~\ref{ex:one-tenth} and this one are:
\begin{onscreen}
interval(0.1) = [0x1.999999999999ap-4]
one_tenth     = [0x1.9999999999999p-4, 0x1.999999999999ap-4]
\end{onscreen}
\end{Ex}

\paragraph{Width, midpoint and radius.}
Three numbers describe an interval $\itv{a}{b}$: its \dfn{width} $b-a$
(\ieee{wid} of IEEE 1788-2015), its \dfn{midpoint} $(a+b)/2$ (\ieee{mid})
and its \dfn{radius} $(b-a)/2$ (\ieee{rad}), so that $\itv{a}{b}$ is the
set of the reals at a distance at most the radius from the midpoint.
Computed on doubles, none of the three is exact in general: the width GAOL
computes is rounded upward, and the midpoint is a double that the
computation has to keep inside the interval, for an unbounded interval
too~\cite{Goualard:TOMS2014}. Chapter~\ref{chap:functions} gives the
methods \code{width()}, \code{midpoint()} and \code{rad()}, and their
values on the empty set and on the unbounded intervals.
\cindex{interval!width}\cindex{interval!midpoint and radius}%

\section{Interval arithmetic and the containment property}
\label{sec:containment}
%%--------------------------------------------------------
\cindex{containment property}\cindex{outward rounding}%
\cindex{rounding!outward|see{outward rounding}}%

An operation on intervals stands for the operation on every pair of points
of its arguments. Its result has to contain all the values the real
operation can take and, to be of use, should contain as little else as
possible: the smallest floating-point interval that does is the hull.

\begin{Def}[Containment property]\label{def:containment}
The \dfn{containment property} of floating-point interval arithmetic
requires that, for any operation ``$\circ$'' and any intervals $I$ and $J$,
$I\circ J$ contains $\{i\circ j\mid i\in I,\ j\in J\}$, the pairs on which
$\circ$ is not defined being left out, and, for any function $f$ of one
variable and any interval $I$, that $f(I)$ contains $\{f(i)\mid i\in I\}$,
the points of $I$ outside the domain of $f$ being left out. The
\dfn{containment principle} asks for the smallest interval that does so, the
hull:
\begin{equation*}
  I\circ J=\hull{\{i\circ j\mid i\in I,\ j\in J\}},\qquad
  f(I)=\hull{\{f(i)\mid i\in I\}}.
\end{equation*}
\end{Def}

\begin{Ex}\label{ex:sqrt-functional}
For the square root, the principle gives
$\sqrt{I}=\hull{\{\sqrt{i}\mid i\in I\}}$. From monotonicity
considerations, the square root of \itv{1}{2} is then
$\sqrt{\itv{1}{2}}=\itv{1}{\sqrt{2}}$ on the real numbers; as $\sqrt2$ is
no double, the interval computed is $\itv{1}{\roundUp{\sqrt{2}}}$, whose
upper bound is the double above $\sqrt2$.
\end{Ex}

Fulfilling the containment property with doubles requires \dfn{outward
rounding}: the lower bound of a result is rounded downward and its upper
bound upward. For example, to add the intervals \itv{a}{b} and \itv{c}{d},
we compute $\itv{\roundDn{a+c}}{\roundUp{b+d}}$, where $\roundDn{r}$ and
$\roundUp{r}$ return the greatest double at most $r$ and the smallest double
at least $r$ (Definition~\ref{def:round-down-up}). These two operations are
performed by switching the rounding direction of the floating-point unit
towards, respectively, $-\infty$ and $+\infty$. With
$P=\{ac,\,ad,\,bc,\,bd\}$ and $Q=\{a/c,\,a/d,\,b/c,\,b/d\}$, the four basic
operations on bounded intervals become
\begin{align*}
  \itv{a}{b}+\itv{c}{d} &= \itv{\roundDn{a+c}}{\roundUp{b+d}},\\
  \itv{a}{b}-\itv{c}{d} &= \itv{\roundDn{a-d}}{\roundUp{b-c}},\\
  \itv{a}{b}\times\itv{c}{d} &= \itv{\roundDn{\min P}}{\roundUp{\max P}},\\
  \itv{a}{b}\div\itv{c}{d} &= \itv{\roundDn{\min Q}}{\roundUp{\max Q}}
    \quad\text{if } 0\notin\itv{c}{d}.
\end{align*}
A divisor that contains zero gives the hull of the quotients by its
non-zero points, which IEEE 1788-2015 defines case by case: an unbounded
interval when the dividend is not \itv{0}{0}, as
$\itv{1}{2}\div\itv{-1}{1}=\itv{-\infty}{+\infty}$, the interval \itv{0}{0} when
it is, and the empty set when the divisor is $\itv{0}{0}$, as
$\itv{1}{2}\div\itv{0}{0}=\emptyset$. The infinite bounds and the zeros of
each operation follow the standard as well; Chapter~\ref{chap:arithmetic}
states the rules.

\begin{Prop}\label{prop:round-neg}
For every real number $r$, $\roundDn{-r}=-\roundUp{r}$.
\end{Prop}
\demo The set \FSet{} is symmetric: $y\in\FSet$ if and only if
$-y\in\FSet$. Hence $\{y\in\FSet\mid y\leq-r\}=\{-z\mid z\in\FSet,\ z\geq
r\}$, and the greatest element of the first set is the opposite of the
smallest element of the second:
$\max\{y\in\FSet\mid y\leq-r\}=-\min\{z\in\FSet\mid z\geq r\}$, that is,
$\roundDn{-r}=-\roundUp{r}$.\boite

\medskip
On most platforms, switching the rounding direction is costly.
Proposition~\ref{prop:round-neg} cuts down the number of switches: nearly
all downward rounding operations can be replaced by upward rounding ones by
negating appropriately twice the operations performed. The lower bound of a
sum, $\roundDn{a+c}$, is $-\roundUp{(-a)+(-c)}$: the sum of the negated
lower bounds, rounded upward, then negated. GAOL stores the lower bound of
an interval negated, as the pair $\langle -a, b\rangle$, and computes both
bounds rounding upward: the rounding direction can then be set towards
$+\infty$ once and for all. The sum of $\langle -a, b\rangle$ and
$\langle -c, d\rangle$ is the pair $\langle\roundUp{-a-c},
\roundUp{b+d}\rangle$, two upward additions, which the kernels of GAOL
compute on the two halves of one register (\pxref{sec:kernels}). How GAOL
sets the rounding direction, and what the upward rounding does to the rest
of the program, is the subject of Chapter~\ref{chap:rounding}
(\pxref{sec:rounding-direction}).
\cindex{rounding!direction}\cindex{lower bound!stored negated}%

\paragraph{Exact operations.} Not every operation needs a rounding. The
negation maps \itv{a}{b} to \itv{-b}{-a}, whose bounds are doubles since
\FSet{} is symmetric; the minimum and the maximum return a bound of an argument,
and the absolute value a bound of its argument, its opposite or $0$; and the
integer functions, floor,
ceiling, truncation, the roundings to the nearest integer and the sign,
return integers that are doubles, the integer part of a double being a
double. These operations are exact: their result is the hull of the exact
result with no rounding at all, whatever the rounding direction in effect.
The four basic operations, the square root and the elementary functions
round: their bounds need a rounding in a chosen direction, and, for the
elementary functions, depend on the mathematical library that computes them
(\exref{sec:math-libraries}).
\cindex{exact operations}%

\section{The dependency problem}\label{sec:dependency}
%%--------------------------------
\cindex{dependency problem}%

The containment principle makes each operation as tight as doubles allow,
but a sequence of operations is not: an expression evaluated on intervals
gives an interval that encloses the range of the expression, and is in
general wider than it. The reason is that each occurrence of a variable
ranges over its interval on its own, as if the other occurrences did not
exist. This is the \dfn{dependency problem} of interval
arithmetic~\cite{Moore:Book1966,Neumaier:Book1990}.

\begin{Ex}\label{ex:dependency}
Let $x=\itv{0}{1}$. The expression $x-x$ is zero for every value of $x$,
but the subtraction of intervals takes its two operands as independent:
$\itv{0}{1}-\itv{0}{1}=\itv{-1}{1}$. In the same way, with $x=\itv{-1}{1}$,
the product $x\times x=\itv{-1}{1}$ combines the value $-1$ of the first
factor with the value $1$ of the second, which one number cannot take at
once, whereas the square, which knows that its two factors are equal, gives
$x^2=\itv{0}{1}$, the range; GAOL provides it as \code{sqr()}
(Chapter~\ref{chap:arithmetic}). The distributivity is weakened in the same
way: for any intervals $x$, $y$ and $z$, the operations being computed
exactly on the reals,
$z\times(x+y)\subseteq z\times x+z\times y$, with equality not in general,
as $\itv{1}{2}\times(\itv{-1}{1}+\itv{3}{4})=\itv{2}{10}$ where
$\itv{1}{2}\times\itv{-1}{1}+\itv{1}{2}\times\itv{3}{4}=\itv{1}{10}$. Computed
with doubles and rounded outward, even the inclusion may fail by one double.
\end{Ex}

The effect grows with the expression. The first program of
Chapter~\ref{chap:overview} (\pxref{sec:first-program}) evaluates the
Goldstein-Price function on the two intervals \itv{-2}{2} of its variables,
and obtains an enclosure of its range that is 231 times as wide as the range:
each occurrence of $x$ and of $y$ in the formula ranges over \itv{-2}{2} on
its own, as if the others did not exist.
\cindex{enclosure}%

\medskip
Cutting $x$ and $y$ into smaller intervals, and taking the hull of the
values of $f$ over the pieces, narrows the enclosure, at the price of more
evaluations: the program \file{examples/03\_dependency\_problem.cpp} of the
sources does it. Writing the expression with fewer occurrences of each
variable, the square in place of a product of a variable by itself for
instance, is the other remedy: the width of an enclosure is a property of
the expression evaluated, not of the function alone.

\section{IEEE 1788-2015}\label{sec:ieee1788-overview}
%%-----------------------
\cindex{IEEE 1788-2015}%

IEEE 1788-2015~\cite{IEEE1788:2015} is the standard for interval
arithmetic. It fixes what the libraries had each decided their own way:
what an interval is, what the operations return at the bounds of their
domains, how the empty set, the infinite bounds and the comparisons behave,
how an interval is written and read, and how tight each operation has to
be. GAOL v5 follows the standard in the operations GAOL provides: the empty
set, the infinite bounds, the domains of the functions, the comparisons and
the interval literals. The few cases where it differs, mostly the
\code{pow} of the namespace \code{gaol}, are stated where the operations
are described (\pxref{sec:pow-namespaces}, Chapter~\ref{chap:ieee1788}),
and the tightness of each operation is documented, as the standard requires
(Chapter~\ref{chap:accuracy}).

\paragraph{The set-based flavor and the bare intervals.} The standard is
organized in \emph{flavors}, and defines one, the \dfn{set-based flavor}:
an interval is a closed connected subset of \RSet, that is, the empty set,
a bounded interval \itv{a}{b} or an unbounded one, and the exact result of
an operation is the set of its values on the points of its arguments that
lie in the domain of the function, which the interval returned encloses
(\exref{sec:accuracy-modes}). The intervals of GAOL are those of the
set-based flavor (\exref{sec:intervals}). The standard also attaches to
its intervals a \emph{decoration}, which records whether the function was
defined and continuous on the whole of its argument along the computation
(Clause 11). GAOL has no decorations: only bare intervals are provided, and
the other parts of the standard it does not provide are listed in
Chapter~\ref{chap:ieee1788} (\pxref{sec:ieee1788-missing}).
\cindex{decorations}\cindex{bare interval}\cindex{set-based flavor}%

\paragraph{The operations.} The standard lists its operations in tables.
Table~9.1 holds the forward elementary functions every implementation
provides, with the domain of each: the arithmetic operations, the powers,
the exponentials and logarithms, the trigonometric and hyperbolic functions
and their inverses, the integer functions, \ieee{abs}, \ieee{min} and
\ieee{max}. Table~10.1 of the standard holds the reverse functions
(\exref{sec:relational-idea}), Table~10.2 the numeric functions of an
interval, Tables~10.3 and 10.4 the comparisons, and Table~10.5 the
recommended functions, which an implementation should provide (10.6). GAOL v5
adds the operations of these tables that GAOL~4 lacked, which the standard
requires or recommends, and the introduction lists them.\newinvfive{} The
relations of Tables~10.3 and 10.4, \ieee{equal}, \ieee{subset},
\ieee{interior}, \ieee{disjoint}, \ieee{precedes}, \ieee{strictPrecedes},
\ieee{less} and \ieee{strictLess}, follow the definitions of the standard
for the empty set and for the infinite bounds too
(Chapter~\ref{chap:relations}).\newinvfive{} The table of
Chapter~\ref{chap:ieee1788} says which operation of the standard each
function of GAOL is.
\cindex{IEEE 1788-2015!relations}%

\paragraph{The conventions.} A few conventions of the standard run through
the whole manual.
\begin{itemize}
\item[$\bullet$] \textbf{The domains.} A function is computed on the part
  of its argument within its domain, which is that of IEEE 1788-2015
  (Tables~9.1 and 10.5): the part outside is left out, and an argument with
  no point of the domain gives the empty set, as $\code{log}(\itv{-4}{0})$
  does.\newinvfive{} Every function gives the empty set for an empty
  argument. GAOL, which has no decorations, says nothing of the part left
  out: $\code{sqrt}(\itv{-1}{4})$ is \itv{0}{2}.
  \cindex{domain of a function}%
\item[$\bullet$] \textbf{The empty set.} The empty set is an interval, the
  result of an operation on an argument outside the domain of the function,
  and an interval whose bounds are in the wrong order, as
  \code{interval(1, 0)}, is empty too. Its numeric functions are those of
  Table~10.2: \ieee{inf} and \ieee{sup} of the empty set are $+\infty$ and
  $-\infty$, where \code{left()} and \code{right()} of GAOL are NaN, and
  \ieee{inf} returns $-0$ for a left bound zero (12.12.8).
\item[$\bullet$] \textbf{The infinite bounds.} An interval is a set of
  real numbers: an infinite bound is no member of it, so that
  \ieee{isMember}$(m, x)$ is \code{false} for an infinite $m$ (10.6.3), and
  the point intervals at $\pm\infty$ do not exist. The relations follow
  these definitions for the infinite bounds too.
\item[$\bullet$] \textbf{The literals.} An interval is written as a
  \dfn{literal} of the standard (12.11): \code{[l, r]}, \code{[a]} for a
  degenerate interval, \code{[empty]} for the empty set and \code{[entire]}
  for the whole line, the case of the letters being ignored; a literal is
  read back as an interval containing the one written (13.2, 13.3). GAOL writes
  and reads them (Chapter~\ref{chap:io}, \pxref{sec:input-format});
  \code{textToInterval} of the namespace \code{gaol\_ieee1788} reads the
  names of the functions of the standard and returns the empty set for a
  string that is no interval (12.1.3), where that of \code{gaol} reads the
  names of GAOL and throws (\pxref{sec:ieee1788-text}).\newinvfive
  \cindex{interval literal}%
\item[$\bullet$] \textbf{The powers.} The \ieee{pow} of Table~9.1 is
  defined for $x>0$, and for $x=0$ when $y>0$; the power with an integer
  exponent is \ieee{pown}$(x, n)$, whose $n$ is an integer rather than an
  interval. The \code{pow} of the namespace \code{gaol} takes the integer
  power for an integer exponent, a negative base included, which is the
  main point where GAOL and the standard differ (\pxref{sec:pow-namespaces},
  Section~\ref{sec:powers}).
\end{itemize}

\medskip
\noindent GAOL names its operations its own way, and the names of the
standard differ: the reverse function \ieee{coshRev}$(c, x)$ is
\code{acosh\_rel(c, x)}, and \ieee{roundTiesToEven}$(x)$ is
\code{round\_ties\_to\_even(x)}. The type \code{interval} and the functions
of GAOL live in the namespace \code{gaol}, and the namespace
\code{gaol\_ieee1788}, which \file{gaol/gaol.h} brings along, gives each
operation of the standard that GAOL provides the name and the argument
order the standard gives it, so that a reader of the standard finds it
without looking for its translation (\pxref{sec:namespaces},
Chapter~\ref{chap:ieee1788}).\newinvfive
\cindex{namespace!gaol\_ieee1788}\cindex{IEEE 1788-2015!names}%

\section{Relational (reverse) functions}\label{sec:relational-idea}
%%--------------------------------------
\cindex{operator!relational}\cindex{relational operator}%
\cindex{relational function}\cindex{reverse function}%

A function on intervals, as the containment principle defines it, maps an
interval of arguments to the hull of its values: this is its
\dfn{functional} version. A second reading of the same function is of use
when the function appears in a relation between variables.

\begin{Ex}\label{ex:sqrt-relational}
The functional square root of \itv{1}{2} is
$\sqrt{\itv{1}{2}}=\itv{1}{\sqrt{2}}$ (Example~\ref{ex:sqrt-functional}).
Another interpretation of the square root is
\begin{equation*}
  \sqrt{_r\,I}=\hull{\{j\in\RSet\mid\exists i\in I\colon j^2=i\}},
\end{equation*}
the \dfn{relational square root}, which permits obtaining both negative and
positive values: $\sqrt{_r\,\itv{1}{2}}=\itv{-\sqrt{2}}{\sqrt{2}}$. This
operator arises when we consider the relation $x^2=y$, which can
alternatively be written $x=\sqrt{_r\,y}$. Here, the \dfn{functional square
root} is not suitable, since it would confine $x$ to $\itv{0}{+\infty}$ and
lose the negative solutions of the relation.
\end{Ex}

Some applications, mainly in the area of constraint programming, require
such operators: a contractor narrows the domains of the variables of a
relation, and needs, for the relation $y=f(x)$, the part of the domain of
$x$ whose image meets the domain of $y$. As a consequence, GAOL offers both
functional and relational versions of the main arithmetic operators, which
IEEE 1788-2015 has since required, as its \dfn{reverse functions}
(Table~10.1).

\medskip
A relational function of one argument, as the relational square root, is
enough when the preimage of an interval has finitely many pieces. For a
periodic function the preimage has infinitely many, and the domain of the
variable has to be taken into account.

\begin{Ex}\label{ex:acos-relational}
Consider the relation $y=\cos x$, where $x$ and $y$ are interval variables.
One would like to express this relation in the equivalent way
$x=\operatorname{acos} y$. However, the function $\operatorname{acos}$
cannot be used, because its result is always included in the interval
\itv{0}{\pi}. What is needed is a relational version of
$\operatorname{acos}$; but since, for any value $y$, there are infinitely
many values $x$ verifying $y=\cos x$, the domain of $x$ has to be taken
into account. GAOL therefore defines a binary operator \code{acos\_rel} as
\begin{equation*}
  \code{acos\_rel}(Y, X)=\hull{\{x\in X\mid\exists y\in Y\colon y=\cos x\}},
\end{equation*}
to be contrasted with the functional arccosine,
\begin{equation*}
  \code{acos}(Y)=\hull{\{x\in\RSet\mid\exists y\in Y\colon
    x=\operatorname{acos} y\}}.
\end{equation*}
Figure~\ref{fig:relational-cos}, in Chapter~\ref{chap:arithmetic}, presents
the different results obtained when computing $\code{acos}(J)$ or
$\code{acos\_rel}(J, I)$, where $J$ and $I$ play the parts of $Y$ and $X$.
An interval $X$ that meets no piece of the
preimage gives the empty set.\newinvfive
\end{Ex}

\begin{Def}[Relational function]\label{def:relational}
Let $f$ be a function of one real variable, $Y$ an interval of values and
$X$ an interval of arguments. The \dfn{relational function} of $f$ with
respect to $X$, its \dfn{reverse function} in the words of IEEE 1788-2015,
is
\begin{equation*}
  f_{\mathrm{rel}}(Y, X)=\hull{\{x\in X\mid\exists y\in Y\colon y=f(x)\}},
\end{equation*}
the hull of the part of $X$ whose image by $f$ meets $Y$. The relational
function of one argument is the case $X=\RSet$.
\end{Def}

GAOL names a relational function after the inverse of the function it
reverses, as \code{acos\_rel} for the cosine, \code{asin\_rel} for the sine
and \code{sqrt\_rel} for the square; the standard names it after the
function itself, \ieee{cosRev}, \ieee{sinRev}, \ieee{sqrRev}, and the
namespace \code{gaol\_ieee1788} gives these names
(Chapter~\ref{chap:ieee1788}). The relational version of a binary
operation, as the relational division \code{div\_rel}, which reverses the
product, takes one more argument. The reference of the relational
arithmetic, with the figure of the relational cosine, is in
Chapter~\ref{chap:arithmetic} (\pxref{sec:relational-arithmetic}).

\section{The elementary functions and the mathematical libraries}
\label{sec:math-libraries}
%%--------------------------------------------------------------
\cindex{elementary functions}\cindex{mathematical library}%
\cindex{CORE-MATH}\cindex{libm}\cindex{mathlib}\cindex{CRlibm}%

The bounds of the four basic operations and of the square root are the
business of the processor: IEEE 754 requires them correctly rounded in the
direction in effect, and the roundings of Definition~\ref{def:round-down-up}
are what the floating-point unit computes. The elementary functions,
$\exp$, $\log$, $\sin$, $\cos$, $\tan$, their inverses and the hyperbolic
functions, are computed by software, by a mathematical library, and the
tightness of their bounds depends on it. GAOL v5 bounds its
elementary functions with CORE-MATH~\cite{CoreMath:2022}, compiled into it
from \file{3rd/math-core}, and takes no elementary function from any other
library: the C library gives it only \code{sqrt}, \code{floor}, \code{fma} and
the functions of \file{fenv.h}, and the integer powers, the roots and
\code{compoundm1} are computed by GAOL itself. Before that, GAOL was built with
one of three: mathlib by default with its \cmd{configure}, CRlibm, which its
meson build took by default, or the mathematical library of the system, whose options the three builds of GAOL
v5 no longer have (Chapter~\ref{chap:changes}).\newinvfive{} The
four come from four moments of the same history and were written with
different goals, which is why they give bounds of different tightness. This
section says where each of them comes from and what separates them, and why
GAOL v5 ended up with the last.

\subsection{The table maker's dilemma}
\cindex{table maker's dilemma}\cindex{correct rounding}%
\cindex{faithful rounding}%

A mathematical library evaluates $\exp$, $\log$, $\sin$\ldots{} in a
format, binary64 here, in which the exact result is not representable: it
has to round a value it only knows approximately. For some arguments the
exact result falls extraordinarily close to the middle of two consecutive
doubles, or to a double for the directed roundings, and deciding which way to round then asks for far more than 53
bits of intermediate precision. No bound on that distance follows from the
definition of the functions: this is the \dfn{table maker's dilemma}, named
by W. Kahan. Two ambitions follow from it, and they are the axis along
which the four libraries differ.

\begin{Def}[Faithful rounding]\label{def:faithful}
A library is \dfn{faithfully rounded}, or ``nearly correct'', when its
error is below one unit in the last place: it returns one of the two doubles
next to the exact value. The libraries of this kind are within about half a
unit in practice. This is enough for most numerical computation.
\end{Def}

\begin{Def}[Correct rounding]\label{def:correct-rounding}
A library is \dfn{correctly rounded} when it returns exactly the double the
rounding direction in effect gives of the exact result:
$\roundNearest{f(x)}$, $\roundDn{f(x)}$ or $\roundUp{f(x)}$. It makes a
program reproducible from one machine and one library to another, and, for
interval arithmetic, it gives the tightest bound a format can hold: a value
correctly rounded upward \emph{is} the upper bound, with nothing to add.
\end{Def}

\noindent IEEE 754-2008 lists the correctly rounded elementary functions
among its recommended operations, which is what the last two libraries
below set out to provide.

\subsection{libm, the mathematical library of the system}

libm is not a library but an interface: the mathematical part of the C
standard library, with one implementation per system. The line starts with
the libm of Unix V7 in the 1970s, then the libm of 4.3BSD in the middle of
the 1980s, written around W. Kahan at Berkeley (Coonen, K. C. Ng, Thomas,
Tang), where the culture of testing the accuracy of elementary functions
was formed. In 1993 Sun published fdlibm, the Freely Distributable libm of
K. C. Ng: portable C, no assembly, aiming at less than one unit in the last
place. fdlibm is the common ancestor of nearly everything that followed:
\file{msun} of FreeBSD, the libm of musl, openlibm of Julia, and part of
that of the glibc.

\medskip
Its design puts speed first and accuracy second. There is one evaluation
phase, an argument reduction and a polynomial or rational approximation
chosen by hand for each function; the error bound announced is empirical,
established by tests rather than proved; the mathematical properties
(monotonicity, symmetry, $\sin^2x+\cos^2x\leq1$) are not guaranteed.
Rounding to nearest is assumed throughout: nothing is promised in the
directed rounding directions, which is precisely what interval arithmetic
needs. Beside these portable implementations live the fast and
non-portable ones: Intel SVML, AMD libm, the \file{optimized-routines} of
ARM, and the vector variants.

\medskip
This is why GAOL could not keep the libm of the system as an option: the
bounds would not be certified, holding only where that libm stays within
about one double of the exact value.

\subsection{mathlib, the IBM Accurate Portable Mathematical Library}

mathlib, also called libultim, was written at IBM Haifa at the end of the
1990s (Abraham Ziv, Moshe Olshansky, Ealan Henis, Anna Reitman). It is the
industrial form of the idea Ziv published in 1991 in \emph{ACM TOMS}: the
\dfn{onion peeling} strategy. The function is evaluated with increasing
precision, a fast phase accurate to about half a unit in the last place
plus a little, then a second one, then a software multiprecision
arithmetic, and the computation stops as soon as a rounding test proves
that the rounding can be decided. Since the hard cases are rare, the
average cost stays close to that of an ordinary libm, while the result is
correctly rounded in practice.

\medskip
Its two limits are the ones that decided its fate. Nothing bounds the
number of iterations: the strategy has no proof and no bound on the worst
case, which can cost thousands of cycles. And it rounds to nearest only, so
a bound drawn from it has to be moved one double outward, which costs one
unit in the last place of tightness: this is what GAOL~4 did, built with
mathlib, for every elementary function but the hyperbolic ones
(\exref{sec:accuracy-modes}).

\medskip
mathlib was contributed to the glibc around 2001 and provided its
\code{exp}, \code{log}, \code{sin}, \code{cos}, \code{atan} and \code{pow}
for some fifteen years. In 2018 the glibc 2.28 removed these
multiprecision slow paths: the pathological execution times were judged
worse than the accuracy was worth. The glibc thereby gave up correct
rounding and went back to about half to one unit in the last place with a
bounded worst case. The version GAOL built was mathlib 2.1.1, under the GNU
LGPL, as Fr\'ed\'eric Goualard distributed it; GAOL v5 carries it no
longer.\newinvfive

\subsection{CRlibm, the proofs}

CRlibm is a project of the LIP at the ENS de Lyon and then of the INRIA
Ar\'enaire team (later AriC), from 2000 to about 2011: Jean-Michel Muller,
Florent de Dinechin, Catherine Daramy-Loirat, David Defour, Christoph
Lauter, with Guillaume Melquiond for the proof tools. Its last release,
1.0beta4, dates from 2010; it is no longer maintained. CRlibm keeps Ziv's
two phases but changes what is guaranteed.
\begin{itemize}
\item[$\bullet$] \textbf{The second phase is bounded.} Lef\`evre and Muller
  had computed, by an exhaustive search made tractable by their algorithms,
  the hardest-to-round cases in double precision: about 120 bits are known
  to be enough. No unbounded multiprecision is needed any more, a
  double-double or triple-double arithmetic suffices, and the worst-case
  execution time is bounded, typically a few times the average.
\item[$\bullet$] \textbf{The error bounds are proved}, and in part checked
  by machine with Gappa. The documentation of CRlibm is as much a proof
  document as a manual, and this is its main scientific contribution.
\item[$\bullet$] \textbf{The four rounding directions are provided
  explicitly}: \code{cos\_rn}, \code{cos\_rd}, \code{cos\_ru},
  \code{cos\_rz}. This is what makes it natural for interval arithmetic,
  and why GAOL took its functions rounded downward and upward directly,
  without moving them outward: built with CRlibm, the \code{exp},
  \code{log}, \code{sin}, \code{cos}, \code{tan}, \code{asin}, \code{acos},
  \code{atan}, \code{sinh} and \code{cosh} of GAOL were the tightest bounds
  of the value at the bound. CORE-MATH gives the same, in the direction in
  effect, for every function and on every system
  (\exref{sec:accuracy-modes}).
\end{itemize}
Its price was the work of a specialist for each function, a narrow coverage
(a dozen functions, binary64 only: neither \code{tanh} nor \code{asinh},
\code{acosh}, \code{atanh}), and an LGPL licence that kept it out of the
libm of the systems. Its direct descendant is Metalibm (de Dinechin,
Kupriianova, Lauter): generating the code and its proof rather than writing
them.

\subsection{CORE-MATH, correct rounding made adoptable}

CORE-MATH~\cite{CoreMath:2022} started in 2021 at Inria (Caramba team,
Nancy), by Paul Zimmermann, also behind MPFR, with Alexei Sibidanov, Tom
Hubrecht, S\'el\`ene Corbineau and others. It draws the lessons of the
three libraries above, and its design is as much a strategy of adoption as
a technique.
\begin{itemize}
\item[$\bullet$] \textbf{The stated goal is correct rounding at least as
  fast as the current libms.} It answers the reason mathlib was removed
  from the glibc: if correct rounding costs nothing, the argument from
  latency falls.
\item[$\bullet$] \textbf{The licence is MIT}, deliberately permissive, so
  that the code can be taken up by the glibc, the libm of LLVM or musl. It
  answers the reason CRlibm was not adopted, and it works: routines of
  CORE-MATH have been integrated upstream. It is also what lets GAOL v5
  take the sources it needs as they are, in \file{3rd/math-core}, and compile
  them into its library.
\item[$\bullet$] \textbf{The coverage is wide}: the functions of C23, in
  binary32, binary64 and binary128, the recent functions included.
\item[$\bullet$] \textbf{The validation is graded}: in binary32 it is
  exhaustive, the $2^{32}$ arguments being enumerated and compared with
  MPFR; in binary64 and beyond, a proved fast phase and an accurate phase,
  with a search for the hard cases.
\item[$\bullet$] \textbf{The four rounding directions are honoured}, the
  one in effect included. This is what GAOL v5 rests on: correctly rounded
  in the upward rounding GAOL computes in, CORE-MATH gives the tightest
  bounds without switching the rounding direction: the interval \code{log} of
  GAOL v5 takes 20.9\,ns where that of GAOL~4.2.3, which calls mathlib, moves
  its value outward and switches the rounding direction twice, takes 46.8\,ns
  (Intel i7-1185G7, Clang 21.1, 17 October 2026, \pxref{sec:comparison}).
\end{itemize}
A proposal to add \code{cr\_*} functions to the C standard comes from the
same team, so that correct rounding can be asked for by the program rather
than imposed on everyone. Beside it, and of an orthogonal design, is RLIBM
(Rutgers, Jay Lim and Santosh Nagarakatte, about 2019--2022): instead of
approximating the real function and then rounding, it looks by linear
programming for a polynomial that directly yields the correct rounding for
every input. Elegant, but practicable mainly for the short formats
(binary32, bfloat16, posits).

\subsection{What separates them}

Table~\ref{tab:math-libraries} sets the four libraries side by side. The
line of the story is plain: libm established that a mathematical library
had to be portable and fast; mathlib showed that correct rounding was
reachable on average, and was rejected on its worst case; CRlibm showed
that it could be proved and bounded, and was not adopted for want of a
licence, of coverage and of maintenance; CORE-MATH takes all of it up and
treats industrial adoption, that is, performance, MIT licence, C23 coverage
and standardisation, as a design constraint of the same rank as accuracy.

\begin{table}[htbp]
  \centering\footnotesize\renewcommand{\arraystretch}{1.15}
  \setlength{\tabcolsep}{4pt}
  \caption{The four mathematical libraries GAOL has been built with} \label{tab:math-libraries}
  \begin{tabular}{L{1.6cm}L{2.8cm}L{2.5cm}L{2.9cm}L{2.9cm}}
    \hline
    & libm (fdlibm, glibc) & mathlib (IBM) & CRlibm & CORE-MATH\\
    \hline\hline
    Time & 1980s, 1993-- & about 1997--2001 & 2000--2011 & 2021--\\
    \hline
    Where from & Berkeley, Sun & IBM Haifa & ENS de Lyon, INRIA
      & INRIA, CERN\\
    \hline
    Goal & speed, below 1 ulp & correct rounding in practice
      & correct rounding proved & correct rounding and speed\\
    \hline
    Method & one phase & unbounded onion peeling
      & two phases, the second bounded by the hard cases
      & proved phases, exhaustive in binary32\\
    \hline
    Evidence & tests & tests & Gappa, published proofs
      & proofs and exhaustive checks\\
    \hline
    Worst case & bounded & unbounded & bounded & bounded\\
    \hline
    Rounding & to nearest & to nearest & the four, explicit
      & the four, the one in effect\\
    \hline
    Formats & all & binary64 & binary64 & binary32, 64, 128\\
    \hline
    Licence & permissive & LGPL & LGPL & MIT\\
    \hline
    State & alive & frozen & abandoned & active, being adopted\\
    \hline
  \end{tabular}

\end{table}

\subsection{What it means for the bounds of GAOL}

The hierarchy is read directly in the tightness of the bounds, which
Chapter~\ref{chap:accuracy} states function by function.
\begin{itemize}
\item[$\bullet$] \textbf{The mathematical library of the system} gave
  bounds that were not certified: the value rounded to nearest moved outward
  by $2\cdot2^{-52}$ relatively and by the smallest normal double, several
  doubles beyond the tightest one.
\item[$\bullet$] \textbf{mathlib} was correctly rounded to nearest, hence
  one double outward: the bounds were accurate within one double.
\item[$\bullet$] \textbf{CRlibm and CORE-MATH} are correctly rounded in the
  direction wanted, hence give the tightest bound of the value at the
  bound, with nothing to add.
\end{itemize}
\noindent This is also why GAOL v5 takes every one of its elementary
functions from CORE-MATH, above all the hyperbolic ones (neither mathlib nor
CRlibm has them all, and the libm of the system is not accurate enough
everywhere for its values moved outward to be bounds), and on every
architecture and with every compiler, \file{gaol/math/gaol\_u128.h} giving
the accurate phases the 128-bit integer where the compiler has none. The
bounds are the same on every architecture and with every compiler, and
there is no mathematical library to install, to find or to link along with
GAOL: CORE-MATH is compiled into it.\newinvfive

\paragraph{Further reading.} The onion peeling strategy is in A. Ziv,
\emph{Fast evaluation of elementary mathematical functions with correctly
rounded last bit}, ACM TOMS 17(3), 1991. The table maker's dilemma, the
hard cases and the two-phase evaluation are in J.-M. Muller,
\emph{Elementary Functions: Algorithms and Implementation}, and in the
\emph{Handbook of Floating-Point Arithmetic}; the floating-point arithmetic
itself is presented by Goldberg~\cite{Goldberg:01} and in the
standard~\cite{IEEE754:2008}. The documentation distributed with the sources
of CRlibm holds the proofs of its error bounds. CORE-MATH is presented
in~\cite{CoreMath:2022} and at \url{https://core-math.gitlabpages.inria.fr/};
the file \file{3rd/README.md} of the sources says what GAOL v5 takes from
it. The NEWS of the glibc 2.28 records the removal of the multiprecision
slow paths inherited from mathlib.

\section{The accuracy modes}\label{sec:accuracy-modes}
%%-----------------------------
\cindex{accuracy!modes}\cindex{tightest}\cindex{accurate}\cindex{valid}%
\cindex{tightness}%

The containment principle asks for the hull of the exact result, which the
basic operations reach but which the elementary functions, computed by a
library, may not. IEEE 1788-2015 therefore defines three accuracy modes
(12.10.1), for an operation $f$, its exact value $f^*$ and an input box
$x$, and requires an implementation to document the tightness of each of
its interval operations (12.10.3).

\begin{Def}[Tightest]\label{def:tightest}
$f$ is \dfn{tightest} when $f(x)$ is the hull, in doubles, of $f^*(x)$:
each bound is the nearest double on its side of the exact bound.
\end{Def}

\begin{Def}[Accurate]\label{def:accurate}
$f$ is \dfn{accurate} when $f(x)$ encloses $f^*(x)$ and is within
\begin{equation*}
  \mathrm{nextOut}\bigl(f_{\mathrm{tightest}}(\mathrm{nextOut}(x))\bigr),
\end{equation*}
$\mathrm{nextOut}$ moving each bound one double outward: roughly, $f(x)$
is no wider than the tightest result for an input one double wider, plus
one double on each side.
\end{Def}

\begin{Def}[Valid]\label{def:valid}
$f$ is \dfn{valid} when $f(x)$ encloses $f^*(x)$.
\end{Def}

The standard requires the basic operations (\ieee{neg}, \ieee{add}, \ieee{sub},
\ieee{mul}, \ieee{div}, \ieee{recip}, \ieee{sqr}, \ieee{sqrt}, \ieee{fma}),
the integer functions and \ieee{abs}, \ieee{min} and \ieee{max} to be
tightest; the other elementary functions, the reverse functions and the
recommended functions to be valid, and recommends them to be accurate
(12.10.2). In this manual, ``within $k$ doubles'' means that each bound is
at most $k$ doubles beyond the tightest one.

\medskip
GAOL bounds its elementary functions with CORE-MATH, which is correctly
rounded in the rounding direction in effect. Computing in the upward
rounding it keeps, GAOL takes the value at a bound as the upper bound, with
nothing to add, and the double below it as the lower one. The upper bound is
the tightest one, and so is the lower bound: the double below the upper bound
is the tightest lower bound when the exact value is no double, and where the
exact value is a double, as for $\log(1)=0$ and $\sin(0)=0$, the operations
return that value itself rather than the double below it
(Remark~\ref{rem:exact-value}). Where an algorithm
evaluates the function at the bounds of the interval, the bounds of the
interval are therefore the tightest ones. Before GAOL v5, built with mathlib,
the default of its \cmd{configure}, the elementary functions but the hyperbolic
ones were correctly rounded to nearest and moved one double outward, which is
at most one double beyond the tightest bound: accurate, and not tightest. The
hyperbolic functions came from the mathematical library of the system, whose
values are sometimes further than one double from the exact ones
(Remark~\ref{rem:accuracy-before}).\newinvfive

\begin{Rem}\label{rem:exact-value}
The double below the correctly rounded upper value is the tightest lower
bound only when the exact value is no double. When it is, as for
$\log(1)=0$, the value correctly rounded upward is the exact value itself,
and an operation that moved it one double down would be accurate, not
tightest. The operations of GAOL return the exact value where it happens,
and Chapter~\ref{chap:accuracy} says where.
\end{Rem}

\medskip
Chapter~\ref{chap:accuracy} documents the tightness of every operation of
\code{interval}, from the algorithm of each operation, as the standard
requires; the tests of \file{tests} check each statement on random and
special arguments (Chapter~\ref{chap:tests}), and the page \file{doc/accuracy.md}
gives, for each, the algorithm and the largest distance the tests
found.\newinvfive{} The statements hold for GAOL as the three builds make
it, whatever the rounding direction the calling code left, on every
architecture and with every compiler: for the SSE2, the NEON and the FPU
intervals, and for the AVX-512 path, whose bounds are those of the SSE2
intervals but for the sign of a zero bound of a product with a double. They hold under the conditions that chapter states: the
code including the headers of GAOL compiled with the flags of interval
arithmetic (\pxref{sec:compiling}), and the floating-point exceptions
masked as a program starts (\pxref{sec:fp-exceptions}). The intervals of
floats are not covered.

\section{Discussion}\label{sec:concepts-discussion}

This chapter has set the vocabulary of the manual. A floating-point
interval is a set of real numbers between two doubles, the empty set and
the unbounded intervals included; the containment property requires every
operation to return an interval that encloses the exact result, and
outward rounding is what enforces it on doubles. The property
$\roundDn{-r}=-\roundUp{r}$ lets GAOL compute both bounds in the upward
rounding, with the lower bound stored negated: a choice that spares the
switches of the rounding direction, at the price of leaving the direction
upward for the rest of the program, which Chapter~\ref{chap:rounding}
examines. The dependency problem is the limit of the method: an enclosure
is as tight as the expression allows, not as the function does.

\medskip
Two choices of GAOL v5 follow from the notions of this chapter. The
standard IEEE 1788-2015 fixes the special cases, the empty set, the
infinite bounds, the domains and the literals, and GAOL v5 follows it in
the operations it provides, under two sets of names; it provides bare
intervals only, without the decorations of the standard. The elementary
functions are bounded with CORE-MATH, whose correct rounding in the
direction in effect gives the tightest bounds without a switch of the
rounding direction, where the libm of the system gave uncertified bounds
and mathlib bounds one double too wide; the three accuracy modes of the
standard, tightest, accurate and valid, are the words in which
Chapter~\ref{chap:accuracy} states what each operation achieves.

\medskip
The next chapter leaves the concepts for the software: how GAOL v5 is
obtained, configured, built and installed on each platform
(Chapter~\ref{chap:installation}).

%% !TEX root = ../manual_gaolv5.tex
% Chapter 2.
%% Chapter 2: Installation. Sources: gaol.tex of f5bf180 lines 258-839,
%% doc/building.md (whole), doc/three-builds.md (whole).
%% Based on the manual of GAOL 4 by Frederic Goualard (gaol.tex).
%%
%% Copyright (c) 2026 ENSTA, France
%%
%% By: Jordan NININ (created 2026-11-06)
%%
%% GAOL is a software distributed WITHOUT ANY WARRANTY. Read the associated
%% COPYING.LIB file for information.
\chapter{Installation}\label{chap:installation}

GAOL v5 comes as source code, to be compiled for the computer it is used on.
The sources of CORE-MATH, which compute the elementary functions, are part of
it and are compiled into the library \file{libgaol} itself: there is no
mathematical library to install, to find or to link along with GAOL, and no
choice to make. The bounds GAOL computes are the same on every architecture
and with every compiler, and they are the tightest ones for the elementary
functions of CORE-MATH, and within a few doubles of them for a few other
operations (Chapter~\ref{chap:accuracy}).

\medskip
GAOL is built with CMake, which is the build to use. The autotools and meson
builds of GAOL are kept for those who use them. On a given machine with a
given compiler, GAOL behaves the same whichever build configured it: the three
give it the same macros and the same flags, and a few options, the same in the
three, are all that the user chooses. The rest of the configuration is decided
by checks of the compiler and of the machine.

\medskip
In this chapter, we describe how GAOL v5 is obtained, built, tested and
installed. The first sections give the platforms of the continuous
integration, where to get the sources, the tools needed, and the compilers and
the options that GAOL refuses because they gave bounds that did not enclose
the exact results, or cannot be shown to enclose them. Section~\ref{sec:three-builds} states what the three builds
agree on. Sections~\ref{sec:cmake} to~\ref{sec:install-meson} describe each
build, its options and the settings of the developers.
Section~\ref{sec:install-options-worth} measures what the options are worth.
The last sections present the targets of the tests and what an installation
holds.

\section{Supported platforms}\label{sec:platforms}
\cindex{platforms}%

The continuous integration of GAOL v5 builds it and runs its tests on the
platforms below. Its jobs, the configurations it refuses or leaves out, and
what each platform showed are described in Chapter~\ref{chap:tests}
(\pxref{sec:ci}).

\begin{itemize}
\item[$\bullet$] \textbf{Linux.} Ubuntu 24.04 and 26.04 on x86-64 and
  64-bit ARM, with GCC and Clang, and the compilers of Ubuntu 22.04 (GCC~11, Clang~14, the latter on x86-64
  only: the three builds refuse it on 64-bit ARM) on Ubuntu 24.04; Debian 12 and 13 on x86-64,
  64-bit ARM and 32-bit ARM (armhf), and Debian 12 on 32-bit x86;
  manylinux\_2\_28 on x86-64 and 64-bit ARM; Alpine (musl) on x86-64 and 64-bit
  ARM; Debian 13 on s390x, ppc64le and riscv64, and Ubuntu 24.04 on POWER9,
  these four under qemu.
\item[$\bullet$] \textbf{macOS.} macOS 15 and 26, on arm64 and x86-64, and
  macOS 27, on arm64, with AppleClang, the Clang of LLVM and GCC.
\item[$\bullet$] \textbf{Windows.} Visual Studio 2022 and 2026 on x86, x64 and
  arm64, with Visual C++ and with clang-cl, and Clang with the target of
  Visual C++; MinGW-w64 GCC 14 and 15 on x64, 12 to 15 on x86; MSYS2 (UCRT64 and
  MINGW32 with GCC, CLANG64 and CLANGARM64 with Clang); Cygwin (GCC, x64).
\end{itemize}

\noindent
The bounds GAOL computes are the same on all of them. On 64-bit ARM
processors, GAOL computes its intervals with NEON instructions, whatever the
compiler: GCC, Clang, AppleClang, clang-cl, Visual C++ and the Clang of
MinGW-w64 (\pxref{sec:neon}).\newinvfive

\section{Getting the software}
\cindex{gaol!web page}%

GAOL v5 is available at \url{https://github.com/Jordan08/GAOL}:
\cmd{git clone https://github.com/Jordan08/GAOL.git} gives the sources in the
directory \file{GAOL}, the root directory where the commands of this chapter are
run. The GAOL of
Fr\'ed\'eric Goualard, which GAOL v5 continues, is available at
\url{https://github.com/goualard-f/GAOL}, and presented in the GAOL section of
\url{https://frederic.goualard.net}.

GAOL comes as source code, to be compiled for the computer it is used on. The
sources of CORE-MATH are part of it (\file{3rd/math-core}): there is nothing
else to download. The archive of the sources that CPack makes from a checkout,
\file{gaol-\var{version}.tar.gz}, holds the same sources
(\pxref{sec:packages}).

\section{Prerequisites}

\subsection{Mandatory tools}
\cindex{tool!mandatory to compile GAOL}%

Two tools are needed to build GAOL:
\begin{itemize}
\item a C++11 and C99 compiler: GCC, Clang, AppleClang, Visual C++ (Visual
  Studio 2022 or later), clang-cl (the Clang of Visual Studio, Clang 16 or
  later), or GCC and Clang for Windows (MinGW-w64, MSYS2, Cygwin);
\cindex{C++!standard}%
\item CMake 3.14 or later, with the build tool of its generator (\cmd{make},
  \cmd{ninja}, Visual Studio); or, for the autotools build, a shell and
  \cmd{make}; or, for the meson build, meson and ninja.
\end{itemize}

\noindent
No mathematical library is needed: CORE-MATH is compiled into GAOL. Neither
flex nor bison is needed: the lexer and the parser they generate are committed
with the sources (Section~\ref{sec:install-maintainers}).

\subsection{Compilers and options refused}\label{sec:refused}
\cindex{compiler!refused}%

Each of the compilers and options below gave bounds not enclosing the exact
results, or cannot be shown to enclose them. The three builds refuse Clang for
32-bit ARM processors and the compilers of the second item when configuring,
the C compiler, which builds CORE-MATH, as well as the C++ compiler, with a
message naming what to use instead. They keep the options
away by giving GAOL the flags of interval arithmetic (\pxref{sec:compiling}).
The header \file{gaol/gaol\_config.h} refuses them again when compiling, for
the code including the headers of GAOL too, all but the second one, which no
macro of the compiler shows; it alone refuses the other compilers of the list,
MinGW-w64, Cygwin without the fused multiply-add instructions, Visual C++ and
clang-cl without \option{/fp:strict}, and the x87 unit, and the SSE2 intervals
on 32-bit Windows compiled without \option{-faligned-new=8}, whose programs
crashed rather than give wrong bounds.

\begin{itemize}
\item[$\bullet$] \textbf{Clang for 32-bit ARM processors.} It does not honour
  the rounding direction there: built by Clang 21, 4556 of 16000 random
  products, squares and cubes did not enclose their exact values. GCC does.
\item[$\bullet$] \textbf{A compiler saying of \option{-frounding-math}
  ``overriding currently unsupported rounding mode on this target''}, as
  Clang 14 for 64-bit ARM. Bounds of \code{pow()} and \code{nth\_root()} did
  not enclose the exact values. Clang 18 honours the rounding direction there.
\item[$\bullet$] \textbf{A MinGW-w64 whose \code{fma()} or \code{round()} is
  wrong.} On x64, the mingw-w64 before version 12 (GCC 11 to 13 of Chocolatey),
  and those of any version linked with \file{msvcrt.dll} rather than the UCRT
  (MSYS2 MINGW64, the cross compilers of Debian and Ubuntu): the two functions are those of the math
  library of mingw-w64, computed in doubles. That \code{fma()} is not
  correctly rounded (10.7~\% of the error-free products \code{fma(a, b, -a*b)}
  are wrong), and CORE-MATH computes with it wherever the compiler has no
  \option{-mfma} (\option{GAOL\_FMA} off; GCC for Windows on x64 and on 32-bit
  x86 takes it by default, and these toolchains stay refused, their
  \code{round()} being wrong too); that \code{round()} depends on the rounding
  direction. Bounds of \code{tan()}, \code{asin()}, \code{atan()} and others
  did not enclose the exact values. On 32-bit x86, the refusal covers mingw-w64 before 11 (GCC 11
  of WinLibs), whose \code{fma()} rounds each of its partial sums to a double,
  and, without \option{-mfma}, mingw-w64 12 and later linked with
  \file{msvcrt.dll} (MSYS2 MINGW32). On ARM it covers, as on x64, mingw-w64
  before 12 and any linked with \file{msvcrt.dll}; the refusal is not tested
  there. A toolchain of
  mingw-w64 12 or later linking the UCRT (MinGW-Builds GCC 14 and 15, MSYS2
  UCRT64, CLANG64 and CLANGARM64) takes the functions of
  \file{ucrtbase.dll}, which pass the tests, and is built and tested.
\item[$\bullet$] \textbf{Cygwin without the fused multiply-add instructions}
  (\option{GAOL\_FMA} off, or a processor without them). The \code{fma()} of
  the C library of Cygwin is not correctly rounded
  (\code{fma(1 + $2^{-30}$, 1 - $2^{-30}$, -1)} gives 0 instead of
  $-2^{-60}$), and \code{std::fma} is a call of it without optimization.
  GAOL calls the instruction instead (\option{-mfma}, \code{\_\_builtin\_fma}),
  which every level of optimization keeps, and Cygwin is built and tested with
  it.
\item[$\bullet$] \textbf{\option{-ffast-math}, \option{-Ofast} and
  \option{/fp:fast}.} The compiler then rounds to nearest and drops the checks
  of NaN and infinities. \option{-Ofast} is not refused where
  \option{-fno-fast-math} turns its fast-math optimizations off (below).
\item[$\bullet$] \textbf{\option{-ffinite-math-only}}, which
  \option{-ffast-math} and \option{-Ofast} turn on. The compiler then takes
  NaN and infinities never to occur, in the inline functions of GAOL too: the
  empty interval has NaN bounds, and \code{([1, 2] \& [3, 4]).is\_empty()} is
  false (GCC 9.4 and 13 at \option{-O0}, \option{-O2} and \option{-O3};
  Clang 18 at \option{-O0}, and at \option{-O2} and \option{-O3} when the
  bounds are \code{volatile}).
\item[$\bullet$] \textbf{Visual C++ without \option{/fp:strict}}, and clang-cl
  from Clang 16 without it (\option{/fp:precise}, their default). They then
  assume rounding to nearest, and may evaluate or rewrite floating-point
  operations accordingly; clang-cl compiles with \option{-fno-rounding-math
  -ffp-contract=on}. No test gave a wrong bound so, but nothing certifies the
  bounds.\newinvfive
\item[$\bullet$] \textbf{Doubles computed on the x87 unit of 32-bit x86
  processors}, in extended precision: without \option{-msse2 -mfpmath=sse},
  without \option{/arch:SSE2} with Visual C++, or without \option{-msse2}
  with clang-cl, which computes them on the x87 unit with \option{/arch:SSE2}.
  CORE-MATH assumes every operation on doubles rounded to a double. Computed
  in extended precision, its results rounded to nearest are rounded twice, and
  175 arguments gave the other neighbour of the exact value with GCC 12 on
  Debian 12 i386; with GCC 9, which rounds to nearest at compile time the
  constants CORE-MATH rounds in the direction in effect, the bounds of
  \code{exp2(-1075)}, \code{expm1(-800)} or \code{atan2()} of a tiny and a
  huge number did not enclose the exact values.
\item[$\bullet$] \textbf{The SSE2 intervals on 32-bit Windows, compiled without
  \option{-faligned-new=8}} (GCC and Clang). The C runtime aligns the memory of
  \code{new} on 8 bytes, and a \code{std::vector} of intervals crashed on its
  first \code{movaps} (\pxref{sec:three-builds}); \file{gaol.pc} and
  \code{gaol::gaol} give the flag.
\end{itemize}

\begin{Rem}
Two reasons for refusing MinGW-w64 are gone with
CORE-MATH.\newinvfive{} The math library of mingw-w64 older than version 12
gave \code{acosh()} near 1 up to 25 million doubles away from the exact
value, and the \code{fesetround()} of mingw-w64 12 runs the instruction
\code{cpuid} at each call. GAOL takes no elementary function from the math
library of the system any more, and no longer changes the rounding direction
for its elementary functions: the older MinGW-w64 are refused for the
\code{fma()} and \code{round()} that CORE-MATH and GAOL still take from the C
library. Before mingw-w64 12 they are those of mingw-w64 in every program;
from 12 on, only in the library of \file{msvcrt.dll}. The header
\file{gaol/gaol\_config.h} tells the two runtimes apart by \code{\_UCRT},
which the headers of a UCRT toolchain define, with GCC and Clang alike. On
32-bit x86, mingw-w64 computes the two functions on the x87 unit: there its
\code{round()} is right, and the \code{fma()} of mingw-w64 11 (MinGW-Builds
GCC 12 and 13), called only with \option{GAOL\_FMA} off, keeps its partial
sums in extended precision and passes the tests, though it is not correctly
rounded.
\end{Rem}

\medskip
\option{-fno-fast-math}, one of the flags of interval arithmetic, turns
\option{-ffast-math} and \option{-ffinite-math-only} off when it comes after
them on the command line, and not when it comes before them: \file{gaol.pc}
and \code{gaol::gaol} give it, and the code including the headers of GAOL is
refused when \option{-ffast-math} or \option{-ffinite-math-only} comes after
\option{-fno-fast-math}, where the compilation stops. With Clang,
\option{-ffast-math} and \option{-Ofast} turn \option{-frounding-math} off
too, which \option{-fno-fast-math} does not turn on again: the remedy is all
the flags of interval arithmetic after the option, not \option{-fno-fast-math}
alone. With GCC, \option{-fno-fast-math} turns off all the fast-math
optimizations of \option{-Ofast} but \option{-fcx-limited-range}, which
changes the arithmetic of complex numbers alone (GCC~9.4), wherever it stands
on the command line, before \option{-Ofast} too: \option{-frounding-math}
stays, and \option{-Ofast} is not refused there. The CMake tests
\code{refused\_finite\_math\_only} and \code{refused\_fast\_math} compile
\file{tests/refused\_options.cpp} with each of the two options, after the
flags of interval arithmetic, and check that the header refuses it, and
\code{refused\_positive} that it compiles without them (GCC and Clang);
\file{tests/refused\_options.sh} does the same in the autotools and meson
builds.

\medskip
No macro of the compiler shows the following options, which
\file{gaol/gaol\_config.h} cannot refuse and which give wrong results all the
same. The code using GAOL is not to be compiled with them.

\begin{itemize}
\item[$\bullet$] \textbf{\option{-funsafe-math-\allowbreak optimizations}, and
  \option{-ffast-math -fno-finite-\allowbreak math-only}, with GCC.} The compiler
  rewrites the addition \code{1.0 + tiny == 1.0}, by which GAOL sees the
  rounding direction, as \code{tiny == 0.0} (GCC 9.4), and the former probe of
  the modes, \code{1.0 + (subnormal + 0.0) == 1.0}, which sees the modes
  flushing the subnormals to zero too, as \code{subnormal == 0.0} (GCC 9.4 from
  \option{-O0}, and GCC 13), so that an operation does not set the direction
  upward again after the code using GAOL left it to nearest, and
  \code{width()} is below the exact width. Clang 18 does not rewrite the first
  addition, but its \option{-funsafe-math-optimizations} implies
  \option{-fno-signed-zeros}: at \option{-O2} and \option{-O3}, the second
  addition loses its \code{+ 0.0} there too. The probe of the modes is now
  \code{!(1.0 + (subnormal + 0.0) > 1.0)}, which GCC 15 does not rewrite
  under \option{-funsafe-math-optimizations} (GCC 9 to 14 were not tried with
  it); GCC 15 and Clang 21 drop its
  \code{+ 0.0} under \option{-fno-signed-zeros} as they did for the earlier
  form.
\item[$\bullet$] \textbf{\option{-fno-signed-zeros}}, with which GCC (9.4 from
  \option{-O0}, and 13) and Clang 18 (at \option{-O2} and \option{-O3}) drop
  the \code{+ 0.0} of the second addition, which then misses the
  flush-to-zero mode (\pxref{sec:ftz}).\newinvfive
\item[$\bullet$] \textbf{\option{-fno-honor-nans} alone, with Clang}, which
  does to the empty interval what \option{-ffinite-math-only} does:
  \code{\_\_FINITE\_MATH\_ONLY\_\_} is 1 only with
  \option{-fno-honor-infinities} too.
\item[$\bullet$] \textbf{\option{-ffast-math} or \option{-Ofast} followed by
  \option{-fno-fast-math}, with Clang.} The macros show nothing, but
  \option{-frounding-math} stays off, and the compiler takes the rounding
  direction to be to nearest. The test \file{tests/rounding\_direction.cpp},
  compiled so by Clang 18, failed 54 of its 292728 checks after
  \option{-ffast-math -fno-fast-math} (operations with an empty operand raised
  the invalid-operation flag), and 126 after \option{-Ofast -fno-fast-math}
  (bounds with subnormal results under flush-to-zero too); with the flags of
  interval arithmetic given again after the option, none. With the FPU
  intervals, the relations in loops are among what fails
  (\code{certainly\_le()}, \code{certainly\_leq()}, \code{certainly\_ge()},
  \code{certainly\_geq()} and the operators \code{<}, \code{<=}, \code{>},
  \code{>=}): at \option{-O3 -march=native}, Clang 18 and 20 vectorize them
  into the signaling comparisons \code{vcmplepd}, \code{vcmpnlepd} and
  \code{vcmpnltpd}, which they make for the empty operands too, whose bounds
  are NaNs (120 of 184 runs of such loops raised the flag, none with
  \option{-frounding-math}; \option{-fno-fast-math} or
  \option{-ffp-contract=off} alone change nothing).
\end{itemize}

\paragraph{Linking with \option{-Ofast}, \option{-ffast-math} or
\option{-funsafe-math-optimizations}.}
\cindex{flush-to-zero}%
The compilation of the headers of GAOL refuses \option{-ffast-math}, but the
link of the program is not seen. GCC and Clang link \file{crtfastmath.o} into
a program linked with one of these three options, on x86 and on ARM Linux,
whose constructor sets the modes that flush the subnormal numbers to zero:
flush-to-zero and denormals-are-zero of the SSE instructions, FZ of ARM
(\pxref{sec:ftz}). Loading a plug-in or a Python module built with
\option{-Ofast} sets them too (built by Clang 18, or by GCC before 13). With
one of them, every operation with a subnormal operand or result gives bounds
that miss the exact result: $\itv{10^{-300}}{10^{-300}} \times
\itv{10^{-20}}{10^{-20}}$ is $\itv{0}{0}$. The \option{-fno-fast-math} of
\file{gaol.pc} and \code{gaol::gaol} does not prevent it: it keeps
\file{crtfastmath.o} out of the link when it comes after \option{-ffast-math},
but not after \option{-Ofast} (GCC 9.4 and 13 and Clang 18 link it all the
same, and \option{-fno-fast-math} silences the error of
\file{gaol/gaol\_config.h} against \option{-ffast-math}, which \option{-Ofast}
would raise) nor, with GCC 9.4 and 13, after
\option{-funsafe-math-optimizations}. Linking with them, or loading code built
with them, makes the bounds wrong whatever the flags the code using GAOL was
compiled with. GAOL v5 defends itself in two ways.\newinvfive

\begin{itemize}
\item[$\bullet$] \textbf{Each operation clears the modes.} On x86 processors,
  and on ARM processors with GCC and Clang and with the NEON intervals of
  Visual C++, each operation that computes bounds clears the modes when it
  starts, in the check that sets the rounding direction upward, which uses a
  subnormal number to see them as well as the direction
  (\code{round\_upward\_if\_needed()} in \file{gaol/gaol\_fpu.h}). It makes
  the check before it reads, compares or copies a bound, the test of the
  empty set aside, the compiler being kept from moving a comparison above it
  (a compiler barrier after the modes are cleared, and an empty asm statement
  for a double passed by value), and with \option{GAOL\_PRESERVE\_ROUNDING}
  makes its result before it sets the modes back. It is the only defence
  against the plug-in, and against the compilers without
  \option{-mno-daz-ftz}; its cost, and the functions that make no check, are
  given in Chapter~\ref{chap:rounding} (\pxref{sec:ftz}).
\item[$\bullet$] \textbf{\option{-mno-daz-ftz} is given to the link}
  (\code{gaol::gaol}, \file{gaol.pc}, the tests) where the compiler accepts
  it, which keeps \file{crtfastmath.o} out of it. Each build checks the
  compiler with a link, as the other flags are checked:
  \code{check\_cxx\_source\_compiles()} with
  \code{CMAKE\_REQUIRED\_LINK\_OPTIONS}, \code{AC\_LINK\_IFELSE} with
  \option{-Werror} and the \env{LDFLAGS} of the user,
  \code{has\_link\_argument()}. \code{gaol::gaol} gives it to a program built
  by GCC 13 or a later one, or by the compiler that built GAOL, of its major
  version and not older: GCC has the option from 11.4 and 12.4 in the series
  11 and 12 only, and a GAOL built by GCC 11.4 gave it to GCC 12.3, which
  stops on it. \file{gaol.pc}, written for that compiler, gives it to every
  program, and a program linked by a compiler that refuses it (Clang 18, a GCC
  before 11.4, GCC 12.0 to 12.3, GCC for ARM) is to be linked without it.
\end{itemize}

\noindent
The modes of other processors, and of ARM with the FPU intervals of Visual
C++, are neither checked nor cleared; GCC links \file{crtfastmath.o} for none
of the other processors of the continuous integration (POWER, s390x, RISC-V).
The tests \file{tests/rounding\_direction.cpp} and
\file{tests/fast\_math\_link.cpp} check the two defences
(\pxref{sec:test-programs}).

\subsection{Tools for maintainers}\label{sec:install-maintainers}

GAOL uses code produced by GNU Flex and GNU Bison for parsing the expressions
used to initialize an interval.
\findex{flex}%
\findex{bison}%
The three builds compile the lexer and the parser committed with the sources
of \file{gaol/parser}, generated by flex 2.6.4 and bison 3.8.2 from
\file{gaol/parser/gaol\_interval\_lexer.lpp} and
\file{gaol/parser/gaol\_interval\_parser.ypp}:
\begin{itemize}
\item \file{gaol/parser/gaol\_interval\_lexer.cpp};
\item \file{gaol/parser/gaol\_interval\_parser.cpp};
\item \file{gaol/parser/gaol\_interval\_parser.h}.
\end{itemize}
\noindent
No build generates them again by itself, and none needs flex nor bison. After
a change of one of the two files, the target \cmd{parser} of the three builds
(\cmd{make parser}, \cmd{ninja parser}) generates the lexer and the parser
again in \file{gaol/parser}, with the flex and the bison of the system, and
they are committed; flex 2.6.4 and bison 3.8.2 give them again byte for byte.
The script \file{gaol/parser/regenerate\_parser.sh} runs the same commands
without a build (the variables \env{FLEX} and \env{BISON} give other
programs):

\begin{cmdshell}
% sh gaol/parser/regenerate_parser.sh
\end{cmdshell}

\subsection{Optional tools}

\begin{itemize}
\item \cmd{flex} and \cmd{bison}, to generate the lexer and the parser again
  (\cmd{make parser});
\item bash, and Python~3 with numpy, for \cmd{make perf} (\pxref{sec:tests});
\item \cmd{gcovr} and Python~3, for the report of the coverage
  (\option{GAOL\_COVERAGE}, \pxref{sec:cmake-options});
\item \cmd{pdflatex}, \cmd{bibtex} and \cmd{makeindex}, to build this manual.
\end{itemize}

\section{The three builds}\label{sec:three-builds}

On a given machine with a given compiler, GAOL behaves the same whichever
build configured it: CMake, \cmd{configure} and meson give it the same macros
and the same flags. The flags and the macros come from the autotools build of GAOL~4, then from
its meson build, and the CMake build gives the same, apart from the errors corrected
(Chapter~\ref{chap:changes}). In every build, by default:

\begin{itemize}
\item[$\bullet$] \textbf{GAOL is compiled in release.} The flags are
  \option{-O3} and the optimizations \cmd{configure} adds
  (\option{-funroll-loops}, \option{-fomit-frame-pointer} and
  \option{-fexpensive-\allowbreak optimizations},
  each where the compiler takes it), \option{-falign-functions=64} for the
  library alone (not for the tests, the examples and the benchmarks),
  \code{NDEBUG}, \option{-std=c++11} (\option{-std=c++14} for Clang with the
  target of Visual C++, whose standard library needs it), hidden visibility
  (\option{-fvisibility=hidden -fvisibility-inlines-hidden}) and
  \option{-Wall -Wconversion}, which GAOL compiles without warnings,
  \option{-Wsign-conversion} of Clang included. Its libraries are compiled in
  release too where a project brings GAOL in without choosing a build type, as
  by a \cmd{configure} calling GAOL's: a CMake project without
  \option{CMAKE\_BUILD\_TYPE}, with a generator of one configuration (with a
  generator of several, as Visual Studio, Release comes first among the
  configurations of GAOL built alone, and a project bringing GAOL in chooses
  its configuration), or a meson project in the default build type of meson,
  with its ninja backend.\newinvfive
\item[$\bullet$] \textbf{With the flags of interval arithmetic}
  (\pxref{sec:compiling}). The code using GAOL, the tests included, is linked
  with \option{-mno-daz-ftz} where the compiler accepts it (GCC 13 and later
  on x86, from 11.4 and 12.4 in the series 11 and 12, and Clang 20 and 21, not
  Clang 18), in \code{gaol::gaol} and in \file{gaol.pc}: each build tests the compiler with
  a link.\newinvfive
\item[$\bullet$] \textbf{SSE2 intervals on x86 processors}
  (\option{-msse2 -msse3}; no flag with Visual C++ and clang-cl for x64, where
  SSE2 is the baseline of the processor). On 32-bit Windows the C runtime
  aligns the memory of \code{new} on 8 bytes, and an SSE2 interval needs 16:
  GCC took it to be aligned on 16, and a \code{std::vector} of SSE2 intervals
  crashed. GCC and Clang (MinGW-w64, MSYS2) are told what the runtime gives
  with \option{-faligned-new=8}, so that they ask for aligned memory for the
  intervals and for what holds one, and \code{gaol::gaol} and \file{gaol.pc}
  give the flag to the code using GAOL too, as it allocates them;
  \file{gaol/gaol\_config.h} refuses the SSE2 intervals compiled without it.
  Visual C++ and clang-cl have no such flag (\option{/Zc:alignedNew} does
  nothing in C++14 with \cmd{cl}, and breaks the STL of clang-cl, and only
  C++17 turns the aligned \code{new} on): their SSE2 intervals keep their two
  bounds as two doubles, whose alignment is 8, and load and store them with
  \code{movupd} (\code{GAOL\_SSE2\_UNALIGNED\_STORAGE} of
  \file{gaol/gaol\_config.h}, which derives it from the target), so that they
  need no aligned memory, and the class allocators that aligned the intervals
  on 16 bytes are not compiled.\newinvfive
\cindex{SSE2}%
\item[$\bullet$] \textbf{NEON intervals on 64-bit ARM processors}
  (\option{GAOL\_SIMD}, \loption{enable-simd}, \option{-Denable-simd}; no
  flag), with every compiler of them: GCC, Clang and AppleClang under Linux
  and macOS, clang-cl and the Clang of MinGW-w64 and MSYS2, and Visual C++
  (\cmd{cl}, whose registers are those of \code{\mbox{\_ReadStatusReg}()} and
  \code{\mbox{\_WriteStatusReg}()}). Each build tests the compiler with the same
  program (CMake and meson for Visual C++, which \cmd{configure} does not
  build with): the intrinsics of \file{<arm\_neon.h>} and the registers of
  the rounding mode. Not with 32-bit ARM, whose NEON has no vector of
  doubles: the FPU computes the intervals there, as before
  (\pxref{sec:neon}). A universal build of macOS
  (\option{CMAKE\_OSX\_ARCHITECTURES} with several architectures, or
  \option{-arch} several times in the flags of \cmd{configure} and meson)
  compiles each source once per architecture with one configuration, and
  takes the intervals of the architecture being compiled, NEON for arm64 and
  SSE2 for x86\_64, where \option{GAOL\_SIMD} allows them
  (\code{GAOL\_UNIVERSAL\_NEON} and \code{GAOL\_UNIVERSAL\_SSE2} of
  \file{gaol\_configuration.h}). \option{GAOL\_SIMD} chooses between the
  vector intervals and the FPU on every processor: \option{GAOL\_ARM64},
  \loption{enable-arm64} and \option{-Denable-arm64} were an option of their
  own for the NEON intervals, independent of \option{GAOL\_SIMD}, and are
  gone.\newinvfive
\cindex{NEON}%
\cindex{universal build}%
\item[$\bullet$] \textbf{The fused multiply-add instructions of the
  processor}, for GAOL and CORE-MATH, where the compiler has a flag for them
  and the machine building runs a program compiled with them (no option of
  the user; \option{-DGAOL\_FMA=OFF} of the developers leaves them out,
  \pxref{sec:install-developers}): \option{-mfma} on x86; with GCC for
  Windows, on x86-64 and on 32-bit x86, \option{-mfma
  -mprefer-vector-width=128} (\pxref{sec:install-gcc-windows});
  \option{/arch:AVX2} with Visual C++ for x64, not for 32-bit x86, where it
  broke the rounding direction; \option{-mfpu=neon-vfpv4 -mfloat-abi=hard} on
  32-bit ARM. The check that the machine runs them matters even when it is
  not a cross-compilation: under Rosetta on macOS 14 (no longer in the CI),
  which had no AVX, the x86\_64 tests stopped on an illegal instruction. The
  exact products of GAOL (\code{fma\_once()}) and the functions of CORE-MATH
  then compute \code{fma()} with one instruction rather than with a call to
  the math library: on an Intel i7-1185G7 with Clang 21.1, \code{pow(x, 3)} of an
  interval takes 5.4~ns rather than 17.4, \code{sin} 1.3 times less and
  \code{cos} 2.5 times less (\pxref{sec:install-options-worth}). The code
  using GAOL is given the flag too (\code{gaol::gaol}, \file{gaol.pc}), and
  compiled for the same processor. \option{-ffp-contract=off} stays, which
  forbids the compiler to contract a multiplication and an addition into a
  fused one, as CORE-MATH asks.\newinvfive
\cindex{fused multiply-add}%
\item[$\bullet$] \textbf{The AVX-512 path} of $x + y$, $x - y$,
  $x \times y$, $x / y$, \code{sqrt(x)}, \code{sqr(x)} and, with FMA, the
  integer powers beyond 1 and $-1$, taken where the rounding direction is
  preserved: the path spares each of these operations the store of that
  direction (\option{GAOL\_PREFER\_AVX512}, a setting of the developers, asks
  for the path without \option{GAOL\_PRESERVE\_ROUNDING} and refuses it with it), and where the
  machine building runs the instructions (\code{GAOL\_AVX512\_ARITHMETIC}),
  which is checked as for the fused multiply-add instructions, the path being
  left out when cross-compiling. The library then needs a processor with
  them, which \code{gaol::init()} checks, and the operations test nothing as
  they run (\pxref{sec:avx512}). Never with GCC for Windows (MinGW-w64,
  Cygwin), which does not align the stack for the AVX-512 registers (GCC bug
  54412; nor for the 256-bit ones of AVX, which
  \option{-mprefer-vector-width=128} keeps MinGW-w64 GCC out of). No flag to
  compile: the functions of the path alone carry the target attribute, which
  Visual C++ has not and needs not, taking the intrinsics of any instruction
  set in any function. Visual C++ has no \code{\_\_builtin\_cpu\_supports()},
  and the one of clang-cl needs the runtime library of LLVM: both tell the
  processor by \code{cpuid} and \code{xgetbv}.\newinvfive
\cindex{AVX-512}%
\item[$\bullet$] \textbf{Without the intervals of floats},
  \code{gaol::intervalf} and \code{gaol::interval2f} (SSE3), unfinished, for
  which no build has an option: only a developer of GAOL compiles them,
  defining \code{GAOL\_FLOAT\_INTERVALS} in \file{gaol/gaol\_config.h}.
\item[$\bullet$] \textbf{With exceptions, the rounding direction left upward,
  and silent}: no line on the standard error when GAOL initializes and cleans
  up, unless \code{GAOL\_VERBOSE\_MODE} (\option{-DGAOL\_VERBOSE\_MODE=ON},
  \loption{enable-verbose-mode}, \option{-Denable-verbose-mode=true}) is asked
  for; the \cmd{configure} of GAOL~4 wrote it by default.\newinvfive
\item[$\bullet$] \textbf{The processor and the system read from the
  compiler.} \code{GAOL\_IX86\_MACOSX}, \code{GAOL\_ARM\_MACOSX}\ldots, the
  sizes of the integer types and the byte order are read from the macros of
  the compiler, in \file{gaol/gaol\_config.h}, rather than from the machine
  building.
\item[$\bullet$] \textbf{The same configuration header.} The three write the
  same macros into \file{gaol/gaol\_configuration.h}, and nothing else
  (\pxref{sec:install-configuration-header}).\newinvfive
\item[$\bullet$] \textbf{The same version file.} The three read the version
  of GAOL from the file \file{VERSION.txt}, the only one to change for a new
  version, by the same rules (\pxref{sec:install-version}).\newinvfive
\item[$\bullet$] \textbf{The same shared library.} The shared library is
  \file{libgaol.so.5.0.0}, whose soname is \file{libgaol.so.5}, where the
  build makes one (\cmd{configure}, meson, CMake with
  \option{BUILD\_SHARED\_LIBS}); it was \file{libgaol-5.0.so.0} with
  \cmd{configure} in the earlier versions of GAOL v5, and meson gave no
  version.\newinvfive
\cindex{soname}%
\end{itemize}

\medskip
In a Debug build (\option{CMAKE\_BUILD\_TYPE=Debug},
\cmd{configure --enable-debug}, \cmd{meson setup --buildtype=debug}), GAOL is
compiled with \option{-g}, without optimization nor \code{NDEBUG}, so that
\code{assert()} stays active. In GAOL~4, \cmd{configure --enable-debug} alone
defined \code{GAOL\_DEBUGGING}, which checked the assertions of GAOL
(\code{GAOL\_ASSERT}), adding \option{-g -ansi -Weffc++ -pedantic} for
\cmd{g++} alone to the optimizations of \loption{enable-optimize}; in the
earlier versions of GAOL v5, the Debug build of CMake did not define it, and
meson defined it with its option \code{enable-debug}.\newinvfive{} \option{GAOL\_DEBUG} of CMake and
\code{enable-debug} of meson build the libraries of GAOL and CORE-MATH so
whatever the build type or the configuration, a project bringing GAOL in
included, as \cmd{configure --enable-debug} does for a \cmd{configure} calling
GAOL's: the code of that project, and the tests and the examples, keep their
own.

\medskip
The directory \file{.github/audit} configures the three builds and compares
the macros GAOL sees, the macros each writes into
\file{gaol/gaol\_configuration.h}, and the flags GAOL is compiled with, those
of its C++ sources and those of a C source of CORE-MATH; the continuous
integration runs it on Ubuntu 24.04 (x86-64 and arm64, with GCC and Clang),
macOS 15 (arm64 and x86-64), Debian 12 on 32-bit x86, Debian 13 on 32-bit ARM
(armhf), and MSYS2 UCRT64 and MINGW32 (\pxref{sec:ci}).\newinvfive

\section{Building and installing with CMake}\label{sec:cmake}
\cindex{CMake}%

In the root directory of the sources:

\begin{cmdshell}
% cmake -S . -B build -DCMAKE_INSTALL_PREFIX=@var^prefix~ \
        -DWITH_TESTS=ON
% cmake --build build --config Release
% cmake --build build --config Release --target test
% cmake --install build --config Release
\end{cmdshell}

\noindent
With Visual Studio and Xcode the target is \cmd{RUN\_TESTS}: \cmd{cmake --build
build --config Release --target RUN\_TESTS}. On Linux with the static library
and the SSE2 intervals, the test step runs 47 tests (the other intervals give
three fewer, the shared library one fewer), of which \code{intervalf} and \code{interval2f} are skipped, the
intervals of floats being compiled only by a developer of GAOL. CMake 3.14 or
later is needed. The build compiles the thirty-eight sources of
CORE-MATH that GAOL uses into \file{libgaol}, which is a static library unless
\option{BUILD\_SHARED\_LIBS} is \code{ON}: there is nothing else to build and
nothing else to install. The build type is \code{Release} unless another is
given. Brought in by a project that gives none (\code{add\_subdirectory},
\code{FetchContent}), GAOL and CORE-MATH are compiled as in \code{Release} all
the same: \option{-O3} (\option{/O2} with Visual C++), \code{NDEBUG} and the
optimizations of \cmd{configure}; the earlier versions of GAOL v5 had
\option{-O3} alone, and CORE-MATH, compiled apart, no optimization. With a
generator of several
configurations (Visual Studio, Xcode, Ninja Multi-Config), GAOL built alone
puts \code{Release} first among them, which Ninja Multi-Config builds by
default, unless \option{CMAKE\_CONFIGURATION\_TYPES} is given; \cmd{cmake
--build} with Visual Studio builds \code{Debug} all the same without
\loption{config}.\newinvfive{} The \code{Debug} build type builds GAOL for
debugging (Section~\ref{sec:three-builds}).

\subsection{The version of GAOL}\label{sec:install-version}
\findex{VERSION.txt}%

The file \file{VERSION.txt} holds the version of GAOL, three numbers without
leading zeros such as \code{5.0.0}, and nothing else: the three builds read
it, and it is the only file to change for a new version.\newinvfive{} It is
not named \file{VERSION}: the root of the sources is on the include path of
the three builds, and on file systems that ignore case (macOS, Windows),
\code{\#include <version>}, which the headers of libc++ write, would find it
instead of the standard header.

\begin{itemize}
\item[$\bullet$] \textbf{CMake} reads it before \code{project()}, which gives
  it to \code{PROJECT\_VERSION}, with \code{gaol\_read\_version()} of
  \file{cmake/gaol\_version.cmake}; a change of \file{VERSION.txt} configures
  again. Where GAOL is the project built, it reads the line
  \code{PACKAGE\_VERSION=} of \file{configure} too, and warns when it
  configures that \file{configure} was generated for another version than
  \file{VERSION.txt} holds, since the archive of the sources that CPack makes
  holds \file{configure} as it is committed. It says nothing for a tree
  without \file{configure}, or with a \file{configure} that has no such line.
\item[$\bullet$] \textbf{meson} reads it in \code{project()}, with Python
  (\pxref{sec:install-meson}), which works with meson 0.53 as with later ones
  (\code{version: files()} needs meson 0.57); a change of \file{VERSION.txt}
  configures again.
\item[$\bullet$] \textbf{\cmd{configure}} reads it when it runs, and
  \cmd{make} runs it again when \file{VERSION.txt} changes
  (\code{CONFIG\_STATUS\_DEPENDENCIES} of \file{Makefile.am}). Only what
  \code{AC\_INIT} writes when autoconf generates \file{configure}
  (\cmd{configure --version}, \cmd{configure --help},
  \cmd{config.status --version} and the head of \file{config.log}) keeps the
  version of the last generation: \file{configure} is committed, and is to be
  generated again before a new \file{VERSION.txt} is committed, with the
  versions of autoconf, automake and libtool of
  Section~\ref{sec:configuration} and \cmd{AUTOHEADER=true autoreconf}
  (without \loption{force} \loption{install}, the other files stay as they
  are committed; \file{configure.ac} includes \file{VERSION.txt} with
  \code{m4\_include}, so that autoconf sees its change despite its cache,
  \file{autom4te.cache}). Until then, \cmd{configure} configures GAOL with the
  version of \file{VERSION.txt} and warns, twice, that it was generated for
  another one, and the continuous integration fails. With
  \loption{enable-maintainer-mode}, \cmd{make} generates \file{configure}
  again itself, and runs no autoheader over
  \file{gaol/gaol\_configuration.h.in}, which is written by hand.
\end{itemize}

\noindent
From \file{VERSION.txt} come the macros \code{GAOL\_MAJOR\_VERSION},
\code{GAOL\_MINOR\_VERSION}, \code{GAOL\_MICRO\_VERSION} and
\code{GAOL\_VERSION} of \file{gaol/gaol\_configuration.h}
(\code{gaol\_core::version\_major}\ldots{} of \file{gaol/gaol\_version.h},
\pxref{sec:version}), the shared library
\file{libgaol.so.\var{major}.\var{minor}.\var{micro}} and its soname
\file{libgaol.so.\var{major}}, the \code{Version} of \file{gaol.pc},
\file{gaolConfigVersion.cmake} (\code{find\_package(gaol \var{version})}
takes a version of the same major number, not older than the one asked for),
the names of the archive and of the packages of CPack, and the version printed
by the manual (\cmd{configure}, meson, and CMake with \option{WITH\_DOC}). Each build refuses a
\file{VERSION.txt} that does not hold three numbers without leading zeros,
which the macros write as C integers, and the three read it by the same
rules, as autoconf does for \cmd{configure --version}. The blanks and empty
lines around the version are ignored, the blanks being the six of ASCII
(space, tab, line feed, vertical tab, form feed and carriage return, the CR of
the line ends of Windows), not those of Unicode such as a no-break space; so
is a UTF-8 byte order mark (the bytes EF BB BF) at the start of the file,
which some editors of Windows write. The message of a build that refuses the
file quotes what it read and gives the first bytes of the file in
hexadecimal, which show what the quotation may hide: CMake and
\cmd{configure} quote a second byte order mark, or a zero-width space, as it
is, and it cannot be seen (meson writes it \code{\textbackslash{}ufeff},
\code{\textbackslash{}u200b}). A file of UTF-16 characters, which Windows
PowerShell 5 writes for a redirection (\code{"5.0.0" > VERSION.txt}), is not
decoded but refused, whether it starts with its byte order mark (FF FE, FE FF)
or not (it then holds NUL bytes, which a text file does not hold): the message
gives its bytes and ends with \emph{the file is UTF-16: save it as UTF-8 or
ASCII}, or, without the mark, \emph{the file holds a NUL byte as UTF-16 does:
save it as UTF-8 or ASCII}. The test \code{version\_file}
(\cmd{ctest -R version\_file}) checks the reading of CMake, and a script of
the continuous integration those of \cmd{configure}, meson and autoconf
(\pxref{sec:scripts}). The edition of this manual is its own, set in
\file{configure.ac}, \file{manual/meson.build} and \file{CMakeLists.txt}.

\subsection{Options}\label{sec:cmake-options}

The options are given to the first command, as
\option{-D\var{option}=\var{value}}. The three builds have the same
options.\newinvfive{} Each had options of its own before, and those of
the fused multiply-add instructions, of the AVX-512 path, of the 128-bit
integer, of the debugging and of the NEON intervals are gone as options: the
configuration decides them (Section~\ref{sec:install-decides}). The standard
options of each tool do the rest. Table~\ref{tab:install-options} gives the
options under their three names, with their defaults.

\begin{table}[htbp]
  \centering\renewcommand{\arraystretch}{1.1}\footnotesize
  \setlength{\tabcolsep}{4pt}
  \caption{The options of the three builds} \label{tab:install-options}
  \begin{tabular}{L{3.7cm}L{3.3cm}L{4.3cm}L{1.9cm}}
    \hline
    \multicolumn{1}{c}{CMake} & \multicolumn{1}{c}{\cmd{configure}} &
    \multicolumn{1}{c}{meson} & Default\\
    \hline\hline
    \code{CMAKE\_BUILD\_TYPE} & \loption{enable-debug} &
    \loption{buildtype=debug} & \code{Release}\\
    \code{CMAKE\_INSTALL\_PREFIX} & \loption{prefix} & \loption{prefix} &
    the system's\\
    \code{BUILD\_SHARED\_LIBS} & \loption{enable-shared} &
    \option{-Ddefault\_library=shared} & off\\
    \code{CMAKE\_POSITION\_\allowbreak INDEPENDENT\_CODE} &
    \loption{enable-pic} (\loption{with-pic} too) &
    \option{-Db\_staticpic=true} & on\\
    \code{CMAKE\_OSX\_ARCHITECTURES} & \option{-arch} several times in
    \env{CFLAGS}, \env{CXXFLAGS} and \env{LDFLAGS} & \option{-arch} several
    times in \code{c\_args}, \code{cpp\_args}, \code{c\_link\_args} and
    \code{cpp\_link\_args} & the machine's\\
    \code{WITH\_TESTS} & \loption{with-tests} & \option{-Dwith-tests=true} &
    off\\
    \code{WITH\_EXAMPLES} & \loption{with-examples} &
    \option{-Dwith-examples=true} & off\\
    \code{WITH\_DOC} & \cmd{make -C manual pdf} (no option) &
    \option{-Dwith-doc=true} & off\\
    \code{GAOL\_SIMD} & \loption{enable-simd} & \option{-Denable-simd=true} &
    on\\
    \code{GAOL\_PRESERVE\_ROUNDING} & \loption{enable-preserve-\allowbreak rounding} &
    \option{-Denable-preserve-\allowbreak rounding=true} & off\\
    \code{GAOL\_COVERAGE} & \loption{enable-coverage} &
    \option{-Denable-coverage=true} & off\\
    \hline
  \end{tabular}

\end{table}

\begin{itemize}
\item[$\bullet$] \textbf{\option{CMAKE\_BUILD\_TYPE}.} \code{Release} by
  default; \code{Debug} builds GAOL without optimization, with debugging
  information and without \code{NDEBUG}, so that \code{assert()} stays
  active, and enables the debugging facilities
  (\pxref{sec:debugging-facilities}). The \code{Release} build is
  \option{-O3}, \code{NDEBUG} and the optimizations of
  Section~\ref{sec:three-builds}. With \cmd{configure},
  \loption{enable-debug} brought into another project with
  \code{AC\_CONFIG\_SUBDIRS} is the \loption{enable-debug} of the
  \cmd{configure} of that project, which passes on its options, for the tests
  of GAOL as well.
\item[$\bullet$] \textbf{\option{CMAKE\_INSTALL\_PREFIX}.} Where
  \cmd{cmake --install} puts GAOL.
\item[$\bullet$] \textbf{\option{BUILD\_SHARED\_LIBS}.} \code{OFF} by default:
  a static library \file{libgaol.a}; \code{ON} builds the shared library
  \file{libgaol.so.5.0.0}, whose soname is \file{libgaol.so.5}, as
  \cmd{configure} and meson do (always a static library with Visual C++).
  Installed in a \var{prefix} that the dynamic loader does not search, the
  shared library is found on Linux by a program run with \env{LD\_LIBRARY\_PATH}
  naming the library directory, \file{\var{prefix}/lib} (or \file{lib64} or
  \file{lib/\var{multiarch}} on the systems whose libraries go there), or linked
  with \option{-Wl,-rpath,} and that directory.
  With \cmd{configure} the static library is built too unless
  \loption{disable-static}, as with \option{-Ddefault\_library=both} of
  meson; \cmd{configure} and meson built both by default in
  GAOL~4.\newinvfive
\item[$\bullet$] \textbf{\option{CMAKE\_POSITION\_INDEPENDENT\_CODE}.}
  \code{ON} by default: GAOL is compiled as position-independent code
  (\option{-fPIC}), so that it can be linked into a shared library, as
  Python bindings; a project building GAOL with \code{FetchContent} that sets
  it is followed. The \cmd{configure} of GAOL~4 compiled the static library
  without it.\newinvfive
\item[$\bullet$] \textbf{\option{CMAKE\_OSX\_ARCHITECTURES}.} A universal
  build of macOS, one configuration compiled for each architecture: SSE2
  intervals for x86\_64 and NEON intervals for arm64, where
  \option{GAOL\_SIMD} allows them, and no fused multiply-add instructions nor
  AVX-512 path. CMake only is checked on a Mac of ours; \cmd{configure} and
  meson are checked by the continuous integration alone. The generator of
  Xcode, which builds arm64 and x86\_64 by itself when no list is given, is
  given the architecture of the machine, as the other generators build for
  it; the list \code{arm64;x86\_64} asks for the universal build.\newinvfive
\item[$\bullet$] \textbf{\option{WITH\_TESTS}.} \code{OFF} by default: builds
  the unit tests of \file{tests}, which \cmd{make test} and \cmd{make check}
  run (\pxref{sec:tests}).
\item[$\bullet$] \textbf{\option{WITH\_EXAMPLES}.} \code{OFF} by default:
  builds the examples of \file{examples}, which \cmd{make check} runs.
\item[$\bullet$] \textbf{\option{WITH\_DOC}.} \code{OFF} by default: adds the
  target \cmd{pdf}, which builds this manual in
  \file{manual/Manual\_GAOLv5.pdf} with pdflatex, bibtex and makeindex
  (Table~\ref{tab:install-targets}).\newinvfive
\item[$\bullet$] \textbf{\option{GAOL\_SIMD}.} \code{ON} by default: the
  intervals are computed with the vector instructions of the processor as
  soon as it has them, their two bounds in one register: SSE2 instructions on
  x86 processors (\option{-msse2 -msse3}; no flag with Visual C++ and
  clang-cl for x64; on 32-bit Windows with GCC and Clang and
  \option{-faligned-new=8}, which the code using GAOL is given too, and with
  Visual C++ and clang-cl, whose intervals keep their two bounds as two
  doubles and need no aligned memory), and NEON instructions on 64-bit ARM
  processors, with every compiler and no flag (\pxref{sec:neon}); the FPU
  computes them otherwise, and on any other processor (32-bit ARM, POWER,
  s390x, RISC-V\ldots). A universal build of macOS takes the NEON intervals
  for its arm64 part and the SSE2 intervals for its x86\_64 part.
  \option{GAOL\_ARM64}, \loption{enable-arm64} and \code{enable-arm64} were
  an option of the NEON intervals, independent of \option{GAOL\_SIMD}, and
  are gone.\newinvfive
\cindex{SSE2}%
\cindex{NEON}%
\item[$\bullet$] \textbf{\option{GAOL\_PRESERVE\_ROUNDING}.} \code{OFF} by
  default: each operation restores the rounding direction it found, rather
  than leaving it upward (\pxref{sec:preserve-rounding}).%
\cindex{rounding!preserving}%
\item[$\bullet$] \textbf{\option{GAOL\_COVERAGE}.} \code{OFF} by default:
  compiles GAOL and its tests with the counters of gcov, and adds the target
  \code{coverage}, which runs the tests and writes the report of their
  coverage, \file{coverage/coverage.html}, and a page of conclusions
  (\cmd{cmake --build build --target coverage}, \cmd{make coverage} of
  \cmd{configure}, \cmd{ninja -C build gaol-coverage} of meson, which keeps
  \code{coverage} for its own report). GCC and Clang, with the \code{Debug}
  build, \cmd{gcovr} and Python~3; the tests have to be built.
\cindex{coverage}%
\end{itemize}

\subsection{What the configuration decides}\label{sec:install-decides}

The rest of the configuration has no option. It is decided by checks of the
compiler and of the machine, the same with the three builds, each told in the
messages of \cmd{cmake}, \cmd{configure} and \cmd{meson setup}. They assume
that the compiler, the architecture and the system that configure and compile
GAOL are those of the code using it: a program compiled and run by the
configuration tells what the machine does, and the headers and \file{gaol.pc}
carry the flags of the code using GAOL.\newinvfive

\begin{itemize}
\item[$\bullet$] \textbf{The fused multiply-add instructions} (was
  \option{GAOL\_FMA}, \loption{enable-fma}, \code{enable-fma}). GAOL and
  CORE-MATH are compiled with them whenever the C++ and the C compiler take
  their flag and the machine building runs a program compiled with them: a
  program that runs the instruction itself, \code{\_\_builtin\_fma()}, and
  the intrinsic \code{\_mm\_fmadd\_sd()} with Visual C++ and clang-cl for
  x64, which a processor without FMA3 stops on. The flags are
  \option{-mfma} on x86 (\option{-mfma -mprefer-vector-width=128} with GCC
  for Windows, on x86-64 and on 32-bit x86), \option{/arch:AVX2} with Visual
  C++ for x64, \option{-mfpu=neon-vfpv4 -mfloat-abi=hard} on 32-bit ARM;
  64-bit ARM, POWER, s390x and RISC-V have the instructions without a flag.
  The library then needs a processor with them (on x86, Intel Haswell and AMD
  Piledriver, 2012--2013, and later; with Visual C++ and clang-cl, whose \option{/arch:AVX2}
  also needs AVX2, Intel Haswell and AMD Excavator, 2013--2015, and later), and the code using GAOL is given the
  flag too, in \code{gaol::gaol} and \file{gaol.pc}; \option{-ffp-contract=off}
  stays. The elementary functions are much faster with them: \code{cos} 2.5
  times, \code{log} 1.6 times, \code{sin} 1.3 times (Intel i7-1185G7,
  Clang~21.1), and with GCC~16.2 for Windows \code{cos} 4.7 times, \code{log}
  3.2 times and \code{sin} 2.0 times (GitHub runner, MSYS2 UCRT64). GCC for
  Windows does not align its stack on 32 bytes for the 256-bit registers that
  \option{-mfma} alone would have it use (GCC bug 54412):
  \option{-mprefer-vector-width=128} keeps it to the 128-bit ones, and the
  objects of GAOL built so by GCC~16.2 have no instruction on a 256-bit
  register (\pxref{sec:install-gcc-windows}). In a cross-compilation, where
  no program can be run, they are left out, with a message. A universal build
  of macOS leaves the flag out, its arm64 part refusing \option{-mfma} while
  it has the instructions without a flag: the x86\_64 part is then built
  without them, runs on any x86\_64 processor, Rosetta 2 included, but
  CORE-MATH and the exact products of GAOL call the \code{fma()} of the
  system library, and the integer powers beyond 1 and $-1$ do not take the
  packed kernel that uses the instruction; a build for x86\_64 alone
  (\option{CMAKE\_OSX\_ARCHITECTURES=x86\_64}) takes \option{-mfma} where the
  machine building runs a program compiled with it.
\cindex{fused multiply-add}%
\item[$\bullet$] \textbf{The AVX-512 path} (was
  \option{GAOL\_PREFER\_AVX512}, \loption{enable-prefer-avx512},
  \code{enable-prefer-avx512}). The four arithmetic operations, the square,
  the square root and, with FMA, the integer powers beyond 1 and $-1$ take
  the AVX-512 instructions, which carry their rounding direction and set
  none, where \option{GAOL\_PRESERVE\_ROUNDING} is on, the compiler takes the
  instructions (a program that uses them, linked) and the machine building
  runs them; otherwise the SSE2 intervals compute them. The rule comes from
  measurements (\pxref{sec:avx512}): the path spares the store of the
  rounding direction that every operation pays with
  \option{GAOL\_PRESERVE\_ROUNDING}, where it gains 2.9 times at the median
  of seven operations, on an Intel i7-1185G7 (no other processor was
  measured); without that option the operations take 20~\% longer at the
  median, the direction being written only when it is not upward already, so
  the path is left out. The library
  then needs a processor with the instructions, as with the fused
  multiply-add instructions, and \code{gaol::init()} stops a program run on
  another one with a message (\pxref{sec:avx512}). In a cross-compilation,
  the path is left out, with a message.
\cindex{AVX-512}%
\item[$\bullet$] \textbf{The 128-bit integer} of the accurate phases of
  CORE-MATH (was \option{GAOL\_U128\_EMULATION}): the type of the compiler
  where it has one, else two 64-bit halves (Visual C++, and GCC for a 32-bit
  target), which \file{gaol/math/gaol\_u128.h} takes by itself and the
  configuration says.
\item[$\bullet$] \textbf{Clang with the target of Visual C++}
  (\cmd{clang++} for \code{x86\_64-pc-\allowbreak windows-msvc} or
  \code{i686-pc-windows-msvc}, with the driver of GCC, not clang-cl). It
  defines \code{\_MSC\_VER}, which the three builds read
  (\code{CMAKE\_CXX\_SIMULATE\_ID} of CMake) to compile GAOL in C++14, the
  standard library of Visual C++ needing it, with the \file{unistd.h} of
  \file{cmake/msvc} for the lexer and without the warnings of the C runtime,
  and, with CMake and meson, to build the static library alone, as with Visual
  C++ (a DLL would export nothing); \cmd{configure} builds the static library
  by default.\newinvfive
\item[$\bullet$] \textbf{The debugging} (was \option{GAOL\_DEBUG},
  \code{GAOL\_DEBUGGING}). GAOL and CORE-MATH are built for debugging,
  without optimization nor \code{NDEBUG}, with debugging information and
  \code{GAOL\_DEBUGGING} (\pxref{sec:debugging-facilities}), for the
  \code{Debug} build type (\option{CMAKE\_BUILD\_TYPE},
  \loption{enable-debug}, \loption{buildtype=debug}); there is no target
  \cmd{perf} in the \code{Debug} build of CMake, nor with
  \option{-Denable-debug=true} of meson (\cmd{configure --enable-debug} and
  \loption{buildtype=debug} of meson keep it). Each configuration keeps its other flags (the C runtime
  of Visual C++, sanitizers\ldots); the tests and the examples, and the code
  of a project bringing GAOL in, keep their build type.
\item[$\bullet$] \textbf{The memory and the flags of the intervals}: the
  aligned memory the SSE2 intervals need on 32-bit Windows, the flags of
  interval arithmetic, the link option \option{-mno-daz-ftz},
  \file{<fenv.h>}, the visibility, the type of the intervals of a universal
  build.
\end{itemize}

\subsection{The settings of the developers}\label{sec:install-developers}

The checks above decide by themselves. A developer who wants to run a path
the machine would not take, or a user who builds a package for any processor,
gives one of the settings of Table~\ref{tab:install-developers}. They are not
options of the user: \cmd{configure --help} lists none of them but
\loption{enable-debug}, which is also the \code{Debug} build of \cmd{configure}, and
\cmd{meson configure} lists them with the mark \emph{developers}, since meson
cannot leave them out. CMake reads them where they apply, without having them
in the cache unless they were given. By default, the three builds report errors by
exceptions (\pxref{sec:exceptions}); one setting of the table makes GAOL abort instead.

\begingroup\renewcommand{\arraystretch}{1.1}\footnotesize
\setlength{\tabcolsep}{4pt}
\begin{longtable}{L{3.5cm}L{3.2cm}L{3.3cm}L{3.8cm}}
  \caption{The settings of the developers}
  \label{tab:install-developers}\\
  \hline
  \multicolumn{1}{c}{CMake} & \multicolumn{1}{c}{\cmd{configure}} &
  \multicolumn{1}{c}{meson} & \multicolumn{1}{c}{Meaning}\\
  \hline\hline
  \endfirsthead
  \multicolumn{4}{l}{\textit{Table~\thetable, continued}}\\*
  \hline
  \multicolumn{1}{c}{CMake} & \multicolumn{1}{c}{\cmd{configure}} &
  \multicolumn{1}{c}{meson} & \multicolumn{1}{c}{Meaning}\\
  \hline\hline
  \endhead
  \hline
  \endfoot
    \option{-DGAOL\_FMA=OFF} & \loption{disable-fma} &
    \option{-Denable-fma=false} & No fused multiply-add instructions: a
    library for any processor of the architecture\\
    \option{-DGAOL\_PREFER\_\allowbreak AVX512=ON}, \code{=OFF} &
    \loption{enable-prefer-\allowbreak avx512}, \loption{disable-prefer-\allowbreak avx512} &
    \option{-Denable-prefer-\allowbreak avx512=true}, \code{=false} & The AVX-512 path
    asked for without \option{GAOL\_PRESERVE\_ROUNDING}, or refused with it
    (a package for any processor); asked for where the compiler cannot take
    it, an error\\
    \option{-DGAOL\_U128\_\allowbreak EMULATION=ON} &
    \code{CPPFLAGS=-DGAOL\_U128\_\allowbreak FORCE\_EMULATION=1} &
    \code{-Dc\_args=\ldots} and \code{-Dcpp\_args=\ldots} with the same macro
    & The two 64-bit halves where the compiler has a 128-bit type: the code of
    Visual C++ and of the 32-bit targets on an ordinary machine\\
    \option{-DGAOL\_DEBUG=ON} & \loption{enable-debug} &
    \option{-Denable-\allowbreak debug=true} & GAOL and CORE-MATH built for debugging
    whatever the build type, a project bringing GAOL in included
    (\option{-Dgaol:enable-\allowbreak debug=true} for a subproject of meson); the tests
    and the examples keep the build type\\
    \option{-DGAOL\_EXCEPTIONS\_\allowbreak ENABLED=OFF} &
    \loption{disable-\allowbreak exceptions} & \option{-Denable-\allowbreak exception=false} & GAOL
    aborts to signal an error instead of raising an exception\\
    \option{-DGAOL\_VERBOSE\_MODE=ON} & \loption{enable-verbose-\allowbreak mode} &
    \option{-Denable-verbose-\allowbreak mode=true} & A line on the standard error when
    GAOL initializes and cleans up; GAOL is silent by default\\
    \option{-DGAOL\_SSE2\_UNALIGNED\_\allowbreak STORAGE=ON} &
    \code{CXXFLAGS=-DGAOL\_SSE2\_\allowbreak UNALIGNED\_STORAGE=1} &
    \code{-Dcpp\_args=\ldots} with the same macro & The SSE2 intervals with
    their two bounds as two doubles, as on 32-bit Windows with Visual C++ and
    clang-cl: that layout on another system. CMake alone writes the macro into the
    installed \file{gaol/gaol\_configuration.h}: given to \cmd{configure} or
    meson by the flags, it reaches what the build compiles, not a program
    compiled later with the installed GAOL, which then has to be compiled with
    the same flag, or its layout differs from the library's; with CMake, a
    program must not define the macro\\
    \code{-DCMAKE\_CXX\_FLAGS=\allowbreak -DGAOL\_RELATION\_\allowbreak
    BOUNDS\_FIRST=0} & \code{CXXFLAGS=-DGAOL\_RELATION\_\allowbreak
    BOUNDS\_FIRST=0} & \code{-Dcpp\_args=\ldots} with the same macro & The
    relations between intervals (\code{<}, \code{<=}, \code{>}, \code{>=})
    test the empty sets before they compare the bounds, as they do on
    RISC-V; \code{=1}, the default elsewhere, compares the bounds first. The
    answers are the same; the order of the tests changes, which a
    processor may run faster or slower, and which decided whether GCC~14 raised
    the invalid-operation flag for an empty operand on RISC-V (the reason for
    its default there). It is to be defined the same in every file that
    includes GAOL (\file{gaol/gaol\_port.h})\\
    \option{-DGAOL\_MACHINE\_RUNS\_FMA=1},
    \option{-DGAOL\_MACHINE\_RUNS\_\allowbreak AVX512=1} &
    \code{gaol\_cv\_machine\_runs\_fma=yes},
    \code{gaol\_cv\_machine\_runs\_\allowbreak avx512=yes} &
    \option{-Dmachine-runs-\allowbreak fma=true}, \option{-Dmachine-runs-\allowbreak avx512=true} &
    In a cross-compilation, or under an emulator, where no program can be
    run: the machine running GAOL has the fused multiply-add instructions or
    runs the AVX-512 ones\\
\end{longtable}
\endgroup

\begin{Rem}
An option that GAOL had and no longer has is refused or warned of: meson
refuses \option{-Denable-arm64}, CMake and \cmd{configure} warn that
\option{GAOL\_ARM64=OFF} and \loption{disable-arm64} are no longer read, and
say \option{GAOL\_SIMD=OFF} and \loption{disable-simd}. A build directory of
CMake configured by an earlier GAOL v5 holds the entries of the options that
are gone (\option{GAOL\_FMA}, \option{GAOL\_DEBUG}\ldots), which stay
harmless, and \option{GAOL\_PREFER\_AVX512}, which the older CMake kept with
the value it had derived: the record \code{GAOL\_PREFER\_AVX512\_DERIVED}
tells it, and both entries go.
\end{Rem}

\subsection{The archive of the sources and the packages}\label{sec:packages}
\cindex{package}%

CPack makes them from the CMake build, in the directory of the build, where
\cmd{make dist} of the autotools build made the archive of the sources:

\begin{cmdshell}
% cmake -S . -B build -DCMAKE_INSTALL_PREFIX=/usr \
        -DGAOL_FMA=OFF -DGAOL_PREFER_AVX512=OFF
% cmake --build build --config Release
% cmake --build build --target package_source
% cmake --build build --target package
\end{cmdshell}

\begin{itemize}
\item[$\bullet$] {\sloppy \textbf{\code{package\_\allowbreak source}} gives
  \file{gaol-\var{version}\allowbreak .tar.gz}: the sources, in
  \file{gaol-\var{version}/}, with \file{configure} and the
  \file{Makefile.in}, which the three builds build. CPack takes the source
  tree as it is, less the files of \code{CPACK\_SOURCE\_IGNORE\_FILES} (the
  repository of git, the builds made within it, the objects) and the
  notes of the work on GAOL v5: made from a clean checkout, the archive holds
  the other files of the commit, \file{configure} as it was generated among
  them. If that was for another version than \file{VERSION.txt} holds,
  \cmd{configure --version} in the archive gives the old one, and CMake warns
  of it when it configures (\pxref{sec:install-version}). The archive takes
  its name from the version CMake read when it configured. After a change of
  \file{VERSION.txt}, \code{package\_source} configures the build directory
  again first, whatever the generator: Ninja does before it runs CPack, and
  CPack does when \file{VERSION.txt} holds another version than the build
  directory was configured for (\file{cmake/gaol\_package\_source.cmake}),
  where the Makefile generators ran CPack at once.\newinvfive\par}
\item[$\bullet$] \textbf{\code{package}} gives
  \file{gaol-\var{version}-\var{system}.tar.gz}, what \cmd{cmake --install}
  installs, in \file{gaol-\var{version}-\var{system}/}: the headers,
  \file{libgaol.a} (or \file{libgaol.so*} with \option{BUILD\_SHARED\_LIBS}),
  the CMake package of GAOL and \file{gaol.pc}, to extract anywhere,
  \file{gaol.pc} finding its directories from where it is. On Linux it gives
  \file{libgaol-dev\_\allowbreak\var{version}\_\allowbreak\var{arch}.deb} too,
  the same files in \file{/usr}, for \cmd{apt install
  ./libgaol-dev\_\var{version}\_\var{arch}.deb}: by default a static library
  and its headers, hence the name and the section \code{libdevel} of the
  development packages of Debian. \option{-DCMAKE\_INSTALL\_PREFIX=/usr}
  gives the library directory of Debian and Ubuntu
  (\file{lib/x86\_64-linux-gnu}); with another prefix, the \file{.deb} puts
  the files in \file{/usr/lib}.
\end{itemize}

\noindent
The packages hold \file{libgaol} as the build compiled it: where the machine
building has the fused multiply-add instructions, only the processors that
have them run it, and the code using it is compiled with \option{-mfma};
\option{-DGAOL\_FMA=OFF} makes packages for any processor of the
architecture that has SSE3 on x86 (\option{-msse3} is given wherever the
compiler takes it). So does \option{-DGAOL\_PREFER\_AVX512=OFF} where
\option{GAOL\_PRESERVE\_ROUNDING} is \code{ON}, whose AVX-512 path makes a
library that only the processors with the AVX-512 instructions run, when the
machine building has them (Visual C++ and clang-cl for x64 too).\newinvfive

\begin{Rem}
On Cygwin, \code{package\_source} did not work in the job of the continuous
integration of 23 October 2026 (CMake 4.4.4, libarchive 3.8): CPack, with
TGZ, TBZ2 and ZIP alike, failed to add \file{README}, a symbolic link to
\file{README.md} in the repository, with ``Couldn't list extended
attributes'', so that the target fails there and the test
\code{cpack\_stale\_configure} says it is skipped. The cause was not found:
\cmd{cmake -E tar}, GNU \cmd{tar} and \cmd{git archive} make an archive of
the checked-out tree, and \code{listxattr()}, \code{llistxattr()} and
\code{flistxattr()} find nothing wrong with its link. \cmd{git archive} makes
the archive from the repository.
\end{Rem}

\section{Building and installing with autotools}\label{sec:configuration}
\findex{configure}%

In the root directory of the sources (for a build outside the sources,
\cmd{configure} is called by its path from another directory):

\begin{cmdshell}
% ./configure --prefix=@var^prefix~ --with-tests
% make
% make test
% make install
\end{cmdshell}

\noindent
\cmd{configure} compiles the thirty-eight sources of CORE-MATH into
\file{libgaol} (\file{3rd/Makefile.am}, whose convenience library
\file{gaol/Makefile.am} puts into \file{libgaol}), with the flags it gives to the C
code of GAOL: there is nothing else to build, nothing else to install and
nothing else to link. There is no \cmd{make dist} nor \cmd{make distcheck}:
the archive of the sources is made with CMake (\pxref{sec:packages}).

\medskip
\file{configure}, \file{aclocal.m4}, the \file{Makefile.in},
\file{ltmain.sh}, \file{m4/*.m4} and the scripts \file{compile},
\file{config.guess}, \file{config.sub}, \file{depcomp}, \file{install-sh},
\file{missing} and \file{test-driver} are committed, so that a checkout
builds without the autotools, and \cmd{make} never regenerates them
(\code{AM\_MAINTAINER\_MODE([disable])}). They are generated by autoconf
2.72, automake 1.18.1 and libtool 2.5.4, with \cmd{AUTOHEADER=true autoreconf
--force --install} followed by \cmd{git checkout INSTALL}, which
\loption{install} replaces by the text of automake; \cmd{AUTOHEADER=true}
keeps \file{gaol/gaol\_configuration.h.in}, written by hand, which autoheader
2.72 overwrites (2.73 leaves it). \file{config.guess} and \file{config.sub}
are those of 2026-05-17 from GNU config
(\url{https://git.savannah.gnu.org/cgit/config.git}), newer than the ones
libtool installs. The versions matter: the files of GAOL 4.2.2 came from
libtool 2.4.2, which linked \file{libgaol} with \option{-flat\_namespace
-undefined suppress} on macOS 11 and later, taking them for Mac OS X 10.0,
and from a \file{config.sub} of 2016, which refused
\code{loongarch64}.\newinvfive{} The file \file{INSTALL} gives the four
commands of each of the three builds, and points to \file{doc/building.md} and
\file{doc/using.md} for the options and the flags.

\medskip
\cmd{configure} accepts the following options, which are those of the three
builds (Table~\ref{tab:install-options}):

\begin{itemize}
\item[$\bullet$] \textbf{\loption{help}.} Displays a list of all options; only
  those described below are options of the user, with the options of libtool and
  automake; the settings of the developers are not listed, but
  \loption{enable-debug}, which is also the \code{Debug} build.
\cindex{help!on configuration}%
\item[$\bullet$] \textbf{\loption{prefix=\var{prefix-dir}}.} The root directory
  where the library is installed. It defaults to \file{/usr/local}.
\item[$\bullet$] \textbf{\loption{libdir=\var{lib-dir}}.} The directory where
  to put the libraries. It defaults to \file{\var{prefix-dir}/lib}.
\item[$\bullet$] \textbf{\loption{includedir=\var{include-dir}}.} The
  directory where to put the header files. It defaults to
  \file{\var{prefix-dir}/include}.
\item[$\bullet$] \textbf{\loption{enable-debug[=\var{yes/no}]}.} Builds GAOL
  as the \code{Debug} build of CMake: with \option{-g}, without optimization
  nor \code{NDEBUG}, and with the debugging facilities
  (\pxref{sec:debugging-facilities}); the \code{Release} build of CMake
  otherwise, with \option{-O3} and \code{NDEBUG}, and the flags
  \option{-funroll-loops}, \option{-fomit-frame-pointer} and
  \option{-fexpensive-optimizations} where the compiler takes them, and
  \option{-falign-functions=64} for the library alone. This option defaults to
  \code{no}.
\cindex{debugging!option}%
\cindex{optimization!configuration option}%
\item[$\bullet$] \textbf{\loption{enable-simd[=\var{yes/no}]}.} The vector
  intervals, SSE2 on x86 processors and NEON on 64-bit ARM processors, as
  \option{GAOL\_SIMD} (\pxref{sec:cmake-options}). This option defaults to
  \code{yes}.\newinvfive
\item[$\bullet$] \textbf{\loption{enable-preserve-rounding[=\var{yes/no}]}.}
  Each operation restores the rounding direction it found
  (\pxref{sec:rounding-direction}). This option defaults to \code{no};\newinvfive{}
  the \cmd{configure} of GAOL~4 had it on by default, although its manual says
  the contrary.
\cindex{rounding!preserving}%
\item[$\bullet$] \textbf{\loption{enable-coverage[=\var{yes/no}]}.} Compiles
  GAOL and its tests with the counters of gcov, and adds the target
  \cmd{make coverage}, as \option{GAOL\_COVERAGE}. This option defaults to
  \code{no}.\newinvfive
\item[$\bullet$] \textbf{\loption{enable-shared[=\var{yes/no}]},
  \loption{enable-static[=\var{yes/no}]}.} The static library
  \file{libgaol.a}, which is all that is built by default, as with CMake;
  \loption{enable-shared} builds the shared library \file{libgaol.so.5.0.0}
  too, whose soname is \file{libgaol.so.5}, as CMake
  (\option{BUILD\_SHARED\_LIBS}) and meson name it (it was
  \file{libgaol-5.0.so.0} in the earlier versions of GAOL v5, and it was built
  by default before them); \loption{enable-pic} (the default, though
  \cmd{configure --help} says \emph{use both}) compiles the code as
  position-independent code.\newinvfive
\item[$\bullet$] \textbf{\loption{with-tests[=\var{yes/no}]}.} Builds the unit
  tests of \file{tests}, which \cmd{make test} and \cmd{make check} run, as
  \option{WITH\_TESTS}. This option defaults to \code{no}.
\item[$\bullet$] \textbf{\loption{with-examples[=\var{yes/no}]}.} Builds the
  examples of \file{examples}, which \cmd{make check} runs, as
  \option{WITH\_EXAMPLES}. This option defaults to \code{no}.
\end{itemize}

\noindent
{\raggedright % the names of the options, in typewriter, do not stretch
As usual, \loption{disable-\var{xxx}} is equivalent to
\loption{enable-\var{xxx}=no}, and \loption{enable-\var{xxx}} to
\loption{enable-\var{xxx}=yes}. The fused multiply-add instructions, the
AVX-512 path, the 128-bit integer and the NEON intervals are decided by checks
of the compiler and of the machine (Section~\ref{sec:install-decides}).
\loption{enable-optimize} (the \code{Release} build, which is the default,
and \option{-O} with \loption{disable-optimize}) and \loption{enable-arm64}
(the NEON intervals follow \loption{enable-simd}) are no longer options; the
\loption{enable-verbose-mode} and \loption{enable-exceptions} of GAOL~4 are
settings of the developers, as \loption{enable-fma} and
\loption{enable-prefer-avx512}, which \cmd{configure --help} does not list
(Section~\ref{sec:install-developers}); the other options of GAOL~4,
\loption{with-mathlib}, \loption{with-mathlib-include},
\loption{with-mathlib-lib}, \loption{with-cppunit-include},
\loption{with-cppunit-lib}, \loption{enable-fast-math}, \loption{enable-asm}
and \loption{enable-relations}, are gone, and \cmd{configure} warns that it
does not know them.\newinvfive{}\par}

\medskip
A universal build of
macOS is configured with \option{-arch} several times in the flags of the
compilers: \code{./configure CFLAGS="-arch arm64 -arch x86\_64"
CXXFLAGS="-arch arm64 -arch x86\_64" LDFLAGS="-arch arm64 -arch x86\_64"};
the dependency tracking of automake, which the compilers do not take with
several \option{-arch}, is turned off then.\newinvfive

\medskip
The targets of the \file{Makefile} in the root directory are \code{all},
which builds \file{libgaol.a} (and the shared library \file{libgaol.so.5.0.0}
with \loption{enable-shared}), \code{test}, \code{check}, \code{perf} and
\code{parser}, those of the three builds (\pxref{sec:tests}),
\code{install}, \code{clean}, \code{distclean} and, with
\loption{enable-coverage}, \cmd{coverage}. \code{clean} erases what
\cmd{make} built, and keeps what \cmd{configure} made, so that \cmd{make}
builds again without it: the \file{Makefile}, \file{config.status},
\file{config.log}, \file{libtool}, \file{gaol.pc},
\file{gaol/gaol\_configuration.h}, the \file{manual/v*/gaol\_version.tex},
the \file{.deps} directories of the dependencies and the \file{.dirstamp} of
the directories of CORE-MATH. \code{distclean} erases them too, the
\file{.deps} directories included, which it left empty in GAOL~4, and gives
the source tree back as git has it, \cmd{configure} to be run again; the
continuous integration checks it.\newinvfive{} \cmd{make -C manual pdf}
builds this manual; no build makes the manual of GAOL~4 again.
\cmd{configure --help} ends with
the address for the bug reports, \code{jordan.ninin@ensta.fr}, and the page of
GAOL v5, \code{https://github.com/Jordan08/GAOL}.

\section{Building and installing with meson}\label{sec:install-meson}
\cindex{meson}%

\begin{cmdshell}
% meson setup build --prefix=@var^prefix~ -Dwith-tests=true
% meson compile -C build
% meson test -C build
% meson install -C build
\end{cmdshell}

\noindent
meson compiles the thirty-eight sources of CORE-MATH into \file{libgaol}
(\file{gaol/meson.build}), as CMake and \cmd{configure} do, and the lexer and
the parser committed with the sources: neither flex nor bison is needed.
\cmd{meson compile} needs meson 0.54; with an older one, as the meson 0.53 of
Ubuntu 20.04, \cmd{ninja -C build} builds GAOL as well. It builds
\file{libgaol.a}, and \file{libgaol.so.5.0.0} with
\option{-Ddefault\_library=shared} (\code{both} for the two), whose soname is
\file{libgaol.so.5}, as \cmd{configure} with \loption{enable-shared} and
CMake name it; with Visual C++, the static library alone, as CMake, and
\code{gaol\_dep} and \file{gaol.pc} no longer give the code using GAOL
\code{\_\_GAOL\_PUBLIC\_\_} defined empty (\pxref{sec:compiling}).\newinvfive{}
The options are given to \cmd{meson setup}, as
\option{-D\var{option}=\var{value}}; the build type gives the optimizations
and the debugging facilities, as with CMake (\loption{buildtype=debug} for
the \code{Debug} build), and the others are those of \cmd{configure}
(Table~\ref{tab:install-meson}).

\begin{table}[htbp]
  \centering\renewcommand{\arraystretch}{1.1}\small
  \caption{The options of the meson build} \label{tab:install-meson}
  \begin{tabular}{lL{8.6cm}}
    \hline
    \multicolumn{1}{c}{Option} & \multicolumn{1}{c}{Default}\\
    \hline\hline
    \code{buildtype} & \code{release}: \option{-O3}, \code{NDEBUG} and the
    optimizations of the \code{Release} build of CMake and \cmd{configure};
    \code{debug} for the \code{Debug} build, without optimization, with
    debugging information and without \code{NDEBUG}. The option
    \code{enable-optimize} is gone. Brought in as a subproject, GAOL is built
    in \code{release} too where the project leaves the build type to its
    default\\
    \code{enable-simd} & \code{true}: the vector intervals, SSE2 or NEON, as
    \option{GAOL\_SIMD}\\
    \code{enable-preserve-rounding} & \code{false}, as
    \option{GAOL\_PRESERVE\_ROUNDING}\\
    \code{enable-coverage} & \code{false}; the target
    \cmd{ninja gaol-coverage}, as \option{GAOL\_COVERAGE}\\
    \code{default\_library} & \code{static}; \code{shared} builds the shared
    library, \code{both} the two, as \option{BUILD\_SHARED\_LIBS} (it was
    \code{both})\\
    \code{b\_staticpic} & \code{true}: position-independent code, as
    \option{CMAKE\_POSITION\_INDEPENDENT\_CODE}\\
    \code{with-tests} & \code{false}; builds the unit tests, which
    \cmd{meson test} and \cmd{ninja check} run (it was \code{with-test})\\
    \code{with-examples} & \code{false}; builds the examples, which
    \cmd{ninja check} runs\\
    \code{with-doc} & \code{false}; the target \cmd{pdf} then builds this
    manual, \file{manual/Manual\_GAOLv5.pdf}, with pdflatex, bibtex and
    makeindex, as \option{WITH\_DOC}\\
    \hline
  \end{tabular}

\end{table}

\noindent
These are the options of the three builds. \code{enable-fma},
\code{enable-prefer-avx512}, \code{enable-debug}, \code{enable-exception},
\code{enable-verbose-mode}, \code{machine-runs-fma} and
\code{machine-runs-avx512} are settings of the developers, which
\cmd{meson configure} lists with the mark \emph{developers}
(Section~\ref{sec:install-developers}). The options \code{enable-arm64} (the
NEON intervals follow \code{enable-simd}), \code{enable-optimize},
\code{enable-relations}, \code{with-test} and \code{check-perf} are gone, and
meson refuses them.\newinvfive{} A universal build of
macOS is configured with \option{-arch} several times in \code{c\_args},
\code{cpp\_args}, \code{c\_link\_args} and \code{cpp\_link\_args}. The
summary of \cmd{meson setup} gives the address for the bug reports,
\code{jordan.ninin@ensta.fr}, and the page of GAOL v5,
\code{https://github.com/Jordan08/GAOL}.

\medskip
\cmd{meson setup} runs Python once, in \code{project()}, to read the version
of GAOL in \file{VERSION.txt} (\pxref{sec:install-version}), and twice more
for the message of a file it refuses: the first of \cmd{python3} and
\cmd{python} that it finds in the \env{PATH} or, when it finds neither, the
Python that runs meson (the case of the \file{meson.exe} of the Windows
installer when no Python is installed). \file{meson.build} names
\cmd{python3} first because meson falls back on its own Python for that name
alone; the continuous integration checks it on Linux with a \cmd{meson setup}
whose \env{PATH} holds no Python.

\begin{Rem}
On Windows, none of what follows was checked, but read in the sources of
meson. The directory \file{WindowsApps} of the profile of the user
(\file{\%USERPROFILE\%\textbackslash{}AppData\textbackslash{}Local%
\textbackslash{}Microsoft\textbackslash{}WindowsApps})
can hold aliases of Python, \file{python3.exe} and perhaps \file{python.exe}
too, which only open the Microsoft Store when Python was not installed from
it. meson 0.53.1 and later (the 0.53.2 of Ubuntu 20.04 and the meson of pip
among them) leave that directory out of their search for programs, as long as
the \env{PATH} names it by that path, not for a profile whose directory
differs from \env{USERPROFILE}. meson 0.53.0 and earlier may take such an
alias, and \cmd{meson setup} then stops on the failure of the command that
reads \file{VERSION.txt}: the remedy is a later meson
(\cmd{pip install meson}), or to turn off the aliases of Python in the
Windows settings (\emph{Manage app execution aliases}).
\end{Rem}

\section{What the options are worth}\label{sec:install-options-worth}

The options of Table~\ref{tab:install-options} that change the code of the
operations are the vector intervals (\option{GAOL\_SIMD}) and the preserved
rounding direction (\option{GAOL\_PRESERVE\_ROUNDING}); the fused
multiply-add instructions, which the configuration decides, change it too.
Table~\ref{tab:install-worth-x86} measures the vector intervals and the fused
multiply-add instructions with
\file{tests/performance.cpp} on an Intel i7-1185G7 with Clang 21.1, in
nanoseconds per operation, the best of nine runs, each program being run in
turn with the others.

\begin{table}[htbp]
  \centering\renewcommand{\arraystretch}{1.1}\footnotesize
  \setlength{\tabcolsep}{4pt}
  \caption{The options on x86: nanoseconds per operation, Intel i7-1185G7, Clang 21.1, the best of nine runs} \label{tab:install-worth-x86}
  \begin{tabular}{L{3.3cm}rrrrrrrr}
    \hline
    Build & $x + y$ & $x \times y$ & \code{sqr} & \code{pow(x, 3)} &
    \code{exp} & \code{log} & \code{sin} & \code{cos}\\
    \hline\hline
    Default (FMA and \option{GAOL\_SIMD} on) & 1.07 & 2.04 & 1.80 & 5.42 &
    18.6 & 21.7 & 31.3 & 36.6\\
    \option{-DGAOL\_FMA=OFF} & 1.14 & 2.06 & 1.81 & 17.4 & 21.6 & 34.5 &
    41.4 & 90.3\\
    \option{-DGAOL\_SIMD=OFF} & 2.30 & 4.32 & 2.46 & 7.65 & 19.9 & 23.6 &
    32.8 & 39.0\\
    \hline
  \end{tabular}

\end{table}

\begin{itemize}
\item[$\bullet$] \textbf{The fused multiply-add instructions pay on the elementary
  functions.} \code{cos} is 2.5 times faster with them,
  \code{pow(x, 3)} 3.2 times, \code{log} 1.6 times, \code{sin} 1.3 times and
  \code{exp} 1.2 times. CORE-MATH computes with \code{\_\_builtin\_fma}
  throughout, which is one instruction where the processor has it and a call
  to the math library where it has not. They are on by default wherever the
  compiler has a flag for them and the machine building runs a program
  compiled with it. They do nothing for the arithmetic, which does not
  multiply and add in one step.\newinvfive
\item[$\bullet$] \textbf{The SSE2 intervals pay on the products and the
  additions.} $x \times y$ takes 2.04~ns rather than 4.32, and $x + y$ 1.07
  rather than 2.30, the two bounds being computed in one instruction. They
  are on by default on x86 processors. The elementary functions do not go
  through them.
\end{itemize}

\medskip
On a 64-bit ARM processor, Table~\ref{tab:install-worth-arm} measures them
with \file{tests/\allowbreak performance.cpp} on an Apple M2 with AppleClang 21.0, the median of nine runs of
each program, 6 November 2026: the first two builds in turn in one session,
the last two in another.

\begin{table}[htbp]
  \centering\renewcommand{\arraystretch}{1.1}\footnotesize
  \caption[The options on 64-bit ARM (Apple M2, AppleClang 21.0)]{The options on 64-bit ARM: nanoseconds per operation, Apple M2, AppleClang 21.0, the median of nine runs; the last column is a line of 13 operations} \label{tab:install-worth-arm}
  \begin{tabular}{L{4.6cm}rrrrrrr}
    \hline
    Build & $x + y$ & $x - y$ & $x \times y$ & $x / y$ & \code{sqr} &
    \code{sqrt} & 13 operations\\
    \hline\hline
    \option{-DGAOL\_SIMD=OFF} (FPU intervals) & 1.15 & 1.15 & 3.26 & 3.22 &
    1.78 & 5.17 & 50.6\\
    Default (NEON intervals) & 0.90 & 0.90 & 1.52 & 2.08 & 1.27 & 2.04 &
    16.6\\
    \option{-DGAOL\_SIMD=OFF} \option{-DGAOL\_PRESERVE\_ROUNDING=ON} & 37.4 &
    37.7 & 43.2 & 43.9 & 42.6 & 46.6 & 519\\
    \option{-DGAOL\_PRESERVE\_ROUNDING=ON} (NEON intervals) & 26.6 & 27.1 &
    27.8 & 35.1 & 27.2 & 35.7 & 370\\
    \hline
  \end{tabular}

\end{table}

\begin{itemize}
\item[$\bullet$] \textbf{The NEON intervals pay on every operation of the
  table.} $x - y$ takes 0.90~ns rather than 1.15, $x \times y$ 1.52 rather
  than 3.26 (of operands of random signs 1.60 rather than 3.32), $x / y$
  2.08 rather than 3.22, \code{sqrt} 2.04 rather than 5.17, and a line of 13
  operations 16.6~ns rather than 50.6. The two bounds are one vector, and the
  operations take them from one to the next in a register. The elementary
  functions, which CORE-MATH computes, go through them only for their
  arguments: no faster, no slower. They are on by default on 64-bit ARM
  processors, with every compiler of them.
\item[$\bullet$] \textbf{The product, the quotient and the square of the FPU
  intervals take no branch on the signs of the bounds}, and $x - y$ reads its
  bounds after the check of the rounding direction.\newinvfive{} In the earlier
  versions of GAOL v5 the row read 3.54 for $x - y$, 3.25 for $x \times y$ (6.39 of operands of random
  signs), 3.98 for $x / y$, 2.07 for \code{sqr} and 49.6 for the line of 13
  operations. The product of a line of operands of one sign, which a processor
  predicts, takes more time than it did (a line of 7 products, 35.0~ns rather
  than 23.3), the largest of four products waiting for two comparisons. With
  \option{GAOL\_PRESERVE\_ROUNDING} the write of FPCR waits for the result as
  well, and the FPU row of that option loses 6 to 7~\% on operands of one sign
  ($x \times y$ 42.4~ns rather than 39.7 for the former code, measured in
  another session than Table~\ref{tab:install-worth-arm}, where it is 43.2)
  and gains on operands of random signs (42.4 rather than 47.5). These are the figures
  of this processor, with AppleClang. The processors of the continuous
  integration that run the FPU intervals (32-bit ARM, Ubuntu with GCC and
  Clang, macOS; and, when they were measured, Visual C++ and clang-cl for x86
  and for arm64, which have the SSE2 and the NEON intervals now), measured in
  turn, give other ratios, losses included (the line of 7 products 0.58 with
  GCC on Ubuntu, 0.57 with clang-cl for arm64); POWER, s390x and RISC-V are
  not measured (\pxref{sec:fpu-intervals}).
\item[$\bullet$] \textbf{The rounding direction is read with one \code{mrs}
  of FPCR} (0.8~ns on this processor) and written with \code{msr} only when
  the direction of the program is not upward: \code{msr} takes 4~ns for the
  value FPCR holds and 11~ns when it changes the direction. With
  \option{GAOL\_PRESERVE\_ROUNDING} an operation of a program that computes
  to nearest writes it twice, which is the 27~ns above (37~ns with the FPU
  intervals, which call \code{fegetround()} and \code{fesetround()}); a
  program that already computes upward writes nothing.
\end{itemize}

\subsection{GCC for Windows}\label{sec:install-gcc-windows}

With GCC for Windows (MinGW-w64, MSYS2, Cygwin) on x86-64 and on 32-bit x86,
the fused multiply-add instructions give
\option{-mfma -mprefer-vector-width=128}; for 32-bit x86 the choice follows
the measurements below.\newinvfive{} The earlier versions of GAOL v5 gave nothing, because of GCC bug
54412: \option{-mfma} turns AVX on, the vectorizer of GCC then writes 256-bit
registers (\code{\%ymm}), and GCC takes the stack to be aligned on 32 bytes
for those it spills there, where the ABI of Windows x64 aligns it on 16, so
that an aligned \code{vmovaps} to such a slot faults. Without \option{-mfma},
CORE-MATH called a software \code{fma()} and the integer powers computed no
pairs of products. \option{-mprefer-vector-width=128}, an option of GCC from
GCC 8, keeps the vectorizer to the 128-bit registers. \cmd{objdump -d
--no-show-raw-insn} found no line with \code{\%ymm} in the objects of GAOL,
CORE-MATH and the tests built so by GCC 16.2 (MSYS2 UCRT64), against 30\,598
lines with \option{-mfma} alone, and none without a flag; the whole suite of
tests passed with the three. That is a count of instructions, and no proof
that the defect cannot occur: the continuous integration makes the count
again on the objects it builds, and the check of the fused multiply-add
instructions with it (\pxref{sec:scripts}). A GCC without the option gets no
flag. The code using GAOL is given the two flags (\code{gaol::gaol},
\file{gaol.pc}); the 256-bit intrinsics it writes itself would stay exposed
to the defect, as with \option{-mavx}.

\begin{table}[htbp]
  \centering\renewcommand{\arraystretch}{1.1}\footnotesize
  \setlength{\tabcolsep}{3pt}
  \caption[GCC for Windows on x86-64: nanoseconds per operation (MSYS2 UCRT64, GCC 16.2)]{GCC for Windows on x86-64: nanoseconds per operation, Windows Server 2022 runner of GitHub, MSYS2 UCRT64 with GCC 16.2, the median of three runs of each variant} \label{tab:install-gcc-windows-x64}
  \begin{tabular}{L{3.5cm}rrrrrrrr}
    \hline
    Flags & $x \times y$ & $x / y$ & \code{pow(x, 3)} & \code{exp} &
    \code{log} & \code{sin} & \code{cos} & \code{floor}\\
    \hline\hline
    No flag (\option{-DGAOL\_FMA=OFF}, what GCC for Windows had) & 3.92 &
    5.04 & 52.31 & 46.44 & 106.43 & 93.46 & 260.50 & 4.24\\
    \option{-mfma -mprefer-vector-\allowbreak width=128} (default) & 3.69 & 4.72 & 8.68 &
    32.58 & 33.63 & 46.61 & 55.44 & 2.19\\
    \hline
  \end{tabular}

\end{table}

\medskip
Table~\ref{tab:install-gcc-windows-x64} measures the two variants with
\file{tests/performance.cpp} on the Windows Server 2022 runner of GitHub,
MSYS2 UCRT64 with GCC 16.2, in nanoseconds per operation, the median of three
runs of each variant, alternated in one job (on these lines, the runs of a
variant differ by 1.16 times at most). \option{-mfma} alone gave the same
times. The lines of 15 additions and of 7 products did not change (within
1.5~\%; the line of 7 divisions took 6.3~\% less). The elementary functions
and \code{pow(x, 3)} gain the most, as with Clang above.

\begin{table}[htbp]
  \centering\renewcommand{\arraystretch}{1.1}\footnotesize
  \caption[GCC for Windows on 32-bit x86: nanoseconds per operation and instructions of the objects]{GCC for Windows on 32-bit x86: nanoseconds per operation and lines of the objects with \code{\%ymm} and \code{vfmadd}, MinGW-w64 for i686 (MinGW-Builds GCC 15.2), Windows Server 2022 runner of GitHub} \label{tab:install-gcc-windows-x86}
  \begin{tabular}{L{4.7cm}rrrrrr}
    \hline
    Flags & $x \times y$ & \code{pow(x, 3)} & \code{log} & \code{cos} &
    \code{\%ymm} & \code{vfmadd}\\
    \hline\hline
    No flag (what GCC for Windows had) & 8.65 & 70.45 & 265.68 & 489.59 & 0 &
    0\\
    \option{-mfma} & 8.76 & 12.23 & 81.24 & 138.19 & 4\,709 & 850\\
    \option{-mfma -mprefer-vector-width=128} & 8.66 & 12.12 & 81.93 & 138.33
    & 0 & 850\\
    \option{-mfma -mprefer-vector-width=128 -mstackrealign} & 8.67 & 12.14 &
    82.12 & 139.20 & 0 & 850\\
    \hline
  \end{tabular}

\end{table}

\medskip
For 32-bit x86, Table~\ref{tab:install-gcc-windows-x86} gives the same
program with MinGW-w64 for i686 (MinGW-Builds, GCC 15.2 and 14.2, Windows
Server 2022 runner of GitHub), built, tested (38 tests in every variant, all
passing, two skipped) and measured with each set of flags in turn, in
nanoseconds per operation (GCC 15.2; an earlier run of the same variants gave
times within 1.5~\%). A third run on another runner, faster, gave the same
ratios (\code{pow(x, 3)} 56.6 against 10.05~ns, \code{log} 235.9 against
64.7, \code{cos} 430.1 against 121.3~ns, with
\option{-mfma -mprefer-vector-width=128}). With GCC 14.2 the ratios are
similar (\code{pow(x, 3)} 67.7 against 12.1~ns, \code{log} 268.6 against
103.0, \code{cos} 490.2 against 137.0 in the third run), and \option{-mfma}
alone leaves 600 lines with \code{\%ymm}. GCC 13.2 (MinGW-w64 13.2 for x86)
leaves 9 lines with \code{\%ymm} in two test objects even with both options:
the check lists them and does not fail on 32-bit x86. $x \times y$, which
does not call \code{fma()}, did not change.

\medskip
The stack of a thread on 32-bit Windows is aligned on 4 bytes (GCC bug
40838). A program that keeps 16 intervals alive across calls in a loop of
$+$, $-$, $\times$, $/$ and \code{sqr}, then in a loop of \code{pow(x, 3)},
\code{pow(x, 5)}, \code{pow(x, -4)}, \code{pow(x, 7)}, \code{exp},
\code{log}, \code{sin} and \code{cos} (which \option{-mfma} sends through the
packed kernel of the integer powers), on the stack of the main thread, of a
\code{CreateThread} thread and of a \code{std::thread}, gave the same results
in the three, without a flag, with \option{-mfma} and with
\option{-mfma -mprefer-vector-width=128}, with GCC 15.2 and 14.2. That does
not prove that a stack aligned on 4 bytes can never meet an aligned access of
16 or 32 bytes. \option{-mstackrealign} changed no time with GCC 15.2 (1~\%
at most); \option{-mincoming-stack-boundary=2} was measured once, with the
times of the base. GAOL gives neither. It gives
\option{-mprefer-vector-width=128} with \option{-mfma} as on x86-64, which
leaves no \code{\%ymm} in the objects at the same speed; that bug 54412 acts
on 32-bit x86 was not shown (GCC 15.2 for 32-bit Linux realigns the stack
itself for the 256-bit spills it makes).

\medskip
These flags are for GCC for Windows on x86-64 and on 32-bit x86: MinGW-w64,
MSYS2 and, since the tests of the three builds name it (\code{CYGWIN} of
CMake, \code{\_\_CYGWIN\_\_} of \cmd{configure} and meson; the AVX-512 path
names \code{\_\_CYGWIN\_\_} too), Cygwin.\newinvfive{} The MSYS environment of MSYS2 (not its MinGW ones) is a
Cygwin as far as these tests go, and gets the same flags. If the user gives a
\option{-mprefer-vector-width=} of their own, the builds order it differently
from GAOL's: CMake and \cmd{configure} put the flags of GAOL after the user's
(\option{CMAKE\_CXX\_FLAGS}, \env{CXXFLAGS}), meson puts them before
\code{cpp\_args}. The last of two such options is the one that counts (as
tried with Clang). \option{-mprefer-vector-width=128} does not undo what
other options of the user turn on. With the Clang that cross-compiles for
x86-64, the objects of GAOL, CORE-MATH and the tests built with
\option{-mavx512f -mfma -mprefer-vector-width=128} held 2 lines with
\code{\%ymm} and 2\,887 with \code{\%zmm}, and those built with
\option{-march=skylake-avx512 -mprefer-vector-width=128} 406 lines with
\code{\%ymm}; a second reader counted none, and 4: the gap is not explained,
and the \code{\%zmm} may come from the explicit intrinsics of the AVX-512 path
rather than from the vectorizer. Nothing was run with GCC but a small loop,
which GCC 15.2 compiled with no \code{\%ymm} and no \code{\%zmm} with
\option{-mavx512f -mfma}, with \option{-march=skylake-avx512} and with
\option{-march=native}, all with \option{-mprefer-vector-width=128}. The code
of GCC itself (\code{ix86\_option\_override\_internal}, version 14.2, read,
not run) applies the preferred width of an architecture only if none was
given, so that \option{-mprefer-vector-width=128} should win over a
\option{-march}. A program built with such options is exposed to GCC bug
54412, as with \option{-mavx}.

\begin{Rem}
Cygwin x86-64 is built and tested by a job of the continuous integration (GCC
14.4, CMake 4.4, Release and Debug; not measured), whose first runs found what Cygwin
needed.\newinvfive{} Its \file{<stdlib.h>} does not declare
\code{posix\_memalign()} under \option{-std=c++11} (the aligned memory is
\code{\_mm\_malloc()} there). Its \code{fma()} is not correctly rounded, and
\code{std::fma} is a call of it without optimization: GAOL calls
\code{\_\_builtin\_fma} (\code{fma\_once()}), which is the instruction with
\option{-mfma} at \option{-O0} and at \option{-O3} (both checked), and
refuses Cygwin without \option{-mfma} (\pxref{sec:refused}). Its
\code{nextafter(-DBL\_MAX, -oo)} is \code{-DBL\_MAX} in the rounding
direction upward, which \code{previous\_float()} and \code{next\_float()} now
tell apart (the lower bound of \code{sinh(-711)} did not enclose the value).
\code{is\_bisectable()} and \code{bisect()} call no \code{nextafter()}: they
tell the bounds by their bits.
\file{tests/core\_math.cpp} and \file{tests/arithmetic.cpp} call
\code{\_\_builtin\_fma} as GAOL does.
\end{Rem}

\subsection{Interprocedural optimization is not offered}
\label{sec:install-ipo}
\cindex{optimization!interprocedural}%

Compiled with it (\option{-flto}), the basic operations of GAOL are 1.1 to
1.5 times faster where the code using GAOL is compiled with it too: an
addition of intervals took 0.73~ns rather than 1.07 with Clang 21.1,
$x \times y$ 1.71 rather than 2.04, \code{sqr(x)} 1.27 rather than 1.80. The
elementary functions are no faster, and a little slower: \code{cos} took
39.2~ns rather than 36.6, \code{sin} 32.6 rather than 31.3, \code{log} 22.7
rather than 21.7. With Clang 18.1, before the operators took their bounds in
a register, the basic operations took two to three times less time with it
and the elementary functions two to three times more. The sources of
CORE-MATH are compiled apart, without it. The standard options
of the tools reach the sources of GAOL, \option{CMAKE\_INTERPROCEDURAL\_OPTIMIZATION},
\option{-Db\_lto=true} of meson and \option{-flto} in \env{CXXFLAGS}, which the
three builds keep away from CORE-MATH alone: they are not to be given with
GCC.

\medskip
With GCC, moreover, it breaks the bounds: GAOL then fails
\file{tests/rounding\_direction.cpp}, which checks that an operation gives
the same result whatever the rounding direction the calling code left. With
\option{GAOL\_PRESERVE\_ROUNDING},
\code{textToInterval("sin(1)+exp(0.1)")} gave an upper bound one double below
the right one when the caller left the direction downward or toward zero: the
compiler moves floating-point operations across the changes of rounding
direction, which \option{-frounding-math} is meant to forbid and which it
does not do across translation units. Compiling the sources of GAOL alone
with it is enough for the test to fail. Clang 18.1 passes the tests with it,
so this is not a property of the option but of the handling of it by one
compiler, which is reason enough not to offer it. An interval that does not
enclose the exact value is worse than a slow one: there is no option for
this, and \file{gaol.pc} and \code{gaol::gaol} carry no such flag.

\section{Tests, examples, performance and the parser}\label{sec:tests}
\cindex{tests}%

Built with \option{WITH\_TESTS} (CMake), \loption{with-tests}
(\cmd{configure}) or \code{with-tests} (meson), the programs of \file{tests}
compare the bounds GAOL computes with the exact results of the operations,
independently of GAOL and of the floating-point environment: the exact
results are computed with integers, or were computed with \code{mpmath}, with 1500 to 7000 bits of
precision. They check, among others, that the basic operations are the
tightest enclosures, that the elementary functions are the tightest or within
one double of them, that every operation gives the same result whatever the
rounding direction the calling code left, that numbers are read exactly and
intervals written rounded outward, that the reverse functions are accurate in
the sense of IEEE 1788-2015 \cite{IEEE1788:2015}, and that each name of
\code{gaol\_ieee1788} is the operation of the standard it names. The unit
tests of GAOL~4, which were in \file{check/} and checked the results of the
operations on chosen intervals with CppUnit, are among them, in \file{tests}
now, and need CppUnit no more.\newinvfive{} What each program checks is
described in Chapter~\ref{chap:tests} (\pxref{sec:test-programs}). The three
builds have the same targets (Table~\ref{tab:install-targets}).

\begin{table}[htbp]
  \centering\renewcommand{\arraystretch}{1.1}\footnotesize
  \caption{The targets of the three builds} \label{tab:install-targets}
  \begin{tabular}{L{3.0cm}L{4.2cm}L{2.8cm}L{3.2cm}}
    \hline
    & \multicolumn{1}{c}{CMake} & \multicolumn{1}{c}{\cmd{configure}} &
    \multicolumn{1}{c}{meson}\\
    \hline\hline
    The unit tests of \file{tests} & \cmd{make test}
    (\cmd{cmake --build build --target test}; \code{RUN\_TESTS} with Visual
    Studio), with \option{WITH\_TESTS} & \cmd{make test}, with
    \loption{with-tests} & \cmd{meson test} or \cmd{ninja test}, with
    \code{with-tests}\\
    The unit tests of \file{tests} and the examples of \file{examples} & \cmd{make check}
    (\loption{target check}), with \option{WITH\_TESTS} and
    \option{WITH\_EXAMPLES} & \cmd{make check}, with \loption{with-tests}
    and \loption{with-examples} & \cmd{ninja check}
    (\cmd{meson test --setup check}), with \code{with-tests} and
    \code{with-examples}\\
    GAOL v5 measured on the benchmark of the comparison of the libraries &
    \cmd{make perf} & \cmd{make perf} & \cmd{ninja perf}\\
    The lexer and the parser regenerated with flex and bison &
    \cmd{make parser} & \cmd{make parser} & \cmd{ninja parser}\\
    The manual of GAOL v5 & \cmd{make pdf}
    (\loption{target pdf}), with \option{WITH\_DOC} & \cmd{make -C manual pdf}
    & \cmd{ninja pdf}, with \code{with-doc}\\
    \hline
  \end{tabular}

\end{table}

\begin{itemize}
\item[$\bullet$] \textbf{The tests.} \cmd{ctest} runs the tests and the
  examples of a CMake build; \cmd{make test} leaves the examples out (label
  \code{example}, CMake 3.17 and later), and runs the compile tests
  \code{refused\_*} and \code{nodiscard\_*} too. With meson before 0.57,
  \cmd{meson test} runs the examples too, and \cmd{meson test --suite unit}
  the tests alone.
\item[$\bullet$] \textbf{The examples} check what they print and fail
  otherwise (\pxref{sec:test-programs}).
\item[$\bullet$] \textbf{\cmd{make perf}} compiles
  \file{doc/compare/code/bench\_gaol.cpp} with the flags and the library of
  the build, runs it on the intervals of the benchmark, and writes its times
  into the tables of \file{doc/compare/performance.md}, with the times of the
  other libraries measured by the last whole run, which
  \file{doc/compare/code/results.csv} keeps: GAOL 4.2.3, filib++,
  libieeep1788, PROFIL/BIAS, Solaris Studio and Boost.Interval are not run
  (\pxref{sec:comparison}); it runs the ten long lines of
  \file{doc/compare/code/run\_lines.sh} too. It needs bash, and Python~3 with
  numpy; with CMake and meson, Linux and macOS (Linux only before), and with
  \cmd{configure} wherever bash is. There is no target \cmd{perf} in the Debug
  build of CMake, nor with \option{-Denable-debug=true} of meson.\newinvfive
\item[$\bullet$] \textbf{\cmd{make parser}} regenerates the lexer and the
  parser of the strings with the flex and the bison of the system, where the
  CMake build finds them (\code{find\_program}, not required;
  \pxref{sec:install-maintainers}).
\item[$\bullet$] \textbf{\cmd{make pdf}} builds this manual with pdflatex,
  bibtex and makeindex, by \file{manual/build-pdf.sh}, which the three builds
  run: with CMake, after \cmd{cmake -S . -B build -DWITH\_DOC=ON}, by
  \cmd{cmake --build build --target pdf}, which CMake had not
  before.\newinvfive{} The PDF is written in \file{manual} of the sources,
  \file{manual/Manual\_GAOLv5.pdf}, and not in the build directory, which
  keeps the files of LaTeX; the manual of GAOL~4 is not built, its PDF,
  \file{manual/v4/Manual\_GAOLv4.pdf}, was built once; the target is not part
  of \code{all}. This manual
  can also be compiled by hand in \file{manual/v5}:
  \cmd{pdflatex manual\_gaolv5.tex}, then \cmd{bibtex manual\_gaolv5},
  \cmd{makeindex manual\_gaolv5.idx -o manual\_gaolv5.ind} and \cmd{pdflatex}
  twice; without \file{gaol\_version.tex}, which the builds write, the title
  page takes the version of \file{VERSION.txt}, the edition 5.0 and today's
  date.
\end{itemize}

\noindent
\file{tests/find\_package} builds the same tests with an installed GAOL, and
\file{tests/fetch\_content} with a GAOL brought into a project by
\code{FetchContent}. \code{gaol\_performance} (\file{tests/performance.cpp})
measures the time per operation of the arithmetic and the elementary
functions of GAOL; its figures are those of
Section~\ref{sec:install-options-worth}.

\section{What is installed}\label{sec:install-installed}

Each build installs:
\begin{itemize}
\item[$\bullet$] \textbf{the headers}, in \file{\var{prefix}/include/gaol};
\item[$\bullet$] \textbf{the library}, \file{libgaol}, which holds CORE-MATH:
  static by default with the three builds, or shared with
  \option{BUILD\_SHARED\_LIBS},
  \loption{enable-shared} of \cmd{configure} and
  \option{-Ddefault\_library=shared} of meson (static only with Visual C++).
  The shared library is \file{libgaol.so.5.0.0}, whose soname is
  \file{libgaol.so.5}, with the three builds;
\item[$\bullet$] \textbf{\file{gaol.pc}}, for pkg-config, in the
  \file{pkgconfig} directory of the library directory (except the CMake build
  with Visual C++);
\item[$\bullet$] \textbf{the CMake package of GAOL}, with CMake, which defines
  the target \code{gaol::gaol}.
\end{itemize}

\noindent
Both \file{gaol.pc} and \code{gaol::gaol} carry the flags that the code
including the headers of GAOL has to be compiled with (\pxref{sec:compiling}).

\subsection{The configuration header}\label{sec:install-configuration-header}
\findex{gaol\_configuration.h}%

Each build writes \file{gaol/gaol\_configuration.h}, which
\file{gaol/gaol\_config.h} includes with every compiler: CMake from
\file{cmake/gaol\_configuration.h.in}, \cmd{configure} from
\file{gaol/gaol\_configuration.h.in} (written by hand, not by autoheader),
meson from \file{gaol/gaol\_configuration\_meson.h.in}. The three define the
same macros, and nothing else, but for the switch of a developer that CMake
alone writes, \code{GAOL\_SSE2\_UNALIGNED\_STORAGE}
(Table~\ref{tab:install-developers}); the continuous integration checks that
they write the same header on Ubuntu 24.04 (x86-64 and arm64, GCC and
Clang), macOS 15 (arm64 and x86-64), Debian 12 (i386), Debian 13 (armhf) and
MSYS2 (UCRT64 and MINGW32) (Section~\ref{sec:three-builds}).\newinvfive{} The sources of GAOL read all
of them but \code{GAOL\_HAVE\_AVX512\_TARGET} and \code{GAOL\_PREFER\_AVX512},
which the three builds use to decide \code{GAOL\_AVX512\_ARITHMETIC} and the
continuous integration prints, and \code{GAOL\_UNIVERSAL\_SSE2} and
\code{GAOL\_UNIVERSAL\_NEON}, which only the configuration header reads, to
define the macros of the architecture being compiled;
\code{GAOL\_USING\_SSE3\_INSTRUCTIONS} is read only by the code of the
intervals of floats, which a developer of GAOL compiles.

\begin{Rem}
Each macro of Table~\ref{tab:install-configuration} but the version,
\code{GAOL\_UNIVERSAL\_SSE2} and \code{GAOL\_UNIVERSAL\_NEON} is 1 where the
table says, and 0 otherwise; those two are defined only in a universal build of
macOS.\newinvfive{} A macro that is 0 now was left undefined in
GAOL~4, so that a header reading it with \code{\#if} warned under
\option{-Wundef}, and a program could not include GAOL under
\option{-Wundef -Werror}. A program tests these macros with \code{\#if} as
well: \code{\#ifdef} is now true of each of the others.
\end{Rem}

\begin{longtable}{L{4.6cm}L{9.4cm}}
  \caption{The macros of \file{gaol/gaol\_configuration.h}}
  \label{tab:install-configuration}\\
  \hline
  \multicolumn{1}{c}{Macro} & \multicolumn{1}{c}{1 when}\\
  \hline\hline
  \endfirsthead
  \multicolumn{2}{l}{\textit{Table~\thetable, continued}}\\*
  \hline
  \multicolumn{1}{c}{Macro} & \multicolumn{1}{c}{1 when}\\
  \hline\hline
  \endhead
  \endfoot
  \hline
  \endlastfoot
  \code{GAOL\_MAJOR\_VERSION}, \code{GAOL\_MINOR\_VERSION},
  \code{GAOL\_MICRO\_VERSION}, \code{GAOL\_VERSION} & Always, from
  \file{VERSION.txt}: \code{5}, \code{0}, \code{0} and \code{"5.0.0"} for
  GAOL 5.0.0 (\code{project()} with CMake and meson, the file read when
  \cmd{configure} runs)\\
  \code{GAOL\_DEBUGGING} & In a Debug build:
  \code{GAOL\_DEBUG\_VERBOSE(level, command)} runs its commands
  (\option{CMAKE\_BUILD\_TYPE=Debug} or \option{GAOL\_DEBUG=ON},
  \loption{enable-debug}, \loption{buildtype=debug} or \code{enable-debug})\\
  \code{GAOL\_EXCEPTIONS\_ENABLED} & GAOL raises exceptions rather than
  abort: by default in the three builds (\option{GAOL\_EXCEPTIONS\_ENABLED=OFF},
  \loption{disable-exceptions}, \code{enable-exception=false} turn it off)\\
  \code{GAOL\_PRESERVE\_ROUNDING} & The operations restore the rounding
  direction they found (\option{GAOL\_PRESERVE\_ROUNDING},
  \loption{enable-preserve-rounding}, \code{enable-preserve-rounding})\\
  \code{GAOL\_PREFER\_AVX512} & The four arithmetic operations, the square,
  the square root and, with FMA, the integer powers beyond 1 and $-1$ are to
  take the AVX-512 instructions: with \option{GAOL\_PRESERVE\_ROUNDING},
  where the compiler takes the path (\code{GAOL\_HAVE\_AVX512\_TARGET});
  \option{GAOL\_PREFER\_AVX512=ON} or \code{OFF},
  \loption{enable-prefer-avx512}, \loption{disable-prefer-avx512} and
  \code{enable-prefer-avx512} decide otherwise\\
  \code{GAOL\_AVX512\_ARITHMETIC} & The same operations take the AVX-512
  instructions: \code{GAOL\_PREFER\_AVX512}, \code{GAOL\_HAVE\_AVX512\_TARGET},
  and a machine building that runs them, not when cross-compiling; the
  library then needs a processor with them (checked by the three builds)\\
  \code{GAOL\_VERBOSE\_MODE} & A line on the standard error when GAOL
  initializes and cleans up (\option{GAOL\_VERBOSE\_MODE=ON},
  \loption{enable-verbose-mode}, \code{enable-verbose-mode})\\
  \code{GAOL\_USING\_SSE2\_INSTRUCTIONS} & The intervals computed with SSE2
  (x86; on 32-bit Windows with GCC and Clang told to ask for aligned memory,
  \option{-faligned-new=8}, and with Visual C++ and clang-cl, whose bounds
  are two doubles there): \option{GAOL\_SIMD}, \loption{enable-simd},
  \code{enable-simd}, on by default\\
  \code{GAOL\_USING\_NEON\_INSTRUCTIONS} & The intervals computed with NEON
  (64-bit ARM, with every compiler; not 32-bit ARM): \option{GAOL\_SIMD},
  \loption{enable-simd}, \code{enable-simd}, on by default
  (\option{GAOL\_ARM64}, \loption{enable-arm64} and \code{enable-arm64} were
  options of their own)\\
  \code{GAOL\_USING\_SSE3\_INSTRUCTIONS} & And compiled with \option{-msse3},
  where the compiler takes it\\
  \code{GAOL\_UNIVERSAL\_SSE2}, \code{GAOL\_UNIVERSAL\_NEON} & A universal
  build of macOS (several architectures in \option{CMAKE\_OSX\_ARCHITECTURES},
  or \option{-arch} several times in the flags of \cmd{configure} and meson):
  the three macros above are then those of the architecture being compiled,
  \code{GAOL\_USING\_SSE2\_INSTRUCTIONS} and
  \code{GAOL\_USING\_SSE3\_INSTRUCTIONS} for x86\_64 and
  \code{GAOL\_USING\_NEON\_INSTRUCTIONS} for arm64, where \option{GAOL\_SIMD}
  allows them; the fused multiply-add and the AVX-512 path stay off, their
  tests failing for one of the architectures\\
  \code{GAOL\_HAVE\_AVX512\_TARGET} & The compiler takes
  \code{\_\_attribute\_\_((target("avx512f")))} and
  \code{\_\_builtin\_cpu\_supports("avx512f")}, which the AVX-512 path needs;
  never GCC for Windows (MinGW-w64, Cygwin), which does not align the stack
  for the AVX-512 registers (GCC bug 54412); Visual C++ for x64 takes the
  instructions in any function and tells the processor by \code{cpuid} and
  \code{xgetbv}, clang-cl has the attribute and does the same (checked by the
  three builds)\\
  \code{GAOL\_HAVE\_VISIBILITY\_OPTIONS} & The library is compiled with
  \option{-fvisibility=hidden} (checked)\\
  \code{GAOL\_HAVE\_FENV\_H} & The compiler has \file{<fenv.h>} (checked;
  required by CMake)\\
\end{longtable}

\noindent
What GAOL needs to know of the processor and the system comes from the
compiler, in \file{gaol/gaol\_config.h}: \code{GAOL\_IX86\_MACOSX} and
\code{GAOL\_ARM\_MACOSX}, the sizes of the integer types
(\code{GAOL\_SIZEOF\_INT}, \code{GAOL\_SIZEOF\_LONG\_LONG\_INT}, from
\file{<limits.h>}) and the order of the bytes (\code{GAOL\_WORDS\_BIGENDIAN}),
the names of the systems and the order of the bytes being 1 or 0.
\code{GAOL\_IX86\_LINUX} and \code{GAOL\_AARCH64\_LINUX} are defined as well,
though no code of GAOL reads them. These macros, and the
\code{GAOL\_USING\_SSE*} and \code{GAOL\_HAVE\_*} of the table, were named
without \code{GAOL\_} before GAOL v5.\newinvfive

\begin{Rem}
The \cmd{configure} and meson builds of GAOL~4 defined some fifty other
macros, of GAOL~4 or of autoconf, which nothing read, and they are
gone:\newinvfive{} the checks of C headers (\code{HAVE\_STDLIB\_H},
\code{HAVE\_UNISTD\_H}, \code{STDC\_HEADERS}\ldots), of functions
(\code{HAVE\_FLOOR}, \code{HAVE\_POW}, \code{HAVE\_MALLOC},
\code{HAVE\_REALLOC}, \code{HAVE\_FINITE}\ldots), of types and sizes
(\code{HAVE\_\_BOOL}, \code{SIZEOF\_LONG}\ldots), \code{const},
\code{inline}, \code{size\_t}, \code{malloc} and \code{realloc}
(\code{rpl\_malloc()} and \code{rpl\_realloc()} of \code{AC\_FUNC\_MALLOC}
and \code{AC\_FUNC\_REALLOC} in a cross-compilation),
\code{CLOCK\_IN\_HEADER}, \code{LT\_OBJDIR}, \code{PACKAGE}, \code{VERSION}
and \code{PACKAGE\_*}, which \file{gaol/gaol\_config.h} undefined, undefining
those of the code including GAOL. \code{HAVE\_LIMITS}, \code{HAVE\_CASSERT},
\code{HAVE\_CLOCK} and \code{GETRUSAGE\_IN\_HEADER} were read, but are true
with every compiler of C++11, and \code{HAVE\_NEXTAFTER} and
\code{HAVE\_ISNAN} were read for Visual C++, which has both: the code that
read them no longer does.
\end{Rem}

\section{Discussion}

This chapter established how GAOL v5 is built: from the sources alone, with
CORE-MATH compiled into \file{libgaol}, by one of three builds that give it
the same macros and the same flags, and whose few options are the same under
three names. The choices it explained are those of a library whose bounds
must be the same everywhere. The compilers and the options that gave bounds
not enclosing the exact results, or that nothing certifies, are refused when
GAOL is configured or compiled, and, but for Clang 14 for 64-bit ARM, again when the code using it is compiled,
and the link with \option{-Ofast} is countered by the clearing of
the flush-to-zero modes and by \option{-mno-daz-ftz}. The fused multiply-add
instructions and the AVX-512 path are decided by what the machine building
runs. The elementary functions gain up to 2.5 times with the fused
multiply-add instructions (4.7 times with GCC for Windows), and the arithmetic
operations 2.9 times with the AVX-512 path where the rounding direction is
preserved; the price is a library bound to the processors that have them, and
a package for any processor turns them off. Interprocedural
optimization, which would gain 1.1 to 1.5 times on the basic operations, is
not offered, because GCC then moves operations across the changes of the
rounding direction.

\medskip
Some points are left open. The refusal of the older MinGW-w64
is not tested on ARM; the processors POWER, s390x and RISC-V are not measured;
the flush-to-zero modes of these processors are neither checked nor cleared;
the Python aliases of Windows were read in the sources of meson, not checked;
and the cause of the failure of \code{package\_source} on Cygwin is not known.
The code using GAOL still has to be compiled with the flags of interval
arithmetic and without the options that no macro shows: the next chapter
shows how a program includes the headers, compiles with these flags, from
CMake, pkg-config or meson, and uses the library.

%% !TEX root = ../manual_gaolv5.tex
% Chapter 3.
%% Chapter 3: Using GAOL. Sources: gaol.tex of f5bf180 lines 840-1108,
%% 1750-1905, 1906-2300 and 5789-5880; doc/using.md lines 1-227, 291-450,
%% 572-630 and 695-1069; examples/tutorial.md (the names of the examples).
\chapter{Using GAOL}\label{chap:overview}
%%=========================

The previous chapter built and installed GAOL. This chapter shows how a
program uses it: what it includes, how it is compiled and linked with the
flags of interval arithmetic that the inline operations of GAOL require, and
which of the two namespaces it opens. It then reviews what GAOL gives to the
algorithms of interval computing, the errors most often made with intervals,
and the names of the other interval libraries next to those of GAOL.

\medskip
In this chapter, we assume that GAOL is installed, and that its library and
its headers are accessible to the compiler (Chapter~\ref{chap:installation}).
The first section walks through a complete program, from the header it
includes to the enclosure it prints. The second section lists the headers a
program may include, and the third gives the flags of interval arithmetic and
the way to take them from CMake, from pkg-config and from meson. The fourth
and the fifth sections present the namespaces \code{gaol} and
\code{gaol\_ieee1788}, and the power, where they differ. The sixth section
tells how an integer is taken, and the seventh what the compiler says of a
result thrown away. The eighth section gathers the tools of interval
algorithms, the ninth summarizes how GAOL reports an error, the tenth reviews
the common errors, and the last one gives the names of IBEX, Codac, C-XSC,
PROFIL/BIAS and Boost.Interval next to those of GAOL.

\section{A first program}\label{sec:first-program}
%%-----------------------
\cindex{first program}%

Let us consider the Goldstein-Price function
\begin{align*}
f(x,y)=&\bigl(1+(x+y+1)^2(19-14x+3x^2-14y+6xy+3y^2)\bigr)\\
&\bigl(30+(2x-3y)^2(18-32x+12x^2+48y-36xy+27y^2)\bigr)
\end{align*}
and let us enclose its range for $x\in\itv{-2}{2}$ and $y\in\itv{-2}{2}$.

\begin{Ex}\label{ex:goldstein-price}
The program below computes the formula of $f$ with intervals in place of
numbers. The line under it is what it prints.
 \begin{example}
 #include <iostream>
 #include <gaol/gaol.h>
 using namespace gaol;

 int main(void)
 {
   interval
     x(-2,2),
     y(-2,2), z;

   z=(1+sqr(x+y+1)*(19-14*x+3*sqr(x)-14*y+6*x*y+3*sqr(y)))*
     (30+sqr(2*x-3*y)*(18-32*x+12*sqr(x)+
        48*y-36*x*y+27*sqr(y)));

   std::cout << "z = " << z << std::endl;
   gaol::cleanup(); @rem^// Always, after the last use of GAOL~
   return 0;
 }
 @outputs^z = [-87881320, 147125080]~
 \end{example}
\end{Ex}

The program includes the header \file{gaol/gaol.h}, which gives access to
every facility of GAOL. The type \code{interval} and the functions of GAOL are
in the namespace \code{gaol\_core}, which a program names through the
namespace \code{gaol}: the \code{using} directive of Line~3 imports it. The
\file{gaol/gaol.h} header of GAOL~4 imported the namespace itself; that of
GAOL v5 no longer does, and a program adds the directive, or names
\code{gaol::interval} and \code{gaol::sqr}
(\pxref{sec:namespaces}).\newinvfive
\cindex{gaol!namespace}%
\cindex{namespace!gaol}%
\cindex{header!and namespace|see{namespace}}%

Nothing has to be called before using GAOL: \code{gaol::init()} runs by
itself when the program starts, before \code{main()}, with every compiler
(\pxref{chap:init}), and each operation of GAOL sets the rounding direction
it needs (\pxref{sec:rounding-direction}).\newinvfive

\begin{Rem}\label{rem:cleanup}
The program calls \code{gaol::cleanup()} right after its last use of GAOL, at
the latest at the end of \code{main()}, whatever the compiler. Unless GAOL is
built to preserve the rounding direction, its initialization sets the rounding
direction upward for the whole program, not only for the operations of GAOL,
and \code{gaol::cleanup()} sets back the one the program started with. GAOL
also calls it itself, but only when the program ends: without the call, the
code after the last use of GAOL would compute its doubles rounded
upward.\newinvfive
\end{Rem}

The function \code{sqr(\var{x})} is the square of \var{x}, and is equivalent
to \code{pow(\var{x},2)}: it is the function that \code{pow(\var{x},2)}
calls, with the same bounds, the sign of a zero bound
included.\newinvfive

\medskip
Let \file{f.cpp} be the name of the file containing the program above. It is
compiled with the flags that pkg-config gives (\pxref{sec:compiling}):
\cindex{compiling!a file using gaol}%

\begin{cmdshell}
% c++ -std=c++11 -O2 $(pkg-config --cflags gaol) -o f f.cpp \
      $(pkg-config --libs gaol)
\end{cmdshell}

\noindent
GAOL is the only library to link: CORE-MATH, which computes its elementary
functions, is compiled into it, and the C++ compiler links the mathematical
library of the system itself.

Executing \file{f} prints the line shown under the program. We then know that
$f(x,y)$ lies in \itv{-87881320}{147125080} for every $x$ and $y$ in
\itv{-2}{2}: \code{z} \emph{encloses} the range of $f$. It is not the range,
which is \itv{3}{1015690.27\ldots}: the minimum is $f(0,-1)=3$, and the
maximum is reached on the edge $y=2$. \code{z} is 231 times wider. Each
occurrence of $x$ and of $y$ in the formula ranges over \itv{-2}{2} on its
own, as if the others did not exist: this is the dependency problem
(\pxref{sec:dependency}), which makes an enclosure wider than the range.
\cindex{dependency problem}%
Cutting $x$ and $y$ into smaller intervals, and taking the hull of the values
of $f$ over the pieces, narrows the enclosure, at the price of more
evaluations; \file{examples/03\_dependency\_problem.cpp} of the sources does
it.

\section{The headers}\label{sec:headers}
%%-------------------
\cindex{headers}%

A program includes \file{gaol/gaol} (or \file{gaol/gaol.h}), and
\file{gaol/parser/gaol\_expression.h} for the expressions, with
\file{gaol/parser/gaol\_expr\_eval.h} to evaluate them. The three builds
install these headers and the headers they include. Each of the installed
headers compiles alone, but for the four that other headers include in the middle
of their code (\file{gaol/core/gaol\_interval\_fpu.h},
\file{gaol/core/gaol\_interval\_neon.h},
\file{gaol/core/gaol\_interval\_sse.h}) or that Visual C++ alone uses
(\file{gaol\_fpu\_msvc.h}). The headers of the sources of GAOL are not
installed:\newinvfive{} \file{gaol/parser/gaol\_interval\_parser.h}, which did
not compile when included, \file{gaol\_init\_cleanup.h},
\file{gaol/math/gaol\_core\_math.h}, \file{gaol/math/gaol\_u128.h},
\file{gaol/math/gaol\_compoundm1.h} and \file{gaol/math/gaol\_roundeven.h}.
Three headers of GAOL~4 are gone: \file{gaol\_parameters.h}, which declared
nothing, \file{gaol\_exact.h} with \file{sysdeps/}, which \file{dtoa.c}
needed, and \file{gaol\_assert.h}, whose \code{GAOL\_ASSERT} is
\code{assert()} now.

\medskip
The headers of GAOL compile without a warning under \option{-Wall -Wextra
-Wold-style-cast -Wundef}, and put nothing into the global namespace.
\file{gaol/gaol\_exceptions.h} declared \code{string} and \code{exception}
there; the header is gone, GAOL throws the exceptions of the standard library
(\pxref{sec:errors-overview}), and a program names \code{std::string} and
\code{std::exception}.\newinvfive{} The macros the headers define start with
\code{GAOL\_} (or \code{gaol\_}), so as not to meet those of the program or of
another library:\newinvfive{} \code{USING\_SSE2\_INSTRUCTIONS},
\code{HAVE\_FENV\_H}, \code{WORDS\_BIGENDIAN}, \code{INLINE},
\code{MEMALIGN}, \code{\_\_HI}\ldots{} are now
\code{GAOL\_USING\_SSE2\_INSTRUCTIONS}, \code{GAOL\_HAVE\_FENV\_H},
\code{GAOL\_WORDS\_BIGENDIAN}, \code{GAOL\_INLINE}, \code{GAOL\_MEMALIGN},
\code{GAOL\_HI}\ldots{} The header \file{<stdlib.h>} is no longer included
with \code{\_XOPEN\_SOURCE} defined again, and \code{nb\_fp\_numbers()}
returns an \code{unsigned long long}, which GAOL~4 named with the macro
\code{ULONGLONGINT}.

\medskip
The doubles \code{pi}, \code{half\_pi}, \code{two\_pi}, \code{pi\_dn},
\code{pi\_up}, \code{half\_pi\_dn}, \code{half\_pi\_up}, \code{ln2\_dn},
\code{ln2\_up}, \code{two\_power\_51} and \code{two\_power\_52} of GAOL~4 are
no longer declared for the program, which found them with \code{using
namespace gaol}, its own \code{pi} being then ambiguous.\newinvfive{} The
bounds of $\pi$ are those of \code{interval::pi()}, and the doubles around
$\pi$ and $\pi/2$ are \code{gaol\_pi\_dn}, \code{gaol\_pi\_up},
\code{gaol\_half\_pi\_dn} and \code{gaol\_half\_pi\_up}, in \code{gaol},
\code{gaol\_ieee1788} and \code{gaol\_core}.

\section{Compiling a program using GAOL}\label{sec:compiling}
%%--------------------------------------
\cindex{flags of interval arithmetic}%
\cindex{compiling!a file using gaol}%

The operations of GAOL on intervals are inline: the code that includes the
headers of GAOL has to be compiled with the \dfn{flags of interval
arithmetic}, not only GAOL itself.\newinvfive{} With GCC and Clang, each where
the compiler takes it, the flags are the following.

\begin{itemize}
\item \option{-frounding-math -fno-fast-math -ffp-contract=off}, so that the
  compiler rounds each operation in the direction in effect, as written, and
  neither evaluates it at compile time in the default rounding nor contracts a
  multiplication and an addition into a fused one.
\item \option{-msse2 -mfpmath=sse} on 32-bit x86, so that doubles are computed
  in double precision rather than on the x87 unit.
\item \option{-msse2 -msse3} on x86 processors with the SSE2 intervals
  (\option{GAOL\_SIMD}).
\item \option{-faligned-new=8} with GCC and Clang on 32-bit Windows
  (MinGW-w64, MSYS2) with the SSE2 intervals, so that the compiler asks for
  aligned memory for the intervals, which the C runtime would align on 8 bytes
  instead of 16.\newinvfive{} \code{gaol::gaol} and \file{gaol.pc} give it to
  the code using GAOL, which allocates the intervals, and
  \file{gaol/gaol\_config.h} checks it. It is an option of C++: GCC warns of a
  C file compiled with it (``valid for C++/ObjC++ but not for C''), and with
  \option{-Werror} the warning is an error. This happens where the flags of
  \file{gaol.pc}, or \code{gaol\_dep} of meson, are given to every file of a
  target, whereas \code{gaol::gaol} of CMake gives them to the C++ files only:
  a project that compiles C files gives the option to its C++ files only,
pkg-config having no flag for one
  language. Visual C++ and clang-cl have no such flag, and need none: their
  SSE2 intervals on 32-bit Windows keep their two bounds as two doubles, whose
  alignment is 8 (\code{GAOL\_SSE2\_UNALIGNED\_STORAGE} of
  \file{gaol/gaol\_config.h}).
\item \option{-mstackrealign} (or \option{-mincoming-stack-boundary=2}) with
  GCC on 32-bit x86, only for code that something calls with the stack aligned
  on 4 bytes, and not on 16 as GCC takes it (bug 40838 of GCC): a caller that
  does not follow the ABI of Linux, a hand-written trampoline or a foreign
  runtime for instance. A function of the program that keeps intervals in
  locals then stops on a \code{movaps} at an address that is not aligned
  (\code{test\_stack\_entered\_aligned\_on\_4\_bytes} of
  \file{tests/misc.cpp}, \pxref{sec:test-programs}).\newinvfive{} The
  programs for 32-bit Windows need neither: the system enters its threads and
  its callbacks with the stack aligned on 4 bytes, and that test passes
  without them with GCC for MinGW-w64, with Clang and with clang-cl, in the
  continuous integration.
\item Nothing more on 64-bit ARM processors with the NEON intervals
  (\option{GAOL\_SIMD}, \pxref{sec:neon}): the first item is all they ask
  for.\newinvfive
\item \option{-ffloat-store} where doubles are still computed on the x87 unit.
\item \option{-mfma} on x86, and \option{-mfpu=neon-vfpv4 -mfloat-abi=hard} on
  32-bit ARM, where GAOL is compiled with the fused multiply-add instructions
  of the processor (\option{GAOL\_FMA}, \pxref{sec:cmake-options}): the code
  using GAOL then runs only on a processor with them, as GAOL itself. With GCC
  for Windows on x86-64 and on 32-bit x86 the flags are \option{-mfma
  -mprefer-vector-width=128}, the second keeping the vectorizer of GCC to the
  128-bit registers, which it would otherwise spill to a stack that Windows
  aligns on 16 bytes, not 32 (GCC bug 54412): the objects of GAOL, CORE-MATH
  and the tests built so by GCC~16.2 have no instruction on a 256-bit
  register. A program that writes 256-bit intrinsics itself is exposed to that
  defect, as with \option{-mavx}, and so is one built with
  \option{-mavx512f} or a \option{-march} that turns on wider registers, which
  the option may not always undo. Cygwin and MSYS get these flags as the other
  GCC for Windows; Cygwin is built and tested by a job of the continuous
  integration, and refused without them.\newinvfive
\end{itemize}

\noindent
With Visual C++ and clang-cl, the flags are \option{/fp:strict},
\option{/arch:AVX2} for x64 with \option{GAOL\_FMA}, and \option{-msse2} with
clang-cl for 32-bit x86, which computes doubles on the x87 unit with
\option{/arch:SSE2}. Each build installs these flags with GAOL, in
\code{gaol::gaol} and \file{gaol.pc}. \file{gaol/gaol\_config.h} refuses code
compiled by Visual C++, or by clang-cl from Clang 16, without
\option{/fp:strict}. GCC and Clang do not tell the code whether
\option{-frounding-math} and \option{-ffp-contract=off} were given: there, it
only refuses what contradicts them, \option{-ffast-math},
\option{-ffinite-math-only} and doubles computed on the x87 unit
(\pxref{sec:refused}). The order of the options matters.
\option{-fno-fast-math} turns the first two off when it comes after them on
the command line, and does nothing when \option{-ffast-math} or
\option{-ffinite-math-only} comes after it, where the compilation stops.
\option{-Ofast}, which turns \option{-ffast-math} on, is refused with it,
except with GCC when \option{-fno-fast-math} is anywhere on the command line.
With Clang, \option{-ffast-math} and \option{-Ofast} also turn
\option{-frounding-math} off, which \option{-fno-fast-math} does not turn on
again: they have to come before the flags of interval arithmetic, not after
them.

\begin{Rem}\label{rem:gcc12-denorm}
With \option{-frounding-math} on x86, GCC 12.1 to 12.3 and 13.1 to 13.2
(Debian 12 has GCC 12.2) initialize wrongly a double member of a static array
of structures given \code{std::numeric\_limits<double>::denorm\_min()}: they
store the long double $2^{-1074}$ there, whose first eight bytes read $-0$ and
whose last two overwrite the next member (\code{\{ denorm\_min(), DBL\_MAX
\}} read $[-0, \mathrm{NaN}]$, the empty set, on 32-bit x86). A double or an
array of doubles initialized so is right, and so are GCC 12.4, 13.3 and
later. A program built with these releases computes such a value at run time,
as \file{tests/rounding\_direction.cpp} does.\newinvfive
\end{Rem}

The link of the program takes one option too, where the compiler accepts it
(GCC 13 and later on x86, and from 11.4 and 12.4 in the series 11 and 12):
\option{-mno-daz-ftz}, which keeps out of the link the file that sets the
flush-to-zero and denormals-are-zero modes when a program is linked with
\option{-Ofast} (\pxref{sec:ftz}). \code{gaol::gaol} and \file{gaol.pc} carry
it too, where the compiler that built GAOL accepts it: the three builds test
the compiler for it with a link.\newinvfive

\subsection{From CMake}\label{sec:compiling-cmake}
%%.....................
\cindex{CMake!using gaol from}%

A CMake project takes everything from the target \code{gaol::gaol} of the
package of GAOL:

\begin{example}
find_package(gaol REQUIRED)
target_link_libraries(my_target PRIVATE gaol::gaol)
\end{example}

\noindent
\code{gaol::gaol} carries the include directory, the flags above, and the
link option above for the code built by GCC~13 or a later one, or by the
compiler that built GAOL, of its major version and not older: another
compiler, as Clang~18, would stop on it, and so would GCC 12.0 to 12.3 with a
GAOL built by GCC 11.4, GCC having the option from 11.4 and 12.4 only in the
series 11 and 12. It no longer carries, for Visual C++, the definition
\code{\_\_GAOL\_PUBLIC\_\_=} of GAOL~4: GAOL being a static library there, its
declarations carry no \code{\_\_declspec(dllimport)}, and a definition left in
a project changes nothing.\newinvfive{} There is no other library to link:
CORE-MATH is compiled into \file{libgaol} itself. A library whose headers
include those of GAOL, as those of Codac do, links \code{gaol::gaol}
\code{PUBLIC}, so that its own users get the flags, and its CMake package
finds GAOL again (\code{find\_dependency(gaol)}). \file{tests/find\_package}
is a project using an installed GAOL this way.

A project can also build GAOL for itself, with the options it wants:

\begin{example}
include(FetchContent)
FetchContent_Declare(gaol GIT_REPOSITORY
    https://github.com/Jordan08/GAOL.git GIT_TAG master)
FetchContent_MakeAvailable(gaol)
target_link_libraries(my_target PUBLIC gaol::gaol)
\end{example}

\noindent
GAOL is then a target of the project, built with it, and \cmd{cmake --install}
of the project installs it with it, CMake package and \file{gaol.pc}
included; nothing is downloaded beyond the sources of GAOL.
\file{tests/fetch\_content} is a project building GAOL this way. GAOL is built
in Release where the project gives no build type, and
\option{-DGAOL\_DEBUG=ON} builds it for debugging whatever the build type or
the configuration of the project, whose own code keeps them
(\pxref{sec:cmake-options}).\newinvfive{} GAOL, a static library, is compiled
as position-independent code (\option{-fPIC}), so that \code{gaol::gaol} can
be linked into a shared library, such as Python bindings. A project that sets
\code{CMAKE\_POSITION\_INDEPENDENT\_CODE}, \code{ON} or \code{OFF}, before
\code{FetchContent\_MakeAvailable(gaol)}, or on the command line, is followed
instead.

\subsection{From pkg-config}\label{sec:compiling-pkgconfig}
%%..........................
\cindex{pkg-config}%

Each build installs \file{gaol.pc} in the \file{pkgconfig} directory of its
library directory (\file{\var{prefix}/lib/pkgconfig}, or \file{lib64} or
\file{lib/\var{multiarch}} rather than \file{lib} on the systems whose
libraries go there), except the CMake build with Visual C++. Its
\code{Cflags} carries the flags above with the include directory, and its
\code{Libs} GAOL itself and the link option above, where the compiler that
built GAOL accepts it. The variable \env{PKG\_CONFIG\_PATH} names the
directory to pkg-config:

\begin{cmdshell}
% export PKG_CONFIG_PATH=@var^prefix~/lib/pkgconfig
% c++ -std=c++11 -O2 $(pkg-config --cflags gaol) program.cpp \
      $(pkg-config --libs gaol)
\end{cmdshell}

\noindent
Built by CMake with its default options on x86-64 Linux, GAOL gives, on one
line:

\begin{cmdshell}
% pkg-config --cflags --libs gaol
-frounding-math -fno-fast-math -ffp-contract=off
-msse2 -msse3 -mfma -I@var^prefix~/include -L@var^prefix~/lib -lgaol
\end{cmdshell}

\noindent
The option \option{-mno-daz-ftz} comes after \option{-lgaol} where the
compiler that built GAOL takes it: \file{gaol.pc} is written for that
compiler, and a program linked by one that refuses the option, as Clang~18, a
GCC before 11.4 or GCC 12.0 to 12.3, stops on it, and is to be linked without
it. There is no \option{-lm}: GAOL still calls functions of the math library of
the system (those of \file{fenv.h}, \code{sqrt()}, \code{floor()},
\code{fma()} where the processor has no such instruction\ldots), which the
C++ compiler links itself.\newinvfive

\subsection{From meson}\label{sec:compiling-meson}
%%.....................
\cindex{meson!using gaol from}%

In a meson project, \code{dependency('gaol')} finds an installed GAOL. A
meson project can also build GAOL for itself, as a subproject, the sources of
GAOL being in its \file{subprojects/gaol}:\newinvfive

\begin{example}
gaol_dep = dependency('gaol', fallback: ['gaol', 'gaol_dep'])
executable('program', 'program.cpp', dependencies: gaol_dep)
\end{example}

\noindent
\code{dependency()} takes an installed GAOL where it finds one, and builds the
subproject otherwise; \code{subproject('gaol').get\_variable('gaol\_dep')}
always builds it. \code{gaol\_dep} carries the include directory, the flags
above and the link option above, as \file{gaol.pc} does: the flags that the
\file{meson.build} of GAOL gives its own targets do not reach those of the
project. \file{tests/meson\_subproject} is a project building GAOL this way.
GAOL could not be a subproject before GAOL v5: meson refuses the
\code{add\_global\_arguments()} its \file{meson.build} called in a subproject.

\medskip
GAOL is built in \code{release}, as CMake and \cmd{configure} build it brought
into another project, where the project leaves the build type to its default,
\code{debug}. meson 1.8.4 and later give a subproject its own default build
type, and a \loption{buildtype=debug} asked for builds GAOL for debugging.
Earlier ones give it the build type of the project, which would have built
GAOL without optimization, and before meson 1.8 with \code{GAOL\_DEBUGGING}:
there, where the options are those of \code{debug}, GAOL compiles its
libraries as in \code{release} all the same, as it cannot tell the default
from a \loption{buildtype=debug} asked for. \code{-Dgaol:enable-debug=true}
(or \code{default\_options: ['enable-debug=true']} in \code{subproject()} or
\code{dependency()}) builds GAOL for debugging, with any meson. Two cases keep
the build type of the project: Visual C++ with meson before 0.57, whose
\code{debug} gives \option{/RTC1}, which \cmd{cl} refuses with \option{/O2},
and the Visual Studio backend of meson (\loption{backend=vs}), which takes the
build type of the project for every target; the default backend, ninja, is not
one.\newinvfive

\section{The namespaces}\label{sec:namespaces}
%%----------------------
\cindex{namespace!gaol}%
\cindex{namespace!gaol\_core}%
\cindex{namespace!gaol\_ieee1788}%

The type \code{interval} of GAOL and its functions are in the namespace
\code{gaol\_core}. A program names them through one of two namespaces, which
\file{gaol/gaol.h} declares and does not open:\newinvfive

\begin{itemize}
\item[$\bullet$] \textbf{\code{gaol}}, GAOL under its own names, which this
  manual uses, to stay compatible with GAOL~4: a program written for GAOL~4
  keeps its names there and, but for the results that change
  (Chapter~\ref{chap:changes}), its results;
\item[$\bullet$] \textbf{\code{gaol\_ieee1788}}, GAOL under the names of IEEE
  1788-2015 (Chapter~\ref{chap:ieee1788}), to follow the standard: its
  operations have the names, the arguments and the results the standard gives
  them.
\end{itemize}

\noindent
The two namespaces hold the same type and compute with the same functions.
The main difference between them is the power, \code{pow} and \code{pown}
(\pxref{sec:pow-namespaces}), and the names that \code{textToInterval} reads.
A program that opens \code{gaol} writes the operations under their names:

\begin{example}
#include <gaol/gaol.h>
using namespace gaol;

interval x(1, 2);
interval y = pow(x, 3) + sin(x);
\end{example}

\noindent
\file{gaol/gaol.h} no longer opens \code{gaol} itself, as GAOL~4 did: a
program that relied on it adds \code{using namespace gaol;}, or names
\code{gaol::interval}.\newinvfive

\medskip
The two namespaces hold \code{interval}, and the functions that have the same
name and the same meaning in both (\code{sin}, \code{exp}, \code{sqrt},
\code{min}\ldots) are the same functions, those of \code{gaol\_core}.
\code{pow} and \code{textToInterval} are the ones whose meaning differs: each
namespace has its own, and none is in \code{gaol\_core}, the namespace where
argument-dependent lookup looks for a call \code{pow(x,~y)} on an interval, so
that each namespace finds its own only. A program opens one of the two
namespaces, not both: with both open, \code{pow(x,~y)} is ambiguous, and so is
\code{textToInterval(s)}, which reads the names of GAOL in \code{gaol} and
those of the standard in \code{gaol\_ieee1788} (\pxref{sec:constructors}).

\medskip
\code{gaol} names the type and the functions of \code{gaol\_core} one by
one.\newinvfive{} They are the names GAOL~4 had in \code{gaol} that GAOL v5
still names there (the class \code{interval}, the expressions with their nodes
and visitor, the functions and the operators of intervals, the functions of
doubles and the rounding direction, \code{round\_downward()},
\code{round\_upward()}, \code{round\_nearest()}, \code{previous\_float()},
\code{next\_float()}, \code{is\_finite()}, the version of the library,
\code{version\_major}, \code{version\_minor}, \code{version\_micro} and
\code{version}, and the doubles around $\pi$ and $\pi/2$ of
\exref{sec:headers}), the functions on intervals GAOL v5 adds (\code{exp2},
\code{log2}, \code{hypot}, \code{compoundm1}, \code{trunc}, \code{sinpi},
\code{erf}, \code{fma}, \code{cancel\_minus},
\code{round\_ties\_to\_even}\ldots{} and the nodes of their expressions),
\code{restore\_rounding()} and \code{exact\_string()}. \code{gaol\_ieee1788}
names the version as \code{gaol} does.

\noindent
Some names are in \code{gaol\_core} only, and a program names them there:
\code{gaol\_core::rnd\_keep}, for instance.\newinvfive{} They are
\code{rnd\_keep()} (GAOL v5 adds it, and had it in \code{gaol} in its 4.3
releases; GAOL~4 has none), \code{NaN\_val}, \code{eval\_stack}, and the
reader of strings, \code{parse\_interval()} and \code{parsing\_names}.

\noindent
What the code of GAOL uses for itself is in \code{gaol\_core} only too, out of
the way of the names of a program that opens \code{gaol}: the functions saving
and setting the rounding direction and the flush-to-zero modes
(\code{rounding\_state}, \code{get\_rounding()}, \code{set\_rounding()}, and
\code{round\_upward\_if\_needed()}, with \code{gaol\_set\_rounding\_x86()},
which GAOL v5 had in \code{gaol} in its 4.3 releases and GAOL~4 had not, and
\code{get\_flush\_modes()}\ldots), the functions computing the \code{pow} of
\code{gaol} (below), and the internals GAOL~4 had in \code{gaol} without its
manual documenting them: the functions \code{gaol\_signbit()},
\code{modulo\_k\_pi()} and the \code{round\_*\_sse()}, the types
\code{uintdouble}, \code{ullidouble}, \code{Interval\_struct} and
\code{prec\_t}, and the variable \code{the\_null\_expr}. A program calling one
of them names it in \code{gaol\_core}: \code{gaol\_core::get\_rounding()},
for instance.

\medskip
The \code{pow} of \code{gaol} takes the integer power for an integer exponent,
a negative base included, and the \code{pow} of IEEE 1788-2015 for any other
exponent (\pxref{sec:powers}). In \code{gaol\_core}, its versions are named
apart: \code{gaol\_pown()}, \code{gaol\_uipow()}, \code{gaol\_pown\_ll()},
\code{gaol\_pow\_real()}, \code{gaol\_pow\_hybrid()}, \code{gaol\_pown\_exp()}
and \code{gaol\_pow\_exp()}, which are not in \code{gaol}.

\section{The powers in the two namespaces}\label{sec:pow-namespaces}
%%----------------------------------------
\cindex{pow!in the two namespaces}%

The power is where \code{gaol} and \code{gaol\_ieee1788} differ.\newinvfive{}
IEEE 1788-2015 has two powers (Table 9.1): \ieee{pown}$(x,n)$, the integer
power, for an integer $n$, defined for every $x$, $x^0$ being $1$ and $0^n$
having no value for $n<0$; and \ieee{pow}$(x,y)$, the real power
$e^{y\log x}$, for an interval $y$, defined for $x>0$, and for $x=0$ when
$y>0$. GAOL~4 had a single \code{pow}, which took the power of a negative
base, and \code{gaol} keeps it: its \code{pow} is a hybrid of the two.

\begin{itemize}
\item In \code{gaol\_ieee1788}, \code{pown(x, n)} is the integer power, for an
  $n$ of any integer type, and \code{pow(x, y)} the real power: only the part
  of $x$ in \itv{0}{+\infty} counts, whatever $y$, and \code{pow(x, p)} with a
  number $p$ is \code{pow(x, [p])}, an integer $p$ included, as the standard
  gives no other meaning to it.
\item In \code{gaol}, \code{pow(x, y)} takes the integer power where the
  exponent is an integer: an integer type, a double that is an integer, or a
  degenerate interval \itv{n}{n} holding an integer, beyond the \code{int}s
  too; any other exponent takes the real power of the standard. There is no
  \code{pown} in \code{gaol}.
\end{itemize}

\begin{table}[htbp]\centering\renewcommand{\arraystretch}{1.1}
  \begin{tabular}{lll}
    \hline
    & \multicolumn{1}{c}{\code{gaol}} &
    \multicolumn{1}{c}{\code{gaol\_ieee1788}}\\
    \hline\hline
    \code{pow(\itv{-4}{-1}, 2)} &
    \itv{1}{16} &
    $\emptyset$ (\code{pown}: \itv{1}{16})\\
    \code{pow(\itv{-2}{3}, 3)} & \itv{-8}{27} & \itv{0}{27}\\
    \code{pow(\itv{0}{0}, 0)} &
    \itv{1}{1} &
    $\emptyset$ (\code{pown}: \itv{1}{1})\\
    \code{pow(\itv{-2}{2}, -1)} & \itv{-\infty}{+\infty} & \itv{0.5}{+\infty}\\
    \code{pow(\itv{-2}{-1}, \itv{10^{10}}{10^{10}})} &
    \itv{1}{+\infty} &
    $\emptyset$\\
    \code{pow(\itv{-1}{4}, 0.5)} & \itv{0}{2} & \itv{0}{2}\\
    \code{pow(\itv{0}{2}, -2)} & \itv{0.25}{+\infty} & \itv{0.25}{+\infty}\\
    \code{pow(\itv{-4}{-1}, \itv{1.5}{2.5})} & $\emptyset$ & $\emptyset$\\
    \hline
  \end{tabular}
  \caption{The power of \code{gaol} and of \code{gaol\_ieee1788} on some
    arguments}\label{tab:pow-namespaces}
\end{table}

Table~\ref{tab:pow-namespaces} shows the two on some arguments. They agree on
$x\subseteq\itv{0}{+\infty}$ and on every exponent that holds no integer
alone; they differ where $x$ has a negative part and the exponent is an
integer, and at $0^0$. The \code{pow} of \code{gaol} is thus not monotone for
the inclusion of its exponent: \code{pow(\itv{-4}{-1}, \itv{2}{2})} is
\itv{1}{16}, and \code{pow(\itv{-4}{-1}, \itv{1.5}{2.5})}, whose exponent
holds $2$, is empty; the \code{pow} of the standard is. Both give the tightest
bounds, but for the integer power: within the ints it comes from exact
products, and is one double beyond the tightest where it is within twice the
bound of the error of its products, never for the square; beyond the ints it
comes from the \code{pow} of CORE-MATH, and is the tightest for an exponent
that is a double and within three doubles of it otherwise, beyond $2^{53}$
(Chapter~\ref{chap:accuracy}). A degenerate integer exponent of
\code{gaol\_ieee1788::pow} takes the integer power too, with its bounds:
\code{pow([0x1.fffffffffffffp+1], [7])} has its lower bound one double below
the tightest.

\medskip
The exponent 2 is the square of the intervals: \code{pow(x, 2)},
\code{pow(x, [2])}, \code{gaol\_ieee1788::pown(x, 2)} and, on the part of $x$
in \itv{0}{+\infty}, \code{gaol\_ieee1788::pow(x, [2])} are \code{sqr(x)}
itself, with its bounds and the sign of its zero bounds, where the integer
power made its square from exact products, the same bounds in a longer time
and, for some intervals holding 0, a zero bound of the other
sign.\newinvfive{} The exception flags it raises are those of \code{sqr(x)}
too. For the empty set, \code{pow(x, 2)}, \code{pow(x, 2u)},
\code{pow(x, 2LL)}, \code{pow(x, 2.0)}, \code{pow(x, [2])},
\code{gaol\_ieee1788::pown(x, 2)} and \code{gaol\_ieee1788::pow(x, [2])}
return as \code{sqr} does, before any check, and leave the rounding direction
as the program had it, and the modes that flush the subnormal numbers to zero
as it had them too; so do the real and interval forms of \code{pow} for any
exponent, an empty exponent included, and the \code{long long} form beyond
the unsigned ints.\newinvfive{} The integer powers of an \code{int} or of an
\code{unsigned int} other than 2 check the direction first and leave it
upward: a program does not count on the direction after an operation that
computes no bound (\pxref{sec:rounding-direction}).

\medskip
The rest of GAOL follows the same split. The expressions \code{pow(e, n)}
and \code{pow(e1, e2)} of \code{gaol} are the hybrid power, \code{pown(e, n)}
and \code{pow(e1, e2)} of \code{gaol\_ieee1788} the two powers of the
standard. The text \code{"pow([-4,-1],2)"} is read as \itv{1}{16} by the
\code{textToInterval} of \code{gaol}, and as the empty set by that of
\code{gaol\_ieee1788}, which reads \code{pow} as the real power and
\code{"pown([-4,-1],2)"} as the integer power, a name \code{gaol} does not
read. The reverse functions are those of the standard in both namespaces:
\code{pow\_rel} and \code{pow\_exponent\_rel} (\ieee{powRev1} and
\ieee{powRev2}) reverse the real power, and \code{nth\_root\_rel}
(\ieee{pownRev}) the integer power, for every integer exponent, where it took
$n\geq 1$ only.\newinvfive

A program written for GAOL~4, or for IBEX and Codac, whose intervals take the
\code{pow} of GAOL, opens \code{gaol}; a program that wants the results of the
standard, or compares GAOL with other libraries of IEEE 1788-2015, as the
comparison of \exref{sec:comparison} does, opens \code{gaol\_ieee1788}.

\begin{Ex}\label{ex:pow-namespaces}
The three powers of \itv{-4}{-1} with the exponent 2, named in their
namespaces.
\begin{example}
#include <iostream>
#include <gaol/gaol.h>

int main(void)
{
  gaol::interval x(-4, -1);
  std::cout << gaol::pow(x, 2) << ' '
            << gaol_ieee1788::pow(x, 2) << ' '
            << gaol_ieee1788::pown(x, 2) << std::endl;
  gaol::cleanup();
  return 0;
}
@outputs^[1, 16] [empty] [1, 16]~
\end{example}
\end{Ex}

\section{Integers}\label{sec:integers}
%%---------------
\cindex{integers!given to GAOL}%

An integer given to GAOL is taken as the integer it is.\newinvfive{}
\code{interval(n)} is the tightest interval containing $n$: \itv{n}{n} where
$n$ is a double, as every \code{int} is, and the two doubles around it
otherwise. \code{interval(9007199254740993LL)}, $2^{53}+1$, is
\itv{2^{53}}{2^{53}+2}, where its conversion to a double gave
\itv{2^{53}}{2^{53}} or \itv{2^{53}+2}{2^{53}+2}, according to the rounding
direction, neither of which contains it, and \code{interval(ULLONG\_MAX)}
contains $2^{64}-1$. The same holds for \code{interval(a, b)} with an integer
bound, which is empty for $a>b$ compared as numbers
(\code{interval(9007199254740993LL, 9007199254740992LL)} is empty, though both
lie between the same two doubles), for the assignment \code{x = n} and
\code{x |= n}, the arithmetic operators (\code{x + n}, \code{n * x},
\code{x += n}\ldots), \code{set\_contains(n)}, \code{expression(n)}, and
\code{numsToInterval()} and \code{isMember()} in \code{gaol\_ieee1788};
\code{x \&= n} is \code{interval(n)} where \code{x} contains $n$, and the empty
set otherwise. The relations and \code{set\_strictly\_contains(n)} compare $n$
as the integer it is: \code{interval(9007199254740992.0) < 9007199254740993LL},
$\itv{2^{53}}{2^{53}}<2^{53}+1$, is \code{true}, though no double lies between
them. The constructors are detailed in \exref{sec:constructors}, and the
relations in \exref{sec:relational-symbols}.

\medskip
Every integer type is taken so, \code{bool} aside, whatever the rounding
direction; an unscoped enumeration, and \code{\_\_int128} outside the GNU
modes of GCC and Clang, still convert to a double, and a \code{long double}
bound next to an integer takes \code{interval(double, double)}. An integer
that is a double gives the result the double gives, and for the types of 53
bits or less (\code{int}, \code{short}, \code{char}\ldots) the same code.
\code{interval(n)} and \code{expression(n)} are explicit, as their
constructors from a double; \code{expression(n)} is \code{expression(double)}
where $n$ is a double.

\medskip
\code{pow(x, n)} of \code{gaol} and \code{pown(x, n)} of \code{gaol\_ieee1788}
take an exponent of any integer type: \code{pow(x, 5L)} did not compile, a
\code{long} being as far from an \code{int} as from an \code{unsigned} and a
\code{double}, and \code{pown(x, 3000000000u)} converted the exponent to a
negative \code{int}. An exponent within the ints is taken as an \code{int},
one within the unsigned ints as an \code{unsigned}, and one beyond them as a
\code{long long} (\code{gaol\_pown\_ll()}), whose power comes from the
\code{pow} of CORE-MATH at the bounds of $x$, as for \code{gaol::pow(x, [n])}
(\itv{-\infty}{+\infty} for a nonempty $x$ before GAOL v5):
\code{pow(interval(1 + 0x1p-30), 1LL << 33)}, about $e^8$, is
\itv{2980.9579759367944}{2980.9579759367948}. Beyond 64 bits (an
\code{\_\_int128}), the exponent is taken as $\pm(2^{63}-1)$ or
$\pm(2^{63}-2)$, of its sign and parity, the power of every double but $0$,
$\pm 1$ and $\pm\infty$ being beyond the doubles there. The same holds for
\code{nth\_root(x, q)} and \code{rootn(x, q)}, whose order beyond the ints and
the unsigned ints gives an enclosure of the roots, computed by \code{pow} with
the exponent $1/q$ ($x^{1/q}$ is within $2^{-21}$ of $1$ for every double
$x>0$ there), for \code{pownRev(c, x, p)}, which gives the hull of the parts
of $x$ in these roots and, for an even $p$, in their opposites (for a negative
$p$, $c$ is cut at 0 first), and for \code{pow(e, n)}, \code{pown(e, n)} and
\code{nth\_root(e, n)} of the expressions. \code{nth\_root(x, 3L)} was
ambiguous; \code{rootn(x, 3000000000u)} and \code{pownRev(c, 3000000000u)}
converted the order to a negative \code{int}, the first giving a wrong
interval (\itv{0.99999999839}{0.99999999946} for \itv{2}{8}, rather than
\itv{1.00000000023}{1.00000000069}) and the second the reverse of a negative
power; \code{rootn(x, 5000000001LL)} and the powers of expressions cut a
\code{long} to an \code{int}; and \code{nth\_root(e, -2)} took the order
4294967294. Beyond the ints, \code{pow(e, n)} and \code{pown(e, n)} of the
expressions are the \code{pow} of GAOL with the exponent \itv{n}{n}, the
powers of the intervals where $n$ is a double, and beyond $2^{53}$ the
product of $x^h$ and $x^l$, $n=h+l$, $h$ a double and $|l|<2^{11}$, no node
holding a wider integer: wider than the power of the interval where $e$
contains 0; beyond the unsigned ints, \code{nth\_root(e, n)} is the root of
the interval but where $e$ contains 0, where it is wider.

\section{A result thrown away}\label{sec:nodiscard}
%%---------------------------
\cindex{result thrown away}\findex{GAOL\_NODISCARD}%

The functions of GAOL do not change the interval they are given: they return
another one. \code{sqrt(x);} alone leaves \code{x} as it was, and so do
\code{x.mid();} and \code{x.emptyset();}, a static function that returns the
empty set. The functions whose result is all they do carry an attribute that
makes the compiler warn about such a call:\newinvfive{} the functions and
operators on intervals, their predicates, the names of \code{gaol\_ieee1788}
and the functions building expressions.

\begin{example}
sqrt(x);        @rem^// warning: the result is thrown away, x is unchanged~
x = sqrt(x);    @rem^// what was meant~
(void)sqrt(x);  @rem^// thrown away on purpose: no warning, but with GCC before 7~
\end{example}

\noindent
The attribute is \code{GAOL\_NODISCARD}, of \file{gaol/gaol\_config.h}. It is
\code{[[nodiscard]]} in C++17 and later. A program is compiled in the standard
its project sets, or else in the one of its compiler, which is C++14 for GCC 6
to 10, Clang 6 to 15 and Visual C++. Before C++17, \code{GAOL\_NODISCARD} is
\code{[[nodiscard]]} too where the compiler takes it there: GCC 7 and later,
and Visual C++ 2019 16.4 and later, whose warning is C4834. Clang, which warns
about \code{[[nodiscard]]} before C++17 with \option{-pedantic}, and GCC
before 7 take \code{\_\_attribute\_\_((warn\_unused\_result))}, whose warning
is \option{-Wunused-result} (on by default); older Visual C++ takes
\code{\_Check\_return\_} of \file{<sal.h>}, which only the code analysis
reports (\option{/analyze}, warning C6031), not the compiler. A cast to void
(\code{(void)sqrt(x);}) silences each of them, but the attribute of GCC
before 7: the result is kept in a variable there. A program that does not
want the warnings defines \code{GAOL\_NODISCARD} empty before including GAOL
(\option{-DGAOL\_NODISCARD=}). The compound assignments (\code{x += y}), which
do change \code{x}, do not carry it.

\section{Tools of interval algorithms}\label{sec:tools}
%%-----------------------------------
\cindex{tools of interval algorithms}%
\findex{width\_enclosure}\findex{inflate}\findex{midrad}\findex{bisect}%
\findex{is\_bisectable}%

The algorithms of interval computing, the enclosure of a range by
subdivision, the interval Newton method, the contractors, the branch and
bound, are written with a few operations beyond the arithmetic, which they
otherwise write for themselves. GAOL v5 provides them.\newinvfive

\begin{itemize}
\item[$\bullet$] \textbf{The width.} \code{x.width\_enclosure()} is the
  interval \code{interval(r) - interval(l)} enclosing the exact width $r-l$ of
  $x=\itv{l}{r}$, whose upper bound is \code{width()}, the width rounded
  upward; it is \itv{\mathrm{MAX}}{+\infty} for an unbounded $x$, as
  \code{textToInterval()} reads \code{inf}, and the empty set for the empty
  set.
\item[$\bullet$] \textbf{The inflation.} \code{x.inflate(r)} is $x$ widened by
  the absolute radius $r$, \code{x + interval(-r, r)}, each bound rounded
  outward; $r=+\infty$ gives the whole line, and $r<0$ or a NaN $r$ the empty
  set.
\item[$\bullet$] \textbf{A midpoint and a radius.} \code{interval::midrad(m, r)}
  is the interval of midpoint $m$ and radius $r$,
  \code{interval(m) + interval(-r, r)}, which encloses $[m-r,\,m+r]$, where
  \code{interval(m - r, m + r)} does not: both bounds are then computed in the
  upward rounding of GAOL, the lower one above $m-r$ where it is inexact. An
  infinite $m$ gives the empty set, as \code{interval(m)} does.
\item[$\bullet$] \textbf{The bisection.} \code{x.bisect(ratio)} is the pair of
  the parts of $x$ on either side of a cut point at the given ratio of its
  width, $0<\mathit{ratio}<1$, which share the cut point and cover $x$;
  \code{bisect(0.5)} cuts where \code{split()} does, at \code{midpoint()}, and
  so does every ratio for a half-line or the whole line. A ratio outside
  $(0,1)$, or NaN, throws \code{std::domain\_error}. \code{x.is\_bisectable()}
  tells whether a double lies strictly between the bounds, both parts being
  then smaller than $x$: not for two consecutive doubles, nor for
  \itv{\mathrm{MAX}}{+\infty} (\pxref{sec:splitting}).
\item[$\bullet$] \textbf{The hull and the intersection.} \code{hull(x, y)} and
  \code{intersect(x, y)} are \code{x | y} and \code{x \& y} as functions, in
  \code{gaol} as in \code{gaol\_core}.
\end{itemize}

\begin{Ex}\label{ex:tools}
The program below uses these tools on an enclosure of $0.1$, on the interval
of midpoint $1$ and radius $10^{-20}$, and on \itv{0}{1} cut at a quarter of
its width.
\begin{example}
#include <iostream>
#include <gaol/gaol.h>
using namespace gaol;

int main()
{
  const interval x = textToInterval("0.1");        @rem^// encloses 0.1~
  const interval y = interval::midrad(1.0, 1e-20); @rem^// [1 - 1e-20, 1 + 1e-20]~
  const std::pair<interval, interval> parts = interval(0.0, 1.0).bisect(0.25);
  std::cout << x.width_enclosure() << ' ' << y << ' ' << parts.first << ' '
            << parts.second << ' ' << hull(x, y) << std::endl;
  gaol::cleanup();
  return 0;
}
\end{example}
\end{Ex}

\paragraph{The reverse functions.}
A contractor shrinks a box without losing a solution of a constraint: it
evaluates the expression forward, from the variables to the result, and goes
backward, each node being intersected with what the others allow. The
backward operations are the relational functions of GAOL, the \code{*\_rel}
functions, which reverse each operation (\pxref{sec:relational-arithmetic}).
Their arguments are in one order, the result first and the operand contracted
last: \code{div\_rel(c, b, x)} contracts a factor of $z=x\cdot y$, and
\code{sqrt\_rel(c, x)} the argument of a square root. The reverse functions of
the power follow it: \code{pow\_rel(K, J, I)} is the hull of the $x$ of $I$
with an $x^y$ in $K$ for a $y$ of $J$, and \code{pow\_exponent\_rel(K, H, I)}
the hull of the $y$ of $I$ with an $x^y$ in $K$ for an $x$ of $H$; they are
\ieee{powRev1}$(b,c,x)$ and \ieee{powRev2}$(a,c,x)$ of IEEE 1788-2015 (Table
10.1), as \code{pow\_rel(c, b, x)} and \code{pow\_exponent\_rel(c, a, x)}, $x$
being \itv{-\infty}{+\infty} when left out.\newinvfive{} Both reverse the
\code{pow} of the standard, that of \code{gaol\_ieee1788::pow}, in either
namespace: $x^y$ is $e^{y\log x}$ for $x>0$, $0$ for $x=0$ and $y>0$, and has
no value elsewhere, so that no $x<0$ is a solution, and
\code{pow\_rel([1, 4], [2], [-3, 3])} is \itv{1}{2}. The \code{pow} of
\code{gaol} takes the integer power for a degenerate integer exponent
\itv{n}{n}, which has negative bases too: it is not one function of $x$ and
$y$, and the reverse of $x^n$ is \code{nth\_root\_rel(K, n, I)}, which takes
every integer $n$ (\pxref{sec:relational-arithmetic}).

\medskip
\code{max\_rel(z, y, x)}, \code{min\_rel(z, y, x)}, \code{sign\_rel(z, x)} and
\code{floor\_rel(z, x)} contract $x$ under the constraints $z=\max(x,y)$,
$z=\min(x,y)$, $z=\mathrm{sign}(x)$ and $z=\lfloor x\rfloor$.\newinvfive{}
Each returns the hull of the $x$ of \code{x} that the function sends into
\code{z}, with some $y$ of \code{y} for \code{max} and \code{min}: the reverse
functions IEEE 1788-2015 defines for every function (10.5.4), with the
arguments in the order of the relational functions of GAOL. \code{max} and
\code{min} being symmetric, \code{max\_rel(z, x, y)} contracts $y$, and
\code{max(x, y) \& z} contracts $z$. IEEE 1788-2015 names and requires only
the reverse functions of its Table 10.1, which these four are not in:
\code{gaol\_ieee1788} has no name for them, and a program that opens it calls
them as \code{max\_rel(z, y, x)} all the same, argument-dependent lookup
finding them in \code{gaol\_core} (\pxref{sec:namespaces}).

\begin{example}
interval x(-5.0, 5.0), z(1.0, 2.0);
interval a = max_rel(z, interval(0.0, 1.0), x);   @rem^// [-5, 2]: y meets z~
interval b = max_rel(z, interval(-1.0, 0.0), x);  @rem^// [1, 2]: y below z~
interval c = max_rel(z, interval(3.0, 4.0), x);   @rem^// empty: y above z~
interval s = sign_rel(interval(-1.0), x);         @rem^// [-5, 0]~
interval f = floor_rel(interval(0.0, 2.0), x);    @rem^// [0, 3]~
\end{example}

\medskip
The interval Newton method encloses the roots of $f$ in $X$ with the
relational division, which solves $(x-m)\,d=-f(m)$ for $d$ in the enclosure of
$f'(X)$ where the division \code{/} does not (\exref{sec:common-errors}, and
\file{examples/05\_interval\_newton.cpp}); dividing by the negative and by the
positive part of $f'(X)$ apart, two \code{div\_rel}, or
\code{mulRevToPair(b, c)}, keeps the gap around $m$. The examples 06 and 07
of \file{examples/} write a branch and bound and a contractor this way, with
\code{bisect()}, \code{is\_bisectable()} and the hull.

\section{Errors, in short}\label{sec:errors-overview}
%%-----------------------
\cindex{errors!overview}%

GAOL reports an error by throwing an exception of the standard library
(\file{<stdexcept>}), where GAOL~4 had classes of its own:\newinvfive{}
\code{std::invalid\_argument} for a string that is no interval, given to
\code{gaol::textToInterval()} or read by \code{operator>{}>}, and
\code{std::domain\_error} for a function called with an argument it does not
take, \code{nb\_fp\_numbers()} with a NaN, \code{bisect()} with a ratio
outside $(0,1)$, or \code{nth\_root(8, 1.5)} in a string. Both derive from
\code{std::logic\_error}, which derives from \code{std::exception}, so that a
handler that knows nothing of GAOL says what went wrong: its \code{what()} is
\code{\var{file}:\var{line}: \var{explanation}}, the name of the source file
of GAOL without its directories and the line that threw, then what went
wrong. Where memory runs out, GAOL throws the \code{std::bad\_alloc} of C++,
the reader of strings included. A build without
\code{GAOL\_EXCEPTIONS\_ENABLED} prints a message and aborts instead of
throwing; the three builds can be asked to do so (a setting of the developers). The classes of GAOL~4, what
became of each, the messages, the warnings and the builds that abort are
detailed in \exref{sec:exceptions}.

\section{Common errors}\label{sec:common-errors}
%%---------------------
\cindex{common errors}%

This section reviews the errors most often made with GAOL, so that the user
recognizes them in a program. Most of them come from the two facts that a
double is not an interval, and that an interval is not a number.

\subsection{Floating-point arithmetic and rounding}
%%.................................................
\cindex{Rounding!to nearest}%

Programming with floating-point numbers is one of the few activities where
the compiler is to be watched. Let us consider the following piece of code:

\begin{example}
#include <gaol/gaol>  @rem^// We do not import the gaol namespace~

using gaol::interval;

int main(void) {
  interval one_tenth(0.1); @rem^BEWARE: wrong !~

  @emph^[some code using one_tenth]~
  gaol::cleanup();
}
\end{example}

\noindent
Though a rational perfectly representable in decimal, \code{0.1} is not
representable in binary with a finite number of bits, and thus requires
rounding. The purpose of the user was to define an interval containing this
value. However, \code{0.1} is rounded \emph{at compile time}, most certainly
to the nearest representable floating-point number \roundNearest{0.1}
(\pxref{sec:rounding-definitions}). As a consequence, \code{one\_tenth} is a
degenerate interval containing only \roundNearest{0.1}, and the containment
property (\pxref{sec:containment}) is violated.

The right way to deal with rational constants that might not be perfectly
representable as floating-point numbers is to write them as text, and to read
them with \code{gaol::textToInterval()} (\pxref{sec:constructors}), so that
they are correctly rounded downward and upward at \emph{run time}:

\begin{example}
#include <gaol/gaol>

using gaol::interval;

int main(void)
{
  @rem^// OK: this is the right way~
  interval one_tenth = gaol::textToInterval("0.1");
  @rem^// Another possible way~
  interval one_tenth2 = gaol::textToInterval("1/10");

  @emph^[some code using one_tenth]~
  gaol::cleanup();
  return 0;
}
\end{example}

\noindent
Now, \code{one\_tenth} is the smallest floating-point interval enclosing
$0.1$. An interval like this one, containing at most two consecutive
floating-point numbers, is a canonical interval
(\pxref{sec:intervals}).
\cindex{interval!canonical|see{canonical interval}}%
\cindex{canonical interval}%
Written in hexadecimal (\pxref{sec:output-format}), which shows the bounds
exactly, the two intervals are:

\begin{onscreen}
interval(0.1) = [0x1.999999999999ap-4]
one_tenth     = [0x1.9999999999999p-4, 0x1.999999999999ap-4]
\end{onscreen}

\subsection{Functions of numbers}
%%...............................
\cindex{sqrt@\code{sqrt} of a number}%

A function called with a number is the function of the C library, even where
the namespace \code{gaol} is open: \code{sqrt(2.0)} is a \code{double},
rounded in the direction in effect, and no enclosure of $\sqrt{2}$. A
\code{double} no longer becomes an interval by itself, the constructor being
explicit, so that \code{interval r = sqrt(2.0);} does not compile; but
\code{interval r(sqrt(2.0));} does, and gives a point that does not contain
$\sqrt{2}$. The enclosure is the function of GAOL on an interval, or the text
read by \code{textToInterval()}:\newinvfive

\begin{example}
interval root2 = textToInterval("sqrt(2)");
cout << boolalpha << interval(sqrt(2.0)).set_contains(root2) << ' '
     << sqrt(interval(2.0)).set_contains(root2);
@outputs^false true~
\end{example}

\noindent
For the same reason, \code{gaol::sin(0.5)} does not compile: it is written
\code{gaol::sin(interval(0.5))}, and \code{sin(0.5)} alone is the \code{sin}
of C. Writing code inside \code{namespace gaol \{\}} has the same effect: the
using-declarations of \code{gaol} are visible there, a call \code{sqrt(2.0)}
finds \code{gaol::sqrt}, which hides \code{::sqrt}, and does not compile. This
was already the case with GAOL~4.

\subsection{Midpoint and radius}
%%..............................
\cindex{interval!midpoint and radius}%

\code{interval(m - r, m + r)}, with doubles, is not an interval enclosing
$[m-r,\,m+r]$: the program computes \code{m - r} and \code{m + r} itself, in
the direction in effect, which is upward until \code{cleanup()}, so that
\code{m - r} may be above the real $m-r$, and to nearest each bound may be on
either side. With \code{m = 1.0} and \code{r = 1e-20}, the interval is
\itv{1}{1+2^{-52}} while the direction is upward, and does not contain
$1-10^{-20}$. \code{interval::midrad(m, r)} (\pxref{sec:constructors}),
or \code{m + interval(-r, r)}, computes both bounds in intervals, and gives
\itv{1-2^{-53}}{1+2^{-52}}, which does.\newinvfive

\begin{example}
double m = 1.0, r = 1e-20;
interval wrong(m - r, m + r);  @rem^// Does not enclose [m - r, m + r]~
interval right = interval::midrad(m, r);  @rem^// Does~
\end{example}

\subsection{The doubles of the program}
%%.....................................

The rounding direction that GAOL leaves applies to the rest of the program,
which computes its own doubles in it: after an operation of GAOL,
\code{1.0 / 3.0} is the double above $1/3$, and \code{printf()} and the
streams of the GNU C library write the decimal digits of a double rounded
upward too, until \code{cleanup()}. \exref{sec:upward-effects} gives what
changes, and what to do.\newinvfive

\subsection{Zero}
%%...............
\cindex{interval!zero}%

\code{interval~x(0,~0)}, \code{interval~x(0)} and \code{interval::zero()}
are \itv{0}{0}. The literal \code{0} is also a null pointer, and the
constructors from one and from two strings of GAOL~4 made these calls
ambiguous; GAOL v5 has no constructor from strings.\newinvfive

\subsection{The sign of a zero bound}\label{sec:zero-sign}
%%...................................
\cindex{zero!sign of a bound}%

A bound that is zero has a sign, which differs between the builds:
\code{sqr(interval(-1, 2))} is \itv{-0}{4} with the SSE2 intervals, which the
default build uses on x86 processors, and with the NEON intervals, which it
uses on 64-bit ARM processors, and \itv{0}{4} with the FPU intervals, those of
the other processors and of the builds with \option{GAOL\_SIMD} off (Visual
C++ and clang-cl have the SSE2 intervals on x86 and on x64);\newinvfive{} \code{interval::zero()} is \itv{-0}{0} with the first
two and \itv{0}{0} with the third. \code{pow(x, 2)} is \code{sqr(x)}, and has
the same sign.\newinvfive{} The NEON and the SSE2 intervals can differ too:
the maximum of 64-bit ARM, which the product takes of its four products, ranks
$+0$ above $-0$, where the one of x86 returns its second operand, so that the
lower bound of \code{interval(-1, 0) * interval(-1, 0)} is $-0$ with the NEON
intervals and $+0$ with the others (\pxref{sec:neon}); the sign of a zero
bound of a product of the FPU intervals is the one the compiler gives
(\pxref{sec:fpu-intervals}). IEEE 1788-2015 and GAOL take the two zeros for
the same number: \code{x.left() == 0.0} is \code{true} whichever it is, and a
point interval at zero is written \code{[0]}. A program that reads the sign,
with \code{std::signbit(x.left())} or \code{1.0 / x.left()}, which is
$-\infty$ or $+\infty$, gets one answer with a build and the other with the
other: it compares the bounds with \code{==}, \code{<} and
\code{<=}.\newinvfive

\subsection{The empty set and the domain of a function}
%%.....................................................
\cindex{empty set!in tests}%
\cindex{domain of a function}%

A function is computed on the part of its argument that lies in its domain
(IEEE 1788-2015), and GAOL, which has no decorations, says nothing of the
rest: \code{sqrt(interval(-1, 4))} is \itv{0}{2}, \code{log(interval(-1, 0))}
is empty, and so are \code{interval(1, 0)} and the interval of one infinite
or NaN double. The empty set then passes the tests as if it were a good
answer: it has no point that contradicts a certainly relation, and it is
included in every interval, so that a box outside the domain of a function is
both inside and outside the set a paving looks for:

\begin{example}
interval e = sqrt(interval(-2.0, -1.0)), z(0.0, 1.0);
cout << boolalpha << e.is_empty() << ' ' << (e < z) << ' '
     << (e > z) << ' ' << z.set_contains(e) << ' '
     << e.set_disjoint(z);
@outputs^true true true true true~
\end{example}

\noindent
A program tests \code{is\_empty()} first, and decides itself what a box
outside the domain of its function means, as \file{examples/08\_sivia.cpp}
does.\newinvfive

\subsection{Relations, sorting and containers}
%%............................................
\cindex{relation!symbols!pitfalls}%

The symbols \code{<}, \code{<=}, \code{>} and \code{>=} are the certainly
relations of IEEE 1788-2015, which do not tell which of two intervals is the
greater, and there is no \code{==}: \code{interval(1,3) < interval(2,4)} and
\code{interval(1,3) >= interval(2,4)} are both \code{false}
(\pxref{sec:relational-symbols}). In C-XSC and in PROFIL/BIAS, \code{x <= y}
means that \code{x} is a subset of \code{y}, and \code{r <= x} that the number
\code{r} is in \code{x}: code ported from them compiles, and changes its
meaning (\pxref{sec:other-libraries}). None of them is an ordering for the
standard library: \code{std::sort}, \code{std::set} and \code{std::max} take
the comparator \code{gaol::lexicographic\_less}
(\pxref{sec:lexicographic-less}).\newinvfive

\subsection{Both namespaces}
%%..........................

A program opens the namespace \code{gaol} or the namespace
\code{gaol\_ieee1788}, not both (\pxref{sec:namespaces}):

\begin{example}
#include <gaol/gaol.h>
using namespace gaol;
using namespace gaol_ieee1788;

gaol::interval x(1,2), y = pow(x,x); @rem^// Does not compile~
\end{example}

\noindent
The compiler reports that the call of \code{pow} is ambiguous. Not every call
is: with an \code{int} or an \code{unsigned int} for exponent, \code{pow(x, 2)}
compiles, and takes the \code{pow} of \code{gaol}, silently, where that of
\code{gaol\_ieee1788} gives another result for a negative base
(\pxref{sec:pow-namespaces}). Only the calls with another type of exponent,
\code{pow(x, 2.0)}, \code{pow(x, 2L)} or \code{pow(x, y)}, are
ambiguous:\newinvfive

\begin{example}
#include <iostream>
#include <gaol/gaol.h>
using namespace gaol;
using namespace gaol_ieee1788;

int main(void)
{
  gaol::interval x(-4,-1);
  std::cout << pow(x, 2) << std::endl; @rem^// The pow of gaol~
  gaol::cleanup();
  return 0;
}
@outputs^[1, 16]~
\end{example}

\subsection{Bisection}
%%....................
\cindex{bisection!of an unbounded interval}%

A loop that cuts an interval until it is narrow,
\code{while (x.width() > eps)}, never ends on an unbounded interval:
\code{bisect(0.5)}, like \code{split()}, cuts \itv{1}{+\infty} at
\code{DBL\_MAX}, and the half \itv{\mathrm{DBL\_MAX}}{+\infty} is cut at
\code{DBL\_MAX} again, its width staying infinite. What a loop stops on is
\code{is\_bisectable()}, which is \code{false} for this half, and for an
interval of two consecutive doubles, the cut of which gives the interval
itself as one half (\pxref{sec:splitting}):\newinvfive

\begin{example}
interval x(1.0, GAOL_INFINITY);
x = x.bisect(0.5).second;
cout << boolalpha << x << ' ' << x.width() << ' '
     << x.is_bisectable();
@outputs^[1.797693134862315e+308, inf] inf false~
\end{example}

\subsection{Newton's method: the division and the relational division}
%%....................................................................
\cindex{Newton's method!division}%

The step $N(X)=m-f(m)/F'(X)$ of the interval Newton method divides with
\code{/}, the set of the quotients $a/b$ of a point of the numerator by a point
other than 0 of the denominator. Where $0\in F'(X)$ and $f(m)=0$, the quotient
$0/F'(X)$ is $\itv{0}{0}$, so that $N(X)=m$ lies inside $X$ and ``proves''
that $X$ holds exactly one root, where it holds seven for $\sin x-x/10$ on
\itv{-10}{10}. A root satisfies a relation, $(x-m)\,d=-f(m)$ for some $d$ in
$F'(X)$, which \code{\%}, or \code{div\_rel()}, the reverse of the product
(\pxref{sec:relational-arithmetic}), solves: it gives the whole line, rightly,
when $d$ and $f(m)$ can both be 0.

\begin{example}
interval fm(0.0), D(-1.1, 0.9);
cout << fm / D << ' ' << fm % D << ' '
     << div_rel(fm, D, interval(-5.0, 5.0));
@outputs^[0] [-inf, inf] [-5, 5]~
\end{example}

\noindent
The functional division is the right one where $F'(X)$ holds no 0.
\file{examples/05\_interval\_newton.cpp} shows the trap and the extended step.

\subsection{Reading text}
%%.......................
\cindex{textToInterval@\code{textToInterval}!errors}%

\code{gaol::textToInterval()} throws \code{std::invalid\_argument} for a text
that is no interval, and \code{std::domain\_error} for a text that is one but
calls a function with an argument it does not take, such as
\code{nth\_root(8, 1.5)}, whose order is not an integer: a loop that catches
the first lets the second escape. Both derive from \code{std::logic\_error}
(Chapter~\ref{chap:errors}).\newinvfive

\begin{example}
try {
  cout << textToInterval("nth_root(8, 1.5)");
} catch (const std::invalid_argument&) {
  cout << "std::invalid_argument";
} catch (const std::domain_error&) {
  cout << "std::domain_error";
}
@outputs^std::domain_error~
\end{example}

\noindent
\code{gaol\_ieee1788::textToInterval()} throws neither: it gives the empty set
for every text that is no interval, as IEEE 1788-2015 has it, so that a
malformed text cannot be told from \code{"[empty]"}:

\begin{example}
cout << gaol_ieee1788::textToInterval("[1, 2") << ' '
     << gaol_ieee1788::textToInterval("[empty]");
@outputs^[empty] [empty]~
\end{example}

\noindent
A null pointer is no text, and the \code{textToInterval()} of both namespaces
and \code{exactToInterval()} throw \code{std::invalid\_argument} for it: none
of the three calls below reaches the reader.

\begin{example}
textToInterval(nullptr);
textToInterval(std::getenv("UNSET"));
gaol_ieee1788::exactToInterval(nullptr);
\end{example}

\noindent
Their argument is a \code{const std::string\&}, which the call builds from the
pointer first, and the C++ standard leaves that undefined (libstdc++ throws
\code{std::logic\_error}); the overloads of a C string take the null pointer
before.\newinvfive

\medskip
Two more texts deserve a word. A sign next to an uncertain number,
\code{10-3.56?1} or \code{10+3.56?1}, throws \code{std::invalid\_argument}: the
lexer reads \code{-3.56?1} (and \code{+3.56?1}) as one token, and
\code{10 -3.56?1} throws too. A blank is needed after the sign:
\code{10 - 3.56?1}, or \code{10- 3.56?1}. Reading a long number takes a time
quadratic in its number of digits: \code{"0."} followed by 20\,000 digits
\code{3} takes 0.017 s, 80\,000 digits 0.28 s, 320\,000 digits 4.4 s (about
45 s expected for a million), and the uncertain form 2.3 times as long. A text
coming from elsewhere can thus block a program.

\subsection{Threads and the global settings}
%%..........................................
\cindex{threads!output of intervals}%
\cindex{interval::format@\code{interval::format}!threads}%

The output format of intervals and their number of digits, set by
\code{interval::\allowbreak format()} and \code{interval::\allowbreak precision()}, are variables of
the program, shared by its threads, and not atomic: setting one while another
thread writes an interval, with \code{operator<{}<}, the conversion to a
\code{std::string} or \code{gaol\_ieee1788::intervalToText()} (which reads
the precision only), is a data race, with the undefined behaviour of C++
(ThreadSanitizer reports it). The program that sets them synchronizes with
the threads that write: it sets them before starting its threads, or takes a
lock of its own around the setting and around each output. Without a setting
in between, several threads write intervals at once safely, and so they read
strings (\pxref{sec:reading-intervals}); \code{exact\_string()} reads neither
setting.\newinvfive

\subsection{The flags of interval arithmetic}
%%...........................................

Compiled without the flags of interval arithmetic (\pxref{sec:compiling}), the
inline operations of GAOL may be evaluated by the compiler at compile time,
rounded to nearest, or contracted into fused multiply-adds, and give bounds
that do not enclose the exact results. \file{gaol/gaol\_config.h} stops the
compilation where it can tell:

\begin{onscreen}
% c++ -ffast-math -c program.cpp
gaol/gaol_config.h: error: #error "GAOL cannot be
compiled with -ffast-math (nor -Ofast): the bounds it
computes would not enclose the exact results (remove
the option, or put the flags of gaol.pc and gaol::gaol
after it)"
\end{onscreen}

\noindent
It stops \option{-ffinite-math-only} in the same way, with which the compiler
takes NaN, the bounds of the empty interval, and infinities never to occur
(\pxref{sec:refused}). Taking the flags from \code{gaol::gaol} or
\file{gaol.pc} avoids the question. The link is not seen by the header: a
program linked with \option{-Ofast} gets the flush-to-zero mode, which makes
the bounds of the operations with subnormal operands and results wrong where
GAOL does not clear it
(\pxref{sec:ftz}).\newinvfive

\section{Coming from another library}\label{sec:other-libraries}
%%--------------------------------------
\cindex{libraries!names in other}%
\cindex{IBEX}%
\cindex{Codac}%
\cindex{C-XSC}%
\cindex{PROFIL/BIAS}%
\cindex{Boost.Interval}%

The names of interval libraries differ, and the same symbol can mean another
relation. Tables~\ref{tab:other-names-ibex} and \ref{tab:other-names-cxsc}
give, for what an algorithm uses, the name of GAOL v5, the one of
\code{gaol\_ieee1788}, and those of IBEX and Codac, which share them, of
C-XSC and of Boost.Interval. A dash means that the sources read have no
function of the kind, not that the algorithm cannot be written.\newinvfive

{\renewcommand{\arraystretch}{1.1}\footnotesize\setlength{\tabcolsep}{3pt}
\begin{longtable}{L{.17\textwidth}L{.24\textwidth}%
    L{.22\textwidth}L{.29\textwidth}}
    \caption{The names of GAOL v5, of \code{gaol\_ieee1788}, and of IBEX and
      Codac}\label{tab:other-names-ibex}\\
    \hline
    & GAOL v5 & \code{gaol\_ieee1788} & IBEX, Codac\\
    \hline\hline
    \endfirsthead
    \hline
    & GAOL v5 & \code{gaol\_ieee1788} & IBEX, Codac\\
    \hline\hline
    \endhead
    Lower bound & \code{x.left()} & \code{inf(x)} & \code{x.lb()}\\
    Upper bound & \code{x.right()} & \code{sup(x)} & \code{x.ub()}\\
    Midpoint, a double & \code{x.midpoint()} & \code{mid(x)} & \code{x.mid()}\\
    Midpoint, an interval & \code{x.mid()} & --- & ---\\
    Width &
    \code{x.width()}, rounded upward; \code{x.width\_enclosure()} &
    \code{wid(x)} &
    \code{x.diam()}\\
    Radius & \code{x.rad()} & \code{rad(x)} & \code{x.rad()}\\
    Smallest, greatest absolute value &
    \code{x.mig()}, \code{x.mag()} &
    \code{mig(x)}, \code{mag(x)} &
    \code{x.mig()}, \code{x.mag()}\\
    Hull, intersection &
    \code{x | y}, \code{x \& y}; \code{hull(x, y)}, \code{intersect(x, y)} &
    \code{convexHull(x, y)}, \code{intersection(x, y)} &
    \code{x | y}, \code{x \& y}\\
    Is empty & \code{x.is\_empty()} & \code{isEmpty(x)} & \code{x.is\_empty()}\\
    Subset &
    \code{x.set\_leq(y)} &
    \code{subset(x, y)} &
    \code{x.is\_subset(y)}\\
    Interior &
    \code{x.set\_le(y)}; \code{y.set\_strictly\_contains(x)} &
    \code{interior(x, y)} &
    \code{x.is\_interior\_subset(y)}\\
    Contains a number &
    \code{x.set\_contains(d)} &
    \code{isMember(d, x)} &
    \code{x.contains(d)}\\
    Disjoint &
    \code{x.set\_disjoint(y)} &
    \code{disjoint(x, y)} &
    \code{x.is\_disjoint(y)}\\
    Equal & \code{x.set\_eq(y)} & \code{equal(x, y)} & \code{x == y}\\
    Inflate by \code{r} & \code{x.inflate(r)} & --- & \code{x.inflate(r)}\\
    Bisect &
    \code{x.bisect(ratio)} &
    --- &
    \code{x.bisect(ratio)}, ratio 0.5 (IBEX) or 0.49 (Codac)\\
    Two-piece division &
    two \code{div\_rel}, or \code{mulRevToPair(b, c)} &
    \code{mulRevToPair(b, c)} &
    \code{div2} (IBEX only)\\
    \hline
  \end{longtable}}

{\renewcommand{\arraystretch}{1.1}\footnotesize\setlength{\tabcolsep}{3pt}
\begin{longtable}{L{.19\textwidth}L{.25\textwidth}%
    L{.25\textwidth}L{.22\textwidth}}
    \caption{The names of GAOL v5, of C-XSC and of
      Boost.Interval}\label{tab:other-names-cxsc}\\
    \hline
    & GAOL v5 & C-XSC & Boost.Interval\\
    \hline\hline
    \endfirsthead
    \hline
    & GAOL v5 & C-XSC & Boost.Interval\\
    \hline\hline
    \endhead
    Lower bound & \code{x.left()} & \code{Inf(x)} & \code{lower(x)}\\
    Upper bound & \code{x.right()} & \code{Sup(x)} & \code{upper(x)}\\
    Midpoint, a double &
    \code{x.midpoint()} &
    \code{mid(x)} &
    \code{median(x)}\\
    Width, rounded upward &
    \code{x.width()} &
    \code{diam(x)} &
    \code{width(x)}\\
    Radius & \code{x.rad()} & --- & ---\\
    Smallest, greatest absolute value &
    \code{x.mig()}, \code{x.mag()} &
    \code{AbsMin(x)}, \code{AbsMax(x)} &
    \code{norm(x)}, the greatest only\\
    Hull, intersection &
    \code{x | y}, \code{x \& y}; \code{hull(x, y)}, \code{intersect(x, y)} &
    \code{x | y}, \code{x \& y} &
    \code{hull(x, y)}, \code{intersect(x, y)}\\
    Is empty & \code{x.is\_empty()} & \code{IsEmpty(x)} & \code{empty(x)}\\
    Subset & \code{x.set\_leq(y)} & \code{x <= y} & \code{subset(x, y)}\\
    Interior & \code{x.set\_le(y)} & \code{x < y}, \code{in(x, y)} & ---\\
    Contains a number &
    \code{x.set\_contains(d)} &
    \code{d <= x}, \code{in(d, x)} &
    \code{in(d, x)}\\
    Disjoint &
    \code{x.set\_disjoint(y)} &
    \code{Disjoint(x, y)} &
    \code{!overlap(x, y)}\\
    Equal & \code{x.set\_eq(y)} & \code{x == y} & \code{equal(x, y)}\\
    Inflate by \code{r} &
    \code{x.inflate(r)} &
    \code{Blow(x, eps)}, relative to the width &
    \code{widen(x, r)}\\
    Bisect &
    \code{x.bisect(ratio)} &
    --- &
    \code{bisect(x)}, at the median only\\
    \hline
  \end{longtable}}

\noindent
\code{overlap(x, y)} of Boost.Interval is a boolean, true when $x$ and $y$
have a point in common: \code{!x.set\_disjoint(y)} in GAOL v5. The
\code{overlap} of GAOL v5 (\code{x.overlap(y)},
\code{gaol\_ieee1788::overlap(x, y)}) is the function of IEEE 1788-2015 and
returns an \code{overlapping\_state}, which does not convert to \code{bool}: a
Boost program that tests \code{overlap(x, y)} does not compile with it.

\medskip
What the symbols \code{<}, \code{<=} and \code{==} mean is where code ported
from another library compiles and changes meaning
(Table~\ref{tab:other-symbols}).

{\renewcommand{\arraystretch}{1.1}\footnotesize\setlength{\tabcolsep}{3pt}
\begin{longtable}{L{.17\textwidth}L{.27\textwidth}%
    L{.27\textwidth}L{.19\textwidth}}
    \caption{The meaning of the symbols \code{<}, \code{<=} and \code{==} in
      each library}\label{tab:other-symbols}\\
    \hline
    & \code{x < y} & \code{x <= y} & \code{x == y}\\
    \hline\hline
    \endfirsthead
    \hline
    & \code{x < y} & \code{x <= y} & \code{x == y}\\
    \hline\hline
    \endhead
    GAOL v5 &
    every point of \code{x} is below every point of \code{y}
    (\ieee{strictPrecedes}) &
    every point of \code{x} is at most every point of \code{y}
    (\ieee{precedes}) &
    not defined: \code{x.set\_eq(y)}\\
    IBEX &
    not defined &
    not defined &
    the same bounds, the empty set equal to itself\\
    Codac &
    a \code{BoolInterval}: true, false, unknown, or empty if an operand is &
    not defined &
    the same bounds, as IBEX\\
    C-XSC &
    \code{x} in the interior of \code{y} &
    \code{x} is a subset of \code{y}; \code{r <= x}: the number \code{r} is in
    \code{x} &
    the same bounds\\
    PROFIL/BIAS &
    \code{x} in the interior of \code{y} &
    \code{x} is a subset of \code{y}; \code{r <= x}: \code{r} is in \code{x} &
    the same bounds\\
    Boost.Interval &
    true if every point of \code{x} is below every point of \code{y}, false if
    none is, else it throws \code{comparison\_error} &
    true if every point of \code{x} is at most every point of \code{y}, false
    if none is, else it throws &
    true if \code{x} and \code{y} are the same point, false if they are
    disjoint, else it throws\\
    \hline
  \end{longtable}}

\noindent
So \code{0.0 <= f} is ``0 is in \code{f}'' in C-XSC and in PROFIL/BIAS, and in
GAOL v5 the certainly relation ``0 is at most every point of \code{f}'': a
program ported from them compiles and no longer tests what it did
(\code{f.set\_contains(0.0)} does). A Boost.Interval program that catches
\code{comparison\_error} to detect an uncertain comparison never sees it with
GAOL v5, whose \code{x < y} is false for overlapping intervals, and the empty
set, which Boost refuses, passes every certainly relation
(\pxref{sec:lexicographic-less}). The names are those of IBEX~2.9.1
(\file{ibex\_Interval.h}), Codac~2 (\file{codac2\_Interval.h}), C-XSC~2.5.4
(\file{interval.hpp} and \file{interval.inl}), PROFIL/BIAS~2.0.8
(\file{Interval.h}) and Boost~1.86 (\file{boost/numeric/interval/}), read in
their sources.

\section{Discussion}
%%------------------

This chapter has shown what a program using GAOL looks like, and what it owes
to the library: one header, the flags of interval arithmetic on every file
that includes it, which CMake, pkg-config and meson pass on, and a call to
\code{gaol::cleanup()} after the last use. The flags are the price of inline
operations: a compiler left free to round to nearest, to fold a constant or
to fuse a product and a sum would break the enclosure, and
\file{gaol/gaol\_config.h} refuses the options it can see, not the absence of
those it cannot. The two namespaces are the answer to two audiences, the
programs written for GAOL~4, IBEX and Codac, and those written to IEEE
1788-2015: the same type and the same functions, under two sets of names,
with the power as their one real difference.

\medskip
The common errors reviewed here all come from the same source: a double is
not an interval, and an interval is not a number. A decimal constant, a
function of a number, a midpoint and a radius, the sign of a zero bound, the
empty set that passes every test, the relations that order nothing, are
each a place where the habits of floating-point programming give a wrong
program that compiles. The tables of names of the other libraries serve the
same purpose for a program ported to GAOL.

\medskip
The next chapter looks at what this one took for granted: the rounding
direction that each operation sets, what it does to the rest of the program,
and the initialization and the cleanup of GAOL.

%% !TEX root = ../manual_gaolv5.tex
% Chapter 4.
%% Chapter 4: The rounding direction and the floating-point environment.
%% Sources: gaol.tex of f5bf180 lines 1109-1749 (the rounding direction, its
%% effects on the program, the NEON intervals, flush-to-zero, the
%% floating-point exceptions), 2301-2408 (initialization and cleanup),
%% 6124-6155 (the rounding functions of the FPU) and 475-598 (the AVX-512
%% path, in the options of CMake); doc/using.md lines 228-290 (initialization
%% and cleanup) and 1070-1972 (the rounding direction, the AVX-512 path, the
%% NEON and the FPU intervals, the control registers of 32-bit x86,
%% flush-to-zero, the floating-point exceptions).

\chapter{The rounding direction and the floating-point environment}
\label{chap:rounding}

An interval library computes its bounds with the rounding direction of the
processor set outward, and that direction is a state of the processor which
the whole program shares. How the library takes charge of this state, and
what it leaves to the program, decides both the correctness of the bounds and
the behaviour of every double the program computes beside its intervals. GAOL
4 left the direction to the user. GAOL v5 sets it itself, in each operation,
and the choice has consequences the user needs to know: the doubles of the
program are rounded upward while GAOL is initialized, the C library converts
its numbers in that direction, and the compiler restores nothing.

\medskip
The rounding direction is not the only state of the floating-point
environment that matters to the bounds. The modes that flush the subnormal
numbers to zero, which a program linked with \option{-Ofast} gets without
asking, make the bounds of some operations wrong, and the floating-point
exceptions, when a program enables them, stop the program in the middle of an
operation of GAOL. The chapter describes the three, and how GAOL defends its
bounds against them.

\medskip
In this chapter, we describe the rounding direction as GAOL v5 handles it,
and what it does to the rest of the program. The first section states how
each operation sets the direction, and which operations leave it alone. The
second section shows what the upward direction does to the doubles of the
program, and what a program does about it; the third describes the build that
preserves the direction of the program. The fourth, fifth and sixth sections
describe the AVX-512 path, which carries its rounding direction in its
instructions, and the NEON and the FPU intervals; the seventh describes the
control registers of the 32-bit x86 processors. The following two sections
treat the flush-to-zero modes and the
floating-point exceptions. The chapter ends with the reference of the
initialization and the cleanup of GAOL, and of its rounding functions. The
design of the kernels, which instructions compute the bounds and what the
check of the direction costs on each processor, is the subject of
Chapter~\ref{chap:computing} (\pxref{sec:kernels}); the options of the builds
are those of Chapter~\ref{chap:installation} (\pxref{sec:cmake-options}).

%%------------------------------------------------------------------------
\section{The rounding direction}\label{sec:rounding-direction}
%%------------------------------------------------------------------------
\cindex{rounding!direction}%
\cindex{rounding!downward}\cindex{rounding!upward}%
\cindex{F@\FSet}\cindex{R@\RSet}%
\cindex{containment property}%
\cindex{rounding!outward|see{outward rounding}}%
\cindex{downward rounding ($\roundDn{\:}$)}%
\cindex{upward rounding ($\roundUp{\:}$)}%
\cindex{outward rounding}%

Floating-point interval arithmetic requires \dfn{outward rounding} to fulfil
the containment property (\pxref{sec:containment}): the sum of \itv{a}{b} and
\itv{c}{d} is computed as \itv{\roundDn{a+c}}{\roundUp{b+d}}, where
\roundDn{r} is the greatest double of \FSet{} below the real $r$ and
\roundUp{r} the smallest one above it (\pxref{sec:rounding-definitions}). The
processor performs these two roundings when its rounding direction is set
towards $-\infty$ and towards $+\infty$ respectively.

\medskip
On most platforms, switching the rounding direction is costly. The number of
switches is cut down with the following property.

\begin{Prop}\label{prop:rounding-negation}
For every real $r$, $\roundDn{-r}=-\roundUp{r}$.
\end{Prop}

\noindent
Nearly every downward rounding is therefore replaced by an upward one, by
negating twice the operation performed. GAOL stores the lower bound of an
interval negated, and computes both bounds rounding upward: the rounding
direction is set towards $+\infty$ once, and no other direction is needed
(\pxref{sec:kernels}).

\medskip
GAOL~4 left the rounding direction to the user, who had to set it towards
$+\infty$ before any computation on intervals: this was its \dfn{trust
rounding mode}. GAOL v5 takes the direction in charge.\newinvfive{} Each
operation that computes a bound checks the direction, and sets it upward when
it is not, whatever the build; the operations of the AVX-512 path, which carry
their direction in the instruction, are the exception (\pxref{sec:avx512}).
The direction is then left upward, unless GAOL is built to preserve it
(\pxref{sec:preserve-rounding}). The bounds are therefore right whatever
rounding direction the calling code left.
\cindex{trust rounding mode}%

\medskip
The check is an addition, $1+2^{-60}$ or one of its variants below, which is above 1 only when rounded
upward, rather than a reading of the rounding direction, which costs far more
on some platforms: under Rosetta~2, reading it made an addition of intervals
take 800~ns rather than 2.8~ns, and it cost far more with 32-bit Visual C++
as well. The check takes three forms, by processor:
\begin{itemize}
\item on x86-64 processors, and on ARM processors with the FPU intervals and
  GCC or Clang, it is
  $1+(2^{-1060}+0)$, which shows the modes that flush the subnormal numbers to
  zero as well (\pxref{sec:ftz});
\item with the NEON intervals of 64-bit ARM processors (\pxref{sec:neon}), it
  is a read of the control register FPCR, which holds both the direction and
  these modes;
\item on 32-bit x86 processors, whose two units that round, the x87 unit and
  the SSE instructions, each have a direction, it is a read of the control
  word of the first and of MXCSR, with \code{fnstcw} and \code{stmxcsr}
  (\pxref{sec:x86-control-registers}).\newinvfive
\end{itemize}
\noindent
Each function of intervals checks the direction once, and the functions it
calls on doubles take it to be upward already. The cost of the check on each
processor is measured in Chapter~\ref{chap:computing}
(\pxref{sec:rounding-by-processor}).

\medskip
Unless GAOL is built to preserve it, the rounding direction stays upward from
the initialization of GAOL on, for the whole program, and \code{cleanup()}
sets back the one that the initialization found: every program calls
\code{gaol::cleanup()} right after its last use of GAOL
(\pxref{chap:init}).\newinvfive{} The direction set back is the one the first
\code{init()} found when GAOL initialized itself: to nearest, as a program
starts; on x86, for the x87 unit and for the SSE instructions each; in the
thread calling \code{cleanup()}. The rest of the floating-point environment is
left as it is, the exception flags raised in particular.
\cindex{rounding!preserving}%

\medskip
Some operations set no direction, and leave the one they found, with or
without \option{GAOL\_PRESERVE\_ROUNDING}: a program that wants a direction
after one of them sets it itself.\newinvfive
\cindex{rounding!left as found}%
\begin{itemize}
\item Those that round nothing and make no check: the negation and the unary
  plus, \code{abs()}, \code{min()}, \code{max()}, \code{sign()}, the roundings
  to an integer (\code{floor()}, \code{ceil()}, \code{trunc()},
  \code{integer()}, \code{round\_ties\_to\_even()}\ldots), \code{\&} and
  \code{|}, \code{hull()} and \code{intersect()}, \code{invabs\_rel()},
  \code{mig()} and \code{mag()}, \code{nth\_root(x, 0)} and
  \code{nth\_root(x, 1)}, the constructors from bounds, the access to the
  bounds, \code{is\_empty()}, the relations and the other predicates,
  \code{exact\_string()}.
\item Those that return before their check. Some return the empty set at once
  for an empty operand, and some do not: \code{y * empty} and
  \code{exp(empty)} leave the direction, where \code{empty + y} and
  \code{sqrt(empty)} set it upward, so that a program counts on neither. And
  \code{hausdorff(x, y)} returns $+\infty$ when a bound is infinite in one of
  the two intervals only, where it sets the direction upward otherwise.
\item The operations of the AVX-512 path, which set none
  (\pxref{sec:avx512}).
\end{itemize}

%%------------------------------------------------------------------------
\section{What the upward rounding does to the program}
\label{sec:upward-effects}
%%------------------------------------------------------------------------
\cindex{rounding!effects on the program}%

The rounding direction is the processor's, and the whole program shares it.
The results of GAOL do not depend on it, since each operation sets it upward
when it is not, in every thread; the doubles the program computes beside its
intervals do.\newinvfive{} Unless GAOL is built with
\option{GAOL\_PRESERVE\_ROUNDING}, they are rounded upward from the
initialization of GAOL, before \code{main()}, until \code{cleanup()}, and
again after each operation of GAOL that follows it and sets the direction. In
a build that takes the AVX-512 path, which without
\option{GAOL\_PRESERVE\_ROUNDING} only the developers ask for, the arithmetic
operations set none, and the direction after one of them is the program's.

\medskip
The compiler does not undo it. With the flags of \code{gaol::gaol},
\option{-frounding-math} in particular, an inexact operation on constants
written in a function, \code{1.0 / 3.0}, is not evaluated when compiling, to
nearest, but when the program runs, in the direction in effect. Only what the
compiler must evaluate when compiling, a \code{constexpr} or the constant
initializer of a static variable, is still rounded to nearest.

\medskip
Table~\ref{tab:upward-effects} gives what the program computes, with GCC~9.4
or Clang~18 and the GNU C library~2.31 on x86-64;
\file{examples/13\_rounding\_environment.cpp} of the sources computes each row
and checks it.

\begin{table}[htbp]\centering
\renewcommand{\arraystretch}{1.1}\footnotesize\setlength{\tabcolsep}{3pt}
  \caption[What the program computes while GAOL is initialized and after \code{cleanup()}]{What the program computes while GAOL is initialized and after \code{cleanup()}, with GCC~9.4 or Clang~18 and the GNU C library~2.31 on x86-64.} \label{tab:upward-effects}
  \begin{tabular}{L{.34\textwidth}L{.29\textwidth}L{.29\textwidth}}
    \hline
    The program computes & While GAOL is initialized & After \code{cleanup()}\\
    \hline\hline
    \code{1.0 / 3.0} & 0.33333333333333337 & 0.33333333333333331\\
    \code{std::stod("0.3") == 0.3} & false & true\\
    \code{std::lrint(2.3)} & 3 & 2\\
    \code{std::printf("\%.2f", 2.675)} & 2.68 & 2.67\\
    \code{std::cout <{}< std::setprecision(3) <{}< 3.14159265358979} & 3.15 & 3.14\\
    0.1 added $10^7$ times & 1000000.0005483569 & 999999.99983897537\\
    20\,000 doubles written with \code{\%.17g}, read back by \code{std::strtod} & 20\,000 come back different & none\\
    \code{TwoSum(1e20, 1e-20)}, the error of a sum & $s+e\neq a+b$ & $s+e=a+b$\\
    \code{TwoSum(a, b)}, 50\,000 pairs of doubles of exponents from $-60$ to 60 & $s+e\neq a+b$ for 10\,717 pairs & none\\
    \code{TwoProd(a, b)}, Dekker's splitting, 50\,000 pairs of doubles of $[1,2)$ & $p+e\neq ab$ for 1\,888 pairs & none\\
    \code{TwoProd(a, b)}, \code{e = std::fma(a, b, -a*b)}, the same pairs & none & none\\
    \hline
  \end{tabular}

\end{table}

The quotient is the double above $1/3$, where the program expects the
nearest. The C library rounds in the direction in effect too, as the C
standard recommends: with the GNU C library, \code{std::rint(-2.7)} is $-2$
where it is $-3$ to nearest, and \code{strtod()}, \code{stod()},
\code{printf()} and the streams round up the decimal digits they read and
write. A double written with 17 digits, which reads back to itself to nearest,
then comes back one double above, and one more at each round trip. Another C
library may convert to nearest whatever the direction: the example reports
what the one it runs with does. Sums are biased: each addition rounds up, so
0.1 added ten million times ends $5.5\cdot10^{-4}$ above the exact
$10^7\times 0.1$, where to nearest it ends $1.6\cdot10^{-4}$ below.

\medskip
The error-free transformations break. TwoSum, the algorithm that recovers the
error of a sum,

\begin{example}
s = a + b;  bb = s - a;  e = (a - (s - bb)) + (b - bb);
\end{example}

\noindent
gives $s+e=a+b$ only to nearest. Rounded upward, the exact error of
$10^{20}+10^{-20}$, which is $10^{-20}-16384$, is no double: \code{s} is
$10^{20}+16384$ and \code{e} is $-16383.999999999998$, whose sum is not
$10^{20}+10^{-20}$. TwoProd, which recovers the error of a product from
Dekker's splitting of the operands, breaks the same way: 1\,888 of the
50\,000 pairs of the table, drawn in $[1,2)$, lose their exactness upward,
where pairs of widely different magnitudes keep it, the error of the product of
the split parts they drop staying below the error of the rounded product. It
is exact in every direction when it is computed with \code{std::fma} instead:
the error of a rounded product is a double, which \code{std::fma} computes
however it rounds, provided it is correctly rounded, by the instruction or by
a call of a C library whose \code{fma()} is (with GCC, \code{std::fma} without
optimization is a call of the C library, whose \code{fma()} is wrong on
Cygwin, where \code{\_\_builtin\_fma} is the instruction with
\option{-mfma}, which GAOL calls itself). Compensated sums, and the
double-double arithmetic built on such algorithms, lose their accuracy the
same way. Built with \option{GAOL\_PRESERVE\_ROUNDING}, the program computes
the right-hand column of the table all the time.

\medskip
\noindent
The upward direction has other consequences:
\begin{itemize}
\item a thread created after the initialization of GAOL starts upward too
  (\code{std::thread} gives a new thread the floating-point environment of the
  thread constructing it), and \code{cleanup()} sets back the direction of the
  calling thread only;
\item the libraries the program calls run upward as well: the \code{exp},
  \code{log} and \code{pow} of the GNU C library~2.31 go from 0.51 ulp to 1.01
  ulp from the exact value, its \code{cbrt} from 2.8 to 3.8 ulp (20\,000
  random arguments each, against mpmath), and its \code{sin}, \code{cos},
  \code{tan} and \code{atan} give the same results;
\item GCC does not model the rounding direction, even with
  \option{-frounding-math} (GCC bug~34678,
  \url{https://gcc.gnu.org/bugzilla/show_bug.cgi?id=34678}), and reuses after
  \code{cleanup()} a double computed before it, still rounded upward, where
  the program computes it again: from \option{-O1}, GCC~9.4 does it for
  \code{a / b} written on both sides of the call, where Clang~18 computes the
  quotient again. A program cannot rely on either.
\end{itemize}

\medskip
\noindent
What a program does about it:
\begin{itemize}
\item[$\bullet$] \textbf{It calls \code{cleanup()} right after the last use of
  GAOL}, before it prints, converts or sums what it computed: its doubles are
  then computed to nearest. A program that goes back to GAOL afterwards calls
  \code{restore\_rounding()} (or \code{std::fesetround(FE\_TONEAREST)}) after
  each use, since an operation of GAOL sets the direction upward again and
  \code{cleanup()} sets it back once only (\pxref{chap:init}).
\item[$\bullet$] \textbf{It reads and writes numbers through GAOL}, not
  through the C library: \code{gaol::textToInterval("0.3")} encloses $3/10$ in
  every direction, where \code{std::stod("0.3")} does not, and
  \code{exact\_string()} and \code{\%a} write a number that reads back to the
  same bits in every direction. GAOL writes its intervals rounded outward
  whatever the direction.
\item[$\bullet$] \textbf{It computes in a block the doubles that need rounding
  to nearest} (an error-free transformation, \code{std::lrint()}, a decimal
  conversion, a call of a library written for that): it saves the direction
  with \code{std::fegetround()}, sets \code{FE\_TONEAREST}, and sets the saved
  one back when the block ends, with no operation of GAOL inside, which would
  set the direction upward again. The example of the sources writes such a
  guard, \code{nearest\_scope}.
\item[$\bullet$] \textbf{It does not count on the compiler} to compute a
  double where the source code does. One that has to be computed on this side
  of a change of direction is computed from operands read again through
  \code{gaol\_core::rnd\_keep()}, which writes its argument to volatile memory
  and returns it: \code{gaol\_core::rnd\_keep(a) / gaol\_core::rnd\_keep(b)}
  after \code{cleanup()} is the quotient to nearest with GCC too, where
  \code{gaol\_core::rnd\_keep(a / b)} may be the one computed before.
\item[$\bullet$] \textbf{Or GAOL is built with
  \option{GAOL\_PRESERVE\_ROUNDING}}, and the program keeps its own direction
  (\pxref{sec:preserve-rounding}).
\end{itemize}

%%------------------------------------------------------------------------
\section{Preserving the rounding direction}\label{sec:preserve-rounding}
%%------------------------------------------------------------------------
\cindex{rounding!preserving}%
\cindex{GAOL PRESERVE ROUNDING@\option{GAOL\_PRESERVE\_ROUNDING}}%

Code that needs its own rounding direction after the operations of GAOL
builds GAOL with the option \option{GAOL\_PRESERVE\_ROUNDING} of CMake, which
is \loption{enable-preserve-\allowbreak rounding} with autotools and
\option{-Denable-preserve-\allowbreak rounding=true} with meson
(\pxref{sec:cmake-options}). It is off by default in the three builds. With
it, each operation reads the rounding direction it finds, sets it upward, and
sets back the one it found when it is done; it also sets back the modes
flushing the subnormal numbers to zero that it cleared (\pxref{sec:ftz}). The
operations of the AVX-512 path set none (\pxref{sec:avx512}). Nothing else
changes: the bounds are the same, and the program still calls
\code{cleanup()}.\newinvfive

\medskip
\noindent
The option changes three other things.
\begin{itemize}
\item[$\bullet$] \textbf{The AVX-512 path.} With it, the build takes the
  AVX-512 path (it was the default of an option,
  \option{GAOL\_PREFER\_AVX512}): where the machine building has the AVX-512
  instructions, the four arithmetic operations, the square, the square root
  and, with FMA, the integer powers beyond 1 and $-1$ take them, which carry
  their rounding direction, set none and store none. The library then needs a
  processor with them.\newinvfive
\item[$\bullet$] \textbf{The initialization.} It leaves the whole
  floating-point environment as it found it, where the default build sets the
  default environment: an exception a static object enabled before it stays
  enabled (\pxref{sec:fp-exceptions}).
\item[$\bullet$] \textbf{\code{restore\_rounding()}.} It does nothing, GAOL
  leaving the direction of the program alone (\pxref{chap:init}).
\end{itemize}

\medskip
\noindent
The direction is written only where it has to be. On the processors whose
control registers GAOL writes itself, an operation reads the direction and the
flush modes once at its start, writes the register only if the direction is
not upward or a flush mode is set, and writes back what it found at its end
only if it differs from what the operation left: a program that computes
upward already, with the flush modes clear, pays no write. These registers
are FPCR on 64-bit ARM and FPSCR on 32-bit ARM with GCC (\pxref{sec:neon}),
and the control word of the x87 unit and MXCSR on x86
(\pxref{sec:x86-control-registers}).

\medskip
The option costs time, less with the AVX-512 path, which it turns on where
the machine has the instructions. Table~\ref{tab:preserve-cost} gives the
factor by which each operation gets longer on an Intel i7-1185G7 with
Clang~21.1.8 (\file{tests/performance.cpp} and the long lines of
\file{doc/compare/code/run\_lines.sh}, the median of 9 interleaved runs of
6 November 2026), with the path, and without it (a build for a machine without
the instructions, or one that refuses the path with
\option{-DGAOL\_PREFER\_AVX512=OFF}, a setting of the developers). The
arithmetic operations take 1.2 to 1.6 times longer with the path, and several
times longer without it; a line of 15 additions of intervals takes 2.1 times
longer with the path and 8.5 times without.

\begin{table}[htbp]\centering\footnotesize
  \caption[The cost of \option{GAOL\_PRESERVE\_ROUNDING}, with and without the AVX-512 path]{How many times longer each operation takes with \option{GAOL\_PRESERVE\_ROUNDING} than without it, with and without the AVX-512 path (Intel i7-1185G7, Clang~21.1.8, \file{tests/performance.cpp}, the median of 9 interleaved runs, 6 November 2026).} \label{tab:preserve-cost}
  \begin{tabular}{lcc}
    \hline
    Operation & With the AVX-512 path & Without it\\
    \hline\hline
    \code{x + y}, \code{x - y} & 1.2 & 7.2\\
    \code{x * y} & 1.6 & 4.3\\
    \code{x / y} & 1.6 & 3.4\\
    \code{sqr(x)} & 1.0 & 4.4\\
    \code{sqrt(x)} & 0.6 (it is faster) & 1.5\\
    \code{pow(x, 3)} & 1.2 & 1.6\\
    \code{exp(x)} & 1.2 & 1.2\\
    \code{log(x)} & 1.1 & 1.0\\
    \code{sin(x)}, \code{cos(x)} & 1.2 & 1.1 to 1.2\\
    a line of 15 additions & 2.1 & 8.5\\
    \hline
  \end{tabular}

\end{table}

%%------------------------------------------------------------------------
\section{The AVX-512 path}\label{sec:avx512}
%%------------------------------------------------------------------------
\cindex{AVX-512}%

The AVX-512 instructions of the x86-64 processors carry a rounding direction
in the instruction itself, the \dfn{embedded rounding}: an operation that
takes them reads no direction, sets none and stores none. GAOL takes them in
the build that preserves the rounding direction, in which every operation of
the other path saves and restores the direction
(\pxref{sec:preserve-rounding}): the path relieves \code{+}, \code{-},
\code{*}, \code{/}, \code{sqrt}, \code{sqr} and the integer powers of that
store.\newinvfive{} The instructions of the kernel are described in
Chapter~\ref{chap:computing} (\pxref{sec:kernels}).

\medskip
\noindent
\textbf{When the build takes it.} The path is decided when GAOL is configured,
as the fused multiply-add instructions are (\pxref{sec:cmake-options}): where
\option{GAOL\_PRESERVE\_ROUNDING} is on, the compiler takes the instructions
and the machine building runs them, the macro \code{GAOL\_AVX512\_ARITHMETIC}
is defined in \file{gaol/gaol\_configuration.h}; a cross-compilation leaves
the path out, and the operations test nothing as they run. With that option
every operation gains, 2.8 times at the median of seven operations, from 1.3
to 4.4 times, on an Intel i7-1185G7, whose AVX-512 units are halved; without
it they cost, 17\,\% at the median, the direction being written only when it
is not upward already, so the path is left out. \option{GAOL\_PREFER\_AVX512},
\loption{enable-prefer-avx512} and \option{-Denable-prefer-avx512} were
options of the user, and are now settings of the developers:\newinvfive{}
\code{=ON} asks for the path without the preserved direction, and
\option{-DGAOL\_PREFER\_AVX512=OFF}, \loption{disable-prefer-avx512} and
\option{-Denable-prefer-avx512=false} refuse it, for a package for any
processor of the architecture. A library built with the path needs a processor
with the instructions, as with the fused multiply-add instructions, and
\code{gaol::init()} stops a program run on another one, with a message on the
standard error, rather than letting its first operation stop on an illegal
instruction (an operation made before GAOL initializes itself, as described
in \exref{chap:init}, still does). GAOL no longer chooses the path as the
program runs, in \code{init()}, with every operation testing the flag it set,
in the default build too.\newinvfive

\medskip
\noindent
\textbf{The operations.} On such a build, \code{x + y}, \code{x - y},
\code{x * y}, \code{x / y}, \code{sqrt(x)} and \code{sqr(x)}, and the
operations \code{+}, \code{-}, \code{*} and \code{/} with a double, are
computed by the rounding direction the instructions carry. So are
\code{pown(x, n)} and \code{pow(x, n)} for an integer $n$ of 2 or more
($n=2$ is \code{sqr(x)}), or of $-2$ or less where \code{x} has no zero,
whose two corners are within the range of the exact products of their kernel;
the exponents 0, 1 and $-1$, and the powers beyond the range, which a corner
near 0 or a large exponent makes, take the path of the direction. These
operations set no rounding direction: the program keeps its own, in the
default build as with \option{GAOL\_PRESERVE\_ROUNDING}, and the other
operations (\code{\%}, the elementary functions, the reading of a number, the
integer powers beyond the range) leave it upward, as they always do.

\medskip
\noindent
\textbf{The flags and the flush modes.} The instructions of the path raise no
floating-point exception flag, the architecture requiring the suppression of
the exceptions with the embedded rounding; the check of the other operations
raises the flags described in \exref{sec:fp-exceptions}. The modes that flush
the subnormal numbers to zero, which the embedded rounding honours as the
other instructions do (measured on an Intel i7-1185G7: with flush-to-zero
set, the difference of \code{3e-308} and \code{2.9e-308} is 0), are cleared
before the operation reads a bound when a program or a plug-in has set them,
as the other operations clear them (\pxref{sec:ftz}), and set back after its
result with \option{GAOL\_PRESERVE\_ROUNDING}: with Clang~18, the difference
of \code{x - y} was once computed after flush-to-zero had been set back, and
\code{[3e-308] - [2.9e-308]} was \itv{0}{0}.\newinvfive

\medskip
\noindent
\textbf{The bounds.} They are those of the other path, bit for bit, but for
the sign of a zero bound in one case.

\begin{Rem}\label{rem:avx512-zero}
\code{x * d}, \code{x *= d} and \code{d * x}, which is computed as
\code{x * d},\newinvfive{} of an interval whose zero bounds have different
signs, \itv{+0}{-0}, are \itv{-0}{+0} with the AVX-512 path and \itv{+0}{-0}
with the SSE2 one, $+0$ and $-0$ being the same number
(\pxref{sec:zero-sign}). The rounding rule is the same, and the tests of
\file{tests/}, which a job of the continuous integration runs with the option
on, check it with \code{==}, which does not see the sign of a zero;
\code{rounding\_direction} checks the bits of that one case in each path, and
every result of the operations of intervals and with a double on the
intervals of two bounds among 0 of either sign, $\pm1$, $\pm2.5$, $\pm3$,
$\pm10^{-310}$, $\pm10^{300}$ and $\pm\infty$ was compared bit for bit.
\end{Rem}

\medskip
\noindent
\textbf{The cost.} The operations of the path cost more time on a processor
that halves its AVX-512 instructions, where they cost less on one that does
not. On an Intel i7-1185G7 (halved), with Clang~21.1.8
(\cmd{gaol\_performance}, the median of 9 runs of 26 October 2026, from
Table~\ref{tab:computing-avx512}), \code{x * y} takes
3.7~ns rather than 2.3 in the default build, \code{x / y} 4.1~ns rather than
2.6 and \code{x + y} 1.5 rather than 1.1, but \code{sqrt(x)} 4.4 rather than
8.8, the instruction rounding each root where the SSE2 path checks the root
of the C library with a product. With \option{GAOL\_PRESERVE\_ROUNDING}, where
the SSE2 path saves and restores the direction, \code{x + y} takes 1.9~ns
rather than 8.4, \code{sqr(x)} 2.0 rather than 8.3 and \code{pown(x, 3)} 7.5
rather than 9.9; the comparison of the libraries (\pxref{sec:comparison}) has
the preserved column of the other operations. The processors that do not
halve them, as the Xeon of the continuous integration, are expected to take
less. Which of the paths a GAOL took,
\file{examples/13\_rounding\_environment.cpp} tells.

%%------------------------------------------------------------------------
\section{The NEON intervals}\label{sec:neon}
%%------------------------------------------------------------------------
\cindex{NEON}\cindex{FPCR}\cindex{FPSCR}%

On 64-bit ARM processors (AArch64), GAOL computes its intervals with NEON
instructions, as it does with SSE2 instructions on x86.\newinvfive{} The two
bounds of an interval are the two lanes of one \code{float64x2\_t}, the lower
bound negated as in the \code{\_\_m128d} of the SSE2 intervals, and \code{+},
\code{-}, \code{*}, \code{/}, \code{sqr()} and \code{sqrt()} are kernels that
take and give their bounds in registers. The sum and the difference are one
vector addition rounded upward; the product is the maximum of the four
products of the bounds, with no branch on their signs; the quotient is one
vector division when the divisor has no zero, and single divisions otherwise;
the square roots of both bounds are one vector square root rounded upward,
the lower bound being that root where it is exact, and the double below it
otherwise (\pxref{sec:kernels}).

\medskip
They are the intervals of GAOL on 64-bit ARM with every compiler (GCC, Clang
and AppleClang under Linux and macOS, clang-cl and the Clang of MinGW-w64 and
MSYS2, and Visual C++), by default: the option is \option{GAOL\_SIMD} of
CMake, \loption{enable-simd} of \cmd{configure} and \code{enable-simd} of
meson, as for the SSE2 intervals (\pxref{sec:cmake-options}), and the macro
\code{GAOL\_USING\_NEON\_INSTRUCTIONS} is then 1 in
\file{gaol/gaol\_configuration.h}. They need no flag beyond those of interval
arithmetic (\pxref{sec:compiling}). \option{GAOL\_ARM64},
\loption{enable-arm64} and \option{-Denable-arm64} were an option of their
own, independent of \option{GAOL\_SIMD}, and are gone.\newinvfive{} 32-bit
ARM, whose NEON has no vector of doubles, does not take them: the option is
ignored there, and the FPU computes the intervals, as before
(\pxref{sec:fpu-intervals}). A universal build of macOS made by CMake takes
them for its arm64 part, and the SSE2 intervals for its x86\_64 part.

\medskip
Visual C++ writes the control register FPCR through the intrinsics
\code{\_ReadStatusReg()} and \code{\_WriteStatusReg()}, where GCC and Clang
have \code{mrs} and \code{msr} in an asm statement; it has no asm statement
to order its operations with respect to those writes, and relies on
\option{/fp:strict} to keep them where the source puts them, which
\file{tests/rounding\_direction.cpp} checks. With \option{/W4}, Visual C++
for arm64 warns (C4324, ``structure was padded due to alignment specifier'')
of each structure that holds an interval, its \code{float64x2\_t} being
aligned on 16 bytes there, and of \code{expr\_eval} of the installed
\file{gaol/parser/gaol\_expr\_eval.h}: the warning is harmless, and no job of
the continuous integration compiles with \option{/W4 /WX} for
arm64.\newinvfive

\medskip
\noindent
\textbf{The control register FPCR.} 64-bit ARM has no instruction that
carries its own rounding direction, as the AVX-512 path has: the direction is
the field RMode of the control register FPCR (bits 23 and 22, upward being
\code{0x400000}), which the scalar and the vector instructions share, and
which holds the modes that flush the subnormal numbers to zero as well (FZ,
bit 24, and from Armv8.7 FIZ, bit 0). The NEON intervals therefore check the
direction by reading FPCR with \code{mrs}, which raises no exception flag, and
write it with \code{msr} only if it is not upward with the flush modes clear,
where the FPU intervals add $1+(2^{-1060}+0)$, which raises the inexact flag
(\pxref{sec:fp-exceptions}). \code{msr} on FPCR is a serializing instruction,
slower than \code{mrs}, which is why it is written when needed only. With
\option{GAOL\_PRESERVE\_ROUNDING}, each operation reads FPCR once, and writes
back the direction and the modes it found at its end only if they differ from
what the operation left: a program that computes upward already, with the
flush modes clear, pays no write. \code{round\_upward()},
\code{round\_nearest()} and \code{round\_downward()} write FPCR as well, and
not through \code{fesetround()}, and only if the direction changes
(\pxref{chap:fpu}). \code{fegetround()} and \code{fesetround()} of the C
library read and write the same bits (the \code{FE\_*} values of the C
libraries of 64-bit ARM are these bits, which a static assertion checks), so
that a program may use either.

\medskip
\noindent
\textbf{The register FPSCR of 32-bit ARM.} The FPU intervals of 32-bit ARM
processors with GCC and a floating-point unit of doubles do the same for
\code{round\_upward()}, \code{round\_nearest()} and \code{round\_downward()},
and for the direction that \option{GAOL\_PRESERVE\_ROUNDING} restores
(\code{GAOL\_RND\_FPSCR\_DIRECT} of \file{gaol/gaol\_config.h}, 1 there;
\option{-DGAOL\_RND\_FPSCR\_DIRECT=0} gives the C library back). The register
is FPSCR, read with \code{vmrs} and written with \code{vmsr}, with the same
field RMode (the \code{FE\_*} values of the C libraries of 32-bit ARM are its
bits, which the same assertion checks). The flags of the exceptions and of
the comparisons are in it too, so that the functions that write it from the
value they read at once (\code{round\_nearest()}, \code{round\_upward()}, the
end of an operation) read it again just before, and no flag raised since is
lost. With \option{GAOL\_PRESERVE\_ROUNDING}, each operation reads FPSCR once
at its start, writes it only if the direction is not upward or a flush mode
is set, and writes back what it found at its end only if it differs. On the
Debian armhf jobs of the continuous integration, which run on the arm64
runners of GitHub, \code{x + y} took 22.5~ns rather than 49 with
\option{GAOL\_PRESERVE\_ROUNDING}, a line of 15 additions 318~ns rather than
721, and \code{midpoint()} 38~ns rather than 63.\newinvfive{} The check of
every operation without the option stays the addition: \code{vmrs} took
longer than the two additions there, \code{x + y} 6.1~ns rather than 4.1 and
a line of 15 additions 87~ns rather than 57, and the 32-bit ARM intervals
keep their time.

\medskip
\noindent
\textbf{The bounds.} They are those of the FPU intervals but for the sign of
a zero bound, which the maximum of 64-bit ARM, ranking $+0$ above $-0$, can
set otherwise than the SSE2 and the FPU intervals do, where the \code{maxpd}
of x86 returns its second operand (\pxref{sec:zero-sign}).

\begin{Rem}\label{rem:neon-zero}
The lower bound of \code{[-1, 0] * [-1, 0]} is $-0$ with the NEON intervals
and $+0$ with the SSE2 ones (the FPU intervals leave it to the compiler,
\pxref{sec:fpu-intervals}), and \code{interval::zero()} is stored
\code{<+0, +0>}, as with SSE2, so that its lower bound is $-0$; $+0$ and
$-0$ being the same number, \code{==} does not see the difference. The tool
\file{tests/tools/arith} compares the NEON and the FPU intervals on a grid of
operands, with the calling code in each of the four rounding directions: it
found no other difference. The same kind of caveat exists for the AVX-512
path (Remark~\ref{rem:avx512-zero}).
\end{Rem}

%%------------------------------------------------------------------------
\section{The FPU intervals}\label{sec:fpu-intervals}
%%------------------------------------------------------------------------
\cindex{FPU intervals}%

The intervals of the processors that have neither SSE2 nor NEON keep their
two bounds as two doubles, the lower one negated as in the SSE2 and the NEON
intervals, and compute with the scalar floating-point unit. They are the
intervals of 32-bit ARM, POWER, s390x, RISC-V and the other processors; of
Windows on 32-bit x86 with a compiler that takes no flag asking for the
aligned memory the SSE2 intervals need (MinGW-w64 with GCC and Clang takes
the SSE2 intervals there, with \option{-faligned-new=8}, and so do Visual C++
and clang-cl, whose intervals need no aligned memory, and 64-bit Windows);
and of the builds that ask for them (\option{-DGAOL\_SIMD=OFF},
\loption{disable-simd} or \option{-Denable-simd=false}), which leave
\code{GAOL\_USING\_SSE2\_INSTRUCTIONS} and
\code{GAOL\_USING\_NEON\_INSTRUCTIONS} at 0 in
\file{gaol/gaol\_configuration.h}.

\medskip
\code{x * y}, \code{x / y} and \code{x \% y} of a divisor without zero, and
\code{sqr(x)} take no branch on the signs of the bounds in the
source:\newinvfive{} GAOL~4 took one of nine cases of the signs for a
product, and went down a chain of up to eight tests for a quotient, which a
processor mispredicts when the signs come at random. The product is the
largest of four products for the upper bound and of four others for the lower
one, each with a factor negated, a product and its opposite not being rounded
the same way where the only direction is upward; the quotient divides each
bound of the dividend by the bound of the divisor that its sign chooses, two
divisions whatever the signs; the square is the largest of two squares for
each bound, the lower one being 0 for an interval that straddles 0. They are
the kernels of the SSE2 and the NEON intervals on two doubles each
(\pxref{sec:kernels}), and give the bounds of the SSE2 ones bit for bit for
the product of finite bounds and the quotient by a divisor without zero
(checked against a model of their kernels, not by running them), but for the
sign of a zero bound of a product (Remark~\ref{rem:fpu-zero}). An interval
with an infinite bound, whose 0 times an infinity is no product to form,
takes a function that sets both factors of such a product to 0 (the nine
cases of GAOL~4 took 12.4~ns for operands with an infinite bound drawn at
random, and it takes 9.2), and a divisor with a zero takes the table of the
NEON intervals, one division at most. \code{x - y} reads the bounds of
\code{y} after the check of the rounding direction, as \code{x + y} does: it
took three times as long as the sum, the bounds being kept across the calls
of the rare path of the check.

\medskip
\noindent
\textbf{The time of the kernels.} Table~\ref{tab:fpu-m2} gives the times
measured with \cmd{gaol\_performance} on an Apple M2 with AppleClang~21
(\option{-DGAOL\_SIMD=OFF}), the median of nine runs of each program, in turn
with the others, on 6 November 2026, in nanoseconds, with the code of GAOL~4
(the tree of 22 October 2026, before the kernels) and with the kernels of GAOL
v5.

\begin{table}[htbp]\centering\footnotesize
  \caption[The FPU intervals on an Apple M2 (\option{GAOL\_SIMD=OFF})]{The FPU intervals on an Apple M2 with AppleClang~21 (\option{-DGAOL\_SIMD=OFF}): the code of GAOL~4 and the kernels of GAOL v5, in nanoseconds (\cmd{gaol\_performance}, the median of nine runs).} \label{tab:fpu-m2}
  \begin{tabular}{lrr}
    \hline
    & GAOL~4 code & GAOL v5 kernels\\
    \hline\hline
    \code{x - y} & 3.54 & 1.15\\
    \code{x * y} & 3.25 & 3.25\\
    \code{x * y}, random signs & 6.39 & 3.31\\
    \code{x / y} & 3.98 & 3.23\\
    \code{x / y}, random signs & 3.17 & 3.21\\
    \code{sqr(x)}, random signs & 2.64 & 1.78\\
    a line of 7 \code{*} & 23.3 & 35.0\\
    a line of 7 \code{/} & 35.1 & 32.9\\
    a line of 13 \code{+ - * /} & 49.6 & 50.4\\
    \hline
  \end{tabular}

\end{table}

\noindent
The operands of \cmd{gaol\_performance} are 1\,024 intervals that the
processor learns by heart, so that a branch on their signs costs less than it
does on other operands. On 100\,000 intervals centered at random (the
benchmark of \file{doc/compare}, \pxref{sec:comparison}), \code{x * y} took
13.9~ns and takes 3.5, \code{x / y} 8.5 and 3.2, \code{sqr(x)} 6.9 and 1.8,
\code{x - y} 3.7 and 2.0, and Shekel 5\,298~ns and 138.

\medskip
A line of products of operands of one sign, which a processor predicts, takes
more time: 7 products in a line took 23.3~ns and take 35.0 (23.7 and 35.2 in
the benchmark of \file{doc/compare}, \code{pos\_mul8}; Apple M2). Each
product waits for the one before it, and the largest of four products waits
for its products and for two comparisons, where one product on a predicted
branch waits for one. The maximum of 64-bit ARM (\code{fmax}, one
instruction, which ranks $+0$ above $-0$) would shorten the wait: it is not
used, being a call with GCC for 32-bit ARM and x86. With
\option{GAOL\_PRESERVE\_ROUNDING} the write of the control register waits for
the result, so that the longer wait of the product shows there as well:
\code{x * y} takes 42.4~ns rather than 39.7 and the line of 7 \code{*} 297
rather than 279, that is 6 to 7\,\% more (\code{sqr(x)} 42.4 rather than
39.5; \code{x / y} 43.6 rather than 43.1; \code{x * y} of random signs 42.4
rather than 47.5, where the branches of the former code were mispredicted;
the median of nine runs of each code, in turn, on 6 November 2026). The operators write their result into
the bounds of \code{*this}, memory of the caller that is also given by address
to the functions of the rare cases, which a call or an asm statement with a
``memory'' clobber that restores the direction finds written: the products
they hold are computed before it, with the direction upward. That is how the
compilers treat it (GCC~15 for 32-bit ARM: the stores precede
\code{bl fesetround}), not a rule of the language, and no test holds it:
\code{GAOL\_RND\_KEEP()} would make it one, for 2.6~ns more with
\option{GAOL\_PRESERVE\_ROUNDING}, and the former product and \code{x + y}
did without.

\medskip
\noindent
\textbf{The processors of the continuous integration.} The processors and
compilers of the continuous integration that took the FPU intervals were
measured in turn (Visual C++ and clang-cl for arm64 take the NEON intervals
since),\newinvfive{} in one run of GitHub Actions, on an experiment branch:
each job builds the code before the kernels and after them, and runs
\cmd{gaol\_performance} five times with each, alternately, with the tools of
\file{.github/scripts/performance-ab.sh} and \file{compare-performance.py}
(\pxref{sec:scripts}). The compilers are GCC~13.3 and Clang~18.1 on
Ubuntu~24.04 for x86-64 (\option{-DGAOL\_SIMD=OFF}), GCC~14.2 on Debian~13
for 32-bit ARM, run natively on an arm64 runner, Visual C++~19.51 for x86 and
arm64, clang-cl~22.1 for x86 and arm64, and AppleClang~21 on macOS~26
(\option{-DGAOL\_SIMD=OFF}). Table~\ref{tab:fpu-ci} gives the ratio of the
median before the kernels to the median after them, so that a ratio above 1
means that the kernels are faster. These are figures of shared runners, of
one run, of the FPU intervals as they were: Visual C++ and clang-cl for x86
have the SSE2 intervals since,\newinvfive{} with their two bounds as two
doubles, and their rows are no longer those of the default build.

\begin{table}[htbp]\centering\scriptsize\setlength{\tabcolsep}{2.5pt}
  \caption[The FPU intervals on the processors of the continuous integration]{The FPU intervals on the processors of the continuous integration: the ratio of the median time of the code of GAOL~4 to the one of the kernels of GAOL v5 (five runs of \cmd{gaol\_performance} each, alternately; a ratio above 1 means that the kernels are faster). The columns are \code{x * y}, \code{x * y} of random signs, \code{x / y}, \code{sqr(x)}, \code{x - y}, a line of 7 \code{*}, a line of 7 \code{/} and a line of 13 \code{+ - * /}; ``preserved'' is the build with \option{GAOL\_PRESERVE\_ROUNDING}. $\dagger$: the spread of the five runs exceeds 1.15, and the ratio is noise; in bold, the ratios below 0.9 of a stable row.} \label{tab:fpu-ci}
  \begin{tabular}{l*{8}{r}}
    \hline
    & \code{x*y} & \code{x*y} rnd & \code{x/y} & \code{sqr} & \code{x-y}
      & 7 \code{*} & 7 \code{/} & 13 ops\\
    \hline\hline
    Ubuntu, GCC 13 & 1.29 & 1.64 & 1.25 & 1.09 & 1.02 & \textbf{0.58} & 1.13 & \textbf{0.76}\\
    Ubuntu, Clang 18 & 1.18 & 1.67 & 1.47 & \textbf{0.85} & 0.99 & 0.98 & 1.05 & 1.01\\
    armhf, GCC 14 & 1.41 & 2.28 & 1.63 & 1.17 & \textbf{0.86} & 1.11 & 1.25 & 1.04\\
    Visual C++, x86 & 1.68$^\dagger$ & 3.47 & 1.11 & 1.39 & 1.06 & 1.43 & 1.16 & 1.22$^\dagger$\\
    Visual C++, x86, preserved & 1.18 & 1.64$^\dagger$ & 1.02$^\dagger$ & 1.09$^\dagger$ & 1.02$^\dagger$ & 1.18 & 1.04 & \textbf{0.53}\\
    Visual C++, arm64 & 1.53 & 2.41 & \textbf{0.82} & 2.36 & 1.08$^\dagger$ & 1.71 & 1.05 & 1.41\\
    clang-cl, x86 & 1.06 & 1.23 & 1.24 & 1.06 & 1.04 & 1.04 & 1.16 & 1.05\\
    clang-cl, arm64 & \textbf{0.77} & 1.75 & 1.29 & \textbf{0.84} & 1.01 & \textbf{0.57} & 1.24 & \textbf{0.86}\\
    macOS, FPU & 0.92$^\dagger$ & 1.98$^\dagger$ & 1.13$^\dagger$ & 1.03$^\dagger$ & 2.49$^\dagger$ & 0.60$^\dagger$ & 1.04$^\dagger$ & 0.83$^\dagger$\\
    \hline
  \end{tabular}

\end{table}

\noindent
The product wins where the signs are random (1.23 to 3.47 times) and, alone,
of positive operands on 32-bit ARM, Visual C++ and Ubuntu. It loses on the
lines of products of positive operands with GCC on Linux (0.58 for the line
of 7 \code{*}, that is 1.7 times slower, where the Apple M2 above gives 23.3
and 35.0~ns, 1.5 times) and with clang-cl for arm64 (0.57), and alone with
clang-cl for arm64 (0.77): there, the branch of the former code is well
predicted and the kernel of eight products adds latency. The macOS runner is
too noisy to say (spreads of 1.30 to 2.71). The quotient wins from 1.11 to
1.63 times on six configurations, is within the noise on two (the direction
preserved, 1.02, and macOS, 1.13) and loses with Visual C++ for arm64 (0.82,
spread 1.04). \code{sqr(x)} is mixed: 0.84 with clang-cl for arm64 and 0.85
with Clang on Ubuntu, 2.36 with Visual C++ for arm64. \code{x - y} does not
change where the rows are stable (0.86 to 1.06), 0.86 on 32-bit ARM where
\code{x + y}, whose code did not change, gives 0.92 in the same job; the 3.54
to 1.15~ns of the M2 does not show. With the direction preserved on Visual
C++ for x86, the line of 13 operations is 1.9 times slower (0.53, spread
1.08), where each operation wins or stays: the cause was not searched. POWER,
s390x, RISC-V and MinGW-w64 for arm64 are not measured.

\medskip
What a compiler makes of the choices of the kernels is its own, and was read
in the assembly on Compiler Explorer, not run. The maxima of the product are
conditional moves with GCC (13 to 16) for 32-bit ARM, arm64 and x86, with
Clang, and with Visual C++ for x86 and arm64, \code{xsmaxcdp} instructions
with GCC~13 for POWER9, and branches on the products with GCC for POWER8,
s390x and RISC-V, which have no conditional move of doubles. The choices of
the quotient and of the square are conditional moves with GCC for POWER and
with Clang and Visual C++ for arm64, and branches on a sign with GCC for
32-bit ARM at \option{-O2} and \option{-O3} (two in the quotient with GCC 13
to 15, four with GCC~16, none in the square; at \option{-Os}, two in the
square with GCC~14 and~15), arm64 (one in the quotient), s390x and RISC-V,
and with Visual C++ for x86.

\begin{Rem}\label{rem:fpu-zero}
The sign of a zero bound of a product is not specified for the FPU
intervals: where two products are zeros of opposite signs, $+0$ and $-0$, the
maximum of the two is a comparison and a selection, and the compiler decides
which of the two it returns. The \code{maxsd} of x86 returns its second
operand, and the compilers that keep the order written do the same (the lower
bound of \code{[-1, 0] * [-1, 0]} was $+0$ on the jobs of the continuous
integration with GCC for 32-bit ARM and for arm64, with Clang for arm64 and
with Visual C++ for x86), but GCC~13 for POWER9 compiles
\code{(a > b) ? a : b} into \code{xsmaxcdp}, whose result for a $+0$ and a
$-0$ is its first operand (the lower bound was $-0$ on the POWER9 jobs under
QEMU, and the assembly of Compiler Explorer shows six \code{xsmaxcdp} in the
product and no other maximum; not run on a processor). The lower bound of
\code{[-1, 0] * [-1, 0]} is therefore $+0$ with the SSE2 intervals, $-0$ with
the NEON ones, and $+0$ or $-0$ with the FPU ones, by the compiler and the
target; so is the upper bound of \code{[-1, 0] * [0, 1]}. $+0$ and $-0$ are
the same number: \code{==} and \code{<} do not see it, and a program does not
read it (\pxref{sec:zero-sign}). The sign of the zero bound of a quotient and
of a square does not depend on a tie, and is fixed: \code{[0, 1] / [-2, -1]}
has the upper bound $-0$ (the zero bound of the dividend is divided by a
number whose sign the bounds of the divisor choose), and
\code{sqr(interval(-1, 2))} has the lower bound $+0$, which a selection of a
constant sets.
\end{Rem}

\medskip
The other bounds and the flags are those of the former code, as
\file{tests/tools/\allowbreak arith} (the four rounding directions of the caller) and 1.2
million random operands of every kind, with and without the mode that flushes
the subnormals to zero, show. \code{x * y} computes every corner product, as
the SSE2 and the NEON intervals do, so that a corner that is not a bound of
the result and underflows raises the underflow exception
(\code{FE\_UNDERFLOW}), which the former code did not, for 135\,808 of the
1.2 million \code{x * y}, and a corner that overflows raises the overflow
exception (\code{FE\_OVERFLOW}) for 23\,672 more, whose result has an
infinite bound all the same; no operation raises the invalid-operation or the
division-by-zero exception it did not (\pxref{sec:fp-exceptions}).

%%------------------------------------------------------------------------
\section{The control registers of 32-bit x86 processors}
\label{sec:x86-control-registers}
%%------------------------------------------------------------------------
\cindex{MXCSR}\cindex{x87 control word}%

A 32-bit x86 processor has two units that round, each with a direction of its
own: the x87 unit, in the field RC of its control word (bits 11 and 10), and
the SSE instructions, which compute the doubles of GAOL, in the field RC of
MXCSR (bits 14 and 13), where the modes that flush the subnormals to zero are
too (FTZ, bit 15, and DAZ, bit 6). Both have to be upward: the square root
and the \code{fma()} of a C library, and a \code{long double}, follow the
direction of the x87 unit (the elementary functions of CORE-MATH read MXCSR,
as on x86-64, after having read the one of the x87 unit with
\code{fegetround()}: the one of glibc, and $-1$ with the C runtime of Windows
when the two differ), and the addition that shows the direction elsewhere
sees the SSE instructions only. Each operation therefore reads the two
registers, the control word with \code{fnstcw} and MXCSR with
\code{stmxcsr}, in one asm statement (inline assembly with Visual C++), where
it called \code{fegetround()} and read MXCSR.\newinvfive{} That call took
about 80~ns with Visual C++ for x86: \code{x + y} took 81 to 100~ns, and a
line of 15 additions 1\,211 to 1\,546~ns, in the step of the continuous
integration that runs \cmd{gaol\_performance}.

\medskip
The registers are written (\code{fldcw}, \code{ldmxcsr}) only when the
direction of one is not upward or a flush mode is set, and the one that
differs only; MXCSR is written once for the direction and the flush modes.
With \option{GAOL\_PRESERVE\_ROUNDING}, an operation reads them once at its
start, and writes back at its end what it found, in the register that differs
from what the operation left: a program that computes upward already, with
the flush modes clear, pays no write. The SSE2 operations (\code{+},
\code{-}, \code{*}, \code{/}, \code{sqr()}\ldots), which use no x87
instruction, keep MXCSR alone with \option{GAOL\_PRESERVE\_ROUNDING} and leave
the x87 unit as it is: they write it only if it is not upward with the flush
modes clear, and give back what they found only if it was
not.\newinvfive{} A program that computes upward already pays no write there
either, \code{x + y} taking 1.7~ns rather than 8.2 and \code{x * y} 2.5
rather than 8.3 (GCC~15.2, an Intel i7-1185G7, the best of 15 runs of a
loop; a program to nearest pays 0.2~ns more, 8.4 for \code{x + y}).
\code{round\_upward()}, \code{round\_nearest()} and \code{round\_downward()}
write the register whose direction is not the one they set (both, as
\code{fesetround()} does, when neither is), and \code{cleanup()} and
\code{restore\_rounding()} set back the direction that \code{init()} found in
each unit (\pxref{chap:init}). The values are those of the hardware
(\code{0x0800} for upward in the control word, \code{0x4000} in MXCSR), not
the \code{FE\_*} constants of \file{<fenv.h>}, which are \code{0x0400} and
\code{0x0800} for downward and upward with glibc and \code{0x0100} and
\code{0x0200} with the C runtime of Windows and with mingw-w64~13 and later.

\medskip
An x86-64 processor has the two units too, and GAOL reads and writes their
registers itself there, as above, where it restores a direction
(\option{GAOL\_PRESERVE\_ROUNDING}, \code{cleanup()},
\code{restore\_rounding()}, and the operations that set the direction to
nearest and back: \code{midpoint()}, \code{chi()}, the writing and the
reading of an interval), rather than call \code{fegetround()} and
\code{fesetround()},\newinvfive{} which write both registers whatever they
hold, and then MXCSR once more for the direction of the SSE instructions:
three writes and a call, where a program that computes to nearest needs one
write and a program that computes upward none. The register that differs is
written only, with \code{fldcw} or \code{ldmxcsr}, which keep the flags;
Visual C++ for x64, which has no inline assembly and uses no x87 unit, reads
and writes MXCSR alone. mingw-w64, the only one that still called them, gave
reasons of its own to stop: from version 13, its \code{fegetround()} joins the
fields RC of the two units by OR, which is the direction of both only where
they agree (to nearest and upward give upward,
\code{FE\_DOWNWARD | FE\_UPWARD} is \code{FE\_TOWARDZERO}), and its
\code{fesetround()} clears the six exception flags of MXCSR when it changes
its value, so that the flags of the program were lost, in the operations with
\option{GAOL\_PRESERVE\_ROUNDING}, and in \code{cleanup()} and
\code{restore\_rounding()} with it or without it.

%%------------------------------------------------------------------------
\section{Flush-to-zero and denormals-are-zero}\label{sec:ftz}
%%------------------------------------------------------------------------
\cindex{flush-to-zero}%
\cindex{denormals-are-zero}%
\cindex{subnormal numbers}%
\cindex{-Ofast@\option{-Ofast}}%
\cindex{crtfastmath.o@\file{crtfastmath.o}}%

\begin{Def}\label{def:ftz}
The \dfn{subnormal numbers} are the nonzero doubles below $2^{-1022}$ in magnitude.
Processors have modes that flush them to zero. On x86, the
\dfn{flush-to-zero} mode (FTZ) and the \dfn{denormals-are-zero} mode (DAZ) of
the SSE instructions, two bits of the control register MXCSR, make every
subnormal result a zero (FTZ) and read every subnormal operand as a zero
(DAZ), in the comparisons too. On ARM, the mode FZ of the control register
(FPCR, FPSCR) does both.
\end{Def}

\noindent
With one of them set, the bounds of every operation with a subnormal result
or operand miss the exact result, whatever the rounding direction:
the product \code{interval(1e-300) * interval(1e-20)}, the square
\code{sqr(interval(1e-160))} and the difference
\code{interval(3e-308) - interval(2.9e-308)} are all \itv{0}{0}.\newinvfive

\medskip
\noindent
\textbf{How a program gets them.} A program linked with one of the options
\option{-Ofast}, \option{-ffast-math} and \option{-funsafe-math-optimizations}
gets them: GCC and Clang then link \file{crtfastmath.o}, on x86 and on ARM
Linux, whose constructor sets them for the whole program. GCC links it last,
and it runs after GAOL initialized itself; Clang links it first, and the
initialization of a static GAOL, which follows, clears the modes again,
unless GAOL is built with \option{GAOL\_PRESERVE\_ROUNDING}, or is a shared
library, initialized before the program. The option \option{-fno-fast-math}
of \file{gaol.pc} and \code{gaol::gaol} does not prevent it: it keeps
\file{crtfastmath.o} out of the link when it comes after
\option{-ffast-math}, but not after \option{-Ofast} (GCC~9.4 and~13 and
Clang~18 link it all the same, and \option{-fno-fast-math} silences the
\code{\#error} of \file{gaol/gaol\_config.h} against \option{-ffast-math},
which \option{-Ofast} would raise), nor, with GCC~9.4 and~13, after the
option \option{-funsafe-math-optimizations}. Loading a plug-in or a Python
module built with \option{-Ofast} sets the modes too, when it is loaded,
where its compiler linked \file{crtfastmath.o} into it: Clang~18 does, GCC
did before version~13.

\begin{Rem}\label{rem:ofast}
Linking a program that uses GAOL with \option{-Ofast}, \option{-ffast-math}
or \option{-funsafe-math-optimizations}, or loading code built with them,
sets the modes, which makes the bounds wrong unless GAOL defends itself
against them, as below. A program that uses GAOL is linked without these
options, or with \option{-mno-daz-ftz}, and the code that needs them is
compiled apart from it.
\end{Rem}

\noindent
GAOL defends itself in two ways.
\begin{itemize}
\item[$\bullet$] \textbf{The check of each operation clears the modes.} On
  x86 processors, and on ARM processors with GCC and Clang and with the NEON
  intervals of Visual C++, each operation that computes bounds clears the
  modes when it starts, in the check that sets the rounding direction upward:
  the check is one comparison with a subnormal number, which shows the modes
  as well as the direction (\pxref{sec:rounding-direction}), or, with the
  NEON intervals, a read of FPCR, which holds both; the modes stay cleared
  for the rest of the program, as the rounding direction stays upward. The
  check comes before the operation reads, compares or copies a bound, the
  test of the empty set aside, and the compiler is kept from making a
  comparison before it, which it would do, not knowing that the check changes
  how a subnormal compares: with GCC~13 at \option{-O3}, \code{asinpi()}
  compared the bounds it had read with 0 before its check, and with the modes
  set \code{asinpi()} of $[100\cdot 2^{-1074}]$ was $[32\cdot 2^{-1074}]$,
  above the exact value. With \option{GAOL\_PRESERVE\_ROUNDING}, each
  operation sets the modes back as it found them, once it has made its
  result, its maxima and minima of bounds included. The bounds of every
  operation that makes the check are then the ones it gives with the modes
  cleared, the tightest ones, and the modes are cleared, or set back, after
  it;\newinvfive{} \file{tests/rounding\_direction.cpp} checks it for each of
  them. It is the only defence against a plug-in, which sets the modes at any
  time, and against the compilers and processors without
  \option{-mno-daz-ftz}.

  The check costs little. On an Intel Xeon of the Cascade Lake generation,
  with Clang~18, an addition with a subnormal operand and result took as long
  as one of normal numbers (1.2~ns in a chain of dependent additions), as on
  an Intel i7-1185G7, and the check as long as the one of the direction
  alone, $1+2^{-60}$, in front of the addition of two SSE2 intervals. Through
  the library (\cmd{gaol\_performance}, medians of 9 interleaved runs),
  \code{x * y} took 4.0~ns rather than 3.4~ns, \code{sqrt} 11.0~ns rather
  than 10.3~ns and \code{pow(x, 3)} 13.7~ns rather than 13.0~ns, and as long
  with the same check made of a normal number: the cost of its second
  addition, not of the subnormal; \code{x + y}, \code{x - y}, \code{x / y},
  \code{sqr}, \code{exp}, \code{log}, \code{sin} and \code{cos} stayed within
  the noise of the machine. Some x86 processors take a microcode assist, of
  the order of a hundred cycles, for an operation with a subnormal operand or
  result, which each operation would pay there. \cmd{gaol\_performance}
  (\file{tests/performance.cpp}) prints the time of the check and of the same
  check made of $2^{-60}$, a normal double: on the 68 jobs of the continuous
  integration that run on processors (x86-64, i386, arm64 and armhf, under
  Linux, Windows and macOS, Rosetta~2 included), the first took 0.35 to
  1.43~ns and the second 0.40 to 1.38~ns, the two differing by $-0.63$ to
  $+0.44$~ns, either way from one job to another: none of these processors
  took such an assist. Under qemu (ppc64le, riscv64, s390x), both took 25 to
  83~ns.
\item[$\bullet$] \textbf{The link option \option{-mno-daz-ftz}.}
  \file{gaol.pc} and \code{gaol::gaol} give \option{-mno-daz-ftz} to the
  link, where the compiler accepts it: GCC~13 and later on x86, and from 11.4
  and 12.4 in the series 11 and 12 (GCC~9.4, 12.3, GCC for ARM and Clang~18
  refuse it, and the builds test the compiler with a link, not with its
  version). It keeps \file{crtfastmath.o} out of the link, so that a program
  linked with \option{-Ofast} does not get the modes at all, its own code
  included.
\end{itemize}

\medskip
\noindent
On the other processors (POWER, s390x, RISC-V, MIPS\ldots), whose FPU
intervals use the generic \code{round\_upward()} of
\file{gaol/gaol\_fpu\_fenv.h} (a call to \code{fesetround(FE\_UPWARD)}, whose
return value is not read), neither defence exists:\newinvfive{} the mode NI
(non-IEEE) of the FPSCR of POWER and the bit FS of the FCSR of MIPS, modes of
the kind of FTZ and DAZ that may flush the subnormal numbers, are not
cleared, and a program that sets them may get bounds that miss the exact
result for subnormal operands and results. What they do depends on the
processor, whether the toolchains of these processors set them in a program
linked with \option{-Ofast} was not checked, and nothing of GAOL was tested
with them set.

\medskip
\noindent
\textbf{The functions that make no check.} The functions that only read,
compare or copy the bounds of an interval make no check, which made each of
them 3 to 4.5~ns slower in a micro-benchmark on an Intel i7-1185G7, two to
three times the constructor or \code{certainly\_le()}: they compare the
bounds as bounds, whatever the modes.\newinvfive{} A mode that reads a
subnormal operand as a zero (DAZ, FZ) can make two different doubles compare
equal, both zero or subnormal, but never puts two doubles in the wrong order:
the doubles are compared first, and only where they compare equal are they
compared as integers, by their bits, which no mode changes. These functions
are the constructor from two bounds, the relations (\code{set\_contains()},
\code{certainly\_le()}, \code{set\_eq()}, \code{is\_zero()},
\code{lexicographic\_less}\dots), the intersection \code{\&} and the hull
\code{|}, \code{max()}, \code{min()}, \code{abs()}, \code{sign()},
\code{floor()}, \code{ceil()}, \code{integer()} and the other roundings to an
integer, \code{invabs\_rel()}, \code{mig()}, \code{mag()},
\code{is\_canonical()}, \code{feven()}, \code{nb\_fp\_numbers()}, and
\code{inf()} and \code{sup()} of \code{gaol\_ieee1788}; the roundings to an
integer round a subnormal bound from its bits, which the \code{floor()} and
\code{ceil()} of musl would return as it is. They give the results they give
with the modes cleared, and leave the modes as they found them, writing
nothing to the control register, with \option{GAOL\_PRESERVE\_ROUNDING} or
without. On the Intel i7-1185G7, with Clang~18 (\cmd{gaol\_performance},
medians of 9 interleaved runs), this left the constructor as fast as before,
and made \code{x <= y}, \code{x.set\_contains(y)}, \code{x \&= y} and
\code{x | y} 0.5 to 0.8~ns slower, and \code{floor()} 1.7~ns. Before, under
denormals-are-zero, \code{max()} of
$\itv{3\cdot2^{-1074}}{100\cdot2^{-1074}}$ and $[200\cdot2^{-1074}]$ was the
first one, \code{abs(interval(-1e-309, -1e-310))} stayed negative,
\code{interval(}$2^{-1073}$\code{, }$2^{-1074}$\code{)} kept its bounds in
the wrong order, and the hull \code{|} of the same two intervals, inline in
the headers of GAOL, depended on how the compiler of the program arranged its
comparisons.

\medskip
\code{midpoint()}, \code{split()} and \code{chi()}, which compute, and the
writing of an interval (\code{operator<{}<}, \code{intervalToText()}) make
the check, which clears the modes, or sets them back after them with
\option{GAOL\_PRESERVE\_ROUNDING} (1.3~ns more for \code{midpoint()}): the
midpoint of $\itv{2\cdot2^{-1074}}{4\cdot2^{-1074}}$ was 0, outside it, and
a subnormal upper bound was written rounded to nearest, below the bound,
where the C library writes a subnormal 0 under denormals-are-zero with gdtoa
(FreeBSD, macOS) and with the Debug runtime of Visual C++
(\file{tests/rounding\_direction.cpp} checks all of them).

\medskip
\noindent
\textbf{What remains.} The reader, which makes no check of the modes, takes
the numbers apart by their bits, and compares the two numbers of
\code{<a, b>}, which have to be the same double, by theirs too,\newinvfive{}
so that under denormals-are-zero \code{<1e-310, 1e-309>} is refused, as
without it (\pxref{sec:input-format}). The unary minus makes no check either,
but it only exchanges the stored bounds, which the modes do not change, as
they do not change the test of the empty set, made before the check of every
operation (an operation with an empty operand may return before its check,
the modes left as it found them), nor the test by which \code{hausdorff()}
returns $+\infty$ before its check when a bound is infinite in one of its two
intervals only. The modes of other processors, and of ARM with the FPU
intervals of Visual C++, are neither checked nor cleared, though the
functions above compare the bounds right there too, and \code{operator<{}<}
writes a subnormal bound outward, as $2^{-1022}$ away from zero and 0 toward
it; GCC links \file{crtfastmath.o} for none of the other processors GAOL is
tested on.

%%------------------------------------------------------------------------
\section{The floating-point exceptions}\label{sec:fp-exceptions}
%%------------------------------------------------------------------------
\cindex{floating-point exceptions}\cindex{exceptions!floating-point}%
\cindex{flags!floating-point}%

A program starts with the floating-point exceptions masked, and GAOL leaves
them as they are: an operation that raises one only sets its flag while the
exception is masked.

\begin{Rem}\label{rem:fp-exceptions-enabled}
A program that enables the floating-point exceptions, with
\code{feenableexcept()} of the GNU C library, \code{\_controlfp\_s()} of
Visual C++ or \code{fesetenv()} with an environment that traps, disables
them while GAOL computes, and enables them again afterwards if it wants them: the
operations of GAOL raise some, as listed below, and the processor then stops
the program (with \code{SIGFPE} on Linux) instead of letting the operation
give its bounds.\newinvfive
\end{Rem}

\noindent
What the operations raise:
\begin{itemize}
\item \code{log(interval(0,~1))}, which is \itv{-\infty}{0}, raises the
  divide-by-zero exception; \code{interval(1e308)*10}, which is
  \itv{\text{DBL\_MAX}}{+\infty}, raises the overflow exception.
\item Reading a number above the largest double (\code{1e400}, \code{1.8e308},
  \code{0x1p1024}) raises no overflow exception. The reader used to give it to
  \code{strtod()}, which computes $+\infty$ and raises it, so that
  \code{1e400*0}, which is \itv{0}{0}, raised the flag that stands for an
  infinite bound, and a program that had enabled the exception died while it
  read. The interval read is \itv{\text{DBL\_MAX}}{+\infty}; an operation on
  it raises the flag where its own result overflows (\code{1e400*2},
  \code{exp(1e400)}), as for any other interval. Nor does reading a number
  below the least normal double (\code{1e-310}, \code{1e-400},
  \code{0x1p-1075}) raise the underflow exception, which \code{strtod()}
  raises when it rounds such a number to a subnormal or to $0$.\newinvfive
\item Almost every operation raises the inexact exception: most bounds are
  rounded results, and where GAOL checks the rounding direction with an
  addition, $1+2^{-60}$ or $1+(2^{-1060}+0)$, whose result is inexact
  (\pxref{sec:rounding-direction}), even an operation whose bounds are exact
  raises it. On x86, the addition of the subnormal $2^{-1060}$ also raises
  the denormal-operand flag of the SSE instructions, which is none of the five
  flags of IEEE~754 (not in \code{FE\_ALL\_EXCEPT}), and its exception stops
  the program where the program enabled it; and with flush-to-zero set, the
  sum $2^{-1060}+0$, flushed to zero, raises the underflow exception
  (\code{FE\_UNDERFLOW}) too, before the check clears the mode. With the NEON
  intervals the check is a read of FPCR (\pxref{sec:neon}), which raises no
  flag: \code{+}, \code{-}, \code{/} and \code{sqrt()} raise the inexact one
  only when a bound they compute is rounded, \code{*}, \code{sqr()} and
  \code{pow(x, 2)}, which is \code{sqr()}, also when a corner product that is
  not a bound of the result is (\code{[-a, 128] * [-a, 4]} for
  $a = 1 + 2^{-52}$ has the exact bounds $-128a$ and 512, and the corner
  $a^2$ is rounded).\newinvfive
\item An empty operand raises no invalid-operation exception, though the
  bounds of the empty interval are NaN: every operation, relation, function
  and output of \code{gaol::interval} and of \code{gaol\_ieee1788} takes it
  without raising that exception. \code{is\_empty()} and the relations
  compare the bounds with the quiet comparisons of \file{<cmath>}
  (\code{std::islessequal()}\dots), each one instruction, as \code{<=} is,
  with GCC and Clang on x86-64, and so does the first comparison of the
  constructor \code{interval(l,~r)}, which gives the empty set for NaN bounds.
  With Visual C++, where these are calls to the C library, which made
  \code{x~*~y} and \code{sqrt()} 34 to 61\,\% slower on x64, GAOL compiles
  them itself into the same instruction, \code{ucomisd}, on x64 and on x86
  with SSE2; on arm64 they stay calls. With the invalid-operation exception
  enabled, \code{is\_empty()} of an empty interval killed the program, and so
  did about 50 operations with an empty operand, which compared its NaN
  bounds: \code{x~\&~y} for an empty \code{y}, \code{sqrt}, \code{exp},
  \code{min}, \code{max}, \code{floor}, \code{set\_contains()},
  \code{set\_disjoint()}, the output of the empty set\dots; and
  \code{interval::emptyset()} itself in a build without optimization. A NaN
  the program gives GAOL, the empty set for \code{interval(d)}, raises it no
  longer either: \code{interval(NAN)}, \code{interval(NAN,~1)},
  \code{x~=~NAN}, \code{x~+=~NAN} (and \code{-=}, \code{*=}, \code{/=},
  \code{\%=}), \code{x~\&=~NAN}, \code{x~<=~NAN} and \code{pow(x,~NAN)}
  compared it with \code{<} or \code{<=}. An empty operand may still raise
  the inexact exception: the sum of two empty intervals does where GAOL
  checks the rounding direction with an addition. The vector comparisons of
  64-bit ARM (\code{fcmgt}, \code{fcmge}) signal the invalid-operation
  exception on a NaN: the NEON operations tell the empty set, the infinite
  bounds and the zeros from the bits of the bounds, and compare no bound as a
  double before they know it is no NaN.\newinvfive
\item A compiler may compile a quiet comparison into a signaling one, or
  compute a comparison that follows a test of the empty set before it, and
  GAOL is written against the cases seen. GCC turns a choice made on a quiet
  comparison into a conditional move with a signaling comparison on 32-bit
  ARM (\code{vcmpe}, GCC 12 to 14; GCC bug 52258,
  \url{https://gcc.gnu.org/bugzilla/show_bug.cgi?id=52258}, is of this kind)
  and on POWER9 (\code{xscmpgedp}, GCC 13 with \option{-mcpu=power9}): there
  \code{is\_empty()}, the intersection and the constructors of the FPU
  intervals tell NaN bounds with \code{std::isunordered()} first, which stays
  quiet, so that the choice a program makes on \code{is\_empty()}
  (\code{x.is\_empty() ? a : b} for an empty \code{x}), \code{x~\&=~y} for an
  empty \code{y} and \code{floor()}, \code{max()}\dots{} of an empty interval
  raise nothing (the armhf jobs of the continuous integration and a POWER9
  job under qemu check them, built for size too, where GCC calls the
  intersection rather than inlining it). GCC~9.4 at \option{-O3} computed the
  comparisons of the bounds of \code{x~\&=~y} before its test of the empty
  set, in a loop of intersections: they are quiet ones. And GCC vectorizes
  loops of comparisons of bounds into vector comparisons that raise the
  exception for an empty element: GCC~9.4 and~13 where AVX is on
  (\option{-mavx}, or \option{-mfma}, which the three builds give where the
  processor has FMA) and GCC 13 to 15 for 32-bit x86 with SSE2 alone
  (MinGW-w64) make the quiet comparisons signaling predicates
  (\code{vcmpnlepd}, \code{cmpnlepd}\dots), GCC~14 for POWER8 uses
  \code{xvcmpgedp}, and GCC 11 to 15 for 64-bit ARM compute the second and
  third comparisons of the constructor for an empty element too. With GCC and
  the FPU or the NEON intervals, on every processor, the bounds
  \code{is\_empty()} compares, those of the relations that do not call it
  (\code{straddles\_zero()}, \code{set\_contains(d)}\dots) and the doubles
  given to the constructors go through an empty asm statement, which keeps
  GCC from vectorizing such loops: loops of \code{x.set\_contains(d)} and of
  \code{x.straddles\_zero()} took 0.8 and 1.3~ns per element rather than
  0.5~ns with GCC~13 on x86-64, about as long as with GCC~9.4, which did not
  vectorize them and took as long as before. A choice the program makes on a
  relation, \code{x.set\_contains(y) ? a : b}, and a loop of relations other
  than those of \file{tests/rounding\_direction.cpp} are not known to raise
  the exception, but nothing keeps a compiler from compiling them
  so.\newinvfive
\item Nonempty operands raised the invalid-operation exception as well, and
  the cases found no longer do. With the SSE2 intervals, the products
  \code{interval(0) * interval::one\_plus\_infinity()} and
  \code{interval::positive() * interval(0)}, which multiply a zero bound by an
  infinite one, did (the FPU intervals gave the same products without it),
  and so did, through them, \code{pow} of the base \code{interval(1)} and the
  exponent \code{interval::one\_plus\_infinity()}; \code{pow} did in every
  build, through the \code{pow} of CORE-MATH, for an exponent of extreme
  magnitude (\code{pow(interval(1, 2), interval(4.9e-324))},
  \code{pow(interval(1, 2), interval(1e-300))}, an exponent interval with a
  bound near \code{DBL\_MAX} such as \itv{10^{300}}{10^{308}}), which the
  CORE-MATH of GAOL v5 no longer does, and for a power next to the largest
  double (\code{pow(interval(2), interval(1023.9999))},
  \code{pow(interval(0x1.000003f82a1fp+0), 3000000001LL)}), where CORE-MATH
  multiplied its first approximation by $2^{1024}$, an infinity, and computed
  $\infty-\infty$ (changed in the vendored sources, see \file{3rd/README.md});
  and with the SSE2 intervals, \code{div\_rel(K, J, I)} did where it divided
  an infinite bound of \code{K} by an infinite bound of \code{J} in a half of
  the register it did not keep ($K=\itv{2}{+\infty}$ and
  $J=\itv{-\infty}{2}$), and raised the division-by-zero exception where it
  divided a bound by a zero one there ($K=\itv{-2}{-2}$ and $J=\itv{-2}{0}$),
  and so did the division, the relational division and the inverse of a
  divisor with a zero bound (\code{interval(1, 2)/interval(0, 2)},
  \code{interval(1, 2)\%interval(-2, 0)}, \code{inverse(interval(0, 2))},
  \code{pow(interval(0, 2), -1)}). This list is not exhaustive.\newinvfive
\item The overflow and the division-by-zero exceptions are no longer raised
  for a result that has no infinite bound. Some functions computed beyond the
  largest double on the way, or divided by zero, and raised them for a finite
  result, the empty set included: \code{sqrt(interval(1,~DBL\_MAX))},
  \code{nth\_root(interval(DBL\_MAX),~3)} and the roots of order 4 and more,
  and the relational \code{sqrt\_rel()} and \code{nth\_root\_rel()} (the
  product that proves a root exact is $2^{1024}$ for the square root of
  \code{DBL\_MAX}), \code{sin(interval(-1,~DBL\_MAX))} and \code{cos()} (the
  width of the interval), \code{tan(interval(DBL\_MAX))} ($x+\pi/2$),
  \code{div\_rel()} (the quotient by a divisor next to $2^{-1074}$),
  \code{atan\_rel()}, \code{atan2\_rel\_x()} and \code{atan2\_rel\_y()},
  \code{pow\_rel()} and \code{pow\_exponent\_rel()} (the powers beyond the
  largest double that their searches try, and the exponents next to
  $2^{-1074}$), and \code{atanh\_rel()} (atanh at $-1$ and $1$). A quotient,
  a product or a power that would overflow is the infinity, or the largest
  double, that the operation gives, without being computed, where it is only
  an intermediate value: the exceptions still tell an infinite bound of the
  result, as above. This list is not exhaustive either, and the tests check
  these functions and the others of the same kind on operands next to the
  limits of the doubles (\file{tests/rounding\_direction.cpp}).\newinvfive
\end{itemize}

\noindent
An empty operand, and the nonempty operands tested, raising no
invalid-operation exception therefore does not make that exception safe to
leave enabled: every exception stays disabled while GAOL computes.\newinvfive

\medskip
The flags tell nothing of the results: after most operations of GAOL,
\code{fetestexcept(FE\_INEXACT)} is raised whatever the result (not after
\code{-x}, \code{abs()}, \code{\&}, \code{|}, \code{floor()}, \code{max()} or
\code{min()}, which round nothing and check no rounding direction, nor after
\code{+}, \code{-}, \code{/} and \code{sqrt()} with the NEON intervals when
their bounds are exact), and \code{FE\_OVERFLOW} or \code{FE\_DIVBYZERO}
stands for an infinite bound of the interval, or an infinite result of
\code{width()} or \code{rad()}, not for an infinite number the program
computed: \code{midpoint()} and \code{mid()} of
\itv{\text{DBL\_MAX}}{\text{DBL\_MAX}}, which overflowed the sum of the
bounds, raise none. IEEE 1788-2015 \cite{IEEE1788:2015} leaves the flags of
its operations unspecified. A program that reads them for its own
computations clears them (\code{std::feclearexcept(FE\_ALL\_EXCEPT)}) just
before, and reads them before its next operation of GAOL.\newinvfive

\medskip
The initialization of GAOL, which runs before \code{main()}, sets the default
environment (\code{fesetenv(FE\_DFL\_ENV)}), which masks every exception,
unless GAOL is built with \option{GAOL\_PRESERVE\_ROUNDING}: the
initialization then leaves the whole environment as it found it, and an
exception a static object enabled before it stays enabled. A program enables
the exceptions in \code{main()} or later, whichever the build, not in a
static object initialized before GAOL. \code{cleanup()} leaves the exceptions
and the flags as they are, and so does every operation, whichever way GAOL is
built: it leaves the exceptions the program enabled enabled, and the flags it
raised raised, and adds the flags it raises itself. With
\option{GAOL\_PRESERVE\_ROUNDING} and the SSE2 intervals, \code{+}, \code{-},
\code{*}, \code{/}, \code{\%}, \code{div\_rel()}, \code{sqr()},
\code{inverse()}, the integer powers and \code{+=~d}, \code{-=~d},
\code{*=~d}, \code{/=~d}, \code{\%=~d} wrote the SSE control register whole,
every exception masked and every flag cleared: an exception the program had
enabled was masked again by the first of these operations, and the flags it
had raised were lost (\file{tests/rounding\_direction.cpp} checks
it).\newinvfive

%%------------------------------------------------------------------------
\section{Initialization and cleanup}\label{chap:init}
%%------------------------------------------------------------------------
\cindex{initialization}\cindex{cleanup}%

GAOL initializes itself when the program starts, before \code{main()}:
\code{init()} need not be called. Since Release 1.0 of GAOL, there is an
automatic initialization and cleanup, so that GAOL works even when the program
calls neither function. Since GAOL v5, the initialization is made with every
compiler, Visual C++ included: each file including the headers of GAOL holds
a static object whose constructor initializes GAOL the first time, as
\file{<iostream>} initializes the standard streams, so that GAOL is
initialized before the static objects that such a file defines after its
\code{\#include}. Each operation also sets the rounding direction upward
itself (\pxref{sec:rounding-direction}), so that \code{init()} no longer has
to be called again after the calling code changed it.\newinvfive

\begin{Rem}\label{rem:cleanup-rounding}
\code{cleanup()}, on the other hand, is to be called by the program right
after its last use of GAOL, at the latest at the end of \code{main()}, with
every compiler.\newinvfive{} Unless GAOL is built to preserve the rounding
direction, the initialization leaves the rounding direction upward for the
whole program, not only for the operations of GAOL, and \code{cleanup()} sets
back the one the program started with. The automatic cleanup does so too, but
only when the program ends, once \code{main()} has returned and the static
objects are destroyed, or when the library is unloaded: until then, the code
after the last use of GAOL (the rest of \code{main()}, the destructors of
static objects, the functions registered with \code{std::atexit()}) computes
its doubles rounded upward.
\end{Rem}

\begin{Ex}\label{ex:cleanup}
The program below computes an interval, calls \code{cleanup()} right after
its last use of GAOL, and then prints a double of its own, which is computed
to nearest. It calls no \code{init()}.
\begin{example}
#include <cstdio>
#include <iostream>
#include <gaol/gaol.h>
using namespace gaol;

int main()
{
  @rem^// No gaol::init(): GAOL initialized itself before main()~
  interval x(1, 2);
  std::cout << pow(x, 3) + sin(x) << std::endl;

  gaol::cleanup();   @rem^// always, right after the last use of GAOL~
  std::printf("%g\n", 1.0 / 3.0);   @rem^// the program's own doubles, to nearest~
  return 0;
}
\end{example}
\end{Ex}

\medskip
Only the first call of \code{cleanup()} sets the direction back: an operation
of GAOL after it sets the direction upward again, and
\code{restore\_rounding()} sets it back as many times as the program
needs.\newinvfive{} What the initialization allocated is freed by the
automatic cleanup only, after the static objects are destroyed, so that the
expressions of \file{gaol/parser/gaol\_expression.h} still alive when the
program calls \code{cleanup()} remain valid.

\medskip
The direction set back is the one GAOL found when it initialized itself. It
is not the one the program started with when an operation of GAOL came
first: in the initialization of a static object of a file that uses GAOL
without including its headers, through a library of the program for
instance, or, with Clang, of an \code{inline} variable or a static member of
a class template, whose initialization C++ does not order with the rest of
the file.

\medskip
In a build that takes the AVX-512 path (\pxref{sec:avx512}), the
initialization stops the program, with a message on the standard error, when
the processor or the system does not run the AVX-512 instructions, even in a
program that computes nothing.\newinvfive

\medskip
The constants of GAOL ($\pi$, the masks of its SSE2 operations) are
initialized when compiling: a static object of the program may compute
intervals, before the files of GAOL are initialized, which with the static
library of the CMake build come after those of the program, except with
MinGW-w64.\newinvfive

\begin{deffun}{bool}{init}{(int \var{dbg\_lvl} = 0)}
Initializes the variable \code{debug\_level} (\pxref{sec:debugging-facilities})
to the value of \var{dbg\_lvl}.

Unless GAOL was built to preserve the rounding direction, the first call notes
the rounding direction it finds, which \code{cleanup()} sets back, then sets
the floating-point environment to its default one, then the rounding
direction towards $+\infty$. In addition, it sets the number of digits to
display for interval bounds to 16. In a build that takes the AVX-512 path,
the first call stops the program, with a message on the standard error, when
the processor or the system does not run the AVX-512 instructions, even in a
program that computes nothing.\newinvfive

Returns \code{true} if the library was not already initialized and
\code{false} otherwise. After its first call, the function sets the debugging
level to its new value and, unless the rounding direction is preserved, the
rounding direction upward again: the C runtime of Windows resets it when the
program starts, after the library initialized itself.
\end{deffun}

\begin{deffun}{bool}{cleanup}{(void)}
Unless GAOL was built to preserve the rounding direction, sets back the
rounding direction that the first call of \code{init()} found: the one of the
program when GAOL initialized itself, to nearest when the program starts. On
x86 processors, the x87 unit and the SSE instructions each get theirs back.
The rest of the floating-point environment is left as it is, the exception
flags raised since in particular. The rounding direction being one of each
thread, it is set back in the thread calling \code{cleanup()}.\newinvfive

The program calls it right after its last use of GAOL: an operation of GAOL
after it sets the rounding direction upward again. The automatic cleanup
calls it when the program ends, or when the library is unloaded, if the
program did not, and frees what the initialization allocated, after the
static objects of the program are destroyed: the expressions of
\file{gaol/parser/gaol\_expression.h} still alive when the program calls
\code{cleanup()} remain valid.

Returns \code{true} the first time it is called and \code{false} afterwards:
only the first call sets back the rounding direction.
\end{deffun}

\begin{example}
#include <gaol/gaol.h>

int main(void)
{
  gaol::init(1); @rem^// First level of debugging requested;~
                 @rem^// without it, GAOL initializes itself~

  @emph^[Some code using interval arithmetic]~

  gaol::cleanup(); @rem^// Always, right after the last use~
                   @rem^// of GAOL: sets back the rounding~
                   @rem^// direction~
  return 0;
}
\end{example}

\begin{deffun}{void}{restore\_rounding}{(void)}
Sets back the rounding direction that the first call of \code{init()} found,
as \code{cleanup()} does at its first call, but as many times as the program
calls it.\newinvfive{} A program that goes back to GAOL after
\code{cleanup()}, whose operations set the direction upward again, calls it
after each use. Like \code{cleanup()}, it sets back the direction of the
calling thread only, and leaves the rest of the floating-point environment as
it is. With \option{GAOL\_PRESERVE\_ROUNDING} it does nothing, GAOL leaving
the direction of the program alone (\pxref{sec:upward-effects}).
\end{deffun}

%%------------------------------------------------------------------------
\section{The rounding functions of the FPU}\label{chap:fpu}
%%------------------------------------------------------------------------
\cindex{rounding!functions}%

GAOL provides functions to set the rounding direction of the floating-point
unit. The operations of GAOL do not need them: each sets the rounding
direction upward itself (\pxref{sec:rounding-direction}).

\medskip
On x86 processors, with GCC, Clang and Visual C++, they write the control
registers of the x87 and SSE units directly (\code{fnstcw}/\code{fldcw},
\code{stmxcsr}/\code{ldmxcsr}), rather than calling \code{fesetround()},
which cost 50 to 250~ns per call with the C runtime of Windows; with Visual
C++ for x64, the SSE register only, the x87 unit being unused there
(\pxref{sec:x86-control-registers}).\newinvfive{} Elsewhere (ARM, POWER,
s390x, RISC-V), they call \code{fesetround()}, except on 64-bit ARM
processors with the NEON intervals, and on 32-bit ARM processors with GCC
(\pxref{sec:neon}), where they write the rounding direction field of the
control register, FPCR or FPSCR, themselves, and only when the direction
changes.\newinvfive

\begin{deffun}{void}{round\_downward}{(void)}
Sets the rounding direction mode towards $-\infty$.
\end{deffun}

\begin{deffun}{void}{round\_nearest}{(void)}
  Sets the rounding direction mode to the nearest/even.
\end{deffun}

\begin{deffun}{void}{round\_upward}{(void)}
  Sets the rounding direction mode to $+\infty$.
\end{deffun}

%%------------------------------------------------------------------------
\section{Discussion}
%%------------------------------------------------------------------------

This chapter has established how GAOL v5 handles the rounding direction: each
operation that computes a bound checks it with an addition and sets it upward
when it is not, so that the bounds are right whatever direction the calling
code left, where GAOL~4 trusted the user to set it. The price of this choice
falls on the doubles of the program, which are rounded upward from the
initialization of GAOL until \code{cleanup()}: the quotients, the decimal
conversions of the C library, the sums and the error-free transformations all
change, and the program answers by calling \code{cleanup()} right after its
last use of GAOL, by reading and writing its numbers through GAOL, and by
computing to nearest in a guarded block what needs it. The build with
\option{GAOL\_PRESERVE\_ROUNDING} removes the price from the program and puts
it on the operations, 1.2 to 1.6 times slower with the AVX-512 path and
several times without it; the control registers are then written only where
they have to be, and the AVX-512 instructions, whose rounding direction is in
the instruction, are taken where the machine has them.

\medskip
The same check defends the bounds against the modes that flush the subnormal
numbers to zero, which a program linked with \option{-Ofast} gets without
asking, and which GAOL clears on x86 and on ARM, while \option{-mno-daz-ftz}
keeps them out of the link where the compiler accepts it. The floating-point
exceptions, on the other hand, are left to the program, which disables them
while GAOL computes: the operations raise the inexact flag almost always, the
overflow and division-by-zero flags for an infinite bound, and no longer the
invalid-operation one for an empty operand in the cases found, but nothing
keeps a compiler from compiling a quiet comparison into a signaling one.
Several points stay open in the sources: the other processors (POWER, s390x,
RISC-V, MIPS) have no defence against their flush modes, and nothing was
tested with them set; the FPU intervals of POWER, s390x, RISC-V and
MinGW-w64 for arm64 are not measured; the slower line of 13 operations with
the direction preserved on Visual C++ for x86 has no explanation yet; and the
processors that do not halve their AVX-512 instructions are expected, not
measured, to take less time on that path.

\medskip
With the floating-point environment in place, the next chapters turn to the
intervals themselves: Chapter~\ref{chap:creation} describes how an interval
is created, from numbers or from text, and assigned.

%% !TEX root = ../manual_gaolv5.tex
% Chapter 5.
%% Chapter 5: Creating intervals. Sources: gaol.tex of f5bf180 lines 2409-2802
%% (interval creation and assignment, assignment combined with an operation,
%% interval constants) and doc/using.md lines 843-898 (integers: the facts
%% about integer bounds).
\chapter{Creating intervals}\label{chap:creation}

An interval of GAOL is an object of the class \code{interval}. A program
creates it from its two bounds, from one number, from nothing, from another
interval, or from the text that writes it; it can also build it from a
midpoint and a radius. Once created, an interval receives new values by
assignment, alone or combined with an operation, and a few intervals of
constant use, the empty set, the whole real line, the enclosures of $\pi$,
are provided ready-made by the class.

\medskip
The way a value enters an interval is not a detail. A \code{double} is taken
as it is, an integer as the integer it is, which may be no double, and a
string as the interval it writes, rounded outward. A decimal constant written
in the program, as $0.1$, is rounded by the compiler before it reaches the
constructor, and an interval built from it does not contain the number the
programmer wrote. The constructors of GAOL v5 answer every input that makes
no interval, bounds in the wrong order, a NaN bound or an infinite point, with
the empty set, and they refuse the silent conversions that GAOL~4 accepted,
from a \code{double} and from a string. Each of these differences is marked
in the text.

\medskip
In this chapter, we describe the creation of intervals and their assignment.
The first section presents the four constructors, the function
\code{textToInterval()}, the static method \code{midrad()} and the intervals
of one infinite point, which GAOL v5 does not have. The second section covers
the straight assignment, the third the assignment combined with an operation
(\code{+=}, \code{\&=}\ldots) and the fourth the interval constants of the
class. The format of the strings that \code{textToInterval()} reads is the
subject of Chapter~\ref{chap:io}.

\section{Constructors}\label{sec:constructors}
\cindex{interval!constructors}%

The class \code{interval} has four constructors. The SSE2 and NEON intervals
also have an explicit constructor from the register that holds the stored
bounds, \code{interval(const \_\_m128d\&)} or
\code{interval(const float64x2\_t\&)}, the lower bound negated in the first
lane, which gives the empty set for a lane of $-\infty$; it serves the code of
GAOL (\pxref{sec:kernels}). A fifth way to create an
interval, from the text that writes it, is the function
\code{textToInterval()}, and a sixth, from a midpoint and a radius, is the
static method \code{interval::midrad()}. We describe the constructors first,
then the two functions, and we end with the intervals of one infinite point.

\subsection{The four constructors}

\begin{itemize}
\item[$\bullet$] \textbf{From the two bounds.} The left and the right bounds
  are given:
  \begin{example}
    interval x(@var^l~,@var^r~);
  \end{example}
  where \var{l} and \var{r} are \code{double}s or of a type that is castable
  into a \code{double}. The result is the empty set when \itv{l}{r} is not an
  interval: a left bound $+\infty$, a right bound $-\infty$, bounds in the
  wrong order ($l>r$) and NaN bounds all give the empty set.\newinvfive{}
  An integer bound of any integer type but \code{bool}, the other bound being
  an integer, a \code{double} or a \code{float}, is taken as the integer it
  is: the left bound is the largest double at most \var{l}, the right one the
  smallest double at least \var{r}. Thus \code{interval(1, 9007199254740993LL)}
  contains $2^{53}+1$, which no double is, and which its conversion to a
  double made $2^{53}$ or $2^{53}+2$ according to the rounding direction. The
  result is the empty set for $l>r$ compared as numbers:
  \code{interval(9007199254740993LL, 9007199254740992LL)} is empty, though
  both bounds lie between the same two doubles.\newinvfive
\cindex{empty!interval: constructing}%

\item[$\bullet$] \textbf{From one number.} A degenerate interval, a point, is
  built from one bound:
  \begin{example}
    interval x(@var^v~);
  \end{example}
  This is equivalent to \code{interval x(\var{v},\var{v});}, and gives the
  empty set for an infinite or NaN \var{v}. The constructor is
  \code{explicit}: a \code{double} no longer converts silently to an
  interval. The statements \code{interval x = 0.5;}, \code{return 0.5;} in a
  function returning an interval, \code{gaol::sin(0.5)} and \code{min(x, 1.0)} do
  not compile, where \code{interval x(0.5)}, \code{interval(0.5)},
  \code{sin(interval(0.5))} and \code{min(x, interval(1.0))} do. The
  assignment (\pxref{sec:assignment}), the operators \code{\&=}, \code{|=}
  and the other assignments combined with an operation
  (\pxref{sec:combined-assignment}), the arithmetic operators and the
  relation symbols (\pxref{sec:relational-symbols}) take a \code{double} as
  they take \code{interval(\var{d})}.\newinvfive{} An integer \var{n} of any
  integer type but \code{bool} gives the tightest interval containing it:
  \itv{n}{n} where $n$ is a double, as every \code{int} is, and the two
  doubles around it otherwise. \code{interval(9007199254740993LL)} is
  \itv{2^{53}}{2^{53}+2}, where the conversion of $2^{53}+1$ to a double gave
  \itv{2^{53}}{2^{53}} or \itv{2^{53}+2}{2^{53}+2}, neither of which contains
  it, and \code{interval(ULLONG\_MAX)} contains $2^{64}-1$. This constructor
  is \code{explicit} too. The same operations take an integer as they take
  \code{interval(\var{n})}; the relations compare it as the integer it is
  (\pxref{sec:relational-symbols}), and \code{\&=} gives
  \code{interval(\var{n})} where the interval contains $n$, the empty set
  otherwise.\newinvfive

\item[$\bullet$] \textbf{Without any bound.} The default constructor builds
  the whole real line:
  \begin{example}
    interval x;
  \end{example}
  This is equivalent to:
  \begin{example}
    interval x(-GAOL_INFINITY,+GAOL_INFINITY);
  \end{example}
  \noindent where \code{GAOL\_INFINITY} is the infinity of the \code{double}
  format (\pxref{sec:floating-point-constants}).

\item[$\bullet$] \textbf{By copy.} The copy constructor duplicates an
  existing interval:
  \begin{example}
    interval x(-12,12), y=x, z(x);
  \end{example}
\end{itemize}

\begin{Rem}
\cindex{integer!as a bound}%
The integers are taken as the integers they are whatever the rounding
direction, and for every integer type but \code{bool}.\newinvfive{} Three
cases still convert to a \code{double}: an unscoped enumeration, an
\code{\_\_int128} outside the GNU modes of GCC and Clang, and a \code{long
double} bound next to an integer, which takes \code{interval(double, double)}.
An integer that is a double gives the result the double gives, and, for the
types of 53 bits or less (\code{int}, \code{short}, \code{char}\ldots),
through the same code. The constructor \code{interval(\var{n})} is
\code{explicit}, as the constructor from a \code{double}. The other
operations of GAOL that take an integer are described in
\exref{sec:integers}.
\end{Rem}

\subsection{Intervals read from a string}

The function \code{textToInterval()} reads a string that writes an interval
in the format of the input (\pxref{sec:input-format}), and returns that
interval. It takes the place of the constructors from strings of GAOL~4.

\begin{deffun}{interval}{textToInterval}{(const std::string\& \var{s})}
  Returns the interval the string \var{s} writes, read with the names of the
  functions of GAOL (\pxref{sec:input-format}). It is in the namespace
  \code{gaol}, under the name IEEE 1788-2015 gives it.
  \begin{example}
interval x = textToInterval("[-23, inf]"),
         y = textToInterval("[5*0.1+dmin, 89*sinh(2.1)]");
  \end{example}
  If the string does not comply with the expected format, the exception
  \code{std::invalid\_argument} is thrown, and \code{std::domain\_error} if it
  calls a function with an argument the function does not take, as
  \code{"nth\_root(8, 1.5)"}, when the library was compiled with exception
  support, which is the default of the three builds. Without exception
  support, the variable \code{errno} is set to \code{-1}, and the function
  \code{gaol\_error()} prints an error message before the program aborts
  (Chapter~\ref{chap:errors}). The \code{textToInterval()} of
  \code{gaol\_ieee1788} reads the names of the functions of the standard, and
  returns the empty set instead (Chapter~\ref{chap:ieee1788}). A null pointer
  is no string: the overload for a \code{const char*}, which
  \code{textToInterval(nullptr)} reaches, throws \code{std::invalid\_argument}
  in both namespaces, and gives the same error without exception support (a
  message, then \code{abort()}).\newinvfive

  GAOL v5 has no constructor from one string. \code{textToInterval()} takes
  the place of the constructor \code{interval(const char*)} of GAOL~4, which
  converted any \code{const char*} implicitly: \code{interval(0)} was
  ambiguous, and \code{x~+~p} compiled for a null \code{p}, and
  crashed.\newinvfive
\end{deffun}

\begin{deffun}{interval}{textToInterval}{(const std::string\& \var{sl}, const std::string\& \var{sr})}
  Returns the interval whose left bound is the one of the interval the string
  \var{sl} writes, and whose right bound is the one of the interval \var{sr}
  writes. The result is the empty set when the lower bound of the first is
  above the upper bound of the second. The comparison is between the
  intervals the two strings write, each rounded outward, not between the
  numbers, where \code{textToInterval("[l,u]")} orders the numbers.
  \begin{example}
 @rem^// x=[-4,4]~
 interval x = textToInterval("[-5,4]+1","[4,6]-[2,3]");
  \end{example}
  It is in the namespace \code{gaol}, and throws \code{std::invalid\_argument}
  as \code{textToInterval(\var{s})} does, for a null pointer in either place
  too. It takes the place of the constructor from two strings of GAOL~4, which
  kept the bounds \itv{2}{1} of \code{interval("2",~"1")}: \code{is\_empty()}
  was true, but its sum with \itv{0}{1} was \itv{2}{2}.\newinvfive
\end{deffun}

\subsection{A midpoint and a radius}

The static method \code{interval::midrad()} builds an interval from its
midpoint and its radius, both given as \code{double}s.

\begin{defstaticmethod}{interval}{interval}{midrad}{(double \var{m}, double \var{r})}
  \begin{operation}
    \op{$\code{midrad}(m,r)=\itv{\roundDn{m-r}}{\roundUp{m+r}}$}
  \end{operation}
  Returns the interval of midpoint \var{m} and radius \var{r}, computed as
  \code{interval(\var{m}) + interval(-\var{r}, \var{r})}. It encloses every
  real of \itv{m-r}{m+r}, where \code{interval(\var{m} - \var{r}, \var{m} +
  \var{r})} rounds both bounds in the same direction, upward unless GAOL is built
  with \option{GAOL\_PRESERVE\_ROUNDING}. A radius $r=0$ gives the point \var{m},
  and $r=+\infty$ the whole line. A negative or NaN \var{r} gives the empty
  set, and so does an infinite or NaN \var{m}, as \code{interval(\var{m})}
  does.\newinvfive
  \begin{example}
cout << interval::midrad(1, 0.25) << ' '
     << interval::midrad(1, -1);
@outputs^[0.75, 1.25] [empty]~
  \end{example}
\end{defstaticmethod}

\subsection{Decimal constants}

Of the three forms in which a number reaches an interval, the \code{double} is
the usual one, and the one that calls for care.

\begin{Rem}
A rational number with a finite expansion in the decimal base may need a
rounding in the binary base. The statement
\begin{example}
  interval x(0.1);
\end{example}
\noindent
does \emph{not} create an interval containing $0.1$. This number has an
infinite expansion in base two, so no significand, whatever its size,
represents it exactly, and the constant $0.1$ is rounded, in all likelihood
to the nearest floating-point number, when the program is compiled
(\pxref{sec:rounding-definitions}, and \exref{sec:common-errors}). The
interval that contains $0.1$ is read from a string:
\begin{example}
interval x = textToInterval("0.1");
\end{example}
\end{Rem}

\subsection{Infinite point intervals}\label{par:inf_inf_note}
\cindex{interval!with infinite bounds}%

The degenerate intervals \itv{+\infty}{+\infty} and \itv{-\infty}{-\infty} do
not exist in GAOL v5.\newinvfive{} IEEE 1788-2015 \cite{IEEE1788:2015} has no
such interval: its constructor fails there and gives the empty set (clauses
10.5.8 and 12.12.7). The constructors of GAOL give the empty set too, as
those of IBEX do, for an infinite point and for all the cases listed above:
\begin{example}
cout << interval(GAOL_INFINITY) << ' '
     << interval(-GAOL_INFINITY,-GAOL_INFINITY) << ' '
     << interval(2,1) << ' '
     << interval(1,GAOL_INFINITY) << endl;
@outputs^[empty] [empty] [empty] [1, inf]~
\end{example}
GAOL~4 built them as intervals of one infinite point, which its operators did
not support. The intervals with one infinite bound, as \itv{1}{+\infty}, are
intervals as any other.

\medskip
Read from a string, the constant \element{inf}, an expression, stands for
\itv{\text{MAX}}{+\infty}, MAX being the largest double, and \element{-inf}
for \itv{-\infty}{-\text{MAX}}. The interval literals whose left bound is
$+\infty$ or whose right bound is $-\infty$ are the empty set, as
\ieee{numsToInterval} of IEEE 1788-2015 has it:
\begin{example}
cout << textToInterval("inf") << ' ' << textToInterval("[inf]") << ' '
     << textToInterval("[-inf, -inf]") << endl;
@outputs^[1.797693134862315e+308, inf] [empty] [empty]~
\end{example}
The string \element{"<-inf, -inf>"}, which GAOL~4 read as
\itv{-\infty}{-\text{MAX}}, throws \code{std::invalid\_argument}: a bound of
\code{<a, b>} must have a double for its value, and the value of \element{-inf}
is \itv{-\infty}{-\text{MAX}}, no point (\pxref{sec:input-format}).\newinvfive

\section{Straight assignment}\label{sec:assignment}
\cindex{assignment}%

An existing interval receives a new value from another interval, from a string,
from a \code{double}, or from an integer (below).

\begin{itemize}
\item[$\bullet$] \textbf{From another interval.} The assignment copies it:
  \begin{example}
    interval x(-12,12),
             y; @rem^// Here, y is @itv^-@infty~^+@infty~~
    y = x; @rem^// Now, y is @itv^-12~^12~~
  \end{example}
\item[$\bullet$] \textbf{From a string.} The string follows the format
  expected for the input (\pxref{sec:reading-intervals}), and
  \code{textToInterval()} reads it:

\begin{example}
interval x;
x = textToInterval("[-inf, 123]");
\end{example}
  GAOL~4 assigned the string itself, \code{x~=~"[-inf,~123]"}. GAOL v5
  refuses this statement, having no constructor from one string.\newinvfive

\item[$\bullet$] \textbf{From a \code{double}.} The operator
  \code{operator=(double)} assigns it as \code{interval(\var{d})}, the
  constructor from one \code{double} being explicit:\newinvfive

\begin{example}
interval x;
x = 1234.5; @rem^// Now, x is @itv^1234.5~^1234.5~~
x = GAOL_INFINITY; @rem^// Now, x is empty~
\end{example}
\end{itemize}

\noindent
The assignments from a string and from a \code{double} follow the rules of
\exref{sec:constructors}: assigning an infinity to an interval empties it. An
integer \var{n} is assigned as \code{interval(\var{n})}, the tightest interval
containing it.\newinvfive

\begin{Rem}
The interface has no method for modifying one bound of an interval: an interval
is created, assigned and replaced as a whole. (The non-const accessors
\code{get\_xmminterval()} and \code{get\_neoninterval()} of the SSE2 and NEON
intervals hand the register of the stored bounds to the code of GAOL; a program
does not write it.)
\end{Rem}

\section{Assignment combined with an operation}\label{sec:combined-assignment}

In this manual, we write \code{self} for the object to which \code{this} is
a pointer in C++.\cindex{this@\texttt{this}}\cindex{self@\texttt{self}} The
assignment operators below combine the value of \code{self} with the value
of their right-hand side, an interval, a \code{double} or an integer, and store the
result in \code{self}.

\defmethodx{interval}{interval\&}{operator\&=}{(const interval\& \var{I})}
\begin{defmethod}{interval}{interval\&}{operator\&=}{(double \var{d})}
  \begin{operation}
    \op{$\code{self}\gets\code{self} \cap I$}
    \op{$\code{self}\gets\code{self} \cap \code{interval}(d)$}
  \end{operation}
Assigns to \code{self} the interval resulting from the intersection of
\code{self} and \var{I} (resp. \code{interval(\var{d})}, empty for an
infinite or NaN \var{d}), which is the empty set when they are
disjoint.\newinvfive

\begin{example}
interval x(-12,12);
x &= interval(-6,23); @rem^// Now, @code^x~ is @itv^-6~^12~~
\end{example}

\end{defmethod}

\defmethodx{interval}{interval\&}{operator\textbar=}{(const interval\& \var{I})}
\begin{defmethod}{interval}{interval\&}{operator\textbar=}{(double \var{d})}
  \begin{operation}
    \op{$\code{self}\gets\hull{(\code{self} \cup I)}$}
    \op{$\code{self}\gets\hull{(\code{self} \cup \code{interval}(d))}$}
  \end{operation}
  Assigns to \code{self} the interval hull of the union of \code{self} and
  \var{I} (resp. \code{interval(\var{d})}, which leaves \code{self} as it is
  for an infinite or NaN \var{d}).\newinvfive

\begin{example}
interval x(-12,12);
x |= interval(-6,23); @rem^// Now, @code^x~ is @itv^-12~^23~~
\end{example}

\end{defmethod}

\defmethodx{interval}{interval\&}{operator+=}{(const interval\& \var{I})}
\begin{defmethod}{interval}{interval\&}{operator+=}{(double \var{d})}
  \begin{operation}
    \op{$\code{self} \gets \code{self} + I$}
    \op{$\code{self} \gets \code{self} + d$}
  \end{operation}

  Assigns to \code{self} the interval resulting from adding \code{self} and
  \var{I} (resp. \var{d}).

\begin{example}
interval x(-12,12);
x += interval(-6,23); @rem^// Now, @code^x~ is @itv^-18~^35~~
\end{example}
\end{defmethod}

\defmethodx{interval}{interval\&}{operator-=}{(const interval\& \var{I})}
\begin{defmethod}{interval}{interval\&}{operator-=}{(double \var{d})}
  \begin{operation}
    \op{$\code{self}\gets \code{self} - \var{I}$}
    \op{$\code{self}\gets \code{self} - \var{d}$}
  \end{operation}
  Assigns to \code{self} the interval resulting from subtracting \var{I}
  (resp. \var{d}) from \code{self}.
\begin{example}
interval x(-12,12);
x -= interval(-6,23); @rem^// Now, @code^x~ is @itv^-35~^18~~
\end{example}
\end{defmethod}

\defmethodx{interval}{interval\&}{operator*=}{(const interval\& \var{I})}
\begin{defmethod}{interval}{interval\&}{operator*=}{(double \var{d})}
\begin{operation}
\op{$\code{self}\gets \code{self} \times\var{I}$}
\op{$\code{self}\gets \code{self} \times\var{d}$}
\end{operation}
Assigns to \code{self} the interval resulting from multiplying \code{self}
and \var{I} (resp. \var{d}).
\begin{example}
interval x(-12,12);
x *= interval(-6,23); @rem^// Now, @code^x~ is @itv^-276~^276~~
\end{example}
\end{defmethod}

\defmethodx{interval}{interval\&}{operator/=}{(const interval\& \var{I})}
\begin{defmethod}{interval}{interval\&}{operator/=}{(double \var{d})}
\begin{operation}
\op{$\code{self}\gets \code{self} / \var{I}$}
\op{$\code{self}\gets \code{self} / \var{d}$}
\end{operation}
Assigns to \code{self} the interval resulting from dividing \code{self} by
\var{I} (resp. \var{d}).
\begin{example}
interval x(-12,12);
x /= interval(-6,23);   @rem^// Now, @code^x~ is @itv^-@infty~^+@infty~~
x /= interval::zero();  @rem^// Now, @code^x~ is @ensuremath^@emptyset~~
\end{example}
\end{defmethod}

\defmethodx{interval}{interval\&}{operator\%=}{(const interval\& \var{I})}
\begin{defmethod}{interval}{interval\&}{operator\%=}{(double \var{d})}
Assigns to \code{self} the interval resulting from dividing \code{self} by
\var{I} (resp. \var{d}), using a \emph{relational division}
(\pxref{sec:relational-arithmetic}).
\begin{example}
interval x(-12,12);
x %= interval(-6,23);   @rem^// Now, @code^x~ is @itv^-@infty~^+@infty~~
x %= interval::zero();  @rem^// Now, @code^x~ is @itv^-@infty~^+@infty~~
\end{example}
\end{defmethod}

\noindent
The arithmetic assignments with a \code{double} (\code{+=}, \code{-=},
\code{*=}, \code{/=}, \code{\%=}) give the result of the same operation with
\code{interval(\var{d})}, choosing the case by the sign of \var{d}; an infinite
or NaN \var{d} gives the empty set, as \code{interval(\var{d})} is empty.\newinvfive

\section{Interval constants}\label{sec:interval-constants}
\cindex{interval!constants}%

Some useful intervals, and some canonical intervals enclosing real constants,
are defined as static methods of the class \code{interval}.
Table~\ref{tab:interval-constants} lists them.

\begin{table}[htbp]
  \centering\renewcommand{\arraystretch}{1.1}
  \caption{The interval constants, static methods of the class \protect\code{interval}} \label{tab:interval-constants}
  \begin{tabular}{lc}
    \hline
    \multicolumn{1}{c}{Function} & \multicolumn{1}{c}{Value}\\
    \hline\hline
    \code{interval::emptyset()} & $\emptyset$\\
    \code{interval::half\_pi()} &
      \itv{\roundDn{\frac{\pi}{2}}}{\roundUp{\frac{\pi}{2}}}\\
    \code{interval::minus\_one\_plus\_one()} & \itv{-1}{1}\\
    \code{interval::negative()} & \itv{-\infty}{0}\\
    \code{interval::one()} & $1$\\
    \code{interval::one\_plus\_infinity()} & \itv{1}{+\infty}\\
    \code{interval::pi()} & \itv{\roundDn{\pi}}{\roundUp{\pi}}\\
    \code{interval::positive()} & \itv{0}{+\infty}\\
    \code{interval::two\_pi()} & \itv{\roundDn{2\pi}}{\roundUp{2\pi}}\\
    \code{interval::universe()} & \itv{-\infty}{+\infty}\\
    \code{interval::zero()} & $0$\\
    \hline
  \end{tabular}

\end{table}

\noindent
The enclosures of $\pi$, $2\pi$ and $\pi/2$ are the tightest ones
(\pxref{sec:accuracy-modes}). The empty set is held as an interval whose two
bounds are NaN.
\cindex{empty!interval: representation}%

\begin{example}
cout << interval::emptyset() << ' ' << interval::pi() << endl;
@outputs^[empty] [3.141592653589793, 3.141592653589794]~
\end{example}

\section{Discussion}

This chapter described the ways an interval of GAOL comes to exist and to
change: four constructors, the function \code{textToInterval()} for the text
of an interval, the static method \code{midrad()} for a midpoint and a
radius, the straight assignment, the assignments combined with an operation,
and the constants of the class. One rule runs through the constructors,
\code{midrad()} and the assignments: an input that makes no interval, bounds in
the wrong order, a NaN bound or an infinite point, gives the empty set rather
than an object the operators do not support; \code{textToInterval()} gives it
too for a literal such as \code{"[2, 1]"} or \code{"[inf]"}, but throws for a
string that is no expression, \code{"<-inf, -inf>"} among them. GAOL v5 follows IEEE 1788-2015 and IBEX there, where GAOL~4 built
intervals of one infinite point.

\medskip
Two choices of GAOL v5 cost the programmer a few characters and remove
ambiguities. The constructors from one \code{double} and from one integer are
explicit, so that a number no longer becomes an interval without being asked
to. The constructors from strings of GAOL~4 are replaced by
\code{textToInterval()}, so that \code{interval(0)} is no longer ambiguous
and a null pointer no longer compiles into a crash. An integer
bound is taken as the integer it is, whatever the rounding direction, which
keeps $2^{53}+1$ inside the interval that encloses it. The decimal constants
remain the responsibility of the programmer: the way to an interval
containing $0.1$ is a string.

\medskip
The intervals created here are compared in the next chapter, which describes
the set relations, the certainly relations and the order of IEEE 1788-2015 on
intervals (Chapter~\ref{chap:relations}).

%% !TEX root = ../manual_gaolv5.tex
% Chapter 6.
%% Chapter 6: Interval relations. Sources: gaol.tex of f5bf180 lines 2803-3492
%% (the chapter "Interval relations" of the current manual) and doc/using.md
%% lines 631-694 (relations, sorting and containers).
\chapter{Interval relations}\label{chap:relations}\label{sec:interval-relations}

Comparing two intervals is not comparing two numbers. An interval stands for
a set of reals, and a relation between two such sets can be read in two ways:
as a relation between the sets themselves, the equality or the inclusion, or
as a relation that every point of the first set bears to every point of the
second. A program of interval analysis needs both readings. It asks whether a
box is included in another one, whether the enclosure of a function lies
certainly below a bound before it prunes a branch, or whether an interval
contains zero before it divides by it. GAOL v5 provides each reading as a
named method of the class \code{interval}, and gives the relation symbols of
C++ one fixed meaning.

\medskip
The relations of GAOL v5 follow IEEE 1788-2015 \cite{IEEE1788:2015}, which
defines the comparisons of intervals, their value for the empty set and for
the infinite bounds, and a primitive, \ieee{overlap}, from which the
comparisons can be derived. GAOL~4 had a larger family of relations and a
build option that chose the meaning of the symbols. GAOL v5 keeps two groups
of relations and no option.

\medskip
In this chapter, we first define the two groups, the set relations and the
certainly relations, and we name the relations of the standard
(Section~\ref{sec:relations-groups}). Sections~\ref{sec:relations-set}
and~\ref{sec:relations-certainly} are the reference of the set relations and
of the certainly relations. Section~\ref{sec:relations-order} covers the
order of IEEE 1788-2015 and the overlap state.
Section~\ref{sec:relational-symbols} gives the meaning of the relation
symbols, Section~\ref{sec:lexicographic-less} explains why they are no
ordering for the containers of the standard library and which comparator to
give them, and Section~\ref{sec:relations-specific} lists the relations that
test one interval: whether it contains zero, is a point, is empty or is
bounded.

\section{Two groups of relations}\label{sec:relations-groups}
\cindex{relation!set}\cindex{relation!certainly}%

Let $I$ and $J$ be two intervals. The relations of GAOL v5 fall into two
groups, according to what they compare.

\begin{Def}\label{def:relations-set}
A \dfn{set relation} compares the intervals $I$ and $J$ as sets of reals.
The equality, for instance, is
\begin{equation*}
I=J \Leftrightarrow(\forall x\in I, \exists y\in J\colon x=y)\wedge
                   (\forall y\in J, \exists x\in I\colon x=y),
\end{equation*}
so that two intervals are equal when they have the same bounds.
\end{Def}

\begin{Def}\label{def:relations-certainly}
A \dfn{certainly relation} holds between $I$ and $J$ when the relation on
the reals holds for every pair of values taken in the intervals. For
instance,
\begin{equation*}
I\leq J \Leftrightarrow(\forall x\in I, \forall y\in J\colon x\leq y),
\end{equation*}
so that $I$ is less than or equal to $J$ when its right bound is at most the
left bound of $J$.
\end{Def}

\begin{Rem}\label{rem:relations-possibly}
GAOL~4 had a third group, the \dfn{possibly relations}, true when at least
one pair of values satisfies the relation on the reals, and a certain
equality, \code{certainly\_eq()} and \code{certainly\_neq()}. They are gone
from GAOL v5, and so are the operators \code{==} and \code{!=} on
intervals.\newinvfive{} \code{certainly\_neq()} was the negation of
\code{certainly\_eq()}, which is ``possibly not equal'', so that \itv{3}{4}
was certainly not equal to \itv{3}{4}. IEEE 1788-2015 has neither. Its
equality is \ieee{equal}, \code{set\_eq()}; that no value of $I$ equals a
value of $J$ is \ieee{disjoint}, \code{set\_disjoint()}; and ``possibly
equal'' is written \code{!set\_disjoint()}.
\end{Rem}
\cindex{relation!possibly (removed)}%

\medskip
The relation symbols $<$, $\leq$, $>$ and $\geq$ are the certainly relations
in every build (\pxref{sec:relational-symbols}).

\medskip
The relations that IEEE 1788-2015 defines (Tables 10.3 and 10.4) follow its
definitions, for the empty set and for the infinite bounds
too:\newinvfive{} \ieee{equal}, \ieee{subset}, \ieee{interior},
\ieee{disjoint}, \ieee{precedes}, \ieee{strictPrecedes}, \ieee{less} and
\ieee{strictLess}. Table~\ref{tab:relations-ieee} gives the method of GAOL
v5 that each one is, with the predicates of the standard on one interval
that this chapter covers; each entry of the reference says which name it
bears. The namespace \code{gaol\_ieee1788} gives the methods the names of
the standard (Chapter~\ref{chap:ieee1788}).
\cindex{IEEE 1788-2015!relations}%

\begin{table}[htbp]
  \centering\small\renewcommand{\arraystretch}{1.1}
  \begin{tabular}{lL{7.2cm}l}
    \hline
    \multicolumn{1}{c}{IEEE 1788-2015} & \multicolumn{1}{c}{GAOL v5} &
    \multicolumn{1}{c}{Symbol}\\
    \hline\hline
    \ieee{equal}$(I, J)$ & \code{I.set\_eq(J)} & \\
    \ieee{subset}$(I, J)$ & \code{I.set\_leq(J)}, \code{J.set\_contains(I)},
      \code{J.set\_geq(I)} & \\
    \ieee{interior}$(I, J)$ & \code{I.set\_le(J)},
      \code{J.set\_strictly\_contains(I)}, \code{J.set\_ge(I)} & \\
    \ieee{disjoint}$(I, J)$ & \code{I.set\_disjoint(J)} & \\
    \ieee{precedes}$(I, J)$ & \code{I.certainly\_leq(J)} & \code{I <= J}\\
    \ieee{strictPrecedes}$(I, J)$ & \code{I.certainly\_le(J)} & \code{I < J}\\
    \ieee{less}$(I, J)$ & \code{I.less(J)} & \\
    \ieee{strictLess}$(I, J)$ & \code{I.strictly\_less(J)} & \\
    \ieee{overlap}$(I, J)$ & \code{I.overlap(J)} & \\
    \hline
    \ieee{isMember}$(d, I)$ & \code{I.set\_contains(d)}, for a finite
      \code{double} \var{d} & \\
    \ieee{isSingleton}$(I)$ & \code{I.is\_a\_double()} & \\
    \ieee{isEmpty}$(I)$ & \code{I.is\_empty()} & \\
    \ieee{isEntire}$(I)$ & \code{I.is\_entire()} & \\
    \ieee{isCommonInterval}$(I)$ & \code{I.is\_common\_interval()} & \\
    \hline
  \end{tabular}
  \caption[The relations of IEEE 1788-2015 and the methods of GAOL v5]{The relations of IEEE 1788-2015 and the methods of GAOL v5 that
    compute them, for two intervals $I$ and $J$ and a double $d$}
  \label{tab:relations-ieee}
\end{table}

\section{Set relations}\label{sec:relations-set}

The set relations compare the intervals as sets of reals: the inclusion, the
inclusion in the interior, the disjointness and the equality. They are
\ieee{subset}, \ieee{interior}, \ieee{disjoint} and \ieee{equal} of
IEEE 1788-2015. Each inclusion has two methods: one where \code{self}, the
interval on which the method is called, is the larger set,
\code{set\_contains()} and \code{set\_strictly\_contains()}, and one where it
is the smaller, \code{set\_leq()} and \code{set\_le()}. \code{set\_geq()} and
\code{set\_ge()} are \code{set\_contains()} and
\code{set\_strictly\_contains()} under the names of an order.

\defmethodx{interval}{bool}{set\_contains}{(const interval\& \var{I}) const}
\begin{defmethod}{interval}{bool}{set\_contains}{(double \var{d}) const}
  \begin{operation}
    \op{$I\subseteq\code{self}$}
    \op{$d\in\code{self}$}
  \end{operation}
  Returns \code{true} if \var{I} (resp.\ $\{\var{d}\}$) is included in
  \code{self}: \ieee{subset}(\var{I}, \code{self}) of IEEE 1788-2015, and
  \ieee{isMember}(\var{d}, \code{self}) for a finite \var{d}. \var{d} shall
  not be a NaN. An integer of any integer type but \code{bool} is compared as
  the integer it is, here and in \code{set\_strictly\_contains()}
  (\pxref{sec:integers}): \itv{2^{53}}{2^{53}} does not contain $2^{53}+1$,
  which its conversion to a double made $2^{53}$ in the rounding to nearest,
  and \itv{2^{53}}{2^{53}+2} contains it strictly.\newinvfive

\begin{example}
interval x(-12,34), y(-12,5);

cout << boolalpha << x.set_contains(y) << ' '
     << x.set_contains(interval::emptyset()) << ' '
     << y.set_contains(x) << ' '
     << interval::emptyset().set_contains(x) << ' '
     << interval::emptyset().set_contains(interval::emptyset())
     << endl;
@outputs^true true false false true~
\end{example}
\end{defmethod}

\defmethodx{interval}{bool}{set\_strictly\_contains}{(const interval\& \var{I}) const}
\begin{defmethod}{interval}{bool}{set\_strictly\_contains}{(double \var{d}) const}
  \begin{operation}
    \op{$I$ interior to \code{self}}
    \op{$d$ in the interior of \code{self}}
  \end{operation}
  Returns \code{true} if \var{I} is interior to \code{self}:
  \ieee{interior}(\var{I}, \code{self}) of IEEE 1788-2015. Each bound of
  \code{self} is then beyond that of \var{I}, an infinite bound being beyond
  the same infinite bound, and the empty set is interior to any interval,
  itself included.\newinvfive{} The second form returns \code{true} if
  \var{d} is strictly between the bounds of \code{self}.

\begin{example}
interval x(-10,12), y(-10, 11), z, t, u(10.5);

cout << boolalpha
     << x.set_strictly_contains(y) << ' '
     << z.set_strictly_contains(t) << ' '
     << x.set_strictly_contains(u) << ' '
     << interval::emptyset().set_strictly_contains(
                                 interval::emptyset()) << ' '
     << u.set_strictly_contains(interval::emptyset())
     << endl;
@outputs^false true true true true~
\end{example}
\end{defmethod}

\begin{defmethod}{interval}{bool}{set\_disjoint}{(const interval\& \var{I}) const}
  \begin{operation}
    \op{$\code{self}\cap I=\emptyset$}
  \end{operation}
  Returns \code{true} if the intersection of \code{self} and \var{I} is
  empty: \ieee{disjoint} of IEEE 1788-2015.

\begin{example}
interval a(2,4), b(6,8);
cout << boolalpha << a.set_disjoint(b) << ' '
     << interval::emptyset().set_disjoint(interval::emptyset());
@outputs^true true~
\end{example}
\end{defmethod}

\begin{defmethod}{interval}{bool}{set\_eq }{(const interval\& \var{I}) const}
  \begin{operation}
    \op{$\forall x\in \code{self} \exists y\in I\colon x=y \wedge
         \forall y\in I, \exists x\in \code{self}\colon y=x$}
  \end{operation}
  Returns \code{true} if the intervals \code{self} and \var{I} are equal
  when considered as sets of reals: \ieee{equal} of IEEE 1788-2015.
\begin{example}
cout << boolalpha << interval(4,8).set_eq(interval(4,8)) << ' '
     << interval::emptyset().set_eq(interval::emptyset()) << endl;
@outputs^true true~
\end{example}
\end{defmethod}

\begin{defmethod}{interval}{bool}{set\_neq}{(const interval\& \var{I}) const}
    \begin{operation}
      \op{$\neg(\code{self}=I)$}
    \end{operation}
    Returns \code{true} if \code{self} and \var{I} are not equal when
    considered as sets of reals.
    \begin{example}
cout << boolalpha << interval::universe().set_neq(interval::emptyset());
@outputs^true~
    \end{example}
\end{defmethod}

\begin{defmethod}{interval}{bool}{set\_le}{(const interval\& \var{I}) const}
  \begin{operation}
    \op{\code{self} interior to $I$}
  \end{operation}
  Returns \code{true} if \code{self} is interior to \var{I}:
  \code{\var{I}.set\_strictly\_contains(self)}, \ieee{interior}(\code{self},
  \var{I}) of IEEE 1788-2015.\newinvfive
\begin{example}
interval a(2,GAOL_INFINITY), b(1,GAOL_INFINITY);
cout << boolalpha << interval(-4.5,3).set_le(interval(-10,10)) << ' '
     << interval::emptyset().set_le(interval(5,6)) << ' '
     << interval::emptyset().set_le(interval::emptyset()) << ' '
     << interval::universe().set_le(interval::universe()) << ' '
     << a.set_le(b);
@outputs^true true true true true~
\end{example}
\end{defmethod}

\begin{defmethod}{interval}{bool}{set\_leq}{(const interval\& \var{I}) const}
  \begin{operation}
    \op{$\code{self}\subseteq I$}
  \end{operation}
  Returns \code{true} if the real set defined by \code{self} is included in
  \var{I}: \ieee{subset}(\code{self}, \var{I}) of IEEE 1788-2015.
\begin{example}
cout << boolalpha << interval(4.5,6).set_leq(interval(4.5,6)) << ' '
     << interval(3.5,9).set_leq(interval(2,6)) << ' '
     << interval::emptyset().set_leq(interval::emptyset());
@outputs^true false true~
\end{example}
\end{defmethod}

\begin{defmethod}{interval}{bool}{set\_ge}{(const interval\& \var{I}) const}
  \begin{operation}
    \op{$I$ interior to \code{self}}
  \end{operation}
  Returns \code{true} if \var{I} is interior to \code{self}:
  \code{set\_strictly\_contains(\var{I})}.
\begin{example}
cout << boolalpha << interval(-10,10).set_ge(interval(-4.5,3)) << ' '
     << interval(1,2).set_ge(interval(1,2));
@outputs^true false~
\end{example}
\end{defmethod}

\begin{defmethod}{interval}{bool}{set\_geq}{(const interval\& \var{I}) const}
  \begin{operation}
    \op{$I\subseteq\code{self}$}
  \end{operation}
  Returns \code{true} if the real set defined by \code{self} contains
  \var{I}: \code{set\_contains(\var{I})}.
\end{defmethod}

\section{Certainly relations}\label{sec:relations-certainly}

A certainly relation holds when every pair of points, one in each interval,
satisfies the relation on the reals. The empty set has no point to
contradict it: the four relations of this section are true when either
interval is empty, as \ieee{precedes} and \ieee{strictPrecedes} of
IEEE 1788-2015 are.\newinvfive{} The four tests of sign that follow them,
\code{certainly\_positive()} and its companions, are the certainly relations
with zero, written as the inclusion of \code{self} in a half line; they are
true of the empty set too.

\begin{defmethod}{interval}{bool}{certainly\_le}{(const interval\& \var{I}) const}
  \begin{operation}
    \op{$\forall x\in \code{self}, \forall y \in \var{I}\colon x<y$}
  \end{operation}
  Returns \code{true} if \code{self} is certainly strictly less than
  \var{I}: \ieee{strictPrecedes} of IEEE 1788-2015, true when either
  interval is empty.\newinvfive
  \begin{example}
cout << boolalpha << interval(4,5).certainly_le(interval(6,9)) << ' '
     << interval(4,5).certainly_le(interval(5,9)) << ' '
     << interval::emptyset().certainly_le(interval(4,6)) << ' '
     << interval(4,6).certainly_le(interval::emptyset());
@outputs^true false true true~
  \end{example}
\end{defmethod}

\begin{defmethod}{interval}{bool}{certainly\_leq}{(const interval\& \var{I}) const}
  \begin{operation}
    \op{$\forall x\in \code{self}, \forall y \in \var{I}\colon x\leq y$}
  \end{operation}
  Returns \code{true} if \code{self} is certainly less than or equal to
  \var{I}: \ieee{precedes} of IEEE 1788-2015, true when either interval is
  empty.\newinvfive
  \begin{example}
cout << boolalpha << interval(4,5).certainly_leq(interval(6,9)) << ' '
     << interval(4,5).certainly_leq(interval(5,9)) << ' '
     << interval(5,9).certainly_leq(interval(4,5)) << ' '
     << interval(4,8).certainly_leq(interval(5,9)) << ' '
     << interval::emptyset().certainly_leq(interval(4,6)) << ' '
     << interval(4,6).certainly_leq(interval::emptyset());
@outputs^true true false false true true~
\end{example}
\end{defmethod}

\begin{defmethod}{interval}{bool}{certainly\_ge}{(const interval\& \var{I}) const}
  \begin{operation}
    \op{$\forall x\in \code{self}, \forall y \in \var{I}\colon x>y$}
  \end{operation}
  Returns \code{true} if \code{self} is certainly strictly greater than
  \var{I}, or if either interval is empty.
  \begin{example}
cout << boolalpha << interval(8,10).certainly_ge(interval(4,8)) << ' '
     << interval(8,10).certainly_ge(interval(4,7)) << ' '
     << interval::emptyset().certainly_ge(interval::emptyset());
@outputs^false true true~
  \end{example}
\end{defmethod}

\begin{defmethod}{interval}{bool}{certainly\_geq}{(const interval\& \var{I}) const}
  \begin{operation}
    \op{$\forall x\in \code{self}, \forall y \in \var{I}\colon x\geq y$}
  \end{operation}
  Returns \code{true} if \code{self} is certainly greater than or equal to
  \var{I}, or if either interval is empty.
  \begin{example}
cout << boolalpha << interval(8,10).certainly_geq(interval(4,8)) << ' '
     << interval::emptyset().certainly_geq(interval::emptyset());
@outputs^true true~
  \end{example}
\end{defmethod}

\begin{defmethod}{interval}{bool}{certainly\_positive}{(void) const}
  \begin{operation}
    \op{$\code{self}\subseteq\itv{0}{+\infty}$}
  \end{operation}
  Returns \code{true} if the left bound of \code{self} is greater than or
  equal to zero, or if \code{self} is empty.
\begin{example}
cout << boolalpha << interval::emptyset().certainly_positive() << ' '
     << interval(4,5).certainly_positive() << ' '
     << interval(-0.0,6).certainly_positive() << ' '
     << interval(-6,0).certainly_positive();
@outputs^true true true false~
\end{example}
\end{defmethod}

\begin{defmethod}{interval}{bool}{certainly\_strictly\_positive}{(void) const}
  \begin{operation}
    \op{$\code{self}\subseteq\left]0,+\infty\right]$}
  \end{operation}
  Returns \code{true} if the left bound of \code{self} is strictly greater
  than zero, or if \code{self} is empty.
\begin{example}
cout << boolalpha << interval::emptyset().certainly_strictly_positive() << ' '
     << interval(4,5).certainly_strictly_positive() << ' '
     << interval(-0.0,6).certainly_strictly_positive() << ' '
     << interval(-6,0).certainly_strictly_positive();
@outputs^true true false false~
\end{example}
\end{defmethod}

\begin{defmethod}{interval}{bool}{certainly\_negative}{(void) const}
  \begin{operation}
    \op{$\code{self}\subseteq\itv{-\infty}{0}$}
  \end{operation}
  Returns \code{true} if the right bound of \code{self} is less than or
  equal to zero, or if \code{self} is empty.
\begin{example}
cout << boolalpha << interval::emptyset().certainly_negative() << ' '
     << interval(4,5).certainly_negative() << ' '
     << interval(-6,0).certainly_negative() << ' '
     << interval(-6,-5).certainly_negative();
@outputs^true false true true~
\end{example}
\end{defmethod}

\begin{defmethod}{interval}{bool}{certainly\_strictly\_negative}{(void) const}
  \begin{operation}
    \op{$\code{self}\subseteq\left[-\infty,0\right[$}
  \end{operation}
  Returns \code{true} if the right bound of \code{self} is strictly less
  than zero, or if \code{self} is empty.
\begin{example}
cout << boolalpha << interval::emptyset().certainly_strictly_negative() << ' '
     << interval(4,5).certainly_strictly_negative() << ' '
     << interval(-6,0).certainly_strictly_negative() << ' '
     << interval(-6,-5).certainly_strictly_negative();
@outputs^true false false true~
\end{example}
\end{defmethod}

\section{The order of IEEE 1788-2015 and the overlap state}
\label{sec:relations-order}
\cindex{IEEE 1788-2015!order}\cindex{overlap}\cindex{Allen's relations}%

Besides the set relations and the certainly relations, IEEE 1788-2015
defines an order on the intervals, \ieee{less} and \ieee{strictLess}
(Table~10.3), which compares the bounds on the same side: $I$ is less than
$J$ when the left bound of $I$ is at most the left bound of $J$ and the right
bound of $I$ is at most the right bound of $J$. The standard also defines a
primitive, \ieee{overlap} (10.6.4), which tells in which of sixteen states
two intervals are, and from which the comparisons can be derived. Three of
the states are for the empty set. The thirteen others are the relations of
Allen between two nonempty intervals $a = [\leftBound{a}, \rightBound{a}]$
and $b = [\leftBound{b}, \rightBound{b}]$, told by their bounds, an infinite
bound being equal to the same infinite bound and $-0$ to $+0$
(Table~\ref{tab:relations-overlap}). The three methods of this section
compute them.

\begin{table}[htbp]
  \centering\renewcommand{\arraystretch}{1.1}
  \begin{tabular}{ll@{\qquad}ll}
    \hline
    \multicolumn{1}{c}{State} & \multicolumn{1}{c}{Bounds} &
    \multicolumn{1}{c}{State} & \multicolumn{1}{c}{Bounds}\\
    \hline\hline
    \code{before} & $\rightBound{a}<\leftBound{b}$ &
    \code{after} & $\rightBound{b}<\leftBound{a}$\\
    \code{meets} &
      $\leftBound{a}<\rightBound{a}=\leftBound{b}<\rightBound{b}$ &
    \code{metBy} &
      $\leftBound{b}<\rightBound{b}=\leftBound{a}<\rightBound{a}$\\
    \code{overlaps} &
      $\leftBound{a}<\leftBound{b}<\rightBound{a}<\rightBound{b}$ &
    \code{overlappedBy} &
      $\leftBound{b}<\leftBound{a}<\rightBound{b}<\rightBound{a}$\\
    \code{starts} &
      $\leftBound{a}=\leftBound{b},\ \rightBound{a}<\rightBound{b}$ &
    \code{startedBy} &
      $\leftBound{a}=\leftBound{b},\ \rightBound{b}<\rightBound{a}$\\
    \code{finishes} &
      $\leftBound{b}<\leftBound{a},\ \rightBound{a}=\rightBound{b}$ &
    \code{finishedBy} &
      $\leftBound{a}<\leftBound{b},\ \rightBound{a}=\rightBound{b}$\\
    \code{containedBy} &
      $\leftBound{b}<\leftBound{a},\ \rightBound{a}<\rightBound{b}$ &
    \code{contains} &
      $\leftBound{a}<\leftBound{b},\ \rightBound{b}<\rightBound{a}$\\
    \code{equals} &
      \multicolumn{3}{l}{$\leftBound{a}=\leftBound{b},\
                          \rightBound{a}=\rightBound{b}$}\\
    \hline
  \end{tabular}
  \caption[The 13 states of Allen returned by \code{overlap()}]{The 13 states of Allen for two nonempty intervals $a$ and $b$,
    as \code{a.overlap(b)} returns them (Table~10.7 of IEEE 1788-2015); each
    state of the left column is the inverse of the one on its right}
  \label{tab:relations-overlap}
\end{table}

\begin{defmethod}{interval}{bool}{less}{(const interval\& \var{I}) const}
  \begin{operation}
    \op{$\leftBound{\code{self}}\leq\leftBound{I}\wedge
         \rightBound{\code{self}}\leq\rightBound{I}$}
  \end{operation}
  Returns \code{true} if each bound of \code{self} is at most the bound of
  \var{I} on the same side: \ieee{less} of IEEE 1788-2015 (Table 10.3). Two
  empty sets are in the relation, and an empty set and a nonempty interval
  are not (Table 10.4).\newinvfive
\begin{example}
interval e = interval::emptyset();
cout << boolalpha << interval(1,2).less(interval(1,3)) << ' '
     << interval(1,4).less(interval(2,3)) << ' '
     << e.less(e) << ' ' << e.less(interval(1,2));
@outputs^true false true false~
\end{example}
\end{defmethod}

\begin{defmethod}{interval}{bool}{strictly\_less}{(const interval\& \var{I}) const}
  \begin{operation}
    \op{$\leftBound{\code{self}}<\leftBound{I}\wedge
         \rightBound{\code{self}}<\rightBound{I}$}
  \end{operation}
  Returns \code{true} if each bound of \code{self} is below the bound of
  \var{I} on the same side, an infinite bound counting as below itself:
  \ieee{strictLess} of IEEE 1788-2015. The empty sets are as for
  \code{less()}.\newinvfive
\begin{example}
interval a(-GAOL_INFINITY,1), b(-GAOL_INFINITY,2);
cout << boolalpha << interval(1,2).strictly_less(interval(1,3)) << ' '
     << a.strictly_less(b) << ' '
     << interval(0,1).strictly_less(interval(2,3));
@outputs^false true true~
\end{example}
\end{defmethod}

\begin{defmethod}{interval}{overlapping\_state}{overlap}{(const interval\& \var{I}) const}
  Returns the state of Table~10.7 of IEEE 1788-2015 that \code{self}, $a =
  [\leftBound{a}, \rightBound{a}]$, and \var{I}, $b = [\leftBound{b},
  \rightBound{b}]$, are in: \ieee{overlap} (10.6.4), a primitive from which
  the comparisons can be derived.\newinvfive\ The result is one of the 16
  values of the type \code{overlapping\_state} (in \code{gaol} and in
  \code{gaol\_ieee1788}), which has the names of the table, in its order,
  from \code{bothEmpty}, 0, to \code{after}, 15. The state number $s$ is the
  bit $s$ of a 16-bit word counted from the left, \code{bothEmpty} being the
  leftmost bit, as in the note of the standard, where \code{disjoint} is the
  mask \code{0xF001}. A state is written to a stream by its name in the
  table, as \code{meets}. \code{bothEmpty}, \code{firstEmpty} and
  \code{secondEmpty} are for an empty $a$ and $b$, an empty $a$ alone and an
  empty $b$ alone. The 13 other states are Allen's, for two nonempty
  intervals, told by their bounds as Table~\ref{tab:relations-overlap} says,
  an infinite bound being equal to the same infinite bound and $-0$ to $+0$.
  A point interval on a bound of the other one is \code{starts},
  \code{finishes}, \code{finishedBy} or \code{startedBy}, or \code{equals}
  for the same point, never \code{meets} nor \code{metBy}, as the standard
  asks. The function compares the bounds as the integers of their bits: it
  raises no floating-point exception, even for the \code{NaN} bounds of an
  empty interval, and the modes that flush the subnormal numbers to zero do
  not change it.
  \begin{example}
interval a(1,3), b(2,4), c(1,2), e = interval::emptyset();
cout << a.overlap(b) << ' ' << c.overlap(interval(2,3)) << ' '
     << interval(2).overlap(c) << ' ' << e.overlap(c) << ' '
     << static_cast<int>(c.overlap(a));
@outputs^overlaps meets finishes firstEmpty 6~
  \end{example}
\end{defmethod}

\section{Relational symbols}\label{sec:relational-symbols}
\cindex{relation!symbols}%

The relation symbols \code{<}, \code{<=}, \code{>} and \code{>=} are the
certainly relations, whichever build configured GAOL. The option of GAOL~4
that chose between the set, certainly and possibly relations is gone, and
\cmd{configure} and meson refuse it (\pxref{sec:refused}).\newinvfive{}
The operators \code{==} and \code{!=} are not defined on intervals: a program
compares intervals with \code{set\_eq()}, \code{set\_disjoint()} or
\code{less()}.

\medskip
An interval and a \code{double} \var{d} compare as the interval and
\code{interval(\var{d})}, which is empty for an infinite or NaN \var{d}:
\code{x < 0.0} is \code{x < interval(0.0)}, and there is still no \code{==}
nor \code{!=}. These overloads are what the implicit conversion of a
\code{double} to an interval gave before the constructor from one
\code{double} became explicit (\pxref{sec:constructors}).\newinvfive

\medskip
An interval and an integer \var{n} of any integer type but \code{bool}
compare as the interval and the integer it is, where it is no double
(\pxref{sec:integers}). A double is below $n$ if and only if it is at most
the largest double below $n$, so that
\code{interval(9007199254740992.0) < 9007199254740993LL}, that is
$\itv{2^{53}}{2^{53}} < 2^{53}+1$, is \code{true}, where the conversion of
$2^{53}+1$ to the double $2^{53}$, in the rounding to nearest, made it
\code{false}. An integer that is a double compares as that
double.\newinvfive

\deffunx{bool}{operator<}{(const interval\& \var{I}, double \var{d})}
\deffunx{bool}{operator<}{(double \var{d}, const interval\& \var{I})}
\begin{deffun}{bool}{operator<}{(const interval\& \var{I1}, const interval\& \var{I2})}
  Returns \code{\var{I1}.certainly\_le(\var{I2})}: \ieee{strictPrecedes} of
  IEEE 1788-2015.
\end{deffun}

\deffunx{bool}{operator<=}{(const interval\& \var{I}, double \var{d})}
\deffunx{bool}{operator<=}{(double \var{d}, const interval\& \var{I})}
\begin{deffun}{bool}{operator<=}{(const interval\& \var{I1}, const interval\& \var{I2})}
  Returns \code{\var{I1}.certainly\_leq(\var{I2})}: \ieee{precedes} of
  IEEE 1788-2015.
\end{deffun}

\deffunx{bool}{operator>}{(const interval\& \var{I}, double \var{d})}
\deffunx{bool}{operator>}{(double \var{d}, const interval\& \var{I})}
\begin{deffun}{bool}{operator>}{(const interval\& \var{I1}, const interval\& \var{I2})}
  Returns \code{\var{I1}.certainly\_ge(\var{I2})}.
\end{deffun}

\deffunx{bool}{operator>=}{(const interval\& \var{I}, double \var{d})}
\deffunx{bool}{operator>=}{(double \var{d}, const interval\& \var{I})}
\begin{deffun}{bool}{operator>=}{(const interval\& \var{I1}, const interval\& \var{I2})}
  Returns \code{\var{I1}.certainly\_geq(\var{I2})}.
\begin{example}
cout << boolalpha << (interval(1,2) < interval(3,4)) << ' '
     << (interval(1,2) <= interval(2,4)) << ' '
     << (interval(1,3) <= interval(2,4)) << ' '
     << (interval(3,4) > interval(1,2));
@outputs^true true false true~
\end{example}
\end{deffun}

\section{Intervals in containers and algorithms}\label{sec:lexicographic-less}
\cindex{lexicographic\_less@\code{lexicographic\_less}}%
\cindex{sorting intervals}%
\cindex{std::set@\code{std::set}}%

The relation symbols do not tell which of two intervals is the greater:
\code{interval(1,3) < interval(2,4)} and \code{interval(1,3) >= interval(2,4)}
are both \code{false}, and \code{false} is not the opposite. The empty set,
which has no point to contradict a certainly relation, is less than every
interval, and greater, itself included: \code{e < e} is \code{true}. A
program ported from a library where \code{x <= y} is the inclusion compiles
and changes its meaning (\pxref{sec:other-libraries}).

\begin{Rem}\label{rem:relations-ordering}
None of the relation symbols is a strict weak ordering, so intervals are not
given to the algorithms and the containers of the standard library that
compare with \code{<} (\code{std::sort}, \code{std::set}, \code{std::map},
\code{std::priority\_queue}, \code{std::max}, \code{std::min},
\code{std::clamp}\ldots). Two overlapping intervals are equivalent for
\code{<}, a relation that is not transitive, and the empty set is less than
itself. The behaviour of these algorithms is then undefined, and differs with
the implementation of the library. A \code{std::set<interval>} that holds
\itv{1}{2} drops \itv{1}{3}: two of the three intervals \itv{1}{2},
\itv{1}{3} and \itv{5}{6} stay. \code{std::sort} of a vector that holds an
empty interval can read past the end of the vector: AddressSanitizer stops
the program with libstdc++ for 17 intervals or more, at some positions of
the empty one. \code{std::max(\itv{1}{3}, \itv{2}{4})} is \itv{1}{3}, its
first argument, where \code{gaol::max} gives \itv{2}{4}.
\end{Rem}

\medskip
\code{gaol::max(x, y)} and \code{gaol::min(x, y)} are the maximum and the
minimum of the standard: the set of the greater (or the lesser) of a point of
\code{x} and a point of \code{y}, and not the greater (or the lesser) of the
two arguments.

\medskip
\code{gaol::lexicographic\_less} is the comparator to give the containers
and the algorithms of the standard library: a strict total order that puts
the empty set first and the other intervals by their left bound, then by
their right bound, $-0$ and $+0$ being the same bound.\newinvfive{} Two
intervals with the same bounds, or two empty sets, are not less than each
other, as a container needs. The comparator is also in \code{gaol\_core}; a
program that opens \code{gaol\_ieee1788} writes its name in full.
Specializing \code{std::less} for intervals is left to the program, so the
comparator is given explicitly.

\begin{Ex}\label{ex:relations-containers}
The program below puts four intervals, the empty set among them, in a
\code{std::set} and in a sorted \code{std::vector} with the comparator
\code{lexicographic\_less}.
\begin{example}
#include <algorithm>
#include <iostream>
#include <set>
#include <vector>
#include <gaol/gaol.h>
using namespace gaol;

int main(void)
{
  const interval boxes[] = { interval(1,2), interval(0,5),
                             interval::emptyset(), interval(1,3) };
  std::set<interval, lexicographic_less> s(boxes, boxes + 4);
  std::vector<interval> v(boxes, boxes + 4);
  std::sort(v.begin(), v.end(), lexicographic_less());
  for (const interval& x : v) std::cout << x << ' ';
  std::cout << s.size() << std::endl;
  gaol::cleanup();
  return 0;
}
@outputs^[empty] [0, 5] [1, 2] [1, 3] 4~
\end{example}
The vector comes out with the empty set first, then \itv{0}{5}, \itv{1}{2}
and \itv{1}{3}, and the set holds the four intervals: none is dropped.
\end{Ex}

\medskip
An algorithm that needs to know which of two boxes to process first, as a
best-first search, orders them by a number of its own, the lower bound of an
enclosure for instance, as \file{examples/06\_global\_optimization.cpp} does.

\section{Interval-specific relations}\label{sec:relations-specific}
\cindex{relation!interval-specific}%

The relations of this section test one interval: its position with respect
to zero, its shape (a point, an integer, a canonical interval, a symmetric
one) and its place among the sets (the empty set, the whole real line, a
bounded interval). Four of them are predicates of IEEE 1788-2015:
\ieee{isSingleton}, \ieee{isEmpty}, \ieee{isEntire} and
\ieee{isCommonInterval}.

\begin{defmethod}{interval}{bool}{straddles\_zero}{(void) const}
  \begin{operation}
    \op{$0 \in \code{self}$}
  \end{operation}
  Returns \code{true} if \code{self} contains zero. The call
  \code{I.straddles\_zero()} is equivalent to \code{I.set\_contains(0.0)}.
\begin{example}
interval x(0,4), y, z(-12,-5);

cout << boolalpha << x.straddles_zero() << ' '
     << y.straddles_zero() << ' '
     << z.straddles_zero() << endl;
@outputs^true true false~
\end{example}
\end{defmethod}

\begin{defmethod}{interval}{bool}{strictly\_straddles\_zero}{(void) const}
  \begin{operation}
  \op{$\{0\}$ interior to \code{self}}
  \end{operation}
  Returns \code{true} if zero is included in the interior of \code{self}.
  The call \code{I.strictly\_straddles\_zero()} is equivalent to
  \code{I.set\_strictly\_contains(0.0)}.
\begin{example}
interval x(0,4), y, z(-12,-5);

cout << boolalpha << x.strictly_straddles_zero() << ' '
     << y.strictly_straddles_zero() << ' '
     << z.strictly_straddles_zero() << endl;
@outputs^false true false~
\end{example}
\end{defmethod}

\begin{defmethod}{interval}{bool}{is\_a\_double}{(void) const}
  Returns \code{true} whenever the left and right bounds of the interval
  are equal: \ieee{isSingleton} of IEEE 1788-2015.
\begin{example}
interval x(-12.5), y;

cout << boolalpha << x.is_a_double() << ' '
     << y.is_a_double() << ' '
     << interval::emptyset().is_a_double() << endl;
@outputs^true false false~
\end{example}
\end{defmethod}

\begin{defmethod}{interval}{bool}{is\_an\_int}{(void) const}
  Returns \code{true} whenever the left and right bounds of the interval
  are equal and castable into an integer (type \code{int}).
\begin{example}
interval x(-12.0), y;

cout << boolalpha << x.is_an_int() << ' ' << y.is_an_int() << ' '
     << interval::emptyset().is_an_int() << ' '
     << interval(3e10).is_an_int() << endl;
@outputs^true false false false~
\end{example}
\end{defmethod}

\begin{defmethod}{interval}{bool}{is\_canonical}{(void) const}
  Returns \code{true} if \code{self} is nonempty and contains at most two
  floating-point numbers.
\begin{example}
interval x(0.0), y(-14,9);

cout << boolalpha << x.is_canonical() << ' '
     << y.is_canonical() << ' '
     << interval::emptyset().is_canonical() << ' '
     << interval::pi().is_canonical();
@outputs^true false false true~
\end{example}
\end{defmethod}

\begin{defmethod}{interval}{bool}{is\_empty}{(void) const}
  \begin{operation}
    \op{$\code{self} = \emptyset$}
  \end{operation}
  Returns \code{true} if \code{self} is an empty interval: \ieee{isEmpty} of
  IEEE 1788-2015. It raises no floating-point exception, the
  invalid-operation one included, though the bounds of the empty set are NaN
  (\pxref{sec:fp-exceptions}).\newinvfive
  \begin{example}
cout << boolalpha << interval(4,5).is_empty() << ' '
     << interval(5,4).is_empty() << ' '
     << interval::emptyset().is_empty();
@outputs^false true true~
  \end{example}
\end{defmethod}

\begin{defmethod}{interval}{bool}{is\_entire}{(void) const}
  \begin{operation}
    \op{$\code{self} = \itv{-\infty}{+\infty}$}
  \end{operation}
  Returns \code{true} if \code{self} is the whole real line: \ieee{isEntire}
  of IEEE 1788-2015.\newinvfive
  \begin{example}
cout << boolalpha << interval::universe().is_entire() << ' '
     << interval::positive().is_entire() << ' '
     << interval::emptyset().is_entire();
@outputs^true false false~
  \end{example}
\end{defmethod}

\begin{defmethod}{interval}{bool}{is\_common\_interval}{(void) const}
  Returns \code{true} if \code{self} is nonempty and bounded:
  \ieee{isCommonInterval} of IEEE 1788-2015.\newinvfive
  \begin{example}
cout << boolalpha << interval(1,2).is_common_interval() << ' '
     << interval::positive().is_common_interval() << ' '
     << interval::emptyset().is_common_interval();
@outputs^true false false~
  \end{example}
\end{defmethod}

\begin{defmethod}{interval}{bool}{is\_zero}{(void) const}
  \begin{operation}
    \op{$\code{self} = \itv{0}{0}$}
  \end{operation}
  Returns \code{true} if \code{self} is equal to the interval containing
  only $0$.
  \begin{example}
cout << boolalpha << interval::zero().is_zero() << ' '
     << interval(0.0,0.0).is_zero() << ' '
     << interval(-0.0,+0.0).is_zero() << ' '
     << interval(0,5).is_zero() << ' '
     << interval::emptyset().is_zero();
@outputs^true true true false false~
  \end{example}
\end{defmethod}

\begin{defmethod}{interval}{bool}{is\_symmetric}{(void) const}
  Returns \code{true} if the left bound of \code{self} is the opposite of
  the right bound. An empty interval is not symmetric.

\begin{example}
cout << boolalpha << interval(-5,5).is_symmetric() << ' '
     << interval::emptyset().is_symmetric() << endl;
@outputs^true false~
\end{example}
\end{defmethod}

\begin{defmethod}{interval}{bool}{is\_finite}{(void) const}
  Returns \code{true} if neither bound is infinite. The empty set, whose
  bounds are NaN, is finite; \code{is\_common\_interval()} is not true of
  it.

\begin{example}
cout << boolalpha << interval(4,GAOL_INFINITY).is_finite() << ' '
     << interval::emptyset().is_finite() << ' '
     << interval(5,80).is_finite();
@outputs^false true true~
\end{example}
\end{defmethod}

\section{Discussion}

This chapter has fixed the vocabulary of the comparisons of GAOL v5. Two
groups remain of the three of GAOL~4: the set relations, which compare the
intervals as sets of reals, and the certainly relations, which hold when
every pair of points is in the relation. The possibly relations,
\code{certainly\_eq()}, \code{certainly\_neq()} and the operators \code{==}
and \code{!=} are gone, since a certain inequality that was the negation of
a certain equality made \itv{3}{4} certainly not equal to itself. Every
relation that IEEE 1788-2015 defines follows the standard, for the empty set
and for the infinite bounds too, and the overlap state tells in which of the
sixteen states of the standard two intervals are, from which the comparisons
can be derived.

\medskip
The relation symbols have one meaning in every build, the certainly
relations, and compare an interval with a \code{double} or an integer as
with the interval of that number, an integer that is no double being compared
as the integer it is. The price of a fixed meaning that follows the
standard is that the symbols order nothing: a program that sorts intervals
or keeps them in a \code{std::set} gives the comparator
\code{gaol::lexicographic\_less}, and an algorithm that chooses the next box
to process orders the boxes by a number of its own. Specializing
\code{std::less} for intervals is left to the program.

\medskip
The relations tell what two intervals are to each other. The next chapter
turns to the operations that build new intervals from them: the functional
arithmetic, the relational arithmetic and the cancellative operations
(Chapter~\ref{chap:arithmetic}).

%% !TEX root = ../manual_gaolv5.tex
% Chapter 7.
%% Chapter 7: Interval arithmetic: the operations. Sources: gaol.tex of f5bf180
%% lines 3493-4445 (chapter "Interval Arithmetic": the introduction and the
%% domains, functional arithmetic, relational arithmetic, cancellative addition
%% and subtraction), doc/using.md lines 695-842 (tools of interval algorithms,
%% the reverse functions of pow, max, min, sign and floor).
\chapter{Interval arithmetic: the operations}\label{chap:arithmetic}

An interval program is written with the operators and the functions of C++,
which GAOL overloads for the type \code{interval}: a sum, a product, a square
root or a cosine of intervals is an interval that contains every value the
operation takes on the real numbers of its arguments. This chapter is the
reference of these operations. The containment property they satisfy, the
roundings that give their bounds and the mathematical libraries that compute
the elementary functions are the subject of Chapter~\ref{chap:concepts}. Here
we state, for each operation, what it computes, on which domain, which
operation of IEEE 1788-2015 \cite{IEEE1788:2015} it is, and what differs from
GAOL~4.

\medskip
The operations are of two kinds. A \emph{functional} operation extends a
function of real numbers to intervals: the square root of an interval is the
hull of the square roots of its points. A \emph{relational} operation inverts
a relation: the relational square root of $J$ with respect to $I$ is the hull
of the $x$ of $I$ whose square is in $J$, the negative ones included. The
relational operations are the reverse functions of IEEE 1788-2015, and the
tools with which the algorithms of constraint programming contract their
domains.

\medskip
The first section recalls the containment principle, the two kinds of
operations and the rule of the domains. The second section is the reference
of the functional arithmetic: the basic operations, the powers and roots, the
exponentials and logarithms, the trigonometric functions, the hyperbolic
functions and the other functions of the standard. The third section is the
reference of the relational arithmetic, and gives its correspondence with the
reverse functions of the standard. The fourth section describes the
cancellative addition and subtraction. The tightness of each result is
stated in Chapter~\ref{chap:accuracy} (\pxref{sec:accuracy-table}), and the
names of the standard are the subject of Chapter~\ref{chap:ieee1788}.

\section{Functional and relational operations}\label{sec:arithmetic-kinds}
\cindex{functional arithmetic}\cindex{relational arithmetic}%

The containment principle of floating-point interval arithmetic
(\pxref{sec:containment}) requires that, for any operation $\circ$ and any
intervals $I$ and $J$,
\begin{equation*}
I\circ J=\hull{\{i\circ j\mid i\in I, j\in J\}}
\end{equation*}
where \hull{} maps any set of real numbers to the smallest floating-point
interval that contains it. The interval square root, for instance, is
\begin{equation*}
\sqrt{I}=\hull{\{\sqrt{i}\mid i\in I\}}
\end{equation*}
and, the square root being increasing, $\sqrt{\itv{1}{2}}=\itv{1}{\sqrt{2}}$.
This is the \dfn{functional square root}: it extends the function
$x\mapsto\sqrt{x}$ to intervals. Another reading of the square root gives
both the negative and the positive values.

\begin{Def}\label{def:arithmetic-relational-sqrt}
The \dfn{relational square root} of an interval $I$ is the hull of the real
numbers whose square is in $I$:
\begin{equation*}
\sqrt{_rI}=\hull{\{j\in\RSet\mid\exists i\in I\colon j^2=i\}}.
\end{equation*}
\end{Def}

\noindent Hence $\sqrt{_r\itv{1}{2}}=\itv{-\sqrt{2}}{\sqrt{2}}$.

\begin{Ex}
The relational square root arises from the relation $x^2=y$ between two
variables, which a program may want to write $x=\sqrt{_ry}$ to compute $x$
from $y$. The functional square root is not suitable there: its result lies
in \itv{0}{+\infty}, and the negative values of $x$ would be lost.
\end{Ex}

Some applications, mainly in constraint programming, need such operators.
GAOL therefore offers both the functional and the relational versions of the
main arithmetic operators (\pxref{sec:relational-idea}); a relational
operator contracts one operand of a relation from the others, which is how
the interval algorithms of Chapter~\ref{chap:overview} use it
(\pxref{sec:tools}).

\paragraph{Domains.}
A function is computed on the part of its argument that lies within its
domain, which is the domain of IEEE 1788-2015 (Tables~9.1 and 10.5): the
part outside is left out, and an argument with no point of the domain gives
the empty set, as \code{log(\itv{-4}{0})} does.\newinvfive{} Every function
gives the empty set for an empty argument. How tight each result is, for each
operation, is stated in Chapter~\ref{chap:accuracy}
(\pxref{sec:accuracy-table}).
\cindex{domain of a function}%

\section{Functional arithmetic}\label{sec:arithmetic-functional}

Each operation of this section is the hull of the image, by the function of
the same name, of the part of its arguments within its domain. The entries
give, for each one, the operation of IEEE 1788-2015 it implements, the
domain where it is not the whole real line, and what GAOL~v5 changed.

\subsection{Basic operations}\label{sec:arithmetic-basic}

\defmethodx{interval}{interval}{operator+}{(void) const}
\deffunx{interval}{operator+}{(const interval\& \var{I}, double \var{d})}
\deffunx{interval}{operator+}{(double \var{d}, const interval\& \var{I})}
\begin{deffun}{interval}{operator+}{(const interval\& \var{I1}, const interval\& \var{I2})}
  Addition of two intervals, or of an interval and a double: \ieee{add} of
  IEEE 1788-2015.
\end{deffun}

\defmethodx{interval}{interval}{operator-}{(void) const}
\deffunx{interval}{operator-}{(const interval\& \var{I}, double \var{d})}
\deffunx{interval}{operator-}{(double \var{d}, const interval\& \var{I})}
\begin{deffun}{interval}{operator-}{(const interval\& \var{I1}, const interval\& \var{I2})}
  Negation (\ieee{neg}), or subtraction (\ieee{sub}) of two intervals, or of
  an interval and a double.
\end{deffun}

\deffunx{interval}{operator*}{(const interval\& \var{I}, double \var{d})}
\deffunx{interval}{operator*}{(double \var{d}, const interval\& \var{I})}
\begin{deffun}{interval}{operator*}{(const interval\& \var{I1}, const interval\& \var{I2})}
  Multiplication of two intervals, or of an interval and a double: \ieee{mul}
  of IEEE 1788-2015. A zero bound times an infinite bound counts as $0$, as
  the product of a real number of the interval by 0.\newinvfive{} \code{d * I}
  is computed as \code{I * d}, and \code{n * I} of an integer \code{n} as
  \code{I * n}: the bounds are the same, with the same signs of a zero bound
  (\pxref{sec:zero-sign}).\newinvfive
\end{deffun}

\deffunx{interval}{operator/}{(const interval\& \var{I}, double \var{d})}
\deffunx{interval}{operator/}{(double \var{d}, const interval\& \var{I})}
\begin{deffun}{interval}{operator/}{(const interval\& \var{I1}, const interval\& \var{I2})}
  Functional division of two intervals, or of an interval and a double:
  \ieee{div} of IEEE 1788-2015. The quotient by an interval containing 0 is
  the hull of the quotients by its nonzero part, and the quotient by
  \itv{0}{0} is the empty set.
\begin{example}
cout << interval(1,2) * interval(-3,4) << ' '
     << interval(0,1) * interval(1,GAOL_INFINITY) << ' '
     << interval(1,2) / interval(0,1) << ' '
     << interval(1,2) / interval(-1,1) << ' '
     << interval(1,2) / interval::zero();
@outputs^[-6, 8] [0, inf] [1, inf] [-inf, inf] [empty]~
\end{example}
\end{deffun}

\deffunx{interval}{operator\%}{(const interval\& \var{I}, double \var{d})}
\deffunx{interval}{operator\%}{(double \var{d}, const interval\& \var{I})}
\begin{deffun}{interval}{operator\%}{(const interval\& \var{I1}, const interval\& \var{I2})}
  Relational division of two intervals, or of an interval and a double
  (\pxref{sec:relational-arithmetic}).
\end{deffun}

\defmethodx{interval}{interval}{inverse}{(void) const}
\begin{deffun}{interval}{inverse}{(const interval\& \var{I})}
  \begin{operation}
    \op{$1/I$}
  \end{operation}
  Inverse of \var{I} (of \code{self} for the method): \ieee{recip} of
  IEEE 1788-2015. The inverse of an interval containing 0 is the hull of the
  inverses of its nonzero part.
\begin{example}
cout << inverse(interval(2,4)) << ' ' << inverse(interval(-1,2));
@outputs^[0.25, 0.5] [-inf, inf]~
\end{example}
\end{deffun}

\begin{deffun}{interval}{sqr}{(const interval\& \var{I})}
  \begin{operation}
    \op{$\var{I}^2$}
  \end{operation}
  Square of \var{I}: \ieee{sqr} of IEEE 1788-2015.
\end{deffun}

\begin{deffun}{interval}{sqrt}{(const interval\& \var{I})}
  \begin{operation}
    \op{$\sqrt{I}$}
  \end{operation}
  Functional square root of \var{I}, on the part of \var{I} in
  \itv{0}{+\infty}: \ieee{sqrt} of IEEE 1788-2015. The bounds are the
  tightest ones, whatever the rounding of the square root of the C
  library.\newinvfive
\begin{example}
cout << sqrt(interval(-4,9)) << ' ' << sqrt(interval(-4,-1));
@outputs^[0, 3] [empty]~
\end{example}
\end{deffun}

\begin{deffun}{interval}{fma}{(const interval\& \var{X}, const interval\& \var{Y}, const interval\& \var{Z})}
  \begin{operation}
    \op{$\hull{\{xy+z\mid x\in X, y\in Y, z\in Z\}}$}
  \end{operation}
  Fused multiply-add: \ieee{fma} of IEEE 1788-2015. Each bound is rounded
  once, and the result is the tightest enclosure.\newinvfive
\begin{example}
cout << fma(interval(1,2), interval(3,4), interval(-1,1));
@outputs^[2, 9]~
\end{example}
\end{deffun}

\subsection{Powers and roots}\label{sec:powers}
\cindex{power}\cindex{root}%

The functions of this subsection are those of the namespace \code{gaol}. The
namespace \code{gaol\_ieee1788} has its own \code{pow}, the \ieee{pow} of
the standard: the two powers are compared in Chapter~\ref{chap:overview}
(\pxref{sec:pow-namespaces}).

\deffunx{interval}{pow}{(const interval\& \var{I}, int \var{n})}
\deffunx{interval}{pow}{(const interval\& \var{I}, unsigned int \var{n})}
\deffunx{interval}{pow}{(const interval\& \var{I}, double \var{p})}
\deffunx{interval}{pow}{(const interval\& \var{I}, long \var{n})}
\deffunx{interval}{pow}{(const interval\& \var{I}, long long \var{n})}
\begin{deffun}{interval}{pow}{(const interval\& \var{I}, const interval\& \var{J})}
  \begin{operation}
    \op{$I^n$}
    \op{$I^p$}
    \op{$I^J$}
  \end{operation}
  Power function. \code{pow} is hybrid in GAOL v5.\newinvfive
  \begin{itemize}
  \item[$\bullet$] \textbf{An integer exponent} takes the \emph{integer
    power}, \ieee{pown} of IEEE 1788-2015: \var{I} to the power \var{n},
    the negative part of \var{I} included, where $x^0=1$ for every $x$, 0
    included, and $x^n$ has no value at 0 for $n<0$. A double \var{p} and a
    degenerate interval $\var{J}=\itv{n}{n}$ whose value is an integer take
    the integer power as well, beyond the \code{int}s too.\newinvfive{} An
    exponent \var{n} of another integer type (\code{long}, \code{long long},
    \code{short}\dots), with which \code{pow} was ambiguous, gives the same
    integer power; beyond 64 bits (an \code{\_\_int128}), it gives the power of $\pm(2^{63}-1)$ or
    $\pm(2^{63}-2)$, with the sign and the parity of \var{n}, every power of
    a double but $0$, $\pm1$ and $\pm\infty$ being beyond the doubles
    there.\newinvfive
  \item[$\bullet$] \textbf{Any other exponent} takes the \emph{real power},
    \ieee{pow} of IEEE 1788-2015, on the part of \var{I} in \itv{0}{+\infty},
    where $0^y$ is $0$ for $y>0$ and has no value for $y\leq0$. An exponent
    \itv{+\infty}{+\infty} or \itv{-\infty}{-\infty}, or an infinite or NaN
    \var{p}, gives the empty set.
  \end{itemize}
  The integer powers are computed from exact products within the
  \code{int}s, $x^{-n}$ included. They are the tightest enclosures, or one
  double beyond where the power is closer to a double than twice a bound of
  their error proved beforehand: $2^{-98}$ for $n=3$, $2^{-88}$ for
  $n=1000$, $2^{-62}$ at most, relatively. They are exact where the power is
  a double. Beyond the \code{int}s, the integer powers come from the power of
  CORE-MATH at the bounds of \var{I}: they are the tightest enclosures where
  $n$ is a double, and within three doubles of them
  otherwise.\newinvfive{} The real power takes the power of CORE-MATH at the
  corners of the box $I\times J$ where it is the least and the greatest. The
  upper bound is the tightest one, and so is the lower bound: it is the power
  itself where that power is a double, as $4^{0.5}$, $9^{1.5}$ or $2^{-3}$,
  and the double below the value of CORE-MATH otherwise. At a corner where
  the base is $0$ or $+\infty$ or the exponent is infinite, the bound is the
  limit of $x^y$ there, $0$ or $+\infty$ ($1$ for a base $1$ or an exponent
  $0$): \code{pow(\itv{4}{+\infty}, 0.5)} is $[2, +\infty]$.\newinvfive
\begin{example}
cout << pow(interval(-4,-1), 2) << ' '
     << pow(interval(-4,-1), interval(2.0)) << ' '
     << pow(interval(-4,9), interval(0.5)) << ' '
     << pow(interval(-4,-1), interval(1.5, 2.5)) << endl;
@outputs^[1, 16] [1, 16] [0, 3] [empty]~
cout << pow(interval(0.0), interval(0.0)) << ' '
     << pow(interval(-2,2), -2) << ' '
     << pow(interval(4), 0.5) << endl;
@outputs^[1] [0.25, inf] [2]~
\end{example}
\end{deffun}

\begin{Rem}\label{rem:arithmetic-pow-monotone}
The result of \code{pow} is not monotone for the inclusion of \var{J}:
\code{pow(\itv{-4}{-1}, \itv{2}{2})} is \itv{1}{16}, and
\code{pow(\itv{-4}{-1}, \itv{1.5}{2.5})} the empty set. The result for a
degenerate integer exponent $[n]$ is not within the result for an exponent
around $n$, and \code{pow} is not one function of $x$ and $y$. GAOL~4
truncated a double exponent \var{p}, and computed the powers of a negative
base on its magnitude.
\end{Rem}

\deffunx{interval}{nth\_root}{(const interval\& \var{I}, unsigned int \var{n})}
\deffunx{interval}{nth\_root}{(const interval\& \var{I}, long long \var{q})}
\deffunx{interval}{nth\_root}{(const interval\& \var{I}, unsigned long long \var{q})}
\begin{deffun}{interval}{nth\_root}{(const interval\& \var{I}, int \var{q})}
  \begin{operation}
    \op{$^{n}\sqrt{I}$}
  \end{operation}
  Computes the \var{n}th functional root of \var{I}: \ieee{rootn} of
  IEEE 1788-2015. For an odd \var{n}, the root of a negative number is the
  opposite of the root of its magnitude, and the root is defined on the whole
  real line; for an even \var{n}, only the part of \var{I} in
  \itv{0}{+\infty} is considered.\newinvfive{} A negative \var{q} gives
  $1/x^{1/|q|}$, defined on $\RSet\setminus\{0\}$ for an odd \var{q} and on
  $\left]0,+\infty\right]$ for an even one. The function returns the empty
  set for $n=0$.

  The bounds are the tightest ones, or up to 2 doubles beyond (4 for a
  negative \var{q}): each is proved with integer powers. The cube root is
  that of CORE-MATH, the tightest one, and the root of a double that is the
  \var{n}th power of a double is that double. An order of another integer
  type (\code{long}, \code{long long}\dots), with which \code{nth\_root} was
  ambiguous, is taken as an \code{int} or an \code{unsigned int} where it is
  one. Beyond them, the roots are enclosed by \code{pow} with the exponent
  $1/q$ taken as an interval, $x^{1/q}$ being within $2^{-21}$ of 1 for every
  double $x>0$ there, after one check of the rounding direction for a
  \code{long long} or an \code{unsigned long long} order.\newinvfive
\begin{example}
cout << nth_root(interval(-8,27), 3) << ' '
     << nth_root(interval(-4,16), 2) << ' '
     << nth_root(interval(4,16), -2) << ' '
     << nth_root(interval(-8,-1), -3);
@outputs^[-2, 3] [0, 4] [0.25, 0.5] [-1, -0.5]~
\end{example}
\end{deffun}

\subsection{Exponentials and logarithms}\label{sec:arithmetic-exp-log}
\cindex{exponential}\cindex{logarithm}%

The functions of this subsection, and of the next two, are bounded with
CORE-MATH at the bounds of the interval, or at the points within it where
the function reaches its extrema. CORE-MATH being correctly rounded, the
value at a bound, computed rounding upward, is the upper bound, and the
double below it is the lower bound. These are the tightest bounds, unless the
exact value is a double; the functions give that double where they know that
it happens ($\exp(0)=1$, $\log(1)=0$, $\sin(0)=0$, $2^3=8$\ldots).\newinvfive
\cindex{CORE-MATH}%

\begin{deffun}{interval}{exp}{(const interval\& \var{I})}
  Exponential of \var{I}.
\end{deffun}

\deffunx{interval}{exp2}{(const interval\& \var{I})}
\begin{deffun}{interval}{exp10}{(const interval\& \var{I})}
  $2^x$ and $10^x$ for $x$ in \var{I}: \ieee{exp2} and \ieee{exp10} of IEEE
  1788-2015.\newinvfive
\end{deffun}

\deffunx{interval}{expm1}{(const interval\& \var{I})}
\deffunx{interval}{exp2m1}{(const interval\& \var{I})}
\begin{deffun}{interval}{exp10m1}{(const interval\& \var{I})}
  $e^x-1$, $2^x-1$ and $10^x-1$ for $x$ in \var{I}, within
  \itv{-1}{+\infty}: \ieee{expm1}, \ieee{exp2m1} and \ieee{exp10m1} of IEEE
  1788-2015 (Table~10.5).\newinvfive
\end{deffun}

\begin{deffun}{interval}{log}{(const interval\& \var{I})}
  Natural logarithm of \var{I}, defined on $\left]0,+\infty\right]$. An
  interval holding no positive number gives the empty set, where GAOL~4 kept
  its part in \itv{0}{+\infty} and gave \itv{-\infty}{-\text{MAX}} for
  \code{log(\itv{-4}{0})}.\newinvfive
\end{deffun}

\deffunx{interval}{log2}{(const interval\& \var{I})}
\begin{deffun}{interval}{log10}{(const interval\& \var{I})}
  Logarithms in base 2 and 10 of \var{I}: \ieee{log2} and \ieee{log10} of
  IEEE 1788-2015.\newinvfive
\end{deffun}

\deffunx{interval}{log1p}{(const interval\& \var{I})}
\deffunx{interval}{log2p1}{(const interval\& \var{I})}
\begin{deffun}{interval}{log10p1}{(const interval\& \var{I})}
  $\log(1+x)$, $\log_2(1+x)$ and $\log_{10}(1+x)$ for $x$ in \var{I}, defined
  on $\left]-1,+\infty\right]$, $-\infty$ being their limit at $-1$:
  \ieee{logp1}, \ieee{log2p1} and \ieee{log10p1} of IEEE 1788-2015
  (Table~10.5).\newinvfive
\begin{example}
cout << exp(interval(0.0)) << ' ' << log(interval(1)) << ' '
     << log(interval(-4,0.0)) << ' '
     << log(interval(-4,1)) << endl;
@outputs^[1] [0] [empty] [-inf, 0]~
cout << exp2(interval(3)) << ' ' << log2(interval(0.25)) << ' '
     << log10(interval(0.0,1000)) << ' '
     << expm1(interval(-1e300,0.0)) << ' '
     << log1p(interval(-2,0.0)) << endl;
@outputs^[8] [-2] [-inf, 3] [-1, 0] [-inf, 0]~
\end{example}
\end{deffun}

\subsection{Trigonometric functions}\label{sec:arithmetic-trigonometric}
\cindex{trigonometric functions}%

\begin{deffun}{interval}{cos}{(const interval\& \var{I})}
  Returns the cosine of \var{I}, within \itv{-1}{1}.
\end{deffun}

\begin{deffun}{interval}{sin}{(const interval\& \var{I})}
  Returns the sine of \var{I}, within \itv{-1}{1}.
\end{deffun}

\begin{deffun}{interval}{tan}{(const interval\& \var{I})}
  Returns the tangent of \var{I}: \itv{-\infty}{+\infty} when \var{I} holds a
  pole. No double is a pole: \code{M\_PI\_2}, the double below $\pi/2$, is
  $6.1\cdot10^{-17}$ short of it. The interval
  \itv{-\code{M\_PI\_2}}{\code{M\_PI\_2}} holds no pole and has a finite
  tangent (the fourth result of the example); its width is the double below
  $\pi$, and the intervals of width between that double and $\pi$ are a
  theoretical case.

  The three functions find the pieces of \var{I} on which they are monotonic
  by dividing its bounds by an enclosure of $\pi$. Where the quotients cannot
  tell, next to an extremum or a pole and at the large magnitudes, they use
  the signs of the derivative at the bounds, which CORE-MATH gives exactly.
  Their bounds are within one double of the tightest ones at every
  magnitude.\newinvfive{} GAOL~4 gave \itv{-1}{1} for
  \code{cos(\itv{2^{60}}{2^{60}})}, and computed the sine as a shifted
  cosine, which widened narrow intervals.
\begin{example}
cout << cos(interval(1,4)) << endl
     << sin(interval(1,2)) << endl
     << tan(interval(1,2)) << endl
     @rem^// 1.5707963267948966 is M_PI_2, the double below pi/2~
     << tan(interval(-1.5707963267948966, 1.5707963267948966)) << endl
     @rem^// 1152921504606846976 is 2 to the 60~
     << cos(interval(1152921504606846976.0)) << endl;
@outputs^[-1, 0.5403023058681398]~
@outputs^[0.8414709848078965, 1]~
@outputs^[-inf, inf]~
@outputs^[-1.633123935319537e+16, 1.633123935319537e+16]~
@outputs^[-0.5567960822766418, -0.5567960822766416]~
\end{example}
\end{deffun}

\begin{deffun}{interval}{acos}{(const interval\& \var{I})}
  Returns the arccosine of \var{I}, on the part of \var{I} in \itv{-1}{1}.
\end{deffun}

\begin{deffun}{interval}{asin}{(const interval\& \var{I})}
  Returns the arcsine of \var{I}, on the part of \var{I} in \itv{-1}{1}.
\end{deffun}

\begin{deffun}{interval}{atan}{(const interval\& \var{I})}
  Returns the arctangent of \var{I}.
\begin{example}
cout << acos(interval(1,3)) << ' ' << asin(interval(1)) << ' '
     << atan(interval(1));
@outputs^[0] [1.570796326794896, 1.570796326794897]~
@outputs^[0.7853981633974482, 0.7853981633974484]~
\end{example}
\end{deffun}

\begin{deffun}{interval}{atan2}{(const interval\& \var{Y}, const interval\& \var{X})}
  Returns the hull of the angles in $(-\pi,\pi]$ of the points of ordinate
  in \var{Y} and abscissa in \var{X}, the point $(0,0)$ left out, as the
  \ieee{atan2} of IEEE~1788-2015 does: \code{atan2([0],[0])} is empty. The
  result is $[-\pi,\pi]$ whenever the box holds points of the half-line
  $y=0$, $x<0$, where the angle is $\pi$, and points below it, where the
  angle is next to $-\pi$. GAOL~4 declared the function, but it raised
  \code{unavailable\_feature\_error}.\newinvfive
\begin{example}
interval x(1,GAOL_INFINITY);
cout << atan2(interval(0.0), interval(0.0)) << endl
     << atan2(interval(-1,1), interval(-1)) << endl
     << atan2(x, x) << endl;
@outputs^[empty]~
@outputs^[-3.141592653589794, 3.141592653589794]~
@outputs^[0, 1.570796326794897]~
\end{example}
\end{deffun}

\deffunx{interval}{sinpi}{(const interval\& \var{I})}
\deffunx{interval}{cospi}{(const interval\& \var{I})}
\begin{deffun}{interval}{tanpi}{(const interval\& \var{I})}
  $\sin(\pi x)$, $\cos(\pi x)$ and $\tan(\pi x)$ for $x$ in \var{I}:
  \ieee{sinPi}, \ieee{cosPi} and \ieee{tanPi} of IEEE 1788-2015
  (Table~10.5).\newinvfive{} Their extrema and poles are at the multiples of
  $1/2$, which are doubles, and are found exactly, at every magnitude.
  \code{tanpi} has no value at the half-integers, which are left out: a
  half-integer alone gives the empty set.
\end{deffun}

\deffunx{interval}{asinpi}{(const interval\& \var{I})}
\deffunx{interval}{acospi}{(const interval\& \var{I})}
\deffunx{interval}{atanpi}{(const interval\& \var{I})}
\begin{deffun}{interval}{atan2pi}{(const interval\& \var{Y}, const interval\& \var{X})}
  $\operatorname{asin}(x)/\pi$, $\operatorname{acos}(x)/\pi$,
  $\operatorname{atan}(x)/\pi$ and $\operatorname{atan2}(y,x)/\pi$:
  \ieee{asinPi}, \ieee{acosPi}, \ieee{atanPi} and \ieee{atan2Pi} of IEEE
  1788-2015 (Table~10.5), within \itv{-1/2}{1/2}, \itv{0}{1},
  \itv{-1/2}{1/2} and \itv{-1}{1}.\newinvfive
\begin{example}
cout << sinpi(interval(1e17)) << ' '
     << sin(interval::pi() * 1e17) << endl;
@outputs^[0] [-1, 1]~
cout << tanpi(interval(0.5)) << ' '
     << tanpi(interval(0.25,0.75)) << ' '
     << asinpi(interval(1)) << ' '
     << atan2pi(interval(1), interval(-1)) << endl;
@outputs^[empty] [-inf, inf] [0.5] [0.75]~
\end{example}
\end{deffun}

\subsection{Hyperbolic functions}\label{sec:arithmetic-hyperbolic}
\cindex{hyperbolic functions}%

\begin{deffun}{interval}{cosh}{(const interval\& \var{I})}
  Returns the hyperbolic cosine of \var{I}.
\end{deffun}

\begin{deffun}{interval}{acosh}{(const interval\& \var{I})}
  Returns the hyperbolic arccosine of \var{I}, on the part of \var{I} in
  \itv{1}{+\infty}.
\end{deffun}

\begin{deffun}{interval}{sinh}{(const interval\& \var{I})}
  Returns the hyperbolic sine of \var{I}.
\end{deffun}

\begin{deffun}{interval}{asinh}{(const interval\& \var{I})}
  Returns the hyperbolic arcsine of \var{I}.
\end{deffun}

\begin{deffun}{interval}{tanh}{(const interval\& \var{I})}
  Returns the hyperbolic tangent of \var{I}, within \itv{-1}{1}.
\end{deffun}

\begin{deffun}{interval}{atanh}{(const interval\& \var{I})}
  Returns the hyperbolic arctangent of \var{I}, defined on
  $\left]-1,1\right[$, $-\infty$ and $+\infty$ being its limits at $-1$ and
  $1$. The part of \var{I} outside is left out, and an interval with no point
  of it gives the empty set, as \code{atanh(\itv{1}{5})} and
  \code{atanh(\itv{-1}{-1})} do.
\begin{example}
cout << cosh(interval(-1,2)) << ' '
     << tanh(interval(20,30)) << endl
     << acosh(interval(0,1)) << ' '
     << atanh(interval(-2,2)) << ' '
     << atanh(interval(1,5)) << endl;
@outputs^[1, 3.762195691083632] [0.9999999999999998, 1]~
@outputs^[0] [-inf, inf] [empty]~
\end{example}
\end{deffun}

\begin{Rem}\label{rem:arithmetic-hyperbolic-coremath}
The six hyperbolic functions are those of CORE-MATH. GAOL~4 took them from
the mathematical library of the system, which is sometimes further than one
double from the exact value.\newinvfive
\end{Rem}

\subsection{Other functions}\label{sec:arithmetic-other}

\begin{deffun}{interval}{hypot}{(const interval\& \var{X}, const interval\& \var{Y})}
  $\sqrt{x^2+y^2}$ for $x$ in \var{X} and $y$ in \var{Y}: \ieee{hypot} of
  IEEE 1788-2015 (Table~10.5), computed at the points of the box nearest to
  the origin and farthest from it.\newinvfive
\end{deffun}

\begin{deffun}{interval}{rsqrt}{(const interval\& \var{I})}
  $1/\sqrt{x}$ for $x$ in \var{I}, defined on $\left]0,+\infty\right]$,
  $+\infty$ being its limit at 0: \ieee{rSqrt} of IEEE 1788-2015
  (Table~10.5).\newinvfive
\begin{example}
cout << hypot(interval(3), interval(4)) << ' '
     << hypot(interval(-3,1), interval(4,5)) << endl
     << rsqrt(interval(4,16)) << ' '
     << rsqrt(interval(0.0,4)) << endl;
@outputs^[5] [4, 5.830951894845301]~
@outputs^[0.25, 0.5] [0.5, inf]~
\end{example}
\end{deffun}

\begin{deffun}{interval}{compoundm1}{(const interval\& \var{X}, const interval\& \var{Y})}
  $(1+x)^y-1$ for $x$ in \var{X} and $y$ in \var{Y}: \ieee{compoundm1} of
  IEEE 1788-2015 (Table~10.5). It keeps the digits of a small $x$ that
  \code{pow(1+x, y) - 1} loses, $1+x$ being rounded and the difference
  cancelling what is left. It is defined for $x>-1$, and for $x=-1$ when
  $y>0$, where it is $-1$. The part of the box outside is left out, and $-1$
  and $+\infty$, the limits at $x=-1$, at $x=+\infty$ and for an infinite
  $y$, are bounds. CORE-MATH has no function for it: the value at a corner
  of the box is computed with integers, 64 bits in a fast phase that settles
  71 to 95\,\% of the pairs, 128 bits and a bound of the error in the others.
  It takes 0.1 microsecond for each of the one or two corners where the
  bounds are, where the fast phase settles it, and 0.6 to 0.9 elsewhere. The
  bounds are the tightest ones, the double itself where the value is one
  (2.375 for $x=1.5-1$ and $y=3$, but also $2^m-1$ for $|m|\le53$, and the
  quotient $-x/(1+x)$ of $y=-1$, divided exactly), but in a case out of
  $10^{17}$ or so, where a bound is a double beyond.\newinvfive
\begin{example}
cout << compoundm1(interval(0.5), interval(2)) << ' '
     << compoundm1(interval(1,3), interval(1,2)) << endl
     << compoundm1(interval(-1), interval(0,1)) << ' '
     << compoundm1(interval(-1,0), interval(-1)) << endl
     << compoundm1(interval(1e-20), interval(3)) << endl;
@outputs^[1.25] [1, 15]~
@outputs^[-1] [0, inf]~
@outputs^[2.999999999999999e-20, 3.000000000000001e-20]~
\end{example}
\end{deffun}

\deffunx{interval}{erf}{(const interval\& \var{I})}
\begin{deffun}{interval}{erfc}{(const interval\& \var{I})}
  The error function $\operatorname{erf}(x)=\frac{2}{\sqrt{\pi}}\int_0^x
  e^{-t^2}\,dt$, increasing from $-1$ to 1, and its complement
  $\operatorname{erfc}(x)=1-\operatorname{erf}(x)$, decreasing from 2 to 0,
  for $x$ in \var{I}: the \code{erf} and \code{erfc} of C, which CORE-MATH
  computes, with the tightest bounds, exact at 0 and at the infinities.
  $\operatorname{erfc}(x)$ is subnormal from $x=26.55$ on, and below the
  least subnormal from $x=27.22$ on, where the lower bound is 0. IEEE
  1788-2015 does not name them, and \code{gaol\_ieee1788} does not hold
  them: a call \code{erf(\var{x})} on an interval finds them all the same,
  by argument-dependent lookup.\newinvfive
\begin{example}
cout << erf(interval(0.0)) << ' ' << erf(interval(-1,1)) << ' '
     << erf(interval(6,7)) << endl
     << erfc(interval(0.0,1)) << ' ' << erfc(interval(28,30)) << endl;
@outputs^[0] [-0.8427007929497149, 0.8427007929497149] [0.9999999999999998, 1]~
@outputs^[0.1572992070502851, 1] [0, 4.940656458412466e-324]~
\end{example}
\end{deffun}

\section{Relational arithmetic}\label{sec:relational-arithmetic}
\cindex{reverse functions|see{relational arithmetic}}%
\cindex{relational arithmetic}%

The relational functions of GAOL are the \dfn{reverse functions} of
IEEE 1788-2015 (10.5.4, Table~10.1), with their arguments in another order:
each computes the hull of the $x$ of an interval \var{I} whose image by the
function is in \var{J}. Their correspondence is given at the end of this
section (Table~\ref{tab:arithmetic-reverse}), and the namespace
\code{gaol\_ieee1788} provides them under their names
(Chapter~\ref{chap:ieee1788}).

\defmethodx{interval}{interval}{operator\% }{(double \var{d}) const}
\defmethodx{interval}{interval}{operator\% }{(const interval\& \var{I}) const}
\begin{deffun}{interval}{operator\%}{(double \var{d}, const interval\& \var{I})}
  \begin{operation}
    \op{$\code{self}\,\%\,\var{I}=\hull{\{z\in \RSet\mid
      \exists x\in \code{self}, \exists y\in I\colon x=yz\}}$}
  \end{operation}
  Relational division: \code{div\_rel(self, \var{I}, \itv{-\infty}{+\infty})}.
\begin{example}
cout << interval(1,2) % interval(0,1) << ' '
     << interval(1,2) % interval(-1,1);
@outputs^[1, inf] [-inf, inf]~
\end{example}
\end{deffun}

\subsection{$(n+1)$-ary relational functions}\label{sec:arithmetic-nary}

Consider the relation $y=\cos x$ between two interval variables $x$ and $y$.
The equivalent form $x=\operatorname{acos}y$ cannot be computed with the
function $\operatorname{acos}$, whose result always lies in \itv{0}{\pi}: a
relational version of $\operatorname{acos}$ is needed. For any value of $y$
in \itv{-1}{1} there are infinitely many $x$ such that $y=\cos x$, so the
domain of $x$ has to be taken into account, and the relational arccosine is
a binary operator.

\begin{Def}\label{def:arithmetic-acos-rel}
The \dfn{relational arccosine} of $Y$ with respect to $X$ is
\begin{equation*}
\code{acos\_rel}(Y,X)=\hull{\{x\in X\mid\exists y\in Y\colon y=\cos x\}}.
\end{equation*}
\end{Def}

\noindent It is to be contrasted with the functional arccosine,
\begin{equation*}
\code{acos}(Y)=\hull{\{x\in\RSet\mid\exists y\in Y\colon
  x=\operatorname{acos} y\}}.
\end{equation*}
Figure~\ref{fig:relational-cos} shows the results of $\code{acos}(J)$ and of
$\code{acos\_rel}(J,I)$ on the same arguments.

\begin{figure}[htbp]
  \centering
  \includegraphics[width=.85\textwidth,bb=11 262 553 470]{relation-cos}
  \caption{The functional arccosine $\code{acos}(J)$ and the relational
    arccosine $\code{acos\_rel}(J, I)$}
  \label{fig:relational-cos}
\end{figure}

\begin{Rem}\label{rem:arithmetic-empty-preimage}
An interval \var{I} that meets no piece of the preimage gives the empty
set.\newinvfive{} The same holds for every relational function of this
section: the result is empty when no $x$ of the interval contracted is in
the relation.
\end{Rem}

\begin{deffun}{interval}{acos\_rel}{(const interval\& J, const interval\& I)}
  \begin{operation}
    \op{$\code{acos\_rel}(J,I)=
      \hull{\{x\in I\mid\exists y\in J\colon y=\cos x\}}$}
  \end{operation}
  Returns the \emph{relational arccosine} of \var{J} with respect to \var{I}.
\end{deffun}

\begin{deffun}{interval}{asin\_rel}{(const interval\& J, const interval\& I)}
  \begin{operation}
    \op{$\code{asin\_rel}(J,I)=
      \hull{\{x\in I\mid\exists y\in J\colon y=\sin x\}}$}
  \end{operation}
  Returns the \emph{relational arcsine} of \var{J} with respect to \var{I}.
\end{deffun}

\begin{deffun}{interval}{atan\_rel}{(const interval\& J, const interval\& I)}
  \begin{operation}
    \op{$\code{atan\_rel}(J,I)=
      \hull{\{x\in I\mid\exists y\in J\colon y=\tan x\}}$}
  \end{operation}
  Returns the \emph{relational arctangent} of \var{J} with respect to
  \var{I}.

  The three functions enclose the multiples of $\pi$ of the pieces of the
  preimage in double-double arithmetic, and add the inverse function of
  \var{J} before rounding the sum: their bounds are within one double of the
  accurate ones.\newinvfive
\begin{example}
cout << asin_rel(interval(0.0), interval(1,7)) << endl
     << atan_rel(interval(1), interval(0,5)) << endl;
@outputs^[3.141592653589793, 6.283185307179588]~
@outputs^[0.7853981633974482, 3.926990816987242]~
\end{example}
\end{deffun}

\begin{deffun}{interval}{acosh\_rel}{(const interval\& J, const interval\& I)}
  \begin{operation}
    \op{$\code{acosh\_rel}(J,I)=
      \hull{\{x\in I\mid\exists y\in J\colon y=\cosh x\}}$}
  \end{operation}
  Returns the \emph{relational hyperbolic arccosine} of \var{J} with respect
  to \var{I}.
\end{deffun}

\deffunx{interval}{asinh\_rel}{(const interval\& J, const interval\& I)}
\begin{deffun}{interval}{atanh\_rel}{(const interval\& J, const interval\& I)}
  \begin{operation}
    \op{$\code{asinh\_rel}(J,I)=
      \hull{\{x\in I\mid\exists y\in J\colon y=\sinh x\}}$}
    \op{$\code{atanh\_rel}(J,I)=
      \hull{\{x\in I\mid\exists y\in J\colon y=\tanh x\}}$}
  \end{operation}
  Return the \emph{relational hyperbolic arcsine} and \emph{arctangent} of
  \var{J} with respect to \var{I}, that is, $\operatorname{asinh}(J)\cap I$
  and $\operatorname{atanh}(J)\cap I$, $\sinh$ and $\tanh$ being one-to-one.
  As $\tanh$ takes its values in $\left]-1,1\right[$, \code{atanh\_rel}
  gives the empty set for a \var{J} with no point of it, as \itv{1}{1}.
\begin{example}
cout << acosh_rel(interval(1,2), interval::universe()) << endl
     << asinh_rel(interval(0,1), interval::universe()) << endl
     << atanh_rel(interval(1), interval::universe()) << endl;
@outputs^[-1.316957896924817, 1.316957896924817]~
@outputs^[0, 0.8813735870195431]~
@outputs^[empty]~
\end{example}
\end{deffun}

\begin{deffun}{interval}{sqrt\_rel}{(const interval\& J, const interval\& I)}
  \begin{operation}
    \op{$\code{sqrt\_rel}(J,I)=\hull{\{x\in I\mid\exists y\in J\colon y=x^2\}}$}
  \end{operation}
  Returns the \emph{relational square root} of \var{J} with respect to
  \var{I} (Definition~\ref{def:arithmetic-relational-sqrt}).
\begin{example}
cout << sqrt_rel(interval(1,4), interval::universe()) << ' '
     << sqrt_rel(interval(1,4), interval(0,10));
@outputs^[-2, 2] [1, 2]~
\end{example}
\end{deffun}

\deffunx{interval}{nth\_root\_rel}{(const interval\& \var{J}, unsigned int \var{n}, \\\hbox{}\hfill const interval\& \var{I})}
\begin{deffun}{interval}{nth\_root\_rel}{(const interval\& \var{J}, int \var{p}, \\\hbox{}\hfill const interval\& \var{I})}
  \begin{operation}
    \op{$\code{nth\_root\_rel}(J,n,I)=
      \hull{\{x\in I\mid\exists y\in J\colon y=x^n\}}$}
    \op{$\code{nth\_root\_rel}(J,p,I)=\hull{\{x\in I\mid x^p\text{ is defined
      and in }J\}}$}
  \end{operation}
  Returns the \emph{relational inverse \var{n}-th power} of \var{J} with
  respect to \var{I}, for $n\geq1$. It takes the roots of
  \code{nth\_root()}.\newinvfive{} For an \code{int} \var{p}, or an integer
  of another type, it is the reverse of $x^p$ for every integer $p$, the
  \ieee{pownRev} of the standard. $x^0$ is $1$ for every $x$, $0$ included,
  so that the result is \var{I} where $1\in J$, and the empty set elsewhere:
  \code{nth\_root\_rel(\itv{-1}{5}, 0, I)} is \var{I}, and
  \code{nth\_root\_rel(\itv{2}{3}, 0, I)} is empty. For $p<0$, $x^p$ is
  $1/x^{|p|}$, which $0$ has not. \var{J} is then cut at $0$, and each part
  gives its roots: the hull of the roots of the whole of \var{J} would keep
  the part of \var{I} between them. The roots reach $0$ where \var{J} is
  unbounded, $0$ being a limit there and no solution:
  \code{nth\_root\_rel(\itv{1}{+\infty}, -2, \itv{0}{0})} is empty. The
  bounds are those of \code{nth\_root(J, p)}, the root and then its inverse:
  up to four doubles beyond the tightest ones for $p<0$. An exponent of type
  \code{double} or \code{float} is refused at compile time, as the ambiguous
  \code{nth\_root(x, 2.0)} is: it was converted to an \code{unsigned int},
  which cut $2.5$ to $2$.
\begin{example}
interval J(1,16);
cout << nth_root_rel(J, 4, interval::universe()) << ' '
     << nth_root_rel(J, 4, interval(-10,0.0));
@outputs^[-2, 2] [-2, -1]~
\end{example}
\begin{example}
cout << nth_root_rel(interval(1,8), -3, interval::universe()) << ' '
     << nth_root_rel(interval(-1,2), -3, interval(-0.5,0.5)) << ' '
     << nth_root_rel(interval(1,GAOL_INFINITY), -2, interval(0,3)) << ' '
     << nth_root_rel(interval(1,2), 0, interval(-5,5));
@outputs^[0.5, 1] [empty] [0, 1] [-5, 5]~
\end{example}
\end{deffun}

\begin{deffun}{interval}{invabs\_rel}{(const interval\& J, const interval\& I)}
  \begin{operation}
    \op{$\code{invabs\_rel}(J,I)=\hull{\{x\in I\mid |x|\in J\}}$}
  \end{operation}
  Returns the \emph{relational inverse absolute value} of \var{J} with
  respect to \var{I}. The negative part of \var{J} holds no magnitude:
  GAOL~4 took it as magnitudes, and gave \itv{-2}{2} where the result is
  empty.\newinvfive
\begin{example}
cout << invabs_rel(interval(1,2), interval::universe()) << ' '
     << invabs_rel(interval(-2,-1), interval(-10,10));
@outputs^[-2, 2] [empty]~
\end{example}
\end{deffun}

\begin{deffun}{interval}{div\_rel}{(const interval\& K, const interval\& J, const interval\& I)}
  \begin{operation}
    \op{$\code{div\_rel}(K,J,I)=
      \hull{\{x\in I\mid\exists z\in K,\exists y\in J\colon z=xy\}}$}
  \end{operation}
  Returns the \emph{ternary relational division} of \var{K} by \var{J} with
  respect to \var{I}, which is empty when no $x$ of \var{I} is in the
  relation.\newinvfive
\begin{example}
cout << div_rel(interval(1,2), interval(-1,1),
                interval(0,10)) << ' '
     << div_rel(interval(5,6), interval(1,2),
                interval(10,20));
@outputs^[1, 10] [empty]~
\end{example}
\end{deffun}

\begin{deffun}{interval}{pow\_rel}{(const interval\& K, const interval\& J, const interval\& I)}
  \begin{operation}
    \op{$\code{pow\_rel}(K,J,I)=
      \hull{\{x\in I\mid\exists z\in K,\exists y\in J\colon z=x^y\}}$}
  \end{operation}
  Returns the \emph{relational base} of the power $z=x^y$: the hull of the
  $x$ of \var{I} with an $x^y$ in \var{K} for a $y$ of \var{J}. It contracts
  the base of $z=x^y$ as \code{div\_rel(K, J, I)} contracts a factor of
  $z=xy$, with the arguments in the same order.\newinvfive{} $x^y$ is the
  \code{pow} of IEEE 1788-2015 (Table~9.1), that of
  \code{gaol\_ieee1788::pow}: $e^{y\log x}$ for $x>0$, $0$ for $x=0$ and
  $y>0$, and no value elsewhere, so that no $x<0$ is a solution, in either
  namespace. The \code{pow} of \code{gaol} takes the integer power for a
  degenerate integer exponent $[n]$, which has negative bases too
  (Remark~\ref{rem:arithmetic-pow-monotone}): the reverse of $x^n$ is
  \code{nth\_root\_rel(K, n, I)}.
\begin{example}
cout << pow_rel(interval(4,9), interval(2), interval::universe()) << ' '
     << pow_rel(interval(2,3), interval(1,GAOL_INFINITY), interval(0,1)) << ' '
     << pow_rel(interval(1,4), interval(2), interval(-3,3));
@outputs^[2, 3] [empty] [1, 2]~
\end{example}
  $x^y\geq2$ for a $y\geq1$ asks $x>1$: the second result is empty.
\end{deffun}

\begin{deffun}{interval}{pow\_exponent\_rel}{(const interval\& K, const interval\& H, \\\hbox{}\hfill const interval\& I)}
  \begin{operation}
    \op{$\code{pow\_exponent\_rel}(K,H,I)=\hull{\{y\in I\mid\exists z\in K,
      \exists x\in H\colon z=x^y\}}$}
  \end{operation}
  Returns the \emph{relational exponent} of the power $z=x^y$, the same
  \code{pow} as \code{pow\_rel()}: the hull of the $y$ of \var{I} with an
  $x^y$ in \var{K} for an $x$ of \var{H}.\newinvfive
\begin{example}
cout << pow_exponent_rel(interval(8), interval(2), interval::universe()) << ' '
     << pow_exponent_rel(interval(3), interval(2,4), interval::universe());
@outputs^[3] [0.792481250360578, 1.584962500721157]~
\end{example}
  The bounds of both functions are the tightest ones, but in one case, given below. The ends of the
  solutions, $c^{1/y}$ and $\log c/\log x$ for bounds $c$, $x$ and $y$ of
  the arguments, are bounded by the largest and the smallest doubles $d$
  whose powers $d^y$, or $x^d$, the \code{pow} of CORE-MATH, correctly
  rounded, proves to be on either side of $c$. A bound of \var{I} where the
  solutions start between it and the next double is kept when its power is
  within one double of $c$, which only knowing whether that power is a
  double would settle:
  \code{pow\_rel(\itv{1+2^{-52}}{1+2^{-52}}, \itv{0.5}{0.5},
  \itv{0}{1+2^{-51}})} is \itv{1+2^{-51}}{1+2^{-51}}, where
  $(1+2^{-52})^2=1+2^{-51}+2^{-104}$ and no $x$ of \itv{0}{1+2^{-51}} is a
  solution: the tightest result for \var{I} one double wider.
\end{deffun}

\begin{Ex}\label{ex:arithmetic-pow-rel}
The two reverse functions of the power contract in turn the base and the
exponent of the relation $z=x^y$, each from the current domains of the two
others.
\begin{example}
interval x(0, 10), y(2, 3), z(4, 9);
x = pow_rel(z, y, x);            @rem^// the x of x with pow(x, y) in z for a y of y~
y = pow_exponent_rel(z, x, y);   @rem^// the y of y with pow(x, y) in z for an x of x~
cout << x << ' ' << y << endl;   @rem^// [1.587401051968199, 3] [2, 3]~
gaol::cleanup();
\end{example}
The base is contracted to \itv{1.587401051968199}{3}, the cube root of 4
being its least value with a power in \itv{4}{9}; the exponent \itv{2}{3} is
not contracted, every $y$ of it having a base in \itv{1.587401051968199}{3}.
\end{Ex}

\deffunx{interval}{max\_rel}{(const interval\& Z, const interval\& Y, const interval\& X)}
\begin{deffun}{interval}{min\_rel}{(const interval\& Z, const interval\& Y, const interval\& X)}
  \begin{operation}
    \op{$\code{max\_rel}(Z,Y,X)=\hull{\{x\in X\mid\exists y\in Y\colon
      \max(x,y)\in Z\}}$}
    \op{$\code{min\_rel}(Z,Y,X)=\hull{\{x\in X\mid\exists y\in Y\colon
      \min(x,y)\in Z\}}$}
  \end{operation}
  Return the \emph{relational maximum} and \emph{minimum} of \var{Z} given
  \var{Y} with respect to \var{X}, the arguments in the order of
  \code{div\_rel}: the result, the other operand, then the operand
  contracted. $\max$ and $\min$ being symmetric, \code{max\_rel(Z,~X,~Y)}
  contracts \var{Y}, and \code{max(X,~Y)~\&~Z} contracts \var{Z}.
  $\max(x,y)$ is in \var{Z} when $x$ and $y$ are at most $\rightBound{Z}$
  and one of them at least $\leftBound{Z}$: the set is empty when \var{Y}
  lies above \var{Z}, $X\cap\itv{-\infty}{\rightBound{Z}}$ when \var{Y}
  meets \var{Z}, and $X\cap Z$ when \var{Y} lies below \var{Z}, the maximum
  being $x$; \code{min\_rel} is its mirror image. The bounds are those of
  \var{X} and \var{Z}: each result is the tightest enclosure.\newinvfive
\begin{example}
cout << max_rel(interval(1,2), interval(3,4), interval(-5,5)) << ' '
     << max_rel(interval(1,2), interval(0,1), interval(-5,5)) << ' '
     << max_rel(interval(1,2), interval(-1,0), interval(-5,5)) << endl
     << min_rel(interval(1,2), interval(2,3), interval(-5,5)) << endl;
@outputs^[empty] [-5, 2] [1, 2]~
@outputs^[1, 5]~
\end{example}
\end{deffun}

\begin{deffun}{interval}{sign\_rel}{(const interval\& Z, const interval\& X)}
  \begin{operation}
    \op{$\code{sign\_rel}(Z,X)=
      \hull{\{x\in X\mid \operatorname{sign}(x)\in Z\}}$}
  \end{operation}
  Returns the \emph{relational inverse sign} of \var{Z} with respect to
  \var{X}: the negative $x$ of \var{X} where \var{Z} holds $-1$, 0 where
  \var{Z} and \var{X} hold it, and the positive $x$ of \var{X} where \var{Z}
  holds 1; the empty set where \var{Z} holds none of them, as
  \itv{0.5}{0.7}. The negative $x$ go up to 0, which their tightest
  enclosure takes although $\operatorname{sign}(0)=0$, the negative doubles
  next to 0 being in the set.\newinvfive
\begin{example}
cout << sign_rel(interval(-1), interval(-2,3)) << ' '
     << sign_rel(interval(-1), interval(0,3)) << ' '
     << sign_rel(interval(0,1), interval(-2,3)) << endl;
@outputs^[-2, 0] [empty] [0, 3]~
\end{example}
\end{deffun}

\begin{deffun}{interval}{floor\_rel}{(const interval\& Z, const interval\& X)}
  \begin{operation}
    \op{$\code{floor\_rel}(Z,X)=\hull{\{x\in X\mid \lfloor x\rfloor\in Z\}}$}
  \end{operation}
  Returns the \emph{relational inverse floor} of \var{Z} with respect to
  \var{X}: the part of \var{X} in
  $\left[\lceil\leftBound{Z}\rceil,\lfloor\rightBound{Z}\rfloor+1\right[$,
  the $x$ whose floor is an integer of \var{Z}, and the empty set when
  \var{Z} holds no integer. That interval is open at its upper end, which
  its tightest closed enclosure takes all the same, the reals just below it
  being in the set: \code{floor\_rel(interval(0), interval(0.5,~2))} is
  \itv{0.5}{1}, and \code{floor\_rel(interval(0), interval(1,~2))} is empty,
  \var{X} starting at that end. Beyond $2^{53}$, where every double is an
  integer, $\lfloor\rightBound{Z}\rfloor+1$ is no double, and the bound is
  the double above it: \code{floor\_rel} of \itv{2^{53}}{2^{53}} ends at
  $2^{53}+2$, and at $+\infty$ for \itv{\text{DBL\_MAX}}{\text{DBL\_MAX}},
  where \var{X} goes on to it. Each result is the tightest
  enclosure.\newinvfive
\begin{example}
cout << floor_rel(interval(0,2), interval(-5,5)) << ' '
     << floor_rel(interval(0), interval(1,2)) << ' '
     << floor_rel(interval(0.5,0.7), interval(-5,5)) << endl;
@outputs^[0, 3] [empty] [empty]~
\end{example}
\end{deffun}

\deffunx{interval}{atan2\_rel\_y}{(const interval\& \var{C}, const interval\& \var{X}, \\\hbox{}\hfill const interval\& \var{Y})}
\begin{deffun}{interval}{atan2\_rel\_x}{(const interval\& \var{C}, const interval\& \var{Y}, \\\hbox{}\hfill const interval\& \var{X})}
  \begin{operation}
    \op{$\code{atan2\_rel\_y}(C,X,Y)=\hull{\{y\in Y\mid\exists x\in X\colon
      \operatorname{atan2}(y,x)\in C\}}$}
    \op{$\code{atan2\_rel\_x}(C,Y,X)=\hull{\{x\in X\mid\exists y\in Y\colon
      \operatorname{atan2}(y,x)\in C\}}$}
  \end{operation}
  Return the \emph{relational} \code{atan2} with respect to its first and to
  its second argument: the hulls of the ordinates and of the abscissas of
  the points of the box $\var{X}\times\var{Y}$ whose angle is in \var{C},
  the point $(0,0)$, which has no angle, left out.\newinvfive{} As in
  \code{div\_rel}, the interval contracted comes last, and the other
  argument of \code{atan2} before it. The angles are those of \code{atan2},
  in $\left]-\pi,\pi\right]$: the half-line $y=0$, $x<0$ has the angle
  $\pi$, and a \var{C} holding $-\pi$ but not $\pi$ leaves it out; the
  angles of \var{C} beyond $\pi$ are no angle of any point. The bounds are
  within two doubles of the tightest ones.
\begin{example}
interval c(1,2), x(-1,1), y(0,1);
cout << atan2_rel_y(c, x, y) << ' '
     << atan2_rel_x(c, y, x) << endl
     << atan2_rel_x(interval(-4,-3), interval(0.0), x)
     << endl;
@outputs^[0, 1] [-0.4576575543602859, 0.6420926159343309]~
@outputs^[empty]~
\end{example}
\end{deffun}

\subsection{The reverse functions of the standard}
\label{sec:arithmetic-reverse-names}

Table~\ref{tab:arithmetic-reverse} gives the correspondence between the
reverse functions of IEEE 1788-2015 and the relational functions of GAOL.
The last argument $x$ of a function of the standard is
\itv{-\infty}{+\infty} when it is left out.

\begin{table}[htbp]
  \centering\renewcommand{\arraystretch}{1.1}
  \caption{The reverse functions of IEEE 1788-2015 and the relational functions of GAOL} \label{tab:arithmetic-reverse}
  \begin{tabular}{ll}
    \hline
    \multicolumn{1}{c}{IEEE 1788-2015} & \multicolumn{1}{c}{GAOL}\\
    \hline\hline
    \ieee{sqrRev}$(c,x)$ & \code{sqrt\_rel(c, x)}\\
    \ieee{absRev}$(c,x)$ & \code{invabs\_rel(c, x)}\\
    \ieee{pownRev}$(c,x,p)$ & \code{nth\_root\_rel(c, p, x)}\\
    \ieee{sinRev}, \ieee{cosRev}, \ieee{tanRev}$(c,x)$ &
      \code{asin\_rel}, \code{acos\_rel}, \code{atan\_rel(c, x)}\\
    \ieee{coshRev}$(c,x)$ & \code{acosh\_rel(c, x)}\\
    \ieee{mulRev}$(b,c,x)$ & \code{div\_rel(c, b, x)}\\
    \ieee{powRev1}$(b,c,x)$ & \code{pow\_rel(c, b, x)}\newinvfive\\
    \ieee{powRev2}$(a,c,x)$ & \code{pow\_exponent\_rel(c, a, x)}\newinvfive\\
    \ieee{atan2Rev1}$(b,c,x)$ & \code{atan2\_rel\_y(c, b, x)}\newinvfive\\
    \ieee{atan2Rev2}$(a,c,x)$ & \code{atan2\_rel\_x(c, a, x)}\newinvfive\\
    \hline
  \end{tabular}

\end{table}

\begin{Rem}\label{rem:arithmetic-reverse-names}
\ieee{mulRevToPair}$(b,c)$, the two parts of \code{div\_rel(c, b)}, is in
the namespace \code{gaol\_ieee1788} under its name
(Chapter~\ref{chap:ieee1788}).\newinvfive{} \ieee{pownRev} takes every
integer $p$, $p\leq0$ included. \code{max\_rel}, \code{min\_rel},
\code{sign\_rel} and \code{floor\_rel} are the reverse functions of $\max$,
$\min$, $\operatorname{sign}$ and $\operatorname{floor}$ as the standard
defines them (10.5.4), which it neither names nor requires, Table~10.1 not
listing them: the namespace \code{gaol\_ieee1788} has no name for them. A
program that opens that namespace calls them as \code{max\_rel(z, y, x)}
all the same, argument-dependent lookup finding them in
\code{gaol\_core}.\newinvfive
\end{Rem}

\section{Cancellative addition and subtraction}
\label{sec:arithmetic-cancellative}
\cindex{cancellative operations}%

The two functions of this section are the \ieee{cancelMinus} and
\ieee{cancelPlus} of IEEE 1788-2015 (10.5.6): \code{cancel\_minus(X, Y)} is
the tightest interval that, added to \var{Y}, gives an interval containing
\var{X}, when such an interval exists.

\begin{deffun}{interval}{cancel\_minus}{(const interval\& \var{X}, const interval\& \var{Y})}
  The tightest interval $Z$ such that $Y+Z$ contains $X$,
  \itv{\leftBound{X}-\leftBound{Y}}{\rightBound{X}-\rightBound{Y}}, which
  exists when \var{X} and \var{Y} are bounded and \var{X} is at least as wide
  as \var{Y}; \itv{-\infty}{+\infty} otherwise, and the empty set for an
  empty \var{X} and a bounded \var{Y}: \ieee{cancelMinus} of IEEE 1788-2015
  (10.5.6). Whether \var{X} is at least as wide as \var{Y} is decided
  exactly, which rounding the differences of the bounds cannot
  do.\newinvfive
\end{deffun}

\begin{deffun}{interval}{cancel\_plus}{(const interval\& \var{X}, const interval\& \var{Y})}
  \code{cancel\_minus(\var{X}, -\var{Y})}: \ieee{cancelPlus} of IEEE
  1788-2015.\newinvfive
\begin{example}
cout << cancel_minus(interval(1,5), interval(0,2)) << ' '
     << cancel_minus(interval(1,2), interval(0,3)) << ' '
     << cancel_plus(interval(1,5), interval(0,2));
@outputs^[1, 3] [-inf, inf] [3, 5]~
\end{example}
\end{deffun}

\section{Discussion}\label{sec:arithmetic-discussion}

This chapter has given the reference of the operations on intervals. Every
functional operation computes the hull of the image of the part of its
argument within the domain of IEEE 1788-2015, and gives the empty set for an
argument with no point of that domain. The elementary functions take their
values from CORE-MATH, which is correctly rounded: their bounds are the
tightest ones, or within a double or two of them, where GAOL~4 depended on
the mathematical library of the system. The power is hybrid: an integer
exponent, whatever its type, takes the integer power \ieee{pown}, with the
negative bases, and any other exponent the real power \ieee{pow}, at the
price of a result that is not monotone for the inclusion of the exponent;
the powers of the two namespaces are compared in Chapter~\ref{chap:overview}.

\medskip
The relational functions are the reverse functions of the standard with the
arguments in the order of GAOL, the result first and the operand contracted
last. GAOL v5 adds the reverse functions of the power and of \code{atan2},
      which the standard names, and those of the maximum, the minimum, the sign and
      the floor, which it defines but does not name, and gives the empty set wherever no point of the operand
is in the relation.\newinvfive{} How tight each operation is, and the
conditions under which the bounds are the same on every machine, are the
subject of Chapter~\ref{chap:accuracy}; how the bounds are computed, of
Chapter~\ref{chap:computing}.

\medskip
Chapter~\ref{chap:functions} describes the functions on intervals that are
not arithmetic operations: the numeric and the integer functions of an
interval, its splitting, and the union and the intersection of intervals.

%% !TEX root = ../manual_gaolv5.tex
% Chapter 8.
%% Chapter 8: Interval functions. Sources: gaol.tex of f5bf180 lines 4446-4903
%% (the numeric functions, the integer functions, the splitting methods, the
%% union and the intersection). No doc/ page is attached to this chapter.
\chapter{Interval functions}\label{chap:functions}
%%==========================

An interval algorithm does more than evaluate expressions on intervals. A
branch-and-bound method measures the boxes it holds, cuts the widest one in
two parts and keeps the parts that its tests do not reject; an interval
Newton step starts from the midpoint of the current interval and intersects
its result with it; a stopping test compares a width with a threshold. The
operations these steps call do not compute the image of a real function:
they read the bounds of an interval and return a number, they round its
elements to integers, they cut it, or they combine two intervals as sets.
This chapter describes them.

\medskip
The numeric functions of IEEE 1788-2015 (Table~10.2, 12.12.8) are those of
this chapter. The namespace \code{gaol\_ieee1788} names them as the standard
does, and returns $+\infty$ and $-\infty$ for the bounds of the empty set,
where \code{left()} and \code{right()} return NaN
(Chapter~\ref{chap:ieee1788}).

\medskip
The first section covers the numeric functions, which return a double read
or computed from the bounds of an interval, or the interval enclosing such a
double: the radius, the width, the midpoint, the magnitude and the
mignitude, the Hausdorff distance and the function \code{chi}, with the
absolute value, the minimum and the maximum of intervals. The second section
covers the integer functions, which round the elements of an interval to
integers. The third section describes the splitting methods, which cut an
interval in two parts, and the fourth the intersection of two intervals and
the hull of their union. A discussion closes the chapter.

\section{Numeric functions}\label{sec:functions-numeric}
%%-------------------------
\cindex{numeric functions}\cindex{width!of an interval}%
\cindex{midpoint}\cindex{radius}\cindex{magnitude}\cindex{mignitude}%
\cindex{Hausdorff distance}%

The quantities of this section describe an interval by a real number. GAOL
computes each of them on the doubles of the bounds, with the rounding stated
in its entry, and gives two of them, the width and the midpoint, in a second
form that returns the interval enclosing the exact value.

\begin{Def}\label{def:functions-numeric}
Let $I=\itv{a}{b}$ be a non-empty interval. The \dfn{width} of $I$ is $b-a$
and its \dfn{midpoint} is $(a+b)/2$. The \dfn{magnitude} of $I$ is the
greatest absolute value of its elements, and its \dfn{mignitude} the
smallest one \cite[chap.~3]{Hansen:book92}: the mignitude is $0$ when
$0\in I$, $a$ when $I$ is strictly positive, and $-b$ otherwise. The
\dfn{signed mignitude} of $I$ \cite[def.~1.3.28]{Stahl:PhD} is its element
of smallest absolute value: $0$ when $0\in I$, $a$ when $I$ is strictly
positive, and $b$ otherwise.
\end{Def}

\begin{Rem}\label{rem:functions-empty}
The empty set has no width, no midpoint and no element. \code{rad()},
\code{width()}, \code{mig()}, \code{smig()}, \code{mag()},
\code{midpoint()} and \code{chi()} return NaN for it, as \code{left()} and
\code{right()} do, its bounds being NaN; \code{hausdorff()} returns NaN
when one of its arguments is empty; \code{width\_enclosure()},
\code{mid()} and \code{inflate()} return the empty set. GAOL~4 returned
$-1$ for the width of the empty set.\newinvfive
\end{Rem}

\begin{Rem}\label{rem:functions-unbounded}
An unbounded interval has an infinite width and no midpoint in \RSet.
\code{width()} and \code{rad()} return $+\infty$ for it, and
\code{width\_enclosure()} the interval \itv{\text{MAX}}{+\infty}, MAX being
the largest double. \code{midpoint()} returns $0$ for the whole line,
$-\text{MAX}$ for \itv{-\infty}{b} and MAX for \itv{a}{+\infty}: the double
at which the splitting methods cut these intervals (\pxref{sec:splitting}).
\end{Rem}

\begin{defmethod}{interval}{double}{rad}{(void) const}
  \begin{operation}
    \op{$\code{rad}(\itv{a}{b})=\min\{r\in\FSet\mid [m-r,m+r]\supseteq\itv{a}{b}\}$}
  \end{operation}
  Returns the radius of \code{self}: the smallest double $r$ such that
  $[m-r,m+r]$ contains \code{self}, $m$ being \code{self.midpoint()}
  (\ieee{rad} of IEEE 1788-2015). The radius of an unbounded interval is
  $+\infty$, and the radius of the empty set is NaN. The radius of a point
  interval of zero is $+0$, whatever the signs of its bounds, \itv{-0}{-0}
  included, as 12.12.8 requires of every numeric function but
  \ieee{inf}.\newinvfive
\end{defmethod}

\begin{defmethod}{interval}{void}{mid\_rad}{(double\& \var{m}, double\& \var{r}) const}
  Sets \var{m} to \code{self.midpoint()} and \var{r} to \code{self.rad()}
  (\ieee{midRad} of IEEE 1788-2015).\newinvfive
\begin{example}
double m, r;
interval(1,4).mid_rad(m, r);
cout << m << ' ' << r << ' ' << interval(1,GAOL_INFINITY).rad();
@outputs^2.5 1.5 inf~
\end{example}
\end{defmethod}

\begin{defmethod}{interval}{double}{width}{(void) const}
  \begin{operation}
    \op{$\code{width}(\itv{a}{b})=\roundUp{b-a}$}
  \end{operation}
  Returns the width of \code{self}, rounded upward: \ieee{wid} of IEEE
  1788-2015. The width of the empty set is NaN, as the standard has it,
  where GAOL~4 returned $-1$, and that of a point zero is $+0$, whatever the
  signs of its bounds.\newinvfive
\begin{example}
cout << boolalpha << interval(4,6).width() << ' '
     << (interval(1,next_float(1)).width()
         == std::numeric_limits<double>::epsilon()) << ' '
     << interval::emptyset().width();
@outputs^2 true nan~
\end{example}
\end{defmethod}

\begin{defmethod}{interval}{interval}{width\_enclosure}{(void) const}
  \begin{operation}
    \op{$\code{width\_enclosure}(\itv{a}{b})
      =\itv{\roundDn{b-a}}{\roundUp{b-a}}$}
  \end{operation}
  Returns an interval enclosing the exact width of \code{self}, of which
  \code{width()} is the upper bound only. It is
  \code{interval(b) - interval(a)}: the point \itv{0}{0} for a point
  interval, and the empty set for the empty set. The width of an unbounded
  interval, $+\infty$, is no real number; the result is then
  \itv{\text{MAX}}{+\infty}, MAX being the largest double, as
  \code{textToInterval("inf")} is.\newinvfive
\begin{example}
cout << interval(1,3).width_enclosure() << ' '
     << interval(-1,1e16).width_enclosure() << endl
     << interval(1,GAOL_INFINITY).width_enclosure();
@outputs^[2] [1e+16, 1.000000000000001e+16]~
@outputs^[1.797693134862315e+308, inf]~
\end{example}
\end{defmethod}

\begin{defmethod}{interval}{interval}{inflate}{(double \var{r}) const}
  \begin{operation}
    \op{$\code{inflate}(\itv{a}{b},r)=\itv{\roundDn{a-r}}{\roundUp{b+r}}$}
  \end{operation}
  Returns \code{self} widened by the absolute radius \var{r}, that is
  \code{self + interval(-\var{r}, \var{r})}. The result encloses every
  $x+e$ with $x$ in \code{self} and $|e|\leq r$, however small $r$ is next
  to the bounds. $r=0$ gives \code{self}, and $r=+\infty$ the whole line.
  The empty set stays empty, and a negative or NaN \var{r} gives the empty
  set.\newinvfive
\begin{example}
cout << boolalpha << interval(1,2).inflate(0.5) << ' '
     << interval(1,2).inflate(1e-20).set_eq(interval(1,2));
@outputs^[0.5, 2.5] false~
\end{example}
\end{defmethod}

\begin{defmethod}{interval}{double}{mig}{(void) const}
  Returns the mignitude of \code{self}
  (Definition~\ref{def:functions-numeric}): \ieee{mig} of IEEE 1788-2015.
  The mignitude of the empty set is NaN.
\begin{example}
cout << interval(4,5).mig() << ' '
     << interval(-6,-3).mig() << ' '
     << interval(-3,8).mig();
@outputs^4 3 0~
\end{example}
\end{defmethod}

\begin{defmethod}{interval}{double}{smig}{(void) const}
  Returns the signed mignitude of \code{self}
  (Definition~\ref{def:functions-numeric}), after Stahl
  \cite[def.~1.3.28]{Stahl:PhD}. The signed mignitude of \itv{a}{b} is $0$
  when the interval contains $0$, $a$ when it is strictly positive, and $b$
  otherwise. The signed mignitude of the empty set is NaN.
\begin{example}
cout << interval(4,5).smig() << ' '
     << interval(-6,-3).smig() << ' '
     << interval(-3,8).smig();
@outputs^4 -3 0~
\end{example}
\end{defmethod}

\begin{defmethod}{interval}{double}{mag}{(void) const}
  Returns the magnitude of \code{self}: \ieee{mag} of IEEE 1788-2015. The
  magnitude of \itv{a}{b} is the greatest absolute value of its elements.
  The magnitude of the empty set is NaN.
\begin{example}
cout << interval(4,5).mag() << ' '
     << interval(-6,-3).mag() << ' '
     << interval(-10,5).mag();
@outputs^5 6 10~
\end{example}
\end{defmethod}

\begin{deffun}{double}{hausdorff}{(const interval\& \var{I1}, const interval\& \var{I2})}
  Returns the Hausdorff distance between the sets \var{I1} and \var{I2},
\begin{equation*}
\code{hausdorff}(I_1,I_2)=\max(|\leftBound{I_1}-\leftBound{I_2}|,
                          |\rightBound{I_1}-\rightBound{I_2}|),
\end{equation*}
  rounded upward: the result is the tightest upper bound of the
  distance.\newinvfive

  Two equal bounds are at distance $0$, infinite ones included: the
  distance of \itv{1}{+\infty} to itself is $0$, and to \itv{2}{+\infty} it
  is $1$, where GAOL~4 returned $+\infty$ for both. The distance is
  $+\infty$ when a bound is infinite in one of the two intervals only, and
  NaN when one of them is empty.\newinvfive
\begin{example}
cout << hausdorff(interval(4,8),interval(5,10)) << ' '
     << hausdorff(interval(1,GAOL_INFINITY),interval(2,GAOL_INFINITY)) << ' '
     << hausdorff(interval(1,GAOL_INFINITY),interval(1,2));
@outputs^2 1 inf~
\end{example}
\end{deffun}

\begin{defmethod}{interval}{double}{midpoint}{(void) const}
  Returns the midpoint of \code{self}: \ieee{mid} of IEEE 1788-2015. For
  $a$ and $b$ two finite doubles, and MAX the largest positive double,
  \code{std::numeric\_limits<double>::max()}, the cases are:
\begin{equation*}
\left\{\begin{array}{ll}
\code{midpoint}(\emptyset{})&=\text{NaN}\\
\code{midpoint}(\itv{-\infty}{+\infty})&=0\\
\code{midpoint}(\itv{-\infty}{b})&=-\text{MAX}\\
\code{midpoint}(\itv{a}{+\infty})&=\text{MAX}\\
\code{midpoint}(\itv{a}{b})&=\roundNearest{(a+b)/2} \\
\end{array}\right.
\end{equation*}
  The midpoint of \itv{a}{b} is rounded to nearest, ties to even. It is the
  half of $a+b$, or the sum of the halves of $a$ and $b$ where $a+b$ may
  overflow \cite{Goualard:TOMS2014}, that is where $|a|$ or $|b|$ is
  $2^{1023}$ or more. This case is tested first, so that \code{midpoint()}
  of \itv{\text{MAX}}{\text{MAX}} raises no overflow
  exception.\newinvfive
\end{defmethod}

\begin{defmethod}{interval}{interval}{mid}{(void) const}
  Returns the tightest interval enclosing the midpoint of \code{self},
  which is included in \code{self}. With the notations of
  \code{midpoint()}, the cases are:
\begin{equation*}
\left\{\begin{array}{ll}
\code{mid}(\emptyset{})&=\emptyset{}\\
\code{mid}(\itv{-\infty}{+\infty})&=\itv{0}{0}\\
\code{mid}(\itv{-\infty}{b})&=\itv{-\text{MAX}}{-\text{MAX}}\\
\code{mid}(\itv{a}{+\infty})&=\itv{\text{MAX}}{\text{MAX}}\\
\code{mid}(\itv{a}{b})&=\itv{\roundDn{(a+b)/2}}{\roundUp{(a+b)/2}}
\end{array}\right.
\end{equation*}
  The midpoint is computed as \code{midpoint()} computes it, with the
  rounding outward. GAOL~4 did not enclose the midpoint of some intervals
  whose bounds are subnormal numbers.\newinvfive
\end{defmethod}

\begin{defmethod}{interval}{double}{left}{(void) const}
  \begin{operation}
    \op{$\code{left}(\itv{x}{y})=x$}
  \end{operation}
  Returns the left bound of \code{self}: NaN for the empty set, whose
  bounds are NaN. \code{gaol\_ieee1788::inf()} returns $+\infty$ there, as
  IEEE 1788-2015 has it.
\end{defmethod}

\begin{defmethod}{interval}{double}{right}{(void) const}
  \begin{operation}
    \op{$\code{right}(\itv{x}{y})=y$}
  \end{operation}
  Returns the right bound of \code{self}: NaN for the empty set.
  \code{gaol\_ieee1788::sup()} returns $-\infty$ there.
\end{defmethod}

\begin{deffun}{interval}{abs}{(const interval\& \var{I})}
  \begin{operation}
    \op{$\text{abs}(I)=\{|x|\in\RSet^+\mid x\in I\}$}
  \end{operation}
  Returns the absolute value of \var{I}.
  \begin{example}
cout << abs(interval(-5,6)) << ' '
     << abs(interval(-4,-2));
@outputs^[0, 6] [2, 4]~
  \end{example}
\end{deffun}

\begin{deffun}{double}{chi}{(const interval\& \var{I})}
  Returns the function \code{chi} of Ratschek and Rokne
  \cite{Ratschek-Rokne:1995}, which characterizes the degree of symmetry
  of an interval:
\begin{equation*}
\text{for } I=\itv{a}{b},\quad\text{chi}(I)=\left\{\begin{array}{ll}
    -1  & \text{if } I = \itv{0}{0}\\
    a/b & \text{if } |a|\leq |b|\\
    b/a & \text{otherwise.}
  \end{array}\right.
\end{equation*}
  It is NaN for the empty set, $1$ for \itv{-\infty}{+\infty}, and $0$ for
  the other intervals with an infinite bound.
\begin{example}
cout << chi(interval(3,6)) << ' '
     << chi(interval(-6,3)) << ' '
     << chi(interval::emptyset()) << ' '
     << chi(interval::universe()) << ' '
     << chi(interval(-5,GAOL_INFINITY));
@outputs^0.5 -0.5 nan 1 0~
\end{example}
\end{deffun}

\begin{deffun}{interval}{min}{(const interval\& \var{I}, const interval\& \var{J})}
  \begin{operation}
    \op{$\min(\itv{a}{b}, \itv{c}{d}) = \itv{\min(a,c)}{\min(b,d)}$}
  \end{operation}
  Returns the minimum of two intervals.
  \begin{example}
cout << min(interval(5,6),interval(3,9)) << ' '
     << min(interval::emptyset(),interval(3,8));
@outputs^[3, 6] [empty]~
  \end{example}
\end{deffun}

\begin{deffun}{interval}{max}{(const interval\& \var{I}, const interval\& \var{J})}
  \begin{operation}
    \op{$\max(\itv{a}{b}, \itv{c}{d}) = \itv{\max(a,c)}{\max(b,d)}$}
  \end{operation}
  Returns the maximum of two intervals.
\begin{example}
cout << max(interval(5,6),interval(3,9)) << ' '
     << max(interval::emptyset(),interval(3,8));
@outputs^[5, 9] [empty]~
\end{example}
\end{deffun}

\section{Integer functions}\label{sec:functions-integer}
%%-------------------------
\cindex{integer functions}\cindex{rounding!to an integer}%

The functions of this section round the elements of an interval to
integers. They are those of IEEE 1788-2015 (Table~9.1), but \code{integer()},
which is GAOL's own; their names in the
standard are given with the operations of Chapter~\ref{chap:ieee1788}.

\begin{Prop}\label{prop:functions-integer}
Each function of this section but \code{integer()} is non-decreasing, so that
      the bounds of its result are its values at the bounds of the interval; those
      of \code{integer()} are the \code{ceil} of the left bound and the \code{floor}
      of the right one. Each is exact.
None of them depends on the rounding direction.
\end{Prop}

\begin{deffun}{interval}{floor}{(const interval\& \var{I})}
  \begin{operation}
    \op{$\itv{\code{floor}(I.\code{left}())}{\code{floor}(I.\code{right}())}$}
  \end{operation}
  The elements of \var{I} rounded down to an integer.
  \begin{example}
cout << floor(interval(4.5,6.5)) << ' '
     << floor(textToInterval("[-10.4,3.5]"));
@outputs^[4, 6] [-11, 3]~
  \end{example}
\end{deffun}

\begin{deffun}{interval}{ceil}{(const interval\& \var{I})}
  \begin{operation}
    \op{\itv{\code{ceil}(I.\code{left}())}{\code{ceil}(I.\code{right}())}}
  \end{operation}
  The elements of \var{I} rounded up to an integer.
  \begin{example}
cout << ceil(interval(4.5,6.5)) << ' '
     << ceil(textToInterval("[-10.4,3.5]"));
@outputs^[5, 7] [-10, 4]~
  \end{example}
\end{deffun}

\begin{deffun}{interval}{integer}{(const interval\& \var{I})}
  \begin{operation}
    \op{\itv{\code{ceil}(I.\code{left}())}{\code{floor}(I.\code{right}())}}
  \end{operation}
  Narrows the bounds down to the closest integers: the empty set when
  \var{I} holds no integer. The bounds of the result are still
  \code{double} numbers, which an integral type may not represent.
  \begin{example}
cout << integer(interval(4.5,6.5)) << ' '
     << integer(interval(4.2,4.8));
@outputs^[5, 6] [empty]~
  \end{example}
\end{deffun}

\begin{deffun}{interval}{sign}{(const interval\& \var{I})}
  The signs of the elements of \var{I}: $-1$ below $0$, $0$ at $0$, $1$
  above (\ieee{sign}).\newinvfive
\end{deffun}

\begin{deffun}{interval}{trunc}{(const interval\& \var{I})}
  The elements of \var{I} rounded toward zero (\ieee{trunc}).\newinvfive
\end{deffun}

\deffunx{interval}{round\_ties\_to\_even}{(const interval\& \var{I})}
\begin{deffun}{interval}{round\_ties\_to\_away}{(const interval\& \var{I})}
  The elements of \var{I} rounded to the nearest integer, halfway values to
  the even one, respectively away from zero (\ieee{roundTiesToEven},
  \ieee{roundTiesToAway}).\newinvfive
  \begin{example}
cout << sign(interval(-2,3)) << ' '
     << trunc(interval(-2.5,3.7)) << ' '
     << round_ties_to_even(interval(0.5,2.5)) << ' '
     << round_ties_to_away(interval(0.5,2.5));
@outputs^[-1, 1] [-2, 3] [0, 2] [1, 3]~
  \end{example}
\end{deffun}

\section{Splitting methods}\label{sec:splitting}
%%-------------------------
\cindex{splitting}\cindex{bisection}\cindex{midpoint!cut at the}%

A branch-and-bound algorithm cuts the intervals it cannot discard. GAOL cuts
an interval \itv{a}{b} at a double $c$ of the interval that both parts
share: the parts \itv{a}{c} and \itv{c}{b} cover the interval exactly.
\code{split()}, \code{split\_left()} and \code{split\_right()} cut at
\code{midpoint()}; \code{bisect()} cuts at a ratio of the width chosen by
the caller; \code{is\_bisectable()} tells whether a cut gives two parts
smaller than the interval.

\begin{Def}\label{def:functions-bisectable}
An interval is \dfn{bisectable} when a double lies strictly between its
bounds; the two parts of a cut at such a double are then both smaller than
the interval. The empty set, a point interval and an interval of two
consecutive doubles are not bisectable, nor are the half-lines
\itv{\text{MAX}}{+\infty} and \itv{-\infty}{-\text{MAX}}, MAX being the
largest double.
\end{Def}

\begin{Rem}\label{rem:functions-unbounded-cut}
A cut point is a double of the interval. An unbounded interval is therefore
cut at the largest double or at its opposite, whatever the ratio, and the
whole line at $0$; the half \itv{\text{MAX}}{+\infty} keeps an infinite
width, so that a loop which cuts an interval until its width falls under a
threshold never ends on it. Such a loop stops on \code{is\_bisectable()},
which is \code{false} for this half (\pxref{sec:common-errors}).\newinvfive
\end{Rem}

\begin{defmethod}{interval}{void}{split}{(interval\& \var{I1}, interval\& \var{I2}) const}
  Splits \code{self} into two parts at \code{midpoint()}, and returns the
  left part in \var{I1} and the right part in \var{I2}. \var{I1} or
  \var{I2} may be \code{self} itself.
  \begin{example}
interval I1a, I2a,
         I3(1.0,next_float(1.0)),
         I1b, I2b;
interval(4,5).split(I1a,I2a);
I3.split(I1b,I2b);
cout << boolalpha << I1a << ' ' << I2a << ' '
     << I1b.set_eq(interval(1.0)) << ' ' << I2b.set_eq(I3) << endl;
@outputs^[4, 4.5] [4.5, 5] true true~
  \end{example}
\end{defmethod}

\begin{defmethod}{interval}{interval}{split\_left}{(void) const}
  \begin{operation}
    \op{$\code{split\_left}(\itv{a}{b})=\itv{a}{m}$}
  \end{operation}
  Splits \code{self} into two parts at $m=\code{midpoint()}$, and returns
  the left part.
  \begin{example}
cout << interval(4,5).split_left();
@outputs^[4, 4.5]~
  \end{example}
\end{defmethod}

\begin{defmethod}{interval}{interval}{split\_right}{(void) const}
  \begin{operation}
    \op{$\code{split\_right}(\itv{a}{b})=\itv{m}{b}$}
  \end{operation}
  Splits \code{self} into two parts at $m=\code{midpoint()}$, and returns
  the right part.
  \begin{example}
cout << interval(4,5).split_right();
@outputs^[4.5, 5]~
  \end{example}
\end{defmethod}

\begin{defmethod}{interval}{std::pair<interval, interval>}{bisect}{(double \var{ratio}) const}
  \begin{operation}
    \op{$\code{bisect}(\itv{a}{b},\rho)=(\itv{a}{c},\itv{c}{b})$}
  \end{operation}
  Returns the two parts of \code{self} on either side of the cut point $c$,
  a double of \itv{a}{b} that both parts share, so that they cover
  \code{self} exactly. For $0<\rho<1$, $c$ is $a+\rho(b-a)$, computed as
  $a(1-\rho)+b\rho$ or, where a bound is $2^{1023}$ or more in magnitude, as
  twice $\frac{a}{2}(1-\rho)+\frac{b}{2}\rho$ kept between $\frac{a}{2}$ and
  $\frac{b}{2}$, so that no overflow exception is raised for finite bounds, and
  kept within \itv{a}{b};
  when \code{is\_bisectable()} is true, $c$ is moreover strictly between
  $a$ and $b$, so that both parts are smaller than \code{self}. $\rho=0.5$
  cuts at \code{midpoint()}, as \code{split()} does, and so does every
  ratio for an unbounded interval: a half-line is cut at the largest double
  or at its opposite, the whole line at $0$. The empty set gives two empty
  sets, and a point interval itself twice. A ratio outside $(0,1)$, or NaN,
  throws \code{std::domain\_error} (\pxref{sec:domain_error}).\newinvfive
  \begin{example}
std::pair<interval, interval> p = interval(0,4).bisect(0.25);
cout << p.first << ' ' << p.second;
@outputs^[0, 1] [1, 4]~
  \end{example}
\end{defmethod}

\begin{defmethod}{interval}{bool}{is\_bisectable}{(void) const}
  Returns \code{true} when a double lies strictly between the bounds of
  \code{self} (Definition~\ref{def:functions-bisectable}), so that
  \code{bisect()} gives two parts, each smaller than \code{self}. Returns
  \code{false} for the empty set, a point interval, two consecutive
  doubles, and the half-lines \itv{\text{MAX}}{+\infty} and
  \itv{-\infty}{-\text{MAX}}, which \code{split()} cuts at their finite
  bound.\newinvfive
  \begin{example}
cout << boolalpha << interval(1,2).is_bisectable() << ' '
     << interval(1,next_float(1)).is_bisectable();
@outputs^true false~
  \end{example}
\end{defmethod}

\section{Union and intersection}\label{sec:functions-union}
%%------------------------------
\cindex{intersection}\cindex{hull!of a union}\cindex{union of intervals}%

Intervals are sets, and two of them can be intersected or joined. The
intersection of two intervals is an interval, possibly empty; their union
is not always one, and GAOL returns its hull, the smallest interval that
contains both. The operators \code{\&} and \code{|} compute them, and the
functions \code{intersect()} and \code{hull()} give them a name.

\begin{Def}\label{def:functions-hull}
Let $I_1=\itv{a}{b}$ and $I_2=\itv{c}{d}$ be two non-empty intervals. Their
\dfn{intersection} $I_1\cap I_2$ is the empty set when $b<c$ or $d<a$, and
\itv{\max(a,c)}{\min(b,d)} otherwise. The \dfn{hull} of their union,
$\hull{(I_1\cup I_2)}$, is \itv{\min(a,c)}{\max(b,d)}.
\end{Def}

\begin{Rem}\label{rem:functions-intersection}
GAOL~4 kept the greatest left bound and the smallest right bound of two
disjoint intervals: $\itv{1}{2}\cap\itv{3}{4}$ was \itv{3}{2}, which the
operations computing on the bounds did not take for empty. GAOL v5 returns
the empty set, \code{interval::emptyset()}, whose bounds are
NaN.\newinvfive
\end{Rem}

\begin{deffun}{interval}{operator\&}{(const interval\& I1, const interval\& I2)}
  \begin{operation}
    \op{$\var{I1}\cap\var{I2}$}
  \end{operation}
  Returns the intersection of \var{I1} and \var{I2}: \ieee{intersection} of
  IEEE 1788-2015. The intersection of disjoint intervals is the empty set
  \code{interval::emptyset()} (Remark~\ref{rem:functions-intersection}).
  \begin{example}
cout << (interval(4,6) & interval(5,9)) << ' '
     << (interval(1,2) & interval(3,4));
@outputs^[5, 6] [empty]~
  \end{example}
\end{deffun}

\begin{deffun}{interval}{operator\textbar}{(const interval\& I1, const interval\& I2)}
  \begin{operation}
    \op{$\hull{(\var{I1}\cup\var{I2})}$}
  \end{operation}
  Returns the hull of the union of \var{I1} and \var{I2}: \ieee{convexHull}
  of IEEE 1788-2015.
  \begin{example}
cout << (interval(3,6) | interval(9,12));
@outputs^[3, 12]~
  \end{example}
\end{deffun}

\deffunx{interval}{hull}{(const interval\& I1, const interval\& I2)}
\begin{deffun}{interval}{intersect}{(const interval\& I1, const interval\& I2)}
  \code{I1 | I2} and \code{I1 \& I2} under a name, for an algorithm written
  with names rather than operators, and for \code{std::accumulate} or a
  pointer to a function. They are in the namespace \code{gaol}
  too.\newinvfive
  \begin{example}
cout << hull(interval(1,2), interval(4,5)) << ' '
     << intersect(interval(1,3), interval(2,5));
@outputs^[1, 5] [2, 3]~
  \end{example}
\end{deffun}

\section{Discussion}
%%-----------------

This chapter has described the functions that read an interval rather than
map it: the measures of an interval, the roundings of its elements to
integers, its cuts, and the set operations on two intervals. Two
conventions run through them. The empty set, whose bounds are NaN, gives
NaN to every function that returns a double and itself to every function
that returns an interval, but the hull of a union, \code{hull(interval::emptyset(), x)}
being \code{x}; GAOL~4 returned $-1$ for the width of the empty set and
\itv{3}{2} for the intersection of \itv{1}{2} and \itv{3}{4}, values that
the operations on the bounds did not take for empty. An unbounded interval
gives $+\infty$ to its width and its radius, and the largest double, or its
opposite, or $0$ for the whole line, to its midpoint and to its cut points,
which have to be doubles of the interval.

\medskip
The roundings follow the containment property (\pxref{sec:containment}).
The width is rounded upward, so that it bounds the exact width from above,
and \code{width\_enclosure()} encloses the exact width when its lower bound
matters; the midpoint is rounded to nearest, with the sum of the halves
where the sum of the bounds would overflow, and \code{mid()} encloses it;
\code{inflate()} rounds outward; the cut point of \code{bisect()} is
computed so that it cannot overflow and stays within the interval. The
integer functions are exact in every rounding direction. The union of two
intervals is given by its hull, the standard's \ieee{convexHull}, which
adds the reals between two disjoint intervals to the result: the
intersection is exact, the union is enclosed.

\medskip
The examples of this chapter print their intervals with the operator
\code{<{}<}, as \code{[4, 4.5]} or \code{[empty]}. The next chapter
describes how GAOL reads and writes intervals, with the formats it accepts
and the number of digits it writes.

%% !TEX root = ../manual_gaolv5.tex
% Chapter 9.
%% Chapter 9: Input and output. Sources: gaol.tex of f5bf180 lines 4904-5531
%% (reading intervals, the input format, writing intervals, converting to
%% strings, the output format, the number of digits, the example); the brief
%% named no doc/ page for this chapter.
%%
%% Copyright (c) 2001 Swiss Federal Institute of Technology, Switzerland
%% Copyright (c) 2002-2016 Laboratoire d'Informatique de
%%                         Nantes-Atlantique, France
%% Copyright (c) 2017-2020 Laboratoire des Sciences du Numérique de Nantes,
%%                         France
%% Copyright (c) 2026 ENSTA, France
%%
%% By: Frederic Goualard, Jordan NININ
%%
%% GAOL is a software distributed WITHOUT ANY WARRANTY. Read the associated
%% COPYING.LIB file for information.
\chapter{Input and output}\label{chap:io}

An interval enters a program and leaves it as text: a literal in a source
file, a line of a data file, a string typed by the user, a result printed on
the screen. A decimal number is seldom a double, and a double is seldom
written exactly with a few decimal digits. Reading and writing must therefore
preserve the containment property (\pxref{sec:containment}): the text read
gives an interval that contains the number or the set it writes, and the text
written for an interval in the formats \element{bounds} and \element{hexa}, or
by \ieee{intervalToText}, read back, gives an interval that contains it. GAOL
v5 does both, whatever the C library and the locale of the
program.\newinvfive

\medskip
In this chapter, we describe how GAOL v5 reads and writes intervals. The
first section covers the reading: the operator \code{>{}>} of the input
streams, then the language of the strings that it and \code{textToInterval()}
accept, with its numbers, the uncertain form of IEEE 1788-2015, its functions,
its bounds and infinities, the limits of the reader and its errors. The second
section covers the writing: the operator \code{<{}<}, the exact text of
\code{exact\_string()}, the conversion to \code{std::string}, the five output
formats, the number of digits, and a complete program. The functions that
read and write with the names of the standard, \ieee{textToInterval},
\ieee{intervalToText}, \ieee{intervalToExact} and \ieee{exactToInterval}, belong to the namespace
\code{gaol\_ieee1788} and are described in Chapter~\ref{chap:ieee1788}
(\pxref{sec:ieee1788-text}).

\section{Reading intervals}\label{sec:reading-intervals}

An interval is read from a stream with the operator \code{>{}>}, as a number
is, and built from a string with \code{textToInterval()}
(Chapter~\ref{chap:creation}). Both read the language of
Section~\ref{sec:input-format}.
\cindex{interval!reading}%

\begin{deffun}{istream\&}{operator>{}>}{(istream\& \var{in}, interval\& \var{I})}
  Reads a line from the input stream \var{in}, and assigns to \var{I} the
  interval it writes. The blanks before the line, the ends of lines
  included, are skipped first, as before a number, even under
  \code{std::noskipws}, which keeps a number from skipping them: an empty
  line, or one holding only blanks, is never a line to read. An interval
  cannot be written on several lines.

  If the line is syntactically ill-formed, \var{I} is set to the empty set.
  A \code{std::invalid\_argument} exception is then thrown
  (\pxref{sec:exceptions}) if the library was compiled with exceptions
  enabled (\pxref{sec:install-developers}); if they were disabled, an error
  message is printed to \code{cerr} and the program aborts. An ill-formed
  line also sets the \code{failbit} of \var{in}, which \code{in.clear()}
  resets before reading on. A second argument of \code{nth\_root} that is not an integer, as in
  \code{nth\_root(8,1.5)}, throws \code{std::domain\_error} instead; a call with
  a wrong number of arguments throws \code{std::invalid\_argument}, and an
  argument outside the domain of a function gives the empty set, as
  \code{sqrt(-1)} does.

  Where no line is left, \var{in} gets its \code{failbit}, \var{I} is left
  as it was, and nothing is thrown, as for a number: \code{while (in >{}> x)}
  reads the lines one after the other, the blank ones skipped, and ends with
  the input, a file ending with an empty line included.\newinvfive{} GAOL~4
  read the empty text left, and threw \code{input\_format\_error}, as it did
  on an empty line, so that \code{in >{}> d >{}> x} over the lines
  \code{1.5} and \code{[1, 2]} threw at the end of the first one.
\begin{example}
std::istringstream in("[1, 2]\n[0.1, 0.2]\n");
interval a, b;
in >> a >> b;
cout << a << ' ' << b;
@outputs^[1, 2] [0.09999999999999999, 0.2000000000000001]~
\end{example}
\begin{example}
std::istringstream in("1.5\n\n[1, 2]\n\n");
double d;
interval x;
in >> d >> x;
cout << d << ' ' << x << ' ';
cout << (in >> x ? "more" : "end");
@outputs^1.5 [1, 2] end~
\end{example}
\end{deffun}

\subsection{The input format}\label{sec:input-format}

A string to be translated into an interval has the syntax below, with the
terminals in lower case and the non-terminals in italic upper case.
\cindex{input format}%

\begin{Verbatim}[commandchars=\@\^\~,fontsize=\small]
@NT^expr~
	: @NT^number~
	| @NT^uncertain_number~ @rem^@hfill// An interval literal, never a bound by itself~
	| @NT^literal~
	| @TM^dmin~  @rem^@hfill// Smallest positive normal floating-point number~
	| @TM^dmax~ @rem^@hfill// Largest positive floating-point number~
	| @TM^pi~
	| @TM^inf~ @rem^@hfill// [MAX, +infinity] as an expression, +infinity as a bound~
	| @TM^infinity~ @rem^@hfill// The same~
	| @TM^empty~ @rem^@hfill// Empty interval~
	| @NT^expr~ @TM^+~ @NT^expr~
	| @NT^expr~ @TM^-~ @NT^expr~
	| @NT^expr~ @TM^*~ @NT^expr~
	| @NT^expr~ @TM^/~ @NT^expr~
	| @TM^-~ @NT^expr~
	| @TM^+~ @NT^expr~
	| @NT^function_call~
	| @TM^(~ @NT^expr~ @TM^)~
	;

@NT^literal~
	: [ @NT^expr~ ]
	| [ @NT^expr~ @TM^,~ @NT^expr~ ]
	| [ @NT^expr~ @TM^,~ ] @rem^@hfill// Right bound +@ensuremath^@infty~~
	| [ @TM^,~ @NT^expr~ ] @rem^@hfill// Left bound -@ensuremath^@infty~~
	| [ @TM^,~ ] @rem^@hfill// Entire real line~
	| [ @TM^entire~ ] @rem^@hfill// Entire real line~
	| [ ] @rem^@hfill// Empty interval~
	| < @NT^expr~ @TM^,~ @NT^expr~ >
	;

@NT^function_call~
	: @NT^name~ @TM^(~ @NT^expr~ @TM^)~
	| @NT^name~ @TM^(~ @NT^expr~ @TM^,~ @NT^expr~ @TM^)~
	| @NT^name~ @TM^(~ @NT^expr~ @TM^,~ @NT^expr~ @TM^,~ @NT^expr~ @TM^)~
	;
\end{Verbatim}

\noindent
Every string is an expression, computed with interval arithmetic. The
operators \TM{*} and \TM{/} bind tighter than \TM{+} and \TM{-}, the four
associate to the left, and a sign binds tighter than all of them:
\code{"1-2-3"} is $-4$ and \code{"2+3*4"} is 14. The
literals of intervals, the numbers, the constants and the functions combine
with the operators in any order, and a literal stands wherever a number may,
in a bound too.\newinvfive{} \code{1+[1,2]}, \code{2*cos([0,1])} and
\code{[cos([0,1]), 2]} are read. \code{[empty]} is the form
\code{[\NT{expr}]} of the expression \code{empty}.

\begin{Rem}
GAOL~4 read a string with two grammars, one for the intervals between
brackets and one for the numbers. It refused \code{1+[1,2]} and
\code{[1,2]+1+[1,2]}, while it read \code{[1,2]*2}. GAOL v5 reads every
string with the one grammar above.\newinvfive
\end{Rem}

Spaces, tabs and newlines are not significant except in numbers, and the
case of letters is ignored, as IEEE 1788-2015 has it for its literals:
\code{[Empty]}, \code{[1, Inf]}, \code{SIN(1)} and \code{1E3} are read,
whatever the locale of the program.\newinvfive{} The sign ``$+$'' before a
number and before \code{inf} is optional.

\paragraph{Numbers.} A \NT{number} is written in decimal, with an optional
point and an optional exponent (\code{12}, \code{0.1}, \code{1e-3}), in
hexadecimal (\code{0x1f}, \code{0x1.3p-1}), in octal (\code{0o17}) or in
binary (\code{0b101}, \code{0b1.1p3}); the digits after \code{0x}, \code{0o}
or \code{0b} have an optional point and an optional binary exponent after a
\code{p}.\newinvfive{} A number that begins with 0 is decimal: \code{010} is
10, as IEEE 1788-2015 has it (9.7.2), and octal is written \code{0o10}, as
in Python.
\cindex{number!reading}%

Numbers are read exactly, whatever their number of digits: a number that is
a double is read as that double, and any other one is replaced by the
tightest interval of doubles containing it.\newinvfive{} GAOL~4 read them
with the \code{strtod()} of the C library, which did not enclose \code{0.1}
where it rounds to nearest in every rounding direction (the C runtime of
Windows, musl on 64-bit ARM processors). The point of a number is a point
whatever the locale of the program, a locale writing a decimal comma
included (as \code{setlocale(LC\_ALL, "")} sets under a French
one).\newinvfive{} GAOL~4 read \code{0.1} as 0 there, \code{strtod()}
stopping at the point. The modes of the processor that flush the subnormal
numbers to zero, flush-to-zero and denormals-are-zero, which a program
linked with \option{-Ofast} has (\pxref{sec:ftz}), do not change what a
number is read as: \code{1e-310} is read as the tightest interval of doubles
containing it, and \code{0} as 0.\newinvfive{} Reading a long number takes a
time quadratic in its number of digits: \code{"0."} followed by $n$ digits
\code{3} takes 0.011~s for $n = 20\,000$, 0.18~s for 80\,000 and 2.8~s for
320\,000 (about 28~s expected for a million), and the uncertain form 2.2 to 2.4
times as long (Intel i7-1185G7, GCC~15.2, the best of five runs,
7 November 2026).\newinvfive

\paragraph{The uncertain form.} An \NT{uncertain\_number} is written as
IEEE 1788-2015 defines it (9.7.4 and 12.11.3).
\cindex{uncertain form}%

\begin{Def}
An \dfn{uncertain number} is a decimal number followed by \code{?}, a radius
in units of its last digit (half a unit when it is left out, and an infinite
radius for a second \code{?}), an optional direction, \code{u} or \code{d},
which keeps only the radius above or below, and an optional exponent.
\end{Def}

\noindent
\code{3.56?1} is \itv{3.55}{3.57}, \code{3.56?1e2} is \itv{355}{357}, and
\code{-10??u} is \itv{-10}{+\infty}.\newinvfive{} An uncertain number is an
interval literal of its own, not a bound by itself: \code{"[5?1]"} and
\code{"[1, 3.56?1]"} throw \code{std::invalid\_argument}, as IEEE 1788-2015
has it (12.11.4), while \code{"3.56?1*2"} is \itv{7.1}{7.14}, and a bound
computed from one is read, \code{"[2*5?1]"} being \itv{8}{12}. A sign written
next to an uncertain number, as \code{10-3.56?1} or \code{10+3.56?1}, is
read as a part of the number, and the text throws
\code{std::invalid\_argument}, as does \code{10 -3.56?1}: a blank is needed
after the sign, \mbox{\code{10 - 3.56?1}}.\newinvfive

\paragraph{Functions.} A \NT{name} is the name of a function of GAOL, with
its arguments, whatever the case of its letters. The names of GAOL are read
by \code{gaol::textToInterval()}, from one string or from two, and by
\code{operator>{}>}; the names of IEEE 1788-2015 are read by
\code{gaol\_ieee1788::textToInterval()} (\pxref{sec:ieee1788-text}), as
each namespace has its own \code{pow}.\newinvfive{} The names of GAOL are
the following.

\begin{center}\small
\begin{tabular}{L{11cm}}
\code{sqr}, \code{sqrt}, \code{cbrt}, \code{pow}, \code{nth\_root},
\code{exp}, \code{exp2}, \code{exp10}, \code{expm1}, \code{exp2m1},
\code{exp10m1}, \code{log}, \code{log2}, \code{log10}, \code{log1p},
\code{log2p1}, \code{log10p1}, \code{sin}, \code{cos}, \code{tan},
\code{asin}, \code{acos}, \code{atan}, \code{atan2}, \code{sinh},
\code{cosh}, \code{tanh}, \code{asinh}, \code{acosh}, \code{atanh},
\code{sinpi}, \code{cospi}, \code{tanpi}, \code{asinpi}, \code{acospi},
\code{atanpi}, \code{atan2pi}, \code{hypot}, \code{compoundm1},
\code{rsqrt}, \code{sign}, \code{ceil}, \code{floor}, \code{trunc},
\code{integer}, \code{round\_ties\_to\_even}, \code{round\_ties\_to\_away},
\code{abs}, \code{min}, \code{max}, \code{inverse}, \code{erf}, \code{erfc},
\code{cancel\_minus}, \code{cancel\_plus}, \code{hull}, \code{intersect},
\code{fma}\\
\end{tabular}
\end{center}

\noindent
Each is the function of GAOL of that name (Chapters~\ref{chap:arithmetic}
and~\ref{chap:functions}); \code{pow} is the \code{pow} of the namespace
\code{gaol} (\pxref{sec:powers}), and \code{cbrt(x)} is
\code{nth\_root(x,3)}. A name that is not in the list, or a call with a
wrong number of arguments, throws \code{std::invalid\_argument}. The second
argument of \code{nth\_root} shall be a point interval that can be evaluated
as an integer, negative ones included; otherwise, a \code{std::domain\_error}
exception is thrown, or an error is reported, depending on whether the
exceptions were enabled at configuration time
(\pxref{sec:install-developers}).\newinvfive{} The integer may be of any size a
double holds: \code{"nth\_root(8,3000000000)"} is
\code{nth\_root(8, 3000000000LL)}, where it was refused beyond the
\code{int}s. One of $2^{63}$ or more in magnitude is taken as
$\pm(2^{63}-2)$: the integer powers are the same, and the enclosure of the
roots of that order holds those of the larger orders. The second argument of
\code{pown} and \code{rootn}, with the names of the standard, is read in the
same way. GAOL~4 read the trigonometric and hyperbolic functions,
\code{exp}, \code{log}, \code{pow}, \code{sqrt} and \code{nth\_root} only.

\paragraph{Bounds.} The notations ``$n$'' and ``[$n$]'' are equivalent, but for
\code{inf} and \code{-inf}, whose literals \code{[inf]} and \code{[-inf]} are
empty (see Infinities below), and for an uncertain number, which is no bound. The
expressions in the bounds are evaluated with interval arithmetic, intervals
included; the left (respectively the right) bound of the result is then
used, depending on the side it appears in: \code{"[cos([0,1]), 2]"} is
\itv{\cos 1}{2}. A bound left out is infinite.

An interval literal \code{[l, u]} whose left bound is above its right bound
is the empty set, as in IEEE 1788-2015: \code{"[2,1]"}, and also
\code{"[9007199254740993, 9007199254740992]"} and
\code{"[1.00000000000000001, 1]"}, whose numbers are compared exactly with
the doubles around them, a number that is not a double lying strictly
between the two doubles around it.\newinvfive{} GAOL~4 read the last two as
the points $2^{53}$ and 1. Two numbers that are not doubles and lie between
the same two doubles are compared exactly too:
\code{"[1.00000000000000002, 1.00000000000000001]"} is the empty set, and
\code{"[1.00000000000000001, 1.00000000000000002]"} is the interval from
the double below to the double above them, as is the same number written
twice. Three kinds of bounds are not told apart: \code{pi} and another
number lying between the same two doubles; a number far beyond the doubles
(about $10^{311}$ or more, or $10^{-331}$ or less, in decimal, about
$2^{1026}$ or more, or $2^{-1081}$ or less, in the other bases); and bounds
that are computed, as \code{1/3}. The literal is then the interval from the
lower bound of the first to the upper bound of the second (case (b) of
12.12.7, which GAOL signals by no exception: \pxref{sec:ieee1788-missing}).

The two bounds in the string \code{<a, b>} must be expressions that
evaluate to the same double whatever the rounding direction; a
\code{std::invalid\_argument} exception is raised otherwise, as for
\code{"<0.1, 0.1>"}, a number that is not a double being no bound of
\code{<a, b>}.\newinvfive{} GAOL~4 read
\code{"<0.1, 0x1.9999999999999p-4>"}, the double just below $0.1$, as the
point $[$\code{0x1.9999999999999p-4}$]$. A bound that is computed is a
double when its value is a point, as \code{0.5}, \code{(0.25+0.25)} and
\code{[0.5]} are; \code{"<[0.1], 0x1.9999999999999p-4>"} and
\code{"<0x1.5555555555556p-2, 1/3>"}, whose first or second bound has the two
doubles around the number for its value, throw \code{std::invalid\_argument}
too, where GAOL~4 read them as the point of the other bound.\newinvfive{}
The two doubles are compared by their bits, $-0$ and $+0$ being the same
bound, so that the mode of the processor that reads the subnormal numbers
as zero (denormals-are-zero) changes nothing: \code{"<0x1p-1074, 0x1p-1073>"},
two subnormal doubles, throws under it too, where a comparison of the bounds as
doubles would take them for one.\newinvfive

\begin{Rem}
GAOL~4 wrote the degenerate intervals in the form \code{<a, b>}, and the
texts it wrote for those whose digits are exact, as \code{<4, 4>}, are still
read. \code{operator<{}<} writes such an interval \code{[a]}, the literal of
IEEE 1788-2015 for a point, and any other with its two bounds
(\pxref{sec:output-format}).
\end{Rem}

\paragraph{Infinities.} An interval literal whose left bound is $+\infty$ or
whose right bound is $-\infty$ is the empty set, as \ieee{numsToInterval} of
IEEE 1788-2015 has it: \code{"[inf]"} and \code{"[-inf, -inf]"} are
empty.\newinvfive{} The expression \code{"inf"} alone is still
\itv{\text{MAX}}{+\infty}, and \code{"-inf"} \itv{-\infty}{-\text{MAX}}
(\pxref{par:inf_inf_note}).

\paragraph{Threads.} Several threads may read strings at once, in parallel:
the lexer is reentrant and the parser pure, and each reading of a string has
a state of its own.\newinvfive{} GAOL~4 kept that state in global variables,
and two threads reading a string each crashed.

\paragraph{Long and nested strings.} A string may be as long as memory
allows: a sum of 100000 terms is read in a time proportional to its length,
and takes no more of the stack of the program than a sum of two
terms.\newinvfive{} The reader of GAOL~4 built a tree with a node for each
operation, which was evaluated and deleted with one call within another per
node: a sum of 100000 terms overflowed the stack and crashed the program.
The operations are now computed as they are read, and the expressions that
a program builds through \file{gaol/parser/gaol\_expression.h} are
evaluated and deleted without recursion too. A token (a number, a name or a
run of blanks) holds $2^{29}$ characters at most, the buffer of the lexer: a
longer one throws \code{std::bad\_alloc}. The nesting of a string is limited
by the stack of the parser, which holds 10000 entries. An open parenthesis
or bracket, or a sign, takes one entry until the expression it opens is
complete; the name and the parenthesis of a call take two; an operand and
the operator or the comma after it take two, until the operand that follows
them is read. About 10000 nested parentheses, brackets or signs, 5000 nested
calls of one argument and 3300 levels of \code{1+(1+(\ldots))} are read; a
string nested deeper is refused as any string that is not an expression,
and throws \code{std::invalid\_argument}.

A \code{std::string}, given to \code{textToInterval()} or read by
\code{operator>{}>}, is read whole: a NUL character is no character of any
token, and a string holding one is not an expression.\newinvfive{} The
reader of GAOL~4 read it up to the NUL, so that \code{"[1,2]"} followed by
a NUL and any text was \itv{1}{2}, and a string of $2^{32}+3$ characters
was read as its first 3, as \code{"1+1"} for \code{"1+1+1+\ldots"}: the
lexer of flex held the length of a string in an \code{int}.

\paragraph{Errors.} A string that is not an expression throws
\code{std::invalid\_argument}. An error found anywhere in the string stops
the reading, and no interval is given: \code{"[nth\_root(8,1.5)]+[1,2]"}
throws, where the two grammars of GAOL~4 gave
\itv{-\infty}{+\infty}.\newinvfive{} Where memory runs out, the reading
throws \code{std::bad\_alloc}, and frees all it allocated.\newinvfive{} In
GAOL~4, flex wrote ``out of dynamic memory'' and ended the program, and the
parser, which read a stack it could not grow as a string nested too deep,
threw no such exception. Without exceptions, an error that the reader finds
as it computes (a wrong call, an uncertain number as a bound, two different
bounds in \code{<a, b>}, or a bound of it that is no double) makes the
reading fail, as a syntax error does: \code{gaol\_core::parse\_interval()}, of
\file{gaol/parser/gaol\_parser.h}, returns
\code{false}, and \code{gaol::textToInterval()} and \code{operator>{}>}
report the error and abort, as for any string that is not an
expression.\newinvfive{} GAOL~4 reported the error, and aborted, inside the
reader.

\begin{Ex}
The program below reads intervals from strings of the language: a literal
with a computed bound, the constants, the form of \code{textToInterval()}
with two strings, a literal divided by a literal, the bounds left out, a
hexadecimal number, the uncertain form, and a literal in a sum and in a
bound. What the printing statements write follows them.
\begin{example}
interval x = textToInterval("[4, 6*7]");
interval y = textToInterval("[-inf, dmax]");
interval z = textToInterval("[3.14,3.15]/8", "[3.14,3.15]/7");
interval t = textToInterval("[3.14,3.15]/[7,8]");
cout << textToInterval("[1,]") << ' ' << textToInterval("[ ]") << ' '
     << textToInterval("[Entire]") << endl;
@outputs^[1, inf] [empty] [-inf, inf]~
cout << textToInterval("[0x1.3p-1, 2/3]") << ' '
     << textToInterval("3.56?1") << endl;
@outputs^[0.59375, 0.6666666666666668] [3.549999999999999, 3.570000000000001]~
cout << textToInterval("1+[1,2]") << ' '
     << textToInterval("[cos([0,1]), 2]") << endl;
@outputs^[2, 3] [0.5403023058681396, 2]~
\end{example}
\end{Ex}

\section{Writing intervals}\label{sec:writing}

An interval is written to a stream with the operator \code{<{}<}, as any
primitive type of C++ is. The text depends on the active format
(Section~\ref{sec:output-format}); the empty set is written the same way in
every format.
\cindex{interval!writing}%

\begin{deffun}{ostream\&}{operator<{}<}{(ostream\& \var{out}, const interval\& \var{I})}
  Prints the interval \var{I} to the output stream \var{out}. The way the
  intervals are displayed depends on the active format
  (Section~\ref{sec:output-format}). Whatever the format, an empty interval
  is displayed as \code{[empty]}, as IEEE 1788-2015 writes the empty set
  (12.11.3).\newinvfive{} The formats \element{width} and \element{center}
  of GAOL~4 wrote \code{empty}, which is no interval literal.
  \cindex{empty!interval: printing}%
  \index[idx]{empty@\code{[empty]}}%
\end{deffun}

\noindent
In the decimal formats, each bound is rounded outward, whatever the C
library: its magnitude is written by the C library, rounded to nearest, then
compared exactly with the bound, and its last digit is moved by one when it
is on the wrong side. The text read back therefore encloses the interval
written.\newinvfive{} GAOL~4 let the C library round each bound in the
direction it set, which the C runtime of Windows and musl do not honour.

\begin{deffun}{std::string}{exact\_string}{(const interval\& \var{I})}
  Returns what \code{operator<{}<} writes in the format \element{hexa}
  (Section~\ref{sec:output-format}), without reading nor changing the global
  output format: \code{"[empty]"} for the empty set, \code{"[a]"} for a
  degenerate interval (\code{"[0x0p+0]"} for a zero of either sign), and
  otherwise the two bounds in the hexadecimal-significand form of IEEE
  1788-2015 (13.4.1), which \code{textToInterval()} reads back as the same
  doubles, a zero degenerate interval as \itv{+0}{+0}, the same set. This is
  the exact text representation of the standard, its
  \ieee{intervalToExact}.\newinvfive{} The text is the same whatever the
  locale of the program, a locale writing a decimal comma included, and
  whatever the C library.
\begin{example}
cout << exact_string(textToInterval("0.1"));
@outputs^[0x1.9999999999999p-4, 0x1.999999999999ap-4]~
\end{example}
\end{deffun}

\subsection{Converting intervals to strings}

The class \code{interval} provides a conversion operator to the standard
type \code{std::string}, which writes the interval in the active format.
\cindex{interval!conversion to \code{string}}%
\index[idx]{string@\code{string}!converting interval to|see{interval}}

\begin{example}
interval I(3,4);
string s = "test line embedding " + string(I) + " as a string";
@rem^// Now, s is "test line embedding [3, 4] as a string"~
\end{example}

\subsection{The output format}\label{sec:output-format}

Intervals may be displayed in five formats, chosen by the static method
\code{interval::format()} described after them.
\cindex{output format}%
\cindex{format!of the output}%

\begin{itemize}
\item[$\bullet$] \element{agreeing}. All the digits that are the same in
the left and in the right bound are printed once, followed by an interval
holding the remaining digits:
\begin{center}
``\element{3.141\textasciitilde[5926, 6001]}'' stands for
``\element{\itv{3.1415926}{3.1416001}}''
\end{center}
Each bound is written with the digits of the precision and the flag
\code{showpoint}, in the global locale of the program. The digits are
shared only by two bounds that are finite, not zero and of the same sign,
written with as many digits before the point and the same exponent, and
that share their first digit that is not 0: the other intervals are printed
as in the format \element{bounds}. The digits of each bound are
kept, but for the zeros that end them, \itv{1.25}{1.2567} being
\code{1.25\textasciitilde[0, 67]}.\newinvfive{} GAOL~4 dropped from both bounds what followed the last
character of the left one that is not a zero, and wrote \itv{1.25}{1.2567}
\code{1.25\textasciitilde[, ]}. It shared the digits of every interval
whose right bound is at most ten times the left one, \itv{1}{10} being
\code{1\textasciitilde[., 0.]}, \itv{4}{6} \code{\textasciitilde[4., 6.]},
without a shared digit, and a zero \code{\textasciitilde[-0., 0.]}, and
printed every interval with a negative bound with its bounds. A point
interval that the digits write exactly is printed as a single number,
without the brackets: \code{interval(4.0)} is \code{4.000000000000000},
where the format \element{bounds} prints \code{[4]}; any other point
interval is printed as any interval is, \code{interval(0.1)} being
\code{0.100000000000000\textasciitilde[0, 1]}. The point zero is printed
\code{[0]}, as in the format \element{bounds} (the formats \element{width}
and \element{center} print it \code{0}, and \element{hexa}
\code{[0x0p+0]}).
\begin{example}
interval::format(interval_format::agreeing);
std::cout << interval(4,6) << ' ' << interval(1.25,1.2567) << ' '
          << interval(-1.2567,-1.25) << std::endl;
@outputs^[4, 6] 1.25@textasciitilde[0, 67] -1.25@textasciitilde[67, 0]~
\end{example}

\item[$\bullet$] \element{bounds}. The bounds are printed between square
brackets. A degenerate interval whose value is written exactly by the
digits displayed is printed as that value alone between square brackets,
the literal of IEEE 1788-2015 for a point (12.11), which
\code{textToInterval()} reads back as the same point:
\begin{example}
interval::format(interval_format::bounds);
std::cout << interval(3,5) << ' ' << interval(4) << std::endl;
@outputs^[3, 5] [4]~
\end{example}
A zero is printed \code{[0]} whatever the signs of its bounds:
\itv{-0}{+0}, which \code{interval::zero()} is with the SSE2 and the NEON
intervals, \itv{+0}{+0} and \itv{-0}{-0} are the same set, which \code{[0]}
and \code{[-0]} are both read back as, and the text is then the same on
every build (the format \element{hexa} prints it \code{[0x0p+0]}). A
degenerate interval whose value has more digits than those displayed is
printed like any other interval, with its two bounds rounded outward:
\code{interval(0.1)} is printed \code{[0.1, 0.1000000000000001]}, which is
read back, as any other text, as an interval containing it. The bounds are
printed with a decimal point, whatever the locale of the stream, so that
the reader, whose decimal point is a point whatever the locale, reads the
text back: under a locale writing a decimal comma, it would take
\code{[-2,5]} for the interval \itv{-2}{5}, and refuses
\code{[0,5, 2,5]}.\newinvfive{} GAOL~4 printed every degenerate interval
between angles, the text of its value rounded downward and upward:
\code{interval(4)} was \code{<4, 4>}, a form of GAOL that is no literal of
the standard, and \code{interval(0.1)} \code{<0.1, 0.1000000000000001>},
which its reader refused (\code{std::invalid\_argument}), the angles
standing for one double only. The reader still takes the angles
(\pxref{sec:input-format}): a text such as \code{<4, 4>}, written by
GAOL~4, is read back.

\item[$\bullet$] \element{width}. The midpoint and the radius are printed,
\code{midpoint()} and \code{rad()}, the \ieee{mid} and \ieee{rad} of IEEE
1788-2015 (12.12.8): the
interval is contained in the midpoint plus or minus the radius, which is
the smallest number that makes it so:
\begin{example}
interval::format(interval_format::width);
std::cout << interval(4,5) << std::endl;
@outputs^4.5 (+/- 0.5)~
\end{example}
The midpoint is written rounded to nearest and the radius rounded upward,
so that a radius that is not zero is never written 0. Rounded upward, a
radius may look noisy with the default 16 digits: \itv{0.1}{0.3} is written
\code{0.2 (+/- 0.1000000000000001)}, its radius being the double nearest to
$0.1$, which is a little above $0.1$; this is intended. A degenerate
interval has no radius: it is written as its midpoint alone. An unbounded
interval has an infinite radius, its midpoint being $0$ or plus or minus
the largest double (\code{[3, inf]} is
\code{1.797693134862316e+308 (+/- inf)}). The midpoint and the radius are
written with the decimal point of the locale of the stream and its grouping
of the digits, the radius without a sign, even with the flag
\code{showpos}, and a midpoint $0$ as \code{0}.\newinvfive{} GAOL~4 wrote
$(a+b)/2$ and $(b-a)/2$ rounded to nearest, which do not contain the
interval (\itv{1}{1+2^{-52}} was \code{1 (+/- 1.11e-16)}), a radius 0 for
\itv{0}{5\cdot10^{-324}}, and \code{inf} for the midpoint of
\itv{10^{308}}{1.7\cdot10^{308}}; its manual and its header described this
number as the width.

\item[$\bullet$] \element{center}. The midpoint alone is printed, as the
format \element{width} prints it.\newinvfive

\item[$\bullet$] \element{hexa}. The left and the right bounds are printed
in the hexadecimal-significand form of IEEE 1788-2015 (13.4.1), which is
exact and which \code{textToInterval()} reads back as the same doubles,
without the round-off error of a binary-to-decimal conversion: this is the
recovery requirement of the standard (13.4), met whatever the precision,
the flags and the locale of the stream. As in the format \element{bounds},
a degenerate interval is printed as its value alone between square
brackets, \code{[0x1.8p+0]} for \code{interval(1.5)}, and a zero
\code{[0x0p+0]} whatever the signs of its bounds, which is read back as
\itv{+0}{+0}, the same set. GAOL~4 printed the sixteen hexadecimal digits
of each double, which its parser did not read.\newinvfive
\cindex{hexadecimal!form}%
\begin{example}
interval::format(interval_format::hexa);
std::cout << textToInterval("0.1") << std::endl;
@outputs^[0x1.9999999999999p-4, 0x1.999999999999ap-4]~
\end{example}
\end{itemize}

\begin{Rem}
The formats \element{bounds} and \element{hexa} are the ones read back as
an input (\pxref{sec:input-format}), whatever the locale of the stream. The
formats \element{width} and \element{center} are for the eye: the midpoint
is written with the digits of the precision, so that the text
\code{c (+/- w)} does not necessarily contain the interval, where the text
of the formats \element{bounds}, \element{agreeing} and \element{hexa} does.
A radius that took the rounding of the midpoint into account would show the
interval \code{0.1}, which is one number, as \code{0.1 (+/- 5.6e-18)}, and
would not show that an interval is narrower than the digits of its
midpoint, which is what the format is for.
\end{Rem}

\medskip
The format is chosen through the following static methods.

\defstaticmethodx{interval}{void}{format}{(interval\_format::format\_t \var{f})}
\begin{defstaticmethod}{interval}{interval\_format::format\_t}{format}{(void)}
  The first form sets the format used by the subsequent printing of
  intervals. The second form returns the format in use, a value of type
  \code{interval\_format::format\_t} (see below, and
  Section~\ref{sec:output-example} for an example of use). The format is
  global to the program, not a property of a stream: it is shared by all the
  threads, and setting it while another thread writes an interval is a data
  race, one of the common errors listed in Chapter~\ref{chap:overview}
  (\pxref{sec:common-errors}).
  \code{exact\_string()} gives the hexadecimal text without changing it.
\end{defstaticmethod}

\begin{deftype}{struct}{interval\_format}
  Structure type used to choose the output format of the intervals. It has
  five possible values of type \element{interval\_format::format\_t}:
\begin{itemize}
\item \element{interval\_format::agreeing}.
\item \element{interval\_format::bounds}.
\item \element{interval\_format::width}.
\item \element{interval\_format::center}.
\item \element{interval\_format::hexa}.
\end{itemize}
\end{deftype}

\begin{Ex}
The program below prints the same interval, \code{interval::pi()}, in the
five formats. What each statement prints is given in the comment under it.
\begin{example}
interval I(interval::pi());

interval::format(interval_format::agreeing);
cout << I << "\n";
@rem^// Prints 3.14159265358979@textasciitilde[3, 4]~
@rem^// The @textasciitilde[] part is dropped if the bounds agree on all digits~
@smallskip
interval::format(interval_format::bounds);
cout << I << "\n";
@rem^// Prints [3.141592653589793, 3.141592653589794]~
@smallskip
interval::format(interval_format::width);
cout << I << "\n";
@rem^// Prints 3.141592653589793 (+/- 4.440892098500627e-16)~
@smallskip
interval::format(interval_format::center);
cout << I << "\n";
@rem^// Prints 3.141592653589793~
@smallskip
interval::format(interval_format::hexa);
cout << I << "\n";
@rem^// Prints [0x1.921fb54442d18p+1, 0x1.921fb54442d19p+1]~
\end{example}
\end{Ex}

\subsection{The number of digits}

The number of digits written for each bound is set by the static method
\code{precision()} of the class \code{interval}. The precision of the
stream, which \code{std::setprecision} sets, is not used for intervals, and
is left as it was: the numbers written after an interval keep it. The width
of the stream (\code{std::setw}) and its adjustment apply to the whole
interval. With \code{std::internal}, the fill goes between the sign of the
midpoint and its digits in the formats \element{width} and \element{center},
and between the sign and the shared digits in the format \element{agreeing},
as it does for a number, and before the \code{[} in the other formats, as
\code{std::right} has it.\newinvfive{} GAOL~4 set the precision of the
stream to that of the intervals and left it so, and \code{std::setw} padded
the \code{[} of an interval alone.
\cindex{precision!of the output}%

\begin{defstaticmethod}{interval}{std::streamsize}{precision}{(void)}
  Returns the number of digits used for printing the bounds of intervals,
  16 by default. Section~\ref{sec:output-example} gives an example.
\end{defstaticmethod}

\begin{defstaticmethod}{interval}{std::streamsize}{precision}{(std::streamsize \var{n})}
  Sets the number of digits used for printing the bounds to \var{n}, and
  returns the number of digits previously used.
  Section~\ref{sec:output-example} gives an example.
\end{defstaticmethod}

\subsection{Example}\label{sec:output-example}

\begin{Ex}
The program below sets the precision and the format, prints $\pi$, reads
the format back, and prints $\pi$ again with sixteen digits in the format
\element{width}. Its output is shown under it.
\begin{example}
#include <iostream>
#include <gaol/gaol.h>
using namespace gaol;

using std::cout;
using std::endl;

int main(void)
{
  interval::precision(4);
  interval::format(interval_format::bounds);
  cout << interval::pi() << endl;

  if (interval::format() != interval_format::bounds) {
    cout << interval::pi() << endl;
  } else {
    interval::precision(16);
    interval::format(interval_format::width);
    cout << interval::pi() << endl;
  }
  gaol::cleanup();
  return 0;
}
@outputs^[3.141, 3.142]~
@outputs^3.141592653589793 (+/- 4.440892098500627e-16)~
\end{example}
\end{Ex}

\noindent
The first call to \code{interval::format()} is unnecessary, since the
default format is \element{interval\_format::bounds}.

The constant \code{interval::pi()} is a predefined canonical interval
containing $\pi$ (\pxref{sec:interval-constants}). Here, its bounds are two
consecutive doubles, $2^{-51}$ apart: the midpoint written is one of them,
and the radius is their distance, the smallest that has the interval within
the midpoint plus or minus it.\newinvfive{} GAOL~4 wrote half of it.

\begin{Rem}
An interval written in decimal and read back gives an interval that
encloses it, whatever the precision chosen, the bounds being rounded
outward; the fewer the digits, the wider. The format \element{hexa} gives
the same interval back.
\end{Rem}

The format \element{width} is useful whenever the number of digits
displayed is insufficient to tell whether the result is a single
floating-point number or an interval too narrow for these digits to show:
      if an interval is not a single
number, the radius displayed is different from zero, whatever its size and
the number of digits, since the radius is rounded upward. A degenerate
interval is displayed without a radius.

\section{Discussion}

This chapter described the two ways an interval crosses the border of a
program. On the way in, GAOL v5 reads every string with one grammar, in
which the literals, the numbers, the constants and the functions combine
with the operators. Its numbers are read exactly, whatever their base, their
number of digits, the locale and the C library, and a number that is not a
double becomes the tightest interval of doubles containing it. Its reader
is reentrant and computes the operations as they are read, so that several
threads may read at once and the length of a string is limited by the
memory alone. On the way out, the decimal bounds are rounded outward
whatever the C library, so that a text read back encloses the interval
written, and the format \element{hexa} gives the interval back exactly.

\medskip
The formats are trade-offs. The decimal formats \element{bounds} and
\element{agreeing} are readable and enclose the interval, at the price of a
bound moved by one unit of its last digit when the text nearest to it lies
on the wrong side; the format \element{hexa} is exact and meets the recovery
requirement of the standard; the formats \element{width} and
\element{center} are for the eye, and their text does not necessarily
contain the interval. The format and the precision are global to the
program, where the width and the adjustment belong to the stream. Reading a
long number takes a time quadratic in its number of digits.

\medskip
The next chapter turns to the namespace \code{gaol\_ieee1788}, which offers
the operations under the names of IEEE 1788-2015, the reading and the
writing of intervals included, and lists where the standard and GAOL
differ.

%% !TEX root = ../manual_gaolv5.tex
% Chapter 10.
%% Chapter 10: The names of IEEE 1788-2015. Sources: gaol.tex of f5bf180 lines
%% 5532-5788 (the chapter "The names of IEEE 1788-2015" of the current
%% manual), doc/using.md lines 451-571 (the names of IEEE 1788-2015, the
%% reverse functions of atan2).
%% Based on the manual of GAOL 4 by Frederic Goualard (gaol.tex).
%%
%% Copyright (c) 2026 ENSTA, France
%%
%% By: Jordan NININ (created 2026-11-06)
%%
%% GAOL is a software distributed WITHOUT ANY WARRANTY. Read the associated
%% COPYING.LIB file for information.
\chapter{The names of IEEE 1788-2015}\label{chap:ieee1788}
%%===================================
\cindex{namespace!gaol\_ieee1788}\cindex{IEEE 1788-2015!names}%

IEEE 1788-2015 \cite{IEEE1788:2015} names the operations of interval
arithmetic, and GAOL names its operations its own way. The reverse function
\ieee{coshRev}$(c,x)$ of the standard is \code{acosh\_rel(c, x)} in GAOL,
\ieee{mulRev}$(b,c,x)$ is \code{div\_rel(c, b, x)}, with its arguments in
another order, and \ieee{roundTiesToEven}$(x)$ is
\code{round\_ties\_to\_even(x)}. A reader of the standard who comes to GAOL
has to look for the translation of each name.

\medskip
The namespace \code{gaol\_ieee1788} spares this search. It gives each
operation of the standard that GAOL provides the name and the argument order
the standard gives it, so that a reader of the standard finds it without
looking for its translation. It is the second of the two namespaces of GAOL
v5 (\pxref{sec:namespaces}); the standard itself is presented in
Chapter~\ref{chap:concepts} (\pxref{sec:ieee1788-overview}).

\medskip
In this chapter, we describe this namespace. The first section presents it
and the way a program opens it. The second lists the operations whose
meaning differs between the standard and GAOL, and which follow the standard
in \code{gaol\_ieee1788}. The third gives the table of the correspondence
between the names of the standard and those of GAOL, then details
\ieee{mulRevToPair}, which has no name in \code{gaol}, and the reverse
functions of \ieee{atan2}. The fourth describes the reading of intervals
with the names of the standard, and the fifth lists what GAOL does not
provide of the standard. The names of the other interval libraries, next to
those of GAOL and of the standard, are compared in
Chapter~\ref{chap:overview} (\pxref{sec:other-libraries}).

\section{The namespace of the standard}\label{sec:ieee1788-namespace}
%%------------------------------------

The namespace \code{gaol\_ieee1788} comes with \file{gaol/gaol.h}. It is not in
\code{gaol}, and it holds
the type \code{interval} as well: one directive lets a program use GAOL under
the names of the standard. The output format is not in it: a program that opens
only \code{gaol\_ieee1788} sets it with
\code{interval::format(gaol::interval\_format::hexa)}.\newinvfive

\begin{Ex}\label{ex:ieee1788-program}
The program below opens \code{gaol\_ieee1788} and computes with the names of
the standard: \code{numsToInterval} builds an interval from its two bounds,
\code{mulRev} is the reverse multiplication, which divides \code{x} by
\itv{2}{2}, \code{sinPi} and \code{rootn} are two functions of Table~10.5 of
the standard, \code{entire} is the interval \itv{-\infty}{+\infty} and
\code{strictLess} the strict order of its Table~10.3. What it prints
follows it.
\begin{example}
#include <iostream>
#include <gaol/gaol.h>
using namespace gaol_ieee1788;

int main(void)
{
  interval x = numsToInterval(1, 2);
  interval y = mulRev(numsToInterval(2, 2), x);   @rem^// x / 2~
  interval z = sinPi(x) + rootn(y, 3);
  bool b = strictLess(x, entire());
  std::cout << std::boolalpha
            << x << ' ' << y << ' ' << z << ' ' << b << std::endl;
  gaol::cleanup();
  return 0;
}
@outputs^[1, 2] [0.5, 1] [-0.2062994740159004, 1] false~
\end{example}
\end{Ex}

\medskip
\noindent
A program opens \code{gaol} or \code{gaol\_ieee1788}, not both
(\pxref{sec:namespaces}). The functions that have the same name and the same
meaning in both are the same functions; the ones whose meaning differs,
\code{pow} and \code{textToInterval}, exist in each namespace, and a program
that opened both would find each call ambiguous.

\begin{Rem}\label{rem:ieee1788-shadow}
A name of the program's own that is also a name of the standard, a constant
\code{inf} for instance, makes every unqualified use of it ambiguous: the
program qualifies both, \code{::inf} and \code{gaol\_ieee1788::inf(x)}.
The same holds for \code{less(x, y)} next to \code{using namespace std;}, where it is
ambiguous with the class template \code{std::less}.
\end{Rem}

\section{Where the standard and GAOL differ}\label{sec:ieee1788-differences}
%%-----------------------------------------
\cindex{IEEE 1788-2015!differences with GAOL}%

The functions of GAOL that have the name and the meaning the standard gives
them (\code{sin}, \code{exp}, \code{sqrt}, \code{min}\ldots) are the same
functions in \code{gaol\_ieee1788}. Where the standard and GAOL differ, the
names of \code{gaol\_ieee1788} follow the standard. The differences are the
following ones.

\begin{itemize}
\item[$\bullet$] \textbf{The power.} \code{pow(x, y)} is the \ieee{pow} of
  Table~9.1, defined for $x>0$, and for $x=0$ when $y>0$. The \code{pow} of
  \code{gaol} takes the integer power for an integer exponent, a negative
  base included (\pxref{sec:powers}). \code{pow(x, p)} with a number $p$ is
  \code{pow(x, [p])}, an integer $p$ included, \code{pow(x, 2)} as well, as
  the standard gives no other meaning to it: the power with an integer
  exponent is \code{pown(x, n)} in the standard, whose $n$ is an integer
  rather than an interval. An integer exponent beyond the \code{int}s gives
  the power of CORE-MATH at the bounds of the part of $x$ in
  \itv{0}{+\infty}, as the \code{pow} of \code{gaol} does there; the \code{pow}
  of \code{gaol} takes the integer power on the negative part of $x$ too. The results of the two powers are
  compared in Chapter~\ref{chap:overview} (\pxref{sec:pow-namespaces}).
\item[$\bullet$] \textbf{The bounds of the empty set.} \code{inf(x)} and
  \code{sup(x)} are $+\infty$ and $-\infty$ for the empty set (Table~10.2),
  where \code{left()} and \code{right()} are NaN. \code{inf} returns $-0$
  for a left bound 0 (12.12.8).
\item[$\bullet$] \textbf{Membership.} \code{isMember(m, x)} is \code{false}
  for an infinite $m$ (10.6.3).
\item[$\bullet$] \textbf{Writing.} \code{intervalToText(x)} writes the
  bounds of \code{x} rounded outward, \code{[l, r]}, or \code{[a]} for a
  degenerate interval whose value the digits write exactly, as \code{[4]},
  or \code{[empty]}, with the digits of \code{interval::precision()} and a
  decimal point. This string is an interval literal of the standard (12.11),
  which \code{textToInterval} reads back as an interval containing
  \code{x} (13.3), whatever the output format that
  \code{interval::format()} sets and the locale of the program.
  \code{operator<{}<} writes the current format instead, which is no
  literal in the formats \element{width} and \element{center}, nor in the
  format \element{agreeing} where the bounds share digits.\newinvfive
\item[$\bullet$] \textbf{Reading.} \code{textToInterval} and
  \code{exactToInterval} read the names of the functions of the standard,
  and return the empty set for a string that is no interval (12.1.3), where
  \code{gaol::textToInterval} reads the names of GAOL and throws
  (\pxref{sec:ieee1788-text}). A null pointer is no string, and they throw
  \code{std::invalid\_argument} for it too (\pxref{sec:constructors}).%
  \newinvfive
\end{itemize}

\begin{Ex}\label{ex:ieee1788-differences}
The program below shows these differences. The \code{pow} of the standard is
empty on a negative base, where \code{pown} gives the integer power; an
integer exponent beyond the \code{int}s gives the power of CORE-MATH at the
bounds; the bounds of the empty set are the infinities; an infinity is a
member of no interval; and a string that is no interval reads as the empty
set.
\begin{example}
cout << pow(interval(-4,-1), interval(2.0)) << ' '
     << pown(interval(-4,-1), 2) << ' '
     << pow(interval(2,3), interval(1e10)) << endl;
@outputs^[empty] [1, 16] [1.797693134862315e+308, inf]~
cout << boolalpha << inf(empty()) << ' ' << sup(empty()) << ' '
     << isMember(GAOL_INFINITY, entire()) << ' '
     << textToInterval("[1, 2") << endl;
@outputs^inf -inf false [empty]~
\end{example}
\end{Ex}

\section{The operations}\label{sec:ieee1788-operations}
%%----------------------
\cindex{IEEE 1788-2015!operations}%

Table~\ref{tab:ieee1788-operations} gives, for each operation of the standard
that GAOL provides, its name in \code{gaol\_ieee1788} and the name GAOL gives
it, grouped as the standard groups them, with the clause or the table of the
standard that defines them. The last argument of the reverse functions
is optional, and \itv{-\infty}{+\infty} when it is left out (10.5.4).

\begingroup\renewcommand{\arraystretch}{1.1}\small
\begin{longtable}{L{6.2cm}L{7.8cm}}
  \caption{The operations of IEEE 1788-2015 in \code{gaol\_ieee1788}, and
    their names in GAOL}\label{tab:ieee1788-operations}\\
  \hline
  \multicolumn{1}{c}{\code{gaol\_ieee1788}} & \multicolumn{1}{c}{GAOL}\\
  \hline\hline
  \endfirsthead
  \multicolumn{2}{l}{\textit{Table~\thetable, continued}}\\*
  \hline
  \multicolumn{1}{c}{\code{gaol\_ieee1788}} & \multicolumn{1}{c}{GAOL}\\
  \hline\hline
  \endhead
  \hline
  \endfoot
  \multicolumn{2}{l}{\textit{Constants and constructors (10.5.2, 10.5.8)}}\\*
  \code{empty()}, \code{entire()} &
    \code{interval::emptyset()}, \code{interval::universe()}\\
  \code{numsToInterval(l, u)} & \code{interval(l, u)}\\
  \code{textToInterval(s)} &
    \code{gaol::textToInterval(s)} with the names of the standard, the
    empty set for a wrong string\\
  \hline
  \multicolumn{2}{l}{\textit{Forward elementary functions (Table~9.1 of the standard)}}\\*
  \code{neg}, \code{add}, \code{sub}, \code{mul}, \code{div} &
    \code{-x}, \code{x + y}, \code{x - y}, \code{x * y}, \code{x / y}\\
  \code{recip(x)} & \code{inverse(x)}\\
  \code{sqr}, \code{sqrt}, \code{fma} & the same\\
  \code{pown(x, p)} &
    \code{gaol::pow(x, p)} for an integer $p$ of any type\\
  \code{pow(x, y)}, \code{pow(x, p)} & the \ieee{pow} of the standard\\
  \code{exp}, \code{exp2}, \code{exp10}, \code{log}, \code{log2},
    \code{log10} & the same\\
  \code{sin}, \code{cos}, \code{tan}, \code{asin}, \code{acos},
    \code{atan}, \code{atan2} & the same\\
  \code{sinh}, \code{cosh}, \code{tanh}, \code{asinh}, \code{acosh},
    \code{atanh} & the same\\
  \code{sign}, \code{ceil}, \code{floor}, \code{trunc} & the same\\
  \code{roundTiesToEven}, \code{roundTiesToAway} &
    \code{round\_ties\_to\_even}, \code{round\_ties\_to\_away}\\
  \code{abs}, \code{min}, \code{max} & the same\\
  \hline
  \multicolumn{2}{l}{\textit{Recommended functions (Table~10.5 of the standard)}}\\*
  \code{rootn(x, q)} & \code{nth\_root(x, q)}\\
  \code{expm1}, \code{exp2m1}, \code{exp10m1} & the same\\
  \code{logp1}, \code{log2p1}, \code{log10p1} &
    \code{log1p}, \code{log2p1}, \code{log10p1}\\
  \code{hypot}, \code{rSqrt} & \code{hypot}, \code{rsqrt}\\
  \code{compoundm1} & the same\\
  \code{sinPi}, \code{cosPi}, \code{tanPi} &
    \code{sinpi}, \code{cospi}, \code{tanpi}\\
  \code{asinPi}, \code{acosPi}, \code{atanPi}, \code{atan2Pi} &
    \code{asinpi}, \code{acospi}, \code{atanpi}, \code{atan2pi}\\
  \hline
  \multicolumn{2}{l}{\textit{Reverse functions (Table~10.1 of the standard)}}\\*
  \code{sqrRev(c, x)}, \code{absRev(c, x)} &
    \code{sqrt\_rel(c, x)}, \code{invabs\_rel(c, x)}\\
  \code{pownRev(c, x, p)} & \code{nth\_root\_rel(c, p, x)}, for every integer $p$,
    $p\leq0$ included\\
  \code{sinRev}, \code{cosRev}, \code{tanRev(c, x)} &
    \code{asin\_rel}, \code{acos\_rel}, \code{atan\_rel(c, x)}\\
  \code{coshRev(c, x)} & \code{acosh\_rel(c, x)}\\
  \code{sinhRev}, \code{tanhRev(c, x)} (not in the standard, named after
    \ieee{coshRev}) & \code{asinh\_rel}, \code{atanh\_rel(c, x)}\\
  \code{mulRev(b, c, x)} & \code{div\_rel(c, b, x)}\\
  \code{powRev1(b, c, x)}, \code{powRev2(a, c, x)} &
    \code{pow\_rel(c, b, x)},
    \code{pow\_exponent\_rel(c, a, x)}\newinvfive\\
  \code{atan2Rev1(x, c, y)}, \code{atan2Rev2(y, c, x)} &
    \code{atan2\_rel\_y(c, x, y)},
    \code{atan2\_rel\_x(c, y, x)}\newinvfive\\
  \code{mulRevToPair(b, c)} (10.5.5) &
    the parts of $c/b$ on either side of 0, each a
    \code{div\_rel} with \code{interval::universe()} for its last argument, in a
    \code{std::pair}\\
  \hline
  \multicolumn{2}{l}{\textit{Cancellative operations and set operations
    (10.5.6, 10.5.7)}}\\*
  \code{cancelMinus}, \code{cancelPlus} &
    \code{cancel\_minus}, \code{cancel\_plus}\\
  \code{intersection(x, y)}, \code{convexHull(x, y)} &
    \code{x \& y}, \code{x | y}\\
  \hline
  \multicolumn{2}{l}{\textit{Numeric functions (Table~10.2 of the standard)}}\\*
  \code{inf(x)}, \code{sup(x)} &
    \code{x.left()}, \code{x.right()}, but for the empty set and a zero bound
    ($-0$ for \code{inf}, $+0$ for \code{sup})\\
  \code{mid}, \code{wid}, \code{rad}, \code{mag}, \code{mig(x)} &
    \code{x.midpoint()}, \code{width()}, \code{rad()}, \code{mag()},
    \code{mig()}\\
  \code{midRad(x, m, r)} & \code{x.mid\_rad(m, r)}\\
  \hline
  \multicolumn{2}{l}{\textit{Boolean functions (10.5.10, Tables~10.3 and
    10.4 of the standard; 10.6.3)}}\\*
  \code{isEmpty(x)}, \code{isEntire(x)} &
    \code{x.is\_empty()}, \code{x.is\_entire()}\\
  \code{equal(a, b)}, \code{subset(a, b)} &
    \code{a.set\_eq(b)}, \code{a.set\_leq(b)}\\
  \code{interior(a, b)}, \code{disjoint(a, b)} &
    \code{a.set\_le(b)}, \code{a.set\_disjoint(b)}\\
  \code{less(a, b)}, \code{strictLess(a, b)} &
    \code{a.less(b)}, \code{a.strictly\_less(b)}\\
  \code{precedes(a, b)}, \code{strictPrecedes(a, b)} &
    \code{a.certainly\_leq(b)}, \code{a.certainly\_le(b)}\\
  \code{isCommonInterval(x)}, \code{isSingleton(x)} &
    \code{x.is\_common\_interval()}, \code{x.is\_a\_double()}\\
  \code{isMember(m, x)} &
    \code{x.set\_contains(m)}, false for an infinite $m$\\
  \code{overlap(a, b)} (10.6.4) &
    \code{a.overlap(b)}, an \code{overlapping\_state}\\
  \hline
  \multicolumn{2}{l}{\textit{Input and output (Clause 13)}}\\*
  \code{intervalToText(x)} &
    the bounds \code{[l, r]}, \code{[a]} or \code{[empty]}, with the
    current precision, whatever the format and locale\\
  \code{intervalToExact(x)} & \code{exact\_string(x)}\\
  \code{exactToInterval(s)} & \code{textToInterval(s)} of \code{gaol\_ieee1788}\\
  \hline
  \multicolumn{2}{l}{\textit{Expressions}}\\*
  \code{pown(e, n)}, \code{pow(e1, e2)} &
    the expressions of \code{pown} and of the \ieee{pow} of the standard,
    declared by \file{gaol/parser/\allowbreak gaol\_expression.h}\\
\end{longtable}
\endgroup

\noindent
Two rows of the table call for a word. \code{overlap(a, b)} is the interval
overlapping function of the standard: it returns an
\code{overlapping\_state}, one of the sixteen states of Table~10.7 under the
names of the table, and the type and the member \code{a.overlap(b)} are in
\code{gaol} too, which has no other name for it
(Chapter~\ref{chap:relations}).\newinvfive{} The expressions \code{pown(e, n)}
and \code{pow(e1, e2)} are computed with the \code{pown} and the \code{pow}
of this namespace; \file{gaol/parser/gaol\_expression.h} declares them with
the other expressions, and \file{gaol/gaol.h} does not include it, as in
GAOL~4. The last two parts of this section detail \code{mulRevToPair} and
the reverse functions of \code{atan2}.

\paragraph{The two-output division.}
\findex{mulRevToPair}%
\code{mulRevToPair(b, c)} returns the set
$\{x \mid \exists\beta\in b,\ \beta x\in c\}$ in two intervals, the pair
$(u,v)$ of the standard: $u$ before $v$ when the set has two parts, which
happens when $c$ does not contain 0 and 0 is strictly between the bounds of
$b$, $(u,\emptyset)$ when it is one interval, and $(\emptyset,\emptyset)$ when
it is empty. The convex hull of $u$ and $v$ is \code{mulRev(b, c)}. Where
$b$ is unbounded on both sides, the two parts share 0.\newinvfive

\begin{Ex}\label{ex:ieee1788-mulrevtopair}
The program below divides \itv{1}{2} by \itv{-1}{1}, whose interior holds 0,
and gets two parts; then by \itv{2}{4}, and gets one part and the empty set.
\begin{example}
std::pair<interval, interval> p = mulRevToPair(numsToInterval(-1, 1),
                                               numsToInterval(1, 2));
std::pair<interval, interval> q = mulRevToPair(numsToInterval(2, 4),
                                               numsToInterval(1, 2));
cout << p.first << ' ' << p.second << ' '
     << q.first << ' ' << q.second << endl;
@outputs^[-inf, -1] [1, inf] [0.25, 1] [empty]~
\end{example}
\end{Ex}

\paragraph{The reverse functions of atan2.}
\findex{atan2Rev1}\findex{atan2Rev2}%
\code{atan2Rev1(x, c, y)} is the hull of the ordinates of the points of the
box $x\times y$ whose angle $\operatorname{atan2}(y,x)$ is in $c$, and
\code{atan2Rev2(y, c, x)} the hull of their abscissas; the point $(0,0)$,
which has no angle, is left out. In \code{gaol}, they are
\code{atan2\_rel\_y(c, x, y)} and \code{atan2\_rel\_x(c, y, x)}
(\pxref{sec:relational-arithmetic}).\newinvfive{} As in \code{mulRev} and
\code{div\_rel}, the interval contracted comes last, and the other argument
of \code{atan2} before it, so that the constraint
$\theta=\operatorname{atan2}(y,x)$ contracts its three variables with the
three lines below.
\begin{example}
theta &= atan2(y, x);
y = atan2Rev1(x, theta, y);
x = atan2Rev2(y, theta, x);
\end{example}

\noindent
The angles are those of \code{atan2}, in $\left]-\pi,\pi\right]$. The
half-line $y=0$, $x<0$ has the angle $\pi$, not $-\pi$, and a $c$ holding
$-\pi$ but not $\pi$ leaves it out: \code{atan2Rev2([0], [-4, -3], [-2, 2])}
is empty. The angles of $c$ beyond $\pi$ are no angle of any point:
\code{atan2Rev1([-1], [3, 4], [-1, 1])} is \itv{0}{-\tan 3}, about
\itv{0}{0.1425}, the points of angles 3 to $\pi$ only. A program whose
angles go around past $\pi$ gives the two sides of the jump apart,
\itv{3}{\pi} and \itv{-\pi}{4-2\pi} for \itv{3}{4}, and takes the hull of
the two results.

\section{Reading with the names of the standard}\label{sec:ieee1788-text}
%%-----------------------------------------------
\findex{textToInterval}\cindex{IEEE 1788-2015!reading}%

\code{textToInterval(\var{s})} reads a string as \code{gaol::textToInterval()}
does (\pxref{sec:input-format}), with the names of the functions of IEEE
1788-2015: those of Tables~9.1 and 10.5 that GAOL provides, whatever the case
of their letters.\newinvfive{} These names are the following ones.

\begin{center}\small
\begin{tabular}{L{11cm}}
\code{neg}, \code{add}, \code{sub}, \code{mul}, \code{div}, \code{recip},
\code{sqr}, \code{sqrt}, \code{fma}, \code{pown}, \code{pow}, \code{exp},
\code{exp2}, \code{exp10}, \code{log}, \code{log2}, \code{log10}, \code{sin},
\code{cos}, \code{tan}, \code{asin}, \code{acos}, \code{atan}, \code{atan2},
\code{sinh}, \code{cosh}, \code{tanh}, \code{asinh}, \code{acosh},
\code{atanh}, \code{sign}, \code{ceil}, \code{floor}, \code{trunc},
\code{roundTiesToEven}, \code{roundTiesToAway}, \code{abs}, \code{min},
\code{max}, \code{rootn}, \code{expm1}, \code{exp2m1}, \code{exp10m1},
\code{logp1}, \code{log2p1}, \code{log10p1}, \code{hypot},
\code{compoundm1}, \code{rSqrt}, \code{sinPi}, \code{cosPi}, \code{tanPi},
\code{asinPi}, \code{acosPi}, \code{atanPi}, \code{atan2Pi}\\
\end{tabular}
\end{center}

\noindent
\code{pow} is the \ieee{pow} of Table~9.1, and the second argument of
\code{pown} and \code{rootn} is an integer. The names of GAOL alone
(\code{nth\_root}, \code{cbrt}, \code{inverse}, \code{integer},
\code{log1p}, \code{round\_ties\_to\_even}\ldots) are no functions there.
They give the empty set, as any string that is no interval: a wrong call
(\code{"pown([2,5],2.5)"}, \code{"sin(1,2)"}) and a string holding a NUL
character included, whether GAOL was configured with exceptions or
not.\newinvfive{} Where memory runs out, \code{textToInterval} throws
\code{std::bad\_alloc}.

\medskip
As \code{pow}, the function \code{textToInterval} is in the namespace
\code{gaol} too, where it reads the names of GAOL and throws
\code{std::invalid\_argument} for a wrong string, and \code{std::domain\_error}
for an order of \code{nth\_root} that is no integer: a program calls the one of
the namespace it opens.

\begin{Ex}\label{ex:ieee1788-text}
The program below reads three strings with the names of the standard: a
power of the standard, a power of a negative base, which the \ieee{pow} of
the standard leaves empty, and a name of GAOL alone, which is no function
there.
\begin{example}
cout << textToInterval("pown([2,5],5)") << ' '
     << textToInterval("pow([-4,-1],2)") << ' '
     << textToInterval("nth_root(8,3)") << endl;
@outputs^[32, 3125] [empty] [empty]~
\end{example}
With \code{using namespace gaol;}, \code{textToInterval("pow([-4,-1],2)")}
is \itv{1}{16}, and \code{textToInterval("pown([2,5],5)")} throws.
\end{Ex}

\section{What is not provided}\label{sec:ieee1788-missing}
%%----------------------------
\cindex{IEEE 1788-2015!operations not provided}\cindex{decorations}%

GAOL has no decorations: only bare intervals are provided (Clause 11). Nor
does it provide the slope functions (Table~10.6), the reduction operations
\ieee{sum}, \ieee{dot}, \ieee{sumSquare} and \ieee{sumAbs} (12.12.12), the
exact reduction operations of 12.13.5, the conversion specifier \var{cs} of
\ieee{intervalToText} (13.3), or the exceptions \ieee{UndefinedOperation} and
\ieee{PossiblyUndefinedOperation} that \ieee{textToInterval} and
\ieee{exactToInterval} signal (12.12.7, 13.4): for a string that is no
interval they return the empty set silently. Where they cannot tell whether
$l \le u$ in a literal \code{[l, u]}, case (b) of 12.12.7 (the bounds of
\pxref{sec:reading-intervals}), they return the interval from the lower bound
of $l$ to the upper bound of $u$, without \ieee{PossiblyUndefinedOperation}.
That interval holds $[l, u]$ where $l \le u$; where $l$ is above $u$ it is a
false positive that the standard permits, and it can miss $l$ or $u$ when a
bound is computed: \code{"[1/3, 0x1.5555555555555p-2]"} is the point
\code{0x1.5555555555555p-2}, below $1/3$. The end of
\file{gaol/gaol\_ieee1788.h} keeps this list.

\begin{Rem}\label{rem:ieee1788-erf}
In the other direction, a function of GAOL that the standard does not name
is not in \code{gaol\_ieee1788} either, but for \code{sinhRev} and
\code{tanhRev}, which are named after \ieee{coshRev}
(Table~\ref{tab:ieee1788-operations}), for \code{centeredForm(m, r)}, the
constructor \code{interval::centered\_form(m, r)} of a midpoint and a radius
(\pxref{sec:constructors}), and for the operators of intervals, the
relational division \code{\%} among them, which the namespace names as
\code{gaol\_core} does. \code{erf} and \code{erfc}, the
error function and its complement, are two functions of GAOL that the
standard does not name.\newinvfive{} A call
\code{erf(x)} on an interval finds them all the same in \code{gaol\_core},
by argument-dependent lookup, and \code{erf(0.5)} remains the \code{erf}
of C.
\end{Rem}

\section{Discussion}
%%------------------

This chapter has presented the namespace \code{gaol\_ieee1788}, which lets a
program use GAOL under the names and the argument orders of IEEE 1788-2015
with a single directive. The functions that the standard and GAOL name and
define alike are shared between the two namespaces. Where they differ, the
power, the bounds of the empty set, the membership of an infinity, the
writing and the reading of intervals, the namespace follows the standard, and Table~\ref{tab:ieee1788-operations}
gives the translation of every other name. Two operations have their
details here: \code{mulRevToPair}, which \code{gaol} does not name, and the
reverse functions of \code{atan2}, whose arguments \code{gaol\_ieee1788}
orders as \ieee{mulRev} does and \code{gaol} as \code{div\_rel} does.

\medskip
The choice of two namespaces rather than one keeps the programs written for
GAOL~4 and, but for the results that change (Chapter~\ref{chap:changes}),
their results, while a program written for the standard, or compared with
another library of the standard, opens \code{gaol\_ieee1788}. Its price is
stated in \exref{sec:ieee1788-namespace}: a program opens one of the two
namespaces, never both, and qualifies the names of its own that are also names
of the standard. What GAOL does not provide of the standard, the decorations, the
slope functions, the reductions, exact or not, the conversion specifier of
\ieee{intervalToText} and the exceptions of the reading, is listed rather
than approximated (\pxref{sec:ieee1788-missing}), and the header keeps the list.

\medskip
Whatever the name a program calls an operation by, the bounds it gets are
the same. The next chapter states how tight they are: the conditions GAOL
guarantees on each bound, and the tightness of each family of functions.

%% !TEX root = ../manual_gaolv5.tex
% Chapter 11.
%% Chapter 11: Accuracy of the operations. Sources: gaol.tex of f5bf180 5881-6006
%% (chapter "Accuracy of the operations"); doc/accuracy.md (whole).
\chapter{Accuracy of the operations}\label{chap:accuracy}
%%=================================
\cindex{accuracy!of the operations}\cindex{tightness}%

The containment property (\pxref{sec:containment}) says that the result of
an operation encloses its exact value; it says nothing of how closely.
IEEE 1788-2015 \cite{IEEE1788:2015} requires an implementation to document
the \dfn{tightness} of each of its interval operations, dividing the
possible inputs into ranges and stating the tightness achieved in each
(12.10.3). This chapter does so for \code{interval}, the inf-sup type of
GAOL on doubles, from the algorithm of each operation. The tests of
\file{tests} check each statement on random and special arguments
(\pxref{sec:test-programs}), and we give, with each operation, the largest
distance they found, or the distance they accept where the text says no more.\newinvfive

\medskip
IEEE 1788-2015 defines three accuracy modes (12.10.1), which
Chapter~\ref{chap:concepts} presents (\pxref{sec:accuracy-modes}). For an
operation $f$, its exact value $f^*$ and an input box $x$:
\begin{itemize}
\item the result is \emph{tightest} when $f(x)$ is the hull, in doubles, of
  $f^*(x)$: each bound is the nearest double on its side of the exact bound;
\item it is \emph{accurate} when $f(x)$ encloses $f^*(x)$ and is within
  \[\text{nextOut}(f_{\text{tightest}}(\text{nextOut}(x))),\]
  nextOut moving each bound one double outward: roughly, $f(x)$ is no wider
  than the tightest result for an input one double wider, plus one double on
  each side;
\item it is \emph{valid} when $f(x)$ encloses $f^*(x)$.
\end{itemize}

\noindent
The standard requires the basic operations (\ieee{neg}, \ieee{add}, \ieee{sub},
\ieee{mul}, \ieee{div}, \ieee{recip}, \ieee{sqr}, \ieee{sqrt}, \ieee{fma}),
the integer functions and \code{abs}, \code{min} and \code{max} to be
tightest; it requires the other elementary functions, the reverse functions
and the recommended functions to be valid, and recommends them to be
accurate (12.10.2). Below, ``within $k$ doubles'' means that each bound is
at most $k$ doubles beyond the tightest one.

\medskip
GAOL bounds its elementary functions with CORE-MATH \cite{CoreMath:2022},
which is correctly rounded in the rounding direction in effect
(\pxref{sec:math-libraries}).\cindex{CORE-MATH!accuracy}
Computing in the upward rounding it keeps, GAOL takes the value at a bound
as the upper bound, with nothing to add, and the double below it as the
lower one. The upper bound is the tightest one, and so is the lower bound unless the exact
value is a double, which the operations give exactly where it happens:
$\log(1)=0$ and $\sin(0)=0$ are exact, and $\arcsin(1)$ is the tightest
enclosure of $\pi/2$. Where an algorithm evaluates the function at the
bounds of the interval, the bounds of the result are therefore the tightest
ones.

\medskip
The first section of this chapter states the conditions under which these
statements hold. The second gives the tightness of each operation in a
table, then, family by family, the algorithm of each operation and the
largest distance from the tightest bounds that the tests found. The third
section explains why the bounds are the same on every machine, and where
the sign of a zero bound differs between the kernels.

\section{The conditions}\label{sec:accuracy-conditions}
%%----------------------------------------------------
\cindex{accuracy!conditions}%

The statements of this chapter hold for GAOL as the three builds of
Chapter~\ref{chap:installation} make it: with the CORE-MATH of
\file{3rd/math-core}, compiled as the builds compile it, and with the code
that includes the headers of GAOL compiled with the flags of interval
arithmetic (\pxref{sec:compiling}). They hold for the SSE2 intervals, the
NEON intervals\newinvfive{} and the FPU intervals alike
(\pxref{sec:kernels}), and for the AVX-512 path (\pxref{sec:avx512}), whose
bounds are those of the SSE2 path, but for the sign of a zero bound of
\code{x * d} (\pxref{sec:accuracy-same-bounds})

\medskip
They hold whatever the rounding direction the calling code left
(\file{tests/\allowbreak rounding\_direction.cpp}), and with the floating-point
exceptions masked, as a program starts: an operation that raises an enabled
exception stops the program rather than give its bounds
(\pxref{sec:fp-exceptions}). They hold on every architecture and with every
compiler alike, CORE-MATH giving the same bits everywhere: there is no other
mathematical library to build GAOL with (\pxref{sec:math-libraries}, and
Chapter~\ref{chap:changes} for what differs from GAOL~4). The intervals of
floats, \code{gaol::intervalf} and \code{gaol::interval2f}, unfinished,
which only a developer of GAOL compiles (\file{gaol/gaol\_config.h}), are
not covered.

\section{The tightness of each operation}\label{sec:accuracy-table}
%%-----------------------------------------------------------------
\cindex{accuracy!table of the operations}%

Table~\ref{tab:accuracy-tightness} gives the tightness of each operation of
\code{interval}, in the families of IEEE 1788-2015, from its algorithm; the
subsections that follow give the algorithm of each operation and what the
tests found. The tables numbered 9.1, 10.1, 10.3, 10.4, 10.5 and 10.7 that this
chapter names are those of IEEE 1788-2015.

\begingroup\renewcommand{\arraystretch}{1.1}\small
\begin{longtable}{L{4.8cm}L{9.2cm}}
  \caption{The tightness of each operation of \code{interval}}%
  \label{tab:accuracy-tightness}\\
  \hline
  \multicolumn{1}{c}{Operations} & \multicolumn{1}{c}{Tightness}\\
  \hline\hline
  \endfirsthead
  \multicolumn{2}{l}{\textit{Table~\thetable, continued}}\\*
  \hline
  \multicolumn{1}{c}{Operations} & \multicolumn{1}{c}{Tightness}\\
  \hline\hline
  \endhead
  \hline
  \endfoot
  \multicolumn{2}{l}{\textit{Basic operations (Table~9.1): tightest
    required}}\\*
  \code{-x}, \code{+}, \code{-}, \code{*}, \code{/}, \code{inverse},
  \code{sqr} & tightest\\
  \code{sqrt} & tightest, whatever the rounding of the square root of the C
  library\\
  \code{fma} & tightest\\
  \hline
  \multicolumn{2}{l}{\textit{Integer and absmax functions: tightest
    required}}\\*
  \code{floor}, \code{ceil}, \code{integer}, \code{sign}, \code{trunc},
  \code{round\_ties\_to\_even}, \code{round\_ties\_to\_away}, \code{abs},
  \code{min}, \code{max} & tightest, and exact, in every rounding
  direction\\
  \hline
  \multicolumn{2}{l}{\textit{Powers, exponentials and logarithms: valid
    required, accurate recommended}}\\*
  \code{pow(x, n)}, integer \var{n} (\ieee{pown}) & the tightest bounds, for a base of either sign, or one double beyond where the power is within $2K$ of a
  double, relatively, $K$ being the bound of the relative error of the power
  ($K=2^{-98}$ for $n=3$, about $2^{-62}$ at most, at $n=-2^{31}$; no such case for $n=2$, which is
  \code{sqr(x)}); this holds below $2^{-968}$ and beyond the largest double
  too; the bounds are exact where the power is a double; the same bounds
  with the SSE2, the NEON and the FPU intervals. Beyond the \code{int}s:
  tightest where $n$ is a double, within three doubles otherwise\\
  \code{pow(x, y)} (\ieee{pow}) & tightest, at the corners of the box and
  at their limits where the base is $0$ or $+\infty$ or the exponent
  infinite; for a degenerate integer exponent within the \code{int}s, the
  integer power of the row above, one double beyond the tightest where it
  is within twice the error of its products
  (\code{pow([0x1.fffffffffffffp+1], [7])} has its lower bound one double
  below the tightest)\\
  \code{exp}, \code{exp2}, \code{exp10}, \code{log}, \code{log2},
  \code{log10} & tightest; exact where the value is a double ($\exp(0)$,
  $\log(1)$, $2^n$, $10^n$ for $0\leq n\leq22$\ldots)\\
  \hline
  \multicolumn{2}{l}{\textit{Trigonometric and hyperbolic functions: valid
    required, accurate recommended}}\\*
  \code{cos}, \code{sin} & tightest where the function is evaluated at the
  bounds, at every magnitude; $-1$ and $1$ exactly at the extrema;
  \itv{-1}{1} from the width $2\underline{\pi}$ on, $\underline{\pi}$ being
  the double below $\pi$\\
  \code{tan} & tightest between two poles, at every magnitude;
  \itv{-\infty}{+\infty} exactly when a pole is within $x$, but for the
  widths above $\underline{\pi}$ and below $\pi$\\
  \code{asin}, \code{acos}, \code{atan}, \code{atan2} & tightest; exactly
  0, or the tightest enclosures of $\pm\pi/4$, $\pm\pi/2$ and $\pi$, where
  the value is one of them\\
  \code{sinh}, \code{cosh}, \code{tanh}, \code{asinh}, \code{acosh},
  \code{atanh} & tightest\\
  \hline
  \multicolumn{2}{l}{\textit{Recommended functions (Table~10.5)}}\\*
  \code{nth\_root(x, n)}, $n>0$ & $n=1$, 2, 3: tightest. Otherwise
  the tightest bounds or up to 2 doubles beyond, proved with
  integer powers; beyond the \code{unsigned int}s, CORE-MATH's \code{pow}
  at the corners, within one double\\
  \code{nth\_root(x, q)}, $q<0$ & valid: the root above, then the
  tightest division, up to 4 doubles beyond the tightest bounds\\
  \code{expm1}, \code{exp2m1}, \code{exp10m1}, \code{log1p},
  \code{log2p1}, \code{log10p1}, \code{hypot}, \code{rsqrt} & tightest\\
  \code{compoundm1} & tightest but in a case out of $10^{17}$ or so, where
  a bound is a double beyond\newinvfive\\
  \code{sinpi}, \code{cospi}, \code{tanpi}, \code{asinpi}, \code{acospi},
  \code{atanpi}, \code{atan2pi} & tightest, at every magnitude\\
  \hline
  \multicolumn{2}{l}{\textit{The error functions, which IEEE 1788-2015 does
    not name}}\\*
  \code{erf}, \code{erfc} & tightest; exact at 0 and at the
  infinities\newinvfive\\
  \hline
  \multicolumn{2}{l}{\textit{Reverse functions (Table~10.1): valid
    required, accurate recommended}}\\*
  \code{sqrt\_rel}, \code{nth\_root\_rel}, \code{acosh\_rel},
  \code{asinh\_rel}, \code{atanh\_rel} & accurate\\
  \code{nth\_root\_rel(c, p, x)}, $p<0$ & valid: the root of each part
  of $c$ and then its inverse, up to four doubles beyond the
  tightest\newinvfive\\
  \code{invabs\_rel} & tightest\\
  \code{max\_rel}, \code{min\_rel}, \code{sign\_rel}, \code{floor\_rel}
  (not in Table~10.1) & tightest\\
  \code{asin\_rel}, \code{acos\_rel}, \code{atan\_rel} & accurate, or one
  double beyond; a bound of $x$ beyond $2^{52}$ is kept\\
  \code{div\_rel}, \code{\%} & accurate; tightest where the bounds of the
  quotients are doubles\\
  \code{pow\_rel}, \code{pow\_exponent\_rel} & tightest, but for a bound of
  $x$ next to the solutions whose power is within one double of $c$, which
  is kept: the tightest result for $x$ one double wider\newinvfive\\
  \code{atan2\_rel\_y}, \code{atan2\_rel\_x} & accurate; within two doubles
  of the tightest bounds where these are not empty\\
  \code{cancel\_minus}, \code{cancel\_plus} & tightest\\
  \hline
  \multicolumn{2}{l}{\textit{Set operations, constructors, numeric and
    boolean functions}}\\*
  \code{\&}, \code{|}, \code{interval(l, u)} & tightest (exact)\\
  \code{textToInterval(s)}, \code{operator>{}>} & tightest for the literals
  of IEEE 1788-2015; valid for expressions, computed with the operations
  above\\
  \code{midpoint()}, \code{width()}, \code{rad()} & as 12.12.8 requires\\
  \code{mag()}, \code{mig()}, \code{smig()} & exact\\
  \code{mid()}, \code{hausdorff()} & tightest\\
  relations, \code{overlap()} & exact, as Tables~10.3, 10.4 and 10.7\\
  \code{operator<{}<} & \code{hexa}: exact, read back as the same interval.
  Decimal formats: valid, and the tightest with the digits asked for where
  the C library rounds to nearest as it should. \element{width} and
  \element{center}: $\itv{c-w}{c+w}$ contains $x$, $w$ being the smallest radius
  that does so. \code{exact\_string()}: exact, read back bit for bit\\
\end{longtable}
\endgroup

\begin{Rem}\label{rem:accuracy-before}
Before GAOL v5, the elementary functions but the hyperbolic ones were
correctly rounded to nearest and moved one double outward, which is at most one double beyond the
tightest bound: accurate, and not tightest. The hyperbolic functions were
taken from the mathematical library of the system, whose values are
sometimes further than one double from the exact ones, and millions of
doubles away for the \code{acosh} of some of them.\newinvfive
\end{Rem}

\subsection{Basic operations}\label{sec:accuracy-basic}

The basic operations of Table~9.1 of the standard must be tightest. Each is
computed from the bounds, rounded upward once (\pxref{sec:kernels}), and
each is tightest; the tests find them so.

\begin{itemize}
\item[$\bullet$] \textbf{Negation.} \code{-x} exchanges the bounds: exact.
\item[$\bullet$] \textbf{Addition and subtraction.} \code{x + y},
  \code{x - y}, \code{+=}, \code{-=}, and their forms with a double, are the
  sums of the bounds rounded upward, the lower bound being stored negated:
  tightest.
\item[$\bullet$] \textbf{Multiplication.} \code{x * y}, \code{*=}, and the
  forms with a double, are the products of the bounds chosen by their signs,
  rounded upward; a zero bound times an infinite one counts as $0$:
  tightest.
\item[$\bullet$] \textbf{Division.} \code{x / y}, \code{/=}, and the forms
  with a double, are the quotients of the bounds chosen by their signs,
  rounded upward; the hull of the quotients when $y$ contains $0$, and the
  empty set when $y$ is $\itv{0}{0}$: tightest.
\item[$\bullet$] \textbf{Inverse.} \code{inverse(x)} and \code{x.inverse()}
  divide $1$ by the bounds, rounded upward: tightest.
\item[$\bullet$] \textbf{Square.} \code{sqr(x)} squares the bounds, rounded
  upward, with $0$ as lower bound when $x$ contains $0$: tightest.
\item[$\bullet$] \textbf{Square root.} \code{sqrt(x)} works on
  $x\cap\itv{0}{+\infty}$. The upper bound $u$ is the square root rounded
  upward, from the square root of the C library checked with a product
  rounded the other way and moved one double when needed; the lower bound
  is $u$ when $u\cdot u$ is exactly the bound of $x$, and the double below
  $u$ otherwise: tightest, whatever the rounding of the square root of the
  C library.
\item[$\bullet$] \textbf{Fused multiply-add.} In \code{fma(x, y, z)}, the
  product $x\cdot y$ is bilinear, so it ranges over the hull of the
  products of the bounds, and $x\cdot y+z$ ranges from the least product
  plus $\leftBound{z}$ to the greatest plus $\rightBound{z}$. Each corner
  is rounded once by \code{std::fma}, upward for the upper bound and
  downward for the lower one, computed as $-\mathrm{fma}(-a, b, -c)$. A
  bound $0$ times an infinite one counts as $0$.\newinvfive{} Tightest,
  over 200\,000 boxes against exact rational arithmetic.
\end{itemize}

\begin{Rem}
The result of \code{std::fma} goes through \code{rnd\_keep()} before it is
negated. Without it, GCC folded $-\mathrm{fma}(-a, b, -c)$ into
$\mathrm{fma}(a, b, c)$, a single instruction rounded upward, and the lower
bound came above the exact value, one double above the tightest lower bound.
\end{Rem}

\subsection{Integer and absmax functions}\label{sec:accuracy-integer}

The integer functions, \code{abs}, \code{min} and \code{max} must be
tightest. None of them rounds: each returns a double that is an integer, or
a bound of its argument, in every rounding direction.

\begin{itemize}
\item[$\bullet$] \textbf{Floor and ceiling.} \code{floor(x)} and
  \code{ceil(x)} apply to each bound: exact.
\item[$\bullet$] \textbf{Integers of an interval.} \code{integer(x)}, which
  the standard does not name, is the hull of the integers in $x$,
  $\itv{\lceil\leftBound{x}\rceil}{\lfloor\rightBound{x}\rfloor}$: exact.
\item[$\bullet$] \textbf{Sign, truncation and roundings.} \code{sign(x)},
  \code{trunc(x)}, \code{round\_ties\_to\_even(x)} and
  \code{round\_ties\_to\_away(x)} take the value at each bound, each
  function being non-decreasing, as \code{ceil} and \code{floor} above, and
  $\emptyset$ for $\emptyset$, which the four test: GAOL holds the empty
  interval as two NaN bounds, which \code{trunc} and the two roundings send
  to themselves, where \code{sign} would give $\itv{0}{0}$, a NaN comparing
  false both to $0$ and above it. \code{trunc} and
  \code{round\_ties\_to\_away} are \code{std::trunc} and \code{std::round},
  which the C++ standard defines without reading the rounding direction,
  and \code{round\_ties\_to\_even} reads the bits
  (\file{gaol/math/gaol\_roundeven.h}), where \code{std::nearbyint} and
  \code{std::rint} would round in the direction in effect, upward in
  GAOL.\newinvfive{} Tightest and exact; the tests find the same result in
  the four rounding directions over 22\,407 doubles: the halfway values,
  the whole numbers and the doubles on either side of them.
\item[$\bullet$] \textbf{Absolute value.} \code{abs(x)} takes the magnitudes
  of the bounds, and $0$ as lower bound when $x$ contains $0$: exact.
\item[$\bullet$] \textbf{Minimum and maximum.} \code{min(x, y)} and
  \code{max(x, y)} apply to the bounds: exact.
\end{itemize}

\subsection{Powers, exponentials and logarithms}\label{sec:accuracy-powers}

These functions must be valid, and the standard recommends them to be
accurate. The exponentials and the logarithms are tightest, being
CORE-MATH's functions at the bounds; the integer power is tightest but one
double beyond in rare cases, and the real power is tightest at the corners
of its box.

\paragraph{The integer power.}\findex{pow}%
\code{pow(x, n)} for an integer $n$ of any type, and \code{pow(x, y)} and
\code{pow(x, d)} for a degenerate integer $y$ or an integer double $d$,
compute \ieee{pown} (\code{gaol\_pown()}, \code{gaol\_uipow()},
\code{gaol\_pown\_ll()}), as does \code{gaol\_ieee1788::pown(x, n)}
(\pxref{sec:powers}). Within the \code{int}s, $n=2$ is \code{sqr(x)}, the
squares of the bounds rounded outward.\newinvfive{} For any other $n$, the
powers of the two corners of $x$ where $x^n$ is the least and the greatest
are computed at once, in the two lanes of an SSE2 register with the FMA
instructions or of a NEON register, or of a pair of doubles otherwise (the
FPU intervals, and SSE2 without FMA), with the same operations in the same
order.\newinvfive{} The power is a binary exponentiation from exact
products: each product is rounded, and its rest, which \code{fma} gives
exactly, is carried along, so that the power is $h+l$ within $K\,|h|$ of
it, with $K=(8nL+16)\cdot2^{-104}$ for an $n$ of $L$ bits, a bound proved
beforehand (\file{gaol/core/gaol\_interval\_core.cpp}). Each bound is then
rounded once, $h+l+K|h|$ upward and $h+l-K|h|$ downward, and $K$ is left
out where every rest is $0$, $x^n$ being a double. The power $x^{-n}$ is
$1/(h+l)$: the inverse of $h$ rounded, corrected with the rest of
$1-q(h+l)$ to the second order and rounded once the same way; $x^{-1}$ is
the inverse. A corner whose powers leave $\itv{2^{-968}}{2^{960}}$ is
computed alone, scaled to $[1, 2)$, its power kept below $2^{501}$ by
powers of two, then scaled back with one rounding in the direction of the
bound, beyond the largest double and among the subnormals too. The bounds
$0$ and $\pm\infty$ give the limits of $x^n$:
$\itv{\mathrm{mag}(x)^n}{+\infty}$ where $x$ holds $0$ for an even $n<0$,
$\itv{-\infty}{+\infty}$ where it holds $0$ inside for an odd one, and
$\emptyset$ for $\itv{0}{0}$.

Beyond the \code{int}s, CORE-MATH's \code{pow} is called at the corners,
as for the real power below, the sign coming from the parity of $n$; for an
$n$ of more than 53 bits that is no double, $x^h\cdot x^l$ is rounded
outward, with $n=h+l$, $h$ a double and $|l|<2^{11}$; beyond 64 bits, $n$
is taken as $\pm(2^{63}-1)$ or $\pm(2^{63}-2)$, of its sign and its parity,
every power of a double but $0$, $\pm1$ and $\pm\infty$ being beyond the
doubles there. GAOL~4 gave $\itv{-\infty}{+\infty}$ beyond the
\code{unsigned int}s, and beyond the \code{int}s for \code{pow(x, y)};
within them, it rounded $1/x^n$ twice, and rounded the products outward, as
it computed every power, where a bound was $0$ or infinite or its power
beyond the doubles.\newinvfive

The result is tightest, or one double beyond where $x^n$ is within $2K$ of
a double, relatively: never for $n=2$, which is \code{sqr(x)}; $K$ is
$2^{-98}$ at $n=3$, $2^{-88}$ at $n=1000$, $2^{-65}$ at $n=2^{31}-1$ and
about $2^{-62}$ at $n=-2^{31}$, where one bound in about 400 is one double
beyond. It is exact where $x^n$ is a double. Beyond the \code{int}s it is
tightest where $n$ is a double, and within three doubles otherwise.

The tests find the tightest bounds at 5\,000 doubles drawn at random, with
exponents from $-30$ to $30$ and over all the doubles, subnormal results and
overflows included, for $n$ from 2 to 9, 16, 31 to 33 and $-1$ to $-4$,
$-7$, $-16$ and $-33$, and over 5\,000 intervals drawn so, for $n$ from 2
to 5, 12, 31 and $-1$ to $-3$, $-6$, $-17$, $x$ holding $0$ included,
against the exact powers; the tightest over 4\,000 intervals near the ends
of the doubles (a bound $0$ or infinite, powers below $2^{-1000}$ or above
$2^{1100}$); the exact value at the 8\,679 powers that are doubles; the
tightest at 13 powers for $n$ from $2^{24}+1$ to $2^{32}-1$, and at 6
corners computed alone, against mpmath. Beyond the \code{int}s, they find
the tightest at the powers of the tests \file{elementary} and
\file{ieee1788} there (odd and even exponents, negative bases, exponents of
64 and 128 bits), the tightest at 5 powers beyond $2^{53}$ where $n$ is a
double, and within three doubles at 4 where it is none, of intervals and of
expressions, against mpmath. They check that $x^2$ is \code{sqr(x)} over
400 intervals near $1$, $3$, $0.75$ and $2^{-500}$, and bit for bit, the
sign of a zero included, in each of its forms (\code{pow(x, 2)},
\code{2u}, \code{2.0}, \code{[2]}, \code{gaol\_ieee1788::pown(x, 2)} and
\code{pow(x, [2])}), over the special bounds and 90\,000 random intervals,
in each direction of the calling code and with each flush mode; that
$(-c)^n$ is $\pm c^n$, and that the bound from a corner is the same whether
the other corner is computed with it or alone, at 1\,088 points near $1$,
$1.5$ and $3$ and at 12\,736 more within 16 doubles of $1$. ITF1788, the
test suite of Chapter~\ref{chap:computing}, gets the tightest in its 163
cases of \ieee{pown}.

\paragraph{The real power.}
\code{pow(x, y)} and \code{pow(x, d)} for a non-integer $y$
(\code{gaol\_pow\_hybrid()}, \code{gaol\_pow\_real()}), and
\code{gaol\_ieee1788::pow(x, y)} for any $y$, compute \ieee{pow} on
$x\cap\itv{0}{+\infty}$: CORE-MATH's \code{pow} at the corners of the box
$x\times y$ where $x^y$ is the least and the greatest, which the places of
the bounds about $1$ and $0$ give, correctly rounded upward. The upper bound
is the value at the upper corner, and the lower bound is the double below
the value at the lower corner, unless $x^y$ is a double there. It is one
for $x=2^p$ and an integer $p\,y$ of $[-1074, 1023]$, and for
$x=c^{2^k}\cdot2^{f\cdot2^k}$ ($c\geq3$ odd) and $y=a/2^k$ ($a>0$, odd if
$k>0$) where $c^a<2^{53}$ and $c^a\cdot2^{f\cdot a}$ is within the doubles,
$y$ being then a multiple of $1/32$ in $(0, 33]$; the lower bound is then
$x^y$ itself, which a test on $x$ and $y$ proves with integers and exact
operations, leaving at once for nearly every $y$, where calling CORE-MATH
downward too would double the cost of the power.\newinvfive{} The result
is $1$ where a corner has the base $1$ or the exponent $0$. At a corner
where the base is $0$ or $+\infty$ or the exponent infinite, it is the
limit of $x^y$ there, $0$ or $+\infty$, which is the bound itself; GAOL~4
computed $\exp(y\log x)$ there, which gave bounds one to hundreds of
doubles beyond.\newinvfive{} $0^y$ is $0$ for $y>0$, and $\emptyset$ for
$y\leq0$. An integer $y$ beyond the \code{int}s, in
\code{gaol\_ieee1788::pow}, takes the corners at the bounds of $x$, $x^y$
being monotone in $x$, where the \code{pow} of \code{gaol} takes
\ieee{pown}; GAOL~4 gave $\itv{-\infty}{+\infty}$.\newinvfive

The result is tightest, at the corners and at their limits, those at the
bounds of $x$ for an integer $y$ beyond the \code{int}s included. For a
degenerate integer $y$ within the \code{int}s, it is the integer power
above, one double beyond the tightest where $x^y$ is within $2K$ of a
double.

\begin{Ex}
\code{pow([4], 0.5)} is $\itv{2}{2}$ and \code{pow([4, +oo], 0.5)} is
$\itv{2}{+\infty}$. \code{gaol\_ieee1788::pow([0x1.fffffffffffffp+1], [7])}
has its lower bound, \code{0x1.ffffffffffff8p+13}, one double below the
tightest, \code{0x1.ffffffffffff9p+13}: the integer power is within $2K$ of
a double there.
\end{Ex}

The tests find the bounds within 1 double at points and over 385 boxes; the
tightest, or one double beyond for one integer power, on the 117 boxes of
the test \file{ieee1788}, the limits at $0$ and $+\infty$ included, against
mpmath; the tightest on the 1\,347 cases of \ieee{pow} of ITF1788; the
tightest at 5 integer exponents beyond the \code{int}s; and the tightest,
both bounds, at the corners of 84\,070 boxes of a base and an exponent
(powers of two, perfect powers, midpoints, at the edges of the doubles too,
and their neighbours), against $x^y$ computed exactly at the 8\,981 where
it is rational, 3\,211 of them with a power that is a double, and compared
exactly at the 12\,540 where it is the $2^k$-th root of a rational,
$k\leq6$.

\paragraph{Exponentials and logarithms.}
\code{exp}, \code{exp2}, \code{exp10}, \code{log}, \code{log2} and
\code{log10} are CORE-MATH's functions at the bounds, correctly rounded
upward: the value at the right bound, and the double below the value at the
left bound, unless that value is a double. The logarithms work on
$x\cap\itv{0}{+\infty}$, give $\emptyset$ when $x$ holds no positive
number, and $\log(0)=-\infty$; the exponentials have a lower bound at least
$0$. The values that are doubles are given exactly: $\exp(0)=1$ and
$\log(1)=0$; for \code{exp2}, \code{exp10}, \code{log2} and
\code{log10},\newinvfive{} $2^x$ for an integer $x$ of $[-1074, 1023]$ and
$10^x$ for an integer $x$ of $[0, 22]$ ($10^{23}$ is no double, and no
negative power of ten is), $\log_2 x$ for a power of two and $\log_{10}x$
for a power of ten of $[10^0, 10^{22}]$. All six are tightest. The tests
find \code{exp} tightest at about 9\,000 doubles of every magnitude, \code{log} tightest at the 75
doubles of the tests and the million intervals of the comparison with the
other libraries (\pxref{sec:comparison}), and the other four tightest over
6\,483 doubles each, and at the values that are doubles.

\subsection{Trigonometric and hyperbolic functions}\label{sec:accuracy-trig}

These functions must be valid, and the standard recommends them to be
accurate. Every one of them is CORE-MATH's function at the bounds of the
pieces where it is monotonic, so each is tightest where the function is
evaluated at the bounds; the periodic functions ask first on which piece
the interval lies. Every elementary function of this section and of the
one before is that of CORE-MATH, whose sources are in \file{3rd/math-core}
(\pxref{sec:math-libraries}).

\paragraph{Cosine and sine.}
For \code{cos(x)}, the bounds divided by an enclosure of $\pi$, rounded
outward, tell whether $x$ lies on one piece where the cosine is monotonic;
rounded inward, they tell whether an extremum is within $x$ for sure. Where
they cannot tell, a bound being within about $|x|\cdot2^{-52}$ of a
multiple of $\pi$, or beyond $2^{52}$, the signs of CORE-MATH's \code{sin}
at the bounds do, exactly, $x$ being narrower than $2\pi$ (and the sign at
the middle of $x$ from $\pi$ on): signs that differ show one extremum, $-1$
or $1$ according to their order. Then CORE-MATH's \code{cos} is taken at
the bounds, correctly rounded upward (the double below the value at the
bound where the cosine is the least), within $\itv{-1}{1}$, and $\cos(0)=1$ exactly. \code{sin(x)}
proceeds as \code{cos}, the bounds divided by an enclosure of $\pi$ minus
$1/2$, the signs of CORE-MATH's \code{cos} where the quotients cannot tell,
and $\sin(0)=0$ exactly.

Both are tightest where the function is evaluated at the bounds, at every
magnitude, and give $-1$ and $1$ exactly at the extrema. They give
$\itv{-1}{1}$ from the width $2\underline{\pi}$ on, $\underline{\pi}$ being
the double below $\pi$, which is then the tightest result: an interval that wide misses an arc shorter than $2.5\cdot10^{-16}$, where the sine and the cosine stay within $10^{-31}$ of $\pm1$. That the test of a width below
$\underline{\pi}$ gives the tightest bounds is measured, not proved: the
tests find the tightest bounds of \code{sin} and \code{cos} over 29\,400
intervals of width about $\pi$ whose bounds are next to extrema of the
sine. They find them tightest for all doubles, and over intervals next to
the extrema and of width about $\pi$ and $2\pi$.

\paragraph{Tangent.}
\code{tan(x)} is $\itv{-\infty}{+\infty}$ when $x$ is wider than
$\underline{\pi}$ or holds a pole: $(x+[\pi/2])/[\pi]$ at the bounds,
rounded outward, tells that no pole is within $x$ when the integer parts
agree, and rounded inward that one is for sure; where they cannot tell, the
signs of CORE-MATH's \code{cos} at the bounds do, exactly, a pole being
within $x$ when they differ. Otherwise the result is CORE-MATH's \code{tan}
at the bounds, correctly rounded upward, and $\tan(0)=0$ exactly. It is
tightest between two poles, at every magnitude, and
$\itv{-\infty}{+\infty}$ exactly when a pole is within $x$, but for the
widths above $\underline{\pi}$ and below $\pi$, a theoretical case. The
tests find it tightest for all doubles, and over intervals next to the
poles, the interval from \code{-M\_PI\_2} to \code{M\_PI\_2} and the
intervals of width $\underline{\pi}$ or next to it between two poles
included.

\paragraph{Inverse trigonometric functions.}
\code{asin(x)} and \code{acos(x)} are CORE-MATH's functions at the bounds
of $x\cap\itv{-1}{1}$, correctly rounded upward; at $-1$, $0$ and $1$ they
give $0$ exactly or the tightest bounds of $\pm\pi/2$ and $\pi$.
\code{atan(x)} is CORE-MATH's \code{atan} at the bounds; at $0$, $\pm1$
and $\pm\infty$ it gives $0$ exactly or the tightest bounds of $\pm\pi/4$
and $\pm\pi/2$. \code{atan2(y, x)} is CORE-MATH's \code{atan2} at the
corners of the box $y\times x$ where the angle is the least and the
greatest, which the signs of the bounds give, correctly rounded upward; it
gives $0$, $\pm\pi/4$, $\pm\pi/2$ and $\pi$ where a corner is on an axis,
on a diagonal of the right half-plane or at an infinity, $\itv{-\pi}{\pi}$
when the box has points on the half-line $y=0$, $x<0$, where the angle is
$\pi$, and points below it, and $\emptyset$ for the box $\{(0, 0)\}$. The
four are tightest, and exactly $0$ or the tightest enclosures of the
multiples of $\pi/4$ named where a bound is one of them. The tests check
them within 1 double, and tightest at the special points.

\paragraph{Hyperbolic functions.}
\code{sinh(x)}, \code{tanh(x)} and \code{asinh(x)} are CORE-MATH's
functions at the bounds, correctly rounded upward, \code{tanh} within
$\itv{-1}{1}$; they give $0$ at $0$ exactly, $\sinh(x)$ beyond the largest
double for $|x|\geq711$, and $\tanh(x)$ beyond $\pm(1-2^{-53})$ for
$|x|\geq20$, where $1-|\tanh(x)|<2^{-54}$. \code{cosh(x)} takes CORE-MATH's \code{cosh} at the bound of largest magnitude for the upper bound, and at the bound of smallest magnitude for the lower one, which is $1$ when $x$ contains $0$, and a value beyond the largest double for
$|x|\geq711$. \code{acosh(x)} and \code{atanh(x)} proceed as \code{sinh},
on $x\cap\itv{1}{+\infty}$ and $x\cap\itv{-1}{1}$ respectively, with
$\mathrm{acosh}(1)=0$ and $\mathrm{atanh}(0)=0$ exactly. The hyperbolic arc
tangent is defined on $(-1, 1)$, $-\infty$ and $+\infty$ being its limits
at $-1$ and $1$, which are bounds but no values: the result is $\emptyset$
when $x$ holds no point of $(-1, 1)$, as for $\itv{1}{1}$, $\itv{1}{5}$,
$\itv{-1}{-1}$ and $\itv{-5}{-1}$.\newinvfive{} The six are tightest, and
exact at the points given. The tests check them within 1 double, tightest
at $0$, at $1$, and beyond $711$ and $20$, and $\emptyset$ for those four
intervals.

\subsection{Recommended functions}\label{sec:accuracy-recommended}

The recommended functions of Table~10.5 of the standard must be valid, and
the standard recommends them to be accurate. All of them are tightest but
the roots of an order other than 1, 2 and 3, which are within 2 doubles of
the tightest, and
\code{compoundm1}, which is tightest but in a case out of $10^{17}$ or so.

\paragraph{Roots.}\findex{nth\_root}%
\code{nth\_root(x, n)}, for $n>0$, computes \ieee{rootn}. For $n=1$ it is
$x$, and for $n=2$ it is \code{sqrt}. For $n=3$ it is CORE-MATH's
\code{cbrt} at the bounds, correctly rounded upward on their magnitudes,
the root of a negative number being the opposite of the root of its
magnitude: the tightest bounds, and exact where the cube root is a double,
which cubing the value tells; it takes 50~ns rather than 189 with the
search below, on an Intel i7-1185G7 with Clang 18.1.\newinvfive{} For the
other orders, on $x$ (odd $n$, the root of a negative number being again
the opposite of the root of its magnitude) or on $x\cap\itv{0}{+\infty}$
(even $n$), the lower bound is the largest double $l$ with $l^n$ rounded
upward at most the bound of $x$, and the upper bound the smallest double
$u$ with $u^n$ rounded downward at least it, which proves them. They are
looked for from CORE-MATH's \code{pow} with the exponent $1/n$ rounded,
after a step of Newton's method, by steps that double then by bisection,
two powers where the start is next to the root; the roots of $0$ and $1$
are $0$ and $1$. The orders 1, 2 and 3 are tightest, and so is any order
where the bound of $x$ is an $n$-th power of a double. Otherwise the bounds are the tightest or up to 2 doubles beyond, whatever the magnitude of $x$: the $n-1$ rounded products of $l^n$ move the root by less than $2^{-52}$ relatively. Two doubles beyond is not accurate where the order is large: $\code{nth\_root([0x1.16ef146ba87dap-473], 1000)}$ has its upper bound two doubles above the tightest one, one above the accurate one. An order beyond the \code{int}s and
the \code{unsigned int}s takes CORE-MATH's \code{pow} at the corners, with
the exponent $1/n$ as an interval, within 1 double of the
tightest.\newinvfive{} The tests find the roots within 2 doubles for all
doubles, 1 double being the largest distance found, among them the $x$
below $2^{-1020}$, subnormal ones included, with $n$ up to $2^{32}-1$,
against exact arithmetic and mpmath, and the tightest on the values of the
test \file{ieee1788} beyond the \code{unsigned int}s.

\code{nth\_root(x, q)}, for $q<0$, is $1/x^{1/|q|}$, the inverse of the
root above: the domain is then $\RSet\setminus\{0\}$ for an odd $q$ and
$(0, +\infty)$ for an even one, which taking the inverse gives, an interval
holding $0$ having for inverse the hull of the values away from it;
\code{nth\_root(x, 0)} is $\emptyset$.\newinvfive{} It is valid, and within 4 doubles of the tightest bounds: the root above, then the tightest division. It is not always accurate: \code{nth\_root([2.25], -127)} has its upper bound 4 doubles above the tightest one, 3 above the accurate one. The tests find that it encloses the root over 1.5 million
values against a reference in \code{long double}, and that it is the
inverse of the positive root over 8\,000 values.

\paragraph{Exponentials minus one and logarithms of one plus.}
\code{expm1(x)}, \code{exp2m1(x)} and \code{exp10m1(x)} are CORE-MATH's
functions at the bounds, correctly rounded upward, the double below the
value at the left bound unless that value is a double, which $b^x-1$ is at
$0$, and at the integers of $[-53, 53]$ for $b=2$ and of $[0, 15]$ for
$b=10$; the result is within $\itv{-1}{+\infty}$.\newinvfive{}
\code{log1p(x)}, \code{log2p1(x)} and \code{log10p1(x)} are CORE-MATH's
functions at the bounds of $x\cap\itv{-1}{+\infty}$, the value at $-1$
being the limit $-\infty$, and $x$ holding no number above $-1$ giving
$\emptyset$; the double below the value at the left bound unless it is a
double, which $\log_b(1+x)$ is at $0$, at $2^k-1$ for $b=2$ ($k$ in
$[-53, 53]$) and at $10^k-1$ for $b=10$ ($k$ in $[0, 15]$)
only.\newinvfive{} The six are tightest. The tests find the three
exponentials tightest over 400\,000 points and intervals each, and the
three logarithms tightest over 40\,000 intervals each, their bounds drawn
among those points and their neighbours.

\paragraph{The functions of $\pi x$.}
For \code{sinpi(x)} and \code{cospi(x)}, the extrema are at the multiples
of $1/2$, which are doubles: doubling the bounds, which is exact, and
taking their ceiling and floor give exactly the integers $t=2x$ within $x$,
and $\sin(\pi x)$ has its maximum $1$ at the $t$ congruent to $1$ modulo
$4$ and its minimum $-1$ at those congruent to $3$ (the cosine at $0$ and
$2$); otherwise the result is CORE-MATH's function at the bounds, exact at
the multiples of $1/2$. Beyond $2^{61}$ two distinct doubles are whole
periods apart, and the result is $\itv{-1}{1}$.\newinvfive{} Both are
tightest, at every magnitude: $\mathrm{sinpi}(10^{17})$ is $0$, where
\code{sin} of $\pi\cdot x$ gave $\itv{-1}{1}$. \code{tanpi(x)} proceeds as
\code{sinpi}, its poles being at the odd $t$: $\itv{-\infty}{+\infty}$
with a pole within $x$, $-\infty$ or $+\infty$ for a pole at a bound, and
$\emptyset$ for a pole alone, where the function has no
value.\newinvfive{} It is tightest. The tests find the three tightest over
400\,000 points and intervals each.

\code{atanpi(x)}, \code{asinpi(x)} and \code{acospi(x)} are CORE-MATH's
functions at the bounds, \code{asinpi} and \code{acospi} on
$x\cap\itv{-1}{1}$, \code{acospi} decreasing; they are exact at $0$ and
$\pm1$ (and at $\pm\infty$ for \code{atanpi}), the only doubles where the
values are doubles (\code{asinpi}$(\pm1/2)=\pm1/6$, \code{acospi}$(1/2)=1/3$ and \code{acospi}$(-1/2)=2/3$ are rational too, but no doubles).\newinvfive{} \code{atan2pi(y, x)} is the analysis of
the box of \code{atan2}, the angles divided by $\pi$: $\itv{-1}{1}$ for a
box crossing the half-line $y=0$, $x<0$, and the corners elsewhere. The
values at the multiples of $1/4$ ($0$, $\pm1/4$, $\pm1/2$, $\pm3/4$, $1$,
and $-1$ as a limit) are exact, and CORE-MATH's \code{atan2pi} gives the
others, which are irrational.\newinvfive{} The four are tightest. The
tests find them so over 400\,000 points and intervals each (40\,000
intervals for \code{asinpi}, and the arguments next to $\pm1$ where the accurate phase of
CORE-MATH's \code{asinpi} shifted a 64-bit integer by 65 to 69 bits, which \file{3rd/README.md} describes),
and \code{atan2pi} over 200\,000 boxes.

\paragraph{Hypotenuse and inverse square root.}
\code{hypot(x, y)} is CORE-MATH's \code{hypot} at the point of the box
nearest to the origin, $(\mathrm{mig}\,x, \mathrm{mig}\,y)$, and at the
farthest, $(\mathrm{mag}\,x, \mathrm{mag}\,y)$. The lower bound is the
value itself where it is a double: writing $a$, $b$ and $h=\sqrt{a^2+b^2}$
as odd integers times powers of two, $a^2+b^2=h^2$ can only hold where $h$
has the power of two of the one of $a$, $b$ whose power is the smaller, and $h^2-b^2=a^2\cdot4^d$
is then compared exactly on 107 bits.\newinvfive{} \code{rsqrt(x)} is
CORE-MATH's \code{rsqrt} at the bounds of $x\cap\itv{0}{+\infty}$,
decreasing, the value at $0$ being the limit $+\infty$, and $x$ holding no
positive number giving $\emptyset$; it is exact at the powers of
$4$.\newinvfive{} Both are tightest. The tests find \code{hypot} tightest
over 200\,000 boxes, among them Pythagorean triples scaled by powers of two
and the doubles next to them, and \code{rsqrt} over 40\,000 intervals.

\paragraph{Compound minus one.}\findex{compoundm1}%
\code{compoundm1(x, y)} computes $(1+x)^y-1$ with the analysis of the box
of \code{pow} for the base $1+x$: the value at the corner each bound is
at, since $(1+x)^y-1$ increases with $y$ for $x>0$ and decreases for
$x<0$, increases with $x$ for $y>0$ and decreases for $y<0$, and is $0$ on
$x=0$ and on $y=0$. The limits $-1$ and $+\infty$ are the bounds at
$x=-1$, at $x=+\infty$ and for an infinite $y$, and the empty set is the
result for $x<-1$ and for $x=-1$ with $y\leq0$, where the function has no
value. CORE-MATH has no \code{compoundm1}, and
$\exp(y\,\mathrm{log1p}(x))-1$ computed with doubles multiplies the
relative error of $\mathrm{log1p}$ by $|y\,\mathrm{log1p}(x)|$: up to
1\,100 doubles from the value near the overflow, over 300 draws with the
\code{expm1} and \code{log1p} of Python. The value at a corner is
therefore computed in \file{gaol/math/gaol\_compoundm1.h} with integers
only. A fast phase in 64-bit fixed point comes first, whose error bound is
proved (27 to 200 units of $2^{-63}$ relatively, 4.6 to 10 times the worst
error found); it settles 95\,\% of the pairs for $|y\,\mathrm{log1p}(x)|$
below $0.35$, 92\,\% from $0.35$ to $1$, 90\,\% from $1$ to $2$, 83\,\%
from $2$ to $4$ and 71\,\% from $4$ to $8$, and none past, nor for
$|\mathrm{log1p}(x)|$ above $8$, $x<-0.99966$ or $x>2\,980$. The other
pairs go on to a 128-bit significand in balls, which add the error of each
operation to a bound of the error. Both phases start from the
$\mathrm{log1p}(x)$ of CORE-MATH, corrected by one step of Newton's method
($d=(1+x)\exp(-y_0)-1$, then $y_0+d-d^2/2+d^3/3$), multiply it by $y$,
then take $\exp-1$ by a series of 27 terms after the reduction by a
multiple of $\log2$. The ball is no wider than $2^{-111}$ relatively for
$|y\,\mathrm{log1p}(x)|$ below $746$ ($2^{-118}$ below $10$, $2^{-120}$
below $1$), and the bounds are the doubles around it. A corner takes
$0.1\,\mu$s where the fast phase settles it and $0.6$ to $0.9\,\mu$s in
the balls, on an Intel i7-1185G7 with GCC 15: $0.12\,\mu$s for a point and
$0.23\,\mu$s for a box of $x$ between $-0.3$ and $0.5$ and $y$ between
$0.1$ and $3$, instead of $2\,\mu$s with the balls alone.

The cases where the value is a double are told apart. For $y=-1$ the value
is the quotient $-x/(1+x)$ of two integers, which is divided exactly, and
the value is the double where the division is exact. The value is the
double $d$ where $(1+x)^y=1+d$, which integers prove for $y=P/2^k$:
$s=1+x$ a power of two with $y\log_2 s$ an integer, or $s$ the $2^k$-th
power of an integer $R$ times a power of two with $R^P=1+d$
($2.375=1.5^3-1$, $1.25$, $0.5=\sqrt{2.25}-1$, $2^m-1$ for $|m|\leq53$);
it is $y\,x$ on the side $y(y-1)$ tells where $y\,x$ is a double and
$x(|y|+2)\leq1/2$, the ball not telling the side where $x\,|y-1|$ is below
$2^{-119}$; and $2^m-1$ is between the double below $2^m$ and $2^m$ for
$m$ past 53, which the ball cannot tell from $m=118$ on.\newinvfive

The function is tightest, but where the value is within $2^{-111}$
relatively of a double that is not the value, none of the cases above
tells it and the fast phase, which tells apart the doubles that are
$2^{-55}$ to $2^{-58}$ apart relatively or more, has not settled the pair:
about one pair $(x, y)$ in $10^{17}$ drawn at random for the balls, 5\,\%
of the pairs going to them, and none among the 5\,854 pairs of the tests;
the result is then one double wider on that side. The contrived values
$y=-1$ and $x=\pm(2^e+2^{2e-1})$ for $e$ from $-51$ to $-40$, where
$1/(1+x)=1-x+x^2-x^3\ldots$ has its cubic term cancelled and is within
$2^{-155}$ of a double, are the tightest ones now, with the division. The
tests find the function tightest over 5\,854 pairs $(x, y)$ of every
magnitude, among them the values that are doubles, the tiny $x$ down to
the subnormals, $y=-1$ (719 values of $x$, the contrived ones among them,
633 of them divided), the edges of the fast phase
($|\mathrm{log1p}(x)|=1/2$ and $8$, $y\,\mathrm{log1p}(x)=8$ and
$\log(2)/2$), the thresholds of the reduction and of the cuts, and 850
boxes (the definition read from the hull of the corners), against mpmath at
1\,500 to 7\,000 bits and fractions; the balls of $\mathrm{log1p}$ and
$\exp-1$ are checked against mpmath to 256 bits.

\subsection{The error functions}\label{sec:accuracy-erf}

IEEE 1788-2015 does not name \code{erf} and \code{erfc}. GAOL computes them
with CORE-MATH's \code{erf} and \code{erfc} at the bounds, correctly
rounded upward, \code{erf} being increasing and \code{erfc} decreasing: the
value at one bound, and the double below the value at the other one unless
that value is a double, which it is at $0$ ($\mathrm{erf}(0)=0$,
$\mathrm{erfc}(0)=1$) and at $\pm\infty$ (the limits $\pm1$ of
$\mathrm{erf}$, $2$ and $0$ of $\mathrm{erfc}$); the results are within
$\itv{-1}{1}$ and $\itv{0}{2}$. $\mathrm{erf}(x)$ is within $2^{-53}$ of
$\pm1$ from $|x|=5.87$ on, where its bounds are $1-2^{-53}$ and $1$, or
$-1$ and $-(1-2^{-53})$; $\mathrm{erfc}(x)$ is subnormal from $x=26.55$
on, and below $2^{-1074}$ from $x=27.22$ on, where its lower bound is $0$.
Neither function is known to take a value that is a double at any other
double, and the sources of CORE-MATH treat no such case.\newinvfive{} Both
are tightest. The tests find them so at the 409 doubles of
\file{elementary\_values.h} (tiny arguments, changes of binade, the
arguments where the values round to $\pm1$, $2$ or $0$ or become subnormal,
and the hard cases of CORE-MATH's \file{erfc.c}) and over the intervals
between them, against mpmath, and over 9\,225 doubles and 80\,000 intervals
against CORE-MATH called in both directed roundings.

\subsection{Reverse functions}\label{sec:accuracy-reverse}
\cindex{reverse functions!accuracy}%

The relational functions of GAOL compute \code{f\_rel(J, I)}, the hull of
the $x$ of $I$ whose image is in $J$: they are the reverse functions of
Table~10.1 of the standard with the arguments in another order
(\pxref{sec:relational-arithmetic}). Each computes the preimage of $J$
with the functions above, rounded outward, and intersects it with $I$: it
is valid, which the standard requires. The test \file{tests/reverse.cpp}
checks, on the cases of the minimal tests of libieeep1788 and on random
ones, that they are also accurate, which the standard recommends for
inf-sup types, the hulls being computed from the definition of each
function, apart from GAOL (\pxref{sec:test-programs}); the test
\file{tests/other\_functions.cpp} checks that they keep the values they
are given, within the numbers of doubles below.

\begin{itemize}
\item[$\bullet$] \textbf{Square root.} \code{sqrt\_rel(c, x)}
  (\ieee{sqrRev}) is accurate; the tests find it accurate, and within 1
  double of the value it keeps.
\item[$\bullet$] \textbf{Absolute value.} \code{invabs\_rel(c, x)}
  (\ieee{absRev}) is tightest: its bounds are those of $J$, of $-J$ and of
  $I$.
\item[$\bullet$] \textbf{Integer power, $p>0$.}
  \code{nth\_root\_rel(c, p, x)} (\ieee{pownRev}) is accurate, its bounds
  being the roots of \code{nth\_root()}; beyond the \code{unsigned int}s it
  is the hull of the parts of $x$ in those roots and, for an even $p$, in
  their opposites.\newinvfive{} The tests find it accurate, within 1 double
  of the tightest, and within 2 doubles of the value it keeps.
\item[$\bullet$] \textbf{Integer power, $p\leq0$.}\newinvfive{} For $p=0$,
  \code{nth\_root\_rel(c, 0, x)} is $x$ where $1$ is in $c$, and
  $\emptyset$ elsewhere: exact. For $p<0$, $c$ is cut at $0$, and the part
  at least $0$ and, for an odd $p$, the part at most $0$ each give the
  roots \code{nth\_root(part, p)}, the root then the tightest division, up
  to 4 doubles beyond the tightest bounds; the result is the hull of their
  parts in $x$, $0$ left out. The hull of the roots of $c$, which joins its
  parts, would keep the part of $x$ between them, and the roots reach $0$
  where $c$ is unbounded, $0$ having no power of negative exponent. The
  result is valid, and within 2 doubles of the accurate bounds of 12.10.1, one
  more than the trigonometric functions: the tests find it $\emptyset$ where
  the hull is $\emptyset$ in the cases written out, and of the 18\,052 random
  cases with $p<0$ of
  \file{tests/reverse\_values.py}, 12\,495 have a nonempty result with
  nonempty accurate bounds, 521 of them one double beyond, 25 two, none
  more. It is at most 2 doubles from the tightest bounds on the table of
  \file{tests/reverse.cpp}, 3 on those random cases (60 of them), and 4 at
  most, found for $c=\itv{2.25}{2.25}$ and $p=-127$. Where $p=-1$ and
  $1/c$ reaches beyond the largest double, $c$ being next to $2^{-1074}$, a
  bound saturates at $\pm$\code{DBL\_MAX}, as \code{div\_rel} does: valid,
  but beyond 4 doubles of the tightest, and not empty where the hull is:
  the result for $c=\itv{2^{-1074}}{2^{-1074}}$, $p=-1$ and
  $x=\itv{0}{\code{DBL\_MAX}}$ is $\itv{\code{DBL\_MAX}}{\code{DBL\_MAX}}$.
\item[$\bullet$] \textbf{Trigonometric functions.} \code{asin\_rel},
  \code{acos\_rel} and \code{atan\_rel} (\ieee{sinRev}, \ieee{cosRev},
  \ieee{tanRev}) are accurate, or one double beyond: the pieces of the
  preimage are $k\pi$, enclosed within about one double from $\pi$ in
  double-double, plus or minus the inverse function of $J$, which is added
  before the sum is rounded; beyond $2^{52}$ a bound of $x$ is kept, and an
  $x$ of a single double is decided by the image of the function. The tests
  find them accurate but for one double, with the bounds of $x$ beyond
  $2^{52}$, and within 6 doubles (4, 5 and 3 found) of the value they keep
  from $1$ to $2^{50}$, away from the points where the inverse function
  magnifies the width of the image.
\item[$\bullet$] \textbf{Hyperbolic functions.} \code{acosh\_rel(c, x)}
  (\ieee{coshRev}) is accurate; the tests find it within 16 doubles of the
  value it keeps. \code{asinh\_rel} and \code{atanh\_rel}, named after
  \ieee{coshRev} (\ieee{sinhRev} and \ieee{tanhRev} are not in Table~10.1),
  are accurate; \code{atanh\_rel(J, x)} is $\emptyset$ where $J$ holds no
  point of $(-1, 1)$, the values of the hyperbolic tangent.\newinvfive{}
  The tests find them within 10 and 26 doubles of the value they keep, and
  $\emptyset$ for $J=\itv{1}{1}$ and $\itv{-1}{-1}$.
\item[$\bullet$] \textbf{Division.} \code{div\_rel(c, b, x)}, and
  \code{c \% b} for $x=\itv{-\infty}{+\infty}$ (\ieee{mulRev}), are
  accurate, and tightest where the bounds of the quotients are doubles. The
  tests find them accurate, the tightest over the cases of \code{\%}, and
  \code{div\_rel} within 2 doubles of the value it keeps.
\item[$\bullet$] \textbf{Powers.}\newinvfive{} \code{pow\_rel(c, b, x)} and
  \code{pow\_exponent\_rel(c, a, x)} (\ieee{powRev1}, \ieee{powRev2}) are
  tightest. The pieces of the solutions over the signs of the other
  argument ($y>0$, $y<0$; $x>1$, $x<1$), with $0^y=0$ for $y>0$, $x^0=1$
  and $1^y=1$, end at $c^{1/y}$ or $\log(c)/\log(x)$, $c$, $x$ and $y$
  being bounds of the arguments; their bounds are the largest double $d$
  with $d^y$ (or $x^d$) proved by CORE-MATH's \code{pow}, correctly
  rounded, to be on one side of $c$, and the smallest one on the other
  side. A bound of $x$ where the solutions start between it and the next
  double, and whose power is within one double of $c$, is kept, which only
  knowing whether that power is a double would settle: the result is then
  the tightest one for $x$ one double wider. The tests find them tightest
  on the 804 cases of ITF1788 and on 1\,244 others, random or whose
  solutions are doubles, within the tightest result for $x$ one double
  wider on 966 cases where a bound of $x$ is next to an end of the
  solutions, and keeping the 54\,305 points of 6\,000 random boxes whose
  power is within $c$.
\item[$\bullet$] \textbf{Maximum, minimum, sign and floor.}\newinvfive{}
  \code{max\_rel(z, y, x)}, \code{min\_rel(z, y, x)},
  \code{sign\_rel(z, x)} and \code{floor\_rel(z, x)}, whose reverse
  functions Table~10.1 does not name, are tightest: their bounds are those
  of $x$ and $z$, $0$, $\lceil\leftBound{z}\rceil$ and
  $\lfloor\rightBound{z}\rfloor+1$, the open ends of the sets of
  \code{sign\_rel} ($0$) and \code{floor\_rel}
  ($\lfloor\rightBound{z}\rfloor+1$) included, the last rounded upward
  beyond $2^{53}$. The tests check them against the definitions on every
  interval of half-integer bounds, on cases written out and by sampling,
  and \code{floor\_rel} at every magnitude, against
  $\lfloor\rightBound{z}\rfloor+1$ computed exactly.
\item[$\bullet$] \textbf{Two-argument arc tangent.}\newinvfive{}
  \code{atan2\_rel\_y(c, b, x)} and \code{atan2\_rel\_x(c, a, x)}
  (\ieee{atan2Rev1}, \ieee{atan2Rev2}) are accurate. The angles of $c$ in
  $(-\pi, \pi]$ are cut at the axes into four cones, whose sides have the
  slopes $\tan\theta$ at the bounds $\theta$ of $c$ (CORE-MATH, correctly
  rounded, the slope of a subnormal $\theta$ being $\theta$ itself), and
  $0$ or $+\infty$ at the axes; the bounds of the box, and their products
  and quotients with the slopes, are rounded outward. The result is within
  2 doubles of the tightest where both are nonempty; a box within one
  rounding of the cone may give a point that the tightest enclosure, the
  empty set, does not hold. The tests find them within 2 doubles of the
  tightest over 1\,969 cases (1\,263 nonempty) against mpmath at 3\,000
  bits, around the multiples of $\pi/4$, the axes, the half-line $y=0$,
  $x<0$, and with zero, subnormal and infinite bounds, and keeping every
  point whose angle is in $c$, over about 600\,000 such points of the
  boxes, among them points of their sides next to the sides of the cones.
\end{itemize}

\begin{Rem}\label{rem:accuracy-reverse-point}
Where a bound of \var{J} is one double from the image of a bound of the
preimage, the reverse functions keep a point of \var{I} whose image is
just outside \var{J}:
\code{sqrt\_rel(\itv{1+2^{-52}}{+\infty}, \itv{-1}{1+2^{-52}})} is
\itv{-1}{1+2^{-52}}, whose tightest enclosure is \itv{1}{1+2^{-52}}. They
are accurate all the same, the widened \var{J} of nextOut holding that
image.
\end{Rem}

\medskip
\ieee{mulRevToPair(b, c)} (10.5.5) is the pair of the two parts, each a
\code{div\_rel} of $c$ by a part of $b$ over $\itv{-\infty}{+\infty}$, the
part of $b$ at most $0$ and the part at least $0$:\newinvfive{} as tight as
\ieee{mulRev}, so the tightest where the bounds of the quotients are
doubles. The test \file{tests/ieee1788.cpp} checks it against exact
quotients, and the 172 cases of ITF1788 pass.

\ieee{cancelMinus(x, y)} and \ieee{cancelPlus(x, y)} (10.5.6) are
\code{cancel\_minus(x, y)} and \code{cancel\_plus(x, y)}, the tightest,
which 12.12.5 requires. Whether the result exists depends on $x$ being at
least as wide as $y$, which rounding the differences of the bounds cannot
decide when they are within a double of each other; each difference is
computed exactly, as the sum of its value rounded to nearest and of its
error, by the TwoSum of Knuth. The result is $\itv{-\infty}{+\infty}$
where there is no value, and $\emptyset$ for an empty $x$ and a bounded
$y$. They are checked over 200\,000 pairs against exact rational
arithmetic, 11\,524 of them of exactly equal widths.\newinvfive

\subsection{Set operations and numeric functions}\label{sec:accuracy-set}

The set operations, the constructors, the numeric functions and the
boolean functions are computed from the bounds, and most of them are
exact. The relations are described in Chapter~\ref{chap:relations}, the
reading and the writing of intervals in Chapter~\ref{chap:io}; here are
their accuracies.

\begin{itemize}
\item[$\bullet$] \textbf{Intersection and hull.} \code{x \& y} and
  \code{x | y} (\ieee{intersection}, \ieee{convexHull}) are taken on the
  bounds: exact.
\item[$\bullet$] \textbf{Constructors.} \code{interval(l, u)} and
  \code{interval(d)} (\ieee{numsToInterval}) store the bounds, and the
  empty set where $l>u$, $l=+\infty$, $u=-\infty$ or a bound is NaN:
  exact.
\item[$\bullet$] \textbf{Reading.} \code{gaol::textToInterval(s)} and
  \code{operator>{}>} (\ieee{textToInterval}) read the numbers, decimal,
  hexadecimal, octal or binary, compared exactly with the doubles around
  them; the rational numbers $p/q$ are divided as intervals; the uncertain
  form is read as its exact decimal bounds; the literals \code{[ ]},
  \code{[empty]}, \code{[entire]}, the infinite and the missing bounds, in
  any case of letters, are read (\pxref{sec:reading-intervals}). GAOL also
  reads expressions, such as \code{"sin(1)+0.1"}, computed with the
  operations above. The reading is tightest for the literals of IEEE
  1788-2015, $p$ and $q$ being doubles. The literal \code{[l, u]} is empty
  where $l$ is above $u$, numbers that are not doubles being compared
  exactly, but for \code{pi} and another number between the same two
  doubles, a number far beyond the doubles, or bounds that are computed,
  which give the interval from the lower bound of $l$ to the upper bound of
  $u$.\newinvfive{} It is valid for expressions. The tests find it
  tightest.
\item[$\bullet$] \textbf{Bounds.} \code{left()} and \code{right()}
  (\ieee{inf}, \ieee{sup}) return the stored bounds, and NaN for the empty
  set, where IEEE 1788-2015 has $+\infty$ and $-\infty$, and a zero lower bound of either sign, where \ieee{inf} gives $-0$ (12.12.8), as \code{gaol\_ieee1788::inf()} does: exact.
\item[$\bullet$] \textbf{Midpoint.} \code{midpoint()} (\ieee{mid}) is
  $(l+u)/2$ rounded to nearest, ties to even, or $l/2+u/2$ where $l+u$ may
  overflow, a bound being $2^{1023}$ or more in magnitude, which is tested
  first, so that $\itv{\code{DBL\_MAX}}{\code{DBL\_MAX}}$ raises no
  overflow exception;\newinvfive{} it is $0$, $-$\code{DBL\_MAX} and
  \code{DBL\_MAX} for the unbounded intervals, and NaN for the empty set:
  as 12.12.8 requires. \code{mid()} is the same midpoint rounded outward:
  tightest.
\item[$\bullet$] \textbf{Width and radius.} \code{width()} (\ieee{wid}) is
  $u-l$ rounded upward, NaN for the empty set; \code{rad()} and
  \code{mid\_rad(m, r)} (\ieee{rad}, \ieee{midRad}) give the greater of
  $m-l$ and $u-m$ rounded upward, $m$ being \code{midpoint()}, $+\infty$
  for the unbounded intervals and NaN for the empty set: as 12.12.8
  requires. \code{mag()}, \code{mig()} and \code{smig()} are taken on the
  bounds: exact. \code{hausdorff(x, y)} is the distance rounded upward:
  tightest.
\item[$\bullet$] \textbf{Relations.} \code{is\_empty()},
  \code{is\_a\_double()} and \code{set\_contains(d)} (\ieee{isEmpty},
  \ieee{isSingleton}, \ieee{isMember}), the set relations \code{set\_eq},
  \code{set\_contains} and \code{set\_leq}, \code{set\_strictly\_contains}
  and \code{set\_le}, \code{set\_disjoint} (\ieee{equal}, \ieee{subset},
  \ieee{interior}, \ieee{disjoint}, as Tables~10.3 and 10.4), the certainly
  relations \code{certainly\_leq}, \code{certainly\_le},
  \code{certainly\_geq} and \code{certainly\_ge} (\ieee{precedes},
  \ieee{strictPrecedes}, true when an interval is empty),
  \code{is\_entire()} and \code{is\_common\_interval()} (\ieee{isEntire},
  \ieee{isCommonInterval}, a common interval being nonempty and bounded,
  the NaN bounds of the empty set being no infinities\newinvfive) and
  \code{less(y)} and \code{strictly\_less(y)} (\ieee{less},
  \ieee{strictLess}, as Table~10.3, \ieee{strictLess} counting an infinite
  bound as below itself, with the empty cases of Table~10.4\newinvfive) are
  decided on the bounds: exact (Chapter~\ref{chap:relations}). The tests
  find them exact, and \code{less} and \code{strictly\_less} as Tables~10.3
  and 10.4 over 20\,000 pairs.
\item[$\bullet$] \textbf{Overlap.} \code{x.overlap(y)} and
  \code{gaol\_ieee1788::overlap(x, y)} (\ieee{overlap}) give the state of
  Table~10.7 among the 16 of \code{overlapping\_state}:
  \element{bothEmpty}, \element{firstEmpty} or \element{secondEmpty} for an
  empty operand, else the order of the four bounds, told as the integers of
  their bits (\code{gaol\_detail::bound\_key()}: $-0$ and $+0$ are the same
  bound, and an infinite bound is the same as itself); a point interval on
  a bound of the other one is \element{starts}, \element{finishes},
  \element{finishedBy}, \element{startedBy} or \element{equals}, never
  \element{meets} nor \element{metBy}.\newinvfive{} Exact. The tests read
  the set specification of Table~10.7 with quantifiers over a finite model
  of the real line, on 31\,449 pairs (the empty set, the 104 intervals of
  a grid of bounds, infinities, zeros of both signs, subnormals and
  neighbours of a double among them, and 20\,000 random pairs whose bounds
  meet); they check \ieee{disjoint}, \ieee{equal} and \ieee{subset} as
  masks of the states, the converse state when the operands are exchanged,
  the 48 cases of ITF1788 without decorations, and the same state with a
  mode that flushes the subnormals to zero.
\item[$\bullet$] \textbf{Writing.} \code{gaol\_ieee1788::intervalToText(x)}
  writes the bounds \code{[l, r]}, \code{[a]} for a point interval whose
  digits are exact (\code{[4]}, and \code{[0]} for a zero of either sign),
  or \code{[empty]}, in decimal with the digits of
  \code{interval::precision()} and a decimal point, whatever the global
  format and the locale.\newinvfive{} \code{operator<{}<} in the format
  \element{bounds}, and in \element{agreeing} for the intervals it writes
  with their bounds (those whose bounds share no nonzero leading digit, and those whose bounds are not both finite, nonzero and of the same sign), writes the same, with the
  digits and the flags of the stream, and a decimal point whatever its
  locale, so that the reader reads the text back; \element{agreeing}
  writes the other intervals with the digits their bounds share
  (\pxref{sec:output-format}). Each bound is written to nearest by the C
  library, compared exactly with the bound, and its last digit moved
  outward when it is on the wrong side. The format \element{hexa} writes
  the bounds in the hexadecimal-significand form of 13.4.1, \code{[a]} for
  a point interval and \code{[0x0p+0]} for a zero. The decimal formats are
  valid, and the tightest with the digits asked for where the C library
  rounds to nearest as it should; the hexadecimal format is exact. The
  tests find less than one unit of the last digit in decimal and the exact
  bits in hexadecimal, read the text of \ieee{intervalToText} back as a
  literal, over every format, precision and a comma locale, and read the
  text that \code{operator<{}<} writes for a point under that locale back
  as the point, or as an interval containing it.
\item[$\bullet$] \textbf{The formats for the eye.} \code{operator<{}<} in
  the formats \element{width} and \element{center} writes the
  \code{midpoint()} to nearest and the \code{rad()} rounded
  upward,\newinvfive{} as \code{c (+/- w)}, or \code{c} for a point
  interval, and \code{[empty]} for the empty set; the digits of $c$ are
  those of the precision, so that the text need not contain the interval.
  As doubles, $\itv{c-w}{c+w}$ contains $x$, $w$ being the smallest radius
  that does so, and a nonzero radius is never written $0$. The tests find
  the radius written at least the exact half-width and \code{rad()} rounded
  upward with less than one unit of its last digit, $\itv{c-w}{c+w}$
  containing $x$, and the midpoint written to nearest with at most half a
  unit of its last digit, over the special values, the subnormals, the
  largest doubles and random intervals, every precision and flag.
\item[$\bullet$] \textbf{Recovery.} \code{exact\_string(x)}, which
  \code{operator<{}<} writes with \element{hexa}, and
  \code{gaol::textToInterval(s)} (\ieee{intervalToExact},
  \ieee{exactToInterval}) meet the recovery requirement of 13.4: an
  interval written in hexadecimal and read again gives the same bounds, bit
  for bit, the empty set, the infinite bounds, a bound $-0$ and the
  subnormals included, but for a point interval of zero, written
  \code{[0x0p+0]} whatever the signs of its bounds and read back as
  $\itv{+0}{+0}$, the same set.\newinvfive{} Exact, over 2\,000 random
  intervals and the values written apart.
\end{itemize}

\section{The same bounds on every machine}\label{sec:accuracy-same-bounds}
%%------------------------------------------------------------------------
\cindex{accuracy!the same bounds on every machine}%
\cindex{zero!sign of a zero bound}%

GAOL is built with CORE-MATH only, whose sources are compiled into the
library (\pxref{sec:math-libraries}): the bounds of this chapter are the
bounds on every machine. They are therefore those of every build of GAOL,
and the test \file{tests/core\_math.cpp} checks them against CORE-MATH
called in the downward and the upward rounding, which are the tightest
bounds there are.

\medskip
The SSE2, the NEON and the FPU intervals (\pxref{sec:kernels}) compute the
integer powers with the same code, from exact products
(\code{pown\_products()} of \file{gaol/core/gaol\_interval\_core.cpp}), and
the square from one product per bound, and give the same bounds. The SSE2
intervals of GAOL~4 multiplied rounded products from the highest bit of
$n$, and their bounds and those of the FPU intervals differed for about half
of the intervals with a bound $0$ or infinite or whose power leaves the range
of the exact products.\newinvfive{} \code{pow(x, 2)} is \code{sqr(x)} in every form of
the integer power with the exponent 2 (\code{pow(x, 2)}, \code{2u},
\code{2.0} and \code{[2]}, \code{gaol\_ieee1788::pown(x, 2)}, and
\code{gaol\_ieee1788::pow(x, [2])} on the part of $x$ in
$\itv{0}{+\infty}$, which is $\itv{0}{0}$ for the part
$\itv{0}{0}$):\newinvfive{} the same function, so the same bits, the sign
of a zero bound being that of \code{sqr} in each build. In the earlier versions of GAOL v5, the integer power
made its square from exact products, which has the same bounds and gave a
zero bound of the other sign for part of the intervals holding $0$:
\code{pow(interval(0, 2), 2)} was $\itv{-0}{4}$ with the NEON intervals and
\code{pow(interval(-1, 2), 2)} was $\itv{-0}{4}$ with the FPU ones, where
\code{sqr} gives $\itv{0}{4}$; with the NEON intervals
\code{pow(interval::universe(), 2)} was $\itv{0}{+\infty}$ and is
$\itv{-0}{+\infty}$, as \code{sqr} has it. The exception flags are those
of \code{sqr} as well (\pxref{sec:fp-exceptions}): the corner products
that are not a bound of the result raise the inexact and the underflow
flags too (\code{pow([-1e-310, 1e154], 2)} raises underflow for the square
of its lower bound, which its lower bound $0$ leaves out), and the
overflow flag is raised when the upper bound is $+\infty$ only, the
smaller of two squares overflowing only with the larger: none of 14
million cases of each build had a finite upper bound with it.

\medskip
The bounds are the same, but the sign of a zero bound may differ between
the kernels, $+0$ and $-0$ being the same number.

\begin{Ex}
\code{sqr(interval(-2.0, 3.0))} is $\itv{-0}{9}$ with the SSE2 and the
NEON intervals, and $\itv{0}{9}$ with the FPU intervals.
\end{Ex}

So may the sign of a zero bound of a product or of a quotient. The SSE2
intervals take the larger of four products, and the quotient of the
corners their signs choose, with no branch, where GAOL~4 chose one product
by the signs of the bounds: \code{[-1, 0] * [-1, 0]} has the lower bound
$+0$ where it had $-0$.\newinvfive{} The FPU intervals take the larger of
four products too (for each bound, the lower one made of products with a
factor negated),\newinvfive{} but with a comparison and a selection that
the compiler compiles as it chooses, not an instruction: where two
products are zeros of opposite signs, the sign of the zero bound of a
product is the compiler's and the target's, $+0$ or $-0$. The lower bound
of \code{[-1, 0] * [-1, 0]} is $+0$ with GCC for 32-bit ARM and for arm64,
with Clang and with Visual C++ for x86, and $-0$ with GCC 13 for POWER9,
which compiles the selection into an \code{xsmaxcdp} that returns its
first operand for a $+0$ and a $-0$. The FPU intervals give the bits of
the SSE2 intervals, checked against a model of their kernels and not by
running them, for the products of finite bounds but for that sign, and for
the quotients by a divisor without zero: \code{[0, 1] / [-2, -1]} has the
upper bound $-0$, where the FPU intervals gave $+0$ before, the zero bound
of the dividend being divided. \code{div\_rel(K, J, I)} of the SSE2
intervals takes the code of the FPU and the NEON intervals where a bound of
$K$ is $2^{500}$ or more in magnitude, or a bound of $J$ that is not $0$
is below $2^{-500}$, the quotients of the vectors being liable to
overflow,\newinvfive{} so that a zero lower bound may have the sign of
that code: \code{div\_rel([0, 0x1.e2ab041732196p-5],
[2\textasciicircum-1074, 0x1.182p+9], [-DBL\_MAX, 0x1.70ef9db22d0e5p+2])}
is $\itv{0}{5.764625}$, a zero lower bound of the other sign than the vector code gives: with a $J$ from $1$, which the vector code takes, the result is $\itv{-0}{\code{0x1.e2ab041732196p-5}}$.

\medskip
The two paths of the SSE2 intervals, with and without the AVX-512 path
(\pxref{sec:avx512}), give the same bits, but
for the sign of a zero bound of \code{x * d}, \code{x *= d} and
\code{d * x} (computed as \code{x * d}) where the zero bounds of $x$ have
different signs, $\itv{+0}{-0}$: it is $\itv{+0}{-0}$ with SSE2 and
$\itv{-0}{+0}$ with AVX-512. The NEON intervals take the larger of four
products too,\newinvfive{} with the maximum of 64-bit ARM, which ranks
$+0$ above $-0$, where the \code{maxpd} of the SSE2 intervals returns its
second operand when the two compare equal: the lower bound of
\code{[-1, 0] * [-1, 0]} is $-0$ with the NEON intervals, where the SSE2
ones give $+0$ and the FPU ones the compiler's, and so may another zero
bound of a product. \code{interval::zero()} is stored
$\langle+0, +0\rangle$ as with SSE2, so that its lower bound is $-0$
(\pxref{sec:kernels}). The tool \file{tests/tools/arith} compares the NEON
and the FPU intervals on a grid of operands, with the calling code in each
of the four rounding directions: it found no other difference.

\section{Discussion}
%%-----------------

This chapter documents the tightness of each operation of \code{interval},
as 12.10.3 of IEEE 1788-2015 requires, from its algorithm, and gives the
largest distance the tests found. The basic operations, the integer
functions and the absmax functions are tightest, as the standard requires.
The elementary functions, the recommended functions and the error
functions are tightest too, where the standard requires validity and
recommends accuracy, because CORE-MATH is correctly rounded in the
direction in effect: the value at a bound and the double below it are the
tightest bounds, and the cases where the exact value is a double are told
apart with integers. The exceptions are stated with their size: the roots
of an order other than 1, 2 and 3 are within 2 doubles, their inverses
within 4, the integer power is one double beyond where $x^n$ is within $2K$ of a double, which is rare but not near $1$, and \code{compoundm1} in about one pair in $10^{17}$.
The reverse functions are valid, and accurate or within a few doubles of the
accurate bounds, as the chapter says of each, which the standard recommends,
and the tests measure how far they are from the values they keep.

\medskip
Several choices trade a little tightness or time for a proof. The real
power proves that its lower bound is a double with integers rather than
call CORE-MATH downward, which would double its cost; \code{compoundm1}
settles most pairs in a fast phase of 64-bit fixed point whose error bound
is proved, and sends the others to balls on 128 bits; the roots of an
order above 3 are proved by the power of the double on each side, at the
price of up to 2 doubles. The bounds are the same on every machine and in
every build, CORE-MATH being compiled into the library and the kernels
computing the integer powers with the same code; only the sign of a zero
bound differs between the SSE2, the NEON and the FPU intervals, and
between the compilers for the FPU ones, which changes no set.

\medskip
What remains is stated as such: the widths of \code{tan} between
$\underline{\pi}$ and $\pi$ are a theoretical case where the result is wider
than the tightest one, and the intervals of floats are not covered. The next chapter
describes how GAOL computes these bounds: the register that holds the
negated lower bound, the SSE2, AVX-512, NEON and FPU kernels, the check of
the rounding direction and its cost, and the comparison with the other
libraries.

%% !TEX root = ../manual_gaolv5.tex
% Chapter 12.
%% Chapter 12: How GAOL computes its intervals. Sources: gaol.tex of f5bf180 lines
%% 1142-1150 and 1369-1423; doc/simd.md (whole); doc/compare/README.md and
%% performance.md (whole); the introductions and summaries of
%% doc/compare/special_cases.md, enclosure.md, itf1788.md and tang2021.md.
\chapter{How GAOL computes its intervals}\label{chap:computing}

The chapters before this one describe what GAOL computes: the operations,
their bounds and their accuracy. This chapter describes how GAOL computes
them, and what each part of the computation costs. An interval library
spends its time in three places: the arithmetic itself, which is one or a
few instructions; the rounding direction, which the processor holds in a
control register that is slow to write; and the elementary functions. Which
of the three weighs most depends on the processor, on the way GAOL is built
and on the program. The measurements of this chapter say how much each
weighs on the machines we had, and they are given in the manner of an
experimental study: each table of measurements names its machine and, where the source gives them, its compiler and its date,
and no number is given that the measurements do not give.

\medskip
In this chapter, we first present the representation of an interval in a
register and the kernels that compute with it on each family of processors,
SSE2, AVX-512, NEON and the floating-point unit, with the rule that makes a
build take the AVX-512 path (Section~\ref{sec:kernels}).
Section~\ref{sec:where-time-goes} profiles the demonstration programs of GAOL
to see where the time of a program goes.
Section~\ref{sec:rounding-by-processor} gives the cost of a write of the
control register on each processor and what each instruction set offers
instead, and Section~\ref{sec:compilers-give} what the compilers give, and
do not give, to a program that needs a rounding direction.
Section~\ref{sec:comparison} compares GAOL v5 with GAOL 4.2.3 and with five
other libraries of interval arithmetic, on the special cases of the standard,
on the enclosure of the elementary functions and on the time of the
operations. Section~\ref{sec:test-suites} reports the run on GAOL v5 of ITF1788,
a suite of test cases for IEEE 1788-2015 written outside the standard, and of
the benchmark of Tang et al., and Section~\ref{sec:new-paths} draws the
verdicts on the paths that were considered and says what a new one has to
show. The chapter closes with a discussion.

%%--------------------------------------------------------------------------
\section{The representation of an interval and the kernels}
\label{sec:kernels}
%%--------------------------------------------------------------------------
\cindex{kernel}\cindex{SSE2}\cindex{NEON}\cindex{AVX-512}%
\cindex{register of an interval}%

On most platforms, switching the rounding direction is costly. The number of
switches is cut down by Proposition~\ref{prop:round-neg}
(\pxref{sec:containment}): GAOL stores the lower bound of an interval negated,
and computes both bounds rounding upward, so that the rounding direction can be
set toward $+\infty$ once and for all.

\begin{Def}\label{def:computing-register}
An interval $\itv{a}{b}$ is kept as the pair \code{<-a, b>}, its lower
bound negated, in one vector of two doubles: a \code{\_\_m128d} on x86 with
the SSE2 intervals, a \code{float64x2\_t} on 64-bit ARM with the NEON
intervals. The FPU intervals keep the same two numbers in two doubles. The
\dfn{kernel} of an operation is the function of the library that computes
the two doubles of its result from the two doubles of each operand, which it
takes and gives in registers, after checking that the rounding direction is
upward and setting it when it is not (the kernels of the AVX-512 path carry
the direction in the instruction, set none, and check only the modes that flush the subnormals to zero, \pxref{sec:avx512}).
\end{Def}

\begin{Ex}
Let $x = \itv{a}{b}$ and $y = \itv{c}{d}$ be kept as \code{<-a, b>} and
\code{<-c, d>}. One vector addition rounded upward gives
$\langle \roundUp{-a-c}, \roundUp{b+d} \rangle$, which is
$\langle -\roundDn{a+c}, \roundUp{b+d} \rangle$ by
Proposition~\ref{prop:round-neg}: the register of $x+y$. The lower
bound is minus the upper bound of the opposite, so that one addition rounded
upward gives both bounds (the accuracy of each operation is the subject of
Chapter~\ref{chap:accuracy}).
\end{Ex}

\medskip
Table~\ref{tab:computing-kernels} says what each build computes with vector
instructions and what stays scalar. The SSE2 intervals are those of x86, the
NEON intervals those of 64-bit ARM, both under the option \option{GAOL\_SIMD}
of CMake (\pxref{sec:neon}); the AVX-512 path is a variant of the SSE2
intervals (\pxref{sec:avx512}); the FPU intervals are those of every other
processor (\pxref{sec:fpu-intervals}).

\begin{table}[htbp]
\centering\footnotesize
\caption[What each build of GAOL computes with vector instructions]{What each build of GAOL computes with vector instructions, and what stays scalar (from \file{doc/simd.md}).} \label{tab:computing-kernels}
\begin{tabular}{|L{3.6cm}|L{4.8cm}|L{4.5cm}|}
\hline
\textbf{Build} & \textbf{Vector kernels} & \textbf{Scalar} \\
\hline
SSE2 (x86, \code{GAOL\_SIMD}) &
\code{+ - * / sqr}, the operations with a double, \code{inverse},
\code{\%=}, \code{-x}, the integer powers with \code{\_\_FMA\_\_}, the
leaves of \code{div\_rel} &
\code{sqrt}, \code{\&=}, \code{|=}, \code{mid}, \code{midpoint}, the
relations, the elementary functions, the reverse functions \\
\hline
AVX-512 path &
\code{+ - * / sqr}, the operations with a double and the integer powers
(with \code{\_\_FMA\_\_}), with the rounding direction in the instruction;
\code{sqrt} as two scalar square roots and \code{/} as two scalar divisions,
with embedded rounding &
as SSE2 but \code{sqrt}; \code{inverse}, \code{\%=} and the leaves of \code{div\_rel} stay the vector kernels of SSE2, which set the direction \\
\hline
NEON (arm64, \code{GAOL\_SIMD}; every compiler) &
\code{+ - * / sqr sqrt}, the operations with a double, \code{\%=},
\code{inverse}, \code{\&=}, \code{|=}, \code{-x}, the integer powers (two
lanes with \code{vfmaq\_f64}) &
\code{mid}, \code{midpoint}, \code{div\_rel}, the relations, the elementary functions, the reverse functions, and \code{sqrt} of an upper bound above $2^{1000}$ \\
\hline
FPU (the rest) & none & everything \\
\hline
\end{tabular}

\end{table}

\begin{Rem}
The two bounds of an interval are one vector of two doubles, and \code{+},
\code{-}, \code{*}, \code{/}, \code{sqr} (and \code{sqrt} where it has a
kernel) are one or a few vector instructions: a wider register does not help
one interval. Only a batch interface, on boxes or arrays of intervals, which
GAOL has not, could use 256 or 512 bits.
\end{Rem}

\subsection{The SSE2 and NEON kernels}

The NEON intervals compute as the SSE2 intervals do on
x86.\newinvfive{} The sum and the difference are one vector addition rounded
upward. The product is the maximum of the four products of the bounds, with
no branch on their signs. The quotient is one vector division when the
divisor has no zero, and single divisions otherwise. With NEON, the square roots of
both bounds are one vector square root rounded upward, the lower bound being
that root where it is exact, and the double below it otherwise. The compilers
that take the NEON intervals, the options that select them and what the user
sees of them are described in Chapter~\ref{chap:rounding} (\pxref{sec:neon}).

\medskip
The products and the quotients were nine cases of the signs of the bounds, and the squares three, told apart by a \code{switch} in GAOL 4.2.3 and in GAOL v5 at
first. When the signs of the operands are random, as those of the benchmark
of Section~\ref{sec:comparison} are, a mispredicted branch costs more than
the operation it chooses. The kernels now compute with no branch on the
signs:\newinvfive{} the larger of four products for a product, one division
of the two lanes by the denominators the signs choose, by masks, for a
quotient, two squares and a mask for a square. The two doubles of each operand
are passed in a register to the function of the library that computes an
operation, which returns the two doubles of its result in one, where the
operators of GAOL 4.2.3 read their operand and wrote their result through
memory, the result going back to memory for the operation that followed.
Each operation still checks that the rounding direction is upward, with one
comparison, and sets it when it is not, where GAOL 4.2.3, in its trust mode, takes it upward on
trust (\pxref{sec:rounding-direction}).

\medskip
The NEON intervals check the direction by reading the control register FPCR
with \code{mrs}, which raises no exception flag, and write it with
\code{msr} only if the direction is not upward with the flush modes clear
(\pxref{sec:ftz}); the FPU intervals add $1+(2^{-1060}+0)$ instead, which
raises the inexact flag (\pxref{sec:fp-exceptions}). \code{msr} on FPCR is a
serializing instruction, slower than \code{mrs}, which is why it is written
when needed only. With \option{GAOL\_PRESERVE\_ROUNDING}, each operation
reads FPCR once and writes back the direction and the modes it found at its
end only if they differ from what the operation left: a program that
computes upward already, with the flush modes clear, pays no write
(\pxref{sec:preserve-rounding}). The bounds of the NEON intervals are those
of the FPU intervals but for the sign of a zero bound (\pxref{sec:zero-sign}),
which the maximum of 64-bit ARM, ranking $+0$ above $-0$, can set otherwise
than the SSE2 and the FPU intervals do, where the \code{maxpd} of x86
returns its second operand. The tool \file{tests/tools/arith} compares the
NEON and the FPU intervals on a grid of operands, with the calling code in
each of the four rounding directions: it found no other difference.

\subsection{When the build takes the AVX-512 path}
\label{sec:computing-avx512-rule}
\findex{GAOL\_PREFER\_AVX512}\findex{GAOL\_PRESERVE\_ROUNDING}%

The AVX-512 path is off in the default build. The build takes it only where
the option \option{GAOL\_PRESERVE\_ROUNDING} is on, where the compiler has the
instructions and where the machine building runs them, and leaves it out
otherwise; the setting of the developers \option{-DGAOL\_PREFER\_AVX512=ON}
asks for it without that option, and \option{=OFF} refuses it
(\pxref{sec:avx512}). A library built with the path needs a processor with the
instructions. The
build decides it the same way with the three build
systems.\newinvfive{} For \code{+}, \code{-}, \code{*} and \code{sqr} the
path puts the two bounds in two of the eight lanes of a 512-bit register,
for the embedded rounding only; \code{/} and \code{sqrt} are scalar
instructions with the embedded rounding, two of them for the two bounds
(\pxref{sec:new-paths}).

\medskip
The rule comes from four builds of one source, \option{GAOL\_PRESERVE\_ROUNDING}
on or off, each with the path (\option{GAOL\_PREFER\_AVX512} on) and without
it, measured by the script \cmd{performance-ab.sh} of \file{.github/scripts/}
with \option{--runs 9 --lines}: \cmd{gaol\_performance} and the long lines of
\file{doc/compare/code/run\_lines.sh}, the two builds run in turn, the median
of 9 runs, in nanoseconds, on an Intel i7-1185G7, whose AVX-512 units are
halved, with Clang 21.1.8, \option{-O3 -mfma}, on 7 November 2026, the
machine doing nothing else. Table~\ref{tab:computing-avx512} gives the
times; the ratio is the time with SSE2 over the time with the path, so that
above 1 the path is faster.

\begin{table}[htbp]
\centering\footnotesize
\caption[The SSE2 intervals against the AVX-512 path]{The SSE2 intervals against the AVX-512 path, with and without \option{GAOL\_PRESERVE\_ROUNDING}, in nanoseconds per operation or per line (median of 9 runs, Intel i7-1185G7, Clang 21.1.8, 7 November 2026).} \label{tab:computing-avx512}
\begin{tabular}{|l|r|r|r|r|r|r|}
\hline
& \multicolumn{3}{c|}{\textbf{Not preserved}}
& \multicolumn{3}{c|}{\textbf{Preserved}} \\
\textbf{Operation} & SSE2 & AVX-512 & Ratio & SSE2 & AVX-512 & Ratio \\
\hline
\code{x + y} & 1.22 & 1.51 & 0.81 & 8.82 & 1.49 & 5.92 \\
\code{x - y} & 1.33 & 1.56 & 0.85 & 9.38 & 1.53 & 6.13 \\
\code{x * y} & 2.75 & 4.45 & 0.62 & 12.14 & 4.46 & 2.72 \\
\code{x / y} & 3.21 & 5.00 & 0.64 & 11.05 & 5.07 & 2.18 \\
\code{sqr(x)} & 1.97 & 2.03 & 0.97 & 8.42 & 2.04 & 4.13 \\
\code{sqrt(x)} & 8.52 & 4.71 & 1.81 & 13.54 & 4.72 & 2.87 \\
\code{pow(x, 3)} & 6.61 & 7.93 & 0.83 & 10.42 & 7.83 & 1.33 \\
\hline
line of 15 \code{+} (\code{sum16}) & 19.67 & 40.66 & 0.48 & 176.0 & 37.92
& 4.64 \\
line of 7 \code{*} (\code{mul8}) & 23.80 & 37.42 & 0.64 & 86.56 & 35.14
& 2.46 \\
line of 7 \code{/} (\code{quot8}) & 26.50 & 46.56 & 0.57 & 86.73 & 43.59
& 1.99 \\
line of 13 \code{+ - * /} (\code{mix13}) & 27.75 & 43.23 & 0.64 & 152.8
& 41.49 & 3.68 \\
\hline
\end{tabular}

\end{table}

The median of the seven operations is 0.83 without
\option{GAOL\_PRESERVE\_ROUNDING}: the operations take 20\,\% longer with
the path, and the ten long lines 72\,\% longer (the table gives four of them;
the others are in the output of the script); their median ratio is 0.58; only
\code{sqrt} gains. It is 2.87 with it: every operation gains, the ten lines by
1.91 to 4.64, with a median of 3.19. The path pays where each operation would otherwise write the
control register twice, and costs where the register is written only if the
direction is not upward already, as it is by default: the rule is therefore
that of \option{GAOL\_PRESERVE\_ROUNDING}. The threshold was a median gain of
5\,\% over \code{+ - * / sqr sqrt} and \code{pown}. One processor was
measured, and not one with full-width AVX-512 units: such a processor is
expected to lose less without \option{GAOL\_PRESERVE\_ROUNDING}, which the
rule does not rely on.

\subsection{CORE-MATH is scalar}
\cindex{CORE-MATH!scalar}\cindex{fused multiply-add}\findex{GAOL\_FMA}%

CORE-MATH \cite{CoreMath:2022}, in \file{3rd/math-core}, is scalar, with no
vector variant; GAOL calls 38 of its binary64 functions, once per bound.
The fused multiply-add, \option{GAOL\_FMA}, pays in these functions and in
the integer powers of GAOL, \code{cos} being 2.5 times faster and
\code{pow(x, 3)} 3.2 times, and does nothing for the arithmetic; the cost of
each build option is given with the installation
(\pxref{sec:install-options-worth}).

%%--------------------------------------------------------------------------
\section{Where the time of a program goes}
\label{sec:where-time-goes}
%%--------------------------------------------------------------------------
\cindex{profile of a program}%

To know which part of GAOL a program waits for, the demonstration programs of
\file{examples/} were profiled. The call counts of
Table~\ref{tab:computing-profile} are exact, from a NEON build with
\option{-finstrument-functions}. The shares of the time are indicative to
about 3 points: one Apple M2 (NEON intervals), the profiler \cmd{sample} of
macOS at 1 ms during 3 s per program, each sample put to the nearest GAOL
function above it, other work running on the machine; the two passes of the
profile differ by up to 3 points, and the table gives the second.

\begin{table}[htbp]
\centering\footnotesize
\caption[Where the time of the demonstration programs goes]{Where the time of the demonstration programs goes: the calls of GAOL and the shares of the time in the arithmetic, in the elementary functions and outside GAOL (Apple M2, NEON intervals, the profiler \cmd{sample} of macOS; from \file{doc/simd.md}).} \label{tab:computing-profile}
\begin{tabular}{|L{2.9cm}|L{2.6cm}|L{2.6cm}|r|r|r|}
\hline
\textbf{Program} & \textbf{Calls} (\code{+ - * /}, \code{sqr}, powers)
& \textbf{Elementary calls} & \textbf{Arith.} & \textbf{Elem.}
& \textbf{Outside} \\
\hline
03 dependency problem & 0.14 M, 26 k & none & 73\,\% & 0 & 16\,\% \\
04 centered form & 10 k & cos 10 k, acos 8 k & 9\,\% & 82\,\% & 6\,\% \\
06 global optimization & 25 M, 5 M, 0.49 M & none & 31\,\% & 0 & 63\,\% \\
08 SIVIA & 1.0 M, 0.5 M & sin 173 k & 14\,\% & 12\,\% & 71\,\% \\
09 estimation & 0.5 M, 68 k & exp 247 k & 20\,\% & 42\,\% & 31\,\% \\
11 ODE & 0.2 M, 8 k, 4.9 k & sin 12 k, cos 3 k, exp 1 k & 29\,\% & 34\,\%
& 30\,\% \\
16 Goldstein-Price & 32 M, 6 M powers & none & 98\,\% & 0 & 2\,\% \\
\hline
\end{tabular}

\end{table}

\begin{itemize}
\item[$\bullet$] \textbf{The solvers.} On the six solvers 06 to 11, the
arithmetic is 28\,\% of the time weighted by the duration of a run (20\,\%
with equal weights), the elementary functions 4\,\% (20\,\%), and the code
outside GAOL (the priority queue, \code{malloc} of the boxes, the program
itself) 62\,\% (45\,\%). The weighted figures are those of 06, which lasts
80\,\% of the total. Inside GAOL, at equal weights, about 40\,\% is
arithmetic, 34\,\% the elementary functions, 10\,\% the reading and writing
of text and 7\,\% \code{mid()}.
\item[$\bullet$] \textbf{The code around CORE-MATH.} The code of GAOL around
CORE-MATH weighs 21 to 66\,\% of the time of an elementary function (median
about 42\,\%): the check of the rounding direction, the reduction of the
argument, the construction of the result. In the code of \code{exp} around
the CORE-MATH call, 43\,\% of the samples fall on the line of the check of
the direction, a \code{mrs} of 0.8 ns: the samples probably land on the
instruction that waits for it. The CORE-MATH calls themselves cost 3 to
13 ns on the M2.
\item[$\bullet$] \textbf{The other operations.} The relations, \code{hull},
\code{intersect}, \code{min} and \code{max} are 3.6\,\% or less in every
program; \code{mid()} is 6\,\% of the global optimization, and
\code{midpoint()} 11 to 12\,\% of the Krawczyk iteration (29 ns on the
M2: two writes of FPCR).
\item[$\bullet$] \textbf{The check of the direction.} The check (one
\code{mrs} of FPCR here), removed in a copy, saves 0.8\,\% (03), 1.4\,\%
(09), 4.7\,\% (06) and 8.3\,\% (16). \option{GAOL\_PRESERVE\_ROUNDING} on the
M2, for a program that computes to nearest, makes 03 and 16 about 16 times
slower, 06 eight times, 09 five times. On x86 with the AVX-512 path the
ratio is not measured for these programs; the operations of the benchmark of Section~\ref{sec:computing-performance} take 0.6 to 1.5 times as long (Table~\ref{tab:computing-times}).
\item[$\bullet$] \textbf{The text.} The text is 41 to 69\,\% of the five
demonstration programs that print (01, 02, 05, 12, 15; 28 to 31\,\% of 14 and
11\,\% of 03), at 0.2 to 0.45 $\mu$s per bound written. Among the solvers it
is 0 to 1\,\% of 06, 08, 09 and 11, 12\,\% of 10 and 22\,\% of 07.
\end{itemize}

\medskip
\noindent On an Intel i7-1185G7, in the benchmark of
Section~\ref{sec:computing-performance} (a million operations on arrays of
intervals, so partly bound by the memory), the time of an operation is 2.2 ns
for \code{+} and \code{-}, 2.30 for \code{*}, 2.52 for \code{/}, 1.87 for
\code{sqr}, 6.30 for \code{sqrt} and 5.12 for \code{pow(x, 3)}, the same
operation on doubles taking 0.96 ns for \code{+}. With
\option{GAOL\_PRESERVE\_ROUNDING} and the AVX-512 path: 2.19, 2.16, 3.42,
3.60, 2.04, 3.75 and 6.18. On the M2, measured with \file{tests/performance.cpp} (the median of nine runs, Table~\ref{tab:install-worth-arm}), in which the i7-1185G7 takes 1.07~ns for \code{x + y} (the best of nine runs,
Table~\ref{tab:install-worth-x86}; 1.22~ns, the median of nine runs, in
Table~\ref{tab:computing-avx512}), the NEON intervals take 0.90 ns for
\code{+}, 1.52 for \code{*}, 2.08 for \code{/}, 1.27 for \code{sqr} and 2.04
for \code{sqrt}, against 1.15, 3.26, 3.22, 1.78 and 5.17 with the FPU
intervals, whose \code{*}, \code{/} and \code{sqr} take no branch on the
signs (1.15, 3.25, 3.98, 2.07 and 4.49 before).

%%--------------------------------------------------------------------------
\section{The rounding direction, by processor}
\label{sec:rounding-by-processor}
%%--------------------------------------------------------------------------
\cindex{rounding direction!cost of a write}\cindex{MXCSR}\cindex{FPCR}%

Except on the AVX-512 path, an operation computes upward after a write of
the control register of the processor when the direction of the program is
not upward already. What a write costs, by family:

\begin{itemize}
\item[$\bullet$] \textbf{x86.} \code{LDMXCSR} takes 1 to 6 cycles of
throughput on the performance cores of Intel and on AMD processors (1.0 on
Zen 5, 6.0 on an Arrow Lake P-core) and 8 on the efficiency cores, according
to uops.info, which measures one instruction at a time and so does not count
the wait for the operations in flight. \code{STMXCSR}, the read, takes 1
cycle on most Intel cores but 14 to 16 on Zen 2 to Zen 5: one more reason,
with those given in \file{gaol/gaol\_fpu.h} (800 ns under Rosetta 2, 80 ns
with 32-bit Visual C++, a single read reused by Clang 18), why the check of
GAOL on x86-64, in the default build, is an addition and not a read. Two writes took 6 ns a
\code{sqr} (8.3 ns, 4.4 times, in the benchmark of
Section~\ref{sec:computing-performance}) on the SSE2 path with \option{GAOL\_PRESERVE\_ROUNDING}, which the AVX-512 path spares.
\item[$\bullet$] \textbf{64-bit ARM.} \code{mrs} of FPCR takes 0.8 ns,
\code{msr} 4 ns for the value FPCR holds and 11 ns when it changes the
direction (Apple M2). The reference manuals of the Neoverse N1, N2, N3 and
of the Cortex-A78 say that a write that changes the control fields
introduces a barrier or a flush.
\item[$\bullet$] \textbf{32-bit ARM, POWER, s390x, LoongArch, MIPS, RISC-V.}
No figure was found in the references.
\end{itemize}

\medskip
\noindent Table~\ref{tab:computing-processors} says what each instruction
set offers instead of a write, read from its reference manual (the versions
are listed at the end of \file{doc/simd.md}): whether an instruction can
carry its own rounding direction, whether it has a fused multiply-add, which
register holds the direction, and what a path of GAOL could be.

\begingroup\footnotesize
\begin{longtable}{|L{2.0cm}|L{3.8cm}|L{2.1cm}|L{2.0cm}|L{2.9cm}|}
\caption[The rounding direction on each processor family]{The rounding direction on each processor family: what the
instruction set offers, and what a path of GAOL could be (from
\file{doc/simd.md}).}\label{tab:computing-processors}\\
\hline
\textbf{Processor} & \textbf{Direction carried by the instruction}
& \textbf{FMA} & \textbf{Register} & \textbf{What a path could be} \\
\hline
\endfirsthead
\multicolumn{5}{l}{\textit{Table~\thetable, continued}}\\*
\hline
\textbf{Processor} & \textbf{Direction carried by the instruction}
& \textbf{FMA} & \textbf{Register} & \textbf{What a path could be} \\
\hline
\endhead
\hline
\endfoot
\endlastfoot
x86-64 SSE2 & no & FMA3, a set of its own: AVX2 does not imply it
& MXCSR (RC, FTZ, DAZ) & the SSE2 path of today \\
\hline
x86-64 AVX-512F & yes, \code{\{er\}}, for the scalar instructions and the
512-bit ones & yes & MXCSR: RC and the exception flags are ignored, FTZ and
DAZ are honoured & the AVX-512 path of today \\
\hline
x86-64 AVX-512VL, AVX10.1, AVX10.2 & no at 128 and 256 bits:
\code{VADDPD ymm} has no \code{\{er\}} in the Intel manual; AVX10.2 had it
for 256 bits in revisions 1 to 3 and removed it in revision 4.0 (8 May
2025), and revision 7.0 (6 April 2026) keeps it for 512 bits only & yes &
MXCSR & none; a vector of two doubles has to go through the 512-bit form,
as the AVX-512 path does \\
\hline
AArch64, NEON, SVE, SME & no: only \code{FCVT*}, \code{FRINT*} and
\code{FCVTX*} name a rounding in the A64 instruction set of March 2026;
\code{FADD}, \code{FMUL}, \code{FDIV}, \code{FSQRT}, \code{FMLA} use FPCR,
scalar, Advanced SIMD and SVE & yes & FPCR (RMode, FZ; FEAT\_AFP adds FIZ,
AH and NEP) & the NEON path of today; SVE has 128-bit vectors on the
Neoverse N2, N3, V2, V3 and the Cortex-A720, 256 on the V1, and adds no
rounding \\
\hline
ARMv7 & no; NEON has no double, and always rounds to nearest & VFPv4 only
& FPSCR (RMode, FZ) & the FPU intervals \\
\hline
POWER (VSX) & no: \code{xsadddp} and \code{xvadddp} use FPSCR.RN; the
\code{RO} bit exists for the quad precision only (round to odd) & yes &
FPSCR (RN, NI) & VSX \code{<-l, r>} with \code{mffscrni}, needs hardware to
measure \\
\hline
s390x & no: \code{ADB} and \code{VFA} have no rounding field; the \code{M3}
field of the conversions is not the arithmetic & yes & FPC (BRM) &
\code{VFADB}, needs hardware to measure \\
\hline
RISC-V, scalar (D) & yes: the \code{rm} field of \code{FADD.D},
\code{FMUL.D}, \code{FDIV.D}, \code{FSQRT.D}, \code{FMADD.D}\ldots;
\code{rup} is 011 & yes & frm (no flush-to-zero mode) & two scalar
instructions, no check, no register to restore; inline assembly only \\
\hline
RISC-V, vector (V) & no: all the vector floating-point operations use frm;
the \code{\_rm} intrinsics compile to \code{fsrmi} and \code{fsrm} around
the operation & yes & frm & none \\
\hline
LoongArch & no & yes & FCSR & none \\
\hline
MIPS & no (read from the Linux headers and the compiler output, not from the
manual) & not verified & FCSR (RM, FS) & none \\
\hline
WebAssembly 3.0 & no: ``Non-default directed rounding attributes are not
supported'' & \code{relaxed\_madd} may not be fused & none & none: only the
rounding to nearest \\
\hline
\end{longtable}
\endgroup

The RISC-V scalar path is the one case of a rounding direction by
instruction that GAOL has not. Its interest is that it needs no check and no
restore, so that \option{GAOL\_PRESERVE\_ROUNDING} would cost nothing, as on
the AVX-512 path. No figure exists for it: no RISC-V board is available. The
boards read about, in secondary sources, are the SpacemiT K1, which has the
vector extension 1.0, the SiFive P550, which has none, and the T-Head C910,
which has the draft 0.7.1.

%%--------------------------------------------------------------------------
\section{What the compilers give}
\label{sec:compilers-give}
%%--------------------------------------------------------------------------
\cindex{FENV\_ROUND@\texttt{FENV\_ROUND}}%
\cindex{-frounding-math@\option{-frounding-math}}%

No compiler gives a rounding direction by instruction from C or C++: the
paths of GAOL are written with intrinsics or inline assembly. What the
compilers offer was tried, or read from their output, as follows.

\begin{itemize}
\item[$\bullet$] \textbf{The pragma \code{FENV\_ROUND}.}
\code{\#pragma STDC FENV\_ROUND} (C23) would ask for a static rounding. On
the six targets tried (x86-64, aarch64, riscv64, ppc64le, s390x,
loongarch64), GCC 14 to 16 and its trunk ignore it, with a warning only
under \option{-Wall} (GCC bug 20785, open since 2005); Clang 18 to its trunk
parses it, warns ``not supported'' and emits the ordinary instruction (the
IR carries \code{round.upward}; no back end of the three tried, x86-64,
aarch64 and riscv64, uses it); Visual C++ warns C4068. It does not exist in
C++. \option{-Werror=unknown-pragmas} makes Clang's ``not supported'' an
error, but not its other forms of the same pragma (``expected identifier''
with the \code{FE\_UPWARD} of \file{<fenv.h>}, which is a number; ``not
supported on this target''), which belong to another group of warnings, and
\option{/we4068} should do it for Visual C++: none of these was tried.
\item[$\bullet$] \textbf{C++.} P2746R7 (Boehm, January 2025) would replace
\code{fesetround} by a type with an explicit mode, without wording. P4231R0
(May 2026, target C++29) unifies it with P3864; no revision since. Nothing
is in C++26, and none proposes a SIMD form.
\item[$\bullet$] \textbf{The intrinsics of the embedded rounding.}
\code{\_mm\_add\_round\_sd} and \code{\_mm512\_add\_round\_pd} (and the
other operations) compile with GCC 14 to its trunk, Clang 18 to its trunk
and Visual C++ 19.20 to 19.51. The mode must be a literal, and the compilers
do not fold a constant computed with \code{\{ru-sae\}}. There is no
\code{\_mm\_add\_round\_pd}, and no \code{\_mm256\_add\_round\_pd} with GCC
15 and 16 or Clang 21 and later; Clang 20 with \option{-mavx10.2} declares
it and compiles it to the encoding of the withdrawn revisions of AVX10.2,
which the current Intel manual makes \code{\#UD} (never run: no processor
with AVX10.2 is on sale, as far as we found). In inline assembly, the GNU
\cmd{as} of GCC 16.2 on Compiler Explorer assembles
\code{vaddpd \{ru-sae\}, \%ymm1, \%ymm0, \%ymm2} (bytes
\code{62 f1 f9 58 58 d1}, against \code{62 f1 fd 58 58 d1} for the 512-bit
form): no assembly of that form may be written.
\item[$\bullet$] \textbf{\option{-frounding-math}.} GCC does not honour it
across a change of the direction (bugs 34678 and 123482, open): on
\code{fesetround(UP); r = a + b; fesetround(old)}, GCC 14 to its trunk moves
the addition after the restore, on x86-64 and on AArch64. Clang keeps it
between. The vector intrinsics are ordinary floating-point operations for
them: with Apple clang 21 on the M2, at \option{-O2} with
\option{-frounding-math} or \option{-ffp-model=strict}, \code{vaddq\_f64}
and \code{vfmaq\_f64} are right, \code{vpaddq\_f64}, \code{vaddvq\_f64} and
\code{vmulxq\_f64} are computed once, in the first direction, and at
\option{-O2} alone everything is (GAOL uses none of these three). What held
on every compiler tried (GCC 14 to its trunk and Clang 18 to its trunk,
compiled; Apple clang 21, run on the M2) is an empty assembly statement on
each operand after the write of the control register, and on the result, or
a copy of it through volatile memory, before the next one. GAOL does that:
\code{reread\_bounds()} and \code{rnd\_reread()} are empty \code{asm}
statements, \code{GAOL\_RND\_KEEP()} is a copy through volatile memory with
SSE2 and the FPU and an empty \code{asm} with NEON, and
\code{GAOL\_RND\_BARRIER()} is a memory barrier. A path written with
intrinsics needs the same. Visual C++ for arm64 has no \code{asm}
statement: its NEON intervals copy the result through volatile memory,
\code{GAOL\_RND\_BARRIER()} and the writes of FPCR are
\code{\_ReadWriteBarrier()}, and the operands rely on \option{/fp:strict}
(checked by the generated code and by \file{tests/rounding\_direction.cpp},
not by the language).
\item[$\bullet$] \textbf{The intrinsics of the other families.} NEON
\code{float64x2\_t} with GCC, Clang and Visual C++ for arm64 from 19.26; SVE
\file{<arm\_sve.h>} with GCC 14 and Clang 18, and in Visual C++ 19.50 and
19.51 only partially (\code{svmla\_f64\_x} becomes an unresolved call
without a diagnostic); VSX \code{vec\_add} with \option{-mvsx}; RVV with
GCC from 14.1 and Clang from 17.0.1. \code{std::experimental::simd} is in
libstdc++ (libc++ has the types but no \code{operator+}), and
\code{std::simd} of C++26 only in GCC 16, for x86-64.
\end{itemize}

\begin{Rem}
A program that relies on \code{\#pragma STDC FENV\_ROUND} compiles and
rounds in the wrong direction, without an error; so do the
\code{\_mm256\_add\_round\_pd} of Clang 20 and the 256-bit \code{\{ru-sae\}}
assembly. A path of GAOL therefore trusts no rounding direction that the
generated code and \file{tests/rounding\_direction.cpp} have not checked.
\end{Rem}

%%--------------------------------------------------------------------------
\section{GAOL compared with the other libraries}
\label{sec:comparison}
%%--------------------------------------------------------------------------
\cindex{libieeep1788}\cindex{filib++}\cindex{PROFIL/BIAS}%
\cindex{Solaris Studio}\cindex{Boost.Interval}\cindex{comparison of libraries}%

Six implementations of interval arithmetic on doubles are compared in
\file{doc/compare/}: libieeep1788, GAOL v5, filib++, the intervals of
Solaris Studio, PROFIL/BIAS and Boost.Interval.
Table~\ref{tab:computing-libraries} says what each is, in which language it
is written, how it computes its bounds and which arithmetic it follows.
GAOL v5 gives its intervals under two namespaces, \code{gaol}, its own
names, compatible with GAOL~4, and \code{gaol\_ieee1788}, the names of IEEE
1788-2015 \cite{IEEE1788:2015}. They compute with the same functions, and
differ mainly in the power: the \code{pow} of \code{gaol} takes the integer
power for an integer exponent, a negative base included, where the
\code{pow} of \code{gaol\_ieee1788} is the real power of the standard,
defined for $x > 0$ only, and \code{pown} the integer power
(\pxref{sec:pow-namespaces}). The directory \file{doc/compare/code/} holds
the scripts and programs, in C++, Fortran, Python and bash, that run the
three parts of the comparison again and rewrite the tables of the reports.

\begingroup\footnotesize
\begin{longtable}{|L{2.5cm}|L{1.9cm}|L{4.7cm}|L{4.0cm}|}
\caption[The six implementations compared]{The six implementations compared: what they are, their language,
how their bounds are computed and the arithmetic they follow (from
\file{doc/compare/README.md}).}\label{tab:computing-libraries}\\
\hline
\textbf{Library} & \textbf{Language} & \textbf{Bounds}
& \textbf{Arithmetic} \\
\hline
\endfirsthead
\multicolumn{4}{l}{\textit{Table~\thetable, continued}}\\*
\hline
\textbf{Library} & \textbf{Language} & \textbf{Bounds}
& \textbf{Arithmetic} \\
\hline
\endhead
\hline
\endfoot
\endlastfoot
libieeep1788, by Marco Nehmeier (University of W\"urzburg), last commit in
2015 & C++11, header-only & Computed by MPFR, correctly rounded & The
set-based flavor of the preliminary IEEE P1788, the prototype of IEEE
1788-2015 \\
\hline
GAOL v5, by Jordan Ninin (ENSTA), which continues the GAOL of Fr\'ed\'eric
Goualard & C++ & Computed with the rounding direction set upward; elementary
functions with CORE-MATH, correctly rounded upward & Set-based, following
IEEE 1788-2015 in most of its special cases, with the deviations the reports
list \\
\hline
filib++ 3.0.2.2 (University of Wuppertal), as IBEX distributes it, which
IBEX computes with by default on Windows & C++, templates and a small
library & Computed with the rounding direction set and restored by each
operation (\code{native\_switched}); elementary functions of its own & The
extended mode of filib++ (\code{i\_mode\_extended\_flag}, as IBEX uses it),
where the empty set and the infinities are handled \\
\hline
Solaris Studio: the \code{interval(8)} type of Sun's Fortran 95 compiler,
\cmd{f90 -xia}, in Solaris Studio 12.4 (2014) & Fortran 95, intervals built
into the compiler & Computed by \code{libsunimath} & The containment sets of
Sun's interval arithmetic (G.~W. Walster), where the infinities are
values \\
\hline
PROFIL/BIAS 2.0.8 (2009), by Olaf Kn\"uppel and Christian Keil (TU
Hamburg-Harburg) & C++ (PROFIL) over C (BIAS) & Computed by the BIAS
routines, which set the rounding direction downward, then upward, then back
to nearest; elementary functions from the libm, moved outward & The interval
arithmetic before IEEE 1788: no empty set, and an argument outside the
domain of a function is an error that aborts the program \\
\hline
Boost.Interval of Boost 1.92.0 (2026), by Herv\'e Br\"onnimann, Guillaume
Melquiond and Sylvain Pion, the interval library C++ users find first &
C++03 templates, header-only & Computed with the rounding direction set
upward by each operation, the lower bounds negated, and set back after it
(the default policies,
\code{save\_state<}\allowbreak\code{rounded\_arith\_opp>}); elementary
functions from the C library, called with the rounding direction set
downward or upward (\code{rounded\_transc\_opp}), a direction glibc does not
always honour & The interval arithmetic before IEEE 1788, without empty set:
its default checking (\code{checking\_strict}) throws an exception where an
operation would create one \\
\hline
\end{longtable}
\endgroup

The comparison has three parts, each with a report of its own in
\file{doc/compare/}:
\begin{itemize}
\item[$\bullet$] \textbf{Special cases.} 290 special cases taken from
GAOL's tests (infinities, zeros, NaN, empty sets, divisions by zero,
\code{pow} and \code{pown}, \code{+=} and the other operators with doubles,
reading from text, midpoints, widths, radii, comparisons and \code{atan2}),
computed by the six libraries and compared with the results of IEEE
1788-2015, GAOL under its own names and under those of the standard
(Section~\ref{sec:computing-special-cases}).
\item[$\bullet$] \textbf{Enclosure.} Whether the elementary functions of
GAOL and Boost.Interval enclose their values, at 100\,000 random arguments
for each function, against mpmath, GAOL under both namespaces
(Section~\ref{sec:computing-enclosure}).
\item[$\bullet$] \textbf{Performance.} The time of a million additions,
subtractions, multiplications, divisions, sines, cosines, powers, one-line
combinations, evaluations of Shekel 5 and five-line blocks, and the time of
long lines of $+$, $-$, $\times$ and $\div$ between 8 intervals, GAOL v5
being called under the names of the standard, as the other libraries of IEEE
1788, and GAOL 4.2.3 under its own
(Section~\ref{sec:computing-performance}).
\end{itemize}

\medskip
\noindent Table~\ref{tab:computing-inshort} sums the three parts up. Its
times were measured on 17 October 2026, every library in the same run,
compiled by Clang 21.1, on the machine of
Section~\ref{sec:computing-performance}.

\begin{table}[htbp]
\centering\footnotesize
\setlength{\tabcolsep}{3pt}
\caption[The comparison in short]{The comparison in short: the special cases, the enclosure of the elementary functions and the time of six operations (17 October 2026, every library in the same run, Clang 21.1, Intel i7-1185G7).} \label{tab:computing-inshort}
\begin{tabular}{|L{2.15cm}|L{1.5cm}|L{2.2cm}|L{1.3cm}|L{1.6cm}|L{1.7cm}|
  L{2.0cm}|}
\hline
& \textbf{libieee\allowbreak p1788} & \textbf{GAOL} & \textbf{filib++}
& \textbf{Solaris Studio} & \textbf{PROFIL/\allowbreak BIAS}
& \textbf{Boost.\allowbreak Interval} \\
\hline
Special cases with IEEE 1788's result & 281 of 281 & 288 of 288 under
\code{gaol\_ieee1788}; 280 under \code{gaol} & 126 of 245 & 150 of 261
& 56 of 223 & 93 of 168 \\
\ldots\ or an interval enclosing it & 0 & 0; 2 & 50 & 41 & 22 & 9 \\
\ldots\ or another result & 0 & 0; 6 & 69 & 70 & 145 & 66 \\
Cases it has no operation for & 8 & 0 & 44 & 28 & 65 & 121 \\
\hline
Elementary functions at random points: results missing the value & ---
& 0 of 1\,600\,000 & --- & --- & --- & 34\,\% (20\,\% with
\code{rounded\_}\allowbreak\code{transc\_std}) \\
\hline
\code{x + y} & 137 ns & 2.2 ns & 5.1 ns & 18 ns & 14 ns & 12 ns \\
\code{x * y} & 171 ns & 2.3 ns & 12 ns & 21 ns & 15 ns & 21 ns \\
\code{sin(x)} & 5.5 $\mu$s & 43 ns & 33 ns & 41 ns & 115 ns & 71 ns \\
\code{log(x)} & 1.9 $\mu$s & 21 ns & 27 ns & 39 ns & 10 ns & 30 ns \\
\code{pown(x, 3)} & 240 ns & 5.1 ns & 18 ns & 145 ns & 35 ns & 20 ns \\
Shekel 5 & 9.0 $\mu$s & 108 ns & 397 ns & 1.7 $\mu$s & 1.1 $\mu$s
& 920 ns \\
\hline
\end{tabular}

\end{table}

\begin{itemize}
\item[$\bullet$] \textbf{libieeep1788} gives the result of IEEE 1788 in
every case it can compute, as tightly as possible, and is 17 to 135 times
slower than GAOL.
\item[$\bullet$] \textbf{GAOL} is the fastest on the arithmetic, the
integer power and most formulas. Under the names of the standard,
\code{gaol\_ieee1788}, it gives IEEE 1788's result in every special case;
under its own, \code{gaol}, in all but 8, all from its hybrid \code{pow},
which takes \code{pown} for integer exponents, a negative base included. Its
elementary functions, CORE-MATH's, give the tightest bounds; its \code{exp}
and \code{log} are faster than filib++'s, its \code{sin} and \code{cos}
slower.
\item[$\bullet$] \textbf{filib++} is the fastest on \code{sin} and
\code{cos}, 1.3 to 1.6 times as fast as GAOL, but 5.3 times as slow on
$\times$, 3.5 times on $\div$ and 2.3 times on $+$ and $-$, and its bounds
of elementary functions and real powers are up to 36 doubles wider than the
tightest. Its extended mode gives other results than IEEE 1788 wherever an
infinity or a division by zero is involved, much as Solaris Studio:
\code{interval(+}$\infty$\code{)} is $\itv{\mathrm{MAX}}{+\infty}$ and
$\itv{0}{0} \times \itv{1}{+\infty}$ is $\itv{-\infty}{+\infty}$; its
\code{operator>>} rounds the bounds it reads to nearest, and its
comparisons with the empty set and the infinities differ from IEEE 1788's.
\item[$\bullet$] \textbf{PROFIL/BIAS} is a library of the interval
arithmetic before IEEE 1788: it has no empty set, and an operation whose
argument leaves the domain of the function, a division by an interval
containing zero included, prints an error and aborts the program. Its
exponential, logarithm and real power, taken from the libm, are the fastest
of the six libraries (9.9, 10.2 and 34 ns), and so is its square root (3.7 ns); its
arithmetic is 6.4 times slower than GAOL's on $+$ and $-$, its \code{sin}
and \code{cos} are the slowest after libieeep1788's (115 and 131 ns), and
its integer power, computed as $\exp(n \log x)$, is $1.1\cdot10^{-5}$ wider
than the tightest on average.
\item[$\bullet$] \textbf{Boost.Interval}, the library of Boost, is the
interval library C++ programmers find first. Its arithmetic gave the
tightest intervals in the benchmark, but its elementary functions are the C
library's, called with the rounding direction set downward or upward, which
glibc does not always honour: at the points of the test of enclosure,
34\,\% of its results miss the value of the function with glibc 2.43, 20\,\%
with its other policy, and \code{atan(}$\itv{-\infty}{+\infty}$\code{)}
misses $\pi/2$. With its default checking it has no empty set: an operation
that would create one throws an exception. It is 5.6 times slower than GAOL
on $+$ and $-$, 1.8 to 11 times on $\times$, $\div$ and the formulas, 3.5
times on the square root, 3.9 times on the integer power and 1.2 to 1.6
times on the other elementary functions.
\item[$\bullet$] \textbf{Solaris Studio} is 1.0 to 1.3 times as fast as GAOL
on \code{sin} and \code{cos}, but slow on integer powers and squares. Its
containment sets give other results than IEEE 1788 wherever an infinity, a
division by zero or an invalid argument is involved:
\code{interval(+}$\infty$\code{)} is $\itv{\mathrm{MAX}}{+\infty}$,
$\itv{1}{2} / \itv{0}{1}$ is $\itv{-\infty}{+\infty}$ and
\code{interval(2, 1)} is $\itv{-\infty}{+\infty}$.
\end{itemize}

\subsection{Performance}
\label{sec:computing-performance}
\cindex{benchmark of the libraries}%

One million operations of each kind are computed by each library on
ordinary intervals: bounded, of widths from $10^{-6}$ to 1, drawn once and
read from the same file by every program. Each operation is timed on the
whole array, its results being stored in another one; the time per
operation is the best time divided by the million. The same operations on
doubles, computed at the midpoints of the intervals without any rounding,
give the cost of the arithmetic itself. Table~\ref{tab:computing-bench-ops}
lists the operations: $a$ and $b$ are intervals centred in $\itv{-10}{10}$,
$p$ in $\itv{1}{10}$, $e$ in $\itv{0.5}{2.5}$, the arguments of Shekel 5 in
$\itv{0}{10}$; their widths are $10^u$, $u$ uniform in $\itv{-6}{0}$
($\itv{-6}{-1}$ for $e$).

\begin{table}[htbp]
\centering\footnotesize
\caption{The operations of the benchmark of \file{doc/compare/performance.md}.} \label{tab:computing-bench-ops}
\begin{tabular}{|l|L{10.5cm}|}
\hline
\textbf{Operation} & \textbf{Computes} \\
\hline
\code{add} & $a + b$ \\
\code{sub} & $a - b$ \\
\code{mul} & $a \times b$ \\
\code{div} & $a / p$ \\
\code{sqr} & $a^2$ (\code{sqr}, \code{square} in Boost.Interval,
\code{x**2} in Fortran) \\
\code{sqrt} & $\sqrt{p}$ \\
\code{exp} & $\exp(a)$ \\
\code{log} & $\log(p)$ \\
\code{sin} & $\sin(a)$ \\
\code{cos} & $\cos(a)$ \\
\code{pow\_int} & $a^3$ (\code{pow(x, int)} in GAOL and Boost.Interval,
\code{pown}, \code{power(x, int)} in filib++, \code{x**3}) \\
\code{pow\_real} & $p^e$ (\code{pow(x, y)}, \code{x**y}; $\exp(e \log p)$
with Boost.Interval, which has no real power) \\
\code{line\_arith} & $(a + b)(a - b) / p$ \\
\code{line\_trig} & $\sin(a) \cos(b) + a^2$ \\
\code{line\_pow} & $\sqrt{p} \cdot a^3 - \exp(b / p)$ \\
\code{shekel5} & Shekel 5: $-\sum_i 1 / (\sum_j (x_j - a_{ij})^2 + c_i)$,
5 terms, 4 variables \\
\code{block5} & five lines: $t_1 = ab + p$; $t_2 = \sin(t_1) \cos(b)$;
$t_3 = a^2 + t_2/p$; $t_4 = \exp(t_2) - b^3$; $t_3 t_4 + \sqrt{p}$ \\
\hline
\end{tabular}

\end{table}

Everything was built and run again with Clang 21.1, \option{-O3 -mfma} and
the flags of interval arithmetic (\pxref{sec:compiling}), so that one
compiler answers for every C and C++ library: GAOL v5 (V5.0.0) built by
CMake in Release with \option{GAOL\_FMA}, the sources of CORE-MATH compiled
into it and the rounding direction not preserved, its default, and called
under the names of IEEE 1788-2015, \code{gaol\_ieee1788}, as the other
libraries of the standard; the same GAOL built with
\option{GAOL\_PRESERVE\_ROUNDING}, which takes the AVX-512 path
(Section~\ref{sec:computing-avx512-rule}), in a column of the table of the
times only; GAOL 4.2.3, the last version of Fr\'ed\'eric Goualard, built
by its \cmd{configure} from his repository (cd0ee1a), static, the rounding
direction not preserved, its verbose mode off, with mathlib 2.1.1 built by
Clang 21.1 with \option{-O3 -mfma -ffp-contract=off}, and called under its
own names; libieeep1788 with MPFR 4.2.1 and GMP 6.3.0; filib++ 3.0.2.2, the
archive IBEX distributes, as
\code{interval<double, native\_switched, i\_mode\_extended\_flag>};
PROFIL/BIAS 2.0.8 in its \code{x86-64-Linux-compat-gcc} configuration;
Boost.Interval of Boost 1.92.0 with its default policies and
\code{rounded\_transc\_opp} for the elementary functions, glibc 2.43. Only
Solaris Studio keeps its own compiler, Sun Fortran 95 8.7 with
\option{-O3 -xia}, by nature. Every library was measured in the same run, on 17 October 2026, and GAOL v5 with the rounding direction preserved again on 18 October, on a laptop with an Intel i7-1185G7 at 3.00 GHz (pinned by
\cmd{taskset -c 2}), Linux 7.0.0-38-generic, Ubuntu 26.04.1; each time is
the best of the runs of three rounds, each program being run in turn with
the others. libieeep1788's times, from a single run of each operation in
each round, are the least steady. Solaris Studio's and PROFIL/BIAS's
operations are calls into libraries, as GAOL v5's \code{+}, \code{-},
\code{*}, \code{/} and \code{sqr} are, the bounds in registers, where the
compiler inlines GAOL 4.2.3's, filib++'s and Boost.Interval's.

\paragraph{What GAOL v5 changes.}
GAOL v5 is compared with GAOL 4.2.3, which it continues and which bounds
the elementary functions with mathlib (\pxref{sec:math-libraries}).
Table~\ref{tab:computing-v5-v4} gives the operations where they differ
most. The elementary functions are faster:\newinvfive{} two things go at
once with mathlib, the outward move of each bound, and the two changes of
rounding direction each function made, to nearest before mathlib and back
upward after. The integer power is computed from exact products, both
bounds at once in an SSE2 register, with an error bound proved beforehand.
The bounds are tighter, not looser: GAOL v5's results are no
wider than libieeep1788's, which computes every bound with MPFR, for every
one of the seventeen operations (excess $0.0$), where GAOL 4.2.3 is wider
on the square root, \code{exp}, \code{log}, \code{sin}, \code{cos}, both
powers and the lines that use them.

\begin{table}[htbp]
\centering\footnotesize
\caption[GAOL v5 against GAOL 4.2.3, in nanoseconds per operation]{GAOL v5 against GAOL 4.2.3, in nanoseconds per operation (Intel i7-1185G7, Clang 21.1, 17 October 2026).} \label{tab:computing-v5-v4}
\begin{tabular}{|l|r|r|l|}
\hline
\textbf{Operation} & \textbf{GAOL v5} & \textbf{GAOL 4.2.3} & \textbf{Ratio} \\
\hline
\multicolumn{4}{|l|}{\emph{The elementary functions and the powers}} \\
\hline
\code{log} & 20.9 ns & 46.8 ns & 2.2 times faster \\
\code{exp} & 18.3 & 35.1 & 1.9 \\
\code{pow(x, y)} & 48.2 & 89.9 & 1.9 \\
\code{pow(x, 3)} & 5.12 & 9.07 & 1.8 \\
line of powers & 38.1 & 61.0 & 1.6 \\
\code{sin} & 43.2 & 68.0 & 1.6 \\
five-line block & 155 & 229 & 1.5 \\
line of sin and cos & 102 & 148 & 1.5 \\
\code{cos} & 54.0 & 62.8 & 1.2 \\
\hline
\multicolumn{4}{|l|}{\emph{The arithmetic}} \\
\hline
\code{mul} & 2.30 ns & 10.0 ns & 4.4 times faster \\
\code{line\_arith} & 5.61 & 18.7 & 3.3 \\
\code{sqr} & 1.87 & 5.81 & 3.1 \\
\code{div} & 2.52 & 7.50 & 3.0 \\
\code{shekel5} & 108 & 189 & 1.7 \\
\hline
\end{tabular}

\end{table}

The arithmetic is faster than GAOL 4.2.3's, which is what the operators of
GAOL v5 are written for (Section~\ref{sec:kernels}). The sums and the
differences, 2.2 ns against 2.4, and the square root, 6.3 against 6.3, take
what they took: a sum is one instruction, and what it waits for is memory,
the benchmark reading two arrays of 16 MB and writing a third, and the
square root is the tightest interval whatever the rounding of the C
library's \code{sqrt}, where GAOL 4.2.3's is wider (excess
$1.4\cdot10^{-14}$). The check of the direction is what a program needs that
changes the rounding direction: the benchmark sums the results rounding to
nearest, and GAOL 4.2.3's own \code{midpoint()} leaves the direction to
nearest. Left so, the operations of GAOL 4.2.3 timed next gave bounds
narrower than libieeep1788's, which did not enclose the exact results; the
benchmark now sets the direction back after its sums. The bounds are the
tightest, as before: the zero bound of a product or of a quotient may have
the other sign (\pxref{sec:zero-sign}).

\paragraph{Preserving the rounding direction.}
The column of Table~\ref{tab:computing-times} where the direction is
preserved (\option{GAOL\_PRESERVE\_ROUNDING} on, with the AVX-512 path)
leaves the sums and the differences as they are, 2.2 ns, and costs 1.5
times on the products and 1.4 times on the quotients (3.4 and 3.6 ns), 1.1 times on
\code{sqr} (2.0 ns against 1.9) and 1.2 times on \code{pow(x, 3)} (6.2 ns
against 5.1): these operations take the AVX-512 instructions, whose rounding
direction is in the instruction, and set none, so that nothing is saved and
restored. The square root, 3.75 ns, is faster than without the preservation
(6.3 ns): the instruction rounds each root in its direction, where the SSE2
path checks the root of the C library with a product. The elementary
functions read the control register of the direction, set it upward and
write it back, once each, and take 1.0 to 1.3 times longer (\code{exp}
23.2 ns against 18.3, \code{log} 20.4 against 20.9, \code{sin}, \code{cos}
and \code{pow(x, y)} 1.1 times). They paid more, as \code{sqr} and the
integer powers did, before the kernels of Section~\ref{sec:kernels}: the two
writes of the control register, each one waiting for the operations in
flight, took 6 ns a \code{sqr} (8.3 ns, 4.4 times); with the C library's
\code{fegetround()} and \code{fesetround()} and the writes of the control
word of the x87 unit as well, \code{pow(x, 3)} took 14 ns more (19.0 ns, 3.7
times; 8.4 ns with the writes of the SSE register alone, where the machine
has no AVX-512) and the elementary functions 1.2 to 1.7 times longer.
Shekel 5 takes 147 ns against 108, and every operation of the table is
faster than with GAOL 4.2.3, which does not preserve the direction either
(Shekel 5: 189 ns). The results are the same: the sums of the midpoints and
of the widths of the million results are those of GAOL v5, to
$2\cdot10^{-16}$ relatively.

\begin{table}[htbp]
\centering\footnotesize
\setlength{\tabcolsep}{3pt}
\caption[Time per operation of the six libraries, in nanoseconds]{Time per operation, in nanoseconds: the same operations on doubles (double), libieeep1788 (p1788), GAOL v5 (v5), GAOL v5 with the rounding direction preserved and the AVX-512 path (v5 pres.), GAOL 4.2.3 (v4), filib++, Solaris Studio (Sun), PROFIL/BIAS (BIAS) and Boost.Interval (Boost). One million operations of each kind on the same intervals, the best of 3 rounds; Intel i7-1185G7, Clang 21.1.8, Ubuntu 26.04.1, 17 October 2026 (the preserved column on 18 October).} \label{tab:computing-times}
\begin{tabular}{|l|r|r|r|r|r|r|r|r|r|}
\hline
\textbf{Operation} & double & p1788 & v5 & v5 pres. & v4 & filib++
& Sun & BIAS & Boost \\
\hline
\code{add} & 0.96 & 137 & 2.23 & 2.19 & 2.38 & 5.10 & 17.6 & 14.2 & 12.4 \\
\code{sub} & 0.92 & 138 & 2.21 & 2.16 & 2.40 & 5.12 & 17.6 & 14.3 & 12.3 \\
\code{mul} & 0.90 & 171 & 2.30 & 3.42 & 10.0 & 12.1 & 21.3 & 14.7 & 21.5 \\
\code{div} & 0.91 & 171 & 2.52 & 3.60 & 7.50 & 8.70 & 20.3 & 14.4 & 18.9 \\
\code{sqr} & 0.59 & 88.4 & 1.87 & 2.04 & 5.81 & 8.25 & 55.5 & 18.5
& 16.2 \\
\code{sqrt} & 1.30 & 107 & 6.30 & 3.75 & 6.26 & 13.4 & 26.6 & 3.73
& 22.3 \\
\code{exp} & 3.44 & 1\,491 & 18.3 & 23.2 & 35.1 & 30.1 & 38.0 & 9.87
& 29.8 \\
\code{log} & 3.32 & 1\,919 & 20.9 & 20.4 & 46.8 & 26.6 & 38.5 & 10.2
& 29.9 \\
\code{sin} & 13.4 & 5\,472 & 43.2 & 46.8 & 68.0 & 33.4 & 41.5 & 115
& 71.2 \\
\code{cos} & 12.0 & 4\,890 & 54.0 & 57.6 & 62.8 & 33.9 & 42.7 & 131
& 65.7 \\
\code{pow\_int} & 13.0 & 240 & 5.12 & 6.18 & 9.07 & 17.8 & 145 & 35.1
& 19.9 \\
\code{pow\_real} & 10.3 & 6\,517 & 48.2 & 51.4 & 89.9 & 64.5 & 125 & 33.6
& 73.2 \\
\code{line\_arith} & 1.14 & 655 & 5.61 & 8.24 & 18.7 & 27.4 & 61.6 & 63.1
& 61.7 \\
\code{line\_trig} & 25.8 & 10\,920 & 102 & 113 & 148 & 90.8 & 166 & 311
& 187 \\
\code{line\_pow} & 19.4 & 2\,401 & 38.1 & 43.7 & 61.0 & 81.9 & 253 & 92.6
& 115 \\
\code{shekel5} & 2.43 & 8\,993 & 108 & 147 & 189 & 397 & 1\,720 & 1\,093
& 920 \\
\code{block5} & 44.5 & 13\,823 & 155 & 171 & 229 & 197 & 464 & 427 & 343 \\
\hline
\end{tabular}

\end{table}

\paragraph{What the timings show.}
On the arithmetic, GAOL is the fastest on the four operations: $+$ and $-$
take GAOL v5 2.2 ns (2.4 for GAOL 4.2.3), 2.3 times as fast as filib++
(5.1 ns) and 6.4 to 8.0 times as fast as PROFIL/BIAS (14 ns) and Solaris
Studio (17.6 ns); $\times$ and $\div$ take GAOL v5 2.3 and 2.5 ns (10.0 and
7.5 for GAOL 4.2.3), against 12.1 and 8.7 ns for filib++, 14.7 and 14.4 ns
for PROFIL/BIAS, 21 and 20 ns for Solaris Studio and 21 and 19 ns for
Boost.Interval. An addition of intervals costs GAOL v5 2.3 times an
addition of doubles: one vector addition for both bounds, and the check that the rounding
direction is upward, where filib++ sets the rounding direction and restores
it, and PROFIL/BIAS calls a function of BIAS that sets it downward, then
upward, then back to nearest, which is why its four arithmetic operations
all take about the same time. The integer power \code{pown(x, 3)} takes
GAOL v5 5.1 ns and GAOL 4.2.3 9.1 ns: filib++'s \code{power(x, 3)} takes
18 ns, Boost.Interval's 20 ns, PROFIL/BIAS's 35 ns and Solaris Studio's
\code{x**3} 145 ns; \code{sqr} takes GAOL v5 1.9 ns, against 8.2 for
filib++, 16 for Boost.Interval, 18 for PROFIL/BIAS and 56 for Solaris
Studio.

\medskip
On the elementary functions, PROFIL/BIAS is the fastest on the square root
(3.7 ns), \code{exp} (9.9 ns), \code{log} (10.2 ns) and the real power
\code{pow(x, y)} (34 ns), which it computes from the libm of the system,
moved outward, without correct rounding: its bounds are among the widest,
and its \code{sin} and \code{cos}, computed by its own argument reduction,
are the slowest after libieeep1788's, 115 and 131 ns. GAOL v5 comes next on
\code{exp} and \code{log}, 18 and 21 ns, faster than Boost.Interval (30 ns
each), filib++ (30 and 27 ns) and Solaris Studio (38 and 39 ns), with
correctly rounded values, where the others are not. filib++ is the fastest
on \code{sin} and \code{cos} (33.4 and 33.9 ns) with its own polynomials,
then Solaris Studio (41 and 43 ns), then GAOL v5 (43 and 54 ns): GAOL pays
there for correct rounding over the whole range, and for dividing the bounds
by an interval enclosing $\pi$ to find where the function is monotonic,
asking CORE-MATH for the signs of the derivative where the division cannot
tell. GAOL v5's \code{sin} became faster than its \code{cos} when CORE-MATH
rewrote \file{sin.c}. On \code{pow(x, y)} GAOL v5 takes 48 ns, against 65
for filib++, 73 for Boost.Interval, 90 for GAOL 4.2.3 and 125 for Solaris
Studio.

\medskip
On the formulas, GAOL is the fastest on the arithmetic line (5.6 ns for GAOL
v5 and 18.7 for GAOL 4.2.3, against 27 for filib++ and 62 to 63 for Solaris
Studio, PROFIL/BIAS and Boost.Interval), on the line of powers (38 and
61 ns, against 82 for filib++, 93 for PROFIL/BIAS and 253 for Solaris
Studio) and on Shekel 5 (108 and 189 ns, against 397 for filib++, 920 for
Boost.Interval, 1\,093 for PROFIL/BIAS and 1\,720 for Solaris Studio: Shekel
5 is 20 squares, 45 additions and subtractions and 5 divisions, and its
squares cost Solaris Studio 56 ns each and PROFIL/BIAS 18). GAOL v5 alone is
the fastest on the five-line block (155 ns, against 197 for filib++, 229 for
GAOL 4.2.3, 427 for PROFIL/BIAS and 464 for Solaris Studio). filib++ is the
fastest on the line of sin and cos (90.8 ns, against 102 for GAOL v5, 148
for GAOL 4.2.3, 166 for Solaris Studio and 311 for PROFIL/BIAS), where its
\code{sin} and \code{cos} weigh most.

\medskip
libieeep1788 is 17 to 135 times slower than GAOL v5: every bound is an MPFR
computation, about 140 ns for an addition and 5 to 7 $\mu$s for \code{sin},
\code{cos} and the real power. Its own documentation warns that its focus is
correctness, not speed. GAOL v5 gives the same bounds as it does, tightest
everywhere, 17 to 135 times faster. Boost.Interval is slower than GAOL v5 on
every operation: 5.6 times on $+$ and $-$, 3.5 to 9.3 times on $\times$,
$\div$, the square, the square root and the integer power, 1.2 to 1.6 times
on \code{exp}, \code{log}, \code{sin}, \code{cos} and the real power, 11
times on the arithmetic line and 8.5 times on Shekel 5. Each of its
operations reads the rounding direction (\code{fegetround}), sets it upward
(\code{fesetround}) and, at its end, sets the one it read back: two changes
of direction where GAOL checks the direction and changes nothing. Its
elementary functions add a change to the downward direction and back for the
lower bound, and its \code{sin} goes through the reduction of
$\cos(x - \pi/2)$ by an interval enclosing $2\pi$. Its bounds are those of
the C library, which are not correctly rounded, and of which 20 to 34\,\%
miss the value of the function with glibc 2.43
(Section~\ref{sec:computing-enclosure}).

\paragraph{The same results.}
The sums of the midpoints of the million results agree to $4\cdot10^{-15}$
relatively, apart from PROFIL/BIAS's and the sines of GAOL 4.2.3 and
Boost.Interval, $1.4\cdot10^{-13}$ and $1.3\cdot10^{-13}$ apart, their sum
being small, 373 for a million values, which makes $5\cdot10^{-17}$ per
result. The arithmetic operations give the tightest intervals in every
library, and so do the square roots, except those of GAOL 4.2.3, filib++ and
PROFIL/BIAS. Of the elementary functions, only GAOL v5's are the tightest on
every operation: on average the others are wider than libieeep1788's by up
to $5.9\cdot10^{-14}$ relatively with GAOL 4.2.3 (the real power, computed
as $\exp(y \log x)$), $2.1\cdot10^{-14}$ to $2.4\cdot10^{-13}$ with filib++ ($1.4\cdot10^{-13}$ for \code{log}, $2.4\cdot10^{-13}$ for the real power), $9.0\cdot10^{-15}$ to $1.2\cdot10^{-13}$ with PROFIL/BIAS, and at most $4.1\cdot10^{-15}$ with Solaris Studio; with Boost.Interval, by
$1.4\cdot10^{-14}$ and $2.5\cdot10^{-14}$ on \code{cos} and \code{sin} and
$3.7\cdot10^{-14}$ on the real power, its \code{exp} and \code{log} being as
tight as libieeep1788's on average. PROFIL/BIAS's integer power is far
wider, $1.1\cdot10^{-5}$ on average: it computes \code{Power(x, n)} as
$\exp(n \log x)$, where the others multiply.

\begin{Rem}
The arithmetic results of PROFIL/BIAS have midpoints that differ by
$8\cdot10^{-15}$ to $4\cdot10^{-14}$ relatively while their widths are the
tightest: the intervals are the same, but its \code{Mid} computes
$\inf + (\sup - \inf)/2$ rounded upward, one double above the midpoint for
4\,\% of the sums of the benchmark, where the others round to nearest. In
the same way, the sums of the widths of Boost.Interval's results of
\code{exp} and \code{log} equal libieeep1788's, while its upper bound of
\code{exp} is one double too low in 2 to 3\,\% of the results of the test of
enclosure: a sum of a million widths cannot show it, and only a check of
each bound against the exact value finds it.
\end{Rem}

\paragraph{Long lines of arithmetic operators.}
The operations above are timed one at a time. A program that computes with
intervals writes lines of them, an operator after another, where what an
operation hands to the next matters as much as what it computes. The lines
of Table~\ref{tab:computing-lines} are expressions of $+$, $-$, $\times$
and $\div$ alone, of 7 to 19 operations between 8 intervals of the same
million, evaluated for each of them, by every library but Solaris Studio,
which has no program for them. The seven first take intervals of both
signs, at random; the last three, \code{pos\_sum16}, \code{pos\_mul8} and
\code{pos\_quot8}, only $p$ and $e$, which are positive, so that the
branches of an implementation that tells the signs apart, as GAOL 4.2.3's
does, are taken the same way from one evaluation to the next. The chains
give the next operation the result of the previous one, the trees
(\code{mul8\_tree}, \code{mix\_tree}) several independent ones.
Table~\ref{tab:computing-lines-x86} gives the times on the x86 machine of
the benchmark, Table~\ref{tab:computing-lines-m2} on an Apple M2.

\begin{table}[htbp]
\centering\footnotesize
\caption[The long lines of arithmetic operators between 8 intervals]{The long lines of arithmetic operators between 8 intervals ($a$, $b$, $p$, $e$ as in Table~\ref{tab:computing-bench-ops}, $s_1$ to $s_4$ in $\itv{0}{10}$).} \label{tab:computing-lines}
\begin{tabular}{|l|L{9.0cm}|r|}
\hline
\textbf{Line} & \textbf{Computes} & \textbf{Operations} \\
\hline
\code{sum16} & $a + b + p + e + s_1 + s_2 + s_3 + s_4 + a + b + p + e + s_1
+ s_2 + s_3 + s_4$ & 15 \\
\code{diff16} & $a - b + p - e + s_1 - s_2 + s_3 - s_4 - a + b - p + e - s_1
+ s_2 - s_3 + s_4$ & 15 \\
\code{mul8} & $a \cdot b \cdot p \cdot e \cdot s_1 \cdot s_2 \cdot s_3 \cdot
s_4$ & 7 \\
\code{mul8\_tree} & $((a \cdot b)(p \cdot e))((s_1 \cdot s_2)(s_3 \cdot
s_4))$ & 7 \\
\code{quot8} & $a / p / e / p / e / p / e / p$ & 7 \\
\code{mix13} & $((((((a + b) \cdot p - e) \cdot s_1 + s_2) \cdot p - s_3) /
e + a) \cdot b - s_4) / p + e$ & 13 \\
\code{mix\_tree} & $((a + b)(p - e) + (s_1 - s_2)(s_3 + s_4)) / (p + e) +
((a - b)(s_1 + s_2) - (p \cdot e) / (p + e)) / (e + p)$ & 19 \\
\code{pos\_sum16} & $p + e + p + e + p + e + p + e + p + e + p + e + p + e
+ p + e$ & 15 \\
\code{pos\_mul8} & $p \cdot e \cdot p \cdot e \cdot p \cdot e \cdot p \cdot
e$ & 7 \\
\code{pos\_quot8} & $p / e / p / e / p / e / p / e$ & 7 \\
\hline
\end{tabular}

\end{table}

\begin{table}[htbp]
\centering\footnotesize
\setlength{\tabcolsep}{3pt}
\caption[Time per evaluation of a line, in nanoseconds]{Time per evaluation of a line, in nanoseconds, with the abbreviations of Table~\ref{tab:computing-times} (one million evaluations of each line, the best of 3 rounds; Intel i7-1185G7, Clang 21.1.8, 17 October 2026).} \label{tab:computing-lines-x86}
\begin{tabular}{|l|r|r|r|r|r|r|r|}
\hline
\textbf{Line} & double & p1788 & v5 & v4 & filib++ & BIAS & Boost \\
\hline
\code{sum16} & 2.45 & 2\,001 & 22.8 & 27.4 & 76.5 & 216 & 179 \\
\code{diff16} & 2.44 & 2\,091 & 23.3 & 27.1 & 76.5 & 215 & 177 \\
\code{mul8} & 2.42 & 1\,166 & 17.9 & 30.9 & 48.6 & 104 & 111 \\
\code{mul8\_tree} & 2.37 & 1\,168 & 15.0 & 29.4 & 50.5 & 105 & 111 \\
\code{quot8} & 3.36 & 1\,178 & 20.6 & 34.0 & 48.6 & 102 & 112 \\
\code{mix13} & 2.38 & 2\,068 & 23.0 & 47.6 & 80.7 & 192 & 185 \\
\code{mix\_tree} & 2.38 & 2\,985 & 34.7 & 64.7 & 116 & 278 & 264 \\
\code{pos\_sum16} & 1.03 & 1\,930 & 23.1 & 24.5 & 73.5 & 214 & 179 \\
\code{pos\_mul8} & 0.89 & 1\,139 & 17.8 & 18.9 & 40.2 & 100 & 105 \\
\code{pos\_quot8} & 3.27 & 1\,160 & 20.2 & 37.6 & 44.3 & 101 & 109 \\
\hline
\end{tabular}

\end{table}

GAOL v5 is the fastest of the libraries on every line, and faster than GAOL
4.2.3 on all ten, 1.06 to 2.1 times as fast: a line of 15 additions takes
22.8 ns, 1.5 ns an addition, against 27.4 ns (1.8 ns), and a line of 7
products 17.9 ns (2.6 ns a product) against 30.9 ns (4.4 ns); the quotients, 20.6 ns for 7 (2.9 ns each) against 34.0 ns (4.9 ns), 1.7 times as fast, and
the lines that mix the four operations, \code{mix13} and \code{mix\_tree},
23.0 and 34.7 ns against 47.6 and 64.7 ns, 2.1 and 1.9 times. The three
lines of intervals of one sign are where GAOL 4.2.3 comes nearest, on the
additions and the products, 1.06 times as fast as it, 23.1 ns against
24.5 ns and 17.8 ns against 18.9 ns: with the signs the same from one
evaluation to the next, its branches are predicted, and what it still pays
is the bounds that go through memory between two operations; its quotients,
on the other hand, take 1.9 times as long (37.6 ns against 20.2 ns). An
operation of a line costs GAOL v5 less than the same operation timed alone,
1.5 ns an addition against 2.2: the two doubles of the result stay in a
register for the next operation, and the benchmark of a single operation,
which reads two arrays and writes a third for each, waits for the memory,
where a line reads eight arrays at most for its 7 to 19 operations. Against
the other libraries, filib++ takes 2.2 to 3.5 times as long as GAOL v5 on
the ten lines, Boost.Interval 5.4 to 8.0 times, PROFIL/BIAS 5.0 to 9.5 times
and libieeep1788 57 to 90 times. The bounds are the tightest: the mean
widths of the results of GAOL v5 are those of libieeep1788 to
$1.1\cdot10^{-16}$ relatively on the lines of additions and exactly on the
others, where GAOL 4.2.3's differ from them by $2.3\cdot10^{-16}$ on the
lines of quotients.

\begin{table}[htbp]
\centering\footnotesize
\caption[The same lines on a 64-bit ARM processor (Apple M2)]{The same lines on a 64-bit ARM processor, in nanoseconds per evaluation: GAOL v5 with the NEON intervals (v5), the same with the rounding direction preserved (v5 pres.) and GAOL 4.2.3 with the FPU intervals (v4). Apple M2, Apple clang 21.0.0, macOS 27.0.1, 6 November 2026, the best of 3 rounds, the processor not pinned.} \label{tab:computing-lines-m2}
\begin{tabular}{|l|r|r|r|r|r|r|}
\hline
\textbf{Line} & \textbf{Ops} & v5 & v5 pres. & v4 & v5 pres. / v5
& v4 / v5 \\
\hline
\code{sum16} & 15 & 9.11 & 385 & 41.9 & 42.3 & 4.6 \\
\code{diff16} & 15 & 9.30 & 390 & 77.6 & 42.0 & 8.3 \\
\code{mul8} & 7 & 10.5 & 187 & 28.3 & 17.8 & 2.7 \\
\code{mul8\_tree} & 7 & 9.95 & 190 & 25.0 & 19.1 & 2.5 \\
\code{quot8} & 7 & 17.6 & 229 & 37.3 & 13.0 & 2.1 \\
\code{mix13} & 13 & 14.3 & 353 & 61.6 & 24.7 & 4.3 \\
\code{mix\_tree} & 19 & 18.5 & 516 & 82.2 & 27.9 & 4.4 \\
\code{pos\_sum16} & 15 & 8.85 & 385 & 39.7 & 43.5 & 4.5 \\
\code{pos\_mul8} & 7 & 9.58 & 187 & 16.1 & 19.6 & 1.7 \\
\code{pos\_quot8} & 7 & 17.6 & 232 & 46.5 & 13.2 & 2.6 \\
\hline
\end{tabular}

\end{table}

On the Apple M2, GAOL v5 takes 8.85 to 18.5 ns for a line of 7 to 19
operations, 0.6 to 2.5 ns an operation, where GAOL 4.2.3, which has no instructions of 64-bit ARM to compute its FPU intervals with, takes 1.7 to
8.3 times as long (its \cmd{configure} was generated again by
\cmd{autoreconf} and told \option{--build=arm-apple-darwin}, and mathlib
2.1.1 configured with \option{--host=aarch64-unknown-linux-gnu}; no line
calls it). With the rounding direction preserved the same lines take 187 to
516 ns, 25.7 to 33.1 ns an operation, 13.0 to 43.5 times as long as without:
the program computes to nearest, so each operation sets the direction upward
and puts nearest back, two \code{msr} of FPCR of 11 ns each that change it,
with the \code{mrs} of 0.8 ns that reads it. The sums of the midpoints and
of the widths of the three agree to $1.1\cdot10^{-16}$ relatively. The
machine was on mains power, the low power mode off, with a load average of
1.7 at the start and 3.0 at the end (the run took 65 s, the build of the
three programs included), and no process above 10\,\% of a core but the interface of
the machine in the readings taken just before and just after the run; the
three rounds of a line differ by 1.4\,\% at the median and by 6.4\,\% at
most, and the best of the three is kept.

\subsection{The special cases}
\label{sec:computing-special-cases}
\cindex{special cases!comparison}%

The special cases of GAOL's tests (\file{tests/*.cpp}) are computed by the six
libraries, GAOL v5 twice, under its own names and under those of IEEE
1788-2015: zeros, infinities and NaN as bounds or operands, empty sets,
divisions by intervals containing zero, arguments at the edges of the
domains, the powers \code{pow} and \code{pown}, the operators \code{+=},
\code{-=}, \code{*=}, \code{/=} with special doubles, the reading of
intervals from text, the numeric functions of intervals, the comparisons
and \code{atan2}. Each result is compared with the result IEEE 1788-2015
defines for the set-based flavor, the tightest interval of doubles, computed
with mpmath, and classed as the result of IEEE 1788, as an interval
enclosing it, wider, as another result (a set not enclosing it, a set where
IEEE 1788 has the empty set, another boolean, or an exception), or as not
applicable when the library has no such operation; IEEE 1788 defines no
result in two cases (a bare number read as an interval, and an n-th root of order 0). Signed zeros are not told apart. Each case is written
once, in \file{doc/compare/code/cases.py}, and translated for each library; the errors
of PROFIL/BIAS, which abort the program, and the exceptions of
Boost.Interval are caught. The counts are those of
Table~\ref{tab:computing-inshort}.

\medskip
libieeep1788 gives the result of IEEE 1788 in each of the 281 cases where
IEEE 1788 defines one and it has the operation; it has no n-th root, no
operators with doubles and no compound assignments. GAOL under the names of
the standard gives the result of IEEE 1788 in each of the 288 cases where
IEEE 1788 defines one.\newinvfive{} Under its own names it gives it in 280
cases, an interval enclosing it in 2, and something else in 6, all from the
hybrid \code{pow(x, y)} of \code{gaol}, which takes the integer power for an
integer exponent, where IEEE 1788's \code{pow} only takes the part of $x$ in
$\itv{0}{+\infty}$ (\pxref{sec:pow-namespaces}): \code{pow([-2], 2.0)} is
$\itv{4}{4}$, \code{pow([0], [0])} is $\itv{1}{1}$,
\code{pow([-1], [2\textasciicircum{}31-1])} is $\itv{-1}{-1}$,
\code{pow([-2, -1], [1e10])}, an integer beyond the ints, is
$\itv{1}{+\infty}$, and \code{pow([-2, 0], [-1])} is $\itv{-\infty}{-0.5}$,
where IEEE 1788 has $\emptyset$; \code{pow([-2, 3], 3)} is $\itv{-8}{27}$,
wider than $\itv{0}{27}$. In GAOL~4, the integer powers beyond the ints were
$\itv{-\infty}{+\infty}$, which made three more cases wider.\newinvfive{} GAOL reads a bare number, \code{textToInterval("0.1")},
as the interval enclosing it, an extension of the literals IEEE 1788 allows
(\pxref{sec:reading-intervals}). Its \code{sin}, \code{cos}, \code{tan},
\code{acos}, \code{acosh}, \code{atan2}, integer powers and n-th roots are
the tightest in every one of these cases, the powers being computed from exact products,
$x^{-n}$ included, and the roots proved with integer powers; across the
half-line $y = 0$, $x < 0$, where the angle jumps from $\pi$ to $-\pi$,
\code{atan2} is $\itv{-\pi}{\pi}$, as in libieeep1788.

\medskip
The other libraries depart from the standard each in its own way. filib++,
in its extended mode, takes the infinities for values, much as Solaris
Studio does: the intervals built from an infinity are
$\itv{\mathrm{MAX}}{+\infty}$, a product of a zero and an infinite bound is
$\itv{-\infty}{+\infty}$, \code{log([0])} is $\itv{-\infty}{-\mathrm{MAX}}$
and \code{atanh([1])} is $\itv{\mathrm{MAX}}{+\infty}$ where IEEE 1788 has
$\emptyset$; its \code{operator>>} reads only the form \code{[l, u]} and
rounds the bounds to nearest, so that \code{[0.1, 0.3]} does not enclose
$\itv{0.1}{0.3}$; its elementary functions are wider than the tightest by
up to 25 doubles and its real powers by up to 36. PROFIL/BIAS, a library of
1993 to 2009, has no empty set: an argument outside the domain of a
function, or a division by an interval containing 0, prints an error and
calls \code{abort()}; the intervals built from an infinity, a NaN or bounds
in the wrong order are not empty, \code{interval(2, 1)} being $\itv{1}{2}$;
its \code{sin} and \code{cos} are wider than the tightest by 4 to 14
doubles, its real powers by 3 to 14. Solaris Studio follows the containment
sets of Sun's interval arithmetic, where $1/0$ is $\{-\infty, +\infty\}$ and
$+\infty$ is a point: \code{interval(2, 1)} and NaN bounds give
$\itv{-\infty}{+\infty}$ rather than $\emptyset$, a division by an interval
containing zero is $\itv{-\infty}{+\infty}$, a single number read as an
interval carries an uncertainty of one unit of its last digit (\code{0.1}
gives $\itv{0}{0.2}$), and its powers are often one double wider than the
tightest. Boost.Interval, with its default checking, throws
\code{std::runtime\_error} wherever the empty set would appear, the empty
set given as an argument included; an infinity is a point, a product of a
zero and an infinity has a NaN bound, its \code{median} overflows for
$\itv{\mathrm{MAX}/2}{\mathrm{MAX}}$, and \code{atan([-}$\infty$\code{, +}%
$\infty$\code{])} is $\itv{-1.5707963267948966}{1.5707963267948966}$, which
does not hold the values of $\arctan$ near $\pm\pi/2$: Boost.Interval asks
the C library for $\arctan(+\infty)$ rounded upward, and glibc rounds it to
nearest, below $\pi/2$, whatever the rounding direction.

\subsection{The enclosure of the elementary functions}
\label{sec:computing-enclosure}
\cindex{enclosure test}%

An interval library promises one thing: the interval it returns contains
the exact result. The special cases check the promise at the edges of the
domains; the test of enclosure checks it in the middle of them. Each of 14
elementary functions, the square root and $x^3$ is computed on the point
intervals $\itv{x}{x}$ of 100\,000 random arguments $x$, half of them
ordinary, half at the edges of the doubles or of the domain, and each result
is compared with $f(x)$, computed by mpmath with 2000 bits (8000 where
$f(x)$ is that close to a double, as $\tanh(709)$ is to 1). A result
$\itv{l}{u}$ misses $f(x)$ when $l > f(x)$ or $u < f(x)$. Boost.Interval
computes its arithmetic in the rounding direction each bound needs, but its
elementary functions are those of the C library, called with the rounding
direction set downward for the lower bound and upward for the upper one:
they enclose $f(x)$ only if the C library honours that direction, which its
documentation says is rarely the case. GAOL v5, whose elementary functions
are CORE-MATH's, correctly rounded in each direction, is the control: none
of its results may miss. Both are compiled by Clang 21 with the flags of
interval arithmetic, Boost 1.92.0 with glibc 2.43 (Ubuntu 26.04.1), under
its two policies: \code{rounded\_transc\_std}, which sets the direction
downward, calls the function for the lower bound, sets it upward and calls
it for the upper one, and \code{rounded\_transc\_opp}, the policy of the
benchmark, which keeps the direction upward and sets it downward only around
the call for the lower bound.

\medskip
GAOL v5 encloses $f(x)$ in every one of the 1\,600\,000 results, under both
namespaces, and gives the tightest interval of doubles around it in all of
them (1\,232 results of $x^3$, where the cube is below $2^{-968}$, came from
products rounded outward and were one or two doubles wider, until these
corners were computed apart, scaled; \pxref{sec:accuracy-powers}). Boost.Interval
misses $f(x)$ in 34\,\% of the results with \code{rounded\_transc\_opp} and
in 20\,\% with \code{rounded\_transc\_std}. Its square root and its $x^3$,
which it computes itself with the rounding of the processor, never miss,
nor do its \code{asinh} and \code{atanh}; every other elementary function
misses, \code{log} once, the others in 0.05\,\% of the results or more, and
\code{asin}, \code{atan}, \code{sinh}, \code{tanh} and \code{acosh} in
100\,\% with \code{rounded\_transc\_opp}. Three causes were found. The
\code{sin}, \code{cos}, \code{tan} and \code{atan} of glibc ignore the
rounding direction: called downward, to nearest and upward, they give the
same double for every argument, which misses $f(x)$ unless $f(x)$ is that
double; Boost.Interval's \code{sin}, \code{cos} and \code{tan} miss less
often, 0.3 to 28\,\%, because they reduce the argument by subtracting a
multiple of an interval enclosing $\pi$ or $2\pi$, which widens most results
by tens to hundreds of doubles. The \code{exp}, \code{cosh}, \code{acos} and
\code{log} of glibc are computed in the direction asked, but not always on
the right side of $f(x)$: called upward, \code{exp} gives a double below
$f(x)$ for 2 to 3\,\% of the arguments, \code{cosh} for up to 7\,\%,
\code{acos} misses in 0.09 to 5\,\%, \code{log} once in the 100\,000. The
\code{asin}, \code{sinh} and \code{tanh} of glibc round the wrong way for
many arguments, the value computed downward being above the one computed
upward, as if the function were computed for $|x|$ and its sign changed
afterwards: with \code{rounded\_transc\_std} the result is then inverted,
its lower bound above its upper one, and the default checking does not see
an inverted interval as empty, since it holds that no interval is.

\begin{Rem}
\code{rounded\_transc\_opp} has a bug of its own. For \code{sin},
\code{tan}, \code{asin}, \code{atan}, \code{sinh}, \code{tanh} and
\code{acosh} it computes the lower bound as the opposite of the opposite of
$f(x)$ computed upward, that is $f(x)$ computed upward again: the lower
bound is the upper one. The trick of the opposite, which works for the
arithmetic, would need $f_{\downarrow}(x) = -f_{\uparrow}(-x)$, and only for
the odd functions, which \code{acosh} is not. Its results of \code{asin},
\code{sinh}, \code{tanh} and \code{acosh} are therefore single doubles,
which miss $f(x)$ whenever it is not a double. The results depend on the C
library: glibc 2.31, under Ubuntu 20.04, missed in 41\,\% and 27\,\% of the
results, \code{asinh}, \code{atanh} and \code{acosh} included.
\end{Rem}

%%--------------------------------------------------------------------------
\section{ITF1788 and the benchmark of Tang et al.}
\label{sec:test-suites}
%%--------------------------------------------------------------------------
\cindex{ITF1788}\cindex{benchmark of Tang et al.}%

Two test suites written outside GAOL were run on GAOL v5 under the names of
the standard, \code{gaol\_ieee1788}, with the default build (SSE2 intervals
on x86, the machine of the tables being an Intel i7-1185G7) and the FPU build. The tests of GAOL
itself are the subject of Chapter~\ref{chap:tests}.

\paragraph{ITF1788.}
ITF1788, the Interval Test Framework for IEEE 1788 of Marco Nehmeier,
Maximilian Kiesner and Oliver Heimlich
(\url{https://github.com/oheim/ITF1788}), holds test cases of the
operations of IEEE 1788-2015: those of libieeep1788, of C-XSC, of FI\_LIB
and of MPFI, which Oliver Heimlich converted, and his own for \ieee{atan2},
\ieee{absRev}, \ieee{powRev1}, \ieee{powRev2}, the constructors and the
exceptions. They are written in a format of its own, one case per line of
its \file{.itl} files, such as
\code{add [1.0,2.0] [3.0,4.0] = [4.0,6.0];}. ITF1788 converts them into test
programs for several libraries; its plugin for GAOL, written for GAOL~4,
calls GAOL by GAOL's own names and leaves out most of the operations GAOL
v5 adds. The generator of \file{doc/compare/itf1788/} reads the
\file{.itl} files itself and writes a C++ program per file, each case being
one call of the operation by the name \code{gaol\_ieee1788} gives it.

\medskip
The suite holds 9542 cases in 19 files. 7415 pass: GAOL v5 gives the result
the case expects, the tightest one, the 1347 cases of \ieee{pow}, the 163
of \ieee{pown}, the 172 of \ieee{mulRevToPair} without decorations, the 48
of \ieee{overlap} without decorations, 802 of the 804 cases of
\ieee{powRev1} and \ieee{powRev2}, and 84 of the 94 cases of \ieee{pownRev}
with an exponent $p \leq 0$ included. No bug of GAOL v5 was found. The 22
failures are of three kinds: 12 results that are valid but not the
tightest, from \ieee{pownRev}, 2 at the largest double and 10 with an
exponent $p < 0$: 11 of them 1 double beyond the tightest bound and accurate in the sense of IEEE 1788-2015, and one 2 doubles beyond, valid but not accurate, the standard requiring of
this operation a valid result and recommending an accurate one; 8 cases where
ITF1788 expects a result that is not the tightest, from the reverse
functions \ieee{sinRev}, \ieee{cosRev} and \ieee{tanRev} of libieeep1788's
tests and \ieee{powRev2} of Oliver Heimlich's, where GAOL v5's result is
narrower than the expected one and encloses the exact result, which mpmath
computes; and 2 cases where ITF1788 expects an interval for a literal
\code{[l, u]} with $l$ above $u$, \code{[1.0000000000000002,1.0000000000000001]}
and the same bounds in hexadecimal, where GAOL v5 compares the numbers
exactly and gives the empty set, as the standard asks of an implementation
that decides $l \leq u$. 418 cases are not provided: eight trigonometric
functions of FI\_LIB and MPFI that are not in the standard (\code{cot},
\code{sec}, \code{csc}, their hyperbolic versions, \code{acot} and
\code{acoth}: 403) and the reduction operations (15). 1687 cases are not
applicable: they check the decorated intervals and NaI, the ill-formed interval of the
standard (11.3), which GAOL v5 does not have (see Chapter~\ref{chap:ieee1788}). The default and the FPU builds give the
same result in every case. The deliberate differences of GAOL v5 from the
standard do not show here: the hybrid \code{pow} is \code{gaol::pow}, while
\code{gaol\_ieee1788::pow} is the \code{pow} of the standard, and the
constructors of intervals with an infinite bound give the empty set where
the cases expect it.

\paragraph{The benchmark of Tang et al.}
The benchmark of interval libraries of Xuan Tang, Zachary Ferguson, Teseo
Schneider, Denis Zorin, Shoaib Kamil and Daniele Panozzo, \emph{A
Cross-Platform Benchmark for Interval Computation Libraries}
\cite{Tang:PPAM2022}, compares filib
(its C version), filib++, Boost.Interval and BIAS on 132 expressions: 28 of
the authors (the four operations, \code{sqrt}, \code{exp}, \code{sin} and
\code{cos}, ten arithmetic expressions of 6 to 10 variables and ten random
expressions with \code{sqrt}, \code{exp}, \code{sin} and \code{cos}) and 104
expressions of FPBench, on the inputs of their preconditions. Each
expression is evaluated on random inputs, each input variable being the
interval of one double; a result is correct when it contains the exact
value of the expression at the input, which the benchmark checks by printing
each result as a query with exact rationals. The paper finds every library
correct on the arithmetic, Boost.Interval wrong on \code{exp}, \code{sin}
and \code{cos}, BIAS wrong on one random expression, and filib++'s
\code{native\_switched} mode wrong on four FPBench expressions; GAOL is
cited but not run. A patch of \file{doc/compare/tang2021/} adds GAOL v5 to
it, and a checker replaces the Mathematica script of the paper.

\medskip
On 100\,000 inputs of each of the 132 expressions, 13\,200\,000 results in
each of the default and FPU builds of GAOL (CMake Release by Clang 21.1.8,
Intel i7-1185G7, Ubuntu 26.04, 14 October 2026), every interval of GAOL
encloses the exact value of the expression the program computes, the eight
single operations give the tightest interval every time (800\,000 of
800\,000), and GAOL's intervals are narrower than those of filib's C
version, the paper's reference, at all but 1\,488 inputs, where they are as
wide. The paper's own criterion fails at 9\,884 inputs, in six
FPBench expressions: each has a decimal constant that is not a double, which
the benchmark's query writes as the decimal while the program computes with
the double. These are a defect of the benchmark, not of GAOL; no bug of GAOL
was found. The two builds give the same tables, line for line. The speed
test of the benchmark, 1000 inputs of each expression each evaluated 10\,000
times, on GAOL and filib's C version built by Clang 21 in the SSE2 build, on
the same machine and day, with a load of about 1.5, is indicative only: GAOL
takes less time than filib's C version on 93 of the 132 expressions, 0.36
to 0.66 times its time on the four operations, 0.40 to 0.61 on the ten
arithmetic expressions, 0.55 for the median expression, and 107.9 s in all
against 131.7 s; it takes more on the elementary functions, 1.54 times on
\code{exp}, 1.77 on \code{sqrt}, 1.98 on \code{sin} and 3.98 on \code{cos},
whose inputs the speed test draws among $\pm(2^{31} - 1)$, where GAOL's
bounds are the tightest and filib's far wider on large arguments.

%%--------------------------------------------------------------------------
\section{The verdicts on new paths}
\label{sec:new-paths}
%%--------------------------------------------------------------------------
\cindex{paths of GAOL!verdicts}\cindex{RISC-V}\cindex{SVE}%

The sections above measured what GAOL has. This section draws the verdicts
on the tracks that were considered for new paths, lists the cheap gains
that need no new instruction set, states what a new path has to show
before it is added, and ends with the things to watch and what was not
measured. Table~\ref{tab:computing-verdicts} gives the verdicts.

\begingroup\footnotesize
\begin{longtable}{|L{2.3cm}|L{2.3cm}|L{5.0cm}|L{3.4cm}|}
\caption{The verdicts on the tracks considered for a path of GAOL (from
\file{doc/simd.md}).}\label{tab:computing-verdicts}\\
\hline
\textbf{Track} & \textbf{Verdict} & \textbf{Why} & \textbf{What would
change it} \\
\hline
\endfirsthead
\multicolumn{4}{l}{\textit{Table~\thetable, continued}}\\*
\hline
\textbf{Track} & \textbf{Verdict} & \textbf{Why} & \textbf{What would
change it} \\
\hline
\endhead
\hline
\endfoot
\endlastfoot
SSE4.1, AVX, AVX2 (a 256-bit path) & no path & A batch of two intervals per
register needs an interface of arrays, and nothing makes one interval
faster; \code{ROUNDPD} rounds to an integer, not as an operation & An
interface of boxes or arrays of intervals, and a program bound by the
computation and not by the memory \\
\hline
FMA3 & in use (\code{GAOL\_FMA}) & It pays in the elementary functions and
the integer powers, none in the basic operations & The exact
transformations that would replace the directed rounding are another study
(the rounding direction left to nearest) \\
\hline
AVX-512 path & keep; the division and the square root are scalar
instructions with the embedded rounding, two for the two bounds & On Intel
the 512-bit division takes 16 cycles of throughput against 4 for the
128-bit one (uops.info), and Zen 5 has full-width units; the square root of
the path takes two scalar embedded roots (3.75 ns against 6.30 with SSE2),
\code{x / y} two scalar divisions, and so do \code{x / d} of a double
(2.8 ns against 3.7 for $d > 0$, 3.0 to 3.3 for $d < 0$, on an Intel
i7-1185G7 with Clang 21), the divisors with a zero (5.5 ns) and the negative
integer powers (\code{pow(x, -3)} 9.1 ns), which the two scalar divisions, in place of the 512-bit one, left within the
noise & A processor where the two scalar divisions lose to the 512-bit one
(not seen on the i7-1185G7; no other processor was measured, Zen 5
included) \\
\hline
AVX-512VL, AVX10.2 at 128 and 256 bits & nothing to build & No embedded
rounding at these widths & A new revision of the Intel manual restoring it
(revision 3.0 of AVX10.2 had it) \\
\hline
NEON (arm64) & keep & Done: 1.3 to 2.5 times faster than the FPU intervals on an M2, 3.2 times on a line of 13 operations & Hardware with
FPCR.AH set (below) \\
\hline
SVE, SVE2, SME & no path & The same FPCR, and 128-bit vectors on most of the
processors & An interface of arrays, on a 256 or 512-bit implementation; a
static rounding in the instruction set \\
\hline
ARMv7 & no vector path & NEON has no double & none \\
\hline
POWER (VSX), s390x (vector facility) & measure first & Two doubles per
register and a write of the control register, as NEON, but no hardware here
and no figure for the write & A machine to run the tests and the benchmark
on \\
\hline
RISC-V, scalar & measure on a board, then write it & The only other rounding
direction by instruction, by inline assembly & A board, and a check that its
cores do not slow the static modes \\
\hline
RISC-V vector, LoongArch, MIPS, WebAssembly & no path & No static rounding
for the vectors, no user, or the rounding to nearest only & An extension
with a static rounding; for WebAssembly, a lifting of its limit \\
\hline
\end{longtable}
\endgroup

\subsection{Cheap gains that need no new instruction set}

The candidates below were seen in the sources and the profiles; the first
and the third are done, and none needs a new path.

\begin{itemize}
\item[$\bullet$] \textbf{\code{pow(x, 2)} through \code{sqr}} (done).
\code{gaol\_pown} took 5.2 ns where \code{sqr} takes 1.3 on the M2 (NEON
intervals; 5.2 and 1.8 with the FPU ones, whose \code{sqr} has no branch on
the signs), and the integer powers are about 40\,\% of Goldstein-Price (6
million calls). It takes \code{sqr} now, for $n = 2$ in every form of the
integer power:\newinvfive{} 1.6 ns with the NEON intervals and 1.9 with the
FPU ones, the call into the library accounting for the 0.3 ns of NEON above
\code{sqr} (0.1 with the FPU ones; the medians of 7 passes alternated with
the library before the change, \file{tests/performance.cpp}, an Apple M2,
AppleClang 21). The other powers cost 0.3 to 0.5 ns more with the FPU
intervals for $n = 3, 4, 5$ and 10 ($+5\,\%$), and less with the NEON ones.
On 14 million cases (3 million intervals, the special bounds, the empty set
and the universe, in the four directions of the calling code and with
flush-to-zero set) the bounds of \code{pow(x, 2)} are the same as before,
the sign of a zero lower bound aside, which is that of \code{sqr(x)} now.
\item[$\bullet$] \textbf{\code{sqrt} on SSE2} is scalar (6.3 ns on the
i7-1185G7, 3.75 with the AVX-512 path and the direction preserved, 2.0 on
the M2 with NEON); the NEON kernel ports to \code{\_mm\_sqrt\_pd}.
\code{sqrt\_rel} follows.
\item[$\bullet$] \textbf{The integer powers on NEON} (done, 1.1 to 1.3 times).
The kernel of \code{pown}, in \file{gaol/core/gaol\_interval\_neon.cpp},
computes on two lanes with \code{vfmaq\_f64}, as the SSE2 one of
\file{gaol\_interval\_sse.cpp} does with \code{\_\_FMA\_\_}, where it was
two scalar lanes with \code{std::fma}; the generic part is in
\file{gaol\_interval\_core.cpp}, and each kind of intervals defines the
operations of its lanes as a class.\newinvfive{} The bounds are the same bit
for bit (4 million powers hashed, and \file{tests/tools/pow/diffpow.py} on
70\,500 lines). The first version gained only 1.03 to 1.07 times, far below
the 1.5 times that was hoped for: \code{pow(x, 3)} took 5.7 ns where it took
6.1 (5 passes alternated by \file{.github/scripts/performance-ab.sh}, an
Apple M2, AppleClang 21). The kernel was not to blame: it is some 8 of the
125 instructions of a call, and its chain of dependent products is no
shorter than that of two scalar lanes (NEON has no vector \code{fnmsub}: each
fused multiply-subtract is a \code{vnegq} and a \code{vfmaq}). The rest was
the code before it, which Clang made longer: with \option{-frounding-math} it
turns the generic comparisons \code{vcgeq\_f64}, \code{vcleq\_f64} and
\code{vceqq\_f64} of the check of the range into scalar ones, 38
instructions where the same masks with \code{vcageq\_f64} and
\code{vcaleq\_f64}, on the absolute values, take 20, and it made the
permutation of the corners for a negative exponent (\code{vextq\_f64}) a
select through a general register. The check is now written with
\code{vcageq\_f64} and \code{vcaleq\_f64}, and the permutation keeps an empty
\code{asm} that leaves it as it is: a call of \code{pow(x, 3)} takes 102
instructions where it took 125. Without \option{GAOL\_PRESERVE\_ROUNDING}, it
takes 4.42 ns where it took 5.84 (1.32 times), and the other exponents from
$-7$ to 100 gain 1.11 to 1.33 times in independent calls (medians of 21
alternations of the old and the new code, an Apple M2, AppleClang 21). With
it, the 27 rows of the micro-benchmark take 0.95 to 0.99 of their time, a call
of \code{pow(x, 3)} taking 36.4 ns where it took 37.2: the preservation itself
adds some 31 ns to a call, whose cause was not looked for. Making the corners
from the stored register rather than from two stores and a load had not
changed the ratio of the first version.
\item[$\bullet$] \textbf{\code{mid()}} takes 8.5 to 12 ns on the M2 for an
addition at 0.9; a two-lane kernel would keep the rare branches.
\item[$\bullet$] \textbf{The argument reduction of \code{sin}, \code{cos}
and \code{tan}} (the quotients by $\pi$ and the \code{floor} and
\code{ceil}) is part of the code of GAOL around CORE-MATH, which is 42 to
66\,\% of the time of \code{sin} in the two programs that call it; no
profile says how much of that is the reduction.
\item[$\bullet$] \textbf{\code{div\_rel} and \code{\%=}} are trees of
branches; the NEON intervals have the cases of the function
\code{divide\_by\_zero\_containing()} as a table for the quotient.
\item[$\bullet$] \textbf{Not for SIMD.} \code{midpoint()} (29 ns on
the M2), \code{bisect()}, \code{chi()} and \code{hausdorff()} spend their
time in changes of the rounding direction, which is the question of
computing without changing it; the elementary functions are limited by
CORE-MATH, and a vector version would be another algorithm.
\end{itemize}

\subsection{What a new path has to show}

Before a path is added, with the rules of GAOL, it has to show:
\begin{enumerate}
\item its time against the path it replaces, on the processor and the
compiler that matter, with the spread, and the operation or the program
that needs it;
\item its bounds: the tightest ones, against references computed apart, in
the four rounding directions of the calling code, with and without
\option{GAOL\_PRESERVE\_ROUNDING}, the subnormals (the flush-to-zero modes
cleared), the NaN and the infinities, the empty set without an
invalid-operation exception, and the sign of a zero bound documented where
it differs;
\item the three build systems alike, a job of the continuous integration
that tells which path it took (\file{.github/scripts/check-intervals.sh},
\pxref{sec:ci}), and the processor it needs written with the platforms
(\pxref{sec:platforms}), with the path it falls back to.
\end{enumerate}

\subsection{Things to watch}

\begin{itemize}
\item[$\bullet$] \textbf{FPCR.AH.} With this bit of FEAT\_AFP (Armv8.7,
present on the Neoverse V3 of Graviton 5 and absent from the M2) set,
\code{FMAX} and \code{FMAXP} (the \code{vmaxq\_f64} and \code{vpmaxq\_f64}
of the NEON kernels) return the second operand for two zeros of different
signs, and when either is a NaN, in which case they also raise the
invalid-operation exception. GAOL clears FZ and FIZ and not AH. The kernels
never give a NaN to a maximum, and the zeros of the maxima of the product
differ only in sign. The square root used to depend on the order of the
operands of a maximum, of the interval with its left bound in its natural
sign and of $\langle +0, +0 \rangle$: a compiler that swapped the operands
would have turned \code{sqrt([-0, r])} into the empty set under AH (Apple
clang 21 and GCC 14 to its trunk keep the order, but Clang 21 merges
\code{vmaxq\_f64(x, y)} and \code{vmaxq\_f64(y, x)} into one instruction:
it takes the intrinsic for commutative). It now clears the lanes whose sign
bit is set, a negative number or $-0$, with a mask, as the part of the
interval in $\itv{0}{+\infty}$:\newinvfive{} the same bits as the maximum
with AH clear, on 20 million pairs of bounds that are no NaN and every pair
of the special values, and no dependence on AH or on the order of the
operands. The Linux and XNU kernels and the glibc, whose sources were read,
do not set AH by default; a program that does (an emulator) is outside what
GAOL supports. Modelled in software on 25.8 million results, the bounds are
those of the kernels, but for the sign of a zero.
\file{tests/rounding\_direction.cpp} runs the operations with AH set where
the processor retains the bit, and says so where it does not (the M2 and the
emulators QEMU 7.2 and 8.2 do not; QEMU 10.0 does, with
\cmd{qemu-aarch64 -cpu max}, where the test runs on 441 pairs of operands
per operation, and fails 189 checks with the operands of the former maximum
exchanged, the square root of \code{[-0, r]} being then a half-NaN
interval).
\item[$\bullet$] \textbf{The modes that make a result approximate or flush
it} are cleared by GAOL where it has a path that reads the control register:
FZ and FIZ on 64-bit ARM (FZ alone on 32-bit ARM), FTZ and DAZ on x86 (the
AVX-512 path does not read MXCSR, and clears them where a test shows them
set; \pxref{sec:ftz}). \code{NI} of FPSCR on POWER (non-IEEE mode, where
results ``may be approximate'') and \code{FS} of the FCSR on MIPS (flush
denormalized results to zero) are states of the same kind that GAOL does not
clear on the FPU intervals of these processors; a path for them would have
to.
\item[$\bullet$] \textbf{The pragma \code{FENV\_ROUND}, the
\code{\_mm256\_add\_round\_pd} of Clang 20 and the 256-bit
\code{\{ru-sae\}} assembly} compile and are wrong
(Section~\ref{sec:compilers-give}).
\item[$\bullet$] \textbf{The NEON quotient on a Neoverse N2} (the arm64
runners of GitHub, Azure Cobalt 100; the A/B of \file{performance.yml} of
7 November 2026, the medians of three runs, the base and the head giving the
same figures): \code{x / y} takes
3.96 ns with Clang, 4.04 with clang-cl, 4.49 with GCC and 4.92 with Visual
C++, where the FPU intervals take 4.29 with GCC, so that the one operation
the NEON intervals do not gain there is the quotient (with GCC \code{x * y}
takes 2.95 ns against 6.37 and \code{x + y} 1.26 against 1.87). It took 5.36
ns with clang-cl and 5.40 with Visual C++ before \code{divide\_upward()}
told the signs of the divisor on general registers (commit c42f95e, 5.35 to
3.96 ns with Clang 18), where the FPU intervals took 4.35 with clang-cl. It
was not the vector division: two scalar divisions in its place
(\code{vdiv\_f64} on each lane, the same bounds) took 6.17 and 6.30 ns, and
the time was in the rest of the kernel, about 150 instructions in
\code{div\_bounds()} of which some 55 were vector ones (counted in the
object of GCC 12.2, the common path being a part of them). On the Apple M2
(7 November 2026, mains power, three runs of each build alternated)
\code{x / y} takes 2.07 ns with the NEON intervals against 3.20 with the FPU
ones, the line of 7 quotients 18.4 ns against 32.4, and \code{x * y}
1.52 ns: the N2 takes 1.9 times as long for the product, a ratio that the
two vector pipes of the N2 against the four of the M2 would explain (not
measured: \cmd{perf} is not allowed on the runners and the machine).
\end{itemize}

\subsection{What was not measured}

No measurement of this chapter on x86 (the SSE2 and AVX-512 paths) was made on another processor than an Intel i7-1185G7, for
the times of Section~\ref{sec:computing-performance} and the divisions of
Table~\ref{tab:computing-verdicts}; no path was timed on RISC-V, POWER, s390x,
LoongArch, MIPS or WebAssembly (the tests of GAOL run on POWER, s390x and
RISC-V under QEMU in the continuous integration, \pxref{sec:ci}): the figures
given for these come from uops.info or the references, and the compilations made to study their instruction sets were made on
Compiler Explorer, without running the code. The calls
were counted on the NEON build only. GCC never ran a program of this study: the behaviour of the
intrinsics and of the pragma under it is read from the assembly. The profile
has one machine and one compiler, and \cmd{sample} is the macOS profiler.
The cost of writing FPSCR (32-bit ARM, POWER), FPC (s390x), FCSR and frm is
not in the references read, nor is the manual of MIPS. Whether FPCR.AH
changes the kernels was not measured on hardware: the M2 has no FEAT\_AFP (the
test runs with AH set under QEMU 10.0).

%%--------------------------------------------------------------------------
\section{Discussion}
%%--------------------------------------------------------------------------

This chapter has shown how GAOL v5 computes an interval: as one vector of two
doubles, \code{<-left, right>}, on which the SSE2 and the NEON kernels
compute both bounds with one or a few instructions rounded upward, with no
branch on the signs, the bounds staying in registers from one operation to
the next. What costs is not the arithmetic but the rounding direction: a
write of the control register takes about 3 ns on x86 and 4 to 11 ns on an
Apple M2, against one or two nanoseconds for the operation, and the choice
of GAOL v5 is to check the direction with one comparison and to write the
register only when the direction is not upward already. The measurements
bear the choice out. GAOL v5 is faster than GAOL 4.2.3 on every line measured and on every
operation but the square root, which takes the same time, 1.06 to 4.4 times
on x86 and 1.7 to 8.3 times on
the M2, and it is the fastest of the six libraries compared on the
arithmetic, the integer power and most formulas, with the bounds of
libieeep1788, which computes every bound with MPFR, 17 to 135 times faster.
Where it is slower, on \code{sin} and \code{cos} against filib++ and Solaris Studio, and on the square root, \code{exp}, \code{log} and the real power against PROFIL/BIAS, whose bounds are among the widest, it pays for a correct rounding over the whole range that the others
do not give: the test of enclosure found no result of GAOL v5 that misses
the value of the function, where 20 to 34\,\% of Boost.Interval's do.

\medskip
The trade-offs are those of the rounding direction. Preserving it costs
no write of the control register on the AVX-512 path, whose instructions carry their own direction,
and that is why a build takes that path with \option{GAOL\_PRESERVE\_ROUNDING}
only: without the preservation the operations take 20\,\% longer with the path at the median, and the
rule relies on one processor, with halved AVX-512 units. On 64-bit ARM, no
instruction carries its direction, and preserving it for a program that
computes to nearest makes a line of operations 13.0 to 43.5 times slower on
the M2. The paths that remain to try are few: the RISC-V scalar
instructions are the only other ones with a rounding direction by
instruction, and they need a board; POWER and s390x need hardware to
measure; the 128 and 256-bit forms of AVX-512 and the vectors of SVE offer
no embedded rounding. The cheap gains left open are a vector square root on
SSE2, a two-lane \code{mid()} and the argument reduction of the
trigonometric functions. Whether FPCR.AH changes the NEON kernels could not
be measured on hardware, the M2 having no FEAT\_AFP; the square root was made independent of it in software, and the other kernels differ under it only in the sign of a zero bound.

\medskip
These measurements rest on the tests, which check that every path gives the
same bounds in the four rounding directions of the calling code, with the
subnormals, the infinities and the empty set, and that a change does not
move the bounds by a double. The next chapter describes what the floating-point
numbers and the errors of GAOL look like to a program; the principles
and the programs of the tests, and the continuous integration that runs them
on every platform, are described in Chapter~\ref{chap:tests}.

%% !TEX root = ../manual_gaolv5.tex
% Chapter 13.
%% Chapter 13: Floating-point numbers, errors and debugging. Sources:
%% gaol.tex of f5bf180 lines 6007-6123 and 6156-6404 (the rounding functions
%% of the FPU, lines 6124-6155, are in Chapter 4), doc/using.md lines 931-998.
\chapter{Floating-point numbers, errors and debugging}%
\label{chap:errors}\label{chap:numbers}

An interval library rests on the floating-point numbers of its machine. The
bounds of the intervals of GAOL are \code{double} values, and a program that
builds an interval bound by bound, or that examines the result of a
computation, needs a few operations on these values: the neighbours of a
double, the number of doubles between two of them, a minimum and a maximum
that tell the two zeros apart. The first part of this chapter describes the
constants and the functions on \code{double} that GAOL defines beside its
intervals, the constants that tell a program the version of the headers it
was compiled with, and two utility functions on integers.

\medskip
The second part is about the errors. GAOL reports an error in one of two
ways, by throwing an exception or by printing a message and aborting, and the
way depends on how the library was built. GAOL v5 throws the exceptions of
the standard library, where GAOL~4 had classes of its own, so that a handler
that knows nothing of GAOL says what went wrong.\newinvfive{} We describe
the mechanism, the macro with which a program adds its own errors, the two
exceptions that GAOL throws and the text of their messages, and the warnings.
The last part presents the debugging facilities of the Debug builds.

\medskip
The chapter is organized as follows.
Section~\ref{sec:floating-point-constants} gives the floating-point constants
and Section~\ref{sec:fp-functions} the floating-point functions.
Section~\ref{sec:version} describes the version information and
Section~\ref{sec:additional-functions} the additional functions.
Section~\ref{sec:errors-handling} presents the error handling, the exceptions
and the warnings, and Section~\ref{sec:debugging-facilities} the debugging
facilities. Section~\ref{sec:errors-discussion} discusses the choices made.

\section{Floating-point constants}\label{sec:floating-point-constants}
\cindex{floating-point constants}\cindex{NaN}\cindex{infinity}%
\findex{GAOL\_NAN}\findex{GAOL\_INFINITY}%

In addition to the constants available through
\code{numeric\_limits<double>}, GAOL defines the two macros of the table below, which expand into \code{double} values.

\begin{center}
  \renewcommand{\arraystretch}{1.1}
  \begin{tabular}{lc}
    \hline
    \multicolumn{1}{c}{\bfseries Constant (\texttt{double})} &
    \multicolumn{1}{c}{\bfseries Value}\\
    \hline\hline
    \element{GAOL\_NAN} & NaN (quiet)\\
    \element{GAOL\_INFINITY} & $+\infty$ \\
    \hline
  \end{tabular}
\end{center}

\begin{Rem}
GAOL~4 also defined the doubles \code{two\_pi}, \code{pi}, \code{half\_pi},
\code{pi\_dn}, \code{pi\_up}, \code{half\_pi\_dn}, \code{half\_pi\_up},
\code{ln2\_dn}, \code{ln2\_up}, \code{two\_power\_51} and
\code{two\_power\_52}, in the namespace of GAOL. A program that opened it with
\code{using namespace gaol} found them among its own names, and a \code{pi}
of its own was ambiguous. GAOL v5 no longer declares them for the code that
uses it, apart from the four doubles around $\pi$ and $\pi/2$, which
\code{gaol} and \code{gaol\_ieee1788} name \code{gaol\_pi\_dn},
\code{gaol\_pi\_up}, \code{gaol\_half\_pi\_dn} and
\code{gaol\_half\_pi\_up}.\newinvfive{} The bounds of $\pi$, $\pi/2$ and $2\pi$ are those of
\code{interval::pi()}, \code{interval::half\_pi()} and
\code{interval::two\_pi()} (\pxref{sec:interval-constants}),
\code{interval::pi().left()} being \roundDn{\pi}; the bounds of $\ln 2$ are
those of \code{log(interval(2.0))}.
\end{Rem}

\section{Floating-point functions}\label{sec:fp-functions}
\cindex{floating-point functions}%

GAOL provides the following functions on \code{double} values. The functions
that set the rounding direction of the floating-point unit are described with
the rounding direction (\pxref{chap:fpu}).

\begin{deffun}{bool}{feven}{(double \var{x})}
  Returns \code{true} whenever \var{x} is even. The function is not meant for
  an infinity or a NaN.

  \begin{example}
cout << boolalpha << feven(3.0) << ' ' << feven(4.0) << ' ' << feven(3.5);
@outputs^false true false~
  \end{example}
\end{deffun}

\begin{deffun}{double}{next\_float}{(double \var{x})}
  Returns the smallest \code{double} greater than \var{x}, the neighbour
  \nextdouble{x} of \var{x}; \code{next\_float(+oo)} is $+\infty$, and \code{next\_float(DBL\_MAX)} is $+\infty$ with every C library.
\end{deffun}

\begin{deffun}{double}{previous\_float}{(double \var{x})}
  Returns the greatest \code{double} smaller than \var{x}, the neighbour
  \prevdouble{x} of \var{x}.
\end{deffun}

\begin{deffun}{bool}{is\_signed}{(double \var{x})}
  Returns \code{true} whenever the sign bit of \var{x} is set: for the negative numbers and $-0$. With \code{using namespace std;} and \code{using namespace gaol;}, as the examples of this manual assume, the name is ambiguous with the class template \code{std::is\_signed}: the call is written \code{gaol::is\_signed(x)}. No provision is made for a
  NaN: a program that wants to test for the negative numbers and $-0$ alone
  has to test also whether \var{x} is a NaN, with \code{std::isnan()}.
\end{deffun}

\begin{deffun}{double}{minimum}{(double \var{x}, double \var{y})}
  Returns the minimum \code{double} value of \var{x} and \var{y}. The
  function is commutative and returns $-0$ when comparing $-0$ and $+0$:
\begin{equation*}
\left\{\begin{array}{lll}
  \min(x,y) &=\min(y,x), &\forall x\neq\text{NaN}, \forall y\neq\text{NaN}\\
  \min(x,\text{NaN}) &=\min(\text{NaN},x)=\text{NaN}, &\forall x\\
  \min(-0,0) &=\min(0,-0)=-0
\end{array}\right.
\end{equation*}
\end{deffun}

\begin{deffun}{double}{maximum}{(double \var{x},double \var{y})}
  Returns the maximum \code{double} value of \var{x} and \var{y}. The
  function is commutative and returns $+0$ when comparing $-0$ and $+0$:
\begin{equation*}
\left\{\begin{array}{lll}
  \max(x,y) &= \max(y,x), &\forall x\neq\text{NaN}, \forall y\neq\text{NaN} \\
  \max(x,\text{NaN})&=\max(\text{NaN},x)=\text{NaN}, &\forall x\\
  \max(-0,0)&=\max(0,-0)=0
\end{array}\right.
\end{equation*}
\end{deffun}

\begin{deffun}{unsigned long long}{nb\_fp\_numbers}{(double \var{a}, double \var{b})}
  Returns the number of floating-point numbers in the interval
  \itv{\var{a}}{\var{b}}. In particular:
  \begin{itemize}
  \item \code{nb\_fp\_numbers(a,next\_float(a)) == 2};
  \item \code{nb\_fp\_numbers(a,a) == 1}.
  \end{itemize}
  The two zeros are the same number: \code{nb\_fp\_numbers(-0.0,0.0)} is 1,
  and \code{nb\_fp\_numbers(-0.0,1.0)}, \code{nb\_fp\_numbers(-1.0,0.0)} and
  \code{nb\_fp\_numbers(0.0,1.0)} are all $4607182418800017409$ (GAOL~4
  returned $13830554455654793217$ for the first two).\newinvfive

  The precondition is that \var{a} be lower than or equal to \var{b}.

  The result is an \code{unsigned long long}. GAOL~4 named its type with the
  macro \code{ULONGLONGINT}, which GAOL v5 no longer defines, as it prefixes
  with \code{GAOL\_} the macros of its headers.\newinvfive

  The function throws a \code{std::domain\_error} exception
  (\pxref{sec:exceptions}) if either \var{a} or \var{b} is a NaN or an
  infinity, or if \var{a} is greater than \var{b}; where the exceptions are
  disabled, it calls \code{gaol\_error} and aborts. It returns no number
  then, where the earlier editions of this manual gave
  \code{numeric\_limits<unsigned long long>::max()}.\newinvfive
\begin{example}
cout << nb_fp_numbers(1.0, next_float(1.0)) << ' '
     << nb_fp_numbers(1.0, 2.0);
@outputs^2 4503599627370497~
\end{example}
\end{deffun}

\section{Version information}\label{sec:version}
\cindex{version of the library}%

The headers of the library define four constants that give the version of
GAOL a program was compiled with, from the macros \code{GAOL\_MAJOR\_VERSION}, \code{GAOL\_MINOR\_VERSION}, \code{GAOL\_MICRO\_VERSION} and \code{GAOL\_VERSION} of \file{gaol/gaol\_configuration.h}, which a program tests with \code{\#if} (\pxref{sec:install-configuration-header}). The
versioning scheme is the one of the Apache Software Foundation, described at
\url{http://apr.apache.org/versioning.html}. The constants are declared in the namespace \code{gaol\_core} and named in
\code{gaol} and \code{gaol\_ieee1788}.\newinvfive

\begin{defconst}{unsigned int}{version\_major}
  Major version of the library: 5 for GAOL v5.
\end{defconst}

\begin{defconst}{unsigned int}{version\_minor}
  Minor version of the library.
\end{defconst}

\begin{defconst}{unsigned int}{version\_micro}
  Patch version of the library.
\end{defconst}

\begin{defconst}{const char *const}{version}
  Version of the library as a string.
\begin{example}
cout << gaol::version;
@outputs^@version~
\end{example}
\end{defconst}

\section{Additional functions}\label{sec:additional-functions}

The following functions are utility functions, not necessarily related to
intervals or to floating-point numbers.

\begin{defTfun}{template <typename T>}{bool}{odd}{(const \var{T}\& \var{x})}
  Returns \code{true} if \var{x} is odd and \code{false} otherwise. The type
  \var{T} may be any type that provides the operator \code{\&} (``bitwise
  \code{and}'') with the same semantics as the one of the \code{int}s.
\end{defTfun}

\begin{defTfun}{template <typename T>}{bool}{even}{(const \var{T}\& \var{x})}
  Returns \code{true} if \var{x} is even and \code{false} otherwise. The type
  \var{T} may be any type that provides the operator \code{\&} (``bitwise
  \code{and}'') with the same semantics as the one of the \code{int}s.
\end{defTfun}

\section{Error handling}\label{sec:errors-handling}
\cindex{errors}\cindex{error handling}%

A program that uses GAOL may see an error reported in two different ways:
\begin{itemize}
\item[$\bullet$] \textbf{By an exception.} GAOL throws an exception of the
  standard library, which the program may catch.
\item[$\bullet$] \textbf{By a message.} GAOL prints a message on the
  standard error, and the program aborts.
\end{itemize}

\noindent
The exceptions are those of the standard library, declared in
\code{<stdexcept>}. GAOL~4 had classes of its own, \code{gaol\_exception},
\code{input\_format\_error}, \code{invalid\_action\_error} and
\code{unavailable\_feature\_error}, which GAOL v5 no longer has
(\pxref{sec:exceptions}).\newinvfive

\medskip
The mechanism in use depends on the way the library is configured
(Chapter~\ref{chap:installation}). With the option
\loption{enable-exceptions=yes} of \cmd{configure}, the default, every error of an operation
is reported by an exception; an AVX-512 build started on a processor without these instructions prints a message and aborts in every build. A build where
\code{GAOL\_EXCEPTIONS\_ENABLED} is 0, obtained with \loption{disable-exceptions}
of \cmd{configure}, \option{-Denable-exception=false} of meson or
\option{-DGAOL\_EXCEPTIONS\_ENABLED=OFF} of CMake, prints a
message on the standard error and aborts instead of throwing. This is a setting
of the developers, for the tests of that path. Relying on exceptions is more in the
spirit of C++, though it may incur some overhead.

\begin{Rem}
In a build that aborts, the reader of strings is an exception to this rule. An error it
finds as it computes, a wrong call (\code{nth\_root(8, 1.5)},
\code{sin(1, 2)}), an uncertain number as a bound, two different bounds in
\code{<a, b>} or a bound of it that is no double make it fail as a syntax
error does: \code{parse\_interval()} returns \code{false}, and
\code{gaol\_ieee1788::textToInterval()} the empty set, as IEEE 1788-2015 has
it (12.12.7).
\code{gaol::textToInterval()} and the operator \code{>{}>} print a message
and abort then, as for any string that is no interval.\newinvfive
\end{Rem}

\medskip
Where memory runs out, GAOL throws the \code{std::bad\_alloc} of C++. The
reader of strings does too, and frees what it allocated; in GAOL~4, flex ended the program with \code{exit(2)} when it could not allocate its buffer, and the parser read a stack it could not grow as a string nested too deep.\newinvfive

\medskip
It is up to the user to comply with this mechanism when adding error
reporting code to a program. GAOL defines the following macro, to be used
whenever one wants to report an error.

\begin{defmacroA}{gaol\_ERROR}{(\var{excep},\var{msg})}
  The behaviour of the macro depends on the value chosen for the option
  \loption{enable-exceptions}. If the exceptions are enabled, an exception of
  the class \var{excep} is thrown with the text
  ``\code{file:line: }\var{msg}'', \code{file} and \code{line} being those of
  the place where the macro is written (\code{\_\_FILE\_\_}, cut after its last
  \code{/} or \code{\textbackslash}, and \code{\_\_LINE\_\_}; \code{error\_text()}
  of \file{gaol/gaol\_common.h}): the path of the machine that built GAOL is not
  in the message. The class \var{excep} derives from \code{std::exception},
  in practice one of \code{<stdexcept>}; the compiler refuses any other type,
  by a \code{static\_assert}, in a build that aborts as in one that throws.
  Otherwise, the program aborts with the message \var{msg}.\newinvfive

\begin{example}
interval x;

@emph^[Code manipulating x]~

if (x.is_empty()) {
   gaol_ERROR(std::domain_error,"Emptiness of one interval");
}
\end{example}
\end{defmacroA}

The function \code{gaol\_error()} is defined as follows.

\deffunx{void}{gaol\_error}{(const char *\var{err})}
\begin{deffun}{void}{gaol\_error}{(const char *\var{file}, int \var{line},
const char *\var{err})}
Displays a message on the standard error output. The ternary version should
be called with the macro \code{GAOL\_FILE\_POS} for its first two parameters
(\pxref{sec:exceptions}). A null \var{file} is written
\code{(unknown file)}, and a null \var{err} is an empty text.\newinvfive
\end{deffun}

\subsection{Exceptions}\label{sec:exceptions}
\cindex{exceptions}\findex{std::invalid\_argument}\findex{std::domain\_error}%
\findex{std::logic\_error}\findex{std::bad\_alloc}%

GAOL throws the exceptions of the standard library, and none of its
own.\newinvfive{} It throws two of them to report an error:
\begin{itemize}
\item \code{std::invalid\_argument}, when a string is no interval
  (\pxref{sec:invalid_argument});
\item \code{std::domain\_error}, when a function is called with an argument
  it does not take (\pxref{sec:domain_error}), as \code{bisect()} is for a
  ratio outside $(0, 1)$.
\end{itemize}

\noindent
Both derive from \code{std::logic\_error}, which derives from
\code{std::exception}. A single handler of \code{std::logic\_error} catches
both, but the handler of one leaves the other alone, and none of them
catches the \code{std::bad\_alloc} that GAOL throws where memory runs out.

\medskip
The \code{what()} of these exceptions is
``\code{file:line: }\var{explanation}'': the name of the source file of GAOL,
without its directories, and the line that threw the exception, then what
went wrong, as in \code{gaol\_interval.cpp:48: Syntax error in interval
initialization: [1, 2}. For the errors that the reader of strings finds in a
call (\code{nth\_root(8, 1.5)}, \code{sin(1, 2)}), the file is that of the
grammar, \code{gaol\_interval\_parser.ypp}, and the line one of that file,
which the directives \code{\#line} of the generated parser give. A handler
of \code{std::exception}, and the message of an exception that nothing
catches, say what went wrong, where GAOL~4 gave \code{std::exception}
whatever the error. The text is valid as long as the exception is.

\begin{Ex}
The program below gives the reader of strings a text that is no interval,
and catches the exception as a \code{std::exception}, knowing nothing of
GAOL.
\begin{example}
try {
  interval x = gaol::textToInterval("[1, 2");
} catch (const std::exception& e) {
  std::cerr << e.what() << '\n';
}
\end{example}
The handler writes \code{gaol\_interval.cpp:48: Syntax error in interval
initialization: [1, 2} on the standard error. Had nothing caught the
exception, the program would have ended with the same text; with libstdc++,
the output is the following.
\begin{onscreen}
terminate called after throwing an instance of 'std::invalid_argument'
  what():  gaol_interval.cpp:48: Syntax error in interval initialization: [1, 2
\end{onscreen}
\end{Ex}

\begin{defmacro}{GAOL\_FILE\_POS}
Expands into the first two arguments of \code{gaol\_error()} and
\code{gaol\_warning()}, the name of the file and the line:

\begin{example}
if (@emph^[some condition]~) {
  gaol_error(GAOL_FILE_POS, "No additional information");
}
\end{example}
\end{defmacro}

\begin{Rem}
What GAOL~4 had is gone: the base class \code{gaol\_exception}, with its
accessors \code{file()}, \code{line()} and \code{explanation()} and the
operator ``\code{<{}<}'' that wrote them, and the classes
\code{input\_format\_error}, \code{invalid\_action\_error} and
\code{unavailable\_feature\_error} of \file{gaol/gaol\_exceptions.h}. The
classes become those of the standard library:
\begin{itemize}
\item \code{gaol::input\_format\_error} becomes
  \code{std::invalid\_argument};
\item \code{gaol::invalid\_action\_error} becomes
  \code{std::domain\_error}.
\end{itemize}
\noindent
What \code{explanation()} gave follows the file and the line in
\code{what()}. \code{unavailable\_feature\_error} has no replacement: \code{atan2()}, which threw it in GAOL~4, is implemented.\newinvfive
\end{Rem}

\subsection{The \code{std::invalid\_argument} exception}%
\label{sec:invalid_argument}

The exception \code{std::invalid\_argument} is thrown whenever a program
attempts to create an interval from an invalid string. This situation may
occur when reading an interval from a stream with the operator \code{>{}>},
or when reading a string with \code{gaol::textToInterval()}, from one string
or from two, a null pointer being no string (\pxref{sec:reading-intervals}).
The example prints what follows the first colon and the blank of
\code{what()}, the explanation, which the file and the line precede.

\begin{example}
try {
  interval x = textToInterval("sin(1)+");
} catch (const std::invalid_argument& e) {
  string text = e.what();
  cout << text.substr(text.find(": ") + 2);
}
@outputs^Syntax error in interval initialization: sin(1)+~
\end{example}

\subsection{The \code{std::domain\_error} exception}\label{sec:domain_error}

The exception \code{std::domain\_error} is thrown in three cases only:
\code{nb\_fp\_numbers()} with NaNs as parameters (\pxref{sec:fp-functions}),
\code{bisect()} with a ratio outside $(0, 1)$ (\pxref{sec:splitting}), and a
string that gives a function an argument it does not take, as
\code{"nth\_root(8, 1.5)"} does.

\subsection{Warnings}\label{sec:warnings}
\cindex{warnings}%

GAOL writes no warning itself; the following function writes one, as
\code{gaol\_error()} writes an error, on the standard error output
\code{std::cerr}, for the code that uses GAOL.

\deffunx{void}{gaol\_warning}{(const char *\var{warn})}
\begin{deffun}{void}{gaol\_warning}{(const char *\var{file}, int \var{line}, const char *\var{warn})}
Prints the message \var{warn} on the standard error output. The second form
should be called with the macro \code{GAOL\_FILE\_POS} for its first two
parameters. A null \var{file} is written \code{(unknown file)}, and a null
\var{warn} is an empty text.\newinvfive
\end{deffun}

\section{Debugging facilities}\label{sec:debugging-facilities}
\cindex{debugging}\findex{GAOL\_DEBUGGING}\findex{NDEBUG}%

The debugging facilities described hereunder are available only when GAOL
has been configured with them. The Debug build of each build enables them:
\option{-DCMAKE\_BUILD\_TYPE=Debug} with CMake, the option
\loption{enable-debug} of \cmd{configure} (\pxref{sec:configuration}), or
\loption{buildtype=debug} with meson; \option{-DGAOL\_DEBUG=ON} with CMake
and \option{-Denable-debug=true} with meson enable them whatever the build
type.\newinvfive{} These builds define \code{GAOL\_DEBUGGING} to 1 in
\file{gaol/gaol\_configuration.h} (with a generator of several configurations, Visual Studio or Xcode, only for the sources of GAOL in the configuration Debug: the header keeps 0, and \code{GAOL\_DEBUG\_VERBOSE} expands to nothing in the program), where it is 0 otherwise, as every option
of that file, so that it is read with \code{\#if} (\code{\#ifdef} is true of
it in every build). The Debug builds compile without \code{NDEBUG} too, so
that \code{assert()} stays active.\newinvfive

\begin{deffun}{int}{debug\_level}{}
  Global variable that remembers the current value of the debugging level.
  The variable is set when the library is initialized, by \code{init()}
  (\pxref{chap:init}). It is declared in the namespace \code{gaol\_core},
  and named in \code{gaol} too.
\end{deffun}

\begin{defmacroA}{GAOL\_DEBUG\_VERBOSE}{(\var{lvl},\var{cmd})}
  Executes \var{cmd} if \var{lvl} is lower than or equal to the current
  debugging level, the variable \code{debug\_level} above, which the macro names unqualified: it is used where \code{gaol} or \code{gaol\_core} is open. The macro was
  named \code{GAOL\_DEBUG} before GAOL v5; this name is now the option of
  CMake that builds GAOL for debugging.\newinvfive

  The macro expands to nothing when the library was not configured with the
  debugging facilities.
\end{defmacroA}

\noindent
A possible use of this macro is the following:

\begin{example}
interval x(-10,10);

@emph^[Some code]~

GAOL_DEBUG_VERBOSE(1,cout << "The value of x is " << x);
x += interval(3.5,4.5);
GAOL_DEBUG_VERBOSE(2,cout << "Now the value of x is " << x);
\end{example}

\noindent
The first message is displayed whenever the debugging level is at least 1,
and the second one only if it is at least 2.

\section{Discussion}\label{sec:errors-discussion}

This chapter gathered what GAOL v5 offers around its intervals. The
constants and the functions on \code{double} let a program reason on the
bounds themselves: the neighbours of a double, the number of doubles an
interval holds, where the two zeros count as one, and a minimum and a
maximum that order $-0$ before $+0$ and propagate a NaN. The doubles of
GAOL~4 that named $\pi$ and $\ln 2$ left the namespace, so that a program
opening it keeps a \code{pi} of its own; the constants of the class
\code{interval} replace them.\newinvfive{} Four constants give the version
of the headers a program was compiled with.

\medskip
On the errors, GAOL v5 made one choice: it throws the exceptions of the
standard library, \code{std::invalid\_argument} for a string that is no
interval and \code{std::domain\_error} for an argument a function does not
take, and none of its own.\newinvfive{} A handler that knows nothing of GAOL
reads in \code{what()} the file and the line of GAOL and what went wrong,
and a program written for GAOL~4 changes the classes it catches
(Chapter~\ref{chap:changes}). The build that aborts instead of throwing
remains available in the three builds, as a setting of the developers
(\option{-DGAOL\_EXCEPTIONS\_ENABLED=OFF} with CMake); in that build, an error that the reader of strings finds
makes the reading fail, so that \code{gaol\_ieee1788::textToInterval()}
returns the empty set as IEEE 1788-2015 has it, rather than end the program.
The debugging facilities belong to the Debug builds, and
\code{GAOL\_DEBUG\_VERBOSE} expands to nothing elsewhere.

\medskip
The next chapter describes the tests of GAOL v5 and its continuous
integration, which exercise the exceptions of this chapter among the rest of
the library, on every platform GAOL supports.

%% !TEX root = ../manual_gaolv5.tex
% Chapter 14.
%% Chapter 14: Tests and continuous integration. Sources: gaol.tex of f5bf180 lines
%% 778-816 (the tests); doc/tests.md, doc/continuous-integration.md,
%% doc/tools.md, examples/examples.md (section 1), examples/tutorial.md,
%% coverage/README.md, manual/check_examples.py (its header comment).
\chapter{Tests and continuous integration}\label{chap:tests}
\cindex{tests}%
\cindex{continuous integration}%

An interval library is trusted for one property: the bounds it computes
enclose the exact result of each operation, on every machine it is built
for. A test that compares GAOL with itself, or with a floating-point library
running in the same environment, cannot establish it, since the reference
would share the roundings and the rounding direction of what it checks. The
tests of GAOL v5 therefore compare the bounds with results computed apart
from the library and from the floating-point environment, with integers,
with exact rational arithmetic, or with mpmath at 2000 bits of precision, and
they run the library in every state a program may leave the processor in.

In this chapter, we describe how GAOL v5 is tested. The first section gives
the principles the whole suite follows: exact references, tests that fail
without the change they check, the floating-point exceptions, the four
rounding directions, the modes that flush the subnormal numbers, the locales,
and the unit tests of GAOL~4. The second section lists the programs of
\file{tests/} and what each of them checks. The third section presents the
example programs of \file{examples/}, which check what they print, and the
script that compiles the examples of this manual. The fourth section
describes the continuous integration: its platforms, the way it measures the
speed of a pull request, the configurations it refuses or leaves out, and
what these platforms showed. The fifth section lists the scripts of the
repository, and the sixth gives the coverage of the tests. The targets that
build and run the tests in the three builds are described with the
installation (\pxref{sec:tests}).

\section{The principles of the tests}\label{sec:test-principles}
\cindex{tests!principles}%

Built with \option{WITH\_TESTS} (CMake), \loption{with-tests}
(\cmd{configure}) or \code{with-tests} (meson), all off by default, the
programs of \file{tests/} compare the bounds GAOL computes with the exact
results of the operations, independently of GAOL and of the floating-point
environment. The exact results are computed with integers, or were computed
with 2000 bits of precision by mpmath. The tests follow the rounding tests
of Codac. They check, among others, that the basic operations are the
tightest enclosures, that the elementary functions are the tightest or
within one double of them, that every operation gives the same result
whatever the rounding direction the calling code left, that numbers are read
exactly and intervals written rounded outward, that the reverse functions
are accurate in the sense of IEEE 1788-2015, and that each name of
\code{gaol\_ieee1788} is the operation of the standard it names.
\code{test\_ibex\_codac} is another kind of test: it checks GAOL against the
cases of the unit tests that IBEX and Codac run on their own intervals,
which are GAOL's. \cmd{make test} runs the tests, and \cmd{make check} runs
the examples of \file{examples/} too (\pxref{sec:tests}).

\medskip
The suite rests on the following principles.

\begin{itemize}
\item[$\bullet$] \textbf{References computed apart.} A reference value is
  never a result of GAOL, nor of the C library running in the same rounding
  direction. It is an integer, an exact rational, or a value computed by
  mpmath with 2000 bits (400 to 7000 for some tables) and written into a
  header that the test includes: \file{elementary\_values.h},
  \file{reverse\_values.h}, \file{atan2\_rel\_values.h},
  \file{pow\_rel\_values.h}, \file{compoundm1\_values.h} and
  \file{extended\_precision\_values.h} are generated by the scripts
  \file{*\_values.py} of \file{tests/}, run by hand when a table changes
  (\pxref{sec:scripts}). Where a function is correctly rounded by CORE-MATH,
  the test computes its tightest bounds by calling CORE-MATH itself in the
  downward and the upward directions, which it sets apart from GAOL.
\item[$\bullet$] \textbf{Tests that fail without the fix.} Each check
  written with a change of GAOL v5 was run against the sources before the
  change, and the documentation of the tests records how many of its checks
  fail there: 10\,900 of the checks under the modes that flush the subnormal
  numbers, and, with the SSE2 intervals, 45 of the 48 calls that compared the NaN
  bounds of an empty operand, which died in a child process that enabled the
  exception, and 25 of the 45 checks of the operations computed before
  \code{main()}. A check that passes with and without the change proves
  nothing, and the mutation tools of \file{tests/tools/} build GAOL with a
  deliberate error that the tests have to catch (\pxref{sec:scripts}).
\item[$\bullet$] \textbf{The floating-point exceptions.} No operation of GAOL
  may raise the invalid-operation, the division-by-zero or the overflow flag
  for a result without an infinite bound, nor lose a flag or a mask the
  program set (\pxref{sec:fp-exceptions}). With glibc, each operation
  checked is also computed in a child process that enabled the exception
  with \code{feenableexcept()}, which must not die on \code{SIGFPE}: the
  comparison of the NaN bounds of an empty interval once did. Where the
  processor does not trap, or keeps no flag, the test says so and skips
  that part.
\item[$\bullet$] \textbf{The four rounding directions.} The bounds must not
  depend on the rounding direction of the calling code: each operation is
  called with the direction set upward, to nearest, downward and toward
  zero (and on x86 with the x87 and the SSE directions differing, six
  states), and has to give the results it gives when called rounding
  upward, and to leave the direction upward, or as it found it with
  \option{GAOL\_PRESERVE\_ROUNDING} (\pxref{sec:rounding-direction}).
\item[$\bullet$] \textbf{The flush-to-zero modes.} Each mode that flushes the
  subnormal numbers to zero that the processor honours (flush-to-zero and
  denormals-are-zero on x86, FZ and FIZ on ARM, \pxref{sec:ftz}) is set just
  before an operation, alone and with the others: the operation has to give
  what it gives with the modes cleared, and to leave them cleared, or set
  back with \option{GAOL\_PRESERVE\_ROUNDING}. A mode the processor keeps
  without honouring it, as an emulator may, is named and skipped.
\item[$\bullet$] \textbf{The locales.} Numbers have to be read as in the C
  locale, and written back bit for bit, under a locale writing a decimal
  comma (\code{fr\_FR.UTF-8}, \code{de\_DE.UTF-8},
  \code{French\_France.1252}) where one is installed, and the names of the
  functions have to be read in any case of letters under the Turkish
  locales of glibc, where \code{std::tolower('I')} is not \code{i}. The
  reading once never ended under a comma locale, and the test, which ctest
  would otherwise let run with no limit, fails after 5 minutes should it
  hang again.
\item[$\bullet$] \textbf{The unit tests of GAOL~4.} The tests of
  \file{check/}, which checked the results of the operations on chosen
  intervals with CppUnit, are in \file{tests/} and run with the checks of
  \file{unit\_tests.h}: CppUnit is needed no more.\newinvfive{} Each suite
  keeps its class and its tests, a check is named after its test and
  describes a failure by its line, its expression and the values compared,
  and a failure no longer ends its test.
\end{itemize}

\begin{Rem}
Under the Debug C runtime of Visual C++, a failed assertion of the runtime
or of the checked iterators of its library, and an invalid parameter, are
written to the standard error, where the test goes on and fails
(\file{tests/gaol\_tests.h}). The runtime showed them in a dialog box, which
nobody closes on a machine of the continuous integration, and the test hung
without output until ctest stopped it after 300~s. That runtime reports such
an assertion when it writes a subnormal number under denormals-are-zero, then
writes 0: \code{operator<<} clears that mode before it writes, which
\code{numbers} checks there too.\newinvfive
\end{Rem}

\begin{Rem}
Built without exceptions (\loption{disable-exceptions} of \cmd{configure},
\option{-Denable-exception=false} of meson), where an error of GAOL ends the
program, the tests leave out the checks of an exception and of a text the
reader refuses.\newinvfive
\end{Rem}

\medskip
\file{tests/find\_package} builds some of the same tests with an installed
GAOL, and \file{.github/scripts/tests.sh} with a GAOL installed by
\cmd{configure} or meson. \file{tests/fetch\_content} builds them with a GAOL
brought into a project by \code{FetchContent}, and
\file{tests/meson\_subproject} with a GAOL brought in as a meson
subproject,\newinvfive{} which meson refused while \file{meson.build} called \code{add\_global\_arguments()}. Its test \code{flags} checks
that \code{gaol\_dep} gives the compilation the flags of the \code{Cflags} of
\file{gaol.pc}, without which the tests pass all the same, and the link the
options of its \code{Libs} (\option{-mno-daz-ftz}, where the compiler takes
it); its test \code{release} checks that GAOL is built in release in the
default build type of meson, debug: \file{tests/release\_flags.py} checks
that the sources of GAOL and CORE-MATH are compiled with \option{-O3} and
\code{NDEBUG}, and that \code{GAOL\_DEBUGGING} is 0. With \loption{debug},
it checks the other way round that \option{GAOL\_DEBUG} of CMake and
\code{enable-debug} of meson build GAOL and CORE-MATH for debugging, GAOL
built alone and brought into another project.

\section{The programs of the test suite}\label{sec:test-programs}
\cindex{tests!programs}%

Table~\ref{tab:test-programs} lists the programs of \file{tests/} and what
each of them checks. The paragraphs that follow give, for each group, the
cases and the numbers of checks worth knowing, and the platforms where a
check applies.

\begingroup\renewcommand{\arraystretch}{1.1}\small
\begin{longtable}{L{4.4cm}L{9.4cm}}
  \caption{The programs of the test suite}\label{tab:test-programs}\\
  \hline
  \multicolumn{1}{c}{Program} & \multicolumn{1}{c}{What it checks}\\
  \hline\hline
  \endfirsthead
  \multicolumn{2}{l}{\textit{Table~\thetable, continued}}\\*
  \hline
  \multicolumn{1}{c}{Program} & \multicolumn{1}{c}{What it checks}\\
  \hline\hline
  \endhead
  \hline
  \endfoot
  \code{arithmetic} & The basic operations, squares, square roots, integer
    powers, \code{fma} and the cancellative operations are the tightest
    enclosures at every magnitude, from every rounding direction, and the
    $n$-th roots are within 2 doubles of them (4 for an inverse root)\\
  \code{elementary} & The elementary functions at doubles, at intervals and
    at intervals whose images are known exactly; \code{pow}, \code{atan2},
    \code{erf}, \code{erfc}, the integer functions\\
  \code{rounding\_direction} & The same bounds from every rounding direction
    and flush mode; the direction left upward or preserved; the control
    registers; the flags and the masks kept; no \code{SIGFPE} on an empty
    operand\\
  \code{fast\_math\_link} & A program linked with \option{-ffast-math}: the
    flush modes that \file{crtfastmath.o} sets\\
  \code{automatic\_cleanup} & (Linux) The rounding direction set back to
    nearest after \code{main()}\\
  \code{static\_initialization} & About 40 operations computed before
    \code{main()} give their bounds of \code{main()}, bit for bit\\
  \code{test\_ibex\_codac} & The cases of the unit tests of IBEX and Codac,
    against the exact results: 3905 checks\\
  \code{numbers} & Reading and writing: the tightest reading of every text,
    the literals of IEEE 1788-2015, the bases 16, 8 and 2, the five output
    formats, the flags of the streams, the locales, the subnormals under
    the flush modes\\
  \code{other\_functions} & Midpoints, widths, radii, distances, splitting,
    \code{max}, \code{min}, the tools of interval algorithms, the
    comparisons of IEEE 1788-2015, the reverse functions of \code{max},
    \code{min}, \code{sign} and \code{floor}\\
  \code{ieee1788} & \code{gaol\_ieee1788} as a program uses it: \code{pow},
    \code{pown}, \code{rootn}, \code{pownRev}, \code{overlap},
    \code{intervalToText}, \code{mulRevToPair}\\
  \code{core\_math} & The bounds against CORE-MATH called in the rounding
    direction in effect; the \code{fma()} and \code{round()} of the C library, and
    \code{gaol\_ldexp()}; the accurate paths of \code{sin} and
    \code{cos}\\
  \code{compoundm1} & $(1+x)^y-1$ from the bottom up: the 256-bit integers,
    the balls, the fast phase, 5\,854 pairs and 850 boxes\\
  \code{expressions} & The lexer, the parser and the evaluation of the
    strings: every node, every wrong string, the exceptions, the deep
    chains, the allocations that fail\\
  \code{u128} & The two 64-bit halves of the 128-bit integer against the
    native type of the compiler\\
  \code{extended\_precision} & The arguments the x87 unit rounds wrongly:
    fails when GAOL is built on it\\
  \code{reverse} & The relational functions against the reverse functions
    of IEEE 1788-2015: valid and accurate; \code{atan2\_rel\_x},
    \code{atan2\_rel\_y}, \code{pow\_rel}, \code{pow\_exponent\_rel}\\
  \code{debugging} & The headers compiled with \code{GAOL\_DEBUGGING}, and
    \code{GAOL\_DEBUG\_VERBOSE} at the level given to \code{gaol::init()}\\
  \code{version\_file} & (CMake) The reading of \file{VERSION.txt}: byte
    order mark, line ends of Windows, UTF-16, NUL bytes\\
  \code{refused\_finite\_\allowbreak math\_only},\newline
    \code{refused\_fast\_math},\newline
    \code{refused\_positive};\newline
    \code{refused\_options.sh} (autotools);\newline
    \code{refused\_options} (meson) & The compilation refused with
    \option{-ffinite-math-only} and \option{-ffast-math}, accepted
    without; \file{refused\_options.sh} and \code{refused\_options} check as well
    the three cases of the tests \code{refused\_aligned\_new*}\\
  \code{refused\_aligned\_new},\newline
    \code{refused\_aligned\_\allowbreak new\_cxx17},\newline
    \code{refused\_aligned\_\allowbreak new\_positive} (CMake) & The SSE2 intervals on
    32-bit Windows refused without \option{-faligned-new=8}\\
  \code{fp\_strict} & (a project of its own, which the Visual Studio jobs
    configure) Visual C++ and clang-cl refused without \option{/fp:strict}\\
  \code{cpack\_stale\_configure} & (CMake) The warning of CPack when the
    committed \file{configure} was generated for another version\\
  \code{automatic\_configuration} & (CMake) The configuration GAOL decides
    by itself, and the settings of the developers\\
  \code{arithmetic\_operators}, \code{assignment}, \code{constants},
    \code{constructor}, \code{float\_functions},
    \code{generic\_functions}, \code{input\_output},
    \code{interval\_functions}, \code{misc}, \code{non\_arithmetic},
    \code{relations}, \code{reverse\_mappings}, \code{trigonometric} & The
    unit tests of GAOL~4: the results of the operations on chosen
    intervals; \code{misc} checks the names of the namespaces and the
    alignment of the intervals\\
  \code{intervalf}, \code{interval2f} & The intervals of floats, where a
    developer compiles them; skipped otherwise\\
  \code{nodiscard\_\allowbreak discard\_cxx11},\newline
    \code{\dots\_cxx14}, \code{\dots\_cxx17};\newline
    \code{nodiscard\_\allowbreak used\_cxx11},\newline
    \code{\dots\_cxx14}, \code{\dots\_cxx17} & A result of GAOL thrown away
    fails the compilation, in the three standards; a result used compiles\\
  \code{gaol\_performance} & (not a test) The time per operation of GAOL's
    arithmetic and elementary functions\\
\end{longtable}
\endgroup

\subsection{The arithmetic and the elementary functions}

\code{arithmetic} checks, on doubles and intervals of every magnitude,
subnormal doubles and overflows included, that sums, differences, products,
quotients, relational divisions, squares, inverses, \code{abs}, \code{min},
\code{max}, \code{\&} and \code{|} are the tightest enclosures, and so are
square roots and integer powers of either sign:\newinvfive{} of doubles for
$n$ from 2 to 9, 16 and 31 to 33 and for $-1$ to $-4$, $-7$, $-16$ and
$-33$, of intervals for $n$ from 2 to 5, 12 and 31 and for $-1$ to $-3$,
$-6$ and $-17$, $x^n$ beyond the largest double or subnormal included, and
for $n<0$ the limits where $x$ holds 0. The suite once allowed up to $2n$
doubles beyond the tightest bounds for $n>0$, and 8 for $n<0$. Every
interval made of 15 special bounds ($-\infty$, the opposite of the largest double, $-2.5$,
$-1$, the opposites of the smallest normal and of the smallest subnormal
number, $-0$, $+0$, the smallest subnormal and normal numbers, $0.1$, 1, 3,
the largest double, $+\infty$), the empty set and the universe, taken two by
two, have to give with \code{+}, \code{-}, \code{*}, \code{/}, \code{\%},
their compound forms (the interval itself as the operand included), with a
double on either side where an interval is a point, and for \code{sqr},
\code{sqrt} and \code{inverse}, the bounds of a reference of plain doubles:
the least of the corners rounded downward and the greatest rounded upward by
the processor, a zero bound times an infinite one counting as 0. An infinite
or NaN double gives the empty set, no operation gives an unexpected empty
set, and the bounds do not depend on the rounding direction of the calling
code, set to each of the four directions before each operation. Inexact
operations on literals the compiler sees (\code{[0.1] + [0.2]},
\code{[1] / [3]}, \code{sqrt([2])}) have to be the tightest enclosures
whatever the direction: an operation computing to nearest fails them, and so
would a compiler that folded the vectors of the library at compile time,
which none does with the \option{-frounding-math} that \code{gaol::gaol}
gives. The sign of three zero bounds the documentation states, that of
\code{[-1, 0]*[-1, 0]} (SSE2 and NEON intervals only), \code{sqr([-1, 2])}
and \code{[0, 1]/[-2, -1]}, has to be the one of the build: it tells a tie
of the maxima of the products, which no other check did, the others
comparing with \code{==}, which does not see it.

\medskip
The powers are checked further. $x^n$ has to be the tightest over 4\,000
intervals near the ends of the doubles, where the products were once rounded
outward and the SSE2 and the FPU intervals gave other bounds for half of
them; 8\,679 powers that are doubles have to be exact, both bounds (an error
bound computed beforehand made \code{pow([2, 3], 3)}
$\itv{8-2^{-50}}{27+2^{-48}}$). $x^2$ has to be \code{sqr(x)}, bit for bit,
the sign of a zero bound included, in each of nine forms of the integer power
with the exponent 2, over the grid of special bounds and 90\,000 random
intervals, each with the calling code in the four rounding directions and,
for the first 3\,000 of each kind, with each flush mode: 4 of the 9 forms
fail without one call of \code{sqr()}, 5 without another. $(-c)^n$ has to be
$c^n$ or $-c^n$ at 1\,088 points near 1, 1.5 and 3 and 12\,736 more. The
powers for $n$ from $2^{24}+1$ to $2^{32}-1$ have to be the tightest against
mpmath: the lower bound was 1962 doubles below the tightest for
\code{pow([1.0000001], 2\textasciicircum{}32 - 1)}; the same for 46 negative $n$ from
$-2^{31}$ to $-2^{24}$, 44 of which fail without the second-order term of
the correction. The $n$-th roots of $x$ below $2^{-1020}$, subnormal
included, have to be within 2 doubles of the tightest bounds (4 for an
inverse root) for $n$ = 4, 5, 33, 65, 129 and 1025 against exact arithmetic
and up to $2^{32}-1$ against mpmath: \code{nth\_root([2\textasciicircum{}-1074], 33)}
was 1983 doubles below. \code{fma(x, y, z)} has to be the tightest over
random boxes, against \code{std::fma} called downward and upward directly:
GCC had folded the negations GAOL rounds its lower bound with.
\code{cancel\_minus} and \code{cancel\_plus} have to be the tightest, and
\code{y + cancel\_minus(x, y)} to enclose $x$, over intervals of the same
width and one double wider, and at the corners where the differences of the
bounds are within a double of each other or beyond the largest double: GCC at
\option{-O3} computed the terms of their TwoSum after the direction was set
back to upward, where it is no longer exact.\newinvfive

\medskip
\code{elementary} checks \code{exp}, \code{log}, \code{sin}, \code{cos},
\code{tan}, \code{asin}, \code{acos}, \code{atan}, \code{sinh}, \code{cosh},
\code{tanh}, \code{asinh}, \code{acosh}, \code{atanh}, \code{sqrt} and
\code{pow} at doubles, at intervals, and at intervals whose images are known
exactly: extrema, poles, domains, \code{log} of an interval holding no
positive number and \code{atanh} of one with no point of $(-1, 1)$ being
empty, \code{exp(0)}, \code{log(1)} and $1^y$ exact. The functions have to
be the tightest where their value is 0, 1, $\pm\pi/4$, $\pm\pi/2$ or $\pi$,
and \code{cosh}, \code{sinh} and \code{tanh} beyond the largest double and
near 1. \code{pow(x, y)} has to be within one double of the tightest bounds
at points, \code{pow([2], [1023.5])} among them, and over 385 boxes of bases
and exponents in every position about the base 1 and the exponent 0, with
the power itself as lower bound where it is a double: \code{pow([4], 0.5)} is
\code{[2]}. \code{atan2} has to be within one double over 324 boxes about
the axes, the tightest over the boxes with infinite bounds or on an axis,
$\itv{-\pi}{\pi}$ across the half-line $y=0$, $x<0$, and empty at $(0, 0)$.
\code{sin}, \code{cos} and \code{tan} have to be within one double at every
magnitude, \code{sin([1e-10])} and \code{cos([2\textasciicircum{}60])} included, and over
649 intervals next to their extrema and poles; \code{tan([-M\_PI\_2, M\_PI\_2])}
has to be the tightest, $\pm1.63\cdot10^{16}$, its width being the double
below $\pi$ and holding no pole, and \code{tan} within one double over 158
intervals whose bounds are next to two consecutive poles.\newinvfive{}
The integer functions
\code{sign}, \code{trunc}, \code{round\_ties\_to\_even} and
\code{round\_ties\_to\_away} have to be exact whatever the rounding
direction, which \code{std::nearbyint} and \code{std::rint} would not be,
over the halfway values, the whole numbers, their neighbours and the
magnitudes beyond $2^{52}$. \code{erf} and \code{erfc} have to be the
tightest against mpmath at 409 doubles (tiny and subnormal arguments,
changes of binade, the arguments past which $\mathrm{erf}(x)$ is within
$2^{-53}$ of $\pm1$ and $\mathrm{erfc}(x)$ is subnormal, and hard cases of
the \file{erfc.c} of CORE-MATH) and over the intervals between
them.\newinvfive{} The values are in \file{elementary\_values.h}, which
\file{elementary\_values.py} generates.

\subsection{The rounding direction and the floating-point \texorpdfstring{\mbox{environment}}{environment}}

\code{rounding\_direction} calls about 110 operations of GAOL's interface
with the rounding direction upward, to nearest, downward and toward zero
(and on x86 with the x87 and the SSE directions differing): they have to
give the results they give rounding upward, and to leave the direction
upward, or as they found it with \option{GAOL\_PRESERVE\_ROUNDING}; among
them the roots of negative order, the integer powers of large exponents and
\code{modulo\_k\_pi()} next to multiples of $\pi$, which compute after one
check of the direction.\newinvfive{} Products and sums have to be the
tightest in a loop that changes the direction before each of them. The 30
forms of \code{+ - * / \%}, \code{sqr}, \code{sqrt} and \code{inverse}
with intervals, with a double and as compound assignments have to be the tightest enclosures on random operands, in a loop that changes the rounding direction before each of them; on ARM they are run $2^{16}$ times each, in the four directions in turn, right after the program wrote the control register itself.
\code{pow} with a subnormal
result, which CORE-MATH rounds in the direction \code{fegetround()} gives,
has to be the tightest in each state, the x87 unit to nearest and the SSE
instructions upward included: the \code{fegetround()} of glibc reads the x87
unit, and the upper bound was below the exact value.\newinvfive

\medskip
The control registers of the processor are checked apart from the library
(\pxref{sec:rounding-direction}). What GAOL writes to the control word of the
x87 unit and to MXCSR on a 32-bit x86 processor is checked on every platform,
with two variables standing for the registers, in the 64 states of the two
rounding fields and the two flush modes, with other bits that must not change
and exception flags raised meanwhile that must not be lost: only the registers
that differ are written, and each is read once; on such a processor, the
registers themselves are checked in the states it can run. The
SSE2 operations with \option{GAOL\_PRESERVE\_ROUNDING}, which keep MXCSR
alone, are checked in its 16 states: MXCSR is written once when the
direction is not upward or a mode is set, and never read or written
otherwise. On x86-64, \code{get\_rounding()}, \code{set\_rounding()} and
\code{set\_rounding\_and\_flush\_modes()} change neither the exception flags
nor any other bit, and the type through which the library reads and writes
the two registers has to read the words the test reads and write the words
it is given, with four sets of the other bits, so that a read that drops the
flags or a write that forces the masks is seen. On 64-bit ARM, FPCR is read
and written with \code{mrs} and \code{msr}, and its bits other than the
direction and the flush modes (the default-NaN mode, FZ16, the trap-enable
bits, and the AH bit of Armv8.7 where the processor retains it) have to be
the program's after each operation.\newinvfive{} With mingw-w64 on x86-64,
the control word of the x87 unit is read with \code{fnstcw} and is part of
the state compared, since the \code{fegetround()} of mingw-w64 13 and later
joins the fields of the two units and its \code{fesetround()} clears the six
flags of MXCSR. Every operation has to keep the exception flags the program
raised (the five of IEEE 754, and on x86 the six of the SSE control
register) and the masks, from each direction, and after
\code{gaol::cleanup()} and \code{gaol::restore\_rounding()}: with
\option{GAOL\_PRESERVE\_ROUNDING}, the SSE2 operations masked the exceptions
again and cleared the flags.\newinvfive

\medskip
The empty set and the special values must raise nothing
(\pxref{sec:fp-exceptions}). \code{is\_empty()} has to be true for six
empty sets and false for three intervals, and to raise no flag; with glibc,
each is also told empty in a child process that enabled the
invalid-operation exception, which must not die: the comparison of the NaN
bounds with \code{<=} did. The operations of the interface with an empty
operand on either side, about 220 calls in 66 checks, among them one for each
of the 48 calls that compared the NaN bounds and five choices a program makes
with \code{is\_empty()} that GCC for 32-bit ARM and for POWER9 made signaling,
have to give the result of the empty set, raise no flag and not die in the
child process: 45 of these 48 calls died with the SSE2 intervals, and all 48
with the FPU ones.\newinvfive{} So do 18 operations given a NaN double, the products of
intervals with bounds in $\{-\infty, -2, -0, 0, 3, +\infty\}$, and the relational
divisions of intervals with bounds in $\{-\infty, -2, -0, 0, 4, +\infty\}$ (the SSE2 intervals divided $+\infty$
by $+\infty$, or a bound by a zero one, in the half of a register they did
not keep), 13 powers with an exponent of extreme magnitude, 10 powers next to
the largest double, where CORE-MATH's \code{pow} computed $\infty-\infty$,
and 65 loops of 64 relations and 4 loops that count the pairs for which a
relation holds, which a compiler vectorizes or if-converts: GCC 9.4 and 13 at
\option{-O3} with \option{-mfma}, GCC 13 to 15 for 32-bit x86, GCC 14 for
POWER8 and GCC 11 to 15 for 64-bit ARM made some of them raise the
exception with the FPU intervals, and Clang 18 and 20 at
\option{-O3 -march=native} without \option{-frounding-math}, so that the
teeth of these loops depend on the processor the code is compiled for. The
midpoints of five intervals with a bound of $2^{1023}$ or more have to be
right and raise no overflow flag, the sum of the bounds having overflowed;
36 calls of functions that computed beyond the largest double on the way
(\code{sqrt()}, \code{nth\_root()}, \code{atan()}, \code{sin()},
\code{pow()}, the relational functions) are run on 20\,000 draws of operands
each, and a result without an infinite bound has to come with neither the
overflow nor the division-by-zero flag;\newinvfive{} 27 cases such as
\code{sqrt([1, DBL\_MAX])}, \code{tan([DBL\_MAX])} and
\code{pow([2\textasciicircum{}-1074], [1.5 2\textasciicircum{}1014])} have to give their tightest
bounds, computed apart, with neither flag, and the same bounds in a child
process that enabled both exceptions.

\medskip
With a mode that flushes the subnormal numbers set, and the direction upward
or not, products, sums, differences, quotients, squares and \code{exp} with
a subnormal operand or result have to be the tightest, computed apart with
exact rationals (and mpmath for \code{exp(-740)}), and the modes have to be
cleared after the operation, or restored with
\option{GAOL\_PRESERVE\_ROUNDING}: a program linked with \option{-Ofast} gets
them from \file{crtfastmath.o}, and \code{[1e-300] * [1e-20]} was
\code{[0, 0]}.\newinvfive{} Then each of the 96 operations that check the
direction, called with each mode set just before it and each direction, has
to give what it gives with the modes cleared and leave them cleared (with
GCC 13 at \option{-O3}, \code{asinpi} of $\itv{100\cdot2^{-1074}}{100\cdot2^{-1074}}$ was above the
exact value and \code{sqrt([-1e-310, 4])} the empty set), and each of the 51
functions that make no check has to leave the modes as it found them, its
doubles compared by their bits; \code{midpoint()}, \code{split()},
\code{chi()} and the writing of an interval, which make the check, have to
clear them: with the sources before that change, 10\,900 of these checks
fail.\newinvfive{} \code{bisect(ratio)} and \code{is\_bisectable()} have to
give, under each mode and direction, what the bits of the bounds give, told
by integers: 2280 of these checks fail with the SSE2 intervals before the
change, and 4008 with \option{GAOL\_PRESERVE\_ROUNDING}. The AVX-512 path is
taken whole where the machine building has the instructions: its bounds are
those of the SSE2 path bit for bit, but for the sign of a zero bound of
\code{x * d} where the zero bounds of $x$ differ in sign, and it leaves the
direction as the program set it (\pxref{sec:avx512}).\newinvfive{}
\code{gaol::cleanup()} has to set back the direction the first
\code{gaol::init()} found, although a static object set it upward before
\code{main()}, and \code{gaol::restore\_rounding()} to set it to nearest
again as many times as it is called.

\medskip
Three programs complete it. \code{fast\_math\_link} is linked with
\option{-ffast-math} (its link only, with GCC and Clang), which links
\file{crtfastmath.o} and the flush modes it sets, unless
\option{-mno-daz-ftz} keeps it out: where the build gives the link that
option, the modes have to be clear when \code{main()} starts; otherwise the
test sets them, as a plug-in built with \option{-Ofast} does, and products,
a difference and \code{exp} with subnormal operands have to be the tightest
and the modes cleared after them.\newinvfive{} \code{automatic\_cleanup}
(Linux only) leaves the rounding direction downward at the end of
\code{main()}, and GAOL's automatic cleanup has to set it back to nearest,
or to leave it downward with \option{GAOL\_PRESERVE\_ROUNDING}: the
direction is read by a function that a constructor of priority 101 registers
with \code{std::atexit()}, before GAOL initializes itself.
\code{static\_initialization} computes about 40 operations before
\code{main()}, in the initialization of a static object: the constants
(\code{pi()}, \code{one()}, \code{universe()}), products and quotients, the
elementary functions and the reader of strings have to give the intervals
they give in \code{main()}, bit for bit. The static library of the CMake
build is initialized after the program, except with MinGW-w64, and GAOL's
constants were computed by its dynamic initialization: 25 of the 45 checks
failed with the SSE2 intervals, 24 with the others, \code{pi()} giving
\code{[-0, 0]}.\newinvfive

\subsection{Reading and writing}

\code{numbers} checks that \code{textToInterval("0.1")} is the tightest
interval enclosing the number read, and the number itself when it is a
double; that the constants are the tightest enclosures of $\pi$, $2\pi$ and
$\pi/2$; that no \code{const char*} nor \code{nullptr} converts to an
interval, the constructor from one string being gone;\newinvfive{} that the
literals of IEEE 1788-2015 (\code{[entire]}, \code{[ ]}, \code{[1,]},
\code{3.56?1}, hexadecimal numbers) are read as the tightest intervals in
any case of letters, and the numbers in base 16, 8 and 2 (\code{0x10},
\code{0xff.8}, \code{0o17}, \code{0b1.1p3}) as the tightest intervals
enclosing what their digits write, computed from the digits by the test,
beyond $2^{64}$, around $2^{53}$ and beyond the doubles, while \code{0x},
\code{0o8} and \code{0b2} are refused (\pxref{sec:reading-intervals}).
Expressions GAOL cannot read or compute have to throw, and, built with
LeakSanitizer, the parser has to free the nodes of all of them. The
intervals written in decimal in the general, scientific and fixed formats
with 1 to 20 digits have to enclose the intervals, each bound less than one
unit of its last digit away, near the powers of ten too; read back, they
have to enclose the intervals written. The format of the agreeing digits
writes the shared digits only for finite, nonzero bounds of the same sign,
with as many digits before the point and the same exponent
(\code{-1.25\textasciitilde[67, 0]}), and the bounds otherwise. A point
interval is written \code{[a]} when the number is the point itself, and
\code{[l, r]} otherwise, neither bound being the point
(\code{[0.1, 0.1000000000000001]} for \code{interval(0.1)}); the largest
doubles and the subnormals, written with all their digits (309 digits, 1074
after the point), are read back as the point.\newinvfive{} In hexadecimal,
the bounds have to be written in the form of IEEE 1788-2015 (13.4.1) and
read back as the same set, which is the recovery requirement of 13.4.
\code{operator<<} has to leave the precision of the stream as it was,
\code{std::setw} to pad the whole interval, and the texts written under the
flags of the stream and a locale of its own have to be those that \code{operator<<} wrote when it still wrote into a
stream of its own, before the text was made in a character string, except that
the bounds format writes in the C locale and the width
format writes the radius without a sign under \code{showpos}: 33 cases of
four formats. The width and center formats have to write the
\code{midpoint()} and the \code{rad()} of the standard, the radius rounded
upward so that midpoint plus or minus it contains the interval, and never 0
for an interval that is not a point: GAOL wrote $(l+r)/2$ and $(r-l)/2$
rounded to nearest, which do not contain the interval.\newinvfive{} The
empty set has to be \code{[empty]} in the five formats.

\medskip
\code{while (in >> x)} has to stop at the end of the input with
\code{failbit} set and the interval unchanged, and a line that is no
interval has to set \code{failbit} and throw \code{std::invalid\_argument};
blank lines have to be skipped, under \code{std::noskipws} too, and the
stream tied to the input flushed before the first character is read, so that
a prompt shows before the user types.\newinvfive{} Numbers with a million
zeros after their point and an exponent of 7 digits have to be read exactly:
the exponent was cut at 100000. Under a locale writing a decimal comma,
numbers of 5000 to 20000 characters, in decimal, in hexadecimal and in the
uncertain form, have to be read as the tightest enclosures, known without
reading them, and the number of 20000 characters has to take at most 10
times the time it takes under the C locale, and 50~ms more: the reader read
the text again for each of its 125 comparisons or so, in a time quadratic
in its length, and took 50 times as long.\newinvfive{} The cache variable
\option{GAOL\_NUMBERS\_TIMEOUT} gives the test more than its 5 minutes on a
slower build (600~s with GCC and the sanitizers on macOS x86\_64). With
each flush mode the processor honours, the numbers from 0 to the least
normal double have to be read as the tightest intervals enclosing them, and
the doubles among them as themselves: the reader compared each number with
the doubles around it as doubles, and denormals-are-zero reads a subnormal
as 0, so that \code{1e-310} was read as the interval from the greatest
subnormal to the least normal double, which does not enclose
it.\newinvfive{} Under such a mode, \code{<a, b>} of two different
subnormals has to be refused, intervals of subnormals written with 16
digits have to be read back as enclosures once the mode is restored, and
the text has to be the one written without the mode, \code{operator<<}
clearing it before it writes.\newinvfive

\subsection[The functions, the standard and the reverse functions]{The functions, the standard and\\ the reverse functions}

\code{other\_functions} checks the midpoints (of subnormal bounds too), the
widths, the radii (that of the point zero, like its width, being $+0$ whatever
the signs of its bounds), the magnitudes and mignitudes, the Hausdorff distances
(of intervals with infinite bounds too), \code{nb\_fp\_numbers()} across
the two zeros, the splitting, the integer parts, and \code{max(x, y)} and
\code{min(x, y)} of every pair of intervals made of special bounds, with the
bits of the zeros and the subnormals, with the flush modes set or not; with
the SSE2 and NEON intervals, every pair of stored lanes made of special
doubles has to give the bits of the formula of the other kinds, and the
quotients of the NEON intervals with a NaN in a lane of an operand the bits
of \code{interval::emptyset()}. The tools of interval algorithms,
\code{width\_enclosure()}, \code{inflate(r)} and
\code{interval::centered\_form(m, r)}, have to be the tightest enclosures of the
exact width, of $\itv{a-r}{b+r}$ and of $\itv{m-r}{m+r}$, computed with exact
rationals; \code{bisect(ratio)} has to give two parts that cover the
interval exactly, cut strictly inside where \code{is\_bisectable()}, for
ratios from $10^{-9}$ to $1-10^{-9}$, and to raise no overflow exception for
finite bounds next to the largest double, \code{bisect(0.5)} being
\code{split()}, and to throw \code{std::domain\_error} for a ratio outside
$(0, 1)$ or NaN.\newinvfive{} The comparisons of IEEE 1788-2015 have to
give the values of its Tables~10.3 and 10.4, \code{==} and \code{!=} must
not compile on intervals, and each name of \code{gaol\_ieee1788} has to be
the operation of the standard it names, which a wrong translation would not
show at compilation: the eight comparisons over 20\,000 pairs, \code{inf},
\code{sup} and the numeric functions against Table~10.2, the reverse
functions with the arguments in the order of the standard,
\code{mulRev(b, c, x)} being \code{div\_rel(c, b, x)}, and \code{pow}
against the \ieee{pow} of Table~9.1, which GAOL's own \code{pow} is not for
a negative base (\pxref{sec:pow-namespaces}). The reverse functions of
\code{max}, \code{min}, \code{sign} and \code{floor}\newinvfive{} have to be
the hulls of their sets, read from the definitions of the functions apart
from the code, on every interval with bounds among $-\infty$, $-2$, $-1.5$,
\dots, 2, $+\infty$ and the two zeros, on 66 cases written out, and on
20\,000 random intervals of every magnitude, where every sampled double the
definition puts in the set has to be in the result.

\medskip
\code{ieee1788} uses \code{gaol\_ieee1788} as a program does, under
\code{using namespace gaol\_ieee1788;} alone: every name of the standard it
provides is called unqualified, which compiles only if none is ambiguous
with a function of \code{gaol\_core} (Chapter~\ref{chap:ieee1788}).
\code{pow} has to be the standard's with an interval, an \code{int} or a
\code{double} exponent, and \code{pown} and \code{gaol::pow} the integer
power for an exponent of any integer type:\newinvfive{}
\code{pown([2], 3000000000u)} is the overflowing power, which the
conversion of the exponent to a negative \code{int} made a tiny one, and an
exponent beyond the long longs or of 128 bits gives the tightest bounds of
mpmath where the power is a double and within three doubles of them
otherwise. \code{rootn()}, \code{gaol::nth\_root()} and \code{pownRev()}
take an order of any integer type, and \code{pownRev(c, x, p)} for $p=0$ and
$p<0$ has to give the exact roots of the reciprocals of $c$ where they are
doubles, with no solution for the part of $x$ in the gap of $c$.
\code{overlap(a, b)}\newinvfive{} has to be the state the set specification
of Table~10.7 gives, whose quantifiers the test reads over a finite model of
the real line, on 31\,449 pairs: exactly one specification holds for two
nonempty intervals, but for \ieee{meets} and \ieee{metBy} with a point
interval on a bound of the other one, which the standard gives to
\ieee{starts}, \ieee{finishes}, \ieee{finishedBy} or \ieee{startedBy};
\code{disjoint}, \code{equal} and \code{subset} have to be the masks of
states 0xF001, 0x8040 and 0xC3C0 of the standard. The bounds of
\code{gaol\_ieee1788::pow} and \code{gaol::pow} on 117 boxes, each reaching
a branch of the \ieee{pow} of Table~9.1, have to be the bounds the table
gives, checked against the exact power computed with 500 bits by mpmath
(\file{tests/tools/pow/checkrows.py}). \code{intervalToText(x)} has to be a
literal of the standard with the digits of the precision of the intervals,
whatever the global format and the locale, which a grammar of the literals
of Tables~9.5 and 12.2 checks, and \code{textToInterval} has to read it back
as an interval containing $x$: under a comma locale the literal
\code{[a]} wrote \code{[-2,5]}, read as $\itv{-2}{5}$.\newinvfive{}
\code{mulRevToPair(b, c)} has to be the pair of 10.5.5, the part below
first, tight on cases whose parts have exact bounds, and every $x=k/4$ of
the set on every $b$ and $c$ with integer bounds in $[-3, 3]$ has to be in
one of the two parts.\newinvfive

\medskip
\code{reverse} checks the relational functions against the reverse functions
of IEEE 1788-2015 (10.5.4, Table~10.1): \code{sqrt\_rel} (\ieee{sqrRev}),
\code{invabs\_rel} (\ieee{absRev}), \code{nth\_root\_rel} (\ieee{pownRev},
for every integer $p$),\newinvfive{} \code{asin\_rel}, \code{acos\_rel},
\code{atan\_rel}, \code{acosh\_rel}, and \code{div\_rel}, \code{\%} and
\code{\%=} (\ieee{mulRev}). Over the cases of the minimal tests of
libieeep1788 and random ones, each result has to be empty when an argument
is, \emph{valid}, enclosing the hull of the $x$ of $I$ whose image is in
$J$, which the test holds the tightest enclosure of, and \emph{accurate},
within $\ieee{nextOut}(\ieee{tightest}(\ieee{nextOut}(\mathit{arguments})))$
(12.10.1), as the standard recommends for inf-sup types. The bounds of
\code{asin\_rel}, \code{acos\_rel} and \code{atan\_rel} are allowed one
double further, their reduction modulo $\pi$ rounding twice, and those of
\code{nth\_root\_rel} for a negative $p$ two: of the 18\,052 random cases
with $p<0$, 12\,495 have a nonempty result with nonempty accurate bounds,
521 were one double beyond them, 25 two, none more
(\pxref{sec:relational-arithmetic}).
The hulls of \file{reverse\_values.h} are computed from the definition of
each function with exact rational arithmetic, and mpmath for the periodic
ones, not from GAOL nor from the results libieeep1788 expects, which the
generator checks are enclosures of them. \code{atan2\_rel\_y} and
\code{atan2\_rel\_x} (\ieee{atan2Rev1}, \ieee{atan2Rev2}),\newinvfive{}
which libieeep1788 does not test, have a table of their own, generated with
mpmath at 3000 bits: 1\,969 cases, 1\,263 of them with points, whose hulls
are computed from the vertices of the polygon of each convex cone in the
box; each result has to be valid, accurate and within 2 doubles of the
tightest bounds, and about 600\,000 sampled points whose angle is in $c$
have to be in both results. \code{pow\_rel} and \code{pow\_exponent\_rel}
(\ieee{powRev1}, \ieee{powRev2})\newinvfive{} are checked on the 804 cases
of \file{pow\_rev.itl} of ITF1788, random ones and ones whose solutions are
doubles, the hulls computed with mpmath at 2000 bits by cells: the generator
found two results of ITF1788 wider than the tightest enclosure, and by
sampling over 6\,000 random boxes, 54\,305 points whose power is within $c$
have to be kept by both functions.

\medskip
\code{test\_ibex\_codac} checks the cases of the unit tests that IBEX and
Codac run on their intervals, which are GAOL's, against the exact results:
3905 checks.\newinvfive{} A case is a row of a table, written with
\code{gaol::interval} and the functions of GAOL, named by its origin
(\code{IBEX TestArith::add05}, \code{Codac Interval\_operations:176}); the
exact result is computed apart by the scripts of
\file{tests/tools/ibex\_codac} with exact rationals and mpmath at 2000 bits,
and the bounds of GAOL have to enclose it with no double between them and
the tightest bounds (two doubles for \code{atan2\_rel\_x} and
\code{atan2\_rel\_y}). Three sections: the forward tests of IBEX (71 of the
139 tests of \file{TestInterval} and 182 of the 343 of \file{TestArith},
and 3 that IBEX never runs; 1393 checks), its backward contractors, its
generalized division and its distances (197 cases, 1306 checks), and the
tests of Codac (1206 checks; of the 945 checks of its three files, 195 are
already in the tests of GAOL). What only IBEX or Codac have, \code{tan06} of
IBEX, whose expectation is wrong, and the distance to the empty set, whose
convention differs, are left out. What the projects expect is compared with
the exact results by the scripts: 1940 comparisons, no disagreement beyond
their tolerances, and 23 expectations of IBEX with an integer bound that is
not a bound of the exact result (\code{tan02} expects $\itv{-1}{1}$ where it
is $\pm(1+2^{-52})$).

\subsection{CORE-MATH, compoundm1 and the 128-bit integer}

\code{core\_math} checks the bounds of the elementary functions against
CORE-MATH itself, which is correctly rounded in the rounding direction in
effect: the tightest bounds of $f$ at a double $x$ are the values it gives
rounding downward and upward, which the test computes by setting the
direction itself. GAOL's bounds have to enclose them, to be equal above and
within one double below, and to be the tightest where the code gives the
exact value itself. Each function is tried at the ends of its domain, at the
values GAOL treats apart, at the powers of two and their neighbours, at the
subnormals and at random doubles. \code{pow(x, y)} has to be the tightest
at the corner of a box, both bounds, over 84\,070 pairs; where $x^y$ is
rational (8\,981 pairs, 3\,211 of them a double) the reference is $x^y$
itself, computed with integers apart from GAOL and from CORE-MATH, whose
values downward and upward have to be its roundings too, which tells a
fault of CORE-MATH from one of GAOL: CORE-MATH's \code{pow} rounded to
nearest the midpoints of its exact phase on 32-bit ARM, and its square roots
on a 32-bit x86 Windows. \code{exp2}, \code{exp10}, \code{log2} and
\code{log10} have to be the tightest, and the exact values where they are
doubles (\code{exp10} of 0 to 22, \code{log10} of a power of ten up to
$10^{22}$). The functions of Table~10.5 but \code{compoundm1}
have to be the tightest over intervals too, the hull of their image computed
from the values at the bounds of the part of the interval in the domain, at
the points nearest to and farthest from the origin (\code{hypot}), and at
every multiple of $1/2$ within for \code{sinpi}, \code{cospi} and
\code{tanpi}.\newinvfive{} \code{sin} is checked, in the four directions,
against the upstream sources of CORE-MATH compiled unchanged at 19 arguments
of its accurate path, which computes with the 128-bit integer GAOL ports,
and \code{cos} at 45 arguments of its accurate path against its roundings
computed with mpmath at 4\,000 bits, which reach the branches of that path.
First of all, the test checks the \code{fma()} and \code{round()} of the C
library, which CORE-MATH calls: \code{fma()} correctly rounded in the four
directions at six triples, and \code{round()} at
$\pm$\code{0x1.fffffffffffffp-2}, $\pm0.5$, $\pm2.5$ and $2^{52}-1/2$;
the \code{fma()} and \code{round()} of
mingw-w64's own math library on x64, which GAOL refuses, fail 21 and 6 of
these checks under wine.\newinvfive{} The \code{gaol\_ldexp()} CORE-MATH
rounds the subnormal values of \code{erfc} with has to be $x\cdot2^n$
rounded once in the four directions at 20\,000 random pairs. \code{pow} at
the corners where the base is 0 or $+\infty$ or the exponent infinite has to
raise neither the division-by-zero nor the invalid-operation exception, and \code{pow} of a power next to the largest double no
invalid-operation exception either: for $y\log x$ from
709.78263 to 709.78267, CORE-MATH's first phase scaled its approximation by
$2^{1024}$ and computed $\infty-\infty$, at 942 of 1\,001 exponents of the
base 2 from $1024-2^{-13}$ in steps of $2^{-24}$.\newinvfive{} The reading
of a number above the largest double (\code{1e400}, \code{0x1p1024}) has to
be $\itv{\mathit{DBL\_MAX}}{+\infty}$ and raise no overflow flag, which
\code{strtod()} raised, and a number below the least normal double no
underflow flag.\newinvfive{} CORE-MATH, called without GAOL in each of the
six directions of x86, has to round in the direction of MXCSR; on a 32-bit
x86 processor its \code{fegetround()} was the C library's, reading the x87
unit (glibc) or giving $-1$ (Visual C++) when the two differ. \code{erf}
and \code{erfc} have to be the tightest over 80\,000 intervals, and at 22
hard-to-round arguments of the \file{erfc.c} of CORE-MATH.

\medskip
\code{compoundm1} checks $(1+x)^y-1$ of IEEE 1788-2015 (Table~10.5), which
CORE-MATH has not and whose bounds GAOL computes with 64-bit fixed point in
a fast phase, and where that does not settle them with a 128-bit
significand and a bound of the error.\newinvfive{} The test goes from the
bottom up: the count of leading zeros, the 256-bit product against the
schoolbook one, the integer square root, the doubles around an integer; the
sums and products of 20\,000 triples of doubles, whose balls have to hold
the exact value; the balls of $\log(1+x)$ for 493 arguments and of $e^t-1$
for 322, which have to hold the value computed to 256 bits by mpmath and be
no wider than $2^{-118}$ relatively; the fast phase by itself, whose center
has to be within $3/10$ of the bound of its error over 20\,000 pairs from
eight domains (the worst error found is 0.22 of it, in 27 million pairs),
and which answers 74 to 96\,\% of the pairs according to the domain; the
exact division of $y=-1$ at the 633 points of 719 where $1+x$ takes 127 bits
or less. Then the bounds at 5\,854 pairs have to be the greatest double at
most the value and the least double at least it, the values computed by
mpmath with 1\,500 to 7\,000 bits, and on 850 boxes the result has to be
the hull of the values at the four corners and of 0 where the box crosses
$x=0$ or $y=0$, with the limits at $x=-1$ and at infinity, read in the
script and not in GAOL.

\medskip
\code{u128} compiles the two 64-bit halves of \file{gaol/math/gaol\_u128.h},
with which GAOL computes the accurate phases of CORE-MATH where the compiler
has no 128-bit integer (Visual C++, and GCC on 32-bit targets), and compares
every operation with the native type at the ends of the ranges, at the
powers of two, at random values and over every shift count from 0 to 127:
the two have to give the same 128 bits. Where the compiler has no native
type, the test checks the identities the operations satisfy instead.
\code{extended\_precision} checks the results that doubles computed on the
x87 unit round wrongly, which is why \file{gaol/gaol\_config.h} refuses
that unit (\pxref{sec:refused}): \code{exp(-0x1.74910d52d3051p+9)} gave 0
rather than $2^{-1074}$ and \code{tanh(0x1.30fc1931f09c9p+4)} 1 rather
than $1-2^{-53}$, their results being rounded twice. The values are computed
by \file{extended\_precision\_values.py} with mpmath at 5000 bits. Built on
the x87 unit, the test fails on Debian 12 i386 with GCC 12 (22 checks) and
on x86-64 with \option{-mfpmath=387} and GCC 9 (76 of its 300 checks).

\subsection{The expressions}

\code{expressions} goes through every node of the tree of
\file{gaol/parser/gaol\_expression.h} and every way a string can be wrong:
the numbers in every form the lexer takes, the operators and the functions
alone and nested, and the strings the parser has to refuse with an exception
rather than an interval. The intervals have to be read wherever a number may
stand (\code{1+[1,2]}, \code{[cos([0,1]), 2]}), which the two grammars of
GAOL~4 refused, an error anywhere has to stop the reading, and the names
GAOL v5 adds (\code{exp2}, \code{log2}, \code{cbrt}, \code{sign},
\code{trunc}, and every function of GAOL) have to be read in any case of
letters, the lexer taking the longest name.\newinvfive{} Each value is
compared with the same computation written in C++, which the other tests
check against the exact results: what is tested is the lexer, the parser
and the evaluation, not the operations. What GAOL throws has to be exactly
\code{std::invalid\_argument} for what the reader refuses and
\code{std::domain\_error} for a function called with an argument it does
not take, told apart by \code{typeid}, with a \code{what()} that is
\code{file:line: explanation}, the file being the name of the source that
threw with no directory (\pxref{sec:exceptions}).\newinvfive{} Long
expressions have to be evaluated and deleted: chains read from a string of
200000 sums or products and of 20000 differences, and chains of 200000
minus signs and 20000 sines, whose evaluation and deletion made one call
within another per node and overflowed the stack of 8~MB; a leaf at the
bottom of a chain of 20000 nodes has to be visited and deleted less than
64~KiB of stack away from the function that starts the evaluation, where the
recursion took 1.9~MB and 0.96~MB (GCC 15.2, x86-64, Release). A string
nested 9000 levels has to be read, and one nested 100000 levels refused, and
a sum of 10000 terms read with no more than 64 allocations live at once,
where the parser built 20000 nodes before the evaluation.\newinvfive{} A
string holding a NUL character has to be no expression, a null pointer has
to throw \code{std::invalid\_argument} with the names of GAOL and of the
standard, and to make \code{parse\_interval()} return \code{false} without
exceptions, where it died of a segmentation fault. \code{[l, u]} with $l$
above $u$ has to be the empty set, two numbers between the same doubles
being compared exactly, as digits, and the bounds of \code{<a, b>} have to
be doubles: \code{<0.1, 0x1.9999999999999p-4>} is refused.\newinvfive{}
Every allocation of a reading has to be able to fail, each in turn: the
reading has to throw \code{std::bad\_alloc} and free all it allocated,
which the test counts with a replacement of the global \code{operator new},
where flex ended the program with \code{exit(2)}. The checks of several
threads reading at once are commented out, the tests running no thread.

\medskip
\code{debugging} compiles GAOL's headers with \code{GAOL\_DEBUGGING}, which
the Debug builds define, and builds and evaluates an expression whose nodes
call \code{GAOL\_DEBUG\_VERBOSE}, which has to run its command at the level
given to \code{gaol::init()} and not above (\pxref{sec:debugging-facilities}):
\file{gaol/parser/\allowbreak gaol\_expression.h} did not compile with
\code{GAOL\_DEBUGGING}, writing on \code{std::cout} without the
header.\newinvfive

\subsection{The builds and the compilers}

\code{version\_file} (a script of CMake) checks that the reading of
\file{VERSION.txt} ignores a UTF-8 byte order mark, the line ends of
Windows and the blanks around the version, refuses anything but three
numbers, and names a UTF-16 file and a NUL byte in its message, whether
\code{file(READ)} cuts the text at the NUL byte (CMake 3.14.7 and 3.16.3)
or keeps it (3.28.3, 4.4.3); \file{.github/scripts/version-file.sh} runs the
same cases on \cmd{configure}, meson and autoconf in the continuous
integration.\newinvfive{} The compile tests \code{refused\_*} compile a
program including \code{<gaol/gaol>} with \option{-ffinite-math-only} and
with \option{-ffast-math}, put after the flags of interval arithmetic, and
the compilation has to fail with the message of \file{gaol/gaol\_config.h},
which the test looks for in the output of the compiler; compiled with
neither, it has to compile (\pxref{sec:refused}). With
\option{-ffinite-math-only} the empty interval is no longer told empty:
\code{([1, 2] \& [3, 4]).is\_empty()} was false with GCC 9.4 and 13, and
with Clang 18. The CMake build copies and touches the source before each
compilation that has to fail (\file{tests/compile\_error.cmake}), since an
object compiled once stayed up to date. \code{refused\_aligned\_new} and \code{refused\_aligned\_new\_cxx17} compile
\file{gaol/gaol\_config.h} alone with \option{-D\_WIN32 -U\_WIN64}, which stand for
32-bit Windows on any system, without the \option{-faligned-new=8} the flags
carry there, in C++11 and in C++17, and have to fail;
\code{refused\_aligned\_new\_positive} compiles it with the option, and has to
succeed: the C runtime of 32-bit
Windows aligns the memory of \code{new} on 8 bytes, an SSE2 interval needs
16, and a \code{std::vector} of SSE2 intervals stopped the program on its
first \code{movaps}.\newinvfive{} \code{cpack\_stale\_configure} configures
copies of the tree and checks that CMake warns when the committed
\file{configure} says another version than \file{VERSION.txt}, naming the
two versions, and that CPack configures the copy again and names the archive
after the new version. \code{automatic\_configuration} configures GAOL
several times in one directory, with the defaults, \option{GAOL\_SIMD=OFF},
the AVX-512 path refused and asked for, the Debug build type, the settings
of the developers and a cross-compilation, and holds
\file{gaol\_configuration.h} and the messages of CMake to what each asks for
and to what the checks found; it builds nothing and takes about ten seconds.
Only the CMake build has these two scripts: \cmd{configure} and meson are
checked by the audit of the three builds. The six tests \code{nodiscard\_*}
compile a program in each of the three standards with
\option{-Werror=unused-result} and \option{-Werror=attributes}
(\option{/we4834}, \option{/we5030} and \option{/we5051} with Visual C++):
throwing away the result of \code{sqrt(x)} has to fail, and using it has to
compile, since \code{GAOL\_NODISCARD} is \code{[[nodiscard]]} before C++17
with GCC 7 and later and Visual C++ 2019 16.4 and later, and the attribute of
Clang with Clang.\newinvfive

\subsection{The unit tests of GAOL~4}

\code{arithmetic\_operators}, \code{assignment}, \code{constants},
\code{constructor}, \code{float\_functions}, \code{generic\_functions},
\code{input\_output}, \code{interval\_functions}, \code{misc},
\code{non\_arithmetic}, \code{relations}, \code{reverse\_mappings} and
\code{trigonometric} check the results of the operations on chosen
intervals, as Fr\'ed\'eric Goualard wrote them in \file{check/}.
\code{reverse\_mappings} draws its random intervals from a seed of its own
rather than from the process identifier, so that its checks are the same at
each run, and checks the infinite ends of \code{atanh\_rel(J, I)} where $J$
ends at $\pm1$.\newinvfive{} \code{input\_output} checks that the reader
refuses \code{<-inf, -inf>} and \code{<inf, inf>}, whose bounds are not one
value. \code{constructor}, \code{assignment}, \code{arithmetic\_operators}
and \code{relations} also check the integers: \code{interval(n)} and
\code{interval(a, b)} around $2^{53}$, $2^{63}$ and $2^{64}$, for integers of
every width and in the four rounding directions, each operator in both
orders, and \code{set\_contains()} at $2^{53}+1$ and $2^{64}-1$, which the
conversion to a double made wrong.\newinvfive{} \code{relations} checks
\code{overlap()} on each of the 16 states of Table~10.7 and
\code{gaol::lexicographic\_less}, a strict total order that sorts a
\code{std::vector} holding an empty interval, which \code{<} reads past the
end of (\pxref{sec:lexicographic-less}).\newinvfive{} \code{misc} checks the
names of \code{gaol}, which names those of \code{gaol\_core} one by one,
with \code{static\_assert}: the names of GAOL~4 that Codac and IBEX use
compile, and what GAOL uses for itself (\code{get\_rounding()},
\code{gaol\_pown()}, \code{rounding\_state}, \code{the\_null\_expr}) is not
found through \code{using namespace gaol} (\pxref{sec:namespaces}); it
checks that the kernels the operators call are in \code{gaol\_detail} and
not in \code{gaol\_core}, where they made a program ambiguous
(\code{test\_kernels\_are\_\allowbreak not\_in\_gaol\_core}); that the heap
holds a type
as aligned as an interval in \code{std::allocator}, \code{new},
\code{std::vector}, \code{std::list}, \code{std::deque}, \code{std::map} and
\code{std::make\_shared} (\code{test\_heap\_alignment}, which on Linux with a
\code{::operator new} aligning on 8 bytes fails without
\option{-faligned-new=8} with GCC 15 and Clang 21 and passes with it); that
intervals at addresses of 8 modulo 16 compute where the SSE2 intervals keep
their bounds as two doubles (\code{test\_intervals\_\allowbreak
aligned\_on\_8\_bytes});
that a function working on local intervals, entered through a trampoline in
assembly with the stack at each of the four alignments of 4 bytes, returns
with the right bounds on 32-bit Windows with GCC or Clang, where the system
calls its threads with the stack aligned on 4 bytes
(\code{test\_stack\_entered\_\allowbreak aligned\_on\_4\_bytes}, which GCC 14 at
\option{-O2} fails on Linux for 32-bit x86 without \option{-mstackrealign});
and which builds with \option{GAOL\_PRESERVE\_ROUNDING} write MXCSR alone
(\code{test\_rounding\_macros\_\allowbreak write\_mxcsr\_alone}).\newinvfive{}
\code{intervalf} and \code{interval2f} test the intervals of floats where a
developer compiles them (\code{GAOL\_FLOAT\_INTERVALS}), and are skipped
otherwise (exit status 77). \file{check/fpu.cpp}, an empty test,
\file{check/essai.cpp} and \file{check/performances.cpp}, the benchmark of
GAOL~4, are gone, \cmd{make perf} replacing the last.\newinvfive

\subsection{The benchmark}
\findex{gaol\_performance}%

\code{gaol\_performance} (\file{tests/performance.cpp}) measures the time per
operation of GAOL's arithmetic and elementary functions, of the quotient by a
positive double, of the constructor, \code{floor()}, \code{x \&= y},
\code{x | y}, \code{max()}, \code{min()}, \code{x \% y},
\code{x.inverse()}, \code{midpoint()}, \code{nth\_root(x, -2)},
\code{nth\_root(x, 3)} and \code{asinpi()}, of two relations, of the same
operations on doubles, and, in its last rows, of lines of 15 additions, 15
additions and subtractions, 7 products, 7 quotients and 13 operations of the
four kinds between 8 and 16 intervals.\newinvfive{} Its title names the
intervals GAOL computes with, NEON, SSE2 or FPU, and four rows give the
product, the quotient, the square and \code{pow(x, 2)} of operands of random
signs, the other rows having positive operands only, which a processor
predicts: the operands are 1\,024 intervals that the predictor learns by
heart, and the product of random signs of the FPU intervals before their
kernels, which chose a branch by the signs as GAOL~4 did, took 6.4~ns on them
and 12.4 on 16\,384 (Apple M2). The table is Markdown, which the jobs of the
continuous integration copy into their summary. After the table, it prints
the time of the check each operation makes, $1+(2^{-1060}+0)$ not above 1,
and of the same check made of $2^{-60}$: a processor that takes a microcode
assist for a subnormal operand shows it there; with the NEON intervals, the
check is the read of FPCR. It is not a test (\pxref{sec:ci-speed}). The
tools of \file{tests/tools/pow/} and \file{tests/tools/arith/}, which
compare the bounds of two builds of GAOL and build GAOL with deliberate
errors, are not tests either (\pxref{sec:scripts}).

\section{The examples}\label{sec:examples}
\cindex{examples}%

The directory \file{examples/} holds 15 programs that use GAOL v5 the way
interval arithmetic is used in practice, as IBEX, Codac, C-XSC, INTLAB and
the textbooks use it, and the example of GAOL~4, 16 programs in all.
\file{examples/CMakeLists.txt}
is a project of its own, which uses an installed GAOL as any user's project
would:

\begin{cmdshell}
% cmake -S examples -B build-examples -DCMAKE_PREFIX_PATH=@var^prefix~
% cmake --build build-examples
% ctest --test-dir build-examples
\end{cmdshell}

\noindent
The three builds of GAOL build them too, with the GAOL they build, where
they are asked to (\option{WITH\_EXAMPLES}, \loption{with-examples},
\code{with-examples}), and \cmd{make check} (\cmd{ninja check} with meson)
runs them with the unit tests (\pxref{sec:tests}).

\medskip
Examples 01 to 15 check what they print:\newinvfive{} each enclosure is
compared with a value computed apart (mpmath, closed forms, exact
fractions), each proof of uniqueness or of emptiness is tested, and each
claim of the output (``4 times too wide'', ``the excess is divided by 100'')
is proved with interval arithmetic. A failed check prints
\code{FAILED: <claim>} and the program returns a failure, so that ctest runs
the examples as tests. The checks test containments and qualitative facts,
never digits: they pass on every build. Example 16 checks its output too,
with the digits of an enclosure that no operation rounds. All 16 examples
compile without a warning (\option{-Wall -Wextra}, GCC 9.4 and Clang 18)
and pass their checks with the four builds tried: the SSE2 intervals, the
FPU intervals (\option{GAOL\_SIMD=OFF}), \option{GAOL\_PRESERVE\_ROUNDING=ON},
and the library compiled by Clang 18. Each runs in less than 0.3~s with the
default build (06 takes 1.3~s with \option{GAOL\_PRESERVE\_ROUNDING}). Their
outputs are the same on every build, apart from the sign of a zero bound
(\code{[-0, 4]} with SSE2, \code{[0, 4]} with the FPU code), from the report
of the rounding policy in 13, and from a digit or two where a program
computes its own doubles (the preconditioner of 10). Every program follows
the rules of GAOL v5: \code{\#include <gaol/gaol.h>}, no call to
\code{gaol::init()}, and \code{gaol::cleanup()} right after the last use of
GAOL. Three headers are shared: \file{box.h}, a vector of intervals with
the names IBEX and Codac use for boxes (\code{widest()},
\code{max\_width()}, \code{mid()}, \code{volume()},
\code{bisect(i, ratio)}, \code{\&}, \code{|}), since GAOL has no box type
and every algorithm on several variables needs one; \file{dual.h}, a
forward-mode automatic differentiation over intervals, with which a function
written once as a template gives its natural extension with \code{interval}
and its gradient enclosure with \code{Dual<N>}; and \file{affine.h}, the
affine forms of example 12. Table~\ref{tab:examples} lists the programs.

\begingroup\renewcommand{\arraystretch}{1.1}\small
\begin{longtable}{L{4.3cm}L{9.5cm}}
  \caption{The examples of \file{examples/}}\label{tab:examples}\\
  \hline
  \multicolumn{1}{c}{Program} & \multicolumn{1}{c}{What it computes and shows}\\
  \hline\hline
  \endfirsthead
  \multicolumn{2}{l}{\textit{Table~\thetable, continued}}\\*
  \hline
  \multicolumn{1}{c}{Program} & \multicolumn{1}{c}{What it computes and shows}\\
  \hline\hline
  \endhead
  \hline
  \endfoot
  \code{01\_first\_steps} & Intervals built every way, their numbers, sets
    and relations: containment, \code{interval(0.1)} against
    \code{textToInterval("0.1")}, \code{mid()} against \code{midpoint()},
    no \code{==}, the certainly relations\\
  \code{02\_decimals\_and\_rump} & 0.1 summed ten times, $\pi$ by Machin,
    Rump's polynomial: decimals are not doubles, \code{sqrt(2.0)} is the C
    library's, the width as an alarm\\
  \code{03\_dependency\_problem} & $x-x$, $x(1-x)$ three ways, $(x-1)^5$,
    subdivision, Goldstein-Price, rotations: the dependency problem,
    single-use forms, \code{sqr} against \code{x*x}, the linear convergence
    of subdivision, the wrapping effect\\
  \code{04\_centered\_form} & $x\cos x$ and Chebyshev $T_5$, natural and
    mean-value forms: automatic differentiation over intervals, the
    quadratic convergence of the centered form, the monotonicity test\\
  \code{05\_interval\_newton} & A root of $\sqrt{x}+(x+1)\cos x$, the 7 roots
    of $\sin x - x/10$, $\pm\sqrt2$: interval Newton, uniqueness proofs,
    the extended division, the trap of $f(m)/F'(X)$ when $0\in F'(X)$\\
  \code{06\_global\_optimization} & The global minima of Goldstein-Price and
    of the six-hump camel: Moore-Skelboe branch and bound, a priority queue
    with an explicit comparator, the cluster effect, mean-value form and
    monotonicity test\\
  \code{07\_contractors} & IBEX's backward arithmetic, $x^2+y^2=1$ with
    $y=e^x-1.5$: forward-backward contraction (HC4Revise) with GAOL's
    reverse functions, their names in GAOL, IEEE 1788 and IBEX\\
  \code{08\_sivia} & Codac's SIVIA example, IBEX lab2 and lab3,
    $\sqrt x+\sqrt y\in\itv{0}{1}$: set inversion, guaranteed area brackets,
    contractors in pavers, the domain trap of a library without
    decorations\\
  \code{09\_parameter\_estimation} & IBEX lab5, range-only localization from
    three beacons: bounded-error estimation by set inversion, measurements
    read as text, a contractor with \code{sqrt\_rel}\\
  \code{10\_krawczyk} & Two circles of IBEX's solver, the roots of
    $\lambda+e^{-\lambda}$: a branch-and-prune solver with the Krawczyk
    operator, why IBEX bisects at 45\,\%, a false proof from a wrong
    Jacobian\\
  \code{11\_ode\_enclosure} & $x'=-x^2$, Codac's Lohner system,
    $\int\exp(-x^2)$: Picard's a priori enclosure, Taylor steps with a
    remainder, the mean-value form against the wrapping of $x+hf(x)$,
    interval Riemann sums\\
  \code{12\_affine\_arithmetic} & $x-x$, $x(1-x)$, $(x-1)^5$, $(x-y)^2$,
    rotations with affine forms: a sound affine arithmetic on GAOL alone,
    when it beats intervals and when it loses to them\\
  \code{13\_rounding\_environment} & The program's own doubles,
    \code{printf}, \code{lrint}, \code{strtod}, threads, \code{cleanup()}:
    what GAOL's rounding direction does to the rest of the program, and how
    to live with it (\pxref{sec:upward-effects})\\
  \code{14\_generic\_programming} & One template evaluated with
    \code{double}, \code{interval} and \code{Dual<1>}: argument-dependent
    lookup, \code{using gaol::pow}, \code{T(0.0)}, the standard algorithms
    and containers with intervals\\
  \code{15\_text\_and\_ieee1788} & Intervals read and written as text, one
    contractor step with the names of IEEE 1788: the literals GAOL reads,
    reading a file, exact output, the namespace \code{gaol\_ieee1788}\\
  \code{16\_Goldstein\_Price} & A million evaluations of the Goldstein-Price
    function, timed: the example of GAOL~4, with the true function; the
    enclosure contains the range, 231 times wider\\
\end{longtable}
\endgroup

\begin{Rem}
The example of GAOL~4 dropped the $+1$ of the Goldstein-Price function,
$(x+y)^2$ instead of $(x+y+1)^2$, and so did the manual of GAOL~4, which
showed the same program and called the enclosure the range of $f$, whose
true range, about $\itv{3}{635255}$, is about 237 times narrower than the
enclosure $\itv{-56254330}{94177270}$. Both now compute the true
function.\newinvfive{} Its enclosure, $\itv{-87881320}{147125080}$ over
$\itv{-2}{2}^2$, contains the range $[3, 1015690.27\ldots]$ and is 231 times
wider than it, for lack of a link between the occurrences of $x$ and $y$: the
dependency problem (\pxref{sec:dependency}).
\end{Rem}

\medskip
A short tutorial, \file{examples/tutorial.md}, writes four of the algorithms
of the examples in programs of a few dozen lines, the condensed forms of
examples 03, 05, 07 and 06: the range of a function by subdivision, the
interval Newton method, a contractor, and a branch and bound. Each program
is compiled with the flags GAOL gives (\pxref{sec:compiling}):
\begin{cmdshell}
% c++ -std=c++11 -O2 $(pkg-config --cflags gaol) prog.cpp -o prog \
      $(pkg-config --libs gaol)
\end{cmdshell}
Its outputs are the same with the SSE2, the NEON and the FPU intervals.
The first program evaluates $x(1-x)$ over
$\itv{0.4}{0.6}$, whose range is $\itv{0.24}{0.25}$: the formula written as
it is gives an enclosure 20 times too wide, $1/4-(x-1/2)^2$, where $x$
occurs once, gives the range itself, and cutting $x$ into $2^k$ pieces with
\code{bisect(0.5)}, whose images are joined by the hull \code{|}, narrows
the enclosure linearly. The second encloses $\pm\sqrt2$ by the interval
Newton method, with \code{div\_rel} on the negative and the positive parts
of $f'(X)$, which keeps the gap around the midpoint that the division
\code{/} would fill. The fourth finds the minimum of $x^4-4x^2+x$ over
$\itv{-3}{3}$ by a branch and bound with a priority queue on the lower
bounds, in 19933 cuts, and prints
$\itv{-5.44419306660636}{-5.444192066610894}$, which contains the minimum
$-5.4441920666108974\ldots$ computed by mpmath; the natural form converges
slowly, and example 06 shows that on the six-hump camel function it needs
25\,057 bisections to bring the enclosure of the minimum within $10^{-3}$,
where the mean-value form with the derivative of \file{dual.h} reaches
$10^{-9}$ in 1\,127, and 261 with a monotonicity test. The third program is
the contractor of Example~\ref{ex:contractor}.

\begin{Ex}\label{ex:contractor}
The forward-backward contractor of $x^2+y^2=1$ evaluates the expression tree
forward, from $x$ and $y$ to the sum, which it intersects with 1; it then
goes backward, and each node is intersected with what the others allow. The
relational functions of GAOL are the backward operations:
\code{sqrt\_rel(a, x)} is the hull of the $x$ of $X$ with $x^2$ in $a$, both
signs included, where \code{sqrt(a)} would lose the solutions with $x<0$
(\pxref{sec:relational-arithmetic}).
\begin{example}
#include <iostream>
#include <gaol/gaol.h>
using namespace gaol;

// The contractor of x*x + y*y = 1: forward-backward evaluation of the tree
// s = a + c, a = sqr(x), c = sqr(y), with s = 1. Returns false when the box
// holds no solution.
bool contract(interval& x, interval& y)
{
  // forward: an enclosure of each node over the box
  interval a = sqr(x);
  interval c = sqr(y);
  interval s = (a + c) & interval(1.0);     // the root must be 1
  if (s.is_empty()) {
    return false;
  }
  // backward: s = a + c gives a = s - c and c = s - a ...
  a &= s - c;
  c &= s - a;
  // ... and a = sqr(x) gives x = +-sqrt(a): sqrt_rel() keeps both signs, and
  // intersects with x itself, where sqrt(a) would keep the positive one only
  x = sqrt_rel(a, x);
  y = sqrt_rel(c, y);
  return !x.is_empty() && !y.is_empty();
}

int main()
{
  interval x(0.8, 3.0), y(-3.0, 3.0);
  std::cout << "box           " << x << ' ' << y << '\n';
  contract(x, y);
  std::cout << "contracted    " << x << ' ' << y << '\n';
  interval u(2.0, 3.0), v(2.0, 3.0);
  std::cout << "box           " << u << ' ' << v << '\n'
            << (contract(u, v) ? "may hold a solution\n"
                                : "proved: no solution\n");
  gaol::cleanup();
  return 0;
}
\end{example}
The program prints the following lines.
\begin{onscreen}
box           [0.8, 3] [-3, 3]
contracted    [0.8, 1] [-0.6000000000000001, 0.6000000000000001]
box           [2, 3] [2, 3]
proved: no solution
\end{onscreen}
The first box is cut to $x$ in $\itv{0.8}{1}$ and $y$ in
$\itv{-0.6}{0.6}$, the bounds being enclosures: the points of the box off
this part are not on the circle. The second box is proved empty of
solutions, its \code{s} being empty. In \code{gaol\_ieee1788},
\code{sqrt\_rel(c, x)} is \ieee{sqrRev(c, x)}, \code{div\_rel(c, b, x)} is
\ieee{mulRev(b, c, x)}, with its arguments in another order, and
\code{nth\_root\_rel(c, n, x)} is \ieee{pownRev(c, x, n)}.
\end{Ex}

\medskip
\findex{check\_examples.py}%
The examples of this manual are checked as well.
\file{manual/check\_examples.py} compiles each example of
the files of \file{manual/v5/chapters/} that shows what it
prints, in its lines \code{@outputs}, with an installed GAOL and the flags
of \cmd{pkg-config gaol}, runs it, and compares what it prints with the
manual, line by line:\newinvfive{} the examples once printed 1 and 0 where
the manual showed \code{true} and \code{false}, and \code{chi()} of the
empty set printed \code{-nan} with GCC where it showed \code{nan}. An
example holding a \code{main()} is compiled as it is; the others are put in
a \code{main()} of their own, after \file{<gaol/gaol>}, \file{<gaol/parser/gaol\_expression.h>} and six headers of the standard library (\file{<cmath>}, \file{<iomanip>},
\file{<iostream>}, \file{<limits>}, \file{<sstream>}, \file{<string>}), with \code{using namespace std} and
\code{using namespace gaol}, or \code{gaol\_ieee1788} in the chapter of its
names. A line printed that is too long for the page may be shown on several
lines of output, each break standing for a blank. Each example is a program
of its own: a format or a precision one sets does not reach the next. A
compiler error points to the line of the example in its file of \file{chapters/}, and
every output line of the manual has to be in an example the script reads;
what the manual shows otherwise, in an \code{onscreen} block after a
program or in a comment, is not checked. The compiler is \env{CXX}, or
\cmd{c++}, and \env{PKG\_CONFIG\_PATH} has to lead to the \file{gaol.pc} of
the GAOL to check; the exit status is 1 when an example does not compile,
fails, or prints something else than the manual shows. It is no test of the
builds: the job Ubuntu 24.04 x86\_64 GCC of the continuous integration runs
it, in C++11, on the GAOL it installs, so that a change of the library that
changes one of the outputs shows.

\section{The continuous integration}\label{sec:ci}

The workflows of \file{.github/workflows/} build GAOL with CMake and run the
tests, with \cmd{make test} (\cmd{cmake -{}-build <build> -{}-target test},
\cmd{RUN\_TESTS} with Visual Studio). They run for each pull request, and
for each push to the branches \code{master}, \code{MATH-CORE} and
\code{configure-clean} only: a push to another branch runs nothing until it
is proposed, which halves the jobs of a pull request;\newinvfive{}
\code{workflow\_dispatch} runs them on any branch. \file{performance.yml} and
\file{manual.yml} are the exceptions: the first runs for the pull requests
that change what it measures only (\pxref{sec:ci-speed}), the second for the
pushes and pull requests that change \file{manual/} or itself. The badges of \file{README.md} show
their state on \code{configure-clean}. Table~\ref{tab:ci-platforms} gives
the systems, the processors and the compilers they build on; the list of
the platforms GAOL supports is in the chapter on the installation
(\pxref{sec:platforms}).

\begin{table}[htbp]
  \centering\renewcommand{\arraystretch}{1.1}\small
  \caption{The platforms of the continuous integration} \label{tab:ci-platforms}
  \begin{tabular}{L{3.2cm}L{3.1cm}L{7.3cm}}
    \hline
    \multicolumn{1}{c}{System} & \multicolumn{1}{c}{Processors} &
    \multicolumn{1}{c}{Compilers and builds}\\
    \hline\hline
    Ubuntu 24.04, 26.04 & x86\_64, arm64 & GCC, Clang; the sanitizers;
      \option{-Wall -Wextra -Werror}; CMake 3.14; GCC 11 and Clang 14 of
      Ubuntu 22.04; GCC 9; a shared library; autotools, meson\\
    Debian 12, 13 (containers) & amd64, arm64, armhf; Debian 12 on i386 &
      GCC; built for size on Debian 13 armhf; autotools and meson on i386
      and armhf\\
    manylinux\_2\_28 & x86\_64, aarch64 & CMake\\
    Alpine (musl) & x86\_64, aarch64 & CMake\\
    Debian 13 under qemu & s390x, ppc64le, riscv64 & CMake; Ubuntu 24.04
      under qemu for POWER9 (GCC 13, \option{-mcpu=power9}), also built
      for size\\
    macOS 15, 26, 27 & arm64, x86\_64 (natively or under Rosetta), a
      universal build & AppleClang, the Clang of LLVM, GCC; CMake,
      autotools, meson, the generator of Xcode; the sanitizers\\
    Windows & x86, x64, arm64 & Visual Studio 2022 and 2026 (Release and
      Debug); clang-cl; MinGW-w64 15 and 14 (x86, x64), 13 and 12 (x86);
      MSYS2 UCRT64 (GCC), CLANG64 and CLANGARM64 (Clang), MINGW32
      (autotools, meson); Cygwin (GCC, x86-64); Clang with the target of
      Visual C++ and the driver of GCC; meson with Visual Studio\\
    \hline
  \end{tabular}

\end{table}

\paragraph{Linux.} The jobs of Ubuntu build with GCC and Clang, also with
the address and undefined behaviour sanitizers, with
\option{-Wall -Wextra -Werror}, with CMake 3.14, with the self-test of the
check of the assembly (\file{tests/tools/asm-order/check\_order.py}
\option{-{}-selftest}, in Python alone), and the check of the library itself
with GCC 13.3 and Clang 18, with and without the rounding direction
preserved,\newinvfive{} with the compilers of Ubuntu
22.04 (GCC 11 and Clang 14 on x86\_64, GCC 11 on arm64), whose images GitHub
retires by April 2027, with GCC 9, which has no
\code{\_\_builtin\_roundeven()}, and as a shared library
(\option{BUILD\_SHARED\_LIBS}, \loption{enable-shared},
\option{-Ddefault\_library=shared}: the static library is the default of the
three builds).\newinvfive{} The test \code{cpack\_stale\_configure} also
runs with \code{CC="ccache gcc"} and a Ninja outside \env{PATH}, which it has
to pass on to the copies of the tree it configures, and the tests run with
the \code{FE\_*} values of Cygwin (newlib, 0 to 3) put over those of glibc
for the sources of CORE-MATH. The containers build Debian 12 and 13 on
amd64, arm64 and armhf, Debian 12 on i386, manylinux\_2\_28 and Alpine
(musl) on x86\_64 and aarch64, Debian 13 under qemu on s390x, ppc64le and
riscv64, and Ubuntu 24.04 (GCC 13) under qemu for POWER9; two jobs build for
size (\code{MinSizeRel}, \option{-Os}) on Debian 13 armhf and for POWER9,
where GCC calls the intersection rather than inlining it, and
\code{rounding\_direction} checks the choices a program makes on
\code{is\_empty()} of an empty interval there too; one job on Debian 13
amd64 tests the empty sets before the bounds in the relations
(\option{-DGAOL\_RELATION\_BOUNDS\_FIRST=0}, the order of RISC-V), where GCC
14 raised the invalid-operation flag for an empty operand.

\paragraph{macOS.} The jobs build macOS 15 and 26 on arm64 and x86\_64,
natively or under Rosetta, macOS 27 on arm64 (on the runner labelled
\code{xcode-27}; there is no job of macOS 27 on x86\_64 under Rosetta, GitHub
having no runner for it), and a universal build (arm64 and x86\_64: the
macros of each part, the tests of the arm64 part natively and the test
programs of the x86\_64 part under Rosetta, and the intervals each part says
it has)\newinvfive{} with CMake, with autotools (\option{-arch} twice in
\env{CFLAGS}, \env{CXXFLAGS} and \env{LDFLAGS}) and with meson, on macOS 15
and 26; the generator of Xcode, with the list \code{arm64;x86\_64} and
without any, where CMake gives Xcode the architecture of the machine; and
the sanitizers with AppleClang, LLVM's Clang and GCC.

\paragraph{Windows.} Visual Studio 2022 and 2026 build on x86, x64 and
arm64, in Release and Debug. On x64 and x86 the intervals are the SSE2 ones
(on x86 their two bounds are two doubles, which need no aligned memory), on
arm64 the NEON ones, with Visual C++ and clang-cl, which a step of each job
checks.\newinvfive{} A job on x86 and two on arm64 ask for the FPU intervals
(\option{GAOL\_SIMD=OFF}), so that their code, and the rounding direction
they keep through the x87 unit and MXCSR, stay covered. Three x64 jobs ask for
the AVX-512 path (\option{GAOL\_PREFER\_AVX512}), two with Visual C++ and one with clang-cl, which they take where the
runner has the instructions, with a warning where it has not; since the
runners have them or not according to the machine GitHub gives, five more
jobs run the unit tests that compute with intervals under Intel SDE, which
emulates a processor with them
(Visual Studio 2022 and 2026, Release and Debug, and clang-cl). clang-cl,
the Clang of Visual Studio, builds x64 in Release and Debug with Visual
Studio 2022, and x86, x64 and arm64 with Visual Studio 2026; on x64 it
defines \code{\_\_x86\_64\_\_} and \code{\_WIN32} but not \code{\_\_WIN32\_\_},
which upstream CORE-MATH tests to read the rounding direction, and on x86
GAOL gives it \option{-msse2}, without which it computes doubles on the x87
unit. MinGW-w64 15 builds x86 and x64, Release and Debug, with the SSE2
intervals on both and \option{-faligned-new=8} on x86, which a step checks;
14 builds x86 and x64, 12 and 13 x86; on x86, GCC 15.2 builds with
\option{GAOL\_PRESERVE\_ROUNDING} and \option{-mfma}, and GCC 13.2 with
\option{GAOL\_FMA=OFF}, the \code{fma()} of mingw-w64 11 on the x87 unit.
MSYS2 builds UCRT64 (GCC, and GCC again with
\option{GAOL\_PRESERVE\_ROUNDING}, with and without \option{GAOL\_FMA}: with
\option{-mfma}, the rounding direction goes through MXCSR alone, which
Linux, clang-cl and Visual C++ x64 with \option{/arch:AVX2} run), CLANG64
(Clang, and Clang again with the direction preserved, which takes the
AVX-512 path where the runner has the instructions, GCC for Windows never
taking it) and CLANGARM64 on the arm64 runners (the NEON intervals). Cygwin
builds GCC for x86-64, Release and Debug, with
\option{-mfma -mprefer-vector-width=128}, all the tests passing but those skipped there (\code{intervalf} and
\code{interval2f}, which no build compiles, and \code{cpack\_stale\_configure},
the CPack of Cygwin making no archive), and a configuration with
\option{GAOL\_FMA=OFF} that \file{gaol/gaol\_config.h} has to
refuse.\newinvfive{} Clang with the target of Visual C++ and the driver of
GCC (\cmd{clang++}, not clang-cl) builds in the environment of Visual Studio
with CMake and meson for x64 and x86, and with autotools for x64, the
library in C++14 there, the standard library of Visual C++ needing it. With
GCC on x64 and x86 in Release, a step reads the objects with \cmd{objdump}
(\file{.github/scripts/check-gcc-windows-fma.sh}):
\file{gaol\_interval\_core.cpp} and the \code{cos} of CORE-MATH
have to hold fused multiply-add
instructions, and no object an instruction on a \code{\%ymm} or \code{\%zmm} register on x64, which GCC would spill to a stack Windows x64 aligns on 16
bytes (GCC bug 54412); on x86 such lines are listed and noticed, not an error.

\paragraph{The three builds.} The workflows also build GAOL with autotools
and meson on Ubuntu (x86\_64, arm64), Debian (i386, armhf), macOS (arm64,
x86\_64), MSYS2 (UCRT64 for x64, MINGW32 for x86) and Cygwin, with meson and
Visual Studio 2022 (x64, x86) and 2026 (x64, arm64), and with meson 0.53.2, the oldest
the meson build takes; run the tests, and build them again with the GAOL
they install, which is the only library installed, CORE-MATH being compiled
into it; build GAOL as a part of another project brought in by
\code{FetchContent} and as a subproject of a meson project, whose code
compiled with \code{gaol\_dep} has to get the flags of \file{gaol.pc}, and
check that GAOL and CORE-MATH are compiled as in Release where the project
gives no build type, and that \option{GAOL\_DEBUG} and \code{enable-debug}
build them for debugging;\newinvfive{} check that the three builds agree on
Ubuntu (x86\_64 and arm64, GCC and Clang), macOS (arm64, x86\_64), Debian (i386,
armhf) and MSYS2 (UCRT64, MINGW32) (the audit of \file{.github/audit},
\exref{sec:scripts}); check that
\cmd{configure}, meson and autoconf read a \file{VERSION.txt} that starts
with a byte order mark or has the line ends of Windows, and that meson reads
it with no Python on \env{PATH}; check that \cmd{make distclean} gives the
source tree back as git has it after the autotools have built, installed
and tested GAOL in it; and check the refusals of \exref{sec:ci-refused}.
Jobs of each build restore the rounding direction
(\option{GAOL\_PRESERVE\_ROUNDING}): Ubuntu x86\_64 GCC and Clang (which
take the AVX-512 path where the runner has it) and arm64 GCC and Clang,
Debian i386 and armhf, macOS arm64, Visual Studio x86, x64 and arm64, clang-cl x64, MinGW-w64 x86, MSYS2 UCRT64, CLANG64 and CLANGARM64, autotools and
meson. Jobs of autotools and of meson on Ubuntu x86\_64 build GAOL without
exceptions and run the examples and the tests, which leave out there the
checks of an exception. The jobs built in Release print the time per
operation in their summary, a table that no other run compares.

\paragraph{Which intervals.} On 64-bit ARM processors the intervals are the
NEON ones by default (\option{GAOL\_SIMD}; \option{GAOL\_ARM64},
\loption{enable-arm64} and \option{-Denable-arm64} were an option of their
own),\newinvfive{} so that the jobs of Ubuntu arm64, macOS arm64, the
aarch64 containers, and those of autotools and meson on arm64 test them
(\pxref{sec:neon}). A few jobs build the FPU intervals instead, so that
their code stays covered, the preserved rounding direction included: the
GCC and the Clang of Ubuntu 24.04 arm64, both with the direction preserved
too, macOS 15 arm64, and autotools and meson on Ubuntu 24.04 arm64. The GCC
and the Clang of Ubuntu 24.04 x86-64 build them too, GCC without and with
the direction preserved and Clang with it: the compound products and
quotients of the FPU intervals hold the direction only by the order of their
stores, which the checks of \code{rounding\_direction} with the direction
preserved make a rule.\newinvfive{} A build that fell back to the FPU
intervals without saying so would keep every test green, so the jobs check
which intervals the configuration took, in a step ``The intervals''
(\file{.github/scripts/check-intervals.sh}, which reads
\code{GAOL\_USING\_\allowbreak NEON\_INSTRUCTIONS} and
\code{GAOL\_USING\_\allowbreak SSE2\_INSTRUCTIONS} in
\file{gaol\_configuration.h} and
fails if they are not the expected ones): the NEON intervals on 64-bit ARM,
the SSE2 intervals on x86-64 and on 32-bit x86 (Visual C++, clang-cl, MinGW-w64, MSYS2 MINGW32, Clang for Visual C++ and Debian 12 i386), and
the FPU intervals where the option is off and on armhf, s390x, ppc64le and
riscv64; two jobs give the SSE2 intervals of Linux the form they have on
x86 with Visual C++ (\option{-DGAOL\_SSE2\_UNALIGNED\_STORAGE=ON}).

\paragraph{Warnings and headers.} Five jobs build the library, the sources
of CORE-MATH included, the tests and the examples with
\option{-Wall -Wextra -Werror}: four with the GCC and the Clang of Ubuntu 26.04
on x86\_64 and on arm64, and one with the AppleClang of macOS 26 arm64. Four
jobs do so with \option{/W4 /WX}, with the Visual C++ of Visual Studio 2026: x64
in Release, and again asking for the AVX-512 path (compiled where the runner has
the instructions), x64 in Debug, and x86 in
Release. On arm64, the check of the installed headers below compiles the header
of the NEON intervals, which \file{gaol/core/gaol\_interval\_core.h} includes,
with the warnings of the headers (\option{-Wold-style-cast},
\option{-Wundef}) that the intrinsics of \file{<arm\_neon.h>} have to stay clear
of.\newinvfive{} The
first five check the headers GAOL installs with
\file{.github/scripts/headers.sh}: each header compiles alone with
\option{-std=c++11}, the oldest standard GAOL takes, and
\option{-Wall -Wextra -Wold-style-cast -Wundef -Werror}, but those that other
headers include in the middle of their code or for Visual C++ only; and the
macros that \file{gaol/gaol} and the two headers of the parser define, the
include guards apart, start with \code{GAOL\_} or \code{gaol\_}.
\option{-Werror} and \option{/WX} go to the targets only
(\option{CMAKE\_COMPILE\_WARNING\_AS\_ERROR}), so that no check of CMake
fails on a warning of its test program. The library is also compiled with
\option{-Wall -Wconversion} in every build, which fails no job.

\paragraph{The locales and the manual.} The jobs on Ubuntu and in the Debian
containers generate \code{fr\_FR.UTF-8} before the tests with
\cmd{localedef} (\file{.github/scripts/comma-locale.sh}), their systems
having no locale writing a decimal comma; macOS has \code{fr\_FR.UTF-8}, and
Windows \code{French\_France.1252}. After the tests, the jobs of Ubuntu, of the Debian containers, of macOS (but the universal builds and the generator of Xcode), of Visual Studio, of MinGW-w64 and of MSYS2 fail if the output of \code{numbers} holds no check under such a locale,
read in \file{Testing/Temporary/LastTest.log}, \file{tests/numbers.log} with
the autotools or \file{meson-logs/testlog.txt} with meson.\newinvfive{}
Alpine and manylinux have no such locale to generate, and check nothing of
the kind. \code{numbers} has 300~s (\option{GAOL\_NUMBERS\_TIMEOUT}), 400 in
the four Linux jobs with the sanitizers, where it took 109 to 287~s and once more
than 300, and 600 in the job of macOS 15 x86\_64 GCC with the sanitizers,
where it took 146 to 367~s on the same code, and once more than 400. The manual of GAOL v5 is
built with the LaTeX of Ubuntu 24.04, by the autotools, the CMake and the
meson builds, when \file{manual/} changes (\file{manual.yml}), and its PDF,
\file{manual/Manual\_GAOLv5.pdf}, is an artefact of the run; the manual of
GAOL~4 is not built; the examples of this manual that show what they
print are compiled by \file{manual/check\_examples.py} in the job Ubuntu
24.04 x86\_64 GCC at each run (\pxref{sec:examples}).

\subsection{The speed of a pull request}\label{sec:ci-speed}

The times of the continuous integration cannot be compared from one run to
the next. The runners are virtual machines shared with other jobs, whose
speed changes by several times between two runs of the same code:
\code{x * y} with Visual C++ x64 took 7.9~ns and then 16.3, and \code{sqrt}
with MSYS2 and GCC 16 6.8~ns and then 170. The step ``Performance'' of the
jobs prints the table of \code{gaol\_performance} in their summary, and
never fails a job: it tells what a platform costs at that moment, and
nothing else. Showing the best of several passes would not help: a pull
request is to be compared with its base, which is not measured at the same
moment.

\medskip
\file{performance.yml} measures two builds in turn, in the same job, so that
what the runner does between two moments reaches both.\newinvfive{} It runs
for a pull request that changes what is measured, not for one that changes
the documentation: \file{gaol/}, \file{3rd/}, \file{tests/performance.cpp},
the files of the benchmark of \file{doc/compare/code/}, \file{CMakeLists.txt}
and \file{cmake/}, and the scripts and the workflow themselves. A job:
\begin{enumerate}
\item checks out the merge commit of the pull request and its first parent,
  the base it is merged into: the difference between the two is the pull
  request alone, whatever the base branch has gained since the pull request
  was started;
\item configures both with CMake in Release and builds \code{gaol\_performance}
  alone, with the options of the job; the programs of the base are kept
  from one run to the next, keyed by the job, the image of the runner, the
  model of the processor, whether it has the AVX-512 instructions
  (processors of one model do not all have them) and the commit of the
  base, so that the next push of the same pull request builds only the pull
  request, and the two builds come from the same tools and the same
  choices;
\item runs \code{gaol\_performance} three times with each build, the base
  first in the first and third runs and the pull request first in the second
  (\file{.github/scripts/performance-ab.sh});
\item writes in the summary of the job the table of the median of each row
  for the base and for the pull request, their ratio (base over pull
  request: above 1, the pull request is faster) and the spread of the runs,
  the largest ratio of the slowest run to the fastest of either build
  (\file{.github/scripts/compare-performance.py}).
\end{enumerate}

\noindent
The jobs are Ubuntu 24.04 x86\_64 with GCC, with Clang, with GCC and the
rounding direction preserved, and with GCC and with Clang and the FPU
intervals; Ubuntu 24.04 arm64 with GCC and with Clang (the NEON intervals)
and with GCC and the FPU intervals; macOS 26 arm64 and macOS 15 x86\_64;
Visual Studio 2026 x64, x86, arm64, x64 with the direction preserved, and x64
with the direction preserved and the AVX-512 path left out, which measures the
macros that write MXCSR alone, and clang-cl for x64, x86 and arm64; MSYS2 UCRT64 and CLANG64; and Debian 13 containers for i386 and armhf, with no cache. The
other platforms (POWER, s390x, RISC-V, MinGW-w64 of Chocolatey, MSYS2 CLANGARM64
and MINGW32) are not compared, nor are the autotools and meson builds:
the three builds give the same flags, which the audit checks.

\medskip
The table never fails a job because of a time. Only a large difference, on
rows that share their code, means something on a shared runner (the same
code has given a ratio of up to 1.17), and a ratio within the spread of the
runs is noise: under the table, the rows whose ratio is 1.5 or more, or
$2/3$ or less, and beyond the spread of the runs are listed, or none is. A
job fails when a build or a program does, and when the tables cannot be
compared: an output with no row, or no row common to the base and the pull
request. The comparison has its test, which needs Python alone and runs
first in each job (\file{.github/scripts/test-performance-tools.py}): it
fails if the median is replaced by the mean or the threshold is inverted. On
Linux, macOS and with Visual Studio the job also runs the ten long lines of
the benchmark of the comparison with the other libraries
(\pxref{sec:comparison}) with \code{gaol\_bench}, the program \cmd{make perf}
builds too: 200\,000 intervals drawn by \file{bench.py}, 3 repeats of each
line, 3 runs of each build, and a second table under the first; the lines
are left out, with a warning, where one of the builds has no
\code{gaol\_bench} that builds or numpy cannot be installed. The job
\code{control-register-performance} of the containers (Debian 13 armhf and
arm64) uses the same scripts to compare two builds of one tree that differ
by a CMake option: the control register of the rounding direction written
directly with \option{GAOL\_PRESERVE\_ROUNDING}, or the functions of the C
library, on the rows that the register can change. The same comparison runs
anywhere with CMake, a C++ compiler and, for the long lines, Python 3 with
numpy:
\begin{cmdshell}
% sh .github/scripts/performance-ab.sh --before-src ../base --after-src . --lines
% python3 .github/scripts/compare-performance.py base perf-before-1.md \
    perf-before-2.md perf-before-3.md -- PR perf-after-1.md perf-after-2.md \
    perf-after-3.md
\end{cmdshell}
\option{-{}-options} and \option{-{}-before-options} give CMake options to
both builds or to the base alone; the script describes them all.

\subsection{Configurations refused or left out}\label{sec:ci-refused}
\cindex{continuous integration!configurations refused}%

Table~\ref{tab:ci-refused} gives the configurations GAOL v5 refuses, and
those the continuous integration leaves out. The builds refuse the compilers
among them when configuring, or \file{gaol/gaol\_config.h} when compiling
(\pxref{sec:refused}), and the last column says which job checks the
refusal. What a configuration gives with GAOL v5 was measured on 21
September 2026, building it past its refusal.

\begingroup\renewcommand{\arraystretch}{1.1}\small
\begin{longtable}{L{2.9cm}L{5.8cm}L{5cm}}
  \caption{The configurations refused or left out}\label{tab:ci-refused}\\
  \hline
  \multicolumn{1}{c}{Configuration} & \multicolumn{1}{c}{Why} &
  \multicolumn{1}{c}{What GAOL v5 does}\\
  \hline\hline
  \endfirsthead
  \multicolumn{3}{l}{\textit{Table~\thetable, continued}}\\*
  \hline
  \multicolumn{1}{c}{Configuration} & \multicolumn{1}{c}{Why} &
  \multicolumn{1}{c}{What GAOL v5 does}\\
  \hline\hline
  \endhead
  \hline
  \endfoot
  Clang for 32-bit ARM (Debian armhf) & Wrong bounds: Clang does not honour
    the rounding direction there. Built by Clang 19 (Debian 13) and Clang 21
    (Debian sid), GAOL v5 fails \code{arithmetic}, \code{elementary},
    \code{reverse}, \code{other\_functions} and \code{rounding\_direction};
    the products are not the tightest in any direction. & Refused by the
    three builds and by \file{gaol/gaol\_config.h}; a job checks that CMake,
    \cmd{configure} and meson refuse it, for the C compiler alone as for the
    C++ compiler alone. GCC builds the armhf jobs.\\
  Clang 14 on 64-bit ARM, and any compiler saying of
    \option{-frounding-math} ``overriding currently unsupported rounding
    mode on this target'' & Wrong bounds: Clang 14 does not honour the
    rounding direction on 64-bit ARM, and says so. GAOL v5 fails
    \code{arithmetic}, \code{elementary}, \code{reverse} and
    \code{other\_functions} in Release, and \code{elementary} in Debug.
    Clang 18 honours it. & Refused by the three builds; a job checks it for
    the C compiler alone, the C++ compiler alone and both. Ubuntu 22.04
    arm64 is built with GCC only.\\
  MinGW-w64 whose \code{fma()} or \code{round()} is wrong: on x64, GCC 11 to
    13 of Chocolatey (mingw-w64 before 12) and the mingw-w64 linked with
    \file{msvcrt.dll} (MSYS2 MINGW64, the cross compilers of Debian and
    Ubuntu); on 32-bit x86, GCC 11 of WinLibs (mingw-w64 9) and, without
    \option{-mfma}, mingw-w64 12 or later linked with \file{msvcrt.dll}
    (MSYS2 MINGW32); on ARM, as on x64 & Wrong bounds. On x64 the
    \code{fma()} of mingw-w64's own math library adds the products of the
    halves of its arguments with four roundings (10.7\,\% of the error-free
    products and 25 to 73\,\% of other triples wrong), and CORE-MATH
    computes with it wherever the compiler has no \option{-mfma}; its
    \code{round()} depends on the direction. GAOL v5 fails
    \code{elementary}, \code{reverse} and \code{core\_math}: bounds of
    \code{tan()}, \code{asin()}, \code{atan()}, \code{sinh()}, \code{cosh()},
    \code{atanh()} and \code{atan2()} of small arguments do not enclose the
    exact values. From mingw-w64 12 on, a toolchain linking the UCRT takes
    the \code{fma()} and \code{round()} of \file{ucrtbase.dll}. & Refused by
    \file{gaol/gaol\_config.h}; seven jobs check the refusal, three of them
    MSYS2. MinGW-w64 GCC 14 and 15 on x64, 12 to 15 on 32-bit x86 (13 also
    with \option{GAOL\_FMA=OFF}), MSYS2 UCRT64 and CLANG64 are built and
    tested; \code{core\_math} checks \code{fma()} and \code{round()}
    first.\\
  Visual C++, and clang-cl from Clang 16, without \option{/fp:strict}
    (\option{/fp:precise}, their default) & Bounds not certified: they then
    assume rounding to nearest, and may evaluate or rewrite floating-point
    operations accordingly, while GAOL computes rounding upward. & Refused
    by \file{gaol/gaol\_config.h}; \code{gaol::gaol} gives
    \option{/fp:strict} to the code linking it, and each Visual Studio job
    checks the refusal (\file{tests/fp\_strict}), those of clang-cl
    included.\\
  The SSE2 intervals on 32-bit Windows, compiled without aligned \code{new}
    & Crash: GCC takes the memory of \code{new} to be aligned on 16 bytes,
    which the C runtime of 32-bit Windows aligns on 8, and a
    \code{std::vector} of SSE2 intervals crashes on its first
    \code{movaps}; with MinGW-w64 14.2 and 15.2 for x86, \code{arithmetic},
    \code{other\_functions}, \code{reverse} and \code{rounding\_direction}
    crash. Visual C++ and clang-cl have no such flag: only
    \option{/std:c++17} turns the aligned \code{new} on. & The three builds
    give \option{-faligned-new=8} to the library and to the code using it
    (\code{gaol::gaol}, \file{gaol.pc}, \code{gaol\_dep}), and
    \file{gaol/gaol\_config.h} refuses the intervals compiled without
    (\code{refused\_aligned\_new*}). Visual C++ and clang-cl keep the two
    bounds as two doubles of alignment 8
    (\code{GAOL\_SSE2\_UNALIGNED\_STORAGE}), which
    \code{test\_intervals\_\allowbreak aligned\_on\_8\_bytes} exercises.\\
  The NEON intervals on 32-bit ARM & Not offered: they need the vectors of
    two doubles of 64-bit ARM, and write the rounding direction into FPCR
    with \code{mrs} and \code{msr}. & The FPU computes the intervals, as
    before, and \file{gaol/gaol\_config.h} stops the compilation of NEON
    intervals elsewhere; the Visual Studio and MSYS2 CLANGARM64 jobs check
    the NEON intervals on arm64, and the armhf jobs the FPU ones.\\
  Cygwin without the fused multiply-add instructions
    (\option{GAOL\_FMA=OFF}, or a processor without them) & Bounds not
    certified: the \code{fma()} of the C library of Cygwin is not correctly
    rounded, \code{fma(1 + 2\textasciicircum{}-30, 1 - 2\textasciicircum{}-30, -1)} giving 0 instead of
    $-2^{-60}$, and CORE-MATH computes with \code{fma()}. & Refused by
    \file{gaol/gaol\_config.h}; the Cygwin job checks the refusal, and
    builds and tests Cygwin with \option{-mfma} (GCC 14.4, Release, three
    tests skipped).\\
  Doubles computed on the x87 unit of 32-bit x86 & Wrong results rounded to
    nearest, and wrong bounds: the x87 unit computes with 64-bit
    significands, and CORE-MATH assumes every operation on doubles rounded
    to a double. With GCC 12 on Debian 12 i386, CORE-MATH rounded to nearest
    gives the other neighbour of the exact value at 175 of the arguments
    compared with SSE2; with GCC 9 on x86-64 (\option{-mfpmath=387}), the
    bounds of \code{exp2(-1075)}, \code{exp10(-400)}, \code{expm1(-800)} and
    of \code{atan2()} of a tiny and a huge number do not enclose the exact
    values. & Refused by \file{gaol/gaol\_config.h}; the builds compile with
    \option{-msse2 -mfpmath=sse} on 32-bit x86, and
    \code{extended\_precision} fails on the x87 unit.\\
  Debian 11 Bullseye (amd64, arm64, armhf) & Cannot be installed: out of
    support since August 2026, its security repository lists packages it no
    longer serves, so that \cmd{g++} cannot be installed in the container.
    GAOL v5 builds and passes its tests there, with GCC 10.2 and CMake
    3.18.4 installed from snapshot.debian.org. & Left out of the continuous
    integration. Debian 12 and 13 are built.\\
  The intervals of floats (\code{gaol::\allowbreak intervalf},
    \code{gaol::\allowbreak interval2f})
    & Unfinished: \code{sqrt(intervalf)} returns its argument and
    \code{interval2f::inverse()} aborts; neither IBEX nor Codac uses them.
    & Compiled by none of the three builds, which have no option for them:
    only a developer of GAOL defines \code{GAOL\_FLOAT\_INTERVALS}. The
    audit still compares it between the builds.\\
\end{longtable}
\endgroup

\subsection{What these platforms showed}\label{sec:ci-showed}

Building and testing GAOL on these platforms, while the workflows were
written and since their first run on 13 September 2026, showed the
following, which GAOL v5 changed; the register of the differences gives each
change with its measures (Chapter~\ref{chap:changes}).

\begin{itemize}
\item[$\bullet$] \textbf{Bounds that were wrong on some platforms only.}
  The hyperbolic functions, taken from the libm of the system: the libms of
  glibc 2.31, musl and MinGW-w64 are sometimes more than one double from the
  exact value; they are now those of CORE-MATH, correctly rounded, on every
  platform.\newinvfive{} The numbers read from text: the C runtime of
  Windows and musl on 64-bit ARM round \code{strtod()} to nearest in every
  direction, and \code{interval("0.1")} did not enclose $1/10$ there;
  numbers are now read exactly, and intervals written in decimal are
  rounded outward by GAOL rather than by the C library.\newinvfive{}
  \code{mid()} with Visual C++ (\option{/O2 /fp:strict}), which rewrote
  $(-0.5)x+y$ into $y-0.5x$: \code{mid([2\textasciicircum{}-1074, MAX])} was wrong in the
  seven Release jobs of Visual Studio. \code{mid()} on 64-bit ARM with
  autotools and meson, which did not give \option{-ffp-contract=off}: the
  compilers fused its multiplications and additions, and it returned a
  single double; the three builds now give the same flags, which the audit
  jobs check on each kind of machine. Square roots with Visual C++ for
  32-bit x86, whose \code{sqrt} rounds to nearest in every rounding
  direction.
\item[$\bullet$] \textbf{Operations that were slow on some platforms only},
  which the table of \code{gaol\_performance} in the summary of each
  Release job showed. Reading the rounding direction before each operation
  made \code{x + y} take 800~ns under Rosetta 2 and 82 to 98~ns with 32-bit
  Visual C++: an addition shows whether it is upward instead.
  \code{fesetround()} cost 130~ns with mingw-w64 13, and 50 to 250~ns with
  Visual C++: on x86 processors the control registers are written directly.
  \option{-ffloat-store}, given to GCC on every target, made \code{x * y}
  take 18.3~ns rather than 4.9 on x86-64: it is given only where doubles are
  still computed on the x87 unit.
\item[$\bullet$] \textbf{Crashes and builds that failed.} \cmd{configure}
  and meson did not install the headers for MinGW and Visual C++, and meson
  did not compile on Linux. The agreement jobs passed while a build had not
  configured in the containers, a pipe hiding the failure: they now fail,
  and print the logs. \cmd{configure -{}-enable-debug} did not compile, the
  Debug jobs of CMake now define \code{GAOL\_DEBUGGING} too, and the test
  \code{debugging} compiles that code. \code{minimum()} and \code{maximum()}
  of two zeros of different signs gave $+0$ and $-0$ with Visual C++ in
  Release, rather than $-0$ and $+0$, which the unit test
  \code{float\_functions} checks, built with Visual C++ since the tests of
  \file{check/} are in \file{tests/}. A shared \file{libgaol} built by
  \cmd{configure} or meson did not export the classes and functions of the
  expressions, and a program building one did not link with it: the shared
  library jobs of the three builds, whose tests link with the installed
  \file{libgaol.so}, build \file{tests/expressions.cpp}.
\item[$\bullet$] \textbf{Reports of the sanitizers that were suppressed at
  first}, and are fixed now: the nodes the parser did not free. Nothing is
  suppressed any more.
\end{itemize}

\section{The scripts of the repository}\label{sec:scripts}
\cindex{scripts}%

Every script of the repository starts with a comment that says how to run it
and what it needs. Tables \ref{tab:scripts-generated} to
\ref{tab:scripts-tools} give what each one is for, so that the right one is
found without opening them all; the scripts of the comparisons with the
other libraries (\file{doc/compare/code/}, \file{doc/compare/itf1788/},
\file{doc/compare/tang2021/}) are described with the comparisons
(\pxref{sec:comparison} and \exref{sec:test-suites}). The files taken from
other projects, which GAOL does not change nor document one by one, are the
scripts of the autotools at the root (\file{config.guess},
\file{config.sub}, \file{configure}, \file{depcomp}, \file{compile},
\file{install-sh}, \file{ltmain.sh}, \file{missing}, \file{test-driver})
and the tools and the CI of CORE-MATH under \file{3rd/math-core/}.

\begingroup\renewcommand{\arraystretch}{1.1}\small
\begin{longtable}{L{5cm}L{8.8cm}}
  \caption{The scripts of the generated files, the manual and the coverage}
  \label{tab:scripts-generated}\\
  \hline
  \multicolumn{1}{c}{Script} & \multicolumn{1}{c}{What it is for}\\
  \hline\hline
  \endfirsthead
  \multicolumn{2}{l}{\textit{Table~\thetable, continued}}\\*
  \hline
  \multicolumn{1}{c}{Script} & \multicolumn{1}{c}{What it is for}\\
  \hline\hline
  \endhead
  \hline
  \endfoot
  \file{gaol/parser/regenerate\_parser.sh} & Regenerates the lexer and the
    parser from the \file{.lpp} and \file{.ypp} sources (flex 2.6.4, bison
    3.8.2), after any change of them (\cmd{make parser})\\
  \file{gaol/float/gaol\_generator\_mul\_i2f.py} & Writes the 256 cases of the
    multiplication of \code{interval2f} by the signs of the operands: a tool for the
    developers (Python~3 and \code{ply}), which no build runs\\
  \file{manual/build-pdf.sh} & Builds the manual v5 with \cmd{pdflatex}, in
    \file{manual/Manual\_GAOLv5.pdf}; \cmd{make pdf} and the target
    \code{pdf} of CMake and meson call it\\
  \file{manual/check\_examples.py} & Compiles and runs each example of the
    manual v5 with an installed GAOL, and compares what it prints with what
    the manual shows\\
  \file{coverage/coverage\_summary.py} & Reads the JSON of gcovr and writes
    \file{coverage/README.md}: lines and branches by file, GAOL apart from
    CORE-MATH, and the lines never run\\
  \file{cmake/gaol\_version.cmake} & \code{gaol\_read\_version()} reads
    \file{VERSION.txt} before \code{project()}, as \file{configure.ac} and
    \file{meson.build} do\\
  \file{cmake/gaol\_package\_source.cmake} & Included by CPack: brings the
    name and the version of the archive of the sources up to date after a
    change of \file{VERSION.txt}\\
\end{longtable}
\endgroup

\begingroup\renewcommand{\arraystretch}{1.1}\small
\begin{longtable}{L{5cm}L{8.8cm}}
  \caption[The audit of the three builds and the checks of the continuous \mbox{integration}]{The audit of the three builds (\file{.github/audit/}) and the
    checks of the continuous integration (\file{.github/scripts/})}
  \label{tab:scripts-ci}\\
  \hline
  \multicolumn{1}{c}{Script} & \multicolumn{1}{c}{What it is for}\\
  \hline\hline
  \endfirsthead
  \multicolumn{2}{l}{\textit{Table~\thetable, continued}}\\*
  \hline
  \multicolumn{1}{c}{Script} & \multicolumn{1}{c}{What it is for}\\
  \hline\hline
  \endhead
  \hline
  \endfoot
  \file{audit/audit.sh} & Configures GAOL with CMake, autotools and meson,
    and writes a JSON file of what each sees (with \file{probe\_build.py})\\
  \file{audit/probe\_build.py}, \file{audit/make\_probe.py} & Read the
    compiler flags and the configuration GAOL sees in the three builds, for
    a C++ source and for a C source of CORE-MATH, through a file that shows
    the configuration once preprocessed\\
  \file{audit/compare.py} & Compares the macros and the flags of the three
    builds, and with \option{-{}-check} fails if they differ\\
  \file{audit/logs.sh} & When the audit fails, prints the end of the logs of
    the three builds, then exits with the same status\\
  \hline
  \file{scripts/tests.sh} & Builds the tests with a GAOL installed by
    autotools or meson (static or shared), then runs them\\
  \file{scripts/headers.sh} & Checks that each installed header compiles
    alone in C++11 with
    \option{-Wall -Wextra -Wold-style-cast -Wundef -Werror}\\
  \file{scripts/\allowbreak check-intervals.sh} & Fails unless the build computes with
    the expected intervals (NEON, SSE2 or FPU)\\
  \file{scripts/\allowbreak check-gcc-windows-fma.sh} & Checks that the objects of \file{gaol\_interval\_core.cpp} and of the \code{cos} of CORE-MATH built by GCC for Windows hold fused multiply-add instructions, and that no object names a 256-bit or 512-bit register (an error on x64, GCC bug 54412; a notice on x86)\\
  \file{scripts/comma-locale.sh} & Generates \code{fr\_FR.UTF-8} with \cmd{localedef} before the tests (\code{generate}), and checks after them that \code{numbers} ran checks under a locale writing a
decimal comma (\code{check}), so that \code{numbers} does not pass without having tried a locale writing a decimal comma\\
  \file{scripts/sde-tests.sh} & Runs the tests that compute with intervals (all
    but \code{numbers}, \code{core\_math} and \code{u128}) under Intel SDE,
    which emulates AVX-512: the runners of GitHub for Windows may have none\\
  \file{scripts/version-file.sh} & Checks how \cmd{configure}, meson and
    autoconf read \file{VERSION.txt} written in several ways\\
  \file{scripts/performance-ab.sh} & Builds and measures two builds in turn
    in the same job: the base and the head of a pull request, or one source
    with two options\\
  \file{scripts/\allowbreak compare-performance.py} & Compares two series of measures
    (median, ratio, spread) and writes a Markdown table for the summary of
    a job\\
  \file{scripts/\allowbreak test-performance-tools.py} & Tests
    \file{compare-performance.py} and the name of the processor, in Python
    alone and without a build\\
\end{longtable}
\endgroup

\begingroup\renewcommand{\arraystretch}{1.1}\small
\begin{longtable}{L{5cm}L{8.8cm}}
  \caption{The scripts of the tests (\file{tests/}) and their tools
    (\file{tests/tools/})}\label{tab:scripts-tools}\\
  \hline
  \multicolumn{1}{c}{Script} & \multicolumn{1}{c}{What it is for}\\
  \hline\hline
  \endfirsthead
  \multicolumn{2}{l}{\textit{Table~\thetable, continued}}\\*
  \hline
  \multicolumn{1}{c}{Script} & \multicolumn{1}{c}{What it is for}\\
  \hline\hline
  \endhead
  \hline
  \endfoot
  \file{elementary\_values.py} & Writes \file{elementary\_values.h}: for
    each $f(x)$, the greatest double below and the least double above,
    computed with mpmath, for \code{elementary}\\
  \file{extended\_precision\_values.py} & Writes
    \file{extended\_precision\_values.h}: the bounding values and the
    nearest double, to compare CORE-MATH and GAOL\\
  \file{compoundm1\_values.py} & Writes \file{compoundm1\_values.h}: the
    doubles around $(1+x)^y-1$ at pairs, the hull over boxes, and
    $\log(1+x)$ and $e^t-1$ to 256 bits, with mpmath and fractions\\
  \file{reverse\_values.py}, \file{atan2\_rel\_values.py},
    \file{pow\_rel\_values.py} & Write the cases of the reverse functions of
    IEEE 1788-2015, of \ieee{atan2Rev1} and \ieee{atan2Rev2}, and of
    \ieee{powRev1} and \ieee{powRev2}, for \code{reverse}\\
  \file{refused\_options.sh} & Checks that autotools and meson refuse
    \option{-ffast-math} and \option{-ffinite-math-only}, and the SSE2
    intervals of 32-bit Windows without \option{-faligned-new=8}, as the tests
    \code{refused\_*} and \code{refused\_aligned\_new*} of CMake do\\
  \file{release\_flags.py} & Checks that a project which brings GAOL in
    without a build type compiles it in release (\option{-O3},
    \code{NDEBUG}, \code{GAOL\_DEBUGGING} at 0)\\
  \file{meson\_subproject/check\_flags.py} & Checks that a meson project
    which takes GAOL as a subproject gets the flags of its \file{gaol.pc}\\
  \file{compile\_error.cmake} & A \cmd{cmake -P} script that compiles a
    source whose compilation has to fail (\code{refused\_*},
    \code{nodiscard\_discard\_cxx*})\\
  \file{cpack\_stale\_configure.cmake}, \file{automatic\_configuration.cmake},
    \file{version\_file.cmake} & The three tests of the CMake build that are
    scripts (\pxref{sec:test-programs})\\
  \hline
  \file{tools/arith/diffarith.py} & Differential test of the arithmetic:
    the bounds and the exception flags of \code{+ - * / \%}, \code{sqr},
    \code{sqrt}, \code{inverse}, \code{div\_rel} and \code{mid} on a grid of
    boxes, under the four rounding directions, with two builds of GAOL\\
  \file{tools/pow/diffpow.py} & Differential test of \code{pow}: the bounds
    of \code{gaol\_ieee1788::pow} and \code{gaol::pow} with two builds, on
    70\,500 boxes and under the four directions\\
  \file{tools/pow/build.py} & Compiles a program of \file{tests/tools/pow}
    against a build of GAOL, with the flags that build gives its own tests\\
  \file{tools/pow/negative\_power\_rows.py} & Writes the rows of the table of
    large negative powers of \file{tests/arithmetic.cpp}: the tightest bounds
    of $x^n$ for $n$ from $-2^{31}+1$ to $-2^{27}$, computed with mpmath on
    2000 bits\\
  \file{tools/pow/gen\_table.py}, \file{tools/pow/checkrows.py} & Read and
    rewrite the table of boxes of \code{ieee1788}, and check its bounds
    against the exact power (mpmath, 500 bits)\\
  \file{tools/pow/mutate.py}, \file{tools/cos/mutate.py},
    \file{tools/x86-csr/mutants.py} & Mutation testing: one deliberate error
    in \code{pow} and the integer powers, in the port of \file{cos.c} to the
    128-bit integer, or in the writes of the control registers of x86, and
    the tests have to fail\\
  \file{tools/pow/pown\_bound.py} & Checks with rationals the bound of the
    error of the integer powers computed beforehand\\
  \file{tools/pow/bench.cpp}, \file{tools/pow/pownhash.cpp},
    \file{tools/pow/pownedge.cpp} & The time per call of each path of
    \code{pow}, and a hash of the bounds of four million integer powers, the
    sign of a zero included\\
  \file{tools/cos/accurate\_values.py} & Computes the roundings of
    $\cos x$ with mpmath and checks the table of the accurate path of
    \code{core\_math}\\
  \file{tools/asm-order/check\_order.py} & Reads the assembly of the library
    (x86-64, GCC and Clang) for an instruction that rounds on the wrong side
    of a write of the control register (GCC bug 34678), and for the
    conversion of a subnormal constant through the x87 unit\\
  \file{tools/ibex\_codac/*.py} & The cases of the tables of
    \code{test\_ibex\_codac}, their exact results (rationals, mpmath), and
    the check that the rows of the test are those the cases give\\
  \file{tools/x86-csr/\allowbreak amalgamate.py} & One source file from three units of
    GAOL, to compile with a compiler that takes one file, such as those of
    Compiler Explorer\\
\end{longtable}
\endgroup

\noindent
The workflows are read by actionlint, which runs \cmd{shellcheck} on the
steps that have a script, and reports nothing; the scripts of
\file{.github/} and of \file{doc/compare/code/} are read by
\cmd{shellcheck}, run from the root of the repository, which reports
nothing either: the lists of flags that the shell has to split carry a
comment with the reason.\newinvfive

\section{Coverage}\label{sec:coverage}
\cindex{coverage}%

The option \option{GAOL\_COVERAGE} of CMake (\loption{enable-coverage} of
\cmd{configure}, \code{enable-coverage} of meson), off by default, compiles
GAOL and its tests with the counters of gcov, and adds the target that runs
the tests and writes the report of their coverage,
\file{coverage/coverage.html} and the page of conclusions
\file{coverage/README.md}: \cmd{cmake -{}-build <build> -{}-target coverage},
\cmd{make coverage}, or \cmd{ninja gaol-coverage}, meson keeping
\code{coverage} for its own report (\pxref{sec:cmake-options}).\newinvfive{}
It needs GCC or Clang, gcovr and Python~3, and the tests (\option{WITH\_TESTS}); it should be the Debug build: with the Release build, the optimizations move the code around, the counts no longer follow the source, and the build warns. The sources of GAOL
(\file{gaol/}) are counted apart from those of CORE-MATH
(\file{3rd/math-core/}), whose accurate phases are reached by a handful of
arguments in a million: counting them together would say nothing of either.
Table~\ref{tab:coverage} gives the report of the run of 4 November 2026.

\begin{table}[htbp]
  \centering\renewcommand{\arraystretch}{1.1}
  \caption{The coverage of the tests on 4 November 2026} \label{tab:coverage}
  \begin{tabular}{lrrrr}
    \hline
    & Lines & Covered & Not run & \%\\
    \hline\hline
    GAOL (\file{gaol/}) & 5466 & 5153 & 313 & 94.3\\
    CORE-MATH (\file{3rd/math-core/}) & 8172 & 6732 & 1440 & 82.4\\
    \hline
  \end{tabular}

\end{table}

\noindent
The tests run 94.3\,\% of the lines of GAOL, above the 80\,\% GAOL v5 aims
at, and 66.5\,\% of its branches are taken both ways. The lines of CORE-MATH
that are never run are its accurate phases, which a handful of arguments in
a million reach, and the paths of the architectures the machine is not: they
are counted among the lines not run, and are checked otherwise: the accurate
path of \code{sin} is compared with the upstream sources compiled unchanged,
and that of \code{cos} with its roundings computed by mpmath
(\pxref{sec:test-programs}), and the other accurate phases are reached by the
tests of their bounds only. File by file, the least
covered are \file{gaol/parser/gaol\_expr\_visitor.h} (9 of 99 lines, 9.1\,\%),
\file{gaol/parser/gaol\_expr\_eval.h} (87.3\,\%),
\file{gaol/parser/gaol\_expression.cpp} (89.6\,\%) and
\file{gaol/core/gaol\_interval\_sse.cpp} (93.9\,\%);
\file{gaol/core/gaol\_interval\_core.cpp}, the largest file, with 2492 lines counted,
is covered at 97.0\,\%, and 18 files of the 26 at 100\,\%. The page lists
the lines the tests never run, file by file.

\section{Discussion}

The tests of GAOL v5 compare its bounds with references that owe nothing to
the library nor to the floating-point environment, and run every operation
in every state a program may leave the processor in: the four rounding
directions, the modes that flush the subnormal numbers, the exceptions
enabled, a locale with a decimal comma, and the initialization of static
objects before \code{main()}. The programs of \file{tests/} hold hundreds
of thousands of such checks, and the documentation records, for many of the changes, how many of the checks written with them fail without them, and the example
programs check what they print, as do the examples of this manual. The continuous integration runs
the suite on three operating systems, seven processor architectures and the
compilers of each, with the three builds, and refuses the configurations
whose bounds it found wrong or uncertified, most of the refusals being checked by a
job of its own.

\medskip
Three choices of the tests and of the continuous integration deserve a word. The teeth of some checks depend
on the processor the code is compiled for: a loop that a compiler
vectorizes raises the invalid-operation flag with one set of instructions
and not with another, so that the jobs on many processors are the test, not
one of them. The speed of a pull request cannot be read from the times of a
shared runner, which change by several times between two runs of the same
code: it is measured by building the base and the head in the same job and
running them in turn, and only a ratio of 1.5 or more, or of 2/3 or less, beyond the spread of
the runs, means something. The coverage report counts the lines of GAOL
apart from those of CORE-MATH, whose accurate phases a handful of arguments
in a million reach; \code{core\_math} checks the accurate paths of \code{sin}
and \code{cos} against values computed apart instead.

\medskip
Some checks are deferred or left out, for the reasons the documentation
gives: the reading of a text of 512~MiB, which found a 32-bit exponent
overflow, has no test, since the text is read through two more copies of
it; the checks of several threads reading or writing at once are commented
out, the tests running no thread; the refusal of the mingw-w64 of ARM whose
\code{fma()} is wrong is not tested; Debian 11 is left out of the
continuous integration, and the intervals of floats are compiled by none of
the builds. The next chapter turns to the programs written for GAOL~4: how
to compile and link them with GAOL v5, which results change, and the
register of the differences by theme.

%% !TEX root = ../manual_gaolv5.tex
% Chapter 15.
%% Chapter 15: Moving a program from GAOL 4 to GAOL v5. Sources: gaol.tex of f5bf180
%% lines 6512-6639 (the chapter of the same name), doc/differences.md (whole).
%%
%% Copyright (c) 2026 ENSTA, France
%%
%% By: Jordan NININ (created 2026-11-06)
%%
%% GAOL is a software distributed WITHOUT ANY WARRANTY. Read the associated
%% COPYING.LIB file for information.
\chapter{Moving a program from GAOL~4 to GAOL v5}\label{chap:changes}
\cindex{GAOL~4!moving from}%
\cindex{migration from GAOL~4}%

GAOL v5 continues the GAOL 4.2.2 of Fr\'ed\'eric Goualard, and most programs
written for GAOL~4 compile and run with it once five things are done: the
namespace is opened, the strings they gave to \code{interval(const char*)}
are read by \code{textToInterval()}, the relations \code{==} and \code{!=} on
intervals are rewritten with \code{set\_eq()}, \code{set\_disjoint()} or
\code{!set\_disjoint()}, the handlers of the exceptions of GAOL~4 are
retargeted on those of the standard library, and the conversions of a
\code{double} or an integer to an interval, explicit now, are written out.
The other changes concern what a program
sees: tighter bounds, the empty set where GAOL~4 built an interval, the
comparisons of IEEE 1788-2015 \cite{IEEE1788:2015}, and texts read and written
otherwise. Each change was made for a reason, and most of them were measured.

\medskip
In this chapter, we gather what a program may have to change and what it
sees differently. The first section lists what concerns the compilation and
the link, the second the results that change. The third section is the
register of the differences, theme by theme: for each change, what GAOL~4
did, what GAOL v5 does, why, and what a program sees, with the measurements
that justified it and the chapter of this manual that describes the new
behaviour. The last section discusses the choices made.

\medskip
A word on the versions. ``GAOL~4'' is the 4.2.2, which GAOL v5 continues;
the benchmarks of Chapter~\ref{chap:computing} use the 4.2.3, the last
version of Fr\'ed\'eric Goualard, which came after it. The releases 4.3.0 to
4.3.2 of the repository of GAOL v5 came next and are not GAOL~4: they already
had \code{rnd\_keep()}, \code{get\_rounding()} and \code{MEMFREE()}
(\code{GAOL\_MEMFREE()} now), which neither the 4.2.2 nor the 4.2.3 has. The
meson build that this chapter counts with GAOL~4 is the one added to the
repository of the 4.2.3 in November 2025, after its last release: the 4.2.2 had
\cmd{configure} only.
GAOL~4 also had classes of exceptions of its own,
\code{input\_format\_error}, \code{invalid\_action\_error} and
\code{unavailable\_feature\_error}, derived from \code{gaol\_exception}. GAOL
v5 throws the \code{std::invalid\_argument} and the \code{std::domain\_error}
of the standard library where GAOL~4 threw the first two, and has none of
these classes (\pxref{sec:exceptions}). Where this chapter says what GAOL~4
threw, it gives the class of GAOL~4.

\section{Compiling and linking}\label{sec:migration-compile}
%%-----------------------------
\cindex{GAOL~4!compiling a program of}%

The changes below are the ones a compiler or a linker reports, and those of
the initialization and the cleanup, which a program makes in its text. Most of
them are described in the section the cross-reference names.

\begin{itemize}
\item[$\bullet$] \textbf{The namespace.} \file{gaol/gaol.h} no longer opens
  the namespace \code{gaol}: a program adds \code{using namespace gaol;}, or
  names \code{gaol::interval} (\pxref{sec:namespaces}). The type and the
  functions are in \code{gaol\_core}, and the symbols of the library name it:
  a program compiled against the headers of GAOL~4 has to be compiled again.
\item[$\bullet$] \textbf{The names GAOL keeps for itself.} \code{gaol} holds
  the names GAOL~4 had there that GAOL v5 still has, except those GAOL uses
  for itself, which a program names in \code{gaol\_core}
  (\pxref{sec:namespaces}): \code{gaol\_signbit()}, \code{uintdouble},
  \code{ullidouble}, \code{Interval\_struct}, \code{the\_null\_expr},
  \code{prec\_t}, \code{modulo\_k\_pi()} and the \code{round\_*\_sse()}, all of
  them internals that the manual of GAOL~4 did not document, and, of the 4.3
  releases of GAOL v5,
  \code{rounding\_state}, \code{get\_rounding()}, \code{set\_rounding()},
  \code{round\_upward\_if\_needed()} and \code{gaol\_set\_rounding\_x86()}.
\item[$\bullet$] \textbf{Visual C++.} The code using GAOL need no longer define
  \code{\_\_GAOL\_PUBLIC\_\_} empty: GAOL is a static library there, and its
  declarations carry no \code{\_\_declspec(dllimport)}. A definition left in
  a project changes nothing. The macro is \code{GAOL\_PUBLIC} now.
\item[$\bullet$] \textbf{The allocator.} The allocator of containers
  \code{aligned\_allocator} is gone, with its header,
  \file{gaol/gaol\_allocator.h}.
\item[$\bullet$] \textbf{One library to link.} GAOL is the only library to
  link: the mathematical library GAOL~4 was built with is no longer used,
  CORE-MATH being compiled into GAOL, and the options of the builds that chose
  it are gone.
\item[$\bullet$] \textbf{The flags of interval arithmetic.} The code
  including the headers of GAOL is compiled with the flags of interval
  arithmetic, which \code{gaol::gaol} and \file{gaol.pc} give
  (\pxref{sec:compiling}); \file{gaol/gaol\_config.h} refuses the options and
  the compilers that give wrong bounds (\pxref{sec:refused}).
\item[$\bullet$] \textbf{The intervals of floats.} \code{gaol::intervalf} and
  \code{gaol::interval2f} are compiled by none of the builds, which have no
  option for them: only a developer of GAOL compiles them, defining
  \code{GAOL\_FLOAT\_INTERVALS} in \file{gaol/gaol\_config.h}.
\item[$\bullet$] \textbf{The bounds at doubles.} The functions of doubles
  rounded downward and upward, \code{exp\_dn()}, \code{sin\_up()},
  \code{ipow\_up()}, \code{ipow\_dn()}, \code{pow\_up()}, \code{pow\_dn()}
  and the others of \file{gaol/gaol\_double\_op.h}, are no longer part of the
  interface: the header is no longer installed. \code{uipow\_upup()} and
  \code{uipow\_dnup()} are no longer declared.
\item[$\bullet$] \textbf{The static constants.} The static intervals
  \code{interval::cst\_one}, \code{cst\_pi} and the other \code{cst\_}
  members are gone: \code{interval::one()}, \code{interval::pi()} and the
  other functions give these constants, right even in the initialization of a
  static object (\pxref{chap:init}).
\item[$\bullet$] \textbf{The initialization.} \code{gaol::init()} need not be
  called again after the program changed the rounding direction: each
  operation sets it upward itself (\pxref{sec:rounding-direction}). It need
  not be called at all, GAOL initializing itself with every compiler, Visual
  C++ included (\pxref{chap:init}).
\item[$\bullet$] \textbf{The cleanup.} The program calls
  \code{gaol::cleanup()} right after its last use of GAOL: unless the
  rounding direction is preserved, it sets back the rounding direction that
  the initialization found, which GAOL~4 left upward when a static object of the
  program had computed before its initialization (\pxref{chap:init}). It
  no longer frees what the initialization allocated, which the automatic
  cleanup frees when the program ends.
\item[$\bullet$] \textbf{The relations.} The operators \code{==} and
  \code{!=} on intervals are gone, with the possibly relations
  (\code{possibly\_eq()}, \code{possibly\_le()}\ldots),
  \code{certainly\_eq()} and \code{certainly\_neq()}, and so is the option
  \loption{enable-relations} of \cmd{configure} and meson: a program compares
  intervals with \code{set\_eq()}, \code{set\_disjoint()} and
  \code{!set\_disjoint()}, and \code{<}, \code{<=}, \code{>} and \code{>=}
  are the certainly relations (Chapter~\ref{chap:relations}). The new
  relations do not give the results of the old operators: \code{==} was
  \code{certainly\_eq()} in every build by default, true for two empty sets and
  for a point interval equal to the upper bound of the other interval
  (\code{[4] == [3, 4]} was true, \code{[3, 4] == [4]} false), and \code{!=} was
  its negation, so that
  \code{==} of $[3, 4]$ and $[3, 4]$ was false, where \code{set\_eq()} is
  true, and \code{!=} was true, where \code{set\_disjoint()} is false. A
  program chooses the relation it means.
\item[$\bullet$] \textbf{The constructors from strings.} The constructor from
  one string, \code{interval(const char*)}, is gone, and with it the
  conversion of a string wherever an interval is expected:
  \code{textToInterval("0.1")} reads a string in its place, and throws
  \code{std::invalid\_argument} where the constructor threw
  \code{input\_format\_error}; \code{x = "[1, 2]"} becomes \code{x = textToInterval("[1, 2]")}, and
  \code{x + "0.1"} becomes \code{x + textToInterval("0.1")}
  (\pxref{sec:constructors}). The constructor from two strings is gone too:
  \code{interval("0.1", "0.3")} becomes \code{textToInterval("0.1", "0.3")},
  which gives the empty set for bounds in the wrong order. \code{interval(0)}
  and \code{interval(0, 0)}, which these constructors made ambiguous, compile.
\item[$\bullet$] \textbf{The constructor from a double.}
  \code{interval(double)} is explicit, and so is the constructor from an
  integer: \code{interval x = 0.5;}, \code{interval x = 5;},
  \code{return 0.5;} in a function returning an interval, \code{min(x, 0.5)}
  and \code{interval r = sqrt(2.0);} do not compile, where GAOL~4 converted the
  double without a word. \code{x = 0.5} and \code{x < 0.0}
  still compile (\pxref{sec:constructors}).
\item[$\bullet$] \textbf{The exceptions.} The classes \code{gaol\_exception},
  \code{input\_format\_error}, \code{invalid\_action\_error} and
  \code{unavailable\_feature\_error} are gone, with
  \file{gaol/gaol\_exceptions.h}: a handler names
  \code{std::invalid\_argument} or \code{std::domain\_error}
  (\pxref{sec:exceptions}).
\item[$\bullet$] \textbf{The headers.} \file{gaol/\allowbreak gaol\_profile.h},
  \file{gaol/\allowbreak gaol\_parameters.h}, \file{gaol/\allowbreak gaol\_limits.h},
  \file{gaol/\allowbreak gaol\_assert.h}, \file{gaol/\allowbreak gaol\_exact.h} and
  \file{gaol/\allowbreak gaol\_flags.h} are gone, and so are
  \file{gaol/\allowbreak gaol\_config\_msvc.h},
  \file{gaol/\allowbreak gaol\_config\_mingw.h},
  \file{gaol/\allowbreak gaol\_version\_msvc.h},
  \file{gaol/\allowbreak gaol\_version\_mingw.h} and
  \file{gaol/\allowbreak gaol\_flt\_output.h} (the source
  \file{gaol/\allowbreak dllmain.cpp} too); \file{gaol/\allowbreak gaol\_init\_cleanup.h} is no longer
  installed, and a program that builds expressions includes
  \file{gaol/parser/\allowbreak gaol\_expression.h}, with
  \file{gaol/parser/\allowbreak gaol\_expr\_eval.h} to evaluate them, where it
  included \file{gaol/\allowbreak gaol\_expression.h} (\pxref{sec:changes-interface}).
\item[$\bullet$] \textbf{The expressions.} \code{eval\_stack} and
  \code{expr\_eval} are no longer copyable, the stack of an evaluation
  growing in the heap (\pxref{sec:changes-io}).
\item[$\bullet$] \textbf{The debugging macro.} \code{GAOL\_DEBUG(level,
  command)} of the debug build is \code{GAOL\_DEBUG\_VERBOSE}: the old name is
  not kept (\pxref{sec:debugging-facilities}).
\item[$\bullet$] \textbf{The options of the builds.} The option
  \code{with-test} of meson is \code{with-tests}. \loption{enable-asm} of
  \cmd{configure} and \code{enable-asm} of meson are gone, GAOL v5 no longer
  calling the assembly they chose: \cmd{configure} warns of an option it does
  not know, and \cmd{meson setup -Denable-asm=true} stops with an error
  (\pxref{sec:changes-builds}).
\end{itemize}

\section{Results that change}\label{sec:migration-results}
%%---------------------------
\cindex{GAOL~4!results that change}%

The changes below are the ones a program sees at run time without any
change of its text. The third section gives their reasons.

\begin{itemize}
\item[$\bullet$] \textbf{Static objects.} The intervals computed in the
  initialization of the static objects of a program linked with the static
  library GAOL are right: \code{interval::pi()} was $[-0, 0]$ there, and
  $[1, 2]/[-1, 1]$ was $[-0, 0]$ rather than $[-\infty, +\infty]$, the
  constants of GAOL being computed after these objects, except with MinGW-w64
  (\pxref{chap:init}).
\item[$\bullet$] \textbf{The bounds are tighter.} The elementary functions
  give the tightest bounds, rather than one double beyond, and are exact
  where their value is a double; the integer powers are
  computed from exact products, and the $n$th roots are proved with powers
  rounded upward and downward; \code{sin}, \code{cos} and \code{tan} are
  within one double of the tightest bounds at every magnitude
  (Chapter~\ref{chap:accuracy}).
\item[$\bullet$] \textbf{The empty set.} The constructors give it for
  \itv{+\infty}{+\infty}, \itv{-\infty}{-\infty}, bounds in the wrong order
  and NaN bounds, and assigning an infinity to an interval empties it
  (\pxref{sec:constructors}); the intersection of disjoint intervals is it,
  with NaN bounds; so are \code{log} of an interval holding no positive
  number, \code{invabs\_rel(J, I)} for a negative \var{J} and \code{div\_rel} with
  no solution in \var{I}. The operators of an interval and an infinite
  \code{double} (\code{x * INFINITY}, \code{x + INFINITY}) give it too.
  \code{width()} of the empty set is NaN rather than $-1$.
\item[$\bullet$] \textbf{Powers.} \code{pow(x, p)} for a double \var{p} no
  longer truncates \var{p}; \code{pow(x, y)} keeps the negative part of
  \var{x} for an integer exponent only, and takes the \ieee{pow} of IEEE
  1788-2015 otherwise (\pxref{sec:powers}). \code{nth\_root} of an odd order
  is defined on the whole real line, and takes a negative order.
\item[$\bullet$] \textbf{The comparisons} follow IEEE 1788-2015 for the empty
  set and the infinite bounds: \code{certainly\_leq()} and its companions are
  true when either interval is empty, and \code{set\_strictly\_contains()}
  takes an infinite bound as beyond the same infinite bound (Chapter~\ref{chap:relations}).
\item[$\bullet$] \textbf{Input and output.} Numbers are read exactly, and the
  literals of IEEE 1788-2015 are read, in any case of letters;
  \code{"[inf]"} is the empty set. Every string is an expression, where an
  interval may stand wherever a number may (\code{"1+[1,2]"}), an uncertain
  number being no bound (\code{"[5?1]"} is refused) and needing a blank
  after a sign next to it (\code{"10-3.56?1"} is refused,
  \code{"10 - 3.56?1"} is read), and an error anywhere
  stops the reading. The reader knows every function of GAOL, and
  \code{gaol\_ieee1788::textToInterval} reads the names of IEEE 1788-2015
  (\code{pown}, \code{rootn}\ldots). The format \code{hexa} writes interval
  literals, which are read back bit for bit; the decimal formats round the
  bounds outward whatever the C library, and the format \code{agreeing} keeps
  the digits of each bound (Chapter~\ref{chap:io}).
\item[$\bullet$] \textbf{New operations.} \code{atan2()} is implemented, and
  the functions of IEEE 1788-2015 GAOL did not provide are added
  (\code{exp2}, \code{hypot}, \code{sign}, \code{trunc}\ldots,
  Chapter~\ref{chap:arithmetic}).
\end{itemize}

\section{The differences, by theme}\label{sec:migration-themes}
%%--------------------------------------------
\cindex{GAOL~4!differences with}%

This section condenses the register of the differences between GAOL~4 and
GAOL v5, theme by theme. The themes are those of the register, regrouped: the
elementary functions first, then the rounding direction, the floating-point
exceptions, the constants, the kernels, the intervals and their functions, the
powers and the roots, the other functions, the relations, the reading and the
writing, the errors, the interface and the builds. Some of the defects told
below were those of the earlier versions of GAOL v5 only, and the text says
so. Each measurement is given with the machine and the compiler it was made
on wherever the register gives them. The timings are those of the step that
made each change, which may be an intermediate one; those of the release
against GAOL 4.2.3 are in Chapter~\ref{chap:computing}.

\subsection{The elementary functions and CORE-MATH}
\label{sec:changes-core-math}
\cindex{CORE-MATH!differences with GAOL~4}%

Every elementary function is bounded with CORE-MATH \cite{CoreMath:2022}, on
every architecture and with every compiler, and it is the only library GAOL
takes its elementary functions from.\newinvfive{} The other mathematical
libraries GAOL~4 could be built with are gone, with their sources, their
headers and the scripts that installed them, and so are the options of the
builds that chose one: there is nothing to choose, to find, to build apart or
to link along with GAOL, whose library holds the functions
(\pxref{sec:math-libraries}). No build names \option{-lm}: the C library
still gives GAOL \code{sqrt()}, \code{floor()}, the functions of
\file{<fenv.h>} and, where the processor has no fused multiply-add
instruction, \code{fma()}, but \cmd{g++} and \cmd{clang++} link the math
library themselves, so that the autotools build (whose \code{LDADD} was
\option{-lm} in GAOL~4), meson and the \file{gaol.pc} the builds install do
not name it, as the CMake build never did.

\medskip
The bounds are the tightest ones rather than one double beyond. CORE-MATH is
correctly rounded in the rounding direction in effect, so GAOL takes the
value at a bound, computed in the upward rounding it keeps, as the upper
bound, and the double below it as the lower one. GAOL~4 took values correctly
rounded to nearest only, so each bound had to be moved one double outward,
and the direction had to be set to nearest and back around every call. On an
Intel i7-1185G7 with Clang 18, the bounds of \code{log} took 30 ns rather
than 62. The accurate phases of \code{log},
\code{sin}, \code{cos}, \code{tan}, \code{atan2} and \code{pow} compute with
a 128-bit integer, which Visual C++ has on no architecture and GCC has on no
32-bit target: \file{gaol/math/gaol\_u128.h} gives them the type of the
compiler where it has one and a structure of two 64-bit halves where it has
none. The two paths were compared with the pristine upstream sources over
about 40 million arguments in the four rounding directions, and give the same
bits; a job of the continuous integration builds the whole of GAOL with the
halves forced (\code{-DGAOL\_U128\_EMULATION=ON}, in \file{linux.yml};
\pxref{sec:ci}).

\medskip
MinGW-w64 is no longer refused for its version alone, both former reasons
being gone: the hyperbolic functions of its math library, which GAOL no
longer uses, and the cost of its \code{fesetround()}, which GAOL no longer
calls for its elementary functions.\newinvfive{} Still refused are the MinGW-w64 whose
\code{fma()} or \code{round()} is wrong, those of its own math library,
computed in doubles: that \code{fma()} is not correctly rounded (10.7\,\% of
the error-free products \code{fma(a, b, -a*b)} were wrong with mingw-w64 11
on x86-64), and CORE-MATH computes with it wherever the compiler has no
\option{-mfma}; that \code{round()} depends on the rounding direction
(\pxref{sec:refused}).

\medskip
The functions of IEEE 1788-2015 that GAOL~4 did not provide are
added.\newinvfive{} \code{exp2}, \code{exp10}, \code{log2} and \code{log10},
required among the forward elementary functions (Table~9.1 of the standard),
are bounded with CORE-MATH's, the tightest bounds, and the value that is a
double is kept as a bound rather than moved: $2^3 = 8$, $10^{22}$,
$\log_2(1/4) = -2$ and $\log_{10}(100) = 2$ are exact. Sixteen functions of
Table~10.5, which the standard recommends, follow: \code{expm1},
\code{exp2m1}, \code{exp10m1}, \code{log1p}, \code{log2p1}, \code{log10p1},
\code{hypot}, \code{rsqrt}, \code{sinpi}, \code{cospi}, \code{tanpi},
\code{asinpi}, \code{acospi}, \code{atanpi}, \code{atan2pi} and
\code{compoundm1}. \code{sinpi}, \code{cospi} and \code{tanpi} find their
extrema and poles exactly, at the multiples of $1/2$, which are doubles:
\code{sinpi([1e17])} is 0, where \code{sin(pi*x)} gave \itv{-1}{1}, the
width of $\pi$ times $10^{17}$ being more than a period.
$\mathrm{compoundm1}(x, y) = (1 + x)^y - 1$ is the one CORE-MATH has no
binary64 version of, and $\exp(y\,\mathrm{log1p}(x)) - 1$ computed with
doubles is up
to 1\,100 doubles from the value: it is computed with integers only, in
64-bit fixed point where that proves the bounds (95\,\% of the pairs for
$|y\,\mathrm{log1p}(x)|$ below 0.35, in 0.1~$\mu$s a corner) and with a 128-bit
significand in balls otherwise (0.6 to 0.9~$\mu$s). \ieee{fma},
\ieee{cancelMinus} and \ieee{cancelPlus} (\code{cancel\_minus} and
\code{cancel\_plus}), which the standard requires the tightest, and
\ieee{less}, \ieee{strictLess}, \ieee{isEntire} and \ieee{isCommonInterval}
are provided too. The error function \code{erf} and
its complement \code{erfc}, which the standard does not name, are bounded
with CORE-MATH's as well, exact at 0 and at the infinities. The integer
functions \code{sign}, \code{trunc}, \code{round\_ties\_to\_even} and
\code{round\_ties\_to\_away} are exact, each being non-decreasing; they take
nothing from CORE-MATH, and \code{round\_ties\_to\_even} reads the bits so as
to ignore the rounding direction, which \code{std::nearbyint} and
\code{std::rint} read (Chapters~\ref{chap:arithmetic} and
\ref{chap:functions}). The namespace \code{gaol\_ieee1788} gives the
operations of the standard GAOL provides, under the names and in the argument
order of the standard (Chapter~\ref{chap:ieee1788}).

\medskip
CORE-MATH reads the rounding direction as GAOL sets it.\newinvfive{} GAOL~4
had no CORE-MATH, and the defects below were found and fixed in the earlier
versions of GAOL v5 that took it. Some
sources of CORE-MATH round a result themselves, in the direction they read,
and two of the ways they read it gave another direction than that of the
doubles GAOL computes, the one of MXCSR. \code{pow} rounds a subnormal
result in the direction \code{fegetround()} gives, which the C library of
Linux reads from the x87 unit: with the x87 unit to nearest and MXCSR upward,
the upper bound of \code{pow} was below the exact value for about half of the
arguments with a subnormal result. GAOL gives the sources a
\code{fegetround()} that reads MXCSR on x86. On 32-bit x86 that spares
\code{cbrt} and \code{asinpi} a call of the C library at each bound, about
80 ns a call with Visual C++ for x86: on the runners of the continuous
integration, \code{nth\_root(x, 3)} takes 179 ns rather than 355 and
\code{asinpi(x)} 794 ns rather than 1005 with Visual C++ for x86, and 120 ns
rather than 324 and 255 ns rather than 438 with clang-cl for x86.
\code{cbrt}, \code{rsqrt} and \code{asinpi} read the rounding field of MXCSR
and turned it into the \code{FE\_*} value of the C library by a shift chosen
from the macros of the platform; with clang-cl and with Cygwin on x86-64,
whose newlib gives \code{FE\_*} the values 0 to 3, the shift was wrong for
the three, and the directed roundings were taken for toward zero
(\code{asinpi} was wrong for about half of its arguments with clang-cl). The
field is 0 to nearest, 1 downward, 2 upward and 3 toward zero on every x86
processor, and the three functions now take it so, whatever the values of
the C library.

\medskip
CORE-MATH's \code{pow} is correctly rounded in every direction. A test of
\code{pow} where $x^y$ is a double found three defects in the earlier
versions of GAOL v5. With Visual C++,
\code{NAN} and \code{INFINITY}, which the UCRT writes for C as products of
\code{1e+300}, were computed when the program runs, in the rounding direction
in effect, so that \code{pow} rounded downward was $-0$ from $2^{-1075}$ to
about $2^{-947}$; they are read from their bits now. On 32-bit x86 the square
root of the C libraries of Windows rounds to nearest in every rounding
direction: CORE-MATH takes its square root from the SSE2 instructions there.
On 32-bit ARM, the conversion of a 54-bit integer rounded to nearest, so that
the upper bound fell below $x^y$ at a midpoint. Moreover, for $y \log x$
from 709.78263 to 709.78267, a power between $1.7975 \cdot 10^{308}$ and
$1.7976 \cdot 10^{308}$, the first phase of \code{pow} computed
$\infty - \infty$ before its second phase gave the right result, and a
program that had enabled the invalid-operation exception died on SIGFPE: the
vendored source sends those arguments to the second phase. The source of
\code{cospi} shifted a signed integer out of its range, undefined behaviour
that the tests run with UBSan reported; it shifts it unsigned now.

\medskip
The cube root \code{nth\_root(x, 3)} is CORE-MATH's \code{cbrt}, the
tightest bounds, exact where the cube root is a double, which cubing the
value tells: it takes 41 ns, where the search by bisection the other roots
use takes 89 (the medians of three runs, Intel i7-1185G7, Clang 18.1; 24 ns
against 77 on an Apple M2, Apple clang 21).\newinvfive{} \code{nth\_root(x, q)} takes
a negative $q$, which the standard recommends (\ieee{rootn} over
$\mathbb{Z}\setminus\{0\}$): the reader of strings converted the exponent to
an unsigned int, and \code{nth\_root(16, -2)} read from a string was the
4294967294-th root of 16 rather than \itv{0.25}{0.25}. The second argument of
\code{nth\_root}, \code{pown} and \code{rootn} read in a string may be any
integer a double holds: \code{nth\_root(8,3000000000)} threw
\code{invalid\_action\_error}, and \code{rootn(8,3000000000)} read with the
names of the standard was the empty set in the earlier versions of GAOL v5,
where the same call from C++ is computed. The reader of strings also takes the new names \code{exp2},
\code{log2}, \code{cbrt}, \code{sign} and \code{trunc}, which gave a syntax
error before (\pxref{sec:input-format}).

\subsection{The rounding direction and the initialization}
\label{sec:changes-rounding}
\cindex{rounding direction!differences with GAOL~4}%

The rounding direction is set upward by each operation that computes a bound
when it is not, in every build; the operations of the AVX-512 path set none,
their instructions carrying their rounding.\newinvfive{} An operation that
rounds nothing (\code{-x}, \code{abs()}, \code{min()}, \code{floor()}, the
predicates) or returns before its check (\code{y * empty},
\code{exp(empty)}) leaves the direction as it found it
(\pxref{sec:rounding-direction}). GAOL~4 has no CMake build; the one of the 4.3
releases of GAOL v5, which derives from that of IBEX, set it once, in
\code{gaol::init()}, and computed in the direction the calling code left:
rounding to nearest, downward or toward zero, 30 operations gave other
bounds, and sums and products did not enclose their exact values (3199 of the
13825 checks of the test of the rounding direction failed with GCC 9.4 on
x86-64). Built with \cmd{configure} or meson, whose option
\loption{enable-preserve-rounding} is on by default in GAOL~4 although its
manual says the contrary, GAOL~4 restored the rounding direction after each
operation, but only the x87 control word on x86-64,
where GAOL computes with SSE (43 operations left the SSE instructions
rounding to nearest or upward, whatever direction they found), and without
the rounding bits of the FPCR register on arm64.

\medskip
With the option \code{GAOL\_PRESERVE\_ROUNDING} of the three builds
(\loption{enable-\allowbreak preserve-\allowbreak rounding} of \cmd{configure} and
\option{-Denable-preserve-rounding} of meson, which GAOL~4 had),
GAOL saves and restores the
direction by reading and writing the control registers itself, a register
only if it differs from what is set back: the x87 control word and MXCSR of
an x86 processor, FPCR with the NEON intervals, FPSCR with the FPU intervals
of 32-bit ARM processors and GCC; the other processors keep
\code{fegetround()} and \code{fesetround()}. Each value computed before the
rounding direction changes is written to a \code{volatile} variable: GCC does
not implement \code{\#pragma STDC FENV\_ACCESS}, and moved the computation of
the midpoint after \code{fesetround()} (GCC bug 34678).\newinvfive{} Whether the
direction is upward is shown by an addition, $1 + 2^{-60}$ with $2^{-60}$
read from \code{volatile} memory, wherever doubles are computed in double
precision; on x86-64 processors and, with GCC and Clang, on ARM ones it is
$1 + (2^{-1060} + 0)$, which shows the modes that flush the subnormal numbers
to zero as well. The NEON intervals of 64-bit ARM processors read the
register FPCR instead, and a 32-bit x86 processor the control registers of
the two units that round. Reading the direction with the SSE register and
\code{fegetround()} rather than by an addition made \code{x + y} take 800 ns
rather than 2.8 ns under Rosetta 2, 82 to 98 ns rather than 2.6 to 4 ns with
32-bit Visual C++, and 48 ns rather than 6.4 ns with 32-bit MinGW-w64 15.2;
and Clang 18 read MXCSR once for a whole loop that changed the rounding
direction.

\medskip
The direction is checked once per operation, but for two cases told
below.\newinvfive{} The bounds of the
elementary functions at doubles are computed by functions that take the
direction to be upward already, as it is after the check of the function of
intervals calling them: in the earlier versions of GAOL v5, \code{exp()} of an
interval checked the direction three times, \code{sin()} and \code{cos()} up to four, the other elementary
functions and \code{pow(x, y)} two or three. \code{exp()} takes 5.7\,\% less
time, \code{sin()} 3.6\,\%, \code{pow(x, y)} 3.4\,\% (Clang 18, Intel
i7-1185G7). The operations calling others check once too: in those versions,
\code{tan()} checked six times, calling \code{width()}, \code{+} and \code{/} on
intervals after its own check, and \code{acos\_rel()}, \code{asin\_rel()}
and \code{atan\_rel()} six times too; they compute on doubles after their one
check, and give the same bounds, bit for bit, but for the quotients of
\code{modulo\_k\_pi()}, now computed on doubles. \code{tan()} takes 22\,\%
less time, \code{acos\_rel()} 24\,\%, \code{sqrt\_rel()} 29\,\%,
\code{pow(x, [3])} 7\,\%, \code{nth\_root(x, -2)} 13\,\%, and
\code{modulo\_k\_pi()} 10 ns rather than 22; with
\code{GAOL\_PRESERVE\_ROUNDING}, where each check saves and restores the
direction, \code{tan()} 74\,\%, \code{acos\_rel()} 83\,\%,
\code{pow(x, [-3])} 20\,\% and \code{nth\_root(x, -3)} 15\,\% (Clang 18,
Intel i7-1185G7). Still checking more than once are \code{nth\_root(x, q)}
for an order beyond 64 bits, which only an \code{\_\_int128} gives, and
\code{nth\_root(x, -2)} where the square root has a vector path of its own (the
AVX-512 path, taken with \code{GAOL\_PRESERVE\_ROUNDING}, and the NEON intervals
without it), which takes \code{sqrt()} and then \code{inverse()}, each
checking, faster than one check before the scalar root. With that option,
\code{GAOL\_RND\_KEEP()} of an interval
keeps its two bounds through two volatile doubles with the SSE2 and FPU
intervals and in the register itself with the NEON ones, where the 4.3
releases copied the interval a byte at a time into volatile memory: \code{pow(x, [-3])} and
\code{nth\_root\_rel()} take 26 to 27\,\% less time, \code{asinh\_rel()} and
\code{atanh\_rel()} 21\,\% (Clang 18, Intel i7-1185G7). \code{tan()} and
\code{pow(x, y)} keep their bounds before setting the direction back, which
they did not, so that GCC, which does not model the rounding direction,
could compute them after; and the functions of doubles \code{exp\_dn()},
\code{sin\_up()} and the thirty others saved the direction they found but
left it upward, a defect the compilers only showed by warning that a variable
was set but not used. They no longer save the direction, and their header,
\file{gaol/math/gaol\_double\_op.h}, is no longer installed.

\medskip
The initialization and the cleanup change.\newinvfive{} GAOL initializes
itself with every compiler, before the static objects of the program: each
file including \file{gaol/gaol\_common.h} holds a static
\code{gaol\_core::gaol\_initializer}, whose constructor initializes GAOL the
first time, as \code{std::ios\_base::Init} does the standard streams. GAOL~4
declared \code{gaol\_init()} with \code{\_\_attribute\_\_((constructor))},
for GCC and Clang only: Visual C++ had no automatic initialization, and the
constructor of the static library ran after those of the program (except with
MinGW-w64, which runs the constructors from the last linked to the first), so
that the static objects of the program computing intervals set the rounding
direction upward before the first \code{gaol::init()} read it, and
\code{gaol::cleanup()} left it upward. \code{gaol\_init()} is gone.
\code{gaol::cleanup()} sets back the rounding direction that the first
initialization found, unless it is preserved: GAOL~4 left it upward after its
use when a static object of the program had computed before the
initialization, even after its automatic cleanup, and a program computing to
nearest after it had to set it back itself. \code{gaol::cleanup()} does so at its
first call only, and an operation of GAOL after it sets the direction upward
again; \code{gaol::restore\_rounding()} sets it back as many times as the
program needs. Only the rounding direction is set back, not the rest of the
environment: the exception flags raised since are the program's.
\code{gaol::cleanup()} no longer frees what the initialization allocated: it
deleted the node of the empty expression while expressions still referred to
it, a local of \code{main()} destroyed after the call or a static object,
whose destructors then wrote into freed memory. The automatic cleanup frees
it when the program ends or the library is unloaded (\pxref{chap:init}).
Following the patch IBEX applies to GAOL, the FPU is initialized with
\code{fesetenv(FE\_DFL\_ENV)} and \code{round\_upward()}, rather than by
writing an x87 control word into \code{fenv\_t}, which on ARM64 left the
rounding to nearest and crashed an ARM64 Mac; and \code{interval::midpoint()}
and \code{operator<<} set the rounding direction back upward, which they did
not when the rounding direction is not preserved.

\medskip
The rounding direction is set on x86 processors by writing the control
registers of the x87 and SSE units (\code{fnstcw}/\code{fldcw},
\code{stmxcsr}/\code{ldmxcsr}) rather than through \code{fesetround()},
which cost 130 ns per call with mingw-w64 13, 50 ns with the C runtime of
Visual C++ for x64 and 250 ns for x86, 8.5 ns with glibc.\newinvfive{} GAOL~4
changed the direction four times for each elementary function of an interval,
where the math library itself took about 10 ns; GAOL v5 no longer changes it
there, CORE-MATH computing upward. On the runners of the continuous
integration, \code{exp()}, \code{log()}, \code{sin()} and \code{cos()} of an
interval took 535 to 630 ns with MinGW-w64 and MSYS2, and take 48 to 209 ns;
250 to 340 ns with Visual C++ for x64, and 45 to 160 ns; 1050 to 1750 ns with
Visual C++ for x86, and 150 to 650 ns. A register is written only if its
direction differs from the one set, and no other bit changes. With mingw-w64
13 and later there is a second reason: its \code{fesetround()} clears the six
exception flags of MXCSR when it changes the register, and its
\code{fegetround()} joins by OR the rounding fields of the two units. On a
32-bit x86 processor GAOL also reads the two registers in line, in one asm
statement, rather than through \code{fegetround()}, about 80 ns with Visual
C++ for x86. With \code{GAOL\_PRESERVE\_ROUNDING}, \code{x.midpoint()} takes
17.5 ns rather than 27.9 (GCC 15.2, Intel i7-1185G7), and 52 ns rather than
131 with Visual C++ for x64 on the runners of the continuous integration;
the bounds are the same. The option also writes the registers only if they
have to change: an operation of the SSE2 intervals read MXCSR and wrote it
twice whatever it held. In a program that computes upward, \code{x + y}
takes 1.7 ns rather than 8.2 and \code{x * y} 2.5 ns rather than 8.3; in a
program that computes to nearest, 0.2 ns more (GCC 15.2, Intel i7-1185G7, the
best of 15 runs of a loop). Visual C++ for x64 writes MXCSR alone with the
option, as GCC, Clang and clang-cl do with the fused multiply-add
instructions: in the earlier versions of GAOL v5, the operations that start
with \code{GAOL\_RND\_ENTER()} called the C runtime, because the macros asked for \code{\_\_FMA\_\_}, which
\cmd{cl} does not define; they take \code{\_\_AVX2\_\_} too now, and
\code{pow(x, 3)} takes 29 ns rather than 92, \code{sin(x)} 71 ns rather than
135 on the runners of the continuous integration (AMD EPYC, Visual Studio
2026).

\medskip
On ARM processors GAOL writes the control register itself too.\newinvfive{} With the NEON
intervals of 64-bit ARM processors the direction is the field RMode of FPCR,
read with \code{mrs} and written with \code{msr} (\pxref{sec:neon}). On
32-bit ARM processors with GCC and a floating-point unit of doubles it is the
field RMode of FPSCR, read with \code{vmrs} and written with \code{vmsr}:
\code{x + y} takes 22.5 ns rather than 49 with \code{GAOL\_PRESERVE\_ROUNDING}
(the Debian armhf jobs of the continuous integration, which run on the arm64
runners of GitHub). FPSCR holds the flags of the exceptions too, so every
write but the one at the start of an operation reads the register again just
before, and no flag raised since is lost. Without the option nothing changes
there: the check of each operation stays the addition, \code{vmrs} costing
more than its two additions (\code{x + y} 6.1 ns rather than 4.1). Elsewhere
(POWER, s390x, RISC-V, MIPS, the FPU intervals of 64-bit ARM, and 32-bit ARM
without GCC or without a floating-point unit of doubles), \code{fesetround()}
still.

\subsection{The floating-point exceptions and the flush-to-zero modes}
\label{sec:changes-exceptions}
\cindex{floating-point exceptions!differences with GAOL~4}%

Operations raise no invalid-operation exception for the empty set, nor for a
NaN the program gives, nor for a zero bound times an infinite one in an SSE2
register.\newinvfive{} The bounds of the empty set are NaN, and
\code{is\_empty()} was \code{!(left() <= right())}, as \code{set\_contains()},
\code{set\_disjoint()} and the other relations compared them with \code{<},
\code{<=}, \code{>} and \code{>=}: comparisons that signal. With the
exception enabled, \code{is\_empty()} of an empty interval and the operations
that call it stopped the program with SIGFPE, and with it masked they raised
the flag. The SSE2 product computed $0 \cdot \infty$ and cleared the NaN
afterwards, and the SSE2 \code{div\_rel} divided an infinite bound by an
infinite one in the half of a register it throws away. The bounds of an empty
operand are now compared with the quiet comparisons of \file{<cmath>}
(\code{std::islessequal()}\ldots), one instruction with GCC and Clang on
x86-64, as is the first comparison of the constructor, which gives the empty
set for NaN bounds: every operation, relation, function and output takes the
empty set without the exception, and so does a NaN the program gives
(\code{interval(NAN)}, \code{x += NAN}, \code{x <= NAN}, \code{pow(x, NAN)}).
The SSE2 divisions, \code{\%}, \code{inverse()} and \code{div\_rel()} raise
no division-by-zero exception for a divisor with a zero bound
(\code{1/[0, 1]}, \code{inverse([0, 2])}), which they raised in GAOL~4, and
\code{midpoint()} of \code{[DBL\_MAX]} raises no overflow exception, a bound
of $2^{1023}$ or more in magnitude telling the case of the sum of the halves
before the bounds are added. What stays: the inexact exception is raised by
almost every operation, and \code{log([0, 1])} and \code{[1e308] * 10}, whose
bounds are infinite, still raise the division-by-zero and the overflow ones.
A program that enables an exception has to disable it while GAOL computes
(\pxref{sec:fp-exceptions}).

\medskip
A result with no infinite bound raises no overflow nor division-by-zero
exception.\newinvfive{} Some functions computed beyond the largest double on
the way, or divided by zero, and raised \code{FE\_OVERFLOW} or
\code{FE\_DIVBYZERO} for a finite result, the empty set included:
\code{sqrt([1, DBL\_MAX])}, \code{nth\_root([DBL\_MAX], 3)} and the roots of
order 4 and more, \code{sqrt\_rel()}, \code{nth\_root\_rel()},
\code{sin([-1, DBL\_MAX])} and \code{cos()} (the width of the interval
overflowed), \code{tan([DBL\_MAX])}, \code{div\_rel()}, \code{atan\_rel()},
the reverse functions of \code{atan2} and of the power, \code{atanh\_rel()}
and, with Visual C++, \code{atan([DBL\_MAX])}. A quotient, a product or a
power that would overflow is replaced by what the operation gives, the
infinity or the largest double, without being computed; the roots of an
argument above $2^{1000}$ are proved on the argument and the root scaled by
powers of 2; \code{sin()} and \code{cos()} compare the half of the width
with $\pi$. A probe of 131 functions and forms on operands next to the limits
of the format found 34 that raised them for a finite result; none does with
the SSE2 intervals, the FPU intervals and the AVX-512 path. The bounds are the
same, bit for bit, up to the sign of a zero bound of \code{div\_rel()} with
the SSE2 intervals.

\medskip
The modes that flush subnormals to zero are cleared by each operation that
computes bounds.\newinvfive{} On x86 the flush-to-zero and denormals-are-zero
modes of MXCSR, on ARM the mode FZ of FPCR (FPSCR), make every subnormal
result a zero and read every subnormal operand as a zero. With one of them
set, the bounds of every operation with a subnormal operand or result miss
the exact result whatever the rounding direction: \code{[1e-300] * [1e-20]}
is \code{[0, 0]}. A program linked with \option{-Ofast}, \option{-ffast-math}
or \option{-funsafe-math-optimizations} has them (\file{crtfastmath.o}), and
so has one that loads a plug-in built so. Each operation that computes
bounds computes $1 + (2^{-1060} + 0)$ (x86-64, and ARM with GCC and Clang),
which shows the modes as well as the direction, and clears them with the
direction; a 32-bit x86 processor and the NEON intervals of 64-bit ARM
processors read the control registers instead, which hold both. With
\code{GAOL\_PRESERVE\_ROUNDING} they are set back as they were, and the
result, maxima and minima of bounds included, is made before. The check comes
before the operation reads, compares or copies a bound, and the compiler is
kept from moving a comparison before it: in the earlier versions of GAOL v5,
with GCC 13 at \option{-O3}, \code{asinpi([100*2\textasciicircum{}-1074])} was
above the exact value, and \code{sqrt([-1e-310, 4])} and \code{atan2} gave the
empty set. At the link, \file{gaol.pc} and
\code{gaol::gaol} give \option{-mno-daz-ftz} to the code using GAOL where the
compiler accepts it (GCC 13 and later on x86, from 11.4 and 12.4 in the series
11 and 12, and Clang 20 and 21, not Clang 18), which keeps \file{crtfastmath.o} out of the program; the check
of the operations is the only defence against a plug-in, and the only one on
the other compilers; on POWER, s390x, RISC-V and MIPS neither defence exists. The cost
is that of the second addition of the check: \code{x * y}, \code{sqrt} and
\code{pow(x, 3)} take 0.6 to 0.7 ns more than with the direction alone (Intel
Xeon Cascade Lake, Clang 18: 4.0 ns rather than 3.4, 11.0 rather than 10.3,
13.7 rather than 13.0). The check of the rounding direction alone,
$1 + 2^{-60}$, is the one left where GAOL cannot clear the modes (POWER,
s390x, RISC-V, MIPS, and the FPU intervals of Visual C++ for ARM)
(\pxref{sec:ftz}).

\medskip
The functions that only read, compare or copy bounds give right results under
these modes. With denormals-are-zero a comparison reads a subnormal as a
zero, and the functions that make no check gave wrong results:
\code{interval(2\textasciicircum{}-1073, 2\textasciicircum{}-1074)} kept its bounds in the wrong order,
\code{abs([-1e-309, -1e-310])} stayed negative, \code{floor([-1e-310])} was
\code{[-0]}, the midpoint of $[2, 4] \cdot 2^{-1074}$ was 0 and
\code{split()} made an empty half of it. The constructor from two bounds, the
relations, \code{\&}, \code{|}, \code{max}, \code{min}, \code{abs},
\code{sign}, the roundings to an integer, \code{mig}, \code{mag},
\code{is\_canonical}, \code{feven}, \code{nb\_fp\_numbers}, and \code{inf}
and \code{sup} of \code{gaol\_ieee1788} now compare the bounds as doubles
first and, where these are equal, by their bits, which no mode changes, and
round a subnormal bound from its bits. \code{bisect(ratio)} and
\code{is\_bisectable()} compare the bounds by their bits too, the cut point
included: $[3 \cdot 2^{-1074}, 5 \cdot 2^{-1074}]$ was not bisectable. The
AVX-512 path reads its operands once the modes are cleared, and makes its
results before they are set back: Clang 18 computed \code{x - y} after, and
\code{[3e-308] - [2.9e-308]} was \code{[0, 0]} (\pxref{sec:avx512}).

\subsection{The constants and the static objects}
\label{sec:changes-constants}
\cindex{constants!differences with GAOL~4}%

The constants of GAOL are initialized when compiling, so that the intervals
computed in the initialization of the static objects of a program are
right.\newinvfive{} They were computed by the dynamic initialization of
GAOL's files, which, with the static library of the CMake build, comes after
the static objects of the program, except with MinGW-w64: an interval
computed there found them 0, and \code{interval::pi()} was $[-0, 0]$,
$[1, 2]/[-1, 1]$ was $[-0, 0]$ rather than $[-\infty, +\infty]$ and
\code{sin([0, 4])} was $[-0, 0]$, without any warning (25 of the 45 checks of
the test of the static initialization failed with the SSE2 intervals, 24 with
the others). The doubles around $\pi$ and $\pi/2$, \code{gaol\_pi\_dn},
\code{gaol\_pi\_up}, \code{gaol\_half\_pi\_dn} and \code{gaol\_half\_pi\_up},
are written in decimal, exactly, rather than read from the unions that give
their bits, which is not a constant expression; they are in \code{gaol\_core}
(GAOL~4 had them in \code{gaol} as \code{pi\_dn}\ldots), and \code{gaol} and
\code{gaol\_ieee1788} name them too. \code{interval::one()}, \code{pi()},
\code{universe()} and the other constants build their interval from these
doubles, which the compiler folds, rather than copy the static intervals
\code{interval::cst\_one}, \code{cst\_pi}\ldots, which are gone
(\pxref{sec:interval-constants}). The masks of the SSE2 intervals are written
as vectors, which the compilers put in read-only memory (GCC and Clang) or in
initialized data (Visual C++), rather than computed with \code{\_mm\_set\_pd()}
and \code{\_mm\_cmpeq\_pd()}.

\medskip
\code{GAOL\_INFINITY} is \code{std::numeric\_limits<double>::infinity()}, a
constant of the compiler, where GAOL~4 took the \code{HUGE\_VAL} of the C
library.\newinvfive{} The UCRT of Windows (SDK 10.0.26100) writes
\code{HUGE\_VAL} as \code{((double)(float)1e+300)}: Visual C++ makes the
conversion when compiling, but clang-cl under \option{/fp:strict} makes it
when the program runs, in the rounding direction of the moment, and got
\code{FLT\_MAX} when that direction was downward or toward zero, with the
overflow flag. \code{interval()} was then $[-\mathrm{FLT\_MAX},
\mathrm{FLT\_MAX}]$, \code{interval(1e300)} was empty and $[1]/[-1, 1]$ did
not contain $10^{300}$ (\pxref{sec:floating-point-constants}).

\subsection{The kernels: SSE2, NEON, AVX-512 and the FPU}
\label{sec:changes-kernels}
\cindex{kernels!differences with GAOL~4}%

On 64-bit ARM processors the intervals are computed with NEON, with every
compiler of them: Visual C++, clang-cl and the Clang of MinGW-w64 and MSYS2 as
well as GCC, Clang and AppleClang under Linux and macOS.\newinvfive{} GAOL~4
computed the intervals of every ARM processor with the FPU, and 32-bit ARM,
whose NEON has no vector of doubles, still does. The two bounds are the two
lanes of a \code{float64x2\_t}, as the SSE2 intervals keep them in a
\code{\_\_m128d}, and \code{+}, \code{-}, \code{*}, \code{/}, \code{sqr()}
and \code{sqrt()} are kernels that take and return them in registers. On an
Apple M2 (AppleClang 21.0), \code{x + y} takes 0.90 ns rather than 1.15,
\code{x * y} 1.52 ns rather than 3.26, \code{x / y} 2.08 ns rather than 3.22,
\code{sqrt(x)} 2.04 ns rather than 5.17 and a line of 13 operations 16.6 ns
rather than 50.6, the second figures being those of the FPU intervals. The
bounds are the same, but for the sign of a zero bound, which the maximum of
64-bit ARM, ranking $+0$ above $-0$, can set otherwise: the lower bound of
$[-1, 0] \cdot [-1, 0]$ is $-0$ with NEON. Each operation checks the rounding
direction with one \code{mrs} (0.8 ns on the M2), which raises no exception
flag; with \code{GAOL\_PRESERVE\_ROUNDING}, an operation takes 27 ns rather
than 37 (\pxref{sec:neon}).

\medskip
The AVX-512 path carries the rounding direction in the
instruction.\newinvfive{} \code{x + y}, \code{x - y}, \code{x * y},
\code{x / y}, \code{sqrt(x)}, \code{sqr(x)}, the operators of an interval and
a double, and the integer powers beyond 1 and $-1$ are computed with the
rounding that the AVX-512 instructions carry in themselves, on x86-64 where
the compiler takes them, never GCC for Windows, which does not align the
stack for the AVX-512 registers (GCC bug 54412). They set no rounding
direction and raise no exception flag. The bounds are those of the SSE2
intervals bit for bit, but for the sign of a zero bound of \code{x * d}
where the zero bounds of the interval differ in sign. The three builds take
the path where \code{GAOL\_PRESERVE\_ROUNDING} is on and the machine building
runs the instructions, where it computes 2.9 times faster at the median on an
Intel i7-1185G7, and leave it out without the option, where the operations take
20\,\% longer;
the library then needs a processor with them, and the initialization stops a
program run on another one with a message (\pxref{sec:avx512}). Visual C++
and clang-cl for x64 take the SSE2 intervals and the AVX-512 path, as GCC and
Clang do, the kernels being \code{\_\_vectorcall} there so that their
\code{\_\_m128d} arguments go in registers.

\medskip
\code{*}, \code{/} and \code{sqr()} of the SSE2 intervals compute with no
branch on the signs of the bounds, and the basic operations hand the bounds
over in registers.\newinvfive{} GAOL~4 chose the product, the quotient and
the square by \code{switch}es and nested branches on the signs of the bounds,
and a branch predicted wrong costs more than the operation it chooses; its
operators read their operands and wrote their result through memory.
\code{x * y} is now the larger of four products, each rounded upward, in the
two lanes of an SSE2 register; \code{x / y} one division of the two lanes by
the denominators that masks drawn from the signs choose; \code{sqr(x)} two
products and a mask. Only an empty operand, an infinite bound or a divisor
holding 0 takes a branch. On an Intel i7-1185G7 with Clang 21, \code{x * y}
takes 2.3 ns rather than 10.0 with GAOL 4.2.3, \code{x / y} 2.5 ns rather
than 7.5, and Shekel 5 108 ns rather than 189 (\pxref{sec:comparison}). The
sign of a zero bound of a product or of a quotient may differ from GAOL~4's,
the value never: \code{maxpd} chooses between a $+0$ and a $-0$ that compare
equal, where GAOL~4 chose one product by the signs; $[-1, 0] \cdot [-1, 0]$
has the lower bound $+0$, where GAOL~4 had $-0$. The FPU intervals multiply,
divide and square without a branch too: GAOL~4 took one of nine cases of the
signs for a product, and went down a chain of up to eight tests for a
quotient. On an Apple M2 (AppleClang 21, \code{GAOL\_SIMD=OFF}), over the
random intervals of the benchmark, \code{x * y} takes 3.5 ns rather than
13.9, \code{x / y} 3.2 rather than 8.5, \code{sqr(x)} 1.8 rather than 6.9,
and Shekel 5 138 ns rather than 298. A line of products of operands of one
sign, which a processor predicts, takes more time (seven \code{*} in a row:
35.0 ns rather than 23.3, and 1.7 times longer with GCC 13 on x86-64 in the
continuous integration), and with \code{GAOL\_PRESERVE\_ROUNDING} the kernels
lose 7\,\% on operands of one sign (\code{x * y} 42.4 ns rather than 39.7) and
gain 11\,\% on operands of random signs (42.4 rather than 47.5). The sign of a zero
bound of such a product is not specified: GCC 13 for POWER9 compiles the
maximum as an \code{xsmaxcdp}, which returns its first operand. \code{x * y}
forms every corner product, so that a subnormal corner that is no bound of
the result raises \code{FE\_UNDERFLOW}, and a corner that overflows
\code{FE\_OVERFLOW} (\pxref{sec:kernels}).

\medskip
Several operations of the kernels were corrected or sped up.\newinvfive{}
\code{x -= x}, \code{x /= x} and \code{x \%= x} enclose their results in the
FPU intervals: GAOL~4 wrote one bound before reading it, and \code{x -= x}
was $[a - b, 2b - a]$ for $x = [a, b]$, \code{[0.1] /= [0.1]} was empty, and
normalizing a row by its pivot, \code{row[j] /= row[i]} down to $j = i$, did
it. A zero bound times an infinite bound counts as 0 in the products of the
FPU intervals, as in those of the SSE2 intervals: $[0, 1] \cdot [1, +\infty]$
is $[0, +\infty]$ in every build, where the FPU intervals gave the empty set.
The operators of an SSE2 interval with a double compute by the sign of the
double rather than through \code{interval(d)}: on an Intel i7-1185G7 (GCC
9.4, CMake Release), \code{x += d} takes 3.9 ns rather than 13.6,
\code{x *= d} 6.3 ns rather than 25.2 and \code{x /= d} 4.1 ns rather than
18.0, giving the same sets. \code{max(x, y)} and \code{min(x, y)} of the SSE2
and NEON intervals take both bounds of the result from one comparison of the
stored bounds: on an Intel i7-1185G7, \code{max(x, y)} takes 2.3 ns rather
than 8.5 with GCC 15.2 and 2.0 ns rather than 5.6 with Clang 21; on 64-bit
ARM they are 1.5 to 2.4 times faster (a Neoverse N2, an Apple M2). The
constructors from a register, \code{interval(const \_\_m128d\&)} and
\code{interval(const float64x2\_t\&)}, give the empty set when a lane is
$-\infty$, which is $l = +\infty$ or $u = -\infty$, the lower bound being
stored negated; both are explicit now.
\code{minimum()} and \code{maximum()} of two zeros give $-0$ and $+0$ with
every compiler: Visual C++ compiled \code{(b <= a) ? b : a} as \code{minsd},
which gives its second operand for two zeros.

\medskip
Square roots are the tightest intervals, as IEEE 1788-2015 requires of its
basic operations: the lower bound is the square root rounded upward when its
square is the argument, the double below it otherwise.\newinvfive{} GAOL~4
computed it as $x/\sqrt{x}$ rounded downward, one double below the tightest
bound for half of the doubles; \code{sqrt()} takes the same time as before
(9.0 ns on an Intel i7-1185G7, GCC 9.4), and is bounded whatever the rounding
of the C library's \code{sqrt}, which Visual C++ for 32-bit x86 rounds to
nearest in every rounding direction. \code{sqrt()} is no longer slowed down
by GCC on Intel processors: GAOL took the smallest subnormal from
\code{std::numeric\_limits<double>::denorm\_min()}, a \code{long double}
constant whose conversion GCC does not fold with \option{-frounding-math},
and the x87 unit converted it with a microcode assist at each use, about
100 ns on an i7-1185G7. \code{sqrt()} took 124 ns rather than 7.8, and
\code{nth\_root(x, 2)} and \code{nth\_root(x, -2)} 113 and 118 ns rather than
6.8 and 10.2 (GCC 15.2, Intel i7-1185G7); the constant is a literal now, and
the bounds are the same, bit for bit. \option{-ffloat-store} is added only
where doubles are still computed on the x87 unit: CMake gave it to GCC on
every target, and GCC then stored every double variable in memory. On an
Intel i7-1185G7 (GCC 9.4, CMake Release), \code{x + y} takes 3.0 ns rather
than 9.1, \code{x * y} 4.9 ns rather than 18.3, \code{sqrt(x)} 8.3 ns rather
than 41.3, and \code{log(x)} 73 ns rather than 110.

\medskip
The SSE2 intervals are computed on 32-bit Windows with GCC, Clang, Visual C++
and clang-cl.\newinvfive{} An SSE2 interval holds a \code{\_\_m128d} and has
to be aligned on 16 bytes, and the C runtime of 32-bit Windows aligns the
memory of \code{malloc()} and \code{new} on 8: a \code{std::vector} of
intervals crashed on its first \code{movaps} (MinGW-w64 14.2 and 15.2 for
x86). GCC and Clang compile the code using GAOL with
\option{-faligned-new=8}, which the three builds give to the library and to
the code that uses it, and \file{gaol/gaol\_config.h} stops the compilation of
the SSE2 intervals on 32-bit Windows without it. Visual C++ and clang-cl have
no such flag: their SSE2 intervals keep the two bounds as two doubles of
alignment 8, loaded with \code{movupd}, so that a program gives no flag. The
memory of the SSE2 intervals is released with the function matching the one
that allocated it (\code{GAOL\_MEMFREE()}): \code{\_aligned\_malloc()} and
\code{\_aligned\_free()} on Windows, so that the SSE2 intervals can be used
with MinGW-w64 and MSYS2 on x64. The placement \code{delete} of
\code{interval}, called when a constructor throws, freed the caller's memory;
it now leaves it alone.

\subsection[The constructors, the empty set and the interval functions]{The constructors, the empty set and the functions of intervals}
\label{sec:changes-intervals}
\cindex{constructors!differences with GAOL~4}%
\cindex{empty set!differences with GAOL~4}%

The constructors give the empty set where their arguments are not an
interval, as IBEX does: \code{interval(+oo)}, \code{interval(-oo)},
\code{interval(+oo, +oo)}, \code{interval(-oo, -oo)}, \code{interval(2, 1)}
and bounds that are NaN are all the empty set, and assigning an infinity to an
interval empties it.\newinvfive{} IEEE 1788-2015 has no interval
$[+\infty, +\infty]$ nor $[-\infty, -\infty]$: its constructor fails there,
and gives the empty set (10.5.8, 12.12.7). GAOL~4 built them as intervals of
one infinite point, whose midpoint, \code{realmax}, was outside them. The
intervals with one infinite bound are unchanged. On an Intel i7-1185G7 (GCC
9.4, CMake Release), \code{interval(a, b)} takes 1.28 ns rather than 1.24 ns:
its bounds are set in a register and behind a branch, since a pair written to
memory or conditional moves cost several nanoseconds per construction
(\pxref{sec:constructors}). The constructor from one \code{double},
\code{interval(double)}, is explicit, and so are \code{expression(double)}
and \code{expression(const interval\&)}: \code{interval x = 0.5;},
\code{return 0.5;} in a function returning an interval, \code{min(x, 0.5)}
and \code{interval r = sqrt(2.0);} do not compile, where GAOL~4 took the
double for a point interval without saying so, the \code{sqrt(2.0)} of the C
library among them. The assignment \code{x = d}, \code{x \&= d} and
\code{x |= d}, the relations between an interval and a \code{double}, the
arithmetic operators of an interval and a \code{double} and
\code{set\_contains(d)} have overloads of their own, which take
\code{interval(d)} as the conversion did: \code{x = 0.5} and \code{x < 0.0}
still compile, and give the same results. \code{interval(double, double)} is
not explicit.

\medskip
Integers are taken as the integers they are.\newinvfive{} An integer given
where GAOL takes a \code{double} was converted to a \code{double}, in the
rounding direction in effect, which an integer beyond $2^{53}$ is not:
\code{interval(9007199254740993LL)}, $2^{53} + 1$, was $[2^{53}, 2^{53}]$ or
$[2^{53} + 2, 2^{53} + 2]$, neither containing it. An integer of any type
but \code{bool} is now taken as the integer it is by the constructors, the
assignments, the arithmetic operators, the relations, \code{set\_contains()},
\code{expression(n)}, and \code{numsToInterval()} and \code{isMember()} of
\code{gaol\_ieee1788}: \code{interval(n)} is the tightest interval
containing $n$, the two doubles around it where $n$ is no double. The
arithmetic operators enclose the exact result, but not always with the
tightest bounds, for an $n$ that is no double: $[1, 2] + (2^{53} + 1)$ is
$[2^{53}, 2^{53} + 4]$. The exponents of \code{pow} and \code{pown}, and the
orders of \code{nth\_root}, \code{rootn} and \code{pownRev}, take an integer
of any type too: GAOL~4 had \code{pow(x, int)} only, so that
\code{pow(x, 3000000000u)} was \code{pow(x, -1294967296)}, and converted the
order of \code{nth\_root} to an \code{unsigned}, so that
\code{nth\_root(e, -2)} of an expression was the order 4294967294.

\medskip
The constructors from strings are gone.\newinvfive{} \code{interval(const
char*)} took any \code{const char*} and converted it implicitly:
\code{interval(0)} was ambiguous, 0 being a null pointer as well as a number,
and \code{x + nullptr} compiled, and crashed. \code{gaol::textToInterval()}
reads the string instead, with the same names of functions, and throws
\code{std::invalid\_argument} for a string that is no interval, and for a
null pointer, through an overload for \code{const char*}, where the overload
for \code{std::string} would build a \code{std::string} from the pointer,
which the standard leaves undefined. \code{interval(const char*, const
char*)} kept the bounds $[2, 1]$ of \code{interval("2", "1")}, which
\code{is\_empty()} called empty but whose sum with $[0, 1]$ was $[2, 2]$;
\code{textToInterval("0.1", "0.3")} takes the left bound of the first string
and the right bound of the second, and gives the empty set for bounds in the
wrong order. A result thrown away gives a warning: \code{sqrt(x);} alone
leaves \code{x} as it was, and the functions whose result is all they do
carry \code{GAOL\_NODISCARD}, \code{[[nodiscard]]} in C++17 and later and
the attribute of the compiler before; a cast to void silences it, and a
program that does not want it defines \code{GAOL\_NODISCARD} empty
(\pxref{sec:nodiscard}).

\medskip
The functions of intervals change on the empty set and on their
rounding.\newinvfive{} The intersection of disjoint intervals is the empty
set of \code{interval::emptyset()}, whose bounds are NaN: GAOL~4 kept the
largest left bound and the smallest right bound, $[3, 2]$ for
$[1, 2] \cap [3, 4]$, which \code{is\_empty()} took for empty but the
operations on the bounds did not, and $([1, 2] \cap [3, 4]) + [0, 1]$ was
$[3, 3]$. \code{width()} of the empty set is NaN, as \ieee{wid} of IEEE
1788-2015, rather than $-1$, and \code{chi()} of the empty set is the
positive NaN, printed \code{nan} with every compiler. \code{hausdorff()}
returns the tightest upper bound of the distance, and 0 for equal bounds,
infinite ones included: GAOL~4 computed \code{fabs(a - c)} in the rounding
direction of the caller and returned $+\infty$ as soon as a bound was
infinite, so that \code{hausdorff([1, +oo], [1, +oo])} was $+\infty$ and a
fixed-point loop on a box with an unbounded side never stopped. An infinite
bound in one interval only still puts them at $+\infty$, as in GAOL~4
(\code{hausdorff([1, +oo], [1, 2])}), and the function then returns before the
rounding direction is checked, 2.5 to 4.4 times as fast as
through the computation of the distance (Intel i7-1185G7, GCC 9.4 and Clang
18). \code{nb\_fp\_numbers()} numbers the doubles by the bits of the
magnitudes, $-0$ and $+0$ being one number: \code{nb\_fp\_numbers(-0.0, 1.0)}
was 13830554455654793217, and is 4607182418800017409, the number of the
doubles from 0 to 1. \code{mid()} returns the tightest interval enclosing the
midpoint, computed as \code{midpoint()} computes it with the rounding
outward: it computed $a/2 + b/2$ rounded upward and subtracted the half of
$b$ for its lower bound, which is not exact when $|b| < 2^{-1021}$, so that
\code{mid([0, 2\textasciicircum{}-1074])} was $[2^{-1074}, 2^{-1074}]$ and 45150 of the
180901 intervals whose bounds are multiples of $2^{-1074}$ from $-300$ to
$300$ times it did not enclose their midpoints. \code{rad()} and
\code{mid\_rad()}, \ieee{rad} and \ieee{midRad} of the standard, are new
(Chapter~\ref{chap:functions}).

\medskip
The tools that algorithms of intervals write for themselves are
new:\newinvfive{} \code{width\_enclosure()}, \code{inflate(r)},\hfill\break
\code{interval::centered\_form(m, r)}, \code{bisect(ratio)}, \code{is\_bisectable()},
and \code{hull(a, b)} and \code{intersect(a, b)}, the functions of
\code{a | b} and \code{a \& b}. The first three compute with the operations
of intervals, each bound rounded outward: \code{interval::centered\_form(m, r)}
encloses $[m - r, m + r]$, which \code{interval(m - r, m + r)} does not, both
of its bounds being computed upward; \code{width\_enclosure()} is the
enclosure of the exact width that \code{width()} bounds from above.
\code{bisect(ratio)} cuts at $l(1 - \mathit{ratio}) + r \cdot \mathit{ratio}$,
which does not overflow (for a bound of $2^{1023}$ or more in magnitude it is
computed from the halves of the bounds) and lies strictly between the bounds when
\code{is\_bisectable()}, and throws \code{std::domain\_error} for a ratio
outside $(0, 1)$ or a NaN; \code{is\_bisectable()} is true when a double
lies strictly between the bounds (\pxref{sec:tools}).

\subsection{The powers and the roots}
\label{sec:changes-powers}
\cindex{powers!differences with GAOL~4}%

The powers with a real exponent were corrected, after the fix of
Codac.\newinvfive{} \code{pow(I, e)} for a floating-point \var{e} called
\code{pow(I, int)}, which truncated the exponent: \code{pow([4], 0.5)}
returned $[1]$. It now computes an integer \var{e} with \code{pow(I, int)}
and any other \var{e} with \code{pow(I, J)}. \code{pow(I, J)} computed the
powers of the negative part of \var{I} on its magnitude: \code{pow([-4,-1],
[0.5])} returned $[-1, 2]$. It is now hybrid: a degenerate integer exponent
$[n]$ takes \code{pow(I, int)}, the \ieee{pown} of the standard, which keeps
the negative part of \var{I} and gives $0^0 = 1$; any other exponent takes
the \ieee{pow} of the standard, on the part of \var{I} in $[0, +\infty]$,
where $0^y$ is 0 for $y > 0$ and has no value for $y \leq 0$:
\code{pow([0], [0.5])} was $[0, 4.9 \cdot 10^{-324}]$ and
\code{pow([0], [-0.5])} was $[\mathrm{MAX}, +\infty]$
(\pxref{sec:pow-namespaces}).

\medskip
\code{pow(I, J)} takes the tightest bounds at the corners of the
box.\newinvfive{} $x^y$ increases with $y$ for $x > 1$ and decreases for
$x < 1$, increases with $x$ for $y > 0$ and decreases for $y < 0$: over a box
of bases above 0 its extrema are at corners, and the pow of CORE-MATH,
correctly rounded upward, is taken there, as the upper bound, and, as the
lower bound, the double below it, or $x^y$ itself where that is a double,
which is proved from the bits of $x$ and $y$: \code{pow([4], 0.5)} is $[2]$,
and \code{pow([3, 4], [2, 3])} is $[9, 64]$. At a corner where $x$ is 0 or
$+\infty$, or $y$ infinite, the bound is the limit of $x^y$, and CORE-MATH
is not called: \code{pow([4, +oo], 0.5)} is $[2, +\infty]$. GAOL~4 computed
$\exp(J \log I)$, which multiplied the relative width of $\log I$, a few
$2^{-52}$, by $|y \log x|$: \code{pow([2], [1023.5])} was 1425 doubles below
the exact value and 748 above, and \code{pow([0, 2\textasciicircum{}-1000], [-oo, -1])}
had a lower bound 617 doubles below $2^{1000}$. Over the million powers of
the benchmark of \file{doc/compare}, its intervals are no wider than those of
libieeep1788, which computes each bound with MPFR, where $\exp(J \log I)$
made them $5.9 \cdot 10^{-14}$ wider relatively. With Clang 18 on an Intel
i7-1185G7, over 100\,000 boxes, \code{gaol::pow(x, J)} with $x = [1/b, b]$
takes 84 ns (GAOL~4: 138 ns), with $x = [a, +\infty]$ 55 ns (GAOL~4: 99 ns)
and with $J = [a, +\infty]$ 54 ns (GAOL~4: 126 ns).

\medskip
Integer powers are the tightest bounds, of either sign.\newinvfive{}
\code{pow(x, n)} of \code{gaol}, \code{pown(x, n)} of \code{gaol\_ieee1788}
and the integer exponents of \code{pow(x, y)} and \code{pow(x, p)} give the
tightest bounds of $x^n$, or one double beyond where $x^n$ is within twice
the error bound of a double ($2^{-98}$ relatively at $n = 3$ and about
$2^{-62}$ at the most; at $n = 3$, one bound in $2^{46}$ is one double beyond
for an $x$ at random), and exact where $x^n$ is a double: \code{pow([2, 3],
3)} is $[8, 27]$. GAOL~4 rounded each product of its binary exponentiation
outward and was up to $n + 1$ doubles from the tightest bounds, 2, 3, 5, 6
and 8 doubles for $n = 3$ to 7, and took $x^{-n}$ for $1/x^n$, two roundings.
Both bounds are computed at once, in the two lanes of an SSE2 or a NEON
register with the FMA instructions, or in a pair of doubles otherwise: each
product is rounded, \code{fma} gives its rest exactly, and the power is
carried as a sum of two doubles whose error is bounded beforehand, so that
each bound is rounded once. A corner whose powers leave
$[2^{-968}, 2^{960}]$ is computed alone, scaled by powers of two: GAOL~4's
$1/x^n$ overflowed before the inversion, and \code{pow([10], -400)} was
$[0, 5.6 \cdot 10^{-309}]$ rather than $[0, 2^{-1074}]$. \code{pow(x, 3)}
takes 5.1 ns against 9.1 with GAOL 4.2.3 (Intel i7-1185G7, Clang 21). The
exponent 2 is \code{sqr(x)} in every form, the squares of the bounds rounded
outward: in the earlier versions of GAOL v5, with the exact products,
\code{pow(x, 2)} took 5.3 ns with the NEON intervals and 5.2 with the FPU intervals, and it takes 1.6 and 1.9 now (Apple
M2, AppleClang 21); the 6 million calls of \code{pow(x, 2)} of the
Goldstein-Price example take 50 ms rather than 77 with the NEON intervals.
The integer powers beyond the ints are computed from the pow of CORE-MATH at
the corners of $x$, the sign from the parity of $n$: the tightest where $n$
is a double, and within three doubles of it beyond $2^{53}$. \code{pow(x, 2)} in all its
forms, the powers with a real or an interval exponent (an empty exponent
included), with a \code{long long} beyond the \code{unsigned} ints, and
\code{gaol\_ieee1788::pow} return the empty set for an empty
operand before they check the rounding direction, as \code{sqr} does, and leave
the direction and the flush modes as the program had them; the other integer
powers, of an \code{int}, an \code{unsigned} or a \code{long long} within their
range, check it first and leave it upward (\pxref{sec:powers}).

\medskip
\code{nth\_root(I, n)} is the \ieee{rootn} of IEEE 1788-2015: for an odd
$n$, it is defined on the whole real line, the root of a negative number
being the opposite of the root of its magnitude, and \code{nth\_root([-8, 27],
3)} encloses $[-2, 3]$; for an even $n$, it takes the roots of the part of
\var{I} in $[0, +\infty]$, as before.\newinvfive{} GAOL~4 took that part for
every $n$, and gave $[-2^{-1074}, 3.000000000000001]$, or the empty set for
\code{nth\_root([-8, -1], 3)}; the root of 0 is 0, rather than
$[-2^{-1074}, 2^{-1074}]$. The $n$th roots are proved with integer powers,
and are the tightest bounds or up to 2 doubles beyond (4 for \ieee{rootn}
with $q < 0$, the inverse of a root), for every $n$ and every double, the
subnormal ones included. GAOL~4 took the power of the mathematical library
with the exponent $1/n$ rounded, moved one double outward: the bounds were up
to 8 doubles from the tightest between $2^{-30}$ and $2^{30}$, and 234 over
all the doubles, and \code{nth\_root([27], 3)} was two doubles wide where it
is $[3]$ now. A bound $l$ is below the root when $l^n$, rounded upward, is at
most $x$, and $u$ above it when $u^n$, rounded downward, is at least $x$; the
search starts from the pow of CORE-MATH brought next to the root by a step of
Newton's method, and goes by steps that double, then by bisection. For an
$x$ below $2^{-1020}$ the proofs take scaled powers, the plain proof losing
bits where a power is subnormal: \code{nth\_root([2\textasciicircum{}-1074], 33)} was 1983
doubles below the tightest lower bound and 28785 above the tightest upper
bound. \code{nth\_root\_rel()} takes the roots of \code{nth\_root()}: it kept
a value within 23 doubles, now 2. A point interval takes 152 ns rather than
130, its only root being looked for once; an interval 194 ns rather than 136
(Intel i7-1185G7, GCC 9.4). \code{pownRev(c, x, p)} for $p \leq 0$, the last
reverse operation of the standard GAOL lacked, is \code{nth\_root\_rel(J, p,
I)}, which takes any integer $p$ now: a negative \code{int} was converted to
a large \code{unsigned int}, a root of order about $4 \cdot 10^9$. For
$p < 0$ the bounds are up to 4 doubles beyond the tightest ones.
\code{uipow()}, the \ieee{pown} for an unsigned exponent, is \code{pow(I,
n)} for an \code{unsigned} $n$ in \code{gaol}, exported with the SSE2
intervals too, where it was inline and did not link; it gives $[1]$ for an
exponent 0, and the empty set for an empty interval.

\subsection{The trigonometric, hyperbolic and other functions}
\label{sec:changes-functions}
\cindex{elementary functions!differences with GAOL~4}%

\code{sin()} is computed as \code{cos()} is: the bounds of the interval
divided by an enclosure of $\pi$, minus $1/2$, tell the pieces where the sine
is monotonic, and the sine of CORE-MATH is taken at the bounds.\newinvfive{}
GAOL~4 computed \code{cos(x - [pi/2])}, whose subtraction widened the
argument by about $2^{-52} \max(2, |x|)$: \code{sin([1e-10])} was
$4.4 \cdot 10^{-16}$ wide, billions of doubles, and \code{sin()} was not
accurate in the sense of IEEE 1788-2015. On an Intel i7-1185G7 (GCC 9.4),
\code{sin()} takes 121 ns rather than 130. \code{sin()}, \code{cos()} and
\code{tan()} are within one double of the tightest bounds at every
magnitude. The bounds divided by an enclosure of $\pi$ only tell on which
pieces the interval lies while the quotients are not within their rounding
errors of an integer: a bound within about $|x| \cdot 2^{-52}$ of an extremum
was taken as reaching it, \code{cos([2\textasciicircum{}60])} was $[-1, 1]$, and
\code{tan()} gave $[-\infty, +\infty]$ next to a pole and beyond $2^{52}$.
CORE-MATH being correctly rounded, and no double but 0 being a multiple of
$\pi/2$, the signs of its sine and cosine at the bounds are the exact ones,
and tell whether an extremum or a pole is within the interval; they are only
asked for next to an extremum or a pole, and at the large magnitudes. The
time is unchanged in the benchmark (113 ns for \code{sin()} and \code{cos()},
Intel i7-1185G7, GCC 9.4), 125 ns rather than 121 for intervals holding an
extremum; at magnitudes from $10^8$ to $10^{15}$, \code{cos()} of a narrow
interval takes 250 ns rather than 224, and \code{tan()} 260 rather than 217.
\code{tan()} is finite over an interval of doubles that holds no pole and
whose width, rounded upward, is at most the double below $\pi$: \code{tan([-M\_PI\_2, M\_PI\_2])} was
$[-\infty, +\infty]$, and is about $\pm 1.63 \cdot 10^{16}$
(Chapter~\ref{chap:accuracy}).

\medskip
\code{atan2(y, x)} is implemented, as the \ieee{atan2} of the standard,
defined on the plane but $(0, 0)$ with values in $(-\pi, \pi]$: GAOL~4
declared it, and its parser read it, but it raised
\code{unavailable\_feature\_error}.\newinvfive{} In each quadrant the angle
is monotonic in $y$ and in $x$, and the atan2 of CORE-MATH is taken at two
corners of the box; a box with points on the half-line $y = 0$, $x < 0$ and
points below it gives $[-\pi, \pi]$, and \code{atan2([0], [0])} is empty. It
takes 100 ns on an Intel i7-1185G7 (GCC 9.4) for boxes within a quadrant,
against 68 ns for \code{atan()}. \code{acos\_rel()}, \code{asin\_rel()} and
\code{atan\_rel()} enclose the multiples of $\pi$ of the pieces of the
preimage with $\pi$ in double-double, one rounding at the magnitude of the
result rather than two: \code{atan\_rel()} was up to 49 doubles from the
tightest bounds over random cases where it is now within 2, and
\code{acos\_rel()} up to 34. An \var{I} of a single double is decided at
every magnitude by whether $f(I)$ meets \var{J}, where a bound beyond
$2^{52}$ was kept as it was (\pxref{sec:relational-arithmetic}).

\medskip
The elementary functions are exact where their value is a double, and the
tightest where it is $\pi/4$, $\pi/2$ or $\pi$.\newinvfive{} $\exp(0) = 1$
and $\log(1) = 0$ exactly, the bounds moved one double outward giving
\code{exp([0])} and \code{log([1])} a width: \code{pow([1], [-oo, +oo])} was
$[0, +\infty]$ rather than $[1]$. \code{log()} takes 1.5 ns more (72.5 ns
rather than 71) and \code{pow(x, y)} 2 ns more (145 ns rather than 143) on
an Intel i7-1185G7 (GCC 9.4). \code{sin(0)}, \code{tan(0)}, \code{asin(0)},
\code{atan(0)}, \code{sinh(0)}, \code{tanh(0)}, \code{asinh(0)} and
\code{atanh(0)} are 0, \code{cos(0)} and \code{cosh(0)} 1, \code{acos(1)}
and \code{acosh(1)} 0, and \code{asin(}$\pm$\code{1)}, \code{acos(0)},
\code{acos(-1)}, \code{atan(}$\pm$\code{1)} and \code{atan(}$\pm$\code{oo)}
are the constants of GAOL: \code{acos([1, 3])} was
$[-2^{-1074}, 2^{-1074}]$, and \code{acosh([1])} with the libm was
$[-3 \cdot 2^{-1074}, 3 \cdot 2^{-1074}]$. The hyperbolic functions
\code{sinh}, \code{cosh}, \code{tanh}, \code{asinh}, \code{acosh} and
\code{atanh} are those of CORE-MATH, the tightest bounds on every system:
GAOL~4 took them from the libm of the system and moved their values one double
outward, while the libms of glibc 2.31, musl and MinGW-w64 are sometimes a
double further, and the \code{acosh} of the MinGW-w64 older than version 12 is
millions of doubles away next to 1. \code{cosh()} and \code{sinh()} are
$[\mathrm{MAX}, +\infty]$ from 711 on, and \code{tanh()} has the double below
1 for lower bound from 20 on; the functions take the same time as before,
within 1.5 ns (Intel i7-1185G7, GCC 9.4).

\medskip
The domains follow the standard.\newinvfive{} \code{log()} gives the empty
set for an interval holding no positive number, \code{log} being defined on
$(0, +\infty)$: GAOL~4 gave $[-\infty, -\mathrm{MAX}]$ for \code{log([-4,
0])} and \code{log([0])}. \code{atanh()} gives the empty set for an interval
holding no point of $(-1, 1)$: \code{atanh([1])} and \code{atanh([1, +oo])}
are empty, and the limits stay the bounds of an interval holding other
points, \code{atanh([0, 1])} being $[0, +\infty]$. \code{div\_rel()} rounds
the lower bound of the quotients downward where \var{K} is below 0 and
\var{J} straddles it, the one branch left without SSE code:
\code{div\_rel([-1], [-3, 1], [0, 10])} was one double above $1/3$, and the
numbers from $1/3$ to the double above it were lost, where \ieee{mulRev} has
to be valid. \code{invabs\_rel()} takes no magnitude from the negative part
of \var{J}: \code{invabs\_rel([-2, -1], [-10, 10])} was $[-2, 2]$ rather than
the empty set. The reverse functions of the power, \code{atan2}, \code{max},
\code{min}, \code{sign} and \code{floor} are new, their arguments in the
order of \code{div\_rel(K, J, I)}: \code{pow\_rel(K, J, I)},
\code{pow\_exponent\_rel(K, H, I)}, \code{atan2\_rel\_y(C, X, Y)},
\code{atan2\_rel\_x(C, Y, X)}, \code{max\_rel(Z, Y, X)},
\code{min\_rel(Z, Y, X)}, \code{sign\_rel(Z, X)} and \code{floor\_rel(Z,
X)}; \code{gaol\_ieee1788} names the first four \ieee{powRev1},
\ieee{powRev2}, \ieee{atan2Rev1} and \ieee{atan2Rev2}, and adds
\ieee{mulRevToPair}. All the reverse operations of the standard are
provided (Chapter~\ref{chap:arithmetic}).

\subsection{The relations}
\label{sec:changes-relations}
\cindex{relations!differences with GAOL~4}%

The comparisons with the empty set and the infinities follow IEEE 1788-2015
(Tables~10.3 and 10.4).\newinvfive{} \code{certainly\_leq()} and
\code{certainly\_le()}, \ieee{precedes} and \ieee{strictPrecedes}, and
\code{certainly\_geq()} and \code{certainly\_ge()}, are true when either
interval is empty: GAOL~4 gave false when the second one only was.
\code{set\_strictly\_contains()} and \code{set\_le()}, \ieee{interior}, take
an infinite bound as beyond the same infinite bound:
\code{interior(Entire, Entire)} and $[2, +\infty]$ interior to $[1, +\infty]$
were false, as \code{set\_le()} of the empty set in the empty set.

\medskip
The possibly relations, \code{certainly\_eq()}, \code{certainly\_neq()},
\code{==} and \code{!=} are removed, and so is the option that chose what the
relation symbols mean.\newinvfive{} \code{certainly\_neq()} was
\code{!certainly\_eq()}, which is ``possibly not equal'': $[3, 4]$ was
certainly not equal to $[3, 4]$, and so was \code{!=} with the certainly
relations, the default of every build. IEEE 1788-2015 has neither the
possibly relations nor a certain equality: its equality is \ieee{equal},
\code{set\_eq()}, and ``for all $x$ and $y$, $x \neq y$'' is
\ieee{disjoint}, \code{set\_disjoint()}; \code{possibly\_eq(y)} was
\code{!set\_disjoint(y)}. \code{<}, \code{<=}, \code{>} and \code{>=} are
the certainly relations in every build, \ieee{strictPrecedes},
\ieee{precedes} and their converses. The set relations of the option were
never taken, \cmd{configure} and meson defining \code{GAOL\_SET\_RELATIONS}
where the header read \code{GAOL\_SET\_RELATION}: they gave the possibly
relations (\pxref{sec:relational-symbols}).

\medskip
\code{gaol::lexicographic\_less} is a total order on the intervals, for the
containers and the algorithms of the standard library that compare with
\code{<} (\code{std::set}, \code{std::map}, \code{std::sort}).\newinvfive{}
\code{<} serves none of them: it is \ieee{strictPrecedes}, true as soon as
either interval is empty, and no strict weak ordering. A
\code{std::set<interval>} that holds $[1, 2]$ drops $[1, 3]$, and
\code{std::sort} of a vector that holds an empty interval can read past its
end. With the comparator, the empty set comes first, then the intervals by
their left bound and by their right bound, $-0$ and $+0$ being the same bound
(\pxref{sec:lexicographic-less}). \code{x.overlap(y)} and \code{gaol\_ieee1788::overlap(x, y)}, the
\ieee{overlap} of IEEE 1788-2015 (10.6.4), are new, with the type \code{overlapping\_state}: each returns one of
the 16 states of Table~10.7, which GAOL~4 had no function for. The state is
read from the bounds as integers in the order of the doubles, so that no
floating-point exception is raised, whatever the flush-to-zero modes. A point
interval on a bound of the other one is never \element{meets} nor
\element{metBy}, as the sentence of 10.6.4 has it
(Chapter~\ref{chap:relations}).

\subsection{Reading and writing intervals}
\label{sec:changes-io}
\cindex{input!differences with GAOL~4}%
\cindex{output!differences with GAOL~4}%

Numbers are read exactly, whatever the locale, the flush-to-zero modes and
the C library.\newinvfive{} GAOL~4 read them with \code{strtod()} rounding
downward and upward, and relied on the inexact flag: the C runtime of Windows
and musl on 64-bit ARM processors round to nearest in every direction and
raise no flag, so \code{interval("0.1")} did not enclose $1/10$ there. Each
number is now compared exactly with the doubles around it, whatever the
number of digits of its significand and of its exponent; the doubles around
the number are bracketed by steps doubling from the double \code{strtod()}
gives, then found by bisection on their bits, two comparisons when
\code{strtod()} gives the nearest double, and at most about 125 whatever it
gives. The time to read a number is quadratic in the number of its digits
(\pxref{sec:reading-intervals}), so that a text from elsewhere can block a
program.
\code{strtod()} reads the decimal point of the C locale: under a locale
writing a decimal comma, it stopped at the \code{.} of \code{0.1}, and
\code{interval("0.1")} was $[0, 0]$; the text is taken apart once, and a
number of 20000 characters, which took over 2 s under such a locale, takes
about 0.02 s. Reading raises no overflow or underflow exception: with glibc,
\code{strtod()} raises them for \code{1e400} or \code{1e-310}, and a program
that had enabled the exception died while it read; a number that may be
beyond the largest double or below $2^{-1022}$ is compared with them first,
and is not given to \code{strtod()}. With denormals-are-zero the reader took
every subnormal for 0, and \code{textToInterval("1e-310")} did not contain
$10^{-310}$; the doubles are taken apart by their bits, which no mode
changes (\pxref{sec:reading-intervals}).

\medskip
Intervals are written in decimal rounded outward, whatever the C
library.\newinvfive{} GAOL~4 set the rounding direction downward, then
upward, and let the C library write each bound: the C runtime of Windows
rounds the magnitude, and wrote $-2/3$ rounded downward \code{-0.6666}, and
musl wrote \code{1.235e+05} for the lower bound of 123456.789; read back,
such a text did not enclose the interval written. The magnitude of each
bound is now written rounding to nearest, compared exactly with the bound,
and its last digit is moved by one when it is on the wrong side; with glibc,
which rounds as asked, the text is the one the C library wrote: 1\,077\,546
bounds written in the general, scientific and fixed formats with 1 to 18 digits
were the same, character for character. A
point interval whose double the digits write exactly is written $[a]$, the
literal of the standard, and a zero of either sign $[0]$: GAOL~4 wrote every
point interval \code{<a, b>}, the text of its double rounded downward and
upward, which the reader refused for \code{interval(0.1)}, throwing
\code{input\_format\_error}, because it takes \code{<a, b>} for two numbers
that are the same double, so that reading back the text of a point interval
failed for about 61\,\% of 11\,788 point intervals tried; the reader still reads \code{<a, b>}, with the numbers
compared by their bits. The format \element{width} writes the midpoint and
the radius of the standard, \code{midpoint()} rounded to nearest and
\code{rad()} rounded upward, which is $+0$ for a point zero of either sign, and
\element{center} the midpoint: GAOL~4 wrote
$(l + r)/2$ and $(r - l)/2$, both rounded to nearest, so that $[1, 1 +
2^{-52}]$ was written \code{1 (+/- 1.110223024625157e-16)}, which stops short
of the upper bound, and $[0, 5 \cdot 10^{-324}]$ had a radius 0. The empty
set is \code{[empty]} in every format. The format \element{bounds} is
written in the C locale whatever the locale of the stream, so that the reader
reads the text again: GAOL~4 wrote \code{[0,25, 0,5]} under a comma locale.
The format \element{agreeing} keeps the digits of each bound: GAOL~4 dropped
from both bounds what followed the last character of the left one that is
not a zero, and $[1.25, 1.2567]$ was written \code{1.25\textasciitilde{}[, ]}, which says
nothing of the right bound; it is \code{1.25\textasciitilde{}[0, 67]} now, and $[1, 10]$
is written with its bounds, where it was \code{1\textasciitilde{}[., 0]}. The hexadecimal
format writes an interval literal, \code{[0x1.999999999999ap-4, ...]}, read
back bit for bit, where GAOL~4 wrote the sixteen hexadecimal digits of each
double, which no parser read; \code{exact\_string(I)} gives that text
without the global output format, and \code{gaol\_ieee1788::intervalToText()}
writes a literal with the digits of \code{interval::precision()}, whatever
the format and the locale. \code{operator<<} leaves the stream as it found
it: GAOL~4 set its precision and left it so, the doubles written afterwards
getting 16 digits, and wrote the \code{[} apart, so that
\code{std::setw(8)} padded it alone. The text is made in a character string,
taking 30 to 45\,\% less time than the \code{std::ostringstream}s of the earlier
versions of GAOL v5 (Clang 18, Release) (\pxref{sec:writing}).

\medskip
\code{operator>>} ends with the input, and skips blank lines.\newinvfive{}
GAOL~4 read the rest of the line where the previous value stopped, so that
\code{in >> d >> x} over two lines threw \code{input\_format\_error}, and so
did \code{while (in >> x)} over a file ending with an empty line: such a loop
always ended with an exception. Where no line is left, the reader sets the
\code{failbit} of the stream, leaves the interval as it was and throws
nothing, as the reading of a double does; a line that is no interval sets the
\code{failbit} too, empties the interval, and \code{std::invalid\_argument}
is thrown. A program that reads on after it calls \code{in.clear()} first.

\medskip
One grammar reads every string as an expression.\newinvfive{} GAOL~4 read a
string with two grammars: the numbers, the constants and the functions of
them built a tree, while the intervals written between brackets, and the
operations on them, were computed at once; Bison chose between the two at
each token (124 shift/reduce and 14 reduce/reduce conflicts), and once a
number had started the tree no interval could come, so that \code{1+[1,2]}
and \code{2*cos([0,1])} were refused while \code{[1,2]*2} was read. The
literals of intervals, the uncertain numbers, \code{empty} and \code{<a,b>}
are now leaves, computed when they are read, as are the calls of functions
and the binary operations; the grammar has no conflict, and an error found
in an action stops the reading at once. An uncertain number is still no
bound of a literal: \code{[5?1]} throws \code{std::invalid\_argument}, the
standard giving it as a string that is not an interval literal. The literals
of IEEE 1788-2015 are read, whatever the case of their letters: \code{[ ]},
\code{[entire]}, the bounds left out (\code{[1,]}, \code{[,]}),
\code{infinity}, the hexadecimal numbers and the uncertain form
(\code{3.56?1}, \code{-10??u}), as the tightest intervals of doubles
enclosing them; \code{[inf]} and \code{[inf, inf]} are the empty set, where
GAOL~4 gave $[\mathrm{MAX}, +\infty]$. Numbers in base 16, 8 and 2 are read,
\code{0x10}, \code{0o17}, \code{0b101} and \code{0b1.1p3}, a number that
begins with 0 being decimal. $[l, u]$ with $l$ above $u$ is the empty set
also where the bounds are numbers that are not doubles: GAOL~4 read
\code{[1.00000000000000001, 1]} as $[1]$. A bound of \code{<a, b>} has to be
a double, and a number that is not one, \code{pi} included, is refused. The
reader knows every function of GAOL on intervals, which GAOL~4 did not read
(\code{exp10}, \code{hypot}, \code{sqr}, \code{abs}, \code{fma}\ldots), and
\code{gaol\_ieee1788::textToInterval()} reads the names of the standard,
whose \code{pow} is the pow of the standard: in the earlier versions of GAOL v5
it read those of GAOL, and \code{pown} gave the empty set. Function names are read in any case of
letters, compared in ASCII and not by \code{std::tolower()}, whose result
depends on the locale. \code{pow(x, n)} read in a string is the \code{pow}
of \code{gaol} for an integer exponent too: the parser took $1/x^{|n|}$ for
$n < 0$, whose $x^{|n|}$ overflowed, and 1 for $n = 0$, which is no value for
an empty $x$ (\pxref{sec:input-format}).

\medskip
The reader is robust.\newinvfive{} The parser frees the nodes of the
expressions it reads: it kept a reference to each node, and the nodes of
every \code{interval("...")} were never freed; the nodes left by a syntax
error are freed too, and reading \code{nth\_root(8, 1.5)}, then the same
text plus 1, crashed, where both now throw \code{std::domain\_error}.
Several threads can read strings at once: the lexer of flex kept its buffer
in globals, and two threads each building an interval from a string crashed;
the lexer is reentrant and the parser pure, and four threads reading strings
at once, which crashed each time with the globals, show nothing under
ThreadSanitizer. Long expressions no longer overflow the stack: a sum of
100000 terms crashed the program, and is read in a time proportional to its
length; the stack of the parser, 10000 entries, bounds the nesting. A
\code{std::string} is read whole, where the lexer stopped at the first NUL
character. An allocation that fails throws \code{std::bad\_alloc}, where
flex ended the program with \code{exit(2)}. Built without exceptions (a setting of the developers), a
wrong call in a string fails the reading as a syntax error does, where an
earlier version of GAOL v5 ended the program inside the reader:
\code{gaol\_ieee1788::textToInterval()} gives the empty set,
\code{parse\_interval()} returns false, and \code{gaol::textToInterval()} and
\code{operator>{}>} report the error and abort. \code{parse\_interval(nullptr)} returns false, where the program
died of a segmentation fault. A newline is a space, where it fell to the
default rule of flex, which wrote it on the standard output. Expressions are
printed as written: a division with \code{/}, where GAOL~4 wrote \code{*},
and a negative number within the parentheses a power of it needs.

\subsection{The errors}
\label{sec:changes-errors}
\cindex{exceptions!differences with GAOL~4}%

GAOL throws \code{std::invalid\_argument} and \code{std::domain\_error}, and
\code{what()} says what went wrong.\newinvfive{} GAOL~4 had a class of
exceptions of its own, \code{gaol\_exception}, and
\code{input\_format\_error}, \code{invalid\_action\_error} and
\code{unavailable\_feature\_error} derived from it, in
\file{gaol/gaol\_exceptions.h}. They are gone, without aliases. A string
that is no interval, or a null pointer, throws \code{std::invalid\_argument};
an argument that a function does not take (\code{nb\_fp\_numbers()} with a
NaN, \code{bisect()} with a ratio outside $(0, 1)$, \code{nth\_root(8, 1.5)}
in a string) throws \code{std::domain\_error}; no operation of GAOL v5 throws
\code{unavailable\_feature\_error}, which only the \code{atan2} of GAOL~4
did. A handler of \code{gaol::input\_format\_error} becomes a handler of
\code{std::invalid\_argument}, one of \code{gaol::invalid\_action\_error} a
handler of \code{std::domain\_error}, and \code{std::logic\_error} catches
both. GAOL~4 left the \code{what()} of \code{std::exception}, the text
\code{std::exception} with libstdc++ and libc++ whatever the error; it is
\code{file:line: explanation} now, the file without its directories.
\code{file()}, \code{line()}, \code{explanation()} and \code{operator<<} of
the exceptions are gone with the classes. A null file name is ``(unknown
file)'' in the message of an exception, in \code{gaol\_warning()} and in
\code{gaol\_error()}: GAOL~4 wrote the null pointer to the stream, which is
undefined, and libstdc++ put the stream in the state \code{bad}, so that the
diagnostics written after it were lost (\pxref{sec:exceptions}).

\subsection{The namespaces, the headers and the interface}
\label{sec:changes-interface}
\cindex{namespaces!differences with GAOL~4}%
\cindex{headers!differences with GAOL~4}%

The type \code{interval}, GAOL's functions and its expressions are in
\code{gaol\_core}; \code{gaol} names them as GAOL~4 did, and
\code{gaol\_ieee1788} as IEEE 1788-2015 does.\newinvfive{} \code{gaol} names
them one by one, with using-declarations, rather than opening
\code{gaol\_core} whole: the names GAOL~4 had in \code{gaol} that GAOL v5
still has, the functions on intervals GAOL v5 adds, \code{restore\_rounding()}
and \code{exact\_string()}. The two namespaces differ mainly in the power,
the standard having two where GAOL~4 had one: the \code{pow} of \code{gaol}
takes the integer power for an integer exponent, a negative base included,
and the real power otherwise, while in \code{gaol\_ieee1788} \code{pow} is
the real power, on the part of $x$ in $[0, +\infty]$, and \code{pown} the
integer power. \code{pow} and \code{textToInterval}, whose meaning differs
between the two, are in each of them and not in \code{gaol\_core}, so that
each namespace finds its own only. \code{gaol} does not name the exceptions,
\code{NaN\_val}, \code{eval\_stack} nor \code{parse\_interval()}, which
GAOL~4 had in \code{gaol}, nor what GAOL uses for itself; a program names
them in \code{gaol\_core}. \code{f\_negate()}, \code{f\_negate\_simple()},
the macro \code{opposite()}, \code{get\_fpu\_cw()} and
\code{reset\_fpu\_cw()}, which GAOL v5 no longer calls, are gone.
\file{gaol/gaol.h} no longer opens \code{gaol}: a program adds \code{using
namespace gaol;} or \code{using namespace gaol\_ieee1788;}, one of the two,
\code{pow(x, y)} being ambiguous with both (\pxref{sec:namespaces}). The
\code{double} constants of \file{gaol/gaol\_port.h}, \code{two\_pi},
\code{pi}, \code{half\_pi}, \code{pi\_dn}, \code{ln2\_dn}\ldots, are no longer
declared for the code using GAOL: a program that opened \code{gaol} found
them among its own names, and a \code{pi} of its own was ambiguous; the
bounds of $\pi$ are those of \code{interval::pi()}, and those of $\ln 2$
those of \code{log(interval(2.0))}.

\medskip
The expressions were repaired.\newinvfive{} \code{pow(e, n)} on an
expression links: the library defined it with an \code{unsigned int}
exponent, where the header declares an \code{int}. \code{expression::operator/=}
links: it was declared and not defined. Empty expressions can be created in
several threads: they all point to one node, whose count of references, a
plain \code{unsigned int}, lost updates, and two threads each creating and
destroying a million empty expressions crashed each time; the references to
that node are no longer counted, and it lives from the initialization of GAOL
to its automatic cleanup. Other expressions still share the nodes of their
copies, whose counts are not atomic: such an expression is not copied or
destroyed in two threads at once.

\medskip
The headers GAOL installs are the ones a program includes.\newinvfive{} The
three builds install the same list: \file{gaol/gaol}, \file{gaol/gaol.h}, the
headers they include, and \file{gaol/parser/gaol\_expression.h} and
\file{gaol/parser/gaol\_expr\_eval.h} for the expressions; a program that
builds expressions includes \file{gaol/parser/gaol\_expression.h}, where it
included \file{gaol/gaol\_expression.h}. They install no header of GAOL's sources:
neither the header of the parser, which did not compile when included, nor
\file{gaol/gaol\_init\_cleanup.h}, both of which GAOL~4 installed, nor the
headers of \file{gaol/math/}, among them \file{gaol\_double\_op.h}, whose
functions of doubles \code{exp\_dn()}, \code{sin\_up()} and the thirty others
leave the interface. Each installed header, but the four of Section~\ref{sec:headers}, compiles
alone with
\option{-std=c++11} and the flags of \file{gaol.pc}, and without a warning
under \option{-Wall -Wextra} and \option{-Wold-style-cast -Wundef}: with
\option{-Wextra}, every assignment of intervals warned that the implicit
assignment was deprecated. The class \code{interval} is declared in
\file{gaol/core/gaol\_interval\_core.h}, and the headers of the SSE2, NEON
and FPU intervals are installed in \file{include/gaol/core/};
\file{gaol/gaol\_interval.h} stays and includes it, so that a program that
includes it, or \file{gaol/gaol}, is not touched by the move. Gone are the
timing of GAOL~4 (\file{gaol/gaol\_profile.h}, the class \code{timepiece}
and \code{get\_time()}, \code{elapsed\_time()}\ldots; a program measures the
time with the clocks of the C++ standard library),
\file{gaol/gaol\_parameters.h} and \file{gaol/gaol\_limits.h} (a program
includes \file{<limits>} itself), \file{gaol/gaol\_assert.h} and
\code{GAOL\_ASSERT} (a program writes \code{assert()}),
\file{gaol/gaol\_exact.h} with \code{get\_inexact()} and
\code{clear\_inexact()}, which the lexer called to know whether
\code{strtod()} had rounded, \file{gaol/gaol\_flags.h} and
\code{s\_nextafter()}, and \file{gaol/gaol\_allocator.h} with
\code{gaol::aligned\_allocator}, which neither GAOL, Codac nor IBEX used.
The \code{operator new} of the SSE2 intervals still aligns their memory.

\medskip
\code{\_\_GAOL\_PUBLIC\_\_} is \code{GAOL\_PUBLIC}, and is empty with Visual
C++ and clang-cl, with which GAOL is always a static library.\newinvfive{}
GAOL~4 declared its functions \code{\_\_declspec(dllimport)} there, so that
the code using GAOL had to define \code{\_\_GAOL\_PUBLIC\_\_} empty; no
header reads it any longer. No header of GAOL declares \code{string} and
\code{exception} in the global namespace, and the macros of GAOL's headers
start with \code{GAOL\_}: \code{INLINE}, \code{MEMALIGN}, \code{\_\_HI},
\code{WORDS\_BIGENDIAN}\ldots are \code{GAOL\_INLINE},
\code{GAOL\_MEMALIGN}, \code{GAOL\_HI}, \code{GAOL\_WORDS\_BIGENDIAN}\ldots
\code{\_XOPEN\_SOURCE} is no longer redefined to 600, and
\code{nb\_fp\_numbers()} returns an \code{unsigned long long}. The macros
that the installed headers read with \code{\#if} are 1 or 0, never undefined,
so that they compile under \option{-Wundef -Werror}: a program has to test
them with \code{\#if} as well, \code{\#ifdef GAOL\_PRESERVE\_ROUNDING} being
true whatever the build.

\subsection{The builds, the compilers and the options}
\label{sec:changes-builds}
\cindex{builds!differences with GAOL~4}%

The three builds agree.\newinvfive{} Before, each had its own idea (GAOL~4 had the first two, and the CMake build
is that of the 4.3 releases of GAOL v5):
\cmd{configure} optimized only when the compiler was named \cmd{g++}
exactly, computed the doubles of 32-bit x86 on the x87 unit, read the
processor of the machine building rather than the target of the compiler, and
refused macOS, Windows, s390x and riscv64; meson did not optimize, compiled
with the assertions of libstdc++, refused i686 and armv7l and linked CRlibm
by default; CMake left out \option{-msse3}, the hidden visibility and
\option{-Wconversion}. \option{-fno-fast-math}, \option{-ffp-contract=off}
and the refused compilers were in CMake only: without
\option{-ffp-contract=off}, the compilers fused the multiplications and
additions of \code{mid()}, which returned a single double, not enclosing the
midpoint of 9735 of the 10000 random intervals of a test, with meson on
macOS arm64 and with autotools on Ubuntu arm64. The three builds now write
the same macros in \file{gaol/gaol\_configuration.h}, offer the same options
(\code{GAOL\_PRESERVE\_ROUNDING}, \code{GAOL\_SIMD}, \code{GAOL\_COVERAGE},
\code{WITH\_TESTS}, \code{WITH\_EXAMPLES}, \code{BUILD\_SHARED\_LIBS} and
the build type, with their counterparts in \cmd{configure} and meson), and
decide the rest by the configuration: the fused multiply-add instructions and
the AVX-512 path by checks that run a program, the 128-bit integer by a check
that compiles. \code{GAOL\_FMA}, \code{GAOL\_PREFER\_AVX512},
\code{GAOL\_U128\_EMULATION}, \code{GAOL\_DEBUG}, \code{GAOL\_EXCEPTIONS\_ENABLED} and
\code{GAOL\_VERBOSE\_MODE} are settings of the developers, no longer options
of the user. The Debug build of each compiles GAOL with \option{-g}, without
optimization nor \code{NDEBUG}; \loption{enable-optimize} of
\cmd{configure} and \option{-Denable-optimize} of meson are gone, the Release build giving the
optimizations. The C sources of CORE-MATH are compiled with
\option{-std=gnu99}. The shared library is \file{libgaol.so.5.0.0} with the
three builds, made when asked, the static one alone by default. GAOL can be a
subproject of meson, and CMake without a build type compiles as in Release,
where, in the earlier versions of GAOL v5, \file{libgaol} had \option{-O3}
alone and CORE-MATH no optimization.
The unit tests of GAOL~4, which were in \file{check/} and needed CppUnit,
are in \file{tests/} and run without it; \cmd{make test} runs them,
\cmd{make check} the tests and the examples, \cmd{make perf} the benchmark
and \cmd{make parser} regenerates the lexer and the parser
(Chapter~\ref{chap:installation}, \pxref{sec:tests}).

\medskip
The compilers and the options that give wrong bounds are
refused.\newinvfive{} Visual C++ and clang-cl without \option{/fp:strict}
are refused, as \option{/fp:fast} is: Visual C++ then assumes rounding to
nearest, and nothing certifies the bounds. clang-cl for 32-bit x86 is given
\option{-msse2}: clang-cl 18 computes the doubles on the x87 unit there,
even with \option{/arch:SSE2}, and CORE-MATH rounds some results twice.
\option{-ffinite-math-only} is refused, as \option{-ffast-math} is; GAOL~4
refused none of them. \option{-ffast-math} and \option{-Ofast} turn it on,
and the compiler then takes NaN and infinities never to occur, in the inline
functions of the headers too: the empty interval has NaN bounds, and
$([1, 2] \cap [3, 4])$.\code{is\_empty()} is false once the compiler has
folded the test away (GCC 9.4 and 13 at \option{-O0}, \option{-O2} and \option{-O3}; Clang 18 at
\option{-O0}, and at \option{-O2} and \option{-O3} when the bounds are
\code{volatile}). With Clang, \option{-ffast-math} and \option{-Ofast} turn
\option{-frounding-math} off too, which \option{-fno-fast-math} does not turn
on again: Clang 18 compiled GAOL's operations as if rounding to nearest (54
of the 292728 checks of the test of the rounding direction failed after
\option{-ffast-math -fno-fast-math}, 126 after \option{-Ofast
-fno-fast-math}, none with all the flags after the option). What no macro
shows is not refused, and the code using GAOL is not to be compiled with it:
\option{-funsafe-math-optimizations} and \option{-fno-signed-zeros}, which
rewrite the addition by which GAOL tells the rounding direction and the
modes. The compilers that do not honour the rounding direction are refused,
the C compiler as well as the C++ compiler: Clang for 32-bit ARM processors,
and any compiler that says of \option{-frounding-math} that the rounding
mode is unsupported on its target, as Clang 14 does for 64-bit ARM. GAOL~4
accepted them; built by Clang 21 for 32-bit ARM, GAOL v5 gave 4556 of 16000
random products, squares and cubes that did not enclose their exact values
(\pxref{sec:refused}).

\medskip
GCC for Windows is given the fused multiply-add instructions (MinGW-w64,
MSYS2 and Cygwin): the three builds give it \option{-mfma
-mprefer-vector-\allowbreak width=128}.\newinvfive{} It had no flag: \option{-mfma}
turns AVX on, GCC then vectorizes with 256-bit registers and spills them to a
stack it takes to be aligned on 32 bytes, where the ABI of Windows x64
guarantees 16 (GCC bug 54412, which crashed the MinGW builds of Codac and is
fixed for GCC 17, commit 8aaa5f6e6207 of 25 September 2026), so that CORE-MATH called a software \code{fma()}. With the
second flag \cmd{objdump} finds no \code{\%ymm} in the objects, against
30598 lines with \option{-mfma} alone (GCC 16.2, MSYS2 UCRT64). On the
runner of GitHub, \code{pow(x, 3)} takes 8.7 ns rather than 52.3,
\code{log(x)} 33.6 rather than 106.4 and \code{cos(x)} 55.4 rather than
260.5; with MinGW-w64 15.2 for 32-bit x86, 12.1 ns rather than 70.5, 81.9
rather than 265.7 and 138.3 rather than 489.6. The library then needs a
processor with the FMA and AVX instructions (Intel Haswell, AMD Piledriver,
2012--2013, and later), as with the other x86 compilers. Cygwin on x86-64 is
built and tested, and refused without the instructions, the \code{fma()} of
its C library not being correctly rounded; with it, \code{previous\_float()}
and \code{next\_float()} give $-\infty$ and $+\infty$ at $\pm$\code{DBL\_MAX}
with every C library, where the \code{nextafter()} of Cygwin computed the
overflow as $x + x$ in the rounding direction of the program, so that the
lower bound of \code{sinh(-711)} did not enclose its value
(\pxref{sec:platforms}).

\medskip
The intervals of floats, \code{gaol::intervalf} and \code{gaol::interval2f},
are compiled by a developer of GAOL only.\newinvfive{} Every build compiled
and installed them, and \file{gaol/gaol} included their headers, while
neither IBEX nor Codac uses them, the manual leaves them undocumented, and
they are unfinished: \code{sqrt(intervalf)} returns its argument, and
\code{interval2f::inverse()} aborts. Their code is within
\code{\#ifdef GAOL\_FLOAT\_INTERVALS}, a macro no build has an option for.
The archive of the sources and the packages are made by CMake: \cmd{cmake
--build <build> --target package\_source} gives \file{gaol-<version>.tar.gz},
and \loption{target package} an archive of what is installed and, on Linux, a
\file{.deb}; \cmd{make dist} and \cmd{make rpm} of the autotools build are
gone, \cmd{make dist} stopping on a file it could not copy. The version is in
\file{VERSION.txt} alone, where GAOL~4 wrote it in \file{configure.ac}, in
\file{meson.build} and in two headers, and \cmd{configure --help} gives the
address for the bug reports. \code{is\_finite()} is \code{std::isfinite()}
in every build, \code{finite()} of the C library not being declared by every
one. \loption{enable-asm} and
\code{HAVE\_ROUNDING\_MATH\_OPTION} are gone, the first choosing an assembly
GAOL v5 no longer calls. CORE-MATH is in the sources, so that GAOL can be
built as a part of another project, brought in by FetchContent, with no
network access beyond its own sources. The CMake build derives from the CMake
build of GAOL in IBEX, with the compilation flags of the fork of GAOL by
Fabrice Le Bars, and the tests of \file{tests/} follow the rounding tests of
Codac (Chapter~\ref{chap:tests}). The Goldstein-Price function of the manual
and of the examples had the factor $(x + y)^2$ where the function has
$(x + y + 1)^2$, as in GAOL~4; both now compute the true function, whose
natural extension over $[-2, 2]^2$ is $[-87881320, 147125080]$, 231 times as
wide as the range, which the example says it encloses.

\section{Discussion}\label{sec:changes-discussion}
%%-------------------

The move from GAOL~4 to GAOL v5 costs a program little: the namespace to
open, the strings to read with \code{textToInterval()}, the relations
\code{==} and \code{!=} to rewrite with \code{set\_eq()},
\code{set\_disjoint()} or \code{!set\_disjoint()}, the handlers of the exceptions of GAOL~4 to
retarget on those of the standard library, and the conversions of a double or an integer
to an interval, explicit now, to write out. What the program then sees is
decided by three choices that run through the register of this chapter. The
first is CORE-MATH as the only source of the elementary functions, which gave
the tightest bounds and removed the mathematical library from the link, at
the price of sources vendored into GAOL and ported to the compilers that lack
a 128-bit integer. The second is the rounding direction taken in charge by
each operation, which made the bounds right whatever the direction the
program computes in, at the price of a check whose cost was measured on each
processor, and whose variants (the control registers written directly, the
AVX-512 path, the NEON register) each removed a part of it. The third is IEEE
1788-2015 as the arbiter wherever GAOL~4 had its own convention: the empty
set for the constructors, the comparisons, the literals read and written, the
domains of \code{log}, \code{atanh} and the roots.

\medskip
Some differences are deliberate losses: the possibly relations and
\code{certainly\_neq()}, which was one of them under another name, the timing of GAOL~4, the intervals of floats, the
headers of the sources. Others are left as they are, and the register says
so: the sign of a zero bound of a product, which depends on the kernel, the
quadratic time of the reading of a number of many digits, the inexact exception
that almost every operation raises, and the processors where neither the
flush-to-zero modes can be cleared nor \option{-mno-daz-ftz} exists. A
program that depended on one of the former behaviours finds here what
replaced it, and in the chapters this chapter refers to how the new
behaviour is defined. The appendices, after the conclusion, give where the
documentation of GAOL v5 lives and how to report what this chapter does not
cover.


%%==========================================================================
%% Conclusion
%%==========================================================================
\cleardoublepage
\pagestyle{CONCLU}
%% !TEX root = ../manual_gaolv5.tex
% The conclusion (unnumbered).
\schapter{Conclusion}\label{chap:conclusion}
\thispagestyle{plain}

This manual described GAOL v5, the continuation of the GAOL of Fr\'ed\'eric
Goualard, as a user's guide and as a reference. We recalled the notions it
rests on, from the roundings of floating-point numbers to the accuracy modes
of IEEE 1788-2015; we described its installation with the three builds, the
flags the code using it is compiled with, the two namespaces under which its
operations are named, and the rounding direction and the floating-point
environment that GAOL v5 takes charge of; we gave each operation of the class
\code{interval} with its examples, the names of the standard, and the
tightness of every operation; we explained how the intervals are computed
and measured, how GAOL is tested, and what a program of GAOL~4 has to change.

\medskip
Three facts sum up what GAOL v5 established. First, the bounds it computes
are the same on every architecture and with every compiler it supports:
every elementary function that CORE-MATH provides is bounded with it, correctly
rounded in the direction in effect, and the integer powers, the roots and
\code{compoundm1} are the tightest bounds or within a few doubles of them; the
tests compare the bounds with their exact
values, computed apart, on every platform of the continuous integration.
Second, the operations follow IEEE 1788-2015 for the empty set, the infinite
bounds, the domains of the functions, the comparisons and the literals; under
the names of the standard, GAOL v5 gives the result of the standard in each of
the 288 special cases of the comparison with the other libraries where the standard
defines one (Chapter~\ref{chap:computing}), and under its
own names in all but those of its hybrid \code{pow}, kept for the programs of
GAOL~4. Third, the rounding direction is no longer left to the user: each
operation checks it and sets it upward when it is not, but the operations of
the AVX-512 path, whose instructions carry the direction, and GAOL can be built
to restore the direction of the program after each operation, at a cost that
the AVX-512 path, where the machine has the instructions, brings back to a
few tens of percent.

\medskip
What GAOL v5 does not provide is stated where it would be expected. GAOL has
no decorations: only bare intervals are provided, and a program that needs
to know that a box lies outside the domain of a function tests the empty
set itself, as Chapter~\ref{chap:overview} explains. The slope
functions, the reduction operations and the exact reduction operations of the
standard are not provided either. The intervals of floats are unfinished, and no build
compiles them. The \code{pow} of the namespace \code{gaol} is a hybrid of the
two powers of the standard, which keeps the results of GAOL~4 for the
programs written for it and for IBEX and Codac; a program that wants the
results of the standard opens \code{gaol\_ieee1788}.

\medskip
What is left to measure concerns the processors on which GAOL v5 could not be timed.
Of the instruction sets read for Chapter~\ref{chap:computing}, the scalar
instructions of RISC-V are the only ones, beside the AVX-512 path of x86,
that carry their rounding direction; a path written with them would need
neither a check nor a restore of the direction, and waits for a board to be
measured on. The vector units of POWER and s390x hold two doubles per
register as NEON does, and wait for hardware as well. Until then, the FPU
intervals compute on these processors, with the same bounds, as the
continuous integration checks under emulation.


%%==========================================================================
%% Bibliography
%%==========================================================================
\cleardoublepage
\pagestyle{BIBLIO}
\addcontentsline{toc}{chapter}{\bibname}
\bibliographystyle{plain}
\bibliography{gaol}

%%==========================================================================
%% Appendices
%%==========================================================================
\appendix
\pagestyle{YO}
\renewcommand{\chaptermark}[1]{\markboth{\appendixname\ \thechapter\ - \ #1}{}}

%% !TEX root = ../manual_gaolv5.tex
% Appendix A.
\chapter{Documentation and references}\label{chap:documentation}

\section{Documentation on GAOL}

The primary reference is this manual. The directory \file{doc} of the sources
holds the pages that give the details of GAOL v5, with the measures behind
its choices, whose content the chapters of this manual take up:

\begin{itemize}
\item \file{doc/building.md}: the CMake, autotools and meson builds, and
  their options (Chapter~\ref{chap:installation});
\item \file{doc/using.md}: the flags of interval arithmetic, GAOL from CMake
  and from pkg-config, the namespaces, the rounding direction and the
  floating-point environment (Chapters~\ref{chap:overview}
  and~\ref{chap:rounding});
\item \file{doc/three-builds.md}: what the three builds agree on, what the
  options are worth, and the compilers and options they refuse
  (Chapter~\ref{chap:installation});
\item \file{doc/tests.md}: what the programs of \file{tests} check
  (Chapter~\ref{chap:tests});
\item \file{doc/accuracy.md}: the tightness of each operation, with its
  algorithm and what the tests found (Chapter~\ref{chap:accuracy});
\item \file{doc/differences.md}: each change GAOL v5 brings to GAOL 4.2.2,
  and where it comes from (Chapter~\ref{chap:changes});
\item \file{doc/continuous-integration.md}: the systems, processors and
  compilers GAOL is built and tested on (Chapter~\ref{chap:tests});
\item \file{doc/simd.md}: what GAOL computes with vector instructions and
  with the fused multiply-add, where the time of a program goes, and the
  rounding direction by processor (Chapter~\ref{chap:computing});
\item \file{doc/math-libraries.md}: the mathematical libraries libm,
  mathlib, CRlibm and CORE-MATH, and what separates them
  (Chapter~\ref{chap:concepts});
\item \file{doc/tools.md}: the scripts of the repository and what each is
  for (Chapter~\ref{chap:tests});
\item \file{doc/compare/README.md}: GAOL compared with libieeep1788,
  filib++, PROFIL/BIAS, Solaris Studio and Boost.Interval, their special
  cases against IEEE 1788-2015, the enclosure of the elementary functions of
  GAOL and Boost.Interval, their performance, and the test suites ITF1788 and
  of Tang et al. (Chapter~\ref{chap:computing});
\item \file{examples/examples.md} and \file{examples/tutorial.md}: the
  example programs, a review of GAOL v5 and a comparison with other interval
  libraries, and a short tutorial (Chapter~\ref{chap:tests});
\item \file{3rd/README.md}: the version of CORE-MATH GAOL takes, what differs
  from it, and the fixes to propose upstream.
\end{itemize}

\noindent
The manual of GAOL~4, by Fr\'ed\'eric Goualard, is kept in \file{manual/v4}; its
PDF, \file{manual/v4/Manual\_GAOLv4.pdf}, was built once, and no build makes it
again.

\section{References}

The following articles, books and standards have inspired in some way or
another the devising of the GAOL library and the writing of this manual.

\begin{itemize}
\item \emph{IEEE Standard for Interval Arithmetic}. IEEE Std
  1788-2015~\cite{IEEE1788:2015}.
\item \emph{The CORE-MATH Project}. Alexei Sibidanov, Paul Zimmermann and
  St\'ephane Glondu~\cite{CoreMath:2022}, and
  \url{https://core-math.gitlabpages.inria.fr/}.
\item \emph{How do you compute the midpoint of an interval?} Fr\'ed\'eric
  Goualard~\cite{Goualard:TOMS2014}.
\item \emph{Fast and Correct SIMD Algorithms for Interval Arithmetic}.
  Fr\'ed\'eric Goualard~\cite{Goualard:PARA08}.
\item \emph{Interval Extensions of Multivalued Inverse Functions}.
  Fr\'ed\'eric Goualard~\cite{Goualard:2008}.
\item \emph{Interval Arithmetic Specification}. Dmitri Chiriaev and G.\:
  William Walster. Draft revised May 1998.
\item \emph{The Extended Real Interval System}. G.\: William Walster. April
  1998.
\item \emph{C++ Interval Arithmetic Programming Reference}. Sun Microsystems,
  Inc. October 2000, revision A.
\item \emph{Interval Arithmetic: From Principles to Implementation}. T.\:
  Hickey, Q.\: Ju, and M.\: H.\: van Emden~\cite{Hickey-Ju-vanEmden:2001}.
\end{itemize}

\section{Reporting bugs}

Bugs and suggestions for improvement of GAOL v5 are submitted as issues
of its repository:

\begin{center}
  \url{https://github.com/Jordan08/GAOL/issues}
\end{center}

\noindent
A report is most useful with the program that shows the problem, the
interval it gives and the one expected, the system, the processor, the
compiler and the options GAOL was built with.

%% !TEX root = ../manual_gaolv5.tex
% Appendix B.
%%
%% Copyright (c) 2026 ENSTA, France
%%
%% By: Jordan NININ (created 2026-11-06)
%%
%% GAOL is a software distributed WITHOUT ANY WARRANTY. Read the associated
%% COPYING.LIB file for information.
\chapter{Library copying}\label{chap:copying}

GAOL is distributed under the GNU Library General Public License,
version~2, reproduced below. The copyright for the initial version, named
\code{cell}\findex{cell} (year 2001), is owned by the \emph{Swiss Federal Institute of
Technology}, Lausanne, Switzerland. The copyright for the following versions
(from 2002 onward) is owned by the \emph{Laboratoire d'Informatique de
Nantes-Atlantique}, then by the \emph{Laboratoire des Sciences du
Num\'erique de Nantes}, France. The files GAOL v5 adds carry the copyright of
ENSTA, France, and are distributed under the same licence as GAOL; the files
of GAOL that GAOL v5 changes keep the copyright of their authors.
\cindex{licence!of GAOL}%

GAOL v5 computes its elementary functions with CORE-MATH~\cite{CoreMath:2022},
whose sources are distributed with GAOL (in \file{3rd/math-core}) and
compiled into its library. CORE-MATH is distributed under the MIT licence,
reproduced after the licence of GAOL.
\cindex{CORE-MATH}\cindex{licence!of CORE-MATH}%

\section{The licence of GAOL}
\centerline{\textbf{GNU LIBRARY GENERAL PUBLIC LICENSE}}
\centerline{Version 2, June 1991}

Copyright \copyright\ 1991 Free Software Foundation, Inc.
59 Temple Place - Suite 330\\
Boston, MA 02111-1307, USA.

Everyone is permitted to copy and distribute verbatim copies
of this license document, but changing it is not allowed.

[This is the first released version of the library GPL\@.  It is
 numbered 2 because it goes with version 2 of the ordinary GPL.]

\subsection*{Preamble}

  The licenses for most software are designed to take away your
freedom to share and change it.  By contrast, the GNU General Public
Licenses are intended to guarantee your freedom to share and change
free software---to make sure the software is free for all its users.

  This license, the Library General Public License, applies to some
specially designated Free Software Foundation software, and to any
other libraries whose authors decide to use it.  You can use it for
your libraries, too.

  When we speak of free software, we are referring to freedom, not
price.  Our General Public Licenses are designed to make sure that you
have the freedom to distribute copies of free software (and charge for
this service if you wish), that you receive source code or can get it
if you want it, that you can change the software or use pieces of it
in new free programs; and that you know you can do these things.

  To protect your rights, we need to make restrictions that forbid
anyone to deny you these rights or to ask you to surrender the rights.
These restrictions translate to certain responsibilities for you if
you distribute copies of the library, or if you modify it.

  For example, if you distribute copies of the library, whether gratis
or for a fee, you must give the recipients all the rights that we gave
you.  You must make sure that they, too, receive or can get the source
code.  If you link a program with the library, you must provide
complete object files to the recipients so that they can relink them
with the library, after making changes to the library and recompiling
it.  And you must show them these terms so they know their rights.

  Our method of protecting your rights has two steps: (1) copyright
the library, and (2) offer you this license which gives you legal
permission to copy, distribute and/or modify the library.

  Also, for each distributor's protection, we want to make certain
that everyone understands that there is no warranty for this free
library.  If the library is modified by someone else and passed on, we
want its recipients to know that what they have is not the original
version, so that any problems introduced by others will not reflect on
the original authors' reputations.

  Finally, any free program is threatened constantly by software
patents.  We wish to avoid the danger that companies distributing free
software will individually obtain patent licenses, thus in effect
transforming the program into proprietary software.  To prevent this,
we have made it clear that any patent must be licensed for everyone's
free use or not licensed at all.

  Most GNU software, including some libraries, is covered by the ordinary
GNU General Public License, which was designed for utility programs.  This
license, the GNU Library General Public License, applies to certain
designated libraries.  This license is quite different from the ordinary
one; be sure to read it in full, and don't assume that anything in it is
the same as in the ordinary license.

  The reason we have a separate public license for some libraries is that
they blur the distinction we usually make between modifying or adding to a
program and simply using it.  Linking a program with a library, without
changing the library, is in some sense simply using the library, and is
analogous to running a utility program or application program.  However, in
a textual and legal sense, the linked executable is a combined work, a
derivative of the original library, and the ordinary General Public License
treats it as such.

  Because of this blurred distinction, using the ordinary General
Public License for libraries did not effectively promote software
sharing, because most developers did not use the libraries.  We
concluded that weaker conditions might promote sharing better.

  However, unrestricted linking of non-free programs would deprive the
users of those programs of all benefit from the free status of the
libraries themselves.  This Library General Public License is intended to
permit developers of non-free programs to use free libraries, while
preserving your freedom as a user of such programs to change the free
libraries that are incorporated in them.  (We have not seen how to achieve
this as regards changes in header files, but we have achieved it as regards
changes in the actual functions of the Library.)  The hope is that this
will lead to faster development of free libraries.

  The precise terms and conditions for copying, distribution and
modification follow.  Pay close attention to the difference between a
``work based on the library'' and a ``work that uses the library''.  The
former contains code derived from the library, while the latter only
works together with the library.

  Note that it is possible for a library to be covered by the ordinary
General Public License rather than by this special one.

\subsection*{Terms and Conditions for Copying, Distribution and Modification}

\begin{enumerate}
\setcounter{enumi}{-1}
\item This License Agreement applies to any software library which
contains a notice placed by the copyright holder or other authorized
party saying it may be distributed under the terms of this Library
General Public License (also called ``this License'').  Each licensee is
addressed as ``you''.

  A ``library'' means a collection of software functions and/or data
prepared so as to be conveniently linked with application programs
(which use some of those functions and data) to form executables.

  The ``Library'', below, refers to any such software library or work
which has been distributed under these terms.  A ``work based on the
Library'' means either the Library or any derivative work under
copyright law: that is to say, a work containing the Library or a
portion of it, either verbatim or with modifications and/or translated
straightforwardly into another language.  (Hereinafter, translation is
included without limitation in the term ``modification''.)

``Source code'' for a work means the preferred form of the work for
making modifications to it.  For a library, complete source code means
all the source code for all modules it contains, plus any associated
interface definition files, plus the scripts used to control compilation
and installation of the library.

  Activities other than copying, distribution and modification are not
covered by this License; they are outside its scope.  The act of
running a program using the Library is not restricted, and output from
such a program is covered only if its contents constitute a work based
on the Library (independent of the use of the Library in a tool for
writing it).  Whether that is true depends on what the Library does
and what the program that uses the Library does.
  
\item You may copy and distribute verbatim copies of the Library's
complete source code as you receive it, in any medium, provided that
you conspicuously and appropriately publish on each copy an
appropriate copyright notice and disclaimer of warranty; keep intact
all the notices that refer to this License and to the absence of any
warranty; and distribute a copy of this License along with the
Library.

  You may charge a fee for the physical act of transferring a copy,
and you may at your option offer warranty protection in exchange for a
fee.

\item You may modify your copy or copies of the Library or any portion
of it, thus forming a work based on the Library, and copy and
distribute such modifications or work under the terms of Section 1
above, provided that you also meet all of these conditions:

\begin{enumerate}
\item[a.]
The modified work must itself be a software library.

\item[b.]
You must cause the files modified to carry prominent notices
stating that you changed the files and the date of any change.

\item[c.]
You must cause the whole of the work to be licensed at no
charge to all third parties under the terms of this License.

\item[d.]
  If a facility in the modified Library refers to a function or a
  table of data to be supplied by an application program that uses the
  facility, other than as an argument passed when the facility is
  invoked, then you must make a good faith effort to ensure that, in
  the event an application does not supply such function or table, the
  facility still operates, and performs whatever part of its purpose
  remains meaningful.

    (For example, a function in a library to compute square roots has
    a purpose that is entirely well-defined independent of the
    application.  Therefore, Subsection 2d requires that any
    application-supplied function or table used by this function must
    be optional: if the application does not supply it, the square
    root function must still compute square roots.)
\end{enumerate}

These requirements apply to the modified work as a whole.  If
identifiable sections of that work are not derived from the Library,
and can be reasonably considered independent and separate works in
themselves, then this License, and its terms, do not apply to those
sections when you distribute them as separate works.  But when you
distribute the same sections as part of a whole which is a work based
on the Library, the distribution of the whole must be on the terms of
this License, whose permissions for other licensees extend to the
entire whole, and thus to each and every part regardless of who wrote
it.

Thus, it is not the intent of this section to claim rights or contest
your rights to work written entirely by you; rather, the intent is to
exercise the right to control the distribution of derivative or
collective works based on the Library.

In addition, mere aggregation of another work not based on the Library
with the Library (or with a work based on the Library) on a volume of
a storage or distribution medium does not bring the other work under
the scope of this License.

\item You may opt to apply the terms of the ordinary GNU General Public
License instead of this License to a given copy of the Library.  To do
this, you must alter all the notices that refer to this License, so
that they refer to the ordinary GNU General Public License, version 2,
instead of to this License.  (If a newer version than version 2 of the
ordinary GNU General Public License has appeared, then you can specify
that version instead if you wish.)  Do not make any other change in
these notices.

  Once this change is made in a given copy, it is irreversible for
that copy, so the ordinary GNU General Public License applies to all
subsequent copies and derivative works made from that copy.

  This option is useful when you wish to copy part of the code of
the Library into a program that is not a library.

\item You may copy and distribute the Library (or a portion or
derivative of it, under Section 2) in object code or executable form
under the terms of Sections 1 and 2 above provided that you accompany
it with the complete corresponding machine-readable source code, which
must be distributed under the terms of Sections 1 and 2 above on a
medium customarily used for software interchange.

  If distribution of object code is made by offering access to copy
from a designated place, then offering equivalent access to copy the
source code from the same place satisfies the requirement to
distribute the source code, even though third parties are not
compelled to copy the source along with the object code.

\item A program that contains no derivative of any portion of the
Library, but is designed to work with the Library by being compiled or
linked with it, is called a ``work that uses the Library''.  Such a
work, in isolation, is not a derivative work of the Library, and
therefore falls outside the scope of this License.

  However, linking a ``work that uses the Library'' with the Library
creates an executable that is a derivative of the Library (because it
contains portions of the Library), rather than a ``work that uses the
library''.  The executable is therefore covered by this License.
Section 6 states terms for distribution of such executables.

  When a ``work that uses the Library'' uses material from a header file
that is part of the Library, the object code for the work may be a
derivative work of the Library even though the source code is not.
Whether this is true is especially significant if the work can be
linked without the Library, or if the work is itself a library.  The
threshold for this to be true is not precisely defined by law.

  If such an object file uses only numerical parameters, data
structure layouts and accessors, and small macros and small inline
functions (ten lines or less in length), then the use of the object
file is unrestricted, regardless of whether it is legally a derivative
work.  (Executables containing this object code plus portions of the
Library will still fall under Section 6.)

  Otherwise, if the work is a derivative of the Library, you may
distribute the object code for the work under the terms of Section 6.
Any executables containing that work also fall under Section 6,
whether or not they are linked directly with the Library itself.

\item As an exception to the Sections above, you may also compile or
link a ``work that uses the Library'' with the Library to produce a
work containing portions of the Library, and distribute that work
under terms of your choice, provided that the terms permit
modification of the work for the customer's own use and reverse
engineering for debugging such modifications.

  You must give prominent notice with each copy of the work that the
Library is used in it and that the Library and its use are covered by
this License.  You must supply a copy of this License.  If the work
during execution displays copyright notices, you must include the
copyright notice for the Library among them, as well as a reference
directing the user to the copy of this License.  Also, you must do one
of these things:

\begin{enumerate}
\item[a.]  Accompany the work with the complete corresponding
  machine-readable source code for the Library including whatever
  changes were used in the work (which must be distributed under
  Sections 1 and 2 above); and, if the work is an executable linked
  with the Library, with the complete machine-readable ``work that uses
  the Library'', as object code and/or source code, so that the user
  can modify the Library and then relink to produce a modified
  executable containing the modified Library.  (It is understood that
  the user who changes the contents of definitions files in the
  Library will not necessarily be able to recompile the application to
  use the modified definitions.)

\item[b.] Accompany the work with a written offer, valid for at
    least three years, to give the same user the materials
    specified in Subsection 6a, above, for a charge no more
    than the cost of performing this distribution.

\item[c.] If distribution of the work is made by offering access to copy
    from a designated place, offer equivalent access to copy the above
    specified materials from the same place.

\item[d.] Verify that the user has already received a copy of these
    materials or that you have already sent this user a copy.
\end{enumerate}

  For an executable, the required form of the ``work that uses the
Library'' must include any data and utility programs needed for
reproducing the executable from it.  However, as a special exception,
the source code distributed need not include anything that is normally
distributed (in either source or binary form) with the major
components (compiler, kernel, and so on) of the operating system on
which the executable runs, unless that component itself accompanies
the executable.

  It may happen that this requirement contradicts the license
restrictions of other proprietary libraries that do not normally
accompany the operating system.  Such a contradiction means you cannot
use both them and the Library together in an executable that you
distribute.

\item You may place library facilities that are a work based on the
Library side-by-side in a single library together with other library
facilities not covered by this License, and distribute such a combined
library, provided that the separate distribution of the work based on
the Library and of the other library facilities is otherwise
permitted, and provided that you do these two things:

\begin{enumerate}
\item[a.] Accompany the combined library with a copy of the same work
    based on the Library, uncombined with any other library
    facilities.  This must be distributed under the terms of the
    Sections above.

\item[b.] Give prominent notice with the combined library of the fact
    that part of it is a work based on the Library, and explaining
    where to find the accompanying uncombined form of the same work.
\end{enumerate}

\item You may not copy, modify, sublicense, link with, or distribute
the Library except as expressly provided under this License.  Any
attempt otherwise to copy, modify, sublicense, link with, or
distribute the Library is void, and will automatically terminate your
rights under this License.  However, parties who have received copies,
or rights, from you under this License will not have their licenses
terminated so long as such parties remain in full compliance.

\item You are not required to accept this License, since you have not
signed it.  However, nothing else grants you permission to modify or
distribute the Library or its derivative works.  These actions are
prohibited by law if you do not accept this License.  Therefore, by
modifying or distributing the Library (or any work based on the
Library), you indicate your acceptance of this License to do so, and
all its terms and conditions for copying, distributing or modifying
the Library or works based on it.

\item Each time you redistribute the Library (or any work based on the
Library), the recipient automatically receives a license from the
original licensor to copy, distribute, link with or modify the Library
subject to these terms and conditions.  You may not impose any further
restrictions on the recipients' exercise of the rights granted herein.
You are not responsible for enforcing compliance by third parties to
this License.

\item If, as a consequence of a court judgment or allegation of patent
infringement or for any other reason (not limited to patent issues),
conditions are imposed on you (whether by court order, agreement or
otherwise) that contradict the conditions of this License, they do not
excuse you from the conditions of this License.  If you cannot
distribute so as to satisfy simultaneously your obligations under this
License and any other pertinent obligations, then as a consequence you
may not distribute the Library at all.  For example, if a patent
license would not permit royalty-free redistribution of the Library by
all those who receive copies directly or indirectly through you, then
the only way you could satisfy both it and this License would be to
refrain entirely from distribution of the Library.

If any portion of this section is held invalid or unenforceable under any
particular circumstance, the balance of the section is intended to apply,
and the section as a whole is intended to apply in other circumstances.

It is not the purpose of this section to induce you to infringe any
patents or other property right claims or to contest validity of any
such claims; this section has the sole purpose of protecting the
integrity of the free software distribution system which is
implemented by public license practices.  Many people have made
generous contributions to the wide range of software distributed
through that system in reliance on consistent application of that
system; it is up to the author/donor to decide if he or she is willing
to distribute software through any other system and a licensee cannot
impose that choice.

This section is intended to make thoroughly clear what is believed to
be a consequence of the rest of this License.

\item If the distribution and/or use of the Library is restricted in
certain countries either by patents or by copyrighted interfaces, the
original copyright holder who places the Library under this License may add
an explicit geographical distribution limitation excluding those countries,
so that distribution is permitted only in or among countries not thus
excluded.  In such case, this License incorporates the limitation as if
written in the body of this License.

\item The Free Software Foundation may publish revised and/or new
versions of the Library General Public License from time to time.
Such new versions will be similar in spirit to the present version,
but may differ in detail to address new problems or concerns.

Each version is given a distinguishing version number.  If the Library
specifies a version number of this License which applies to it and
``any later version'', you have the option of following the terms and
conditions either of that version or of any later version published by
the Free Software Foundation.  If the Library does not specify a
license version number, you may choose any version ever published by
the Free Software Foundation.

\item If you wish to incorporate parts of the Library into other free
programs whose distribution conditions are incompatible with these,
write to the author to ask for permission.  For software which is
copyrighted by the Free Software Foundation, write to the Free
Software Foundation; we sometimes make exceptions for this.  Our
decision will be guided by the two goals of preserving the free status
of all derivatives of our free software and of promoting the sharing
and reuse of software generally.
\end{enumerate}

\subsection*{NO WARRANTY}

\begin{enumerate}
\setcounter{enumi}{14}
\item BECAUSE THE LIBRARY IS LICENSED FREE OF CHARGE, THERE IS NO
WARRANTY FOR THE LIBRARY, TO THE EXTENT PERMITTED BY APPLICABLE LAW\@.
EXCEPT WHEN OTHERWISE STATED IN WRITING THE COPYRIGHT HOLDERS AND/OR
OTHER PARTIES PROVIDE THE LIBRARY ``AS IS'' WITHOUT WARRANTY OF ANY
KIND, EITHER EXPRESSED OR IMPLIED, INCLUDING, BUT NOT LIMITED TO, THE
IMPLIED WARRANTIES OF MERCHANTABILITY AND FITNESS FOR A PARTICULAR
PURPOSE\@.  THE ENTIRE RISK AS TO THE QUALITY AND PERFORMANCE OF THE
LIBRARY IS WITH YOU\@.  SHOULD THE LIBRARY PROVE DEFECTIVE, YOU ASSUME
THE COST OF ALL NECESSARY SERVICING, REPAIR OR CORRECTION.

\item IN NO EVENT UNLESS REQUIRED BY APPLICABLE LAW OR AGREED TO IN
WRITING WILL ANY COPYRIGHT HOLDER, OR ANY OTHER PARTY WHO MAY MODIFY
AND/OR REDISTRIBUTE THE LIBRARY AS PERMITTED ABOVE, BE LIABLE TO YOU
FOR DAMAGES, INCLUDING ANY GENERAL, SPECIAL, INCIDENTAL OR
CONSEQUENTIAL DAMAGES ARISING OUT OF THE USE OR INABILITY TO USE THE
LIBRARY (INCLUDING BUT NOT LIMITED TO LOSS OF DATA OR DATA BEING
RENDERED INACCURATE OR LOSSES SUSTAINED BY YOU OR THIRD PARTIES OR A
FAILURE OF THE LIBRARY TO OPERATE WITH ANY OTHER SOFTWARE), EVEN IF
SUCH HOLDER OR OTHER PARTY HAS BEEN ADVISED OF THE POSSIBILITY OF SUCH
DAMAGES.
\end{enumerate}

\begin{center}\textbf{END OF TERMS AND CONDITIONS}\end{center}

\subsection*{How to Apply These Terms to Your New Libraries}

  If you develop a new library, and you want it to be of the greatest
possible use to the public, we recommend making it free software that
everyone can redistribute and change.  You can do so by permitting
redistribution under these terms (or, alternatively, under the terms of the
ordinary General Public License).

  To apply these terms, attach the following notices to the library.  It is
safest to attach them to the start of each source file to most effectively
convey the exclusion of warranty; and each file should have at least the
``copyright'' line and a pointer to where the full notice is found.

\begin{quote}
{\em one line to give the library's name and a brief idea of what it does.}\\
Copyright (C) {\em year} {\em name of author}\\

    This library is free software; you can redistribute it and/or
    modify it under the terms of the GNU Lesser General Public
    License as published by the Free Software Foundation; either
    version 2 of the License, or (at your option) any later version.

    This library is distributed in the hope that it will be useful,
    but WITHOUT ANY WARRANTY; without even the implied warranty of
    MERCHANTABILITY or FITNESS FOR A PARTICULAR PURPOSE\@.  See the GNU
    Lesser General Public License for more details.

    You should have received a copy of the GNU Lesser General Public
    License along with this library; if not, write to the Free Software
    Foundation, Inc., 59 Temple Place, Suite 330, Boston, MA  02111-1307  USA
\end{quote}

Also add information on how to contact you by electronic and paper mail.

You should also get your employer (if you work as a programmer) or your
school, if any, to sign a ``copyright disclaimer'' for the library, if
necessary.  Here is a sample; alter the names:

\begin{quote}
  Yoyodyne, Inc., hereby disclaims all copyright interest in the
  library `Frob' (a library for tweaking knobs) written by James Random Hacker.

  {\em signature of Ty Coon}, 1 April 1990\\
  Ty Coon, President of Vice
\end{quote}

That's all there is to it!


\section{The licence of CORE-MATH}

The sources of CORE-MATH, in \file{3rd/math-core}, which GAOL compiles into
its library, carry the copyright notices of their authors, as each file states
(Inria, Paul Zimmermann, Alexei Sibidanov, St\'ephane Glondu and others), and
are distributed under the following licence (MIT licence), whose copyright
notice is theirs:

\begin{verbatim}
Permission is hereby granted, free of charge, to any person
obtaining a copy of this software and associated documentation
files (the "Software"), to deal in the Software without
restriction, including without limitation the rights to use,
copy, modify, merge, publish, distribute, sublicense, and/or
sell copies of the Software, and to permit persons to whom the
Software is furnished to do so, subject to the following
conditions:

The above copyright notice and this permission notice shall be
included in all copies or substantial portions of the Software.

THE SOFTWARE IS PROVIDED "AS IS", WITHOUT WARRANTY OF ANY KIND,
EXPRESS OR IMPLIED, INCLUDING BUT NOT LIMITED TO THE WARRANTIES
OF MERCHANTABILITY, FITNESS FOR A PARTICULAR PURPOSE AND
NONINFRINGEMENT. IN NO EVENT SHALL THE AUTHORS OR COPYRIGHT
HOLDERS BE LIABLE FOR ANY CLAIM, DAMAGES OR OTHER LIABILITY,
WHETHER IN AN ACTION OF CONTRACT, TORT OR OTHERWISE, ARISING
FROM, OUT OF OR IN CONNECTION WITH THE SOFTWARE OR THE USE OR
OTHER DEALINGS IN THE SOFTWARE.
\end{verbatim}


%%==========================================================================
%% Index
%%==========================================================================
\cleardoublepage
\pagestyle{INDEX}
\addcontentsline{toc}{chapter}{Index}
\printindex[idx]

\end{document}

%%% Local Variables:
%%% mode: latex
%%% TeX-master: t
%%% End:
